% Options for packages loaded elsewhere
\PassOptionsToPackage{unicode}{hyperref}
\PassOptionsToPackage{hyphens}{url}
%
\documentclass[
  10pt,
]{article}
\usepackage{amsmath,amssymb}
\usepackage{iftex}
\ifPDFTeX
  \usepackage[T1]{fontenc}
  \usepackage[utf8]{inputenc}
  \usepackage{textcomp} % provide euro and other symbols
\else % if luatex or xetex
  \usepackage{unicode-math} % this also loads fontspec
  \defaultfontfeatures{Scale=MatchLowercase}
  \defaultfontfeatures[\rmfamily]{Ligatures=TeX,Scale=1}
\fi
\usepackage{lmodern}
\ifPDFTeX\else
  % xetex/luatex font selection
  \setmainfont[]{DejaVu Serif}
  \setmonofont[]{DejaVu Sans}
\fi
% Use upquote if available, for straight quotes in verbatim environments
\IfFileExists{upquote.sty}{\usepackage{upquote}}{}
\IfFileExists{microtype.sty}{% use microtype if available
  \usepackage[]{microtype}
  \UseMicrotypeSet[protrusion]{basicmath} % disable protrusion for tt fonts
}{}
\makeatletter
\@ifundefined{KOMAClassName}{% if non-KOMA class
  \IfFileExists{parskip.sty}{%
    \usepackage{parskip}
  }{% else
    \setlength{\parindent}{0pt}
    \setlength{\parskip}{6pt plus 2pt minus 1pt}}
}{% if KOMA class
  \KOMAoptions{parskip=half}}
\makeatother
\usepackage{xcolor}
\usepackage[margin=1in]{geometry}
\usepackage{color}
\usepackage{fancyvrb}
\newcommand{\VerbBar}{|}
\newcommand{\VERB}{\Verb[commandchars=\\\{\}]}
\DefineVerbatimEnvironment{Highlighting}{Verbatim}{commandchars=\\\{\}}
% Add ',fontsize=\small' for more characters per line
\newenvironment{Shaded}{}{}
\newcommand{\AlertTok}[1]{\textcolor[rgb]{1.00,0.00,0.00}{\textbf{#1}}}
\newcommand{\AnnotationTok}[1]{\textcolor[rgb]{0.38,0.63,0.69}{\textbf{\textit{#1}}}}
\newcommand{\AttributeTok}[1]{\textcolor[rgb]{0.49,0.56,0.16}{#1}}
\newcommand{\BaseNTok}[1]{\textcolor[rgb]{0.25,0.63,0.44}{#1}}
\newcommand{\BuiltInTok}[1]{\textcolor[rgb]{0.00,0.50,0.00}{#1}}
\newcommand{\CharTok}[1]{\textcolor[rgb]{0.25,0.44,0.63}{#1}}
\newcommand{\CommentTok}[1]{\textcolor[rgb]{0.38,0.63,0.69}{\textit{#1}}}
\newcommand{\CommentVarTok}[1]{\textcolor[rgb]{0.38,0.63,0.69}{\textbf{\textit{#1}}}}
\newcommand{\ConstantTok}[1]{\textcolor[rgb]{0.53,0.00,0.00}{#1}}
\newcommand{\ControlFlowTok}[1]{\textcolor[rgb]{0.00,0.44,0.13}{\textbf{#1}}}
\newcommand{\DataTypeTok}[1]{\textcolor[rgb]{0.56,0.13,0.00}{#1}}
\newcommand{\DecValTok}[1]{\textcolor[rgb]{0.25,0.63,0.44}{#1}}
\newcommand{\DocumentationTok}[1]{\textcolor[rgb]{0.73,0.13,0.13}{\textit{#1}}}
\newcommand{\ErrorTok}[1]{\textcolor[rgb]{1.00,0.00,0.00}{\textbf{#1}}}
\newcommand{\ExtensionTok}[1]{#1}
\newcommand{\FloatTok}[1]{\textcolor[rgb]{0.25,0.63,0.44}{#1}}
\newcommand{\FunctionTok}[1]{\textcolor[rgb]{0.02,0.16,0.49}{#1}}
\newcommand{\ImportTok}[1]{\textcolor[rgb]{0.00,0.50,0.00}{\textbf{#1}}}
\newcommand{\InformationTok}[1]{\textcolor[rgb]{0.38,0.63,0.69}{\textbf{\textit{#1}}}}
\newcommand{\KeywordTok}[1]{\textcolor[rgb]{0.00,0.44,0.13}{\textbf{#1}}}
\newcommand{\NormalTok}[1]{#1}
\newcommand{\OperatorTok}[1]{\textcolor[rgb]{0.40,0.40,0.40}{#1}}
\newcommand{\OtherTok}[1]{\textcolor[rgb]{0.00,0.44,0.13}{#1}}
\newcommand{\PreprocessorTok}[1]{\textcolor[rgb]{0.74,0.48,0.00}{#1}}
\newcommand{\RegionMarkerTok}[1]{#1}
\newcommand{\SpecialCharTok}[1]{\textcolor[rgb]{0.25,0.44,0.63}{#1}}
\newcommand{\SpecialStringTok}[1]{\textcolor[rgb]{0.73,0.40,0.53}{#1}}
\newcommand{\StringTok}[1]{\textcolor[rgb]{0.25,0.44,0.63}{#1}}
\newcommand{\VariableTok}[1]{\textcolor[rgb]{0.10,0.09,0.49}{#1}}
\newcommand{\VerbatimStringTok}[1]{\textcolor[rgb]{0.25,0.44,0.63}{#1}}
\newcommand{\WarningTok}[1]{\textcolor[rgb]{0.38,0.63,0.69}{\textbf{\textit{#1}}}}
\setlength{\emergencystretch}{3em} % prevent overfull lines
\providecommand{\tightlist}{%
  \setlength{\itemsep}{0pt}\setlength{\parskip}{0pt}}
\setcounter{secnumdepth}{-\maxdimen} % remove section numbering

\usepackage{fvextra}
\DefineVerbatimEnvironment{Highlighting}{Verbatim}{breaklines,fontsize=\scriptsize,commandchars=\\\{\}}
\emergencystretch=3em
\ifLuaTeX
  \usepackage{selnolig}  % disable illegal ligatures
\fi
\usepackage{bookmark}
\IfFileExists{xurl.sty}{\usepackage{xurl}}{} % add URL line breaks if available
\urlstyle{same}
\hypersetup{
  pdftitle={Intrinsic zero-Hessian recognition on arbitrary-field analytic curves},
  pdfauthor={Unverified research draft; author self-audit only},
  hidelinks,
  pdfcreator={LaTeX via pandoc}}

\title{Intrinsic zero-Hessian recognition on arbitrary-field analytic
curves}
\author{Unverified research draft; author self-audit only}
\date{3 October 2026}

\begin{document}
\maketitle

\section{Result, conventions, and proof
map}\label{result-conventions-and-proof-map}

This submission claims a curve-scope theorem and a ground-field probe
consequence. It is partial progress on P06.2, not a solution of the
original arbitrary-dimensional atom. No prior draft, native note, review
disposition, production run, or hash is a mathematical premise. The
arguments needed from the supplied drafts are reconstructed below, and
the primary source statements and proof bodies are reproduced in this
document. Standard analytic foundations are explicitly stated and cited.
The claims remain unverified mathematical drafts; typesetting is not
verification.

Write \(D=d'd''\). All test forms and all current equations are
\textbf{ordinary CLD} forms and currents on ordinary topological opens.
For a pure-dimensional analytic space \(Z\), \([u]=u\delta_Z\) uses
canonical integration with the actual generic Artinian lengths,
including nilpotents. On a curve the equation means \[
             \int_Z u(D\psi)_{\mathrm{Rad}}=0
             \quad\text{for every }\psi\in A_c^{0,0}(Z).
\] Only the smooth top form has been identified with a signed Radon
measure; \(D[u]\) is not assumed to be a measure. All functions under
recognition are finite, continuous, and real valued. Real PL means the
real extension of the integral G-local PL class, not ordinary CLD
smoothness or rational slopes.

\textbf{Theorem R (intrinsic arbitrary-field curve recognition).} Let
\(K\) be any complete nontrivially valued field. Let \(X\) be any good,
separated, strict, boundaryless \(K\)-analytic space pure of dimension
one, and let \(U\subseteq X\) be any ordinary open. If
\(h:U\to\mathbb R\) is continuous and \(D[h]=0\), then \(h\) is harmonic
in GR Definition 3.13 and classically pluriharmonic in GR Definition
7.1. In particular, for \textbf{every} complete analytic extension
\(L/K\), every \(L\)-analytic curve \(Y\) (separated and pure of
dimension one), and every analytic map \(f:Y\to U_L\), the two functions
\(\pm h\circ\pi_{L/K}\circ f\) are GR-subharmonic on \(Y\). On every
smooth strict connected such probe over a nontrivially valued field,
they are Thuillier-subharmonic and their common finite function is
harmonic. Moreover, \(h\) is G-locally real PL and is GJR harmonic:
every residual class
\(L_h(x)\in\operatorname{Pic}(\operatorname{red}(U,x))_{\mathbb R}\) is
numerically trivial. Singularities, nonreduced structures, imperfect
fields, inseparable residue fields, and nonproper nonalgebraic curves
are permitted. No scalar stability assumption is imposed.

\textbf{Theorem P (ground-field probes).} Let \(K\) have the same field
hypotheses, and let \(X\) be good, separated, strict, boundaryless, and
pure of any dimension \(n\geq1\). For every ordinary open
\(U\subseteq X\), continuous finite \(h:U\to\mathbb R\) with \(D[h]=0\),
and analytic \(K\)-map \(f:Y\to U\) from a separated smooth strict
\(K\)-analytic curve, \(h\circ f\) is Thuillier harmonic, GR-harmonic,
classically pluriharmonic on \(Y\), G-locally real PL, and GJR harmonic.
Here smooth includes absence of absolute analytic boundary. This
statement concerns maps over the ground field; it does not identify
\(h\) as classically pluriharmonic on a higher-dimensional \(U\).

The proof of R has three separate operations. Part T proves current
transport to the completion of an algebraic extension on actual pure
curves. Part C recognizes on algebraically closed fields by actual local
algebraic models and the curve comparison, with component lengths and
punctures handled explicitly. Part B proves that the completion of an
algebraic closure is an admissible field and checks descent of the
recognized functions with full quantifiers. Part Q reconstructs
closed-curve restriction and the actual local-image theorem, then
deduces P. The scalar proof in T is dimension one; Q does not create a
current equation after a transcendental extension of a
higher-dimensional ambient space.

The controlling original question is Problem
\texttt{prob:frontier-pluriharmonic} in the supplied
\texttt{ch72\_major\_open\_problems.tex}: complete nontrivially valued
\(K\), good separated strict boundaryless \(X\), and recognition of
every continuous zero ordinary Hessian current, in arbitrary dimension.
The accompanying \texttt{appM\_ch18\_frontiers.tex} distinguishes
recognition, scalar-current ascent, and two-sign positivity
regularization. These texts specify scope, not proof inputs. R removes
the ground-field restriction from the supplied algebraically closed
curve recognition. P adds ground-field smooth strict probes over
arbitrary \(K\). The full atom and P05.2 remain open in this submission;
\texttt{candidate=false}.

\section{Part T. Completed-algebraic transport:
statement}\label{part-t.-completed-algebraic-transport-statement}

\textbf{Transport lemma.} For a good separated strict boundaryless pure
\(K\)-analytic curve \(Z\), an algebraic extension \(M/K\) with its
extending absolute value, and \(L=\widehat M\), the actual scalar
projection satisfies \[
                  D[u]=0\quad\Longleftrightarrow\quad D[u\circ\pi]=0.
\] It holds on every ordinary open, with the original lengths. A
subsidiary lemma, not requiring dimension one, is that every compact
ordinary smooth scalar test on \(Z_L\) is exactly the pullback of a
compact smooth scalar test on \(Z_E\) for some finite embedded
\(K\subseteq E\subseteq M\). The following T2--T13 prove these
assertions. Equation numbers inside T are local to this proof. All
source excerpts referenced there are supplied in Appendix S; the
historical misselected Section 15 blocks are excluded.

For a finite Galois stage the approximation assertion proved below is
that an invariant compact ordinary scalar test \(\beta\) has downstairs
compact smooth approximants \(\alpha_j\) supported in one compact \(H\),
with all errors supported in \(\pi^{-1}H\), such that \[
 \big\|(D(\beta-\pi^*\alpha_j))_{\mathrm{Rad}}\big\|_{\mathrm{TV}}\to0.\tag{2}
\] Set \[
 B_u(\psi)=\int_Zu(D\psi)_{\mathrm{Rad}}.\tag{3}
\] Equation (1) denotes the equivalence in the transport lemma, while
(2) and (3) are local equation labels for T. The proof handles finite
separable, finite inseparable, mixed embedded finite stages, then the
completion.

\subsection{T2. Exact analytic dependencies and source
checks}\label{t2.-exact-analytic-dependencies-and-source-checks}

CLD refers to Chambert-Loir--Ducros, \emph{Formes différentielles
réelles et courants sur les espaces de Berkovich}, arXiv:1204.6277v2, 25
July 2025. GR refers to Gubler--Rabinoff, \emph{Classically psh and
pluriharmonic functions on Berkovich spaces}, arXiv:2506.13548v2, 17
September 2025. Numbering is that of these cached primary snapshots,
whose statements and proof bodies, not just the former drafts'
descriptions, were inspected.

\begin{enumerate}
\def\labelenumi{\arabic{enumi}.}
\item
  \textbf{Affinoid foundations.} CLD 11.1.1--2: an affinoid spectrum is
  compact Hausdorff; its algebra is noetherian with closed ideals; the
  map to its prime spectrum is surjective. CLD 11.1.5: every point of an
  affinoid has rational affinoid neighbourhoods inside any prescribed
  ordinary open, with the point satisfying all associated strict
  inequalities. Rational systems have positive real radii and
  denominator/numerators with no common zero (T12 uses this stronger
  neighbourhood clause of 11.1.5). CLD 11.1.7--8: a good space has an
  affinoid neighbourhood basis; ordinary opens are analytic domains;
  analytic sections form a sheaf also for ordinary opens, and
  \(\mathcal O(\mathcal M(A))=A\). CLD 11.1.9: affinoid fibre products
  and scalar extensions use completed tensor products, and scalar
  projection is surjective. CLD 11.4.4: a finite algebra with its
  finite-module Banach topology is affinoid. CLD 11.7.7: over a
  nontrivially valued field, every point of a good strict space has a
  strict affinoid neighbourhood. Applying this to ordinary opens gives
  such a neighbourhood inside the open. Particular tensor, orbit,
  compactness, shrinking and unit claims are proved in T3.
\item
  \textbf{Dimension and boundary.} CLD 11.2.2: separated analytic spaces
  are topologically separated. CLD 11.2.3: strict affinoid dimension is
  Krull dimension, dimension is preserved by complete scalar extension,
  and domains preserve local analytic dimension. If the minimal-prime
  components of \(\mathcal M(A)\) are \(X_i\), then
  \(\dim_xX=\sup_{x\in X_i}\dim X_i\). Pure dimension means constant
  local dimension. GR 2.11(5), equivalently (5'), states for a morphism
  \(Y\to Z\) and any analytic extension \(L/K\) that
  \(\partial(Y_L/Z_L)\subseteq\pi^{-1}(\partial(Y/Z))\). Hence a
  boundaryless space remains boundaryless. GR 2.11(3),(4) gives
  boundarylessness for an ordinary open in a boundaryless space: its
  relative interior is its ordinary topological interior and interiors
  compose. T3 proves purity component by component; equality of total
  dimensions alone is not the proof.
\item
  \textbf{Ordinary forms.} CLD Corollary 15.1.2 states that a function
  is smooth near \(x\) if it has a representation
  \(\varphi(t_{f_1},\ldots,t_{f_m})\) on a neighbourhood of \(x\), where
  the \(f_i\) are analytic units and \(\varphi\) is smooth on a
  Euclidean neighbourhood of the coordinate value. CLD Proposition
  15.1.1(b) gives ordinary open moment descriptions for forms in every
  bidegree. CLD 6.2.1--2 defines ordinary forms and says they form an
  ordinary sheaf stable under wedge, \(J,d',d''\). Moment pullback
  commutes with the exterior calculus (CLD 3.3.12). In a moment a
  logarithm is an affine coordinate, so \(Dt_f=0\) as a \textbf{form}.
  CLD 3.4.4: for a G-tropical chart \(f\) and a coefficient form \(a\),
  its pullback is zero in the G-form sheaf exactly when \(a\) restricts
  to zero on \(f(V)\) for every compact domain \(V\). This identifies
  the ordinary form sheaves when the analytic units, domains and moment
  images are the same; it is not descent through scalar projection.
\item
  \textbf{Compact cutoffs only.} CLD 15.2.2: a good topologically
  separated space is locally holomorphically separated. CLD Corollary
  15.2.5: on a locally holomorphically separated analytic space, for any
  compact \(F\) and ordinary open neighbourhood \(U\) of \(F\), there is
  an ordinary smooth function equal to one near \(F\) with compact
  support contained in \(U\). Its proof uses local differential
  separation (15.2.4) and the smooth bump result 6.2.4. T4 constructs
  the only partition statement needed, for compact \(F\) and finite
  covers. The overbroad fixed-index countable-at-infinity reading of CLD
  15.2.9 is not used.
\item
  \textbf{Canonical top measures.} CLD Theorem 15.7.1: a pure
  \(n\)-dimensional topologically separated analytic space with its
  canonical calibration and analytic boundary is a tropical space of
  level \(n\). CLD 15.7.2: the induced structure on a domain agrees with
  that domain's analytic structure. CLD 5.4.4: every compact tropical
  chart is pleasant, and its tropical polyhedron is the top-dimensional
  part of its compact moment image. CLD 14.3.2--5, 14.5.1--2: canonical
  calibrations are standard and positive, and on a top moment cell \(Q\)
  the calibration is \(d_Q\tau_{Q,*}|e_1\wedge\cdots\wedge e_n|\), where
  \(\tau_Q\) is the affine inverse of an injective torus projection and
  \(d_Q\) is the sum of the actual finite-flat degrees of the skeletal
  cells above \(Q\). Those degrees retain fibre-algebra dimensions,
  including lengths. CLD 5.8.3--5 constructs a top form's measure by
  contraction on its supporting skeleton. CLD Proposition 5.10.1(c,d):
  for a pleasant chart \(f\) and a Borel \((n,n)\) coefficient form
  \(a\), the induced measure is locally finite exactly when \(a\) is
  locally integrable on the calibrated tropical polyhedron, and then
  \(f_*((f^*a)_{\mathrm{Rad}})=\langle a,\mu_{\Pi_f}\rangle_{\mathrm{Rad}}\);
  the corresponding integrals agree. In particular a smooth top form is
  finite on a compact moment domain. Lower-dimensional cells have no top
  measure. Compact moment images are finite polyhedra by the
  polyhedrality assertion of CLD Theorem 12.2.4. The dimension bound is
  specifically CLD 12.2.10(1)--(2):
  \(\dim_{\mathrm{trop},x}f\le d_K(x)\le\dim_xX\) and, for compact
  \(X\), \(\dim f_{\mathrm{trop}}(X)=\sup_x\dim_{\mathrm{trop},x}f\).
  Thus compact curve images have dimension at most one. We do not use a
  claim that a moment image always has dimension equal to its source;
  constant moments rule out that reading. CLD 7.3.5 defines integration
  currents for locally integrable multipliers; a continuous finite
  multiplier is locally integrable. T8 proves the carrier
  total-variation estimate needed here, before any possibly cancelling
  signed pushforward.
\item
  \textbf{Scalar projection has no field-degree factor.} CLD Proposition
  14.3.6: on a level-\(n\) analytic space, if \(\Sigma\) has canonical
  calibration \(\mu\) and \(\Sigma_L\) its scalar extension has
  canonical calibration \(\mu_L\), then \(\pi_*\mu_L=\mu\). CLD
  Proposition 15.3.1: for a tropically Borel \((n,n)\) G-form
  \(\omega\), \(\pi_*\omega_{L,\mathrm{Bor}}=\omega_{\mathrm{Bor}}\);
  for a locally integrable such form, \[
             \pi_*\omega_{L,\mathrm{Rad}}
                       =\omega_{\mathrm{Rad}}.          \tag{4}
  \] Here is the primary proof mechanism, retaining its lengths. On a
  convenient skeletal cell choose a torus projection injective on the
  cell. Its finite-flat degree downstairs equals the sum of its degrees
  at all upstairs skeletal lifts. This is CLD Lemma 13.3.3(a): degrees
  at \(y\) equal the sum of degrees over \(y\) after complete base
  change and choice of a target lift. The lemma reduces to a fibre; a
  zero-dimensional local fibre algebra is Artinian, and its tensor
  decomposes into its local Artinian factors. The sum of their
  dimensions is the original dimension after scalar extension.
  Positive-dimensional fibres give infinite degrees on both sides. Thus
  the cell volume vectors sum to the downstairs vector, proving 14.3.6.
  Contracting the form and taking Radon measures proves 15.3.1. We apply
  (4) only to smooth top forms, then pair with \(h\); this never assumes
  pullback of a continuous current.
\item
  \textbf{Same-carrier field restriction has a degree factor.} CLD
  Remark 14.4.3: for a finite extension \(F/K\), an \(F\)-analytic space
  \(Z\) of level \(n\), and its restricted \(K\) structure \(Z_0\),
  dimensions and skeletal subsets agree and \(\mu_{Z_0}=[F:K]\mu_Z\).
  The proof uses the finite-flat degree \([F:K]\) of the projection of
  the split torus on this restricted carrier. CLD 15.4.1 and Proposition
  15.4.3: for a morphism \(f:Y\to Z\) of level-\(n\) spaces finite off a
  Zariski closed set of dimension \(<n\), a locally integrable \((p,n)\)
  G-form \(\omega\), \(p\in\{n-1,n\}\), has locally integrable pullback
  and
  \(f_*((f^*\omega)_{\mathrm{Rad}})=\mathrm{dc}(f)\omega_{\mathrm{Rad}}\).
  Its proof first proves the calibrated volume-vector identity on the
  inverse skeletal set (14.4.2), then contracts and integrates. We use
  only an actual isomorphism on \(Z_0\), whose calibrated degree is one:
  both local Artinian rings are identical. Cancelling the same-carrier
  factor \([F:K]\) proves invariance over \(F\).
\end{enumerate}

The stated calibration and form statements apply to the actual,
potentially singular and nonreduced structure. They require no
normalization and no quasi-smoothness. No weakly smooth GJR descent,
positive-current recognition, or catalogue theorem is used as an
ordinary form theorem.

\subsubsection{T2.1. Direct proofs of the neighbourhood and
compact-cutoff facts used
here}\label{t2.1.-direct-proofs-of-the-neighbourhood-and-compact-cutoff-facts-used-here}

The finite proof uses the primary foundations above. For the two
topological localization facts decisive at completion, the following
direct constructions supply their proof bodies as well.

On \(\mathcal M(A)\) the ordinary topology is initial for all
evaluations \(|a|\) (CLD 11.1.1). A prescribed neighbourhood of a point
therefore contains a finite intersection of conditions
\(|a_j|>\ell_j>0\) and \(|b_k|<u_k\), where thresholds can be chosen
strictly between the point's value and the excluded endpoint. Missing
lower or upper bounds impose no condition. Put \(g=\prod_j a_j\), with
\(g=1\) when there are no lower bounds. Choose \(C>0\) so that
\(C^{-1}<\prod_j\ell_j\). The rational conjunction \[
 |1|\le C|g|,\qquad
 |\prod_{j'\ne j}a_{j'}|\le\ell_j^{-1}|g|\ (\forall j),\qquad
 |b_kg|\le u_k|g|\ (\forall k)                         \tag{R}
\] has no common zero among its denominator and numerators, because one
numerator is \(1\). It is a rational analytic domain by the universal
rational-domain construction of CLD 11.1.5. Its validity locus lies in
the originally prescribed neighbourhood if the thresholds were chosen
inside that neighbourhood's intervals: its first inequality excludes
\(g=0\), so division of norms gives all chosen lower and upper bounds.
The original point satisfies every strict inequality, including the
first. Non-strict thresholds in (R) were chosen inside the original open
intervals, so validity of the weak system suffices for containment. This
proves the strong rational-neighbourhood clause needed in T12 using the
initial topology and the rational-domain construction. Real positive
thresholds suffice even when the value group is discrete.

Here is also a direct compact-cutoff construction in every good
Hausdorff analytic space occurring in the proof. For a point \(x\) in an
ordinary open \(U\), take an affinoid neighbourhood \(V_0\) and within
its ordinary interior a rational neighbourhood \(R\) of \(x\) contained
in \(U\), with strict inequalities \(|f_i(x)|<r_i|g(x)|\). Its
denominator satisfies \(g(x)\ne0\). Choose \(b_i\) with
\(|f_i(x)|/|g(x)|<b_i<r_i\) and choose \(0<a<|g(x)|\). On the ordinary
nonzero locus of \(g\) in \(V_0\), a smooth function of \(\log|g|\)
equals one near \(x\) and vanishes where \(|g|\le a\). For each \(i\),
choose a smooth real function of \(\log|f_i/g|\) that is one when
\(|f_i/g|\le b_i\) and zero when \(|f_i/g|\ge r_i'\) for some
\(b_i<r_i'<r_i\). Extend that factor as the constant one at zeros of
\(f_i\): near such a zero its ratio is below \(b_i\), so the extension
is locally constant and ordinary smooth. On its nonzero locus the ratio
is an analytic unit, giving the required ordinary moment description.
Multiply the factors with the denominator cutoff. At points where
\(g=0\), extend by zero; the denominator cutoff already vanishes on a
neighbourhood of each such point.

The support of this product lies in the compact set in \(V_0\) defined
by \(|g|\ge a\) and \(|f_i|\le r_i'|g|\) for all \(i\). That compact set
lies in the strict locus of \(R\), hence in
\(U\cap\operatorname{int}_X(V_0)\). Extend the product by zero outside
\(V_0\): every point outside the compact support has a zero
neighbourhood, so the extension is ordinary smooth. It equals one near
\(x\). Taking squares gives nonnegative cutoffs. For compact
\(F\subset U\), finitely many positive loci of these point cutoffs cover
\(F\). Their sum \(s\) is positive near \(F\). Compose with a real
smooth function that vanishes near zero and equals one on a
neighbourhood of the compact positive set \(s(F)\). This gives a compact
cutoff equal to one near \(F\) and supported in \(U\). This reconstructs
the compact CLD 15.2.5 conclusion in our actual good Hausdorff scope. T4
then proves its finite-cover partition specialization. No
countable-at-infinity fixed-index compact-support assertion is used.

In particular all cutoffs are obtained from ordinary analytic-unit
moments, flat constant extensions near analytic zeros and compact
extension by zero. No weakly smooth tropicalization is substituted.

\subsection{T3. Finite fields, fibres, analytic structures, and
purity}\label{t3.-finite-fields-fibres-analytic-structures-and-purity}

We first treat a finite Galois extension \(E/K\), with group \(G\) and
\(d=[E:K]=|G|\). The elementary finite field facts used below can be
justified as follows. A finite separable extension has as many distinct
embeddings in an algebraic closure as its degree: extend embeddings
successively by the distinct roots of each separable minimal polynomial.
The field \(K\) is infinite. For two separable generators \(u,v\), all
but finitely many \(c\in K\) make the values of \(u+cv\) distinct for
every pair of distinct embeddings; if their \(v\) values agree their
\(u\) values already differ. Its minimal polynomial then has at least as
many roots as the extension degree, so \(u+cv\) generates the extension.
Induction gives a primitive element. The splitting field of its
separable minimal polynomial is finite separable. Every embedding of
this splitting field permutes the roots and is onto, so it is Galois and
its automorphism group has its degree as order. Its fixed field is
\(K\): the embeddings of any intermediate separable simple field extend
successively to the splitting field, so a fixed element has a minimal
polynomial of degree one. This also provides the normal closure used in
T10.

Finite-dimensional norms over a complete non-Archimedean field are
equivalent to the coefficient maximum norm in a fixed basis. Here is a
proof. Induct on dimension. The span of all but the last vector is
complete by induction, hence closed. The last vector has positive
distance from this span, bounding its coefficient in any combination.
The ultrametric inequality and induction then bound all other
coefficients. The reverse bound is the triangle inequality. Consequently
finite-dimensional normed spaces are complete and their coordinate maps
bounded. Existence of an absolute value on a finite field extension can
also be seen by giving the finite algebra its finite-module Banach
topology, using CLD 11.4.4, and taking a point of its nonempty affinoid
spectrum. Its kernel is zero since the algebra is a field. Uniqueness
follows from norm equivalence: compare two extending norms on \(z^m\)
and take \(m\)th roots. Hence all finite field automorphisms preserve
the absolute value, and every finite field here is complete.

For a \(K\)-affinoid algebra \(A\), tensor the bounded coordinate maps
of \(E\) with \(A\). The tensor norm upper bound and these coordinate
maps identify its topology with a finite sum of copies of \(A\). Already
complete, this proves \[
             B=A\widehat\otimes_K E=A\otimes_K E.        \tag{5}
\] This is the actual algebra, including nilpotents. Taking the kernel
of \(E\to\bigoplus_{g\in G}E\), \(z\mapsto(gz-z)_g\), and tensoring over
the field gives \(B^G=A\). Invariant units descend as units since their
inverses are invariant too. There is no field averaging or division by
\(d\). An affinoid element nonzero at every analytic point is a unit:
otherwise a maximal ideal containing it and spectral surjectivity give a
zero.

The fibre over \(x\) is \(\mathcal M(\mathcal H(x)\otimes_K E)\). A
primitive separable polynomial factors without repeated factors over
\(\mathcal H(x)\), so the tensor is a finite product of fields. Each has
one analytic point by the existence and uniqueness argument above. The
\(G\) action on these finitely many points is transitive: two orbits
would supply a nontrivial invariant idempotent, but the invariant
algebra is \(\mathcal H(x)\). Stabilizers can be nontrivial; no local
covering-space description is claimed or used.

Inverse images of affinoid ordinary interiors give goodness upstairs.
Strict presentations base change to strict presentations. Base change of
the closed diagonal gives separatedness. Boundarylessness follows from
T2(2). In particular all carriers are Hausdorff. The inverse image of
any compact \(F\subset X\) is compact: cover \(F\) by finitely many
interiors of affinoid neighbourhoods \(V_i\); its inverse image is
closed in the finite union of the compact \(V_{i,E}\).

\textbf{Full-fibre shrinking.} In a downstairs strict affinoid
neighbourhood \(V_0\) of \(x\), let \(N\subset V_{0,E}\) be a relative
ordinary open containing the whole fibre. The compact complement has
closed image excluding \(x\). Choose a smaller strict affinoid
neighbourhood \(V\) of \(x\) in the complement of that image and inside
\(\operatorname{int}_X(V_0)\). Then \(V_E\subset N\). Thus finitely many
unit or strict norm conditions holding on the whole fibre persist on one
neighbourhood. This includes conditions on elements of zero seminorm at
\(x\) and never replaces \(A\) by its reduction.

For completeness, purity is checked at minimal components. In a strict
affinoid \(\mathcal M(A)\) in \(X\), each minimal-prime component has
dimension one: choose a point with kernel that prime by spectral
surjectivity and apply the local dimension formula. Let \(J\) be the
nilradical. It is nilpotent since \(A\) is noetherian. The injection
\(A/J\hookrightarrow\prod_P\operatorname{Frac}(A/P)\) stays injective
after tensoring with a separable field. Each field tensor is a product
of fields, hence \((A/J)\otimes E\) is reduced and the nilradical of
\(B\) is exactly \(J\otimes E\). The finitely many kernels of its maps
to those field factors have zero intersection. Their product is zero, so
a minimal prime contains one of these kernels, and by minimality equals
it. Its contraction is the corresponding minimal \(P\) of \(A\). The
finite integral injection \(A/P\hookrightarrow B/Q\) preserves Krull
dimension. One elementary justification is lying over and going up from
the determinant argument for a finite module over a local ring:
\(\mathfrak pD_{\mathfrak p}\ne D_{\mathfrak p}\), since otherwise the
determinant of \(I-M\), \(M\) with entries in \(\mathfrak p\), is a unit
annihilating that nonzero module. Choose a maximal ideal above it and
then repeat after quotienting by a chosen prime. For incomparability
quotient by the smaller prime and invert the nonzero elements of the
base domain. The resulting finite integral domain over a field is a
field, since a least-degree monic equation for a nonzero element
expresses its inverse. A strictly larger prime of the same contraction
could not survive, nor could a nonzero element of a domain vanish under
that localization. Going up lifts chains and incomparability makes
contractions strict. The dimensions agree. Every minimal component
upstairs therefore has dimension one, proving purity. The same proof
works for any finite separable stage. Reduction here is a dimension
check only; all measures below use \(B\) itself.

Finally consider the same carrier \(Z=X_E\) with its restricted
\(K\)-analytic structure \(Z_0\). A finite \(E/K\) is \(K\)-affinoid; an
\(E\)-affinoid presentation is therefore \(K\)-affinoid under finite
field restriction. The finite-module topology is unchanged. Rational
domains are defined by the same elements and inequalities and have the
same analytic rings. They give a common local basis and the same unit
germs and compact moment images. The moment nullity criterion of Section
2(3) identifies the ordinary form sheaves, their exterior calculus and
supports. Each \(g\in G\) is a bounded \(K\)-algebra automorphism of
(5), gluing to a \(K\)-analytic automorphism of \(Z_0\). Its isomorphism
pushforward formula, followed by cancellation of the same-carrier factor
\(d\), gives for a smooth compact top form \(\omega\) \[
             g_*((g^*\omega)_{\mathrm{Rad},E})
                            =\omega_{\mathrm{Rad},E}.    \tag{6}
\] This is different from (4). Both are valid with actual generic
lengths.

\subsection{T4. Compact partitions and full-fibre scalar
descriptions}\label{t4.-compact-partitions-and-full-fibre-scalar-descriptions}

Let compact \(F\) have an ordinary open cover. At every point of \(F\),
apply T2(4)'s bump theorem to the singleton in a chosen cover member.
Squaring the bump makes it nonnegative, still equal to one near that
point and compactly supported there. Select finitely many whose positive
loci cover \(F\), and group them by cover index as \(\psi_i\). Then
\(s=\sum_i\psi_i>0\) near \(F\). Take a bump \(\chi=1\) near \(F\)
compactly supported in \(\{s>0\}\) and put \[
                \theta_i=\chi\psi_i/s                  \tag{7}
\] there, extending by zero. Products, smooth composition, and
reciprocal on \(s>0\) are ordinary smooth operations in moment
coordinates. The compact support of \(\chi\) lies inside \(\{s>0\}\), so
extension by zero is smooth also at its complement. Each support is
compact and contained in its cover member. The sum is one near \(F\).
This proves the needed finite compact partition without any noncompact
fixed-cover-index assertion.

Let \(\beta\) be an invariant ordinary smooth scalar function near a
whole fibre over \(x\). Its fibre has finitely many points. Choose
pairwise disjoint ordinary neighbourhoods \(N_i\) of them with local
representations \(\beta=\varphi_i(t_{f_{i,1}},\ldots)\), using Corollary
15.1.2. Shrink \(N_i\) so the coordinate values are in a relatively
compact Euclidean region on which \(\varphi_i\) is smooth. Full-fibre
shrinking gives a strict affinoid \(V\) with \(V_E\subset\bigcup_iN_i\).
The pieces \(C_i=V_E\cap N_i\) are disjoint ordinary opens and closed in
\(V_E\). They are compact. Analytic gluing gives idempotents \(e_i\)
which are one on \(C_i\) and zero on the other pieces, and gives a
global unit \(\widetilde f_{i,k}\) equal to \(f_{i,k}\) on \(C_i\) and
one elsewhere. Its inverse glues in the same way. These are elements of
the actual affinoid algebra \(B\), not formal topological selectors.

Choose \(a\in K^\times\) with \(|a|\ne1\). The analytic unit
\(s_i=e_i+a(1-e_i)\) has \(t_{s_i}=0\) on \(C_i\) and
\(t_{s_i}=-\log|a|\) elsewhere. Choose a real smooth \(q\) equal to one
near zero and zero near \(-\log|a|\). Extend \(\varphi_i\) by a
Euclidean cutoff equal to one near its compact coordinate image. Thus \[
 \beta=\sum_i q(t_{s_i})\varphi_i(t_{\widetilde f_{i,1}},\ldots)
                          =\Phi(T)\quad\text{on }V_E,   \tag{8}
\] for one global unit tuple \(T\) and a Euclidean smooth \(\Phi\). All
group translates of these units exist on \(V_E\); adjoining them gives
one fixed tuple. The real average \(|G|^{-1}\sum_g\Phi(T\circ g)\)
equals \(\beta\) there. This construction works for every orbit and
stabilizer, including singleton fibres. It does not assume exact
invariant ordinary scalar descent.

\subsection{T5. Polarization removes the orbit-correlation
difficulty}\label{t5.-polarization-removes-the-orbit-correlation-difficulty}

For a smooth Euclidean function on a neighbourhood of a compact box,
there are real polynomials converging through order two on that box.
Here is a proof of precisely the version used. Extend it by a compact
Euclidean cutoff, convolve with a Gaussian, and differentiate under the
integral. Uniform continuity of the derivatives through order two and
the Gaussian tail give uniform derivative convergence on the box. For
fixed convolution width, the Gaussian and its first two derivatives are
uniform limits of their Taylor polynomials on the compact difference of
the box and the original coefficient support. Integrate these
polynomials to get polynomial approximants through order two, then
choose widths and truncation degrees diagonally.

For \(k\ge1\) and real symbols \(z_1,\ldots,z_k\), including repeated
symbols, inclusion-exclusion and the multinomial expansion give \[
 z_1\cdots z_k=\frac1{k!}\sum_{J\subseteq\{1,\ldots,k\}}
            (-1)^{k-|J|}\left(\sum_{j\in J}z_j\right)^k. \tag{9}
\] Indeed each term missing one symbol cancels, and the only surviving
degree-\(k\) term containing every symbol is their product with
coefficient \(k!\). For unit logarithms, a sum is the logarithm of the
corresponding product unit. Consequently the real group average of
\textbf{every} scalar polynomial in an arbitrary tuple of unit
logarithms is a finite real linear combination of constants and terms \[
        T_{f,k}=\frac1d\sum_{g\in G}t_{gf}^{k},          \tag{10}
\] with actual product units \(f\). This recovers mixed orbit
correlations; no claim that separate symmetric sums for the original
individual units suffice is used. The number of product units may grow
with polynomial degree. T9 handles this explicitly; no common
neighbourhood for infinitely many product units is assumed.

\subsection{T6. Actual root valuations and all coefficient-zero
cases}\label{t6.-actual-root-valuations-and-all-coefficient-zero-cases}

For a unit \(f\in B^\times\) set \[
 P_f(Z)=\prod_{g\in G}(Z-gf)
             =Z^d+\sum_{i=0}^{d-1}c_iZ^i.               \tag{11}
\] Its coefficients lie in \(B^G=A\); \(c_0\) is a unit. At an analytic
point \(y\), the roots \(gf(y)\) are actual nonzero elements of
\(\mathcal H(y)\) with multiplicities indexed by \textbf{all} group
elements. This includes repeated roots and arbitrary stabilizers.
Neither a separable orbit polynomial nor characteristic-zero
coefficients are required.

We prove the root-valuation rule rather than importing Newton-polygon
recognition. For a polynomial \(\prod_{j=1}^d(Z-r_j)\) over any valued
field with nonzero roots, sort \(u_j=-\log|r_j|\) in nondecreasing
order. The coefficient of \(Z^{d-k}\) is a sum of products of \(k\)
roots. Every summand has valuation at least \(u_1+\cdots+u_k\), so its
coefficient has valuation at least that number, allowing infinity. If
\(u_k<u_{k+1}\), the subset of the \(k\) roots of smallest valuation is
unique and its product has strictly greater norm than every other
summand. Ultrametricity makes the coefficient valuation exactly that
sum. Equality also holds at \(k=0,d\). These equality points are
precisely the ends of constant-valuation blocks. The convex polygon
joining them lies below all the coefficient points and has its
successive slopes equal to the root valuations in reversed sign when
indexed by coefficient exponent. Thus the lower convex hull of \[
          (i,v(c_i))\ (0\le i\le d),\qquad v(c_d)=0      \tag{12}
\] has each segment \(i<j\) contributing \(j-i\) roots of valuation
\((v(c_i)-v(c_j))/(j-i)\). Repeated roots, ties, and every cancellation
are covered by the strict-gap argument. No division occurs in the valued
field.

Near a fixed downstairs \(x\), choose \(M\) so all \(|t_{gf}|<M\) on its
finite fibre and retain this inequality by full-fibre shrinking. Then
\(|t_{c_0}|<dM\). If \(c_i(x)=0\) for \(0<i<d\), shrink further until \[
                      |c_i|<\exp(-dM-1).                \tag{13}
\] Its valuation lies strictly above the chord from \((0,t_{c_0})\) to
\((d,0)\), since that chord's height is at most \(dM\). Hence it cannot
belong to the lower hull and may be omitted. If \(c_i(x)\ne0\), shrink
to make it a downstairs unit: choose a strict affinoid neighbourhood
inside its nonzero locus, then use the unit criterion of T3. Keep
\(c_0,c_d=1\) in all cases. The retained index set
\(I\subseteq\{0,\ldots,d\}\) therefore includes both ends. Its other
coefficients are actual units; all omitted coefficients obey (13)
throughout this neighbourhood. Nilpotent coefficients are included in
the zero-seminorm case. No coefficient with value zero is logarithmed,
and small coefficients are not simply assumed harmless.

\subsection{\texorpdfstring{T7. A \(C^1\) piecewise polynomial and its
tangential
mollification}{T7. A C\^{}1 piecewise polynomial and its tangential mollification}}\label{t7.-a-c1-piecewise-polynomial-and-its-tangential-mollification}

For arbitrary real heights \(z_i\), \(i\in I\), with \(z_d=0\), form the
lower hull of the finitely many points \((i,z_i)\). If consecutive
vertices are \(i<j\), put \(m=j-i\) and \(u=(z_i-z_j)/m\). Define \[
               H_{I,k}(z)=\sum_{\text{hull segments }i<j}
                                 m\,u^k.                \tag{14}
\] Collinear vertices may be retained or omitted: (14) is unchanged.
There are finitely many possible hull chains. Each is selected by linear
weak inequalities in the heights, so (14) is a polynomial on each region
of a finite polyhedral subdivision of the whole height space. Equations
(11)--(13) prove at \textbf{every} point upstairs that \[
               T_{f,k}=d^{-1}H_{I,k}(t_{c_i}).            \tag{15}
\]

We need more than continuity of this formula. For a general smooth
one-variable \(\phi\), the segment term \(m\phi((z_i-z_j)/m)\) has
gradient \(\phi'(u)(e_i-e_j)\) and Hessian \[
            \frac{\phi''(u)}m(e_i-e_j)\otimes(e_i-e_j).  \tag{16}
\] If a vertex is inserted or deleted on a collinear segment, all the
affected \(u\) equal one common value. The gradients then telescope:
\(\sum\phi'(u)(e_i-e_j)\) is exactly the gradient of the single merged
segment. This remains true if several vertices are collinear. At a
height configuration on a chamber boundary, every limiting hull chain
refines the unique lower hull by such collinear vertices. Values and
gradients therefore agree for every limiting chamber. A finite piecewise
smooth function with these matching values and gradients is \(C^1\):
integrate the matching, continuous candidate gradient on short line
segments subdivided by the finitely many chamber faces; segments
contained in faces follow by a small translation and continuity. The
usual differentiability remainder follows from continuity of that
gradient. Thus (14) is \(C^1\).

On any compact neighbourhood, slopes \(u\) are bounded. Each Hessian
(16) is bounded by \(2\sup|\phi''|/m\) in operator norm; there are at
most \(d\) segments. For \(\phi(u)=u^k\), there is a finite bound for
all chamber Hessians on that neighbourhood. The gradient is locally
Lipschitz by integrating these Hessians along line segments, using its
continuity at faces as above. The weak Hessian is the bounded piecewise
Hessian, with no singular face measure: integrate by parts on each
polyhedral region; the boundary contributions from the continuous
gradient cancel in pairs. Lower-dimensional intersections give no
additional boundary term. This proves all assertions used in
mollification; bounded third or fourth derivatives are unnecessary.

Choose a nonnegative even \(C^\infty\) kernel \(\rho\) supported in the
unit ball with integral one. Let \(H_\varepsilon=H*\rho_\varepsilon\).
The local gradient Lipschitz bound yields uniformly on fixed compacts
\(H_\varepsilon\to H\) and \(\nabla H_\varepsilon\to\nabla H\); in fact
the errors are \(O(\varepsilon^2)\) and \(O(\varepsilon)\). The first
bound follows from the first-order remainder and the kernel's zero first
moment. Its Hessian is the convolution of the weak Hessian and is
uniformly bounded on those compacts.

\textbf{Tangential second derivative lemma.} Let an open segment \(L\)
lie in the relative interior of a face \(F\) of the subdivision of
\(H\), with tangent vector \(v\) in the affine direction of \(F\). At
every point \(z\in L\), \[
        \nabla^2H_\varepsilon(z)(v,v)
                  \longrightarrow (H|_L)''(z;v,v).       \tag{17}
\] To prove this, only chambers incident to \(F\) meet a sufficiently
small ball about \(z\). Let their polynomial formulas be \(Q_a\). Each
agrees with \(H\) on a relatively open neighbourhood in \(F\), so
\(\nabla^2Q_a(z)(v,v)\) is the same tangential derivative of that
restriction. For points \(w\) tending to \(z\) through any such chamber,
the corresponding tangential Hessian tends to this common value.
Convolving the bounded, piecewise Hessian proves (17). If the line is
mapped to one point in height space, the tangent is zero and the
statement is immediate. This proof deals with segments \textbf{entirely
contained in collision walls}; they are not dismissed as measure zero. A
compact moment segment maps affinely to height space and can be
subdivided at finitely many face preimages, so the conclusion holds
almost everywhere on every such segment, including those of deficient
height-map rank.

\subsection{T8. Calibrated differentiated approximation of a fixed
trace}\label{t8.-calibrated-differentiated-approximation-of-a-fixed-trace}

Fix \(f,k\) and one neighbourhood on which (13)--(15) hold. The
functions \[
                b_\varepsilon=d^{-1}H_\varepsilon(t_{c_i})
                                                               \tag{18}
\] are ordinary smooth downstairs, using actual unit moments. On a
compact upstairs domain add to one moment tuple the units \(gf\), the
pullbacks of all retained coefficients, and any fixed cutoff's units.
Its image is a finite polyhedral set. Refine its one-dimensional cells
by the height-fan preimages and by a common skeletal calibration
subdivision. On each segment, (15) is an identity of functions of affine
coordinates. Its right side restricts to a polynomial on the segment,
whose first and second tangent derivatives are therefore those of the
ordinary upstairs polynomial \(T_{f,k}\).

The errors \(e_\varepsilon=T_{f,k}-\pi^*b_\varepsilon\) and their first
derivatives converge uniformly tangentially. Their second tangential
derivatives are uniformly bounded and converge almost everywhere on
every segment by (17). No differentiation of (15) as a nonsmooth
ordinary section is performed; the derivatives are compared on actual
moment cell tangents.

Here is why this is a total-variation statement on the actual carrier.
For the enlarged tuple \(t_1,\ldots,t_r\) put
\(\zeta=\sum_a d't_a\wedge d''t_a\), with the positive top-order
convention. On a segment with tangent \(v\), a coefficient matrix
\(A=(A_{ab})\) contracts to \(\sum A_{ab}v_av_b\) times its positive
calibration; \(\zeta\) contracts to \(\sum v_a^2\) times the same
calibration. Hence \(|A_{ab}|\le M\) gives \[
       |\eta_{\mathrm{Rad}}|\le rM\zeta_{\mathrm{Rad}},    \tag{19}
\] by \(\big(\sum|v_a|\big)^2\le r\sum v_a^2\). Equivalently use bounded
tangential coefficients in unit tangent frames on the finitely many
segments. The estimate is made on the supporting skeletal cells before
pushforward. T2(5) gives finite total mass of \(\zeta\) on the compact
domain and a constant finite positive density on each moment segment.
Thus bounded tangential error converging almost everywhere has carrier
total variation tending to zero by bounded convergence on those finite
calibrated segments. Zero-dimensional faces carry no top measure.
Multiplicities and generic lengths multiply both sides of (19); none is
removed or replaced by a point count.

For any \textbf{fixed} ordinary smooth cutoff \(\lambda\) supported in
the neighbourhood, the product rule for \(D(\lambda e_\varepsilon)\)
uses the scalar error, its two first derivatives, and its mixed second
derivative, together with fixed bounded derivatives of \(\lambda\). The
preceding convergence and (19) show \[
       \big\|(D((\lambda\circ\pi)e_\varepsilon))_{
                          \mathrm{Rad},E}\big\|_{\mathrm{TV}}
                                 \longrightarrow0.       \tag{20}
\] Cover a fixed compact cutoff support by finitely many compact moment
neighbourhoods to use this argument; their enlarged tuples and
calibration refinements are fixed during this limit. Signed fibre
cancellation is not used to infer (20).

For a fixed finite combination of traces (10), T6--T8 apply near every
downstairs point, with finitely many units, omitted coefficients and
strict inequalities at that point. If a fixed downstairs cutoff \(\chi\)
has compact support in the original unit domain, cover its support by
finitely many of these local neighbourhood interiors, and choose an
inner compact partition \(\theta_a\) by (7). Let \(b_{a,\varepsilon}\)
be the local combination of (18). Put \[
        a_\varepsilon=
              \sum_a\widetilde{\chi\theta_a b_{a,\varepsilon}},     \tag{21}
\] where tilde means extension by zero. This is ordinary smooth and
supported in \(\operatorname{supp}\chi\). A product supported compactly
inside its ordinary open is identically zero on a neighbourhood of each
point outside that open, so this extension is smooth and has no boundary
derivative term. The sum-one identity near \(\operatorname{supp}\chi\)
and (20), for this finite fixed cover, prove \[
 \big\|(D((\chi\circ\pi)T-\pi^*a_\varepsilon))_{
                      \mathrm{Rad},E}\big\|_{\mathrm{TV}}\to0       \tag{22}
\] for the given invariant scalar polynomial \(T\). All supports here
lie in \(\pi^{-1}(\operatorname{supp}\chi)\), independently of
\(\varepsilon\). There is no assertion of positive-test preservation.

\subsection{T9. Every compact ordinary scalar test; the two limits and
support}\label{t9.-every-compact-ordinary-scalar-test-the-two-limits-and-support}

Let invariant \(\beta\in A_c^{0,0}(X_E)\) be arbitrary, and set
\(F=\pi(\operatorname{supp}\beta)\). It is compact. Cover \(F\) by
finitely many ordinary interiors \(U_i\) of strict affinoids \(V_i\) on
which (8) holds, and choose outer cutoffs \(\chi_i\) from (7), compactly
supported in \(U_i\), with sum one near \(F\). Put
\(C=\bigcup_i\operatorname{supp}\chi_i\). It is a fixed compact set
containing \(F\).

On \(V_{i,E}\) approximate the fixed Euclidean coefficient \(\Phi_i\)
through order two on a box containing the fixed unit moment image by
polynomials \(P_{i,j}\). Choose the box to contain all translated tuple
images as well. Let \(T_{i,j}\) be the real group average of
\(P_{i,j}(T_i)\). Since \(\beta\) is invariant, it is also the average
of \(\Phi_i(T_i)\), so \(\beta-T_{i,j}\) tends to zero in coefficient
\(C^2\) on a \textbf{fixed upstairs} moment tuple containing all the
translates. Include fixed moment descriptions of \(\chi_i\circ\pi\) on a
finite compact cover of its support. The product rule and (19) give \[
 \big\|(D((\chi_i\circ\pi)(\beta-T_{i,j})))_{
                     \mathrm{Rad},E}\big\|_{\mathrm{TV}}\to0.       \tag{23}
\]

For each fixed \(j\), use (9) to express \(T_{i,j}\) as finitely many
traces. Apply (21)--(22) with outer cutoff \(\chi_i\) to obtain
\(a_{i,j}\) supported in \(\operatorname{supp}\chi_i\) and with \[
 \big\|(D((\chi_i\circ\pi)T_{i,j}-\pi^*a_{i,j}))_{
                     \mathrm{Rad},E}\big\|_{\mathrm{TV}}<1/j       \tag{24}
\] (or \(1/(ij)\) after any fixed finite enumeration). Each
fixed-polynomial local cover, inner partition and moment refinement is
chosen \textbf{before} its smoothing limit. It may depend on \(j\). Its
derivative bounds may grow with \(j\); smoothing is chosen after those
bounds. This causes no problem because (23) is evaluated with the fixed
outer cutoffs, before the inner localization of \(T_{i,j}\). The varying
inner cutoffs never multiply the polynomial approximation error in (23),
and every approximate support remains inside the same outer \(C\).
Neither a degree-uniform estimate nor an infinite simultaneous unit
shrink is asserted.

Set \(\alpha_j=\sum_i a_{i,j}\). Globally
\(\beta=\sum_i(\chi_i\circ\pi)\beta\), since the sum is one near its
support and it vanishes off that support. Equations (23)--(24) and the
finite triangle inequality prove (2). The support of every
differentiated error is contained in \(\pi^{-1}(C)\), compact by T3.
Thus for \textbf{every} continuous finite multiplier \(h\), \[
       B_{h\circ\pi}(\beta-\pi^*\alpha_j)\longrightarrow0,
                                                               \tag{25}
\] using its boundedness on that compact set. This proves the actual
ordinary scalar test gate for curves, not merely \(C^0\) density of
invariant functions.

\subsection{T10. Forward Galois ascent and actual normal-closure
descent}\label{t10.-forward-galois-ascent-and-actual-normal-closure-descent}

For a general compact ordinary scalar test \(\beta\) upstairs, (6),
ordinary exterior functoriality on the restricted carrier, and
invariance of \(h\circ\pi\) show
\(B_{h\circ\pi}(g^*\beta)=B_{h\circ\pi}(\beta)\). Its real average
\(\beta^{\mathrm{av}}=d^{-1}\sum_g g^*\beta\) has the same pairing. Its
support is in a compact finite group saturation. Real averaging does not
divide by \(d\) in \(K\), including when \(\operatorname{char}K\)
divides \(d\).

For every downstairs compact test \(\alpha\), its pullback is compact.
The scalar projection formula (4), with \textbf{no degree factor}, gives
\[
  B_{h\circ\pi}(\pi^*\alpha)
     =\int_{X_E}(h\circ\pi)(\pi^*D\alpha)_{\mathrm{Rad},E}
     =\int_X h(D\alpha)_{\mathrm{Rad},K}=B_h(\alpha).     \tag{26}
\] Smooth exterior functoriality suffices; no continuous-current
pullback has been presumed. If the downstairs current is zero,
approximate \(\beta^{\mathrm{av}}\) by (2). Its approximate pairings
vanish by (26), and (25) proves that its pairing, hence that of every
upstairs test, vanishes. This proves forward ascent for every finite
Galois extension, including wild ones, on all curves in the stated
scope.

Conversely, (26) alone proves reverse descent whenever the upstairs
current is zero. This calculation works for every finite separable
extension, without a group, by compactness, purity, and (4). It is an
actual proof of the descent needed below, not a use of an earlier note.

For arbitrary finite separable \(E/K\), take the finite Galois normal
closure \(N/K\) constructed in T3, with its unique extending absolute
value and completeness. Forward Galois ascent gives zero current on
\(X_N\). The space \(X_E\) retains every ambient hypothesis and pure
dimension one by T3, and \(N/E\) is finite separable. There is a
canonical analytic identification \((X_E)_N=X_N\): on affinoids it is
the associative tensor isomorphism
\((A\otimes_K E)\otimes_E N=A\otimes_K N\) with the equivalent
finite-module topologies. The pullback of \(h\circ\pi_E\) is the
original pullback on \(X_N\). Apply the just-proved reverse calculation
to \(N/E\), testing each compact ordinary function on \(X_E\), to obtain
zero current on \(X_E\). This proves (1) at every finite separable
stage. It is the finite dependency, before the completed stage in T13.
Ordinary opens retain the stated assumptions and the same argument
applies to their scalar inverse images.

\subsection{T11. Finite inseparability and arbitrary finite algebraic
stages}\label{t11.-finite-inseparability-and-arbitrary-finite-algebraic-stages}

\subsubsection{T11.1. Frobenius, topology and exact scalar
tests}\label{t11.1.-frobenius-topology-and-exact-scalar-tests}

Let \(k\) be complete, nontrivially valued, of characteristic \(p>0\),
and let \(F/k\) be finite purely inseparable. Choose \(q=p^e\) with
\(F^q\subset k\). T3's bounded-coordinate argument applies to
\textbf{every} finite extension, and gives
\(A\widehat\otimes_k F=A\otimes_k F\) with its finite-module topology.
Consequently every \(u=\sum a_i\otimes\lambda_i\) satisfies \[
                    u^q=\sum a_i^q\lambda_i^q\in A.       \tag{27}
\] This is an identity in the actual algebra, including nilpotents.

Scalar projection is surjective by the affine foundation in T2; T12.1
also gives a direct proof for all fields used here. If two points of
\(\mathcal M(A\otimes F)\) project to \(x\), (27) forces their values on
every \(u\) to equal \(|u^q(x)|^{1/q}\). They are the same seminorm. The
projection is thus a continuous bijection from a compact space to a
Hausdorff space, hence a homeomorphism. On a good separated space these
affine homeomorphisms glue on the ordinary interiors of downstairs
affinoid neighbourhoods. No covering-space assertion is used.

For a unit \(f\in(A\otimes F)^\times\), (27) gives \(f^q=g\in A\). The
function \(g\) has no analytic zero downstairs, by surjectivity, so it
is an actual unit by T3's unit criterion. On the upstairs carrier
\(t_f=q^{-1}t_g\). Now let \(\beta\in A_c^{0,0}(X_F)\) be arbitrary. Its
support has compact image \(C\) under the homeomorphism. For each point
of \(C\), select an ordinary local description of \(\beta\), then a
downstairs affinoid neighbourhood \(V\) whose full inverse image lies in
that description open. The inverse image is \(V_F\), so all coordinate
units restrict to elements of its actual affinoid algebra. Taking their
\(q\)th powers gives an ordinary downstairs scalar description
\(\varphi(q^{-1}t_{g_1},\ldots,q^{-1}t_{g_m})\) on
\(\operatorname{int}_X V\). Equality of numerical smooth scalar
functions is equality of their ordinary \((0,0)\) sections, as explained
in T12.3. Cover \(C\) by finitely many such opens, apply the compact
partition of T4, multiply these local descriptions by its cutoffs, and
extend by zero. This produces \(\alpha\in A_c^{0,0}(X)\) with
\(\beta=\pi_F^*\alpha\) everywhere. Compact support inside each open
makes the extension ordinary smooth; no derivative is extended through
an exposed boundary.

Goodness, strictness, separatedness and boundarylessness are preserved
as in T3. Purity is also preserved, and here is a proof that does not
assume equality of total dimensions gives local purity. Given an
upstairs point \(y\), a downstairs affinoid neighbourhood of its image
has dimension one, and its scalar extension has dimension one, giving
the upper local bound. If local dimension at \(y\) were zero, take this
affinoid upstairs and remove all its positive-dimensional irreducible
components, none of which contains \(y\). An affinoid neighbourhood
inside the remaining ordinary open has dimension zero by CLD 11.2.3's
component/domain formula. By the homeomorphism choose a downstairs
affinoid neighbourhood whose full inverse image is inside that open. Its
downstairs dimension is one (it is a domain in the pure curve), and its
upstairs total dimension is one by scalar dimension invariance, a
contradiction. Existence of the smaller downstairs affinoid
neighbourhood uses goodness. This argument only checks dimension; it
never reduces the measure.

On these pure curves \(D\beta=\pi_F^*D\alpha\), and the smooth
top-measure identity (4) proves \[
             B_{h\circ\pi_F}(\beta)=B_h(\alpha).         \tag{28}
\] All supports are compact. The converse is obtained by taking
\(\beta=\pi_F^*\alpha\) for an arbitrary compact downstairs scalar test.
Thus finite purely inseparable zero-Hessian ascent and descent are
proved here directly. No earlier inseparable manuscript is a premise.

\subsubsection{T11.2. The finite mixed-field
decomposition}\label{t11.2.-the-finite-mixed-field-decomposition}

Let \(E/K\) be any finite algebraic extension. In characteristic zero
every irreducible polynomial has nonzero derivative, and is separable,
so T3--T10 apply. In characteristic \(p\), choose finitely many
generators \(a_1,\ldots,a_r\) of \(E/K\). For each irreducible minimal
polynomial, repeatedly remove common factors \(p\) from its exponents
until it has the form \[
                f_i(Z)=g_i(Z^{p^{e_i}}),\qquad g_i'\ne0.
                                                               \tag{29}
\] The polynomial \(g_i\) is irreducible: a factorization of \(g_i\)
would give a factorization of \(f_i\) by substitution. Since \(g_i\) is
irreducible and its derivative is nonzero of smaller degree,
\(\gcd(g_i,g_i')=1\); its roots are distinct. Therefore
\(b_i=a_i^{p^{e_i}}\) is separable over \(K\). The compositum
\(S=K(b_1,\ldots,b_r)\subset E\) is finite separable: adjoining a root
of a separable polynomial remains separable over an intermediate field
because its minimal polynomial divides that polynomial; extending
embeddings successively gives degree-many embeddings of the tower. An
intermediate simple field has degree-many embeddings as well: each
embedding has at most the remaining extension degree lifts, while the
entire tower has its full degree of embeddings. Thus its minimal
polynomial is separable. This proves the field facts used here without
assuming separability of an arbitrary algebraic stage.

With \(e=\max e_i\), every \(a_i^{p^e}\in S\), and taking \(p^e\)th
powers in a rational expression in the \(a_i\) shows
\(E^{p^e}\subset S\). Hence \(E/S\) is finite purely inseparable. Both
fields are complete under the prescribed absolute value by the finite
norm proof in T3. This gives the \textbf{actual embedded} tower
\(K\subset S\subset E\), not an abstract factorization of degrees.

The finite tensor associativity with bounded inverse maps, proved in
T12.1, identifies \((X_S)_E=X_E\) as analytic spaces, with the actual
algebra and not its reduction. \(X_S\) has all ambient hypotheses and
pure curve dimension by T3. Finite separable ascent to \(S\), then the
proved Frobenius ascent to \(E\), proves (1) at every finite algebraic
stage. Reverse descent uses (4) directly. Generic lengths are those of
the actual tensors throughout. No degree factor is introduced into
scalar projection.

\subsection{T12. Completion, rational neighbourhoods and every compact
scalar
test}\label{t12.-completion-rational-neighbourhoods-and-every-compact-scalar-test}

\subsubsection{T12.1. Fields, completed tensors, surjectivity and
compactness}\label{t12.1.-fields-completed-tensors-surjectivity-and-compactness}

Write \(M=\bigcup_E E\), where \(E\) ranges over its finite
subextensions of \(K\). It is a directed union: a finite list of
algebraic elements generates a finite field, by successive
multiplication of finite degrees. Each \(E\) is complete by T3. We never
regard \(M\) as a complete analytic field unless it happens already to
be complete. Its completion \(L\) is a complete valued field. For a
\(K\)-affinoid algebra \(A\), the image of \(A\otimes_K M\) is dense in
\(A_L=A\widehat\otimes_K L\). Indeed finite algebraic tensors are dense
by definition of completion; a tensor \(\sum a_j\otimes\lambda_j\) can
be approximated by replacing its finitely many \(\lambda_j\in L\) with
\(m_j\in M\), with tensor error at most
\(\max_j\|a_j\||\lambda_j-m_j|\). Approximation statements refer to the
\textbf{image}, so no injectivity or isometry of an uncompleted infinite
tensor is presumed.

Here is the associative tensor argument actually required. For finite
\(E/K\) with a basis \(e_i\), the bounded coordinate maps give
\(A_E=A\otimes_K E\) with its finite-module topology. The maps on
algebraic tensors \[
 (a\otimes e)\otimes\lambda\mapsto a\otimes e\lambda,
 \qquad a\otimes\lambda\mapsto(a\otimes1)\otimes\lambda
                                                               \tag{30}
\] extend continuously to inverse maps
\((A_E)\widehat\otimes_E L\simeq A\widehat\otimes_K L\). For the first
map, the bounded multiplication map
\((b,\lambda)\mapsto\sum_i a_i\otimes e_i\lambda\), where
\(b=\sum_i a_i\otimes e_i\), is bounded in \(b\) by the coordinate
estimates and in \(\lambda\) by the tensor inequality. It is balanced
over \(E\). The reverse map is bounded by the tensor inequality and the
bounded inclusion \(A\to A_E\). These maps are inverse on dense
elementary tensors, hence inverse on the completions. The same proof
works when \(A\) is affinoid with arbitrary radii, and for two finite
stages. Thus all stage projections and inverse images of rational
domains used below are actual analytic base changes. No associativity
claim for an infinite uncompleted Banach tensor is needed.

For clarity, scalar surjectivity in the present completed-algebraic
scope also has a direct seminorm proof. Given \(x\in\mathcal M(A)\), put
\(H=\mathcal H(x)\). Embed the algebraic extension \(M\) into an
algebraic closure \(\overline H\) over \(K\). Such an embedding is
obtained by successively adjoining roots, or by maximizing partial
embeddings and adjoining a root of the minimal polynomial if the domain
is not all of \(M\). Equip \(\overline H\) with the extending absolute
value: the finite extension existence and uniqueness proof of T3 makes
the norms on its finite \(H\)-subfields compatible. On every finite
\(K\)-subfield of \(M\) this value is the specified one, again by
uniqueness. Completing \(\overline H\) gives a complete field \(F\) into
which the isometric embedding \(M\to F\) extends to \(L\to F\). For
\textbf{every} tensor representation with \(\lambda_j\in L\), evaluation
satisfies \[
 \left|\sum_j a_j(x)\lambda_j\right|
       \le \max_j\|a_j\||\lambda_j|.
                                                               \tag{31}
\] Taking the infimum over all representations bounds it by the full
tensor seminorm. It extends to \(A_L\) and gives an analytic point over
\(x\). This proof handles the separated quotient in the tensor
completion and uses no claimed isometry of \(A\otimes M\). The identical
proof over \(E\) gives surjectivity \(X_L\to X_E\) and surjectivity on
every scalar-extended affinoid domain. On good spaces these affine facts
cover the whole space.

The inverse image of a compact \(C\subset X_E\) is compact: finitely
many ordinary interiors of affinoid neighbourhoods cover \(C\). Its
inverse image is closed in the finite union of their compact scalar
extensions, since the spaces are Hausdorff. This works for the infinite
completed projection. It does not assert a finite fibre, a finite Galois
cover or any infinite group averaging.

For a nonzero \(\lambda\in L\), choose \(m\in M\) with
\(|m-\lambda|<|\lambda|\). Then \(|m|=|\lambda|\). Thus
\(|L^\times|=\bigcup_E|E^\times|\). This observation handles any
value-group radius if desired, but the rational argument below allows
\textbf{all positive real radii}, so it needs no strictifying extension
and no approximation of a radius.

\subsubsection{T12.2. Exact rational and strict-locus
descent}\label{t12.2.-exact-rational-and-strict-locus-descent}

Let \(A\) be affinoid over a finite stage \(E_0\), set
\(B=A\widehat\otimes_{E_0}L\), and let a rational domain of
\(\mathcal M(B)\) be defined by \[
             |f_i|\le r_i|g|\quad(1\le i\le s),
             \qquad r_i>0,                              \tag{32}
\] where \(g,f_1,\ldots,f_s\) have no common analytic zero. Empty
systems can instead be handled by the whole affinoid with the harmless
condition \(|1|\le2|1|\). Compactness gives \[
 \delta=\min_{y\in\mathcal M(B)}
             \max\{|g(y)|,r_1^{-1}|f_1(y)|,\ldots,
                                  r_s^{-1}|f_s(y)|\}>0.
                                                               \tag{33}
\] Choose simultaneous finite-stage approximants \(g',f_i'\) in one
\(A_E\) with Banach errors less than \(\delta\min(1,r_1,\ldots,r_s)\).
This is possible by T12.1. Their evaluation errors satisfy \[
                 |g'-g|<\delta,
                 |f_i'-f_i|<r_i\delta                 \tag{34}
\] at every point upstairs. Neither vanished nor nilpotent functions are
logarithmed, and none is assumed to have small Banach norm merely
because its spectral norm is zero. They are approximated in the actual
Banach algebra just like the other elements.

Both conjunctions, (32) and the conjunction of its strict versions, are
\textbf{unchanged} by (34). To check all boundary cases, first suppose
\(|g(y)|\ge\delta\). Then \(|g'(y)|=|g(y)|\). For each \(i\) with
\(|f_i(y)|\ge r_i|g(y)|\), (34) forces \(|f_i'(y)|=|f_i(y)|\). For each
\(i\) with \(|f_i(y)|<r_i|g(y)|\), it instead forces
\(|f_i'(y)|<r_i|g'(y)|\). Hence each weak inequality, each equality, and
each strict inequality is preserved at that point. If \(|g(y)|<\delta\),
(33) supplies an \(i\) with \(r_i^{-1}|f_i(y)|\ge\delta\). Its numerator
norm is unchanged, while \(|g'(y)|<\delta\). This inequality remains
strictly false, so both conjunctions remain false. These arguments also
cover \(g(y)=0\) and any number of zero numerators.

The weighted maximum for the approximants is everywhere positive
upstairs. Surjectivity to \(\mathcal M(A_E)\) then shows the
denominator/numerators have no common zero there. They generate the unit
ideal: otherwise a maximal ideal containing them gives an analytic
common zero. Thus they define a genuine rational analytic domain \(R_E\)
by CLD 11.1.5. Its scalar extension has exactly the set (32). Analytic
domain uniqueness, or the universal rational-domain presentation by
adjoining the bounded ratios, identifies it with the given domain
\(R_L\) as an analytic domain. The actual rational localization algebra,
with all nilpotents, is retained. The strict locus
\(O_E=\{|f_i'|<r_i|g'|\ \forall i\}\) is ordinary open in the ambient
affinoid, lies in its rational domain, and has inverse image
\textbf{exactly} the original strict locus \(O_L\). No invariance of
absolute analytic interiors or equality of all topological interiors is
asserted.

The stronger rational-neighbourhood statement of CLD 11.1.5 says that
for any \(y\) and any prescribed ordinary open in an affinoid, one can
choose (32) with \(R_L\) inside that open and \(y\) in its strict locus.
Thus the preceding construction supplies finite-stage ordinary
neighbourhoods of each chosen point, without controlling an entire
infinite scalar fibre. Only finitely many such choices are subsequently
needed on a compact support.

\subsubsection{T12.3. Exact finite-stage local
descriptions}\label{t12.3.-exact-finite-stage-local-descriptions}

Let \(\beta\) be ordinary smooth on \(X_L\) near \(y\), and let \(x\) be
its image in \(X\). Choose a downstairs affinoid neighbourhood \(V_0\)
and an ordinary open \(U_0\subset V_0\) containing \(x\). The inverse
image \(U_{0,L}\) is open. Intersect it with an ordinary moment
description open of \(\beta\); on that intersection write
\(\beta=\varphi(t_{f_1},\ldots,t_{f_m})\), with actual analytic units
and a smooth real coefficient defined on an open Euclidean set
containing the coordinate image. This can always be arranged locally by
shrinking the description near \(y\).

Inside \(V_{0,L}\) select a rational neighbourhood \(R_L\) of \(y\)
contained in this intersection, with \(y\) in its strict locus. T12.2
gives \(R_L=(R_{E_0})_L\) for a finite stage \(E_0\), and a descended
strict locus. Because \(R_L\subset U_{0,L}\), surjectivity on
\(R_{E_0}\) gives \(R_{E_0}\subset U_{0,E_0}\). Its strict locus is
therefore an ordinary open of the whole \(X_{E_0}\), not just an open of
an exposed affinoid boundary.

On the compact rational affinoid \(R_L\), put \(\delta_j=\min|f_j|>0\).
Apply density in the actual rational affinoid algebra, now viewed over
\(E_0\), to select a common finite \(E\supset E_0\) and
\(u_j\in\mathcal O(R_E)\) with \(\|u_j-f_j\|<\delta_j\) upstairs.
Ultrametricity gives \[
                     |u_j(y')|=|f_j(y')|
                                 \quad(y'\in R_L).     \tag{35}
\] By surjectivity \(u_j\) has no zero on \(R_E\), so it is an actual
unit by the unit criterion. Its inverse exists in that affinoid algebra;
no approximation of an inverse or reduction of the algebra is necessary.
Surjectivity also makes the finite-stage moment image of the \(u_j\)
equal to the upstairs image of the \(f_j\), so the same Euclidean
coefficient \(\varphi\) is defined on it.

On the inverse image \(O_E\) of the strict locus define the ordinary
smooth function \(b_E=\varphi(t_{u_1},\ldots,t_{u_m})\). Then
\(q_E^*b_E=\beta\) on \(O_L=q_E^{-1}O_E\). This is an exact
scalar-section identity. To justify the section claim precisely, CLD
3.4.6 embeds the \((0,0)\) G-form sheaf in numerical functions: on each
compact chart the nullity criterion 3.4.1/3.4.4 says that a zero
numerical coefficient is the zero section; compact charts G-cover the
space. Hence equal numerical scalar sections coincide, and ordinary
scalar sections are included in this sheaf. Alternatively, the combined
moment \((f,u)\) has image in the diagonal \(\{t_f=t_u\}\) by (35); the
difference of the two scalar coefficient functions restricts to zero
there, and moment nullity gives the same conclusion. No universal
moment-dimension equality is used. Smooth differentiation consequently
commutes with this \textbf{exact identity}, giving \(D\beta=q_E^*Db_E\)
on that open.

\subsubsection{T12.4. Compact scalar tests and extension by
zero}\label{t12.4.-compact-scalar-tests-and-extension-by-zero}

Now let \(\beta\in A_c^{0,0}(X_L)\) be arbitrary and put
\(F=\operatorname{supp}\beta\). If \(F\) is empty, take \(\alpha=0\) at
any finite stage. Otherwise T12.3 gives a neighbourhood
\(O_{i,L}=q_{E_i}^{-1}O_{i,E_i}\) and a descended ordinary function
\(b_{i,E_i}\) near each point of \(F\). Choose finitely many covering
\(F\). Enlarge to a common finite field \(E\subset M\) and base change
all these rational domains, strict loci and unit lists to \(E\).
Associativity (30) shows \(O_{i,L}=q_E^{-1}O_{i,E}\) and
\(\beta=q_E^*b_{i,E}\) there.

The compact set \(F_E=q_E(F)\) is covered by the finitely many
\(O_{i,E}\). Apply T4's explicit compact partition, built only from CLD
15.2.5, to obtain ordinary smooth \(\theta_i\) with compact support in
\(O_{i,E}\) and sum one near \(F_E\). Put \[
                \alpha=\sum_i\widetilde{\theta_i b_{i,E}},
                                                               \tag{36}
\] where tilde means extension by zero to \(X_E\). Each product is
ordinary smooth on its open, with support in the compact
\(\operatorname{supp}\theta_i\subset O_{i,E}\). Since that compact set
is closed in the Hausdorff \(X_E\), every point outside it has an
ordinary zero neighbourhood. Therefore its extension is ordinary smooth
at every exterior point as well. This proves (36) lies in
\(A_c^{0,0}(X_E)\), with no boundary derivative term.

On the inverse image of its defining open each summand pulls back to
\((\theta_i\circ q_E)\beta\). Off that open its pullback is zero, and
\(\theta_i\circ q_E\) is also zero. Thus, \textbf{globally}, \[
 q_E^*\alpha=(\sum_i\theta_i\circ q_E)\beta=\beta.       \tag{37}
\] On \(F\) the sum is one. Outside \(F\) the function \(\beta\) is
zero, so the equality also holds there regardless of that sum. No
descent or saturation of the original support is assumed. No agreement
of the separate \(b_{i,E}\) on downstairs overlaps needs to be presumed;
the finite sum is itself a smooth section, and (37) proves precisely its
required pullback. Its support is contained in a finite union of compact
cutoff supports. The inverse image of that union is compact by T12.1.
Differentiation commutes with the exact ordinary smooth identity (37),
so every compact upstairs scalar test and its Hessian are covered. This
is stronger than \(C^0\) density and does not use a differentiated
approximation topology at completion.

\subsubsection{T12.5. Ambient geometry and local purity after
completion}\label{t12.5.-ambient-geometry-and-local-purity-after-completion}

Goodness follows because inverse images of downstairs affinoid
neighbourhoods are affinoid and contain inverse images of their ordinary
interiors. Strict presentations base change to strict presentations. The
diagonal base change preserves separatedness. Boundarylessness follows
from GR 2.11(5'), as stated in T2; ordinary open restrictions retain it
by GR 2.11(3),(4).

We check local purity for \(X_L\), rather than just total dimension.
Around the image of \(y\in X_L\) take a downstairs strict affinoid
neighbourhood \(V_0\) of dimension one. Then \(V_{0,L}\) has dimension
one, giving the upper local bound. Suppose local dimension at \(y\) were
zero. In \(V_{0,L}\) discard its finitely many dimension-one
minimal-prime components, none of which contains \(y\). In the remaining
ordinary open choose a rational affinoid neighbourhood \(R_L\) of \(y\)
inside the inverse image of \(\operatorname{int}_X V_0\), with strict
inequalities at \(y\). The component/domain dimension formula in CLD
11.2.3 gives \(\dim R_L=0\). T12.2 descends \(R_L\) to a nonempty
\(R_E\) at a finite field, inside \(X_E\). The actual \(X_E\) is pure
one-dimensional by T3 and 11. Its domain \(R_E\) has dimension one,
including if its chosen real radii make the domain non-strict. Scalar
dimension invariance and analytic domain identification give
\(\dim R_L=\dim(R_E)_L=\dim R_E=1\), a contradiction. Hence every
upstairs point has dimension one. This argument keeps the original
strict ambient space and does not infer an endpoint theorem from a
non-strict radial core.

All top measures below are on the \textbf{actual} pure curves just
proved to exist. Their multiplicities need not remain numerically
unchanged point by point. They are the finite-flat fibre-algebra
dimensions in the canonical calibration. For reference, the usual
generic-length identity is obtained by tensoring a composition series of
a generic Artinian stalk with the flat upstairs generic stalk: each
residue-field quotient becomes the generic Artinian algebra of the
base-changed reduction. If its length is \(\mu\) and the original length
is \(\ell\), the length upstairs is \(\mu\ell\). This is CLD 13.2.8 with
\(\mathcal F=\mathcal O_X\); its proof body is included in Appendix S.
We use the calibration and the Artinian degree-sum proof of (4), not a
replacement by geometric point counts or a supposed constant extension
degree for an infinite extension.

\subsection{T13. Completed forward ascent and actual reverse
descent}\label{t13.-completed-forward-ascent-and-actual-reverse-descent}

Assume \(D[h]=0\) downstairs. Given \textbf{any}
\(\beta\in A_c^{0,0}(X_L)\), T12.4 supplies finite \(E\) and
\(\alpha\in A_c^{0,0}(X_E)\) with \(\beta=q_E^*\alpha\) exactly. Put
\(h_E=h\circ\pi_E\). Finite algebraic ascent, proved in T3--T11, gives
\(D[h_E]=0\). T12.5 verifies the pure curve and ambient hypotheses for
both \(X_E\) and \(X_L\). Smooth functoriality and the \textbf{\(L\)-
and \(E\)-calibrated} projection formula (4) yield \[
 \begin{aligned}
 B_{h\circ\pi}(\beta)
 &=\int_{X_L}(h_E\circ q_E)(q_E^*D\alpha)_{\mathrm{Rad},L}\\
 &=\int_{X_E}h_E(D\alpha)_{\mathrm{Rad},E}
 =0.
 \end{aligned}                                             \tag{38}
\] The relevant supports lie in the compact inverse image of
\(\operatorname{supp}\alpha\); the continuous finite multiplier is
bounded there. The signed-measure pushforward formula is therefore
legitimate. Since \(\beta\) was arbitrary, the upstairs zero-current
equation follows. There is no pullback of an arbitrary continuous
current in this proof and no field-degree factor in (38).

Conversely, suppose the upstairs equation holds. For an arbitrary
\(\psi\in A_c^{0,0}(X)\) its pullback has compact support by T12.1.
Apply smooth functoriality and (4) directly for \(L/K\): \[
             B_h(\psi)=B_{h\circ\pi}(\pi^*\psi)=0.        \tag{39}
\] Thus the downstairs equation holds. This reverse calculation is
included for completeness, not counted as new forward progress.
Ordinary-open localization follows from the ambient checks and the same
arguments on that open; no global quasicompactness or paracompactness is
added. This proves (1) with all its stated scope.

\section{Part C. Algebraically closed recognition, with the actual
structures}\label{part-c.-algebraically-closed-recognition-with-the-actual-structures}

Here \(k\) is complete, algebraically closed, and nontrivially valued.
The transport proof T2 already supplies affinoid foundations, compact
ordinary cutoffs, finite compact partitions, ordinary moment
presentations, and canonical calibrated integration. We need the
following additional exact source inputs.

CLD 11.2.5 states that affinoid irreducible components are finite in
number, that an irreducible analytic space is pure, that its proper
closed analytic subsets have smaller dimension, and that components of
analytic-domain intersections retain their dimensions. CLD 11.3.1
supplies finite normalization of a reduced analytic space, locally the
integral closure in its total quotient ring, and persistence of \(S_j\),
regularity, and normality on domains. CLD 11.4.8--9 identifies
quasi-smoothness with geometric regularity and, on good spaces,
smoothness with quasi-smoothness plus relative interior. CLD 11.4.1.3
gives \[
 \operatorname{Int}(V/Z)\cap f^{-1}\operatorname{Int}(Z)=\operatorname{Int}(V)
\] when the target is separated. Domain relative interior is ordinary
topological interior; finite maps are boundaryless. CLD 16.6.3 states
\(\delta_Z=\delta_{[Z]}\): the integration current is the sum over
reduced irreducible components with their actual generic lengths. Its
proof contracts canonical calibrations, not point-count measures. CLD
15.2.1(b) says the support of a \((p,q)\) form avoids a closed analytic
subspace of dimension less than \(\max(p,q)\). Thus a curve top form
vanishes in a neighbourhood of each rigid point, but a scalar test need
not vanish there.

The strict affinoid Nullstellensatz used here is BGR,
\emph{Non-Archimedean Analysis}, 6.1.2/1: over a complete nontrivially
valued field, every maximal ideal of a strict affinoid has residue field
finite over the ground field. A zero-dimensional strict affinoid is
consequently a finite-dimensional Artinian algebra over the field.
Indeed its semisimple quotient is a finite product of these fields, and
its nilpotent radical has a finite filtration with finite-dimensional
quotients. Its points are finitely many rigid points, with their actual
residue degrees and lengths. This standard foundation is stated and
cited exactly; no BGR PDF retrieval is claimed.

The calibrated specialization input is CLD 14.7.2 with \textbf{target
dimension \(d=1\)}: a compact pure \(r\)-dimensional source satisfying
\(S_1\), a regular pure one-dimensional target, a morphism pure of
relative dimension \(r-1\), a unit moment \(b\), and a rigid target
point have locally constant calibrated \((r-1)\)-dimensional fibre
tropicalization. For target \(\mathbb A^1\) and zero, the top calibrated
tropicalization of \((b,f)\) for sufficiently negative \(\log|f|\) is
its zero-fibre tropicalization times the standard negative interval. The
theorem permits source boundary. The complete theorem and proof,
including 14.7.1's boundary estimate and 14.7.3's finite-flat degree
argument, appear in Appendix S. No unsupported \(S_1\) inference in CLD
16.7.5 is used.

\subsection{C1. Lifting compact tests through closed
immersions}\label{c1.-lifting-compact-tests-through-closed-immersions}

If \(i:C\hookrightarrow V\) is a closed embedding of good separated
spaces, every compact ordinary test on \(C\) lifts to a compact ordinary
test of the same bidegree on \(V\). Locally express the test by unit
germs and a smooth Euclidean superform. The quotient sheaf
\(i_*\mathcal O_C=\mathcal O_V/\mathcal I\) lifts each germ. Its nonzero
value at the point permits an ordinary shrinking on which the lift is a
unit. Shrink also into the coefficient domain and into the given
presentation on \(C\). Thus one local ambient form restricts to the
given test everywhere on that intersection, including its zero regions.

Choose finitely many such opens covering the compact support. A compact
ambient cutoff \(\chi=1\) near that support can be supported inside
their union. The finite compact partition of T4 on
\(\operatorname{supp}\chi\) gives \(\theta_i\) supported in the
respective opens with sum one nearby. The form
\(\widetilde\eta=\chi\sum_i\theta_i\eta_i\), extended by zero, is
compact ordinary smooth. Its restriction is \(\eta\): wherever a term
contributes it restricts to \(\eta\), and the factors sum to one near
its support. Off the support the restriction is zero. Therefore \(i_*\)
is injective on currents, by evaluation on these lifts. This proves
injectivity, not pullback of arbitrary continuous currents.

\subsection{C2. Scalar tests at rigid points and removal of rigid
punctures}\label{c2.-scalar-tests-at-rigid-points-and-removal-of-rigid-punctures}

Every ordinary smooth scalar germ at a rigid point \(a\) is locally
constant, also over a nonclosed imperfect field. For a unit germ \(v\),
let \(q=|v(a)|>0\) and let \(P(T)=T^e+\sum_{j<e}c_jT^j\) be the minimal
polynomial of \(v(a)\) over the ground field. All conjugate roots have
norm \(q\), with inseparable repetitions, by the unique extending
valuation proved in T3. Hence \(|c_j|\leq q^{e-j}\) and \(|c_0|=q^e\).
On the ordinary neighbourhood \(|P(v)|<q^e\), neither \(|v|>q\)
(dominant leading term) nor \(|v|<q\) (dominant constant term) is
possible. Thus \(|v|=q\) there. Intersect for the finitely many units of
a scalar moment description. All its coordinates are constant. This
proves the assertion for every ordinary scalar section.

On a reduced good separated strict boundaryless pure curve \(V\), at any
rigid \(a\) there are compact ordinary smooth cutoffs \(\chi_m=1\) near
\(a\) whose supports shrink to \(a\), with \[
             \int_V D\chi_m=0,
             \|(D\chi_m)_{\mathrm{Rad}}\|_{\mathrm{TV}}\leq M
\] for one fixed finite \(M\). Here is a construction, including its
weights. Choose a reduced strict affinoid neighbourhood \(N\) of \(a\).
In its algebra choose \(f\in\mathfrak m_a\) outside every minimal prime.
Such a choice is possible: no one-dimensional component is contained in
the rigid point, and finite prime avoidance applies. Explicitly, express
generators of \(\mathfrak m_a\) and avoid the finitely many proper
linear subspaces of their coefficients over the infinite ground field. A
product of nonzero linear forms cannot vanish everywhere on that
coefficient space. Reduced rings have zero intersection of minimal
primes, so this \(f\) is a non-zero-divisor.

Its zero locus on \(N\) is pure zero-dimensional and consists of
finitely many rigid points. Remove the others and shrink inside
\(\operatorname{int}_V N\). At \(a\) the map \(f\) has relative
dimension zero. Zariski upper semicontinuity (CLD 11.4.5, Ducros's local
fibre-dimension theorem) gives an ordinary neighbourhood on which the
relative dimension is zero. Choose a smaller strict affinoid
neighbourhood \(N'\) there. Its zero fibre has support exactly \(a\),
with some length \(\ell>0\). Reduced noetherian rings satisfy \(S_1\):
an element outside all minimal primes is regular, and prime avoidance in
a positive-dimensional local ring produces one in its maximal ideal. The
property persists on analytic domains. Thus CLD 14.7.2 applies to
\(N'\), with \(r=d=1\), target \(\mathbb A^1\), and trivial auxiliary
moment. Its exact negative interval has calibration \[
                      d_0|e_1|,
                  \qquad d_0=\ell[\mathcal H(a):k].
\] The degree is the full degree, not its separable part: CLD 13.3.7 and
14.4.3 compute zero-dimensional canonical mass from fibre algebra
dimension.

Choose a real smooth \(\chi\), one for \(s\leq-1\), zero for \(s\geq0\),
valued in \([0,1]\). On \(N'\setminus\{a\}\) set
\(\chi_m=\chi(\log|f|+m)\) and extend as one at the zero and zero
outside \(N'\). These extensions are ordinary smooth. For any
neighbourhood \(H\) of \(a\) with closure inside
\(\operatorname{int}_V N'\), compactness gives
\(\min_{N'\setminus H}|f|>0\). For sufficiently large \(m\) the support
lies in \(H\) and avoids the domain boundary. The same argument for any
neighbourhood proves shrinking support. On the transition region
\(D\chi_m=\chi''(\log|f|+m)d'\log|f|\wedge d''\log|f|\). Take compact
nonnegative smooth \(b\geq|\chi''|\) on \(\mathbb R\). The positive form
with coefficient \(b\) dominates its absolute measure. Calibrated
integration gives total mass \(d_0\int b\), independent of \(m\), and
\(\int D\chi_m=d_0\int\chi''=0\). Stokes on the boundaryless curve gives
the same zero. No measure has been assigned to an unknown Hessian.

Now suppose \(g\) is continuous and \(D[g]=0\) off a locally finite
rigid set. Localize to one puncture \(a\). For a compact scalar test
\(\psi\), local constancy just proved implies \(\psi-\psi(a)\theta\) has
support disjoint from \(a\), where \(\theta=1\) near \(a\) is a compact
smooth cutoff. Thus \(D[g](\psi)=c\psi(a)\) with \(c=D[g](\theta)\) as a
functional on all ordinary scalar tests, without presuming order zero.
Evaluate on \(\chi_m\): \[
 |c|=\left|\int(g-g(a))D\chi_m\right|
 \leq M\sup_{\operatorname{supp}\chi_m}|g-g(a)|\longrightarrow0.
\] Continuity and shrinking support prove \(c=0\). This removes all
punctures locally. This argument is valid over every complete
nontrivially valued field.

\subsection{C3. Component weights and normalization
transport}\label{c3.-component-weights-and-normalization-transport}

On the ordinary interior \(O\) of a strict affinoid curve with algebra
\(A\), associated-cycle integration gives
\([g]=\sum_i m_i(i_i)_*[g|_{Z_i}]\), where \(Z_i\) are reduced
components and \(m_i=\operatorname{length}(A_{P_i})>0\) are the original
generic lengths. Distinct components intersect in a finite rigid set. On
the ordinary open where only \(Z_i\) remains, \(D[g]=0\) implies
\(m_i(i_i)_*D[g|_{Z_i}]=0\). C1 and the actual positive \(m_i\) imply
the intrinsic reduced-component equation there. C2 removes its rigid
intersections, giving \(D[g|_{Z_i}]=0\) on all of \(Z_i\cap O\).
Multiplicities may be unequal; they have been separated locally, not
divided by a global scalar. Combining the equations with generic lengths
also proves the equation on the reduction. This works pointwise on
ordinary neighbourhoods, so the reduced-curve equation holds globally
without choosing global component labels.

Let \(\nu:\widetilde Z\to Z\) be the normalization of a reduced strict
curve. On a reduced affinoid \(A\), its finite integral closure
\(\overline A\) lies in the total quotient ring. A common regular
denominator \(q\) for finitely many module generators gives
\(q\overline A\subseteq A\) and \(A[q^{-1}]=\overline A[q^{-1}]\). Thus
normalization is an actual analytic isomorphism away from a locally
finite rigid set. Its preimage is rigid and locally finite. Finite maps
give good separated strict boundaryless pure curves here. On the
isomorphism open, ordinary forms, \(D\), and calibration transport the
equation to \(g\circ\nu\). C2 removes the upstairs punctures, giving
\(D[g\circ\nu]=0\) throughout. This proves normalization transport
without inverting a finite-current direct image.

\subsection{C4. Local algebraic model and Wanner
comparison}\label{c4.-local-algebraic-model-and-wanner-comparison}

The exact embedding input is van der Put, \emph{The class group of a
one-dimensional affinoid space}, Ann. Inst. Fourier 30(4) (1980),
Theorem 1.1: every one-dimensional normal connected affinoid over a
\textbf{complete algebraically closed} non-Archimedean field is an
affinoid subspace of a nonsingular complete algebraic curve. His
affinoids are strict quotients of Tate algebras. The theorem's proof and
field assumption appear in Appendix S. The Berkovich-domain formulation
is explicitly used in Thuillier's proof of Theorem 2.3.8, printed page
53: a pure rig-smooth strict affinoid, after a finite separable
extension, is an affinoid domain of a smooth proper algebraic curve.
Here the extension is trivial. Affinoid-domain compatibility of rigid
and Berkovich geometry is Berkovich (1993), Section 1.6; the actual
algebra and domain structure are unchanged. No general category
equivalence on nonparacompact spaces is needed.

Every point \(z\) of a smooth strict boundaryless analytic curve has a
strict affinoid neighbourhood \(Q\). Select its connected component
containing \(z\). It is still a neighbourhood and is normal, reduced,
pure one-dimensional. Its algebra is a domain: distinct minimal
components of a noetherian normal ring are disjoint and yield
idempotents, and connectedness selects one. The embedding gives
\(j:Q\simeq Q'\subseteq C^{\mathrm{an}}\) for a smooth proper algebraic
curve over this same algebraically closed field. This is a domain
embedding, so it does not yet identify an ordinary open. Because both
ambient curves are separated and boundaryless, the boundary equality
above gives \[
 \operatorname{int}_Z Q=\operatorname{Int}(Q),\qquad
 \operatorname{int}_{C^{\mathrm{an}}}Q'=\operatorname{Int}(Q').
\] Absolute interior is intrinsic to the analytic space and preserved by
\(j\). It therefore identifies actual ordinary neighbourhoods of \(z\).
This treats every analytic point, not merely rigid points.

Wanner, arXiv:1801.04713v3, Theorem 4.7 has exactly the following scope:
\(C\) is a smooth proper algebraic curve over an algebraically closed
complete nontrivially valued field, \(W\subset C^{\mathrm{an}}\) is
ordinary open, and \(v:W\to\mathbb R\) is continuous. If \(v\) is
current-psh in her Definition 3.16, then it is Thuillier-subharmonic.
The field and curve assumptions are in her global conventions; ``psh''
here is positive CLD Hessian current, not the classical curve-test
class. The statement and its complete proof, and the definition of its
test forms, appear in Appendix S. Every one of her tests is an ordinary
analytic CLD test: its algebraic moment coordinates are analytic units
on that ordinary open, with the same smooth superform and compact
support. Canonical integration of this form is the CLD calibrated
integral. Hence a zero ordinary current gives her positivity for
\textbf{both} signs. Only this inclusion of test classes is used, not a
reverse identification of all analytic tests with algebraic ones.
Analytic isomorphisms and domain restriction preserve that integral.

Transport a continuous zero-current function on a smooth curve to the
ordinary algebraic-model open just constructed. Compact tests pull back
and extend by zero in the original open; \(D\) and integration are
unchanged. Wanner's theorem applies with all its hypotheses, so both
signs are Thuillier-subharmonic. They are harmonic by the Dirichlet
comparison criterion: on a strict affinoid domain the unique harmonic
solution with the same boundary values dominates the function by
subharmonicity; applying the same comparison to its negative gives
equality. This is GR 3.12, with Thuillier Lemma 2.3.9; its proof body
appears below. GR 3.4.2 identifies finite Thuillier-subharmonicity with
GR-subharmonicity on smooth strict curves, and GR 3.5(1) gives locality.
Thus the function is GR-harmonic on the entire smooth curve.

In this finite two-sign situation one can also check the universal part
of GR Definition 3.4 directly, avoiding any inference of current
regularization: the Dirichlet argument first makes the function
Thuillier-harmonic. Thuillier Propositions 2.3.18 and 2.3.19, with their
statements and proofs included below, preserve already harmonic
functions under every non-Archimedean field extension and every morphism
between smooth strict curves. Apply them after scalar extension and then
to each connected smooth strict curve probe. A harmonic function and its
negative are subharmonic: their graph maximum principle on each strict
affinoid gives comparison with every harmonic boundary majorant. Thus
both signs satisfy every test in Definition 3.4. This explains exactly
why Thuillier recognition on the local model yields GR recognition with
the field quantifiers; a zero-current equation itself is never used to
pull back to these probes.

For a reduced singular curve, C3 transports the intrinsic equation to
its finite normalization. Since \(k\) is perfect, the normalization is
rig-smooth (GR Lemma 3.9(1)); boundarylessness makes it smooth. The
previous paragraph recognizes the lifted function. GR 3.14(7) descends
the \textbf{recognized} harmonic function through this finite surjective
curve map. For a nonreduced curve C3 first gives the equation on its
reduction; the reduced result recognizes it there, and GR 3.14(7)
descends through the finite surjective nilpotent closed immersion.
Neither descent divides out a current nor suppresses original nilpotent
weights.

The exact published function-theory inputs used here, with their proofs
in Appendix S, are: GR Definition 3.13 (both signs GR-subharmonic),
3.14(1,6,7) (sheafiness, scalar equivalence for every analytic
extension, pullback under curve maps and converse for finite surjective
maps), 4.1.2 (on curves classical psh equals GR-subharmonicity), and
3.15 (on \textbf{any strict analytic curve} over a nontrivially valued
field, harmonic is equivalent to real PL and all residual classes
numerically trivial). The latter uses GJR Definition 7.2, Remark 7.3(1),
Proposition 15.7, and scalar/proper semipositivity descent; their
relevant statements and proof bodies are also supplied. These are exact
external function-theory theorems, not consequences deduced from a
current pullback. They permit nonreduced curves. This completes
algebraically closed intrinsic recognition.

There is a source-numbering distinction in GR 3.15's printed proof: it
cites GJR's earlier Lemma 5.3(3) and Corollary 5.11. In the cached
arXiv:2111.05741v3 (7 March 2025), the matching complete statements are
\textbf{Lemma 5.4(3)} (PL equivalence under the reduction immersion) and
\textbf{Corollary 5.12} (PL equivalence under completion of an algebraic
field extension, for divisible coefficients). Their statements and
proofs are included below together with Proposition 5.11, on finite flat
surjective PL descent. This is an explicit match of statements, not
identification of different numbers by assumption. GR's bibliography
also prints 2111.05747; the title and authors match the supplied primary
PDF, whose actual identifier is 2111.05741v3. The numerical-triviality
part uses GJR Proposition 7.6(1',2',3') on proper, scalar, and reduction
functoriality of both signs. Its hypotheses are good strict spaces,
already real-PL functions, and the stated proper or scalar maps; these
hold after real-PL descent. The rig-smooth comparison is Proposition
15.7: on a strict affinoid semistable curve, the weighted sum of
outgoing slopes equals the degree of the model residue line bundle, so
harmonicity is equivalent to numerical triviality. Its complete supplied
proof reduces real slopes to integral ones using a rational basis of
that slope equation. This matches precisely the real coefficient
convention used here.

\section{Part B. The completed algebraic closure and full
descent}\label{part-b.-the-completed-algebraic-closure-and-full-descent}

Equip an algebraic closure \(\overline K\) with the extending absolute
value. It exists consistently on finite subextensions: existence and
uniqueness are proved in T3, and uniqueness makes restrictions
compatible. Let \(C=\widehat{\overline K}\). It is a complete valued
field containing \(K\) isometrically, so its valuation is nontrivial. It
is \textbf{algebraically closed}, not just perfect; the following
argument proves this without assuming spherical completeness or
separability.

For a monic polynomial \(F\in C[T]\) of degree \(d>0\), approximate its
coefficients by coefficients in \(\overline K\), obtaining monic \(F_m\)
with errors tending to zero arbitrarily fast. All coefficients are
bounded, so every root of every \(F_m\) has norm at most one fixed
\(B\geq1\): if a root had norm larger than both 1 and all coefficient
norms, the leading term would dominate every other term. Arrange
successive coefficient differences so that
\(\max_i|a_{i,m+1}-a_{i,m}|B^{d-1}\leq\epsilon_m^d\), with
\(\epsilon_m\downarrow0\). Choose a root \(r_1\) of \(F_1\) in
\(\overline K\). After choosing \(r_m\), evaluate \(F_{m+1}(r_m)\); its
norm is at most \(\epsilon_m^d\). Factoring
\(F_{m+1}=\prod_{j=1}^d(T-s_j)\) in \(\overline K\) shows that some
\(s_j\) has \(|r_m-s_j|\leq\epsilon_m\). Set it to \(r_{m+1}\).
Ultrametricity makes this a Cauchy sequence. Its limit \(r\in C\)
satisfies \(F(r)=0\), by continuity and uniform coefficient convergence
on the bounded ball. This holds for every monic polynomial. Divide by
the linear factor and induct on degree to conclude algebraic closedness.
The argument works also for inseparable polynomials.

Given R's \(U,h\), apply T to \(M=\overline K\). T12.5 checks that the
actual \(U_C\) is good separated strict boundaryless and locally pure of
dimension one, not merely of total dimension one. Its actual nilpotent
lengths are used in T13. Thus \(D[h_C]=0\) unconditionally. Part C
recognizes \(h_C\) as GR-harmonic. GR Proposition 3.14(6) then gives
GR-harmonicity of \(h\). Apply 4.1.2 to both signs for classical
pluriharmonicity. Apply 3.15 over the original \(K\) for G-local real PL
and GJR harmonicity.

For clarity about the full quantifiers in this descent, let \(L/K\) be
\textbf{any} complete analytic extension, possibly transcendental. Embed
\(\overline K\) into an algebraic closure of \(L\) over \(K\); extend
partial algebraic embeddings by adjoining roots. The norm agrees on
finite \(K\)-subfields by uniqueness. Complete that algebraic closure to
a field \(F\). The same root argument proves that \(F\) is algebraically
closed and complete; the isometric embedding extends from
\(\overline K\) to \(C\). Thus \(F\) is a joint analytic extension of
\(C\) and \(L\). The extension \(F/C\) is not asserted to be algebraic.

The converse scalar assertion in GR 3.5(7), whose complete proof is
reproduced below, uses precisely this joint-field argument. For a smooth
strict connected \(L\)-curve probe \(Y\to U_L\), pull the recognized
\(\pm h_C\) to \(F\) and to \(Y_F\). They are subharmonic. Subharmonic
descent on a smooth strict curve is checked by compact comparison: for a
strict affinoid domain \(V\subset Y\) and harmonic boundary majorant
\(b\), its scalar pullback is harmonic (Thuillier 2.3.18), the upstairs
absolute boundary maps onto \(\partial V\) (GR 2.12), and comparison
upstairs therefore implies comparison downstairs by scalar surjectivity.
Components cause no \(-\infty\) issue since our functions are finite.
GR's proof supplies the general curve version, and 4.1.2 then uses both
scalar functoriality and curve-map functoriality to cover \textbf{every}
analytic curve \(Y\), without smoothness, strictness, or nonreducedness
restrictions on the probe. Equivalently one can use GR 4.2(5) and
7.2(6); their exact proofs are in Appendix S. This is descent of a
recognized function. No equation on \(U_L\) or on a probe is presumed at
this stage. R is proved. \# Part Q. Closed-curve restriction and
ground-field probes

The statements and arguments in this part are over an arbitrary complete
nontrivially valued field. They do not use perfection. In addition to
the foundations in T and C, we use these exact CLD calculus inputs, with
complete primary bodies in Appendix S. CLD 7.2.8--12 gives coefficient
Radon measures and locally bounded Borel multiplication for positive
currents. CLD 7.2.7 and 7.3.1--3 gives positivity of integration and
preservation of positivity when wedging a strongly positive smooth form.
CLD 7.7.1--8 gives weak continuity of differentiation, smooth Leibniz,
ordinary-open restriction, and compact direct-image compatibility.
Canonical integration on a boundaryless pure space is doubly closed. CLD
5.11.1 is Green's formula on compact calibrated domains, and 5.9.5 gives
finite-domain inclusion-exclusion for boundary integration. CLD
13.7.1--2 gives purity and the associated effective divisor cycle for
the zero space of a regular analytic function. They are used on reduced
\(S_1\) component sources only.

Smooth uniform approximation of a continuous function on a \textbf{fixed
compact} set, needed in the positive-limit lemma below, follows directly
from T4. Cover the compact by finitely many ordinary opens on which its
oscillation is less than \(\epsilon\), select representative values, and
multiply them by a nonnegative compact subordinate partition of sum one
near that compact. Their sum is ordinary smooth and differs by at most
\(\epsilon\) there. Thus no noncompact fixed-index compact-support
partition theorem is used.

For Poincare--Lelong here we supply a proof with \(S_1\) already
imposed; we do not import the reduction in CLD 16.7.5 claiming that one
regular element forces an arbitrary local ring to be \(S_1\). The only
local dimension inference used from that discussion is valid: near a
zero of a regular function on a pure \(r\)-space, the fibre dimension is
\(r-1\); upper semicontinuity (Ducros, \emph{Variation de la dimension
relative}, Theorem 4.9, stated in CLD 11.4.5) and the lower bound CLD
11.5.2 imply pure relative dimension \(r-1\) after an ordinary
shrinking. Every compact domain in that neighbourhood retains this
relative dimension and \(S_1\). The exact d=1 calibrated theorem stated
in C then applies. Generic lengths are constant on analytic-domain
restrictions (CLD 13.2.6).

The arguments below reconstruct the positive cycle cuts and their local
weights. C1 already proves closed-embedding test lifting, and C2 already
proves rigid puncture removal. Numbered lemmas and equations within Q
are local to that part.

\textbf{Closed-curve restriction lemma.} Let \(X\) be good, separated,
strict, boundaryless, pure of dimension \(n\geq1\), let \(\Omega\) be
ordinary open, and let \(h\) be continuous finite on \(\Omega\) with
\(D(h\delta_\Omega)=0\). For every closed analytic subspace
\(C\hookrightarrow X\) pure of dimension one,
\(D(h|_{C\cap\Omega}\delta_{C\cap\Omega})=0\) intrinsically. Both
structures may be nonreduced. The same assertion holds for a curve
closed in an ordinary open of \(X\). The proof follows, keeping the
component lengths at every stage. \#\# Elementary reducedness and prime
avoidance

A reduced noetherian ring \(R\) satisfies \(S_1\). Indeed let
\(P_1,\ldots,P_t\) be its minimal primes. If \(u\) avoids them and
\(uv=0\), the domain \(R/P_i\) makes \(v\) vanish in every quotient, so
\(v\in\bigcap_iP_i=0\). In a positive-dimensional localization the
maximal ideal is not contained in any minimal prime; finite prime
avoidance supplies a non-zero-divisor in that maximal ideal. Its depth
is therefore at least one. Zero-dimensional localizations impose only
depth at least zero. CLD 11.3.1 carries this property to every analytic
domain of a reduced component. It does not make the \textbf{original}
nonreduced or embedded-prime ambient source \(S_1\).

Here is also the particular avoidance argument needed in the cuts. Let
\(A\) be a strict affinoid algebra, \(J=(g_1,\ldots,g_s)\) an ideal, and
\(P_1,\ldots,P_t\) primes such that \(J\not\subset P_i\) for every
\(i\). The coefficient sets \[
 L_i=\{(a_1,\ldots,a_s)\in K^s:\ \textstyle\sum a_jg_j\in P_i\}
\] are proper \(K\)-linear subspaces. A finite union of proper linear
subspaces cannot cover \(K^s\) when \(K\) is infinite: choose a nonzero
linear form vanishing on each subspace, take their product, and use that
a nonzero polynomial over an infinite field cannot vanish on all \(K^s\)
(induction on the number of variables). A nontrivially valued field is
infinite, for the powers of an element of nonunit absolute value are
distinct. Thus there is an actual \(f=\sum a_jg_j\in J\) outside all
\(P_i\). No separability, geometric point transitivity, or smoothness is
involved.

\subsection{Continuous multiplication of positive weak limits in every
degree}\label{continuous-multiplication-of-positive-weak-limits-in-every-degree}

\textbf{Lemma 1.} On a good separated space, if positive currents
\(T_m,T\) of bidimension \((p,p)\) satisfy \(T_m\to T\) weakly and \(u\)
is continuous, then \(uT_m\to uT\) weakly.

\textbf{Proof.} Fix a compact \((p,p)\) ordinary smooth test \(\omega\).
Positive currents are symmetric, so replace the test by
\((\omega+J\omega)/2\). Use finitely many ordinary moment charts near
its compact support and a nonnegative compact subordinate partition
\(\theta_i\) with sum one there. On each chart choose a Euclidean smooth
representative and symmetrize that representative before estimating it.
With moment coordinates \(s_1,\ldots,s_b\) put
\(Q=\sum_jd's_j\wedge d''s_j\). The matrix of \(Q^p\) on
\(\Lambda^p\mathbb R^b\) is \(p!\) times the identity, with the
stipulated CLD coefficient sign. A symmetric representative of
\(\omega\) has a symmetric coefficient matrix \(M(s)\), bounded in
operator norm on the compact moment image of
\(\operatorname{supp}\theta_i\). Choose \(c_ip!\geq\sup\|M(s)\|\). Then
\(c_iQ^p\pm\omega\) are middle-positive wherever \(\theta_i\) is
nonzero. If \(p=0\), use \(Q^0=1\) and the scalar bound; if \(p>b\) the
represented test is zero.

Extend \(\beta_i=\theta_ic_iQ^p\) by zero after pullback, and put
\(\beta=\sum_i\beta_i\). It is compact and positive, and
\(\beta\pm\omega\) are positive. For any positive \(S\), its coefficient
measures consequently satisfy \[
 |\mu_{S,\omega}|\leq\mu_{S,\beta}.
\] The masses \(T_m(\beta)\) are bounded because they converge to
\(T(\beta)\). Choose an ordinary smooth \(v\) approximating \(u\)
uniformly within \(\varepsilon\) on \(\operatorname{supp}\beta\). For
fixed \(v\), weak convergence gives \(T_m(v\omega)\to T(v\omega)\). The
respective errors on replacing \(v\) by \(u\) are at most
\(\varepsilon T_m(\beta)\) and \(\varepsilon T(\beta)\). Let \(m\) tend
to infinity and then \(\varepsilon\) tend to zero. Antisymmetric tests
vanish before and after scalar multiplication. This proves the claim. In
particular, for a \textbf{fixed} positive current and a compact test,
uniform scalar errors are bounded by a finite coefficient-measure mass.
This bound alone would not make a differentiated pairing zero; our
subsequent cuts also use a closed smooth wedge and its ambient equation.
\(\square\)

\subsection{Poincare--Lelong on an S1
source}\label{poincarelelong-on-an-s1-source}

\textbf{Lemma 3.} Let \(M\) be good, separated, boundaryless, pure of
dimension \(r\geq1\), and satisfying \(S_1\). Let \(f\) be analytic and
stalkwise a non-zero-divisor. With the log current extended from the
unit locus, \[
 D[\log|f|]=\delta_{\operatorname{div}_M(f)}.              \tag{3}
\]

\textbf{Proof.} The actual zero fibre is pure of dimension \(r-1\).
Compact top tests have support disjoint from it, so the log current is
unambiguously defined by integration on the unit locus (CLD 15.2.1(b),
16.7.2). Outside the zero set the unit logarithm is ordinary smooth and
its Hessian is zero. Near each zero, the relative-dimension foundations
stated above supply an ordinary neighbourhood on which
\(f:M\to\mathbb A^1_K\) is pure of relative dimension \(r-1\). The
already imposed \(S_1\) condition persists on all compact analytic
domains. A compact partition reduces the formula to such neighbourhoods.

Fix a symmetric compact \((r-1,r-1)\) test \(\omega\) in one of them.
Choose finitely many compact analytic domains \(V_i\) tropicalizing it,
whose ordinary interiors cover its support, and write
\(V=\bigcup_iV_i\). For each nonempty index set \(E\), the compact
analytic domain \(V_E=\bigcap_{i\in E}V_i\) is pure and \(S_1\). Its
fibres are analytic domains in the original fibres, so restriction
preserves the pure relative dimension. Take a unit moment \(b_E\) on
\(V_E\) and a superform \(\alpha_E\) representing \(\omega\). CLD 14.7.2
with \(d_{\rm target}=1\) gives, for a common sufficiently small
\(R>0\), \[
 (b_E,f)_{\rm trop}(V_E\setminus V_{E,0})^{(r)}
     \big|_{\log|f|<\log R}
 =P_E\times(-\infty,\log R),                              \tag{4}
\] including product calibration. Here
\(P_E=b_{E,\rm trop}(V_{E,0})^{(r-1)}\) has the \textbf{actual
zero-fibre} calibration, and the last factor has standard primitive
calibration. Empty domains contribute zero.

Choose a smooth \(u_1\) equal to one sufficiently far to the left and
zero for \(s\geq\log t\), with \(0<t<R\), and put \(u_2=1-u_1\). Extend
\(u_1(\log|f|)\) by one at zero, and \(u_2\) by zero there. Write
\(\omega_1=u_1\omega\) and \(\omega_2=u_2\omega\). For \(\omega_2\) both
proposed pairings vanish, as its support lies in the unit locus. The
compact support of \(d''\omega_1\), of maximum degree \(r\), avoids the
zero fibre. Choose \(0<s<t\) smaller than the minimum of \(|f|\) on that
support and with \(u_1=1\) near \(\log s\). On
\(V_{\geq s}=V\cap\{|f|\geq s\}\) the logarithm is a smooth unit
logarithm. Green's formula gives \[
 \langle D[\log|f|],\omega_1\rangle
 =\int_{V_{\geq s}}\log|f|D\omega_1
 =-\int_{V_{\geq s}}^\partial d''\log|f|\wedge\omega_1.     \tag{5}
\] The interior Hessian term is zero. The other boundary term contains
\(d''\omega_1\) and is zero: its support lies in the original interior
of \(V\) and away from the new inner boundary.

Boundary inclusion-exclusion reduces (5) to each \(V_{E,\geq s}\).
Because \(u_1\) is zero near \(t\), compute its boundary integral on
\(V_E\cap\{s\leq|f|\leq t\}\). Equation (4) identifies its top
calibrated image with \(P_E\times[\sigma,\tau]\), where
\(\sigma=\log s\), \(\tau=\log t\). The inner face has negative fibre
calibration and contributes \(-\int_{P_E}\alpha_E\), because
\(u_1(\sigma)=1\). The outer face contributes zero. A side face is a
fibre face of dimension at most \(r-2\) times the interval. Contracting
the last-coordinate \(d''\) leaves an \((r-1)\)-form coming from that
fibre face, so its contribution is zero by dimension. For \(r=1\) there
are no such side faces. Thus the integral is \(-\int_{V_{E,0}}\omega\).
Ordinary inclusion-exclusion on the zero fibre proves \[
 \langle D[\log|f|],\omega_1\rangle
 =\int_{V_0}\omega
 =\langle\delta_{\operatorname{div}(f)},\omega_1\rangle.
\] Associated-cycle integration is the exact analytic/calculus
foundations stated in T, C, and Q; the zero quotient itself need not be
reduced or \(S_1\). Symmetry suffices for all tests. This proves (3),
without the excluded regular-element-to-\(S_1\) sentence. \(\square\)

\subsection{Cutting positive cycles while preserving the ambient
equation}\label{cutting-positive-cycles-while-preserving-the-ambient-equation}

\textbf{Lemma 4.} Let \(U\) be good, separated, strict and boundaryless.
Let \[
 S=\sum_i m_i(i_i)_*\delta_{Z_i},\qquad m_i>0,
\] be a finite effective integration cycle of dimension \(r\geq2\),
where the \(Z_i\) are reduced pure closed analytic subspaces of \(U\).
Suppose \(h\) is continuous and \(D(hS)=0\). If \(f\) is an ambient
analytic function nonzero on every \(Z_i\) component, let \[
 S_f=\sum_i m_i(i_i)_*\delta_{\operatorname{div}_{Z_i}(f|_{Z_i})}.
\] Then \(S_f\) is an effective cycle of pure dimension \(r-1\) and \[
 D(hS_f)=0.                                               \tag{6}
\] It is the stated sum of component divisor cycles, with their actual
lengths. It is not identified with an arbitrary quotient scheme's
associated cycle.

\textbf{Proof.} Each reduced source is \(S_1\) and boundaryless, and
\(f|_{Z_i}\) is regular. To see the latter on analytic domains, a
nonzero function on a reduced irreducible component avoids its minimal
prime; local analytic rings and domains inherit regularity of that
element, by faithful flat analytification/domain restriction or the
Cartier sheaf property in CLD 13.6--7. Lemma 3 applies component by
component.

Let \(\kappa\) be a nonnegative even smooth kernel of integral one
supported in \([-1,1]\), and put \(\rho(s)=\int\max(s-t,0)\kappa(t)dt\).
This is smooth convex, zero for \(s\leq-1\), and \(s\) for \(s\geq1\).
Define \[
 w_m=\rho(\log|f|+m)-m
\] on the unit locus, and set it equal to \(-m\) near zero. The
definitions agree on \(|f|<e^{-m-1}\), so \(w_m\) is ordinary smooth.
Its Hessian \(\alpha_m=Dw_m\) is doubly closed and strongly positive; on
the unit locus it is \[
 \rho''(\log|f|+m)d'\log|f|\wedge d''\log|f|.
\] Set \(S_m=\alpha_m\wedge S\). These are positive currents of
bidimension \((r-1,r-1)\). On each component \(S_m=D(w_mS)\) by
boundarylessness. For every compact top test on that component, its
support avoids the zero fibre, and eventually \(w_m=\log|f|\) there.
Hence \(w_m\delta_{Z_i}\to[\log|f|]_{Z_i}\) weakly. Lemma 3 and weak
continuity of differentiation give \(S_m\to S_f\) after summing the
finitely many components and using closed-embedding direct image.

Coefficient measures give \(hS_m=\alpha_m\wedge(hS)\). The legitimate
smooth product rule, using both closed differentials of \(\alpha_m\),
gives \(D(hS_m)=\alpha_m\wedge D(hS)=0\). Lemma 1, in degree \(p=r-1\),
now gives \(hS_m\to hS_f\). Differentiate the limit to obtain (6). The
last cut has \(r=2\), so its positive limit has bidimension \((1,1)\)
and its differentiated tests are \((0,0)\); no negative degree occurs.
\(\square\)

\subsection{Proper effective cuts containing a prescribed analytic
curve}\label{proper-effective-cuts-containing-a-prescribed-analytic-curve}

Fix \(x\in C\cap\Omega\). Choose a strict affinoid neighbourhood
\(Q=\mathcal M(A)\subset\Omega\) with \(x\) in
\(U=\operatorname{Int}(Q/X)\). This strict neighbourhood exists by the
exact analytic/calculus foundations stated in T, C, and Q applied inside
\(\Omega\). Every component of \(Q\) has dimension \(n\), with the
original ambient generic length, and \(C_Q=C\cap Q\) is a closed pure
one-dimensional strict affinoid subspace. If it is empty there is
nothing to prove.

Let \(Z_i\) be the reduced components of \(Q\) and
\(m_i=\operatorname{length}(A_{P_i})>0\) their actual generic lengths.
On \(U\), \[
 \delta_\Omega|_U=S_0:=\sum_i m_i\delta_{Z_i\cap U},
 \qquad D(hS_0)=0.                                       \tag{10}
\] This is associated-cycle integration and analytic-domain
compatibility. Embedded primes do not add support components and are not
assigned artificial multiplicity one. At subsequent stages only the
constructed effective cycles are used.

Let \(J=\mathcal I_{C_Q}(Q)\subset A\). Suppose \(S_j\) is a finite
effective \textbf{affinoid} cycle of pure dimension \(d=n-j>1\),
restricted to \(U\), whose support on \(Q\) contains the support of
\(C_Q\). If \(P\) is the prime defining a current reduced component of
dimension \(d\), \(J\subset P\) would imply that this entire component
is a closed subspace of the one-dimensional \(C_Q\), contradicting
\(d>1\). Thus choose \(f_{j+1}\in J\) nonzero on every component of
\(S_j\), by the explicit finite avoidance argument. Define the next
cycle on \(Q\) by the weighted sum of its component divisor cycles and
their reduced irreducible decompositions. CLD 13.7.1 and the exact
analytic/calculus foundations stated in T, C, and Q make \textbf{every}
resulting component dimension \(d-1\), with positive integer
coefficient. There are finitely many on the affinoid. Since \(f_{j+1}\)
vanishes on \(C_Q\), every point of \(C_Q\) in the support of \(S_j\) is
in the support of the next divisor cycle. This preserves the
containment. We do not assert a scheme embedding of \(C_Q\) in each
individual reduced component.

Lemma 4 on the boundaryless ordinary open \(U\) gives \(D(hS_{j+1})=0\)
there. The functions are analytic on all \(Q\) and hence on \(U\); each
reduced component intersected with \(U\) is a closed boundaryless
subspace of \(U\). Domains disjoint from \(U\) can stay in the affinoid
cycle and contribute zero on this open. At the next step avoidance still
treats all affinoid components. No operation uses tests meeting the
boundary of \(Q\).

After exactly \(n-1\) cuts, or none if \(n=1\), we have \[
 S_{n-1}=\sum_{k=1}^t b_k\delta_{D_k},\quad b_k\in\mathbb Z_{>0},
 \qquad D(hS_{n-1})=0\quad\text{on }U,                    \tag{11}
\] where the \(D_k\) are distinct reduced pure one-dimensional closed
affinoid subspaces of \(Q\), and their union contains \(C_Q\). This is a
constructed cycle. An arbitrary iterated quotient scheme might have
embedded primes; neither its \(S_1\) property nor equality of its
associated cycle with (11) is assumed.

Every reduced irreducible component of \(C_Q\) equals some \(D_k\).
Indeed it is an irreducible closed subset contained in the finite union
of the \(D_k\), hence contained in one of them. A proper closed subset
of an irreducible analytic curve has dimension less than one, whereas
this component has dimension one. Equality follows, with the reduced
structures.

Distinct \(D_k\) intersect in a zero-dimensional closed strict affinoid,
since a proper closed subset of an irreducible curve has dimension zero.
Thus \[
 F_Q=\bigcup_{k\ne l}(D_k\cap D_l)
\] is a finite set of rigid points. On the ordinary ambient open
\(U\setminus\bigcup_{l\ne k}D_l\), only \(b_k\delta_{D_k}\) remains in
(11). Its closed embedding into that open, coefficient-measure
multiplication, and compact direct-image differentiation give \[
 0=b_k(i_k)_*D(h|_{D_k\cap U}\delta_{D_k\cap U})
 \quad\text{off }F_Q.
\] C1 in bidegree \((0,0)\) and division by the \textbf{actual} positive
\(b_k\) give zero intrinsic Hessian on \(D_k\cap U\setminus F_Q\). This
is an equation on an actual ordinary open, not its extension by density.
the puncture-removal lemma C2 removes the finite rigid points, proving
it on all \(D_k\cap U\).

For the original possibly nonreduced \(C_Q\), its canonical current is
the sum of the reduced component currents weighted by its own actual
generic lengths, by CLD 16.6.3. Every one of those multipliers has zero
intrinsic Hessian, so differentiation and closed-embedding direct image
give (A) on \(C\cap U\). The opens \(C\cap U\) constructed for
\textbf{every analytic point} \(x\) cover \(C\cap\Omega\). Intrinsic
currents are local; (A) follows everywhere. This proves the closed-curve
restriction lemma in the stated arbitrary-field scope. \(\square\)

\subsection{Q-local images. Exact hypotheses and
statement}\label{q-local-images.-exact-hypotheses-and-statement}

CLD 11.4.4 says finite analytic maps to an affinoid are equivalent to
finite Banach algebra maps; their finite-module quotient topologies are
determined by the ring module structure. For a nontrivially valued
ground field and a strict affinoid base, a Banach algebra finite as a
ring algebra is finite Banach. CLD 11.4.5 says the locus of relative
fibre dimension at least a fixed integer is Zariski closed, and a
morphism is finite at a point exactly when its local relative dimension
is zero and it has no relative boundary there. Finite at a point means
finite from an analytic neighbourhood to an analytic neighbourhood.
These are pointwise criteria; the different entire fibre criterion with
a compactness hypothesis, CLD 11.6.6, is not used. The exact statements
and Berkovich's local-finiteness proof bodies appear below. The boundary
equality of C shows no relative boundary for a smooth curve map to a
separated boundaryless target.

\textbf{Local-image lemma.} For an analytic map \(f:Y\to X\) from a
separated smooth strict curve to a good separated strict boundaryless
analytic space, near every analytic source point either the map is
constant with value a rigid point, through a finite field over \(K\), or
there are actual strict affinoid neighbourhoods \(W,V\) on which the map
is finite, with a reduced closed pure curve image \(C\hookrightarrow V\)
and \(W\to C\) finite surjective. The source neighbourhood can be chosen
connected and normal and the map then lifts finitely to the
normalization of \(C\). No perfection, separability, or finite flatness
assumption occurs. The following proves the statement.

\subsection{The affinoid branch lemma, including Banach
factorization}\label{the-affinoid-branch-lemma-including-banach-factorization}

\textbf{Image branch lemma.} Let \(K\) be complete nontrivially valued,
\(\mathcal M(B)\) a connected normal strict affinoid pure curve, and
\(f:\mathcal M(B)\to\mathcal M(A)\) an analytic map to a strict
affinoid. Either its bounded algebra map factors through a finite field
\(E/K\) and a rigid point, or \(\dim_y f=0\) at every analytic source
point.

\textbf{Proof.} As in C4, connected normality makes \(B\) an integral
domain, of Krull dimension one. Fix an actual bounded Tate presentation
\(K\langle T_1,\ldots,T_N\rangle\twoheadrightarrow A\) and write
\(b_j=f^*(T_j)\in B\).

If every \(b_j\) is algebraic over \(K\), then
\(E=K[b_1,\ldots,b_N]\subset B\) is finite-dimensional over \(K\) and a
domain, hence a finite field. Every point of \(\mathcal M(B)\) restricts
to the unique extending absolute value on \(E\). Thus
\(\rho_B|_E=|\cdot|_E\). Since \(B\) is reduced its Banach norm is
equivalent to \(\rho_B\); \(E\) is therefore a closed complete
finite-dimensional subspace of \(B\) with its usual finite-field
topology. For every Tate series, polynomial partial sums evaluated at
the \(b_j\) converge in \(B\) under the bounded map and belong to \(E\).
Closedness puts the limit in \(E\). Consequently the \textbf{entire}
bounded map, not just its polynomial subalgebra, factors through \(E\).
Boundedness to \(E\) follows from \(|e|_E=\rho_B(e)\leq C\|e\|_B\) and
the bound of the original map; descent to \(A\) uses its quotient norm.
Its image contains every \(b_j\), so \(A\to E\) is surjective. The
inclusion \(E\to B\) is bounded by the same norm equivalence. Hence
\(f\) factors analytically through the singleton \(\mathcal M(E)\),
mapping to a rigid target point. No separability is required for
\(E/K\).

Otherwise some \(b=b_j\) is transcendental over \(K\). Let \(a\) be any
rigid source point. The element \(b(a)\in\mathcal H(a)\) is algebraic
over \(K\); let \(P_a\in K[t]\) be its monic minimal polynomial. The
element \(P_a(b)\in B\) is nonzero by transcendence, vanishes at \(a\),
and is not a unit. A quotient of a noetherian one-dimensional domain by
a nonzero nonunit has dimension zero: every prime containing it is
nonzero and therefore maximal. The stated strict Nullstellensatz and the
artinian argument make \(B/(P_a(b))\) finite over \(K\), retaining its
possible nilpotents.

Put \(x=f(a)\). The isometric embedding
\(\mathcal H(x)\hookrightarrow\mathcal H(a)\) makes \(x\) rigid, with
finite residue field over \(K\). In the actual analytic fibre
\(\mathcal M(B)\times_{\mathcal M(A)}\mathcal M(\mathcal H(x))\), the
coordinate \(b\) equals the evaluated target coordinate. Its polynomial
\(P_a\) is zero there, since it is zero after the injective
residue-field map into \(\mathcal H(a)\). The universal property of a
closed zero subspace therefore factors the fibre as a closed subspace of
\[
                 \mathcal M(B/(P_a(b)))_{\mathcal H(x)}.
\] This last space is finite over \(\mathcal H(x)\) and
zero-dimensional. Equivalently, the completed fibre algebra is the
quotient of the scalar-extended finite zero algebra by the target-value
relations. Thus \(\dim_a f=0\). This argument compares fibres over their
actual field and applies also when that field extension is inseparable.

The locus \(F=\{y:\dim_y f\geq1\}\) is Zariski closed by the local
finiteness foundations stated here. On a strict affinoid any nonempty
such locus is the zero set of a proper ideal and has a rigid point by
the affinoid foundations stated in C and T. Such a point would
contradict the preceding calculation. Hence \(F\) is empty and
\(\dim_y f=0\) for \textbf{every} analytic \(y\). This uses Zariski
upper semicontinuity, not an extension by ordinary-topological density.
\(\square\)

\subsection{Proof of the local image
theorem}\label{proof-of-the-local-image-theorem}

\textbf{Proof of the local-image lemma.} Fix \(y\in Y\) and \(x=f(y)\).
Choose a strict affinoid neighbourhood \(U_1\) of \(x\). Inside the
ordinary open \(f^{-1}(\operatorname{int}_X U_1)\) choose a strict
affinoid neighbourhood of \(y\) and then its connected component \(Q\)
containing \(y\). Smoothness makes its algebra normal and domain
invariance makes it pure one-dimensional, so the image branch lemma
applies.

In its constant case, the bounded factorization through \(E\) restricts
to an ordinary neighbourhood of \(y\) and proves case 1. In the other
case the local relative dimension of the original \(f\) at \(y\) is
zero: the restriction fibre is an analytic domain neighbourhood in the
original fibre, and domain restriction preserves local dimension. Since
the target \(X\) is separated, the local finiteness foundations stated
here gives \[
                    \partial Y
                     =\partial(Y/X)\cup f^{-1}(\partial X).
\] The source is boundaryless, so \(y\) is in the relative interior of
\(f\). The finite-at-a-point criterion supplies analytic neighbourhoods
\(V'\to U'\) on which \(f\) is finite. This step establishes no
compactness of the original entire fibre and uses no such compactness.

Take a strict affinoid neighbourhood \(U\) of \(x\) inside an ordinary
neighbourhood contained in \(U'\). Its inverse image \(W\) in \(V'\) is
affinoid and finite over \(U\). It contains an ordinary neighbourhood of
\(y\) in \(Y\): intersect an ordinary neighbourhood contained in \(V'\)
with \(f^{-1}(\operatorname{int}_X U)\). Select the connected component
\(V\) of \(W\) containing \(y\). It is open and closed in the affinoid,
still a neighbourhood of \(y\), and still finite over \(U\). As a source
domain it is regular normal and pure one-dimensional, hence its algebra
\(B\) is a domain.

For completeness, the finite affinoid over the strict
\(A=\mathcal O(U)\) is strict here. Choose finitely many module
generators of its reduced algebra \(B\). Multiply each by a suitable
nonzero scalar from \(K\) so its spectral radius is at most one. Their
powers are bounded by the equivalence of reduced spectral and Banach
norms. They and the target Tate generators give a bounded surjection
from a strict Tate algebra onto \(B\). Its quotient topology is
equivalent to the finite-module topology: expressions as sums of the
chosen module generators bound the quotient norm by a fixed multiple of
the finite-module quotient norm, and the bounded surjection gives the
reverse bound. the local finiteness foundations stated here identifies
this finite-module topology with the existing Banach topology. Thus
\(V\) is strict. The same argument applies to the finite reduced
normalization algebra used below.

Let \(I=\ker(A\to B)\) and \(A_0=A/I\). It is a strict affinoid domain
algebra, with an injective bounded map to \(B\), and \(B\) is finite
integral over \(A_0\). The closed subspace
\(C=\mathcal M(A_0)\hookrightarrow U\) is reduced. Injective integral
extensions preserve Krull dimension, by going-up and incomparability, so
\(\dim A_0=\dim B=1\). The prime kernel and the affinoid foundations
stated in C and T make \(C\) irreducible and pure one-dimensional.

To check surjectivity at \textbf{all analytic points}, let \(c\in C\)
have kernel \(\mathfrak p\subset A_0\). Lying-over for the integral
injection gives a prime of \(B\) over \(\mathfrak p\). Thus
\(B\otimes_{A_0}\kappa(\mathfrak p)\) is a nonzero finite-dimensional
algebra, and extending its field to \(\mathcal H(c)\) leaves it nonzero.
Finite-module completion identifies this algebra with the finite
analytic fibre algebra \(B\widehat\otimes_{A_0}\mathcal H(c)\) (the
local finiteness foundations stated here); finite-dimensional vector
spaces over a complete field are already complete. A maximal ideal gives
a finite residue field extension of \(\mathcal H(c)\) with its extending
absolute value, hence an analytic point of the fibre. This proves that
\(V\to C\) is finite surjective. It is not an inference merely from
surjectivity of the two separate maps to algebra spectra.

Let \(\overline A_0\) be the finite integral closure of \(A_0\) in its
fraction field. It exists by the affinoid foundations stated in C and T.
The injection into the domain \(B\) embeds \(\operatorname{Frac}(A_0)\)
in \(\operatorname{Frac}(B)\). An element of \(\overline A_0\) is
integral over \(A_0\), hence integral over \(B\), and belongs to \(B\)
by normality. This gives the algebraic lift \(\overline A_0\to B\). It
is bounded: choose finitely many \(A_0\)-module generators for
\(\overline A_0\); boundedness of \(A_0\to B\) and the norms of their
images bound the module map from a free presentation, and hence its
quotient norm. the local finiteness foundations stated here gives the
finite Banach topology on \(\overline A_0\). Further, the same
\(A_0\)-module generators of \(B\) generate \(B\) over
\(\overline A_0\), so the lift is finite as a ring map; the strict
nontrivial finite-Banach assertion of the local finiteness foundations
stated here makes it finite analytic. This is the asserted lift to the
normalization. It does not require a separable map, a geometrically
reduced image, or finite flatness. \(\square\)

\subsection{Q-conclusion. The probe consequence and its
limits}\label{q-conclusion.-the-probe-consequence-and-its-limits}

At an arbitrary source point \(y\) of the smooth strict curve \(Y\) in
Theorem P, the local-image result just proved has two cases. In the
rigid constant case \(h\circ f\) is locally constant. Otherwise it
supplies an actual strict affinoid target neighbourhood \(V\) of
\(x=f(y)\), a reduced closed pure curve \(C\hookrightarrow V\), and a
finite map onto that image from a source affinoid neighbourhood. Put
\(V^o=\operatorname{int}_U V\). It is an ordinary neighbourhood of \(x\)
and is good separated strict boundaryless pure of dimension \(n\). Its
closed subspace \(C^o=C\cap V^o\) is a good separated strict
boundaryless pure curve. The current equation on \(U\) restricts to
\(V^o\), with compact ordinary tests extended by zero. The closed-curve
restriction proof above supplies the intrinsic actual equation
\(D[h|_{C^o}]=0\).

Apply R to \(C^o\) over this same possibly imperfect field. It
recognizes \(h|_{C^o}\) as GR-harmonic, with no scalar stability or
smoothness assumption on the image. On the ordinary source neighbourhood
\(\operatorname{int}_Y W\cap f^{-1}(V^o)\), the actual map factors
through \(C^o\). GR 3.14(7) pulls back the \textbf{recognized} harmonic
function. Consequently \(h\circ f\) is harmonic near \(y\). The source
is smooth strict, so GR 3.4.2 and 3.12 identify this with Thuillier
harmonicity. Sheafiness treats every analytic source point and proves P.
GR 4.1.2 and 3.15 then give the stated classical and real-PL/GJR
conclusions on that curve.

The local-image proof also constructs the finite lift to the
normalization. Over an imperfect field its normalization is not asserted
smooth; R makes that unnecessary. The restriction proof has already
checked component weights, singular intersections, rigid punctures, and
nilpotents before recognition. No finite-current direct image is
inverted, no graph-product current is assumed, and no admissible cover
replaces ordinary interiors.

\section{Scope ledger, failed routes, and
reproduction}\label{scope-ledger-failed-routes-and-reproduction}

R covers every field and every pure analytic curve in the original
strict boundaryless scope. The completed-algebraic current gate is
proved in T; recognition and full function-theory scalar descent are
separately proved in C and B. The result extends the algebraically
closed recognition evidence and the catalogue's algebraically closed
smooth-proper-algebraic comparison, without imposing scalar stability on
the input function. It retains the ordinary CLD equation, actual
nilpotent lengths, purity, and every ordinary open. P is an additional
ground-field consequence on a pure ambient space in any dimension. The
source in P is a \textbf{smooth strict} curve, matching the reduced
testing class GR 4.1.3. No claim about arbitrary non-strict ground-field
source curves is added.

Neither statement solves original P06.2 in dimension at least two. On
such an ambient space, complementary Hessian tests have positive
bidegree; the finite separable proof T5--T9 approximates scalar tests on
curves only. Exact compact scalar-test descent after completion cannot
fill that missing finite differential-form gate. P tests over the actual
ground field and therefore does not cover probes into \(X_L\) for
arbitrary transcendental \(L/K\). The universal quantifiers of R apply
to \textbf{curve targets}, not to the higher-dimensional ambient of P.
No counterexample in the original strict ambient scope is constructed.
The non-strict invisible-boundary polycylinder is excluded by the
original strict/boundaryless hypotheses. P05.2 is not addressed.

The optional assertion \(D[h_L]=0\) under \textbf{every} analytic
extension on curves is not claimed here. GR harmonicity and classical
pluriharmonicity of \(h_L\) do follow for every extension. However, this
submission has not reconstructed a harmonic-to-\textbf{ordinary analytic
CLD} zero-current theorem against every ordinary analytic test in that
upstairs scope. Wanner's reverse comparison on algebraic test charts
alone would not prove that larger test statement. The proved ordinary
current transport here remains the completed-algebraic theorem T. This
explicitly separates a remaining optional implication from the proved
recognition theorem.

No two-sign ordinary smooth decreasing psh regularization is proved for
arbitrary-field arbitrary analytic curves or higher-dimensional targets.
Weak smoothness, real PL, and GJR harmonicity do not give ordinary
smooth harmonic approximants. The known nonapproximation phenomenon for
harmonic approximants is unaffected. The convex cutoffs in Q regularize
a divisor logarithm; they are not positive approximants of \(h\) or
\(-h\). The mollification in T controls a fixed differentiated test; it
is not positivity-preserving regularization of the input function.

The authoritative file is \texttt{P06.2-argument.md}. It contains all
reconstructed proofs and the essential primary bodies below.
\texttt{P06.2-partial.tex} is a standalone typesetting of it;
\texttt{P06.2-partial.pdf} is built twice. The exact five-key
\texttt{status.json} records the partial scope and the concrete gaps.
\texttt{self-audit.md} and \texttt{research-notes.md} record the
fallible author audit and discarded routes.
\texttt{source-manifest.json}, \texttt{evidence/source-checks.json}, and
\texttt{evidence/commands.json} record input hashes, newline extraction
checks, and actual commands. No old artifact or research database is
modified.

Rerun in this fresh directory with \texttt{python3\ prepare.py},
followed by \texttt{python3\ build-and-check.py}. The programs use
installed local tools and bounded standard-library diagnostics; no
additional model, subagent, network, installation, or resource-heavy
computation is used. No supervised compute facility was exposed, and
none was needed for these text/build operations. The checks bind
excerpts and outputs to bytes and test limited finite lemmas; they do
not establish mathematical correctness, verification, priority, or the
full atom. Native notes are returned through the required response
schema; no local alternate notes file is treated as canonical memory.

\section{Appendix S. Primary bodies, extracted by newline
boundaries}\label{appendix-s.-primary-bodies-extracted-by-newline-boundaries}

The following are primary source data. Statements and complete proof
bodies are included so applications can be checked in this document.
Quotation does not assert correctness of every unused sentence. In
particular the universal moment-dimension reading of CLD 12.2.4, the
broad fixed-index compact-support reading of 15.2.9, and the unsupported
S1 reduction in 16.7.5 are unused. Form-feed page markers alone are
removed.

\subsection{CLD.txt, newline lines
14838--15060}\label{cld.txt-newline-lines-1483815060}

\begin{Shaded}
\begin{Highlighting}[]

\NormalTok{ § 11.1. GÉOMÉTRIE ANALYTIQUE : CONVENTIONS, NOTATIONS,}
\NormalTok{                       PREMIERS RAPPELS}

\NormalTok{   On fixe pour toute la suite de l’article un corps ultramétrique complet k ; sa valeur}
\NormalTok{absolue peut être triviale. On note k ◦ son anneau de valuation. Nous désignerons par ek}
\NormalTok{                                                         1}
\NormalTok{le corps gradué résiduel de k (10.8.1). Sa composante k de degré 1 est le corps résiduel}
\NormalTok{                                                        e}
\NormalTok{standard de k, c’est{-}à{-}dire le corps résiduel de l’anneau local k ◦ .}
\NormalTok{   Dans ce texte, la notion d’espace k{-}analytique sera à prendre au sens de Berkovich,}
\NormalTok{et plus précisément de [9]. Rappelons qu’un polyrayon est une famille (le plus souvent}
\NormalTok{finie) de nombres réels strictement positifs ; la longueur d’un polyrayon est le cardinal}
\NormalTok{de son ensemble d’indices. Faisons quelques rappels pour la commodité du lecteur.}
\NormalTok{11.1.1. — La construction fondamentale sur laquelle repose la géométrie de Berkovich}
\NormalTok{est la suivante. Soit A un anneau de Banach, c’est{-}à{-}dire un anneau A muni d’une norme}
\NormalTok{sous{-}multiplicative ∥ · ∥ (on inclut dans la définition le fait que ∥1∥ ⩽ 1, ce qui entraîne}
\NormalTok{que ∥1∥ = 1 dès que A est non nul) pour laquelle il est complet. À un tel anneau A on}
\NormalTok{associe son spectre analytique M (A), qui est défini comme l’ensemble des semi{-}normes}
\NormalTok{multiplicatives bornées sur A, c’est{-}à{-}dire des applications φ : A → R+ vérifiant les}
\NormalTok{propriétés suivantes :}

\NormalTok{  – on a φ(0) = 0 et φ(f + g) ⩽ φ(f ) + φ(g) pour tous f, g ∈ A (propriété de semi{-}}
\NormalTok{norme) ;}
\NormalTok{  – on a φ(1) = 1 et φ(f g) = φ(f )φ(g) pour tout f, g ∈ A (multiplicativité) ;}
\NormalTok{  – il existe un nombre réel c tel que φ(f ) ⩽ c∥f ∥ pour tout f ∈ A (caractère borné de}
\NormalTok{φ).}

\NormalTok{Notons qu’on impose que φ(0) = 0 et φ(1) = 1 : si A est nul, M (A) est donc vide.}
\NormalTok{Cet ensemble M (A) est muni de la topologie la moins fine pour laquelle les applications}
\NormalTok{φ 7→ φ(f ) sont continues pour tout f ∈ A. L’application φ 7→ (φ(a))a∈A identifie topo{-}}
\NormalTok{logiquement M (A) à un fermé de l’espace produit a∈A [0, ∥a∥], si bien que M (A) est}
\NormalTok{                                                   Q}
\NormalTok{compact. Si φ appartient à M (A), son noyau est un idéal premier de A, et l’application}
\NormalTok{298                   CHAPITRE 11. RAPPELS DE GÉOMÉTRIE ANALYTIQUE}

\NormalTok{φ 7→ ker φ de M (A) vers Spec A est continue. Son image contient tous les idéaux maxi{-}}
\NormalTok{maux de A (c’est ce qui est réellement démontré dans la preuve du théorème 1.2.1 de}
\NormalTok{[8]) ; en particulier M (A) est non vide dès que A est non nul.}
\NormalTok{   En général, on note typiquement plutôt x que φ un élément de M (A), en y pensant}
\NormalTok{comme à un point. Soit x ∈ M (A). La semi{-}norme multiplicative x induit une valeur}
\NormalTok{absolue sur le corps des fractions de l’anneau intègre A / ker x, corps des fractions dont}
\NormalTok{le complété est noté H (x). L’application composée A → A / ker x → H (x) sera notée}
\NormalTok{f 7→ f (x), et la valeur absolue du corps complet H (x) sera notée |·| ; avec ces conventions}
\NormalTok{on a x(f ) = |f (x)| pour toute f ∈ A et nous utiliserons sauf exception la notation de}
\NormalTok{droite. Cette construction est fonctorielle : tout morphisme borné A → B d’anneaux de}
\NormalTok{Banach induit par composition une application continue de M (B) vers M (A).}
\NormalTok{   La semi{-}norme spectrale de A est la semi{-}norme ρ commutant aux puissances définie}
\NormalTok{par les formules}
\NormalTok{                     ρ(f ) = sup |f (x)| = lim inf ∥f n ∥1/n = inf ∥f n ∥1/n}
\NormalTok{                              x∈M (A)            n→∞                n→∞}


\NormalTok{(voir [8], 1.3, pour les deux dernières égalités). Nous désignerons par A◦ (resp. A◦◦ )}
\NormalTok{l’ensemble des éléments f de A tels que ρ(f ) ⩽ 1 (resp. ρ(f ) \textless{} 1), et par A     e l’anneau}
\NormalTok{gradué résiduel de A (10.8.1) relativement à ρ. Sa composante A     e1 est l’anneau résiduel}
\NormalTok{            ◦    ◦◦}
\NormalTok{standard A /A .}
\NormalTok{   Supposons maintenant que A est une k{-}algèbre et que sa norme est une norme de k{-}}
\NormalTok{espace vectoriel. Dans ce cas λ 7→ |λ(x)| définit pour tout x ∈ M (A) une valeur absolue}
\NormalTok{sur k bornée relativement à la valeur absolue structurale de k, ce qui entraîne que}
\NormalTok{|λ(x)| = |λ| pour tout λ : la semi{-}norme x induit donc la valeur absolue structurale de k,}
\NormalTok{et (H (x), |·|) est une extension de (k, ·). Comme la valeur absolue de k est ultramétrique,}
\NormalTok{celle de H (x) l’est aussi (une valeur absolue |·| étant ultramétrique si et seulement si}
\NormalTok{|n| ⩽ 1 pour tout n ∈ Z). On a donc dans ce cas |(f + g)(x)| ⩽ max(|f (x)|, |g(x)|) pour}
\NormalTok{tout f, g dans A, puis ρ(f + g) ⩽ max(ρ(f ), ρ(g)) (comme on le voit en faisant varier x).}
\NormalTok{11.1.2. — Soit n un entier. Pour tout polyrayon r = (r1 , . . . , rn ), nous désignerons}
\NormalTok{         −1             −1}
\NormalTok{      P 1 T1 , . . . ,νrn Tn \} la sous{-}algèbre}
\NormalTok{par k\{r}
\NormalTok{                                         ν}
\NormalTok{                                                   de k[[T1 , . . . , Tn ]] formée des séries formelles}
\NormalTok{                                                  ν1      νn}
\NormalTok{                                                                                                 P}
\NormalTok{f =      ν∈N n a ν T    telles que a ν r   = a ν 1 . . . rn converge vers 0 lorsque |ν| =}
\NormalTok{                                                r                                                  |νj |}
\NormalTok{tend vers +∞. On la munit de la norme multiplicative ultramétrique définie par ∥f ∥ =}
\NormalTok{max|aν |rν . Pour abréger les notations, on notera aussi k\{r−1 T \} cette algèbre. C’est une}
\NormalTok{algèbre de Banach dont tous les idéaux sont fermés et de type fini.}
\NormalTok{   Soit maintenant A une k{-}algèbre de Banach commutative (par convention, dans cet}
\NormalTok{ouvrage, les normes de k{-}algèbres seront toujours ultramétriques). On dit que A est une}
\NormalTok{algèbre k{-}affinoïde s’il existe un entier n, un polyrayon (r1 , . . . , rn ) et une surjection de}
\NormalTok{k\{r−1 T \} vers A telle que la norme de A soit équivalente à la norme quotient correspon{-}}
\NormalTok{dante. On dit que A est strictement k{-}affinoïde si l’on peut de plus prendre rj = 1 pour}
\NormalTok{tout j. Pour que A soit strictement k{-}affinoïde, il suffit en fait que l’on puisse prendre}
\NormalTok{les rj de torsion modulo |k × | ([20] 6.1.5/4).}
\NormalTok{   Une k{-}algèbre de Banach A est dite affinoïde s’il existe une extension valuée com{-}}
\NormalTok{plète K de k telle que A soit une algèbre K{-}affinoïde.}
\NormalTok{   Si A est une k{-}algèbre affinoïde l’application continue naturelle de M (A) vers Spec A}
\NormalTok{est surjective ([9], Proposition 2.1.1).}
\NormalTok{   Soit A une k{-}algèbre affinoïde et soit ρ sa semi{-}norme spectrale. L’algèbre A est}
\NormalTok{noethérienne, et ses idéaux sont fermés. Pour tout f ∈ A on a ρ(f ) = 0 si et seulement}
\NormalTok{     § 11.1. GÉOMÉTRIE ANALYTIQUE : CONVENTIONS, NOTATIONS, PREMIERS RAPPELS             299}

\NormalTok{si f est nilpotent. En particulier, ρ est une norme si et seulement si A est réduite et si}
\NormalTok{c’est le cas, ρ est équivalente à la norme structurale de A.}
\NormalTok{Exemple 11.1.3. — Soit r = (ri ) un polyrayon de longueur finie. Soit T P        = (Ti ) une}
\NormalTok{famille d’indéterminées et soit kr la k{-}algèbre constituée des séries formelles I∈Zn aI T I}
\NormalTok{à coefficients}
\NormalTok{           P dans    k, où |aI | · rI tend vers zéro quand |I| tend vers l’infini, munie de}
\NormalTok{la norme      aI T 7→ max|aI | · rI . La k{-}algèbre normée kr est une k{-}algèbre affinoïde :}
\NormalTok{                  I}

\NormalTok{elle est en effet isomorphe à k\{r−1 T, rS\}/(ST − 1). Sa norme est multiplicative ([48],}
\NormalTok{1.2.1). Supposons de plus que r soit k{-}libre, c’est{-}à{-}dire que les ri soient Q{-}linéairement}
\NormalTok{indépendants dans Q ⊗Z (R×           ×}
\NormalTok{                                + /|k |). Dans ce cas kr est un corps ([48], 1.2.2), et c’est}
\NormalTok{donc à la fois une extension complète de k et une algèbre k{-}affinoïde.}
\NormalTok{Remarque 11.1.4. — Soit A = k\{r−1 T \} une algèbre k{-}affinoïde, où r = (ri ) est un}
\NormalTok{polyrayon de longueur finie. Choisissons un ensemble J d’indices tel que le polyrayon}
\NormalTok{s := (rj )j∈J soit k{-}libre et quiQsoit maximal pour cette propriété. Chacun des ri est}
\NormalTok{alors de torsion modulo |k × | · j∈J riZ , si bien que la ks {-}algèbre de Banach A⊗}
\NormalTok{                                                                                 b k ks est}
\NormalTok{strictement ks {-}affinoïde.}
\NormalTok{11.1.5. — Soit A une algèbre k{-}affinoïde ; posons X = M (A ). Une partie V de X}
\NormalTok{est appelée un domaine affinoïde de X si le foncteur qui envoie une k{-}algèbre de Ba{-}}
\NormalTok{nach affinoïde B sur l’ensemble des morphismes bornés de A vers B tels que l’image}
\NormalTok{de M (B) soit contenue dans V est représentable par une algèbre k{-}affinoïde AV . Si c’est}
\NormalTok{le cas, M (AV ) ≃ V .}
\NormalTok{   Un système rationnel d’inéquations sur X est une combinaison booléenne d’inégalités}
\NormalTok{qui est de la forme}
\NormalTok{                     |f1 | ⩽ λ1 |g| et |f2 | ⩽ λ2 |g| et . . . et |fn | ⩽ λn |g|}
\NormalTok{où g et les fi sont des fonctions analytiques sur X sans zéro commun, et où les λi sont}
\NormalTok{des réels.}
\NormalTok{  Si β est un système rationnel d’inéquations sur X, on notera β̊ le système défini à partir}
\NormalTok{de β en remplaçant toutes inégalités larges par les inégalités strictes correspondantes.}
\NormalTok{  Soit}
\NormalTok{                    β : |f1 | ⩽ λ1 |g| et |f2 | ⩽ λ2 |g| et . . . et |fn | ⩽ λn |g|}
\NormalTok{un système rationnel d’inéquations sur X.}
\NormalTok{  On qualifiera de rationnel le lieu de X défini par un tel système d’inéquations. C’est}
\NormalTok{un domaine affinoïde de X, avec}
\NormalTok{                          AV ≃ k\{T1 /λ1 , . . . , Tn /λn \}/(fi − Ti g)i .}
\NormalTok{Notons que si g s’annulait en un point de V , il en irait de même de chacune des fi ; en}
\NormalTok{conséquence, g|V est inversible.}
\NormalTok{  Tout point de X possède une base de voisinages rationnels, en particulier une base}
\NormalTok{de voisinages affinoïdes. Plus précisément, pour tout x ∈ X et tout voisinage ouvert U}
\NormalTok{de x, il existe un système rationnel d’inéquations β dont le lieu de validité est contenu}
\NormalTok{dans U , et tel que x satisfasse β̊.}
\NormalTok{11.1.6. — Le théorème de Gerritzen{-}Grauert affirme que tout domaine affinoïde de X}
\NormalTok{est réunion finie de domaines rationnels (pour le cas non strict, voir [47], lemme 2.4 ;}
\NormalTok{pour une preuve directe, voir aussi [128], théorème 3.1). Il en résulte que si x est un}
\NormalTok{point de X situé sur un domaine affinoïde V de X, le corps résiduel complété H (x) ne}
\NormalTok{dépend pas du fait que l’on voie x comme appartenant à V = M (AV ) ou à X = M (A).}
\NormalTok{300                CHAPITRE 11. RAPPELS DE GÉOMÉTRIE ANALYTIQUE}

\NormalTok{11.1.7. — À partir des espaces affinoïdes, Berkovich définit dans [9] la catégorie des}
\NormalTok{espaces k{-}analytiques par un procédé de recollement un peu délicat qu’il serait trop}
\NormalTok{long de décrire ici en détail ; nous allons tout de même rappeler quelques faits. Suivant}
\NormalTok{Berkovich, si X est un espace topologique et si (Xi ) est une famille de parties de X,}
\NormalTok{nous dirons que les Xi constituent}
\NormalTok{                           S         un G{-}recouvrement de X si tout point x de X possède}
\NormalTok{un voisinage de la forme i∈J Xi où J est fini et où x ∈ Xi pour tout i ∈ J. Si l’on}
\NormalTok{ne s’astreint à recoller que des espaces strictement affinoïdes, on obtient la classe des}
\NormalTok{espaces strictement k{-}analytiques.}
\NormalTok{   Soit X un espace k{-}analytique. C’est en particulier un espace topologique dont tout}
\NormalTok{point possède une base de voisinages compacts (mais X n’est pas nécessairement lo{-}}
\NormalTok{calement compact puisqu’un espace topologique localement compact est par définition}
\NormalTok{séparé, et que ce n’est pas forcément le cas de X). L’espace X possède par ailleurs une}
\NormalTok{classe distinguée de partie compactes, ses domaines affinoïdes (si X = M (A ), les do{-}}
\NormalTok{maines affinoïdes de X sont précisément ceux décrits au 11.1.5 ci{-}dessus). Un domaine}
\NormalTok{analytique de X est une partie de X qui est G{-}recouverte par les domaines affinoïdes}
\NormalTok{qu’elle contient. Tout domaine affinoïde de X est un domaine analytique de X ; plus}
\NormalTok{généralement, tout compact de X qui est une union finie de domaines analytiques de}
\NormalTok{X est un domaine analytique. Et tout ouvert de X (et notamment X lui{-}même) est un}
\NormalTok{domaine analytique. En particulier, tout point x de X possède un voisinage compact}
\NormalTok{qui est une réunion finie de domaines affinoïdes contenant x. Les espaces définis dans [8]}
\NormalTok{sont ceux qui sont qualifiés de bons dans [9], c’est{-}à{-}dire ceux dont tout point admet un}
\NormalTok{voisinage affinoïde (et donc une base de voisinages affinoïdes.) Tout domaine analytique}
\NormalTok{de X est une partie localement fermée de X : l’intersection d’une famille finie de do{-}}
\NormalTok{maines analytiques est un domaine analytique ; la réunion d’une famille localement finie}
\NormalTok{de domaines analytiques fermés est un domaine analytique.}
\NormalTok{   Tout domaine analytique V de X hérite d’une structure canonique d’espace analy{-}}
\NormalTok{tique, et d’un morphisme V ,→ X dont l’application sous{-}jacente est l’inclusion. Tout}
\NormalTok{morphisme f : Y → X se factorisant ensemblistement par V se factorise en fait unique{-}}
\NormalTok{ment par la structure analytique de V .}
\NormalTok{11.1.8. La G{-}topologie. — Soit X un espace k{-}analytique. En plus du site topolo{-}}
\NormalTok{gique sur X on définit un second site plus fin, noté XG . Ses objets sont les domaines}
\NormalTok{analytiques de X, ses flèches sont les inclusions, et ses familles couvrantes sont les G{-}}
\NormalTok{recouvrements. La topologie de ce site est appelée la G{-}topologie de X.}
\NormalTok{  Comme tout ouvert de X est un domaine analytique et comme tout recouvrement}
\NormalTok{ouvert est un G{-}recouvrement, on dispose d’un morphisme de site canonique π : XG →}
\NormalTok{X. L’unité d’adjonction id → π∗ π ∗ est un isomorphisme ([9], p. 26).}
\NormalTok{  Le site XG est annelé. Nous noterons OX le faisceau structural de XG . Sauf mention du}
\NormalTok{contraire, lorsque nous parlerons de faisceaux cohérents sur X, il s’agira plus précisément}
\NormalTok{de OX {-}module cohérents sur XG . Précisons que le faisceau OX ainsi que sa restriction}
\NormalTok{au site topologique usuel de X sont cohérents ([49], lemme 0.1 ; pour la correction d’une}
\NormalTok{petite erreur dans la preuve, signalée par Jérôme Poineau, voir [52], 1.1.1). Si X = M (A)}
\NormalTok{alors OX (X) = A.}
\NormalTok{  Tout morphisme f : Y → X entre espaces analytiques est en particulier un morphisme}
\NormalTok{de sites annelés : si V est un domaine analytique de Y alors f −1 (V ) est un domaine}
\NormalTok{analytique de X et l’on dispose d’un morphisme f ∗ : OX (V ) → OY (f −1 (V )). Si X est}
\NormalTok{un espace k{-}affinoïde M (A) et si Y est un espace affinoïde M (B) sur une extension de}
\NormalTok{     § 11.1. GÉOMÉTRIE ANALYTIQUE : CONVENTIONS, NOTATIONS, PREMIERS RAPPELS               301}

\NormalTok{k, l’application f 7→ f ∗ établit une bijection entre l’ensemble des k{-}morphismes de Y}
\NormalTok{vers X et l’ensemble des morphismes bornés de k{-}algèbres de A vers B.}
\NormalTok{   Si X = M (A) est un espace k{-}affinoïde et si M est un A{-}module de type fini, le}
\NormalTok{faisceau Mf sur XG associé au préfaisceau V 7→ M ⊗A OX (V ) est cohérent. Le foncteur}
\NormalTok{M 7→ M f établit une équivalence entre la catégorie des A{-}modules de type fini et celle}
\NormalTok{des faisceaux cohérents sur X, de quasi{-}inverse F 7→ F (X) (cf. [9] ; ce résultat est dû}
\NormalTok{à Kiehl dans le cas strictement affinoïde).}
\NormalTok{   Lorsque X est bon, le foncteur π∗ induit une équivalence de catégories entre faisceaux}
\NormalTok{cohérents sur XG et faisceaux cohérents sur X, et qu’un faisceau cohérent F sur XG est}
\NormalTok{localement libre si et seulement si c’est le cas de π∗ F ([9], 1.3.4 ; notons que la définition}
\NormalTok{de faisceau cohérent qu’il utilise dans cette proposition est ad hoc, mais qu’elle coïncide}
\NormalTok{bien en fait avec la notion usuelle au vu de la cohérence des faisceaux structuraux sur}
\NormalTok{les deux sites annelés en jeu).}
\NormalTok{11.1.9. Produits fibrés, extension des scalaires. — La catégorie des espaces k{-}}
\NormalTok{analytiques admet des produits fibrés. Si A → B et A → C sont deux morphismes}
\NormalTok{d’algèbres k{-}affinoïdes, M (B) ×M (A) M (C) = M (B ⊗     b A C).}
\NormalTok{   Si L est un extension complète de k on dispose d’un foncteur d’extension des scalaires}
\NormalTok{X 7→ XL de la catégorie des espaces k{-}analytiques vers celle des espaces L{-}analytiques.}
\NormalTok{Si X = M (A) alors XL = M (A⊗       b k L). En général on dispose d’un morphisme naturel}
\NormalTok{de sites annelés de XL vers X, qui est surjectif (cf. [48], 0.5).}
\NormalTok{   Supposons X bon, soit y un point de XL et soit x son image sur X. Le morphisme}
\NormalTok{naturel de Spec OXL ,y vers Spec OX,x est plat à fibres localement d’intersection complète,}
\NormalTok{et à fibres géométriquement régulières si L est analytiquement séparable sur k ; c’est le}
\NormalTok{théorème [52], th. 2.6.5. La notion d’extension analytiquement séparable est définie}
\NormalTok{dans [49] (définition 1.6). Une extension finie de k est analytiquement séparable si et}
\NormalTok{seulement si elle est séparable ; si r est un polyrayon de longueur finie et k{-}libre, kr est}
\NormalTok{une extension analytiquement séparable de k ; si k est parfait, toute extension complète}
\NormalTok{de k est analytiquement séparable.}
\NormalTok{   On définit la catégorie des espaces analytiques au{-}dessus de k comme la catégorie dont}
\end{Highlighting}
\end{Shaded}

\subsection{CLD.txt, newline lines
15287--15355}\label{cld.txt-newline-lines-1528715355}

\begin{Shaded}
\begin{Highlighting}[]
\NormalTok{topologie de V induite par la topologie de Zariski de X, et en général strictement plus}
\NormalTok{fine.}
\NormalTok{   Si E est une partie de X et si (Xi ) est un G{-}recouvrement de X par des domaines}
\NormalTok{analytiques, alors E est un fermé (resp. un ouvert) de Zariski de X si et seulement si}
\NormalTok{E ∩ Xi est un fermé (resp. un ouvert) de Zariski de Xi pour tout i ([49], Prop. 4.2.1).}
\NormalTok{306                  CHAPITRE 11. RAPPELS DE GÉOMÉTRIE ANALYTIQUE}

\NormalTok{   Une immersion fermée (1) est un morphisme d’espaces k{-}analytiques f : Y → X se}
\NormalTok{factorisant par un isomorphisme entre Y et un sous{-}espace analytique fermé de X. Une}
\NormalTok{immersion fermée préserve les corps résiduels complétés.}
\NormalTok{   Un morphisme d’espaces k{-}analytiques f : Y → X est dit séparé si sa diagonale}
\NormalTok{δ : Y → Y ×X Y est une immersion fermée, et X est dit séparé si le morphisme structural}
\NormalTok{de X vers M (k) est séparé. La séparation dans ce sens entraîne la séparation topolo{-}}
\NormalTok{gique, mais la réciproque est fausse : le recollement de deux copies du disque unité le}
\NormalTok{long du domaine affinoïde décrit par l’égalité |T | = 1 est topologiquement séparé car le}
\NormalTok{domaine de recollement est fermé, mais n’est pas séparé.}
\NormalTok{11.2.3. Dimension. — On dispose d’une théorie de la dimension en géométrie ana{-}}
\NormalTok{lytique, due à Berkovich, au sujet de laquelle on pourra consulter le chapitre 2 de [8], la}
\NormalTok{fin du paragraphe 1.2 de [9] ou encore l’article [48].}
\NormalTok{   Soit X un espace k{-}analytique. Sa dimension (k{-}analytique) est définie comme suit.}
\NormalTok{Commençons par le cas affinoïde. Écrivons X = M (A) et choisissons une extension}
\NormalTok{complète L de k telle que AL := A⊗        b k L soit strictement L{-}affinoïde. La dimension de}
\NormalTok{Krull de AL ne dépend alors pas de L ([48], 1.5) et est appelée dimension de X. (Si X est}
\NormalTok{strictement k{-}affinoïde sa dimension est donc simplement sa dimension de Krull relative}
\NormalTok{à sa topologie de Zariski.) On a par construction dim Xred = dim X et dim XL = dim X}
\NormalTok{pour toute extension complète L de k.}
\NormalTok{   Si Y est un fermé de Zariski de X, le choix d’une structure de sous{-}espace analytique}
\NormalTok{fermé sur Y , c’est{-}à{-}dire d’un idéal I de A définissant Y , fait de Y un espace affinoïde}
\NormalTok{(identifié à M (A/I)). Sa dimension ne dépend pas du choix de I, et nous parlerons donc}
\NormalTok{de la dimension de Y sans spécifier I.}
\NormalTok{   Soit X1 , . . . , Xn les fermés de Zariski de X correspondant aux idéaux premiers mini{-}}
\NormalTok{maux de A. Si V est un domaine affinoïde de X on a alors dim V = supi,V ∩Xi ̸=∅ dim Xi .}
\NormalTok{En particulier dim V ⩽ dim X = supi dim Xi . Si x est un point de X on définit dimx X}
\NormalTok{comme le minimum des dim V où V est un domaine affinoïde de X contenant x ; c’est}
\NormalTok{aussi supi,x∈Xi dim Xi . Il vient dim X = supi dim Xi = supx∈X dimx X.}
\NormalTok{   Ne supposons plus X affinoïde. Si x est un point de X, on définit la dimension locale}
\NormalTok{dimx X de X en x comme le minimum des dim V où V parcourt l’ensemble des domaines}
\NormalTok{affinoïdes de X contenant x, et l’on pose alors dim X = supx∈X dimx X. Cette définition}
\NormalTok{est compatible avec la précédente lorsque X est affinoïde. Si V est un domaine analytique}
\NormalTok{de X on a dimx V = dimx X pour tout x ∈ V , si bien que dim V ⩽ dim X. Si d est un}
\NormalTok{entier on dit que X est purement de dimension d si dimx X = d pour tout x ∈ X. Si}
\NormalTok{c’est le cas X est ou bien vide, ou bien de dimension d.}
\NormalTok{   Pour tout point x de X on pose}
\NormalTok{(11.2.3.1)}
\NormalTok{                                                   1}
\NormalTok{      dk (x) = degtr(H      (x)/k̃) = degtr(H        k 1 ) + dimQ Q ⊗Z (|H (x)× |/|k × |) ,}
                                                                                         
\NormalTok{                         \^{}                  \^{}  (x) /e}
\NormalTok{où la seconde égalité provient de 10.1.9 ; Berkovich le note d(κ(x)/k) ou d(H (x)/k)}
\NormalTok{L’intérêt majeur de cet invariant est l’égalité dim(X) = supx∈X dk (x) ([9], Lemma 2.5.2).}
\NormalTok{   Soit Y un fermé de Zariski de X. Munissons{-}le d’une structure de sous{-}espace analy{-}}
\NormalTok{tique fermé. La dimension de l’espace analytique correspondant ne dépend que de Y , et}
\NormalTok{(1)}
\NormalTok{   Dans l’expression « immersion fermée » ou dans l’expression « sous{-}espace analytique fermé », l’ad{-}}
\NormalTok{jectif « fermé » fait implicitement référence à la topologie de Zariski. Ainsi, en général, un domaine}
\NormalTok{analytique fermé V de X n’est pas un sous{-}espace analytique fermé de X, et le morphisme canonique}
\NormalTok{V ,→ X n’est pas une immersion fermée.}
\NormalTok{                 § 11.2. TOPOLOGIE DE ZARISKI, THÉORIE DE LA DIMENSION                     307}

\NormalTok{pas de la structure en question (par exemple au vu de l’égalité ci{-}dessus, ou par réduc{-}}
\NormalTok{tion au cas affinoïde via la description des domaines analytiques de Y ). On pourra donc}
\NormalTok{parler de la dimension de Y sans spécifier la structure dont on le munit.}


\NormalTok{Remarque 11.2.4. — Insistons sur le fait que la dimension considérée ici est en fait}
\NormalTok{une dimension relative au corps de base k, même si nous ne rappelons pas cette dépen{-}}
\NormalTok{dance. Considérons par exemple l’algèbre de Banach kr associée à un polyrayon k{-}libre}
\NormalTok{(r1 , . . . , rn ) (exemple 11.1.3). C’est une extension complète de k. En tant qu’algèbre kr {-}}
\NormalTok{affinoïde elle est stricte, et sa dimension de Krull est nulle. Par conséquent la dimension}
\end{Highlighting}
\end{Shaded}

\subsection{CLD.txt, newline lines
15567--15604}\label{cld.txt-newline-lines-1556715604}

\begin{Shaded}
\begin{Highlighting}[]
\NormalTok{compact si l’image réciproque de toute partie compacte est compacte. Comme tout point}
\NormalTok{possède un voisinage compact, cela revient à dire que f est topologiquement propre et}
\NormalTok{topologiquement séparé.}
\NormalTok{  Suivant [9] (1.5.3), un morphisme est dit propre s’il est compact et sans bord.}
\NormalTok{11.4.4. Morphismes finis. — Soit A une algèbre k{-}affinoïde. Une A{-}algèbre de Ba{-}}
\NormalTok{nach finie est une A{-}algèbre de Banach B qui est finie comme A{-}algèbre, et dont la}
\NormalTok{norme est équivalente à la norme quotient induite par une présentation finie An → B}
\NormalTok{du A{-}module B. Dans ce cas la norme de B est en fait équivalente à la norme quotient}
\NormalTok{associée à n’importe quelle présentation finie du A{-}module B (c’est une conséquence de}
\NormalTok{[8], Prop. 2.1.9). Toute A{-}algèbre de Banach finie est k{-}affinoïde (combiner [20], Prop.}
\NormalTok{6.1.3/4 et [8], Cor. 2.1.8). Le foncteur d’oubli de la norme induit une équivalence de}
\NormalTok{catégories entre A{-}algèbres de Banach finies et A{-}algèbres finies(2) ([8], Prop. 2.1.12).}
\NormalTok{   Soit f : Y → X un morphisme entre espaces k{-}analytiques. Si X et Y sont affinoïdes,}
\NormalTok{on dit que f est fini si f ∗ fait de O(Y ) une O(X){-}algèbre de Banach finie ; cela revient}
\NormalTok{à demander que f soit sans bord ([8], Cor. 2.5.13). En général (si l’on ne suppose plus}
\NormalTok{X et Y affinoïdes) les assertions suivantes sont équivalentes : ([9], Cor. 1.3.7).}
\NormalTok{  (i) Il existe un G{-}recouvrement (Xi ) de X par des domaines affinoïdes tels que pour}
\NormalTok{tout i, l’image réciproque f −1 (Xi ) soit affinoïde et que le morphisme induit de f −1 (Xi )}
\NormalTok{vers Xi soit fini}
\NormalTok{  (ii) Pour tout domaine affinoïde V de X l’image réciproque f −1 (V ) est affinoïde et le}
\NormalTok{morphisme induit de f −1 (V ) vers V est fini.}

\NormalTok{  Si |k × | ̸= \{1\} et si A est strictement k{-}affinoïde, on déduit de [8], Prop. 2.1.1 et [20], Thm. 6.1.3/1}
\NormalTok{(2)}

\NormalTok{que toute A{-}algèbre de Banach qui est finie comme A{-}algèbre est une A{-}algèbre de Banach finie. Par}
\NormalTok{contre si k est trivialement valué et si 0 \textless{} r \textless{} s \textless{} 1, le morphisme continu naturel de k\{T /s\} vers}
\NormalTok{k\{T /r\} est bijectif (sa source et son but s’identifient tous deux à k[[T ]]) mais sa réciproque n’est pas}
\NormalTok{bornée, si bien que k\{T /s\} n’est pas une k\{T /r\}{-}algèbre de Banach finie).}
\NormalTok{312                CHAPITRE 11. RAPPELS DE GÉOMÉTRIE ANALYTIQUE}

\NormalTok{On dit que f est fini si elles sont satisfaites ; cette définition est compatible avec la}
\NormalTok{précédente si X et Y sont affinoïdes. Un morphisme fini est sans bord.}
\NormalTok{  La classe des morphismes finis est stable par composition et changement de base}
\NormalTok{analytique sur k, et la propriété d’être fini est G{-}locale au but.}
\NormalTok{  Soit y un point de Y . On dit que f est fini en y s’il existe un voisinage analytique}
\NormalTok{U de f (y) dans X et un voisinage analytique V de y dans f −1 (U ) tel que f induise un}
\NormalTok{morphisme fini de V vers U . Si c’est le cas le point y appartient à Int(Y /X), et si de}
\end{Highlighting}
\end{Shaded}

\subsection{CLD.txt, newline lines
16122--16143}\label{cld.txt-newline-lines-1612216143}

\begin{Shaded}
\begin{Highlighting}[]
\NormalTok{  Supposons réciproquement que (X, \^{}   x) soit l’image réciproque d’un ouvert affine de}
\NormalTok{P \^{} 1 e1 . Il résulte de la preuve du lemme 3.5.1 de [52] qu’il existe un entier n, un}
\NormalTok{ H (x) /k}
\NormalTok{point y du polydisque unité D de dimension n et un morphisme sans bord de (X, x) vers}
\NormalTok{(D, y). Lorsque la valeur absolue de k n’est pas triviale, y possède une base de voisinages}
\NormalTok{strictement affinoïdes dans D, si bien que (X, x) possède un représentant strictement}
\NormalTok{affinoïde.}
\NormalTok{11.7.7. — Supposons que la valeur absolue de k n’est pas triviale. On déduit alors}
\NormalTok{des résultats rappelés au 11.7.6 que (X, x) est strictement k{-}analytique si et seulement}
\NormalTok{   \^{}}
\NormalTok{si (X, x) est strict, cf. [52, Lemma 3.5.2]. Mentionnons incidemment que si (X, x) est}
\NormalTok{                         \^{}}
\NormalTok{à la fois bon et strict, (X, x) est à la fois affine et strict, donc est l’image réciproque}
\NormalTok{d’un ouvert affine de P \^{} 1 e1 (10.5.5) ; il découle alors de 11.7.6 que le germe (X, x) est}
\NormalTok{                        H (x) /k}
\NormalTok{strictement affinoïde. Autrement dit tout point d’un bon espace strictement k{-}analytique}
\NormalTok{admet un voisinage strictement affinoïde ce qui, à notre connaissance, ne figurait pas}
\NormalTok{jusqu’ici explicitement dans la littérature.}
\NormalTok{   Notons que l’hypothèse que la valeur absolue de k n’est pas triviale est indispensable}
\NormalTok{à la validité des énoncés ci{-}dessus. Supposons en effet que |k × | = \{1\}, et soit X l’ouvert}
\NormalTok{de A1k défini par les inégalités 0 \textless{} |T | \textless{} 1 ; il est sans bord, si bien que pour tout x ∈ X}
\NormalTok{              \^{}}
\end{Highlighting}
\end{Shaded}

\subsection{CLD.txt, newline lines
3750--3920}\label{cld.txt-newline-lines-37503920}

\begin{Shaded}
\begin{Highlighting}[]
\NormalTok{3.3.8. — Soit ω une G{-}forme (resp. une forme) de type (p, q) à coefficients dans G}
\NormalTok{sur X. Si V est un domaine de X, une G{-}tropicalisation de ω sur V est la donnée d’une}
\NormalTok{carte G{-}tropicale (f : V → E, P ) et d’une forme α ∈ GPp,q (P ) telle que ω = f ∗ α.}
\NormalTok{  On dira que ω est G{-}tropicale si elle admet une G{-}tropicalisation sur X.}
\NormalTok{  Si V est un domaine de X tel que ω possède une G{-}tropicalisation sur V , nous dirons}
\NormalTok{aussi que V G{-}tropicalise ω.}
\NormalTok{  Soit V un domaine compact de X qui G{-}tropicalise ω, et soit (f V → E, P, α) une}
\NormalTok{G{-}tropicalisation de ω|V . Comme f (V ) est un sous{-}espace paralinéaire de E, (f : V →}
\NormalTok{E, f (V ), α|f (V ) ) est encore une G{-}tropicalisation de ω|V ; une telle G{-}tropicalisation sera}
\NormalTok{simplement notée (f : V → E, α) : lorsque V est compact, l’omission du sous{-}espace}
\NormalTok{paralinéaire P signifie qu’il est égal à f (V ).}

\NormalTok{Remarque 3.3.9. — Soit H une classe de coefficients contenant G . Observons que}
\NormalTok{les G{-}formes à coefficients dans G s’identifient à des G{-}formes à coefficients dans H .}
\NormalTok{   En effet, le G{-}préfaisceau Gpfp,q sur XG s’identifie à un sous{-}G{-}préfaisceau de Hpfp,q , de}
\NormalTok{sorte que le G{-}faisceau GXp,qG sur X s’identifie alors à un sous{-}G{-}faisceau de HXp,qG}
\NormalTok{                                                                                       . Ces}
\NormalTok{identifications sont compatibles avec les constructions faites dans ce paragraphe.}

\NormalTok{3.3.10. — Le faisceau GXG := GX0,0G est un faisceau en R{-}algèbres et les faisceaux GXp,qG}
\NormalTok{sont des GXG {-}modules. Leur somme directe GX∗,∗G = p,q GXp,qG est un faisceau en algèbres}
\NormalTok{                                                   L}

\NormalTok{bigraduées commutatives : si ω est de type (p, q) et ω ′ de type (p′ , q ′ ), alors}
\NormalTok{                                                    ′     ′}
\NormalTok{                                ω ′ ∧ ω = (−1)(p+p )(q+q ) ω ∧ ω ′ .}
\NormalTok{3.3.11. — Le G{-}préfaisceau Gpf∗,∗ est naturellement muni d’un opérateur J qui s’étend}
\NormalTok{au faisceau GX∗,∗G . Ce sont des morphismes d’algèbres. On dit qu’une section ω de GXp,qG}
\NormalTok{est symétrique si p = q et si Jω = ω.}
\NormalTok{    Soit H une classe de coefficients telles que les dérivées partielles des fonctions de G}
\NormalTok{existent et appartiennent à H . On dispose alors d’opérateurs différentiels d′ et d′′ de}
\NormalTok{Gpf∗,∗ dans Hpf∗,∗ ; ils induisent des opérateurs différentiels de GX∗,∗G dans HX∗,∗}
\NormalTok{                                                                                  G}
\NormalTok{                                                                                     , caractérisés}
\NormalTok{par les relations}
\NormalTok{                               d′ f ∗ α = f ∗ d′ α, d′′ f ∗ α = f ∗ d′′ α,}
\NormalTok{                                 § 3.4. NULLITÉ DES G{-}FORMES                             81}

\NormalTok{pour toute carte G{-}tropicale (f : U → E, P ) et toute forme α de type (p, q) à coeffi{-}}
\NormalTok{cients dans G sur P . Les différentielles vérifient la règle de Leibniz : si ω et ω ′ sont}
\NormalTok{respectivement de type (p, q) et (p′ , q ′ ), on a}
\NormalTok{(3.3.11.1)               d′ (ω ∧ ω ′ ) = d′ ω ∧ ω ′ + (−1)p+q ω ∧ d′ ω ′ ,}
\NormalTok{et de même pour d′′ (ω ∧ ω ′ ).}
\NormalTok{   Si G ⊂ C 2 , ces opérateurs vérifient d′ d′ = d′′ d′′ = 0.}
\NormalTok{3.3.12. — Cette construction est fonctorielle. Plus précisément, si f : Y → X est un}
\NormalTok{morphisme d’espaces G{-}tropicaux, on dispose de morphismes de faisceaux}
\NormalTok{                                     f ∗ : GX∗,∗G → f∗ (GY∗,∗}
\NormalTok{                                                           G}
\NormalTok{                                                              )}
\NormalTok{sur XG qui sont des morphismes d’anneaux gradués, commutent aux dérivations d′ et}
\NormalTok{d′′ (lorsqu’elles existent) et à l’opérateur J.}
\NormalTok{   Lorsque le contexte rend évident le morphisme f , on se permettra parfois de noter}
\NormalTok{ω|Y l’image dans GYp,q G}
\NormalTok{                         (Y ) d’une G{-}forme ω ∈ GXp,qG .}
\NormalTok{3.3.13. — Soit (p, q) un couple d’entiers naturels. Soit Σ une partie paralinéaire de X}
\NormalTok{et soit j l’inclusion de Σ dans X. C’est un morphisme d’espaces G{-}tropicaux. On dispose}
\NormalTok{donc d’un morphisme de faisceaux de FXp,qG dans j∗ FΣp,qG . On notera ω 7→ ω|Σ l’image}
\NormalTok{d’une forme ω et on dira que c’en est la restriction à Σ.}
\NormalTok{   Concrètement, cette restriction se calcule de la façon suivante.}
\NormalTok{   Soit V un domaine de X, soit (f : V → E, P ) une carte G{-}tropicale sur V , et soit}
\NormalTok{α ∈ FPp,q (P ). Alors Σ ∩ V est une partie paralinéaire de V , l’application f : Σ ∩ V → E}
\NormalTok{est paralinéaire et son image est contenue dans P , l’image réciproque f ∗ α de α est}
\NormalTok{une forme de type (p, q) sur l’espace paralinéaire Σ ∩ V . Par définition de l’ensemble}
\NormalTok{Fpfp,q (V ) comme limite inductive, la forme f ∗ α ne dépend que de la classe de f ∗ α dans}
\NormalTok{cet ensemble. Cette construction définit un morphisme de préfaisceaux de FX,pf     p,q}
\NormalTok{                                                                                       dans}
\NormalTok{j∗ FΣp,q .}
\NormalTok{   Par passage aux faisceaux associés, on en déduit le morphisme de restriction de FXp,qG}
\NormalTok{dans j∗ FΣp,qG .}
\NormalTok{3.3.14. — Soit f : Y → X un morphisme d’espaces G{-}tropicaux, soit T une partie}
\NormalTok{paralinéaires de Y et soit Σ une partie paralinéaire de X telle que f (Σ) ⊂ T. Pour toute}
\NormalTok{forme ω ∈ FXp,qG (X), on a (f ∗ ω)|T = (f |T )∗ (ω|Σ ).}
\NormalTok{   Nous appliquerons notamment cette remarque lorsque Y est un domaine de X et}
\NormalTok{T = Σ ∩ Y , ou bien lorsque Y = X et T est une partie polyédrale de Σ.}



\NormalTok{                          § 3.4. NULLITÉ DES G{-}FORMES}

\NormalTok{Proposition 3.4.1. — Soit X un espace G{-}tropical compact, soit f : X → E une carte}
\NormalTok{G{-}tropicale, soit P un sous{-}ensemble paralinéaire de E contenant f (X) et soit α ∈}
\NormalTok{GPp,q (P ). Les assertions suivantes sont équivalentes :}
\NormalTok{   a) La restriction de α au sous{-}espace paralinéaire f (X) est nulle ;}
\NormalTok{   b) La section globale f ∗ α de Gpfp,q est nulle.}
\NormalTok{Démonstration. — L’hypothèse (1) entraîne que l’image de α dans la carte tropicale}
\NormalTok{(f, f (X)) est nulle, de sorte que (2) vaut. Supposons que (2) est vérifiée et prouvons (1).}
\NormalTok{82                          CHAPITRE 3. ESPACES G{-}TROPICAUX}

\NormalTok{   On peut remplacer P par f (X). Par définition, il existe une carte G{-}tropicale g : X →}
\NormalTok{F , une morphisme affine p : F → E tel que p ◦ g = f , et un sous{-}espace paralinéaire Q}
\NormalTok{de F contenant g(X) tels que p(Q) ⊂ P et p∗ α = 0 dans GQp,q (Q). On peut aussi}
\NormalTok{remplacer Q par g(X) ; on a alors p(Q) = P .}
\NormalTok{   Le morphisme de Q sur P induit par p possède G{-}localement des sections. D’après}
\NormalTok{le corollaire 1.8.5, il existe une décomposition cellulaire (Ci ) de P et, pour tout i, une}
\NormalTok{section de ptrop sur Ci à valeurs dans Q. Puisque p∗ α = 0, on a donc α|Ci = 0 pour}
\NormalTok{tout i. Par suite, α = 0.}

\NormalTok{Corollaire 3.4.2. — Soit X un espace G{-}tropical compact, soit ω ∈ Gpfp,q (X) et soit}
\NormalTok{(Xi ) un G{-}recouvrement de X. Supposons que pour tout i, la restriction de ω à Xi soit}
\NormalTok{nulle ; alors ω = 0.}
\NormalTok{Démonstration. — Tout G{-}recouvrement admet un G{-}recouvrement plus fin formé de}
\NormalTok{domaines compacts. Comme X est compact, cela permet de supposer que le recouvre{-}}
\NormalTok{ment (Xi ) est formé d’un nombre fini de domaines compacts. Par définition du préfais{-}}
\NormalTok{ceau Gpfp,q , il existe une carte G{-}tropicale f : X → E, un sous{-}espace paralinéaire P de E}
\NormalTok{contenant f (X), et une forme α ∈ GPp,q (P ) tels que ω = f ∗ α.}
\NormalTok{   Il résulte que de la proposition 3.4.1 que α|f (Xi ) = 0 pour tout i. Comme f (X) est la}
\NormalTok{réunion des f (Xi ), la forme α est nulle sur f (X) (lemme 1.9.2). On a donc ω = f ∗ α =}
\NormalTok{0.}

\NormalTok{Remarque 3.4.3. — L’hypothèse que X est compact est cruciale pour le corol{-}}
\NormalTok{laire 3.4.2.}
\NormalTok{   Prenons pour espace G{-}tropical l’ensemble X = R, muni de la topologie discrète, et}
\NormalTok{pour carte l’application identique f : R → R. Soit α la forme d′ x sur R. Soit g : X → Rn}
\NormalTok{une carte qui domine la carte f , et soit p : Rn → R une application affine telle que}
\NormalTok{f = p ◦ g. Alors, p est surjective, donc p∗ α ̸= 0. Cela prouve que la section f ∗ α de Gpf1,0}
\NormalTok{n’est pas nulle. En revanche, pour tout t ∈ R, le singleton Vt = \{t\} est un domaine de X}
\NormalTok{de tropicalisation Pt = \{t\} et α|Pt = 0, si bien que α|Vt = 0. Comme ces domaines Vt}
\NormalTok{forment un G{-}recouvrement de X, cela montre que le préfaisceau Gpf1,0 n’est pas séparé.}
\NormalTok{Proposition 3.4.4. — Soit X un espace G{-}tropical, soit f : X → E une carte G{-}}
\NormalTok{tropicale, soit P un sous{-}espace paralinéaire de E contenant f (X) et soit α ∈ FPp,q (P ).}
\NormalTok{   Les propriétés suivantes sont équivalentes :}
\NormalTok{   a) Pour tout domaine compact V de X, la restriction de α au polytope f (V ) est nulle ;}
\NormalTok{   b) Pour tout domaine compact V de X, la restriction de f ∗ α à V est nulle ;}
\NormalTok{   c) L’image de f ∗ α dans FXp,qG (X) est nulle.}
\NormalTok{Démonstration. — D’après la proposition 3.4.1, les assertions (1) et (2) sont équiva{-}}
\NormalTok{lentes.}
\NormalTok{   Supposons (1). Comme X admet un G{-}recouvrement par des domaines compacts,}
\NormalTok{l’assertion (3) s’en déduit.}
\NormalTok{   Supposons (3). Il existe donc un G{-}recouvrement (Xi ) tel que f ∗ α|Xi = 0 en tant que}
\NormalTok{section du préfaisceau Fpfp,q . Soit V un domaine compact de X. La restriction de f ∗ α}
\NormalTok{à Xi ∩ V est nulle pour tout i ; comme la famille (Xi ∩ V ) est un G{-}recouvrement de V ,}
\NormalTok{le corollaire 3.4.2 entraîne que f ∗ α|V = 0, d’où l’assertion (2). Cela termine la preuve}
\NormalTok{de la proposition.}
\NormalTok{                                § 3.4. NULLITÉ DES G{-}FORMES                                83}

\NormalTok{Corollaire 3.4.5. — Soit X un espace G{-}tropical compact, soit f : X → E une carte}
\NormalTok{G{-}tropical et soit P le sous{-}espace paralinéaire f (X). Soit α ∈ GPp,q (P ), et soit ω la}
\NormalTok{G{-}forme f ∗ α. Notons K le support de ω dans X et H celui de α dans f (X). On a}
\NormalTok{H = f (K).}

\NormalTok{Démonstration. — Soit x un point de X. Soit Q un voisinage de f (x) dans P qui est}
\NormalTok{                                             −1}
\NormalTok{un polytope compact. Le domaine W = ftrop       (Q) de X est un voisinage de x et ω|W =}
\NormalTok{       ∗}
\NormalTok{(f |W ) (α|Q ).}
\NormalTok{   Supposons que x appartienne à K. Le compact K rencontre alors l’intérieur de W , ce}
\NormalTok{qui entraîne que ω|W ̸= 0, et donc que α|Q ̸= 0. Il s’ensuit que H rencontre Q. Comme}
\NormalTok{Q peut être choisi arbitrairement petit, f (x) appartient à H, d’où l’inclusion f (K) ⊂ H.}
\NormalTok{   Réciproquement, supposons que f (x) ∈ H. Par définition de H, la restriction à Q de}
\NormalTok{la forme α est non nulle. D’après la proposition 3.4.4, la restriction de ω à f −1 (Q) est}
\NormalTok{non nulle. Par conséquent, f −1 (Q) rencontre K, donc Q = f (f −1 (Q)) rencontre f (K).}
\NormalTok{Comme Q est un voisinage arbitrairement petit de f (x), celui{-}ci est adhérent à f (K).}
\NormalTok{Puisque f (K) est compact, f (x) appartient à f (K).}

\NormalTok{3.4.6. — À toute carte G{-}tropicale (f : V → E, P ) sur un domaine V de X et à toute}
\NormalTok{fonction u sur P , associons la fonction u◦f sur V . On définit ainsi un morphisme de pré{-}}
\NormalTok{faisceaux de Fpf0,0 vers le G{-}faisceau des fonctions sur X. Il découle de la proposition 3.4.4}
\NormalTok{que le morphisme de FX0,0G dans le G{-}faisceau des fonctions sur X est injectif, morphisme}
\NormalTok{par lequel nous identifierons systématiquement une (0, 0){-}G{-}forme à une fonction.}
\NormalTok{   Une section globale du faisceau GX0,0G est une fonction sur X qui est G{-}localement de la}
\NormalTok{forme φ(f1 , . . . , fn ) où les fi sont des fonctions paralinéaires sur X, et où φ appartient}
\NormalTok{à la classe G . De telles fonctions seront appelées G{-}localement tropicales de classe G .}

\NormalTok{Définition 3.4.7. — Soit A la classe des fonctions lisses. Les sections de AX0,0}
\NormalTok{                                                                              G}
\NormalTok{                                                                                 sont}
\NormalTok{appelées fonctions G{-}lisses sur X.}

\NormalTok{   Les fonctions G{-}lisses sont donc les fonctions sur X sont G{-}localement de la forme}
\NormalTok{φ(f1 , . . . , fn ) où les fi sont des fonctions paralinéaires et où φ est une fonction de}
\NormalTok{classe C ∞ .}
\end{Highlighting}
\end{Shaded}

\subsection{CLD.txt, newline lines
5289--5314}\label{cld.txt-newline-lines-52895314}

\begin{Shaded}
\begin{Highlighting}[]

\NormalTok{Lemme 4.4.5. — Soit f : Y → X une immersion par morceaux d’espaces paralinéaires}
\NormalTok{de niveau n. Soit α une forme{-}volume à coefficients boréliens sur X. On a l’égalité}
\NormalTok{(f ∗ α)Bor = f ∗ (αBor ) de mesures de Borel régulières positives sur Y . Si α est localement}
\NormalTok{intégrable, alors f ∗ α est localement intégrable et l’on a l’égalité (f ∗ α)Rad = f ∗ (αRad ) de}
\NormalTok{mesures de Radon sur Y .}

\NormalTok{   Comme f est une immersion par morceaux, l’image inverse de la mesure de Borel}
\NormalTok{régulière αBor par f est définie G{-}localement (cf. A.3.13) ; de même pour les mesures de}
\NormalTok{Radon.}
\NormalTok{Démonstration. — L’espace Y est G{-}recouvert par des cellules C telles que f induise un}
\NormalTok{isomorphisme de C sur son image, laquelle possède une immersion paralinéaire dans Rn .}
\NormalTok{L’assertion découle alors de la construction de ces mesures.}

\NormalTok{Lemme 4.4.6. — Soit f : Y → X une immersion par morceaux compacte d’espaces}
\NormalTok{paralinéaires de niveau n. Soit α une forme{-}volume à coefficients boréliens sur Y . On}
\NormalTok{a une inégalité (f∗ α)Bor ⩽ f∗ (αBor ) de mesures de Borel régulières positives sur X, qui}
\NormalTok{est une égalité si f est injective. Si α est localement intégrable, alors f∗ α est localement}
\NormalTok{intégrable et l’on a l’égalité (f∗ α)Rad = f∗ (αRad ) de mesures de Radon sur X.}

\NormalTok{Démonstration. — En raisonnant localement, on se ramène au cas où X est paracom{-}}
\NormalTok{pact. Soit (C , D) un couple de décompositions cellulaires qui est adapté à f . Après}
\NormalTok{restriction à toute cellule D de D, les assertions découlent directement des définitions}
\NormalTok{des mesures et des images directes.}
\NormalTok{108    CHAPITRE 4. CALIBRAGES ET INTÉGRATION SUR LES ESPACES PARALINÉAIRES}
\end{Highlighting}
\end{Shaded}

\subsection{CLD.txt, newline lines
5449--5483}\label{cld.txt-newline-lines-54495483}

\begin{Shaded}
\begin{Highlighting}[]
\NormalTok{forme ω sont positifs.}
\NormalTok{   Les (n, n){-}formes localement µ{-}intégrables forment un sous{-}faisceau en BXG {-}module}
\NormalTok{de BX n,n}
\NormalTok{        G}
\NormalTok{          , et l’application ω 7→ ⟨ω, µ⟩ est linéaire.}
\NormalTok{   On dira enfin que le calibrage µ est localement intégrable si toute forme à coefficients}
\NormalTok{localement bornés est localement µ{-}intégrable. Sur une cellule, cela revient à dire que}
\NormalTok{le calibrage est produit d’un calibrage standard par une fonction localement intégrable.}
\NormalTok{Les calibrages G{-}localement constants sont localement intégrables.}

\NormalTok{Exemple 4.5.7. — Lorsque X est un sous{-}espace affine de dimension n d’un espace}
\NormalTok{affine V , ω une forme lisse sur X et µ un vecteur{-}volume sur X, on retrouve la notion}
\NormalTok{d’intégrale définie dans le paragraphe 1.5.3 à partir des idées de Lagerberg.}
\NormalTok{   D’une certaine manière, les constructions de ce chapitre visaient simplement à « G{-}}
\NormalTok{faisceautiser » cette formule élémentaire de Lagerberg sur un espace affine.}

\NormalTok{Proposition 4.5.8. — Soit f : Y → X une immersion par morceaux propre d’espaces}
\NormalTok{paralinéaire de niveau n.}
\NormalTok{   a) Pour toute (n, n){-}forme ω sur X et tout calibrage µ de Y , on a l’égalité f∗ ⟨f ∗ ω, µ⟩ =}
\NormalTok{⟨ω, f∗ µ⟩ de formes{-}volume n{-}dimensionnelles sur X.}
\NormalTok{   b) Pour toute (n, n){-}forme ω sur Y et tout calibrage µ de X, on a l’égalité ⟨f∗ ω, µ⟩ =}
\NormalTok{f∗ ⟨ω, f ∗ µ⟩ de formes{-}volume n{-}dimensionnelles sur X.}

\NormalTok{Démonstration. — En raisonnant localement, on se ramène au cas où X est paracom{-}}
\NormalTok{pact. Il existe alors un couple (C , D) de décompositions cellulaires qui est adapté à f .}
\NormalTok{                  § 4.6. CALIBRAGES ET INTÉGRALES DES (n − 1, n){-}FORMES                  111}

\NormalTok{Les assertions indiquées se vérifient alors, cellule par cellule, à partir des constructions}
\NormalTok{des intégrales et de la définition du calibrage f∗ µ (resp. du calibrage f ∗ µ).}



\NormalTok{     § 4.6. CALIBRAGES ET INTÉGRALES DES (n − 1, n){-}FORMES}

\NormalTok{4.6.1. — Soit G une classe de coefficients. Soit X un espace paralinéaire de niveau n,}
\end{Highlighting}
\end{Shaded}

\subsection{CLD.txt, newline lines
6308--6332}\label{cld.txt-newline-lines-63086332}

\begin{Shaded}
\begin{Highlighting}[]
\NormalTok{Lemme 5.3.4. — Si la forme ω est G{-}tropicalisée par une carte G{-}tropicale f de do{-}}
\NormalTok{                (n)}
\NormalTok{maine X, alors Σf supporte fortement ω.}

\NormalTok{Démonstration. — Le support de ω est contenu dans l’ensemble des points x ∈ X en}
\NormalTok{                                                                    (n)}
\NormalTok{lesquels f est de dimension tropicale ⩾ n qui, par définition, est Σf .}
\NormalTok{                                                                  (n)}
\NormalTok{  Soit T un sous{-}espace paralinéaire fermé de X qui contient Σf . Démontrons que ω|T}
\NormalTok{est le prolongement par zéro de ω|Σ(n) . Soit C une décomposition cellulaire en f {-}cellules}
\NormalTok{                                     f}
\NormalTok{                             (n)}
\NormalTok{de T qui est compatible à Σf . Il suffit de prouver que pour toute cellule C de C      telle}
\NormalTok{                        (n)}
\NormalTok{que C̊ est disjoint de Σf , on a ω|C = 0.}
\NormalTok{  Par hypothèse, f est de dimension tropicale \textless{} n en tout point de C̊, donc il existe une}
\NormalTok{cellule D, compacte, d’intérieur non vide, de C̊ telle que dim(f (D)) \textless{} n. Comme C est}
\NormalTok{une f {-}cellule, on dim(f (C)) = dim(f (C̊)) = dim(f (D)). D’après la proposition 3.4.4, on}
\NormalTok{a donc ω|C = 0.}
\NormalTok{                          (n)}
\NormalTok{  Cela démontre que Σf supporte fortement ω.}

\NormalTok{Proposition 5.3.5. — Soit n un entier et soit X un espace G{-}tropical vertébré de}
\NormalTok{niveau n. Soit (p, q) un couple d’entiers naturels tel que max(p, q) = n et soit ω une}
\NormalTok{G{-}forme de type (p, q) sur X. Supposons que X est paracompact, ou que le support de ω}
\end{Highlighting}
\end{Shaded}

\subsection{CLD.txt, newline lines
6400--6414}\label{cld.txt-newline-lines-64006414}

\begin{Shaded}
\begin{Highlighting}[]
\NormalTok{                                      f}
\NormalTok{composition cellulaire finie en f {-}cellules.}
\NormalTok{                                                                         (n)}
\NormalTok{   Lorsque f est une carte agréable, nous noterons Πf l’ensemble f (Σf ) et l’appellerons}
\NormalTok{le polyèdre tropical de la carte G{-}tropicale agréable f . Il est purement de dimension n}
\NormalTok{                                                               (n)}
\NormalTok{et l’application f induit une immersion par morceaux de Σf sur Πf .}
\NormalTok{5.4.4. — Si X est compact, toute carte G{-}tropicale est agréable et l’on a Πf = f (X)(n)}
\NormalTok{(proposition 5.2.5).}
\NormalTok{Proposition 5.4.5. — Soit X un espace G{-}tropical vertébré de niveau n, soit f : X →}
\NormalTok{E une carte G{-}tropicale agréable et soit P un sous{-}espace paralinéaire P de E. Consi{-}}
\NormalTok{dérons la carte G{-}tropicale fP : f −1 (P ) → E sur le domaine f −1 (P ) (considéré comme}
\NormalTok{130                      CHAPITRE 5. ESPACES G{-}TROPICAUX CALIBRÉS}

\NormalTok{espace G{-}tropical vertébré de niveau n) déduite de f par restriction. La carte fP est}
\end{Highlighting}
\end{Shaded}

\subsection{CLD.txt, newline lines
6753--6812}\label{cld.txt-newline-lines-67536812}

\begin{Shaded}
\begin{Highlighting}[]
\NormalTok{            = ⟨ω|T , µT ⟩                  d’après (3.3.14)}
\NormalTok{            = ⟨(ω|Σ )|T , (µΣ )|T ⟩        d’après (3.3.14) et la définition d’un calibrage}
\NormalTok{            = ⟨ω|Σ , µΣ ⟩|T = ω}
\NormalTok{                              e |Y}
\NormalTok{parce que l’accouplement entre (n, n){-}formes et calibrages est un morphisme de faisceaux}
\NormalTok{sur le site XG . Le lemme est ainsi démontré.}

\NormalTok{Proposition 5.8.3. — Il existe un unique morphisme de faisceaux FXn,n            G}
\NormalTok{                                                                                    → VXnG}
\NormalTok{sur XG , noté ω 7→ ω e , qui prolonge l’application définie en (5.8.1) lorsque ω possède}
\NormalTok{un support fort. Ce morphisme est FX0,0G {-}linéaire. Si ω est une G{-}forme positive, la}
\NormalTok{forme{-}volume ωe est positive. Si ω est à coefficients dans une classe G , il en va de même}
\NormalTok{de la forme{-}volume ωe.}
\NormalTok{Démonstration. — Compte tenu du lemme 5.8.2, la première assertion résulte de ce que}
\NormalTok{toute forme sur un domaine paracompact possède un support fort, et que tout domaine}
\NormalTok{est G{-}recouvert par ses domaines paracompacts. Les autres propriétés découlent alors}
\NormalTok{de la construction.}

\NormalTok{5.8.4. — Explicitons cette construction lorsque ω est tropicale. Soit f : X → E une}
\NormalTok{carte G{-}tropicale et soit α une (n, n){-}forme sur un sous{-}espace paralinéaire de E conte{-}}
\NormalTok{                                         (n)}
\NormalTok{nant f (X) telle que ω = f ∗ α. Alors, Σf est un support fort de ω (lemme 5.3.4) et on}
\NormalTok{                                         (n)}
\NormalTok{a l’égalité de formes{-}volume sur Σf            :}
\NormalTok{                                                ∗}
\NormalTok{(5.8.4.1)                                ω̃ = ⟨ftrop α, µΣ(n) ⟩.}
\NormalTok{                                                            f}

\NormalTok{Si, de plus, la carte G{-}tropicale f est agréable, la proposition 4.5.8 entraîne alors que}
\NormalTok{(5.8.4.2)                       ftrop,∗ (ω̃) = ⟨α, f∗ µΣ(n) ⟩ = ⟨α, µΠf ⟩.}
\NormalTok{                                                        f}
\NormalTok{                                   § 5.9. MESURES DE BORD                                  137}

\NormalTok{ 5.8.5. — Soit ω une forme de type (n, n) sur XG à coefficients tropicalement boréliens.}
\NormalTok{ Nous venons de lui associer une forme{-}volume n{-}dimensionnelle ω̃. Nous noterons alors}
\NormalTok{ ωBor la mesure de Borel sur X associée à ω̃ par la construction du §5.7.1.}
\NormalTok{   Nous dirons que ω est localement intégrable si cette mesure est localement finie, et}
\NormalTok{ que ω est intégrable si cette mesure est finie. Lorsque ω est localement intégrable, nous}
\NormalTok{ noterons ωRad la mesure de Radon associée à ω̃ ; on a |ωRad | = ωBor . Lorsque ω est une}
\NormalTok{ G{-}forme positive, la mesure ωRad est positive.         R           R}
\NormalTok{   Si f est une fonction borélienne sur X, on notera X f |ω| et X f ω les intégrales de f}
\NormalTok{Rcontre les mesures ωBor et ωRad lorsqu’elles existent. Cela définit en particulier l’intégrale}
\NormalTok{  X}
\NormalTok{    ω pour toute G{-}forme de type (n, n) à support compact qui est localement intégrable}
\NormalTok{(par exemple lisse).}
\NormalTok{   Les G{-}formes localement intégrables constituent un sous{-}espace vectoriel de F n,n (XG )}
\NormalTok{contenant les G{-}formes lisses. Les G{-}formes intégrables en sont un sous{-}espace vectoriel}
\NormalTok{qui contient les G{-}formes lisses à support compact.}
\NormalTok{5.8.6. — Soit u une fonction localement intégrable sur X et soit ω une G{-}forme lisse}
\NormalTok{de type (n, n) sur X. Alors, u est localement intégrable par rapport à la mesure ωBor ,}
\NormalTok{de sorte que uωRad est une mesure de Radon sur X.}
\NormalTok{   L’application (u, ω) 7→ uωRad est bilinéaire ; elle induit une application linéaire de}
\NormalTok{Lloc}
\NormalTok{   1}
\NormalTok{     (XG ) ⊗ A n,n (XG ) dans l’ensemble des mesures de Radon sur X. Cette application}
\NormalTok{se factorise par L1loc (XG ) ⊗ A n,n (XG ). On la notera bien sûr α 7→ αRad .}
\NormalTok{   Cette construction s’étend en des morphismes de G{-}faisceaux}
\NormalTok{                  Lloc,X}
\NormalTok{                    1}
\NormalTok{                        G}
\end{Highlighting}
\end{Shaded}

\subsection{CLD.txt, newline lines
6934--6977}\label{cld.txt-newline-lines-69346977}

\begin{Shaded}
\begin{Highlighting}[]

\NormalTok{               § 5.10. INTÉGRALES ET TROPICALISATIONS}

\NormalTok{  Soit X un espace G{-}tropical calibré de niveau n.}
\NormalTok{Proposition 5.10.1. — Soit f : X → E une carte G{-}tropicale agréable, soit Πf le}
\NormalTok{polyèdre tropical de la carte f et soit µΠf son calibrage.}
\NormalTok{  Soit p ∈ \{n−1, n\} et soit α une (p, n){-}forme à coefficients boréliens sur un sous{-}espace}
\NormalTok{paralinéaire P de E contenant f (X).}
\NormalTok{  a) Si p = n, on a l’égalité}
\NormalTok{                                   ftrop,∗ (f ∗ α)Bor ⩾ ⟨α, µΠf ⟩Bor}
\NormalTok{                             (n)}
\NormalTok{de mesures de Borel sur Πf .}
\NormalTok{  b) Si p = n − 1, on a l’inégalité}
\NormalTok{                                   ftrop,∗ (f ∗ α)Bor ⩾ ⟨α, µΠf ⟩Bor}
\NormalTok{                             (n)}
\NormalTok{de mesures de Borel sur Πf .}
\NormalTok{  c) La mesure (f ∗ α)Bor est localement finie si et seulement si la forme α|Πf est locale{-}}
\NormalTok{ment intégrable sur Πf ; dans ce cas, on a l’égalité}
\NormalTok{                                   ftrop,∗ (f ∗ α)Rad = ⟨α, µΠf ⟩Rad}
\NormalTok{de mesures de Radon sur Πf .}
\NormalTok{   d) En particulier, si α est intégrable sur Πf , la forme f ∗ α est intégrable sur X, on a}
\NormalTok{l’égalité}
\NormalTok{                                   Z         Z}
\NormalTok{                                        ∗}
\NormalTok{(5.10.1.1)                            f α=       ⟨α, µΠf ⟩}
\NormalTok{                                        X           Πf}

\NormalTok{lorsque p = n et l’égalité}
\NormalTok{                                      Z ∂          Z ∂}
\NormalTok{                                            ∗}
\NormalTok{(5.10.1.2)                                  f α=         ⟨α, µΠf ⟩}
\NormalTok{                                        X           Πf}

\NormalTok{lorsque p = n − 1.}
\NormalTok{Démonstration. — Posons ω = f ∗ α. D’après la relation (5.8.4.2) lorsque p = n et la}
\NormalTok{relation (5.9.4.2) lorsque p = n−1, on a l’égalité (ftrop )∗ (ω̃) = ⟨α, µΠf ⟩ de formes{-}volume}
\NormalTok{p{-}dimensionnelles sur Πf . Le lemme 4.4.6 entraîne alors que f∗ (ω̃Bor ) ⩾ ⟨α, µΠf ⟩Bor .}
\NormalTok{Loreque p = n, on vérifie, en revenant aux définitions et en utilisant la positivité du}
\NormalTok{140                   CHAPITRE 5. ESPACES G{-}TROPICAUX CALIBRÉS}

\NormalTok{calibrage canonique, que l’on a en fait l’égalité, d’où (1). En appliquant le lemme 4.4.6}
\NormalTok{aux mesures de Radon correspondantes, on en déduit (2), puis (3) par intégration.}

\NormalTok{5.10.2. — La formule (5.10.1.1) définit explicitement l’intégrale d’une forme tropicale.}
\end{Highlighting}
\end{Shaded}

\subsection{CLD.txt, newline lines
7070--7101}\label{cld.txt-newline-lines-70707101}

\begin{Shaded}
\begin{Highlighting}[]


\NormalTok{      § 6.2. FORMES LISSES SUR UN ESPACE G{-}TROPICAL AFFINÉ}

\NormalTok{  Soit G une classe de coefficients.}
\NormalTok{Définition 6.2.1. — Soit (p, q) un couple d’entiers naturels. Soit ω une G{-}forme de}
\NormalTok{type (p, q) à coefficients dans G sur un domaine V de X.}
\NormalTok{  On dit que ω est tropicale s’il existe une carte tropicale qui G{-}tropicalise ω.}
\NormalTok{  On dit que ω est une forme (de type (p, q) à coefficients dans G ) si pour tout point x}
\NormalTok{de V , il existe un voisinage ouvert U de x dans V sur lequel ω est tropicale.}
\NormalTok{  Lorsque G est la classe des fonctions lisses, on parle de forme lisse sur l’espace G{-}}
\NormalTok{tropical affiné X. Si de plus p = q = 0, on parle de fonction lisse.}
\NormalTok{6.2.2. — Les formes lisses forment un sous{-}G{-}préfaisceau en R{-}espaces vectoriels}
\NormalTok{de AXp,q qui est un faisceau pour la topologie usuelle. Ces faisceaux sont stables par les}
\NormalTok{opérateurs d′ , d′′ , J et par le produit extérieur.}
\NormalTok{6.2.3. — Un espace G{-}tropical affiné (X, LX ) est dit affinement séparé si pour tout}
\NormalTok{couple (x, y) de points de X avec x ̸= y, il existe une fonction affine f sur X telle que}
\NormalTok{f (x) ̸= f (y). A fortiori, les fonctions lisses sur X séparent les points. On dira alors que}
\NormalTok{X est différentiellement séparé.}
\NormalTok{   Si X est affinement (resp. différentiellement) séparé, il en est de même de tout domaine.}
\NormalTok{   Un espace G{-}tropical affiné (X, LX ) est dit localement affinement (resp. différentielle{-}}
\NormalTok{ment) séparé si chacun de ses points possède un voisinage affinement (resp. différentiel{-}}
\NormalTok{lement) séparé.}
\NormalTok{   Cela entraîne que le faisceau des fonctions lisses sépare localement les points au sens}
\NormalTok{du §A.2.2.}
\NormalTok{Corollaire 6.2.4. — Soit (X, LX ) un espace G{-}tropical affiné qui est localement dif{-}}
\NormalTok{férentiellement séparé. Soit K une partie compacte de X et soit U un voisinage ouvert}
\NormalTok{de K. Il existe une fonction lisse sur X qui vaut 1 au voisinage de K et dont le support}
\NormalTok{est compact et contenu dans U .}
\NormalTok{Démonstration. — C’est une conséquence du corollaire A.2.4.}

\NormalTok{Théorème 6.2.5 (Stone–Weierstraß). — Soit (X, LX ) un espace G{-}tropical affiné}
\end{Highlighting}
\end{Shaded}

\subsection{CLD.txt, newline lines
8130--8168}\label{cld.txt-newline-lines-81308168}

\begin{Shaded}
\begin{Highlighting}[]
\NormalTok{forme positive est positive.                             R∂}
\NormalTok{   On définit aussi le courant [ω]∂ par la formule α 7→ X ω∧α défini à l’aide de l’intégrale}
\NormalTok{de bord des G{-}formes. Il est d’ordre 0 et de bidimension (n − p − 1, n − q). Lorsque ω}
\NormalTok{est une forme lisse, on a [ω]∂ = δX ∂}
\NormalTok{                                      ∧ ω.}
\NormalTok{   Il découle de la définition de la topologie des espaces de courants que les applications}
\NormalTok{ω 7→ [ω] et ω 7→ [ω]∂ sont continues.}
\NormalTok{7.3.5. — Nous allons maintenant généraliser cette construction aux « G{-}formes » qui}
\NormalTok{sont combinaisons linéaires de G{-}formes usuelles à coefficients dans un anneau de fonc{-}}
\NormalTok{tions non nécessairement tropicales.}
\NormalTok{   Soit V un domaine de X et soit U un voisinage ouvert de V dans X dans lequel V}
\NormalTok{est fermé. Soit u ∈ L1loc (V ), ω ∈ AXp,q}
\NormalTok{                                       G}
\NormalTok{                                          (V ) et α ∈ Acn−p,n−q (U ) ; la G{-}forme ω ∧ α|V est à}
\NormalTok{support compact}
\NormalTok{            R     dans V , de sorte que u est intégrable pour la mesure associée. L’appli{-}}
\NormalTok{cation α 7→ V u(ω ∧ α)Rad est un courant d’ordre 0 sur U . La section correspondante de}
\NormalTok{                              § 7.3. COURANTS D’INTÉGRATION                              165}

\NormalTok{π ∗ Dp,q (V ) ne dépend pas du choix de U . On définit ainsi un morphisme de G{-}faisceaux}
\NormalTok{                                L1loc,X ⊗ AXp,q}
\NormalTok{                                             G}
\NormalTok{                                                → π ∗ (Dp,q )X .}
\NormalTok{Par adjonction, on en déduit un morphisme de faisceaux}
\NormalTok{                            π∗ (L1loc,X ⊗ AXp,q}
\NormalTok{                                             G}
\NormalTok{                                                ) → π∗ π ∗ (Dp,q )X .}
\NormalTok{Comme l’unité id → π∗ π ∗ est un isomorphisme de foncteurs, on en déduit un morphisme}
\NormalTok{de faisceaux}
\NormalTok{                               π∗ (L1loc,X ⊗ AXp,q}
\NormalTok{                                                G}
\NormalTok{                                                   ) → (Dp,q )X .}
\NormalTok{   L’image d’une section ω sera notée [ω]. La notation est cohérente avec le cas des G{-}}
\NormalTok{formes lisses ; via l’isomorphisme de la proposition 3.4.11, elle est aussi cohérente avec}
\NormalTok{les cas des G{-}formes à coefficients tropicalement boréliens localement bornés. Ce courant}
\NormalTok{est d’ordre 0 et de bidimension (n − p, n − q).}
\NormalTok{   Par composition avec l’homomorphisme canonique}
\NormalTok{                               1}
\NormalTok{                         π∗ (Lloc,X ⊗ AXp,q}
\end{Highlighting}
\end{Shaded}

\subsection{CLD.txt, newline lines
16809--16927}\label{cld.txt-newline-lines-1680916927}

\begin{Shaded}
\begin{Highlighting}[]
\NormalTok{Définition 12.2.1. — Soit X un espace k{-}analytique ; on appelle moment sur X un}
\NormalTok{morphisme de X vers un tore.}

\NormalTok{   Soit f : X → T un moment sur un espace k{-}analytique X. On désignera par ftrop : X →}
\NormalTok{Ttrop l’application composée X → T → Ttrop . On dira que l’application ftrop est la}
\NormalTok{tropicalisation de f ; c’est une application continue.}

\NormalTok{Remarque 12.2.2. — Si f : T → T ′ est un morphisme affine entre tores, la nota{-}}
\NormalTok{                                                                  ′}
\NormalTok{tion ftrop désigne donc aussi bien l’application induite Ttrop → Ttrop que sa composée}
\NormalTok{        ′     ′}
\NormalTok{T → T → Ttrop avec l’application de tropicalisation de T . Cette ambiguïté n’aura, en}
\NormalTok{pratique, guère de conséquences.}

\NormalTok{Lemme 12.2.3. — Soit X un espace analytique, soit x un point de X et soit V un}
\NormalTok{domaine analytique de X contenant x. Soit f : V → Gm une fonction analytique in{-}}
\NormalTok{versible. On suppose que X est presque bon en x. Il existe alors un voisinage ouvert U}
\NormalTok{de x et une fonction analytique inversible g sur U telle que |f − g| \textless{} |f | sur U ∩ V ; en}
\NormalTok{particulier, ftrop et gtrop coïncident sur U ∩ V .}
\NormalTok{                      § 12.2. MOMENTS, TROPICALISATIONS, POLYÈDRES                         335}

\NormalTok{Démonstration. — Soit U un voisinage ouvert de x dans X et g ∈ OX (U ) tels que}
\NormalTok{|f (x) − g(x)| \textless{} |f (x)| ; il en existe car κ(x) est dense dans H (x) et |f (x)| \textgreater{} 0. L’en{-}}
\NormalTok{semble des points de U ∩ V en lesquels |f − g| \textless{} |f | est un voisinage de x dans V ; quitte}
\NormalTok{à restreindre U , on peut supposer que cette inégalité vaut dans U ∩ V tout entier. Quitte}
\NormalTok{à restreindre encore U , on peut supposer que g est inversible sur U , d’où le lemme.}

\NormalTok{Théorème 12.2.4. — Soit X un espace k{-}analytique compact et soit f : X → T un}
\NormalTok{moment. Le sous{-}espace compact ftrop (X) est un polyèdre de Ttrop de dimension dim(X).}
\NormalTok{   L’assertion a été établie par Berkovich dans le cas où X est G{-}localement algébrisable,}
\NormalTok{cf. [11], cor. 6.2.2. Pour le cas général, on pourra se reporter à [46], §3.30, ou à [50],}
\NormalTok{th. 3.2.}
\NormalTok{Corollaire 12.2.5. — Soit X un espace k{-}analytique compact et soit V une partie}
\NormalTok{fermée de X munie d’une structure faiblement analytique de dimension d. Soit f : X → T}
\NormalTok{un moment. Alors, ftrop (V ) est une partie paralinéaire de dimension ⩽ d de Ttrop .}
\NormalTok{Démonstration. — On se ramène immédiatement au cas où il existe un morphisme}
\NormalTok{g : X → Y et un point y ∈ Y tel que V = g −1 (y). Alors, le produit fibré X ×Y H (y)}
\NormalTok{permet de considérer V comme un espace H (y){-}analytique compact, de dimension ⩽ d,}
\NormalTok{et f |V : V → TH (y) comme un moment sur cet espace. On donc dim(ftrop (V )) ⩽ d.}

\NormalTok{Corollaire 12.2.6. — Soit X un espace k{-}analytique compact de dimension n. Soit}
\NormalTok{f : X → T un moment. Si X est affinoïde, alors ftrop (∂(X)) est une partie paralinéaire}
\NormalTok{de dimension ⩽ n − 1 de Ttrop ; en genéral, ftrop (∂(X)) est contenu dans une telle partie}
\NormalTok{paralinéaire.}
\NormalTok{Démonstration. — Compte tenu de l’exemple 11.10.2 selon lequel ∂(X) est une partie}
\NormalTok{faiblement analytique de dimension ⩽ dim(X) − 1, cela résulte du corollaire précédent.}
\NormalTok{Il s’agit en fait du théorème 3.2 de [50], démontré de la même façon, même si la notion}
\NormalTok{de structure faiblement analytique n’avait pas été dégagée.}

\NormalTok{12.2.7. — Soit X un espace k{-}analytique et f : X → T un moment. Soit x ∈ X. On}
\NormalTok{appellera degré de transcendance résiduel (gradué) de f en x le degré de transcendance}
\NormalTok{résiduel torique du point f (x) de T . Si χ1 , . . . , χn est une base de X(T ), ce degré est}
\NormalTok{le degré de transcendance sur ek du sous{-}corps gradué de H         \^{}  (x) engendré par e}
\NormalTok{                                                                                       k et les}
\NormalTok{χi\^{}◦ f (x). Si T = Gn et f = (f1 , . . . , fn ), c’est donc le degré de transcendance sur e}
\NormalTok{                     m                                                                    k du}
\NormalTok{sous{-}corps gradué de H \^{} (x) engendré par e k et les f]}
\NormalTok{                                                      i (x).}
\NormalTok{   Le degré de transcendance résiduel de f en x est majoré à la fois par dk (x) et par n.}
\NormalTok{Il vaut n si et seulement si f (x) ∈ S(T ).}
\NormalTok{Proposition 12.2.8. — Soit X un espace k{-}analytique et soit x un point de X. Soit}
\NormalTok{f = X → T un moment ; posons ξ = ftrop (x). Notons d le degré de transcendance}
\NormalTok{résiduel de f en x.}
\NormalTok{   Il existe un voisinage analytique compact U de x dans X tel que pour pour tout voi{-}}
\NormalTok{sinage analytique compact V de x contenu dans U , on ait les propriétés suivantes :}
\NormalTok{   i) Les germes de polyèdres (ftrop (U ), ξ) et (ftrop (V ), ξ) coïncident ;}
\NormalTok{   ii) La dimension locale dimξ (ftrop (V )) est majorée par d, et a fortiori par dk (x).}
\NormalTok{336           CHAPITRE 12. PARTIES SQUELETTIQUES D’UN ESPACE ANALYTIQUE}

\NormalTok{   iii) Si x ̸∈ ∂(X), alors ftrop (V ) est purement de dimension d en ξ. Si, de plus, d = n,}
\NormalTok{alors ftrop (V ) est un voisinage de ξ dans Ttrop .}
\NormalTok{   iv) Si dimξ (ftrop (V )) = n, alors f (x) ∈ S(T ) ; inversement, si f (x) ∈ S(T ) et si de}
\NormalTok{plus x ∈ / ∂(X), alors dimξ (ftrop (V )) = n.}
\NormalTok{Démonstration. — L’existence d’un voisinage analytique U vérifiant les i), ii) et iii) est}
\NormalTok{le théorème 3.4 de [50]. Montrons (iv). Si dimξ (ftrop (V )) = n, il résulte de (ii) que n ⩽ d,}
\NormalTok{donc que d = n et que f (x) ∈ S(T ). Réciproquement, si f (x) ∈ S(T ) alors d = n et si}
\NormalTok{de plus x ∈}
\NormalTok{          / ∂X il résulte de (iii) que dimξ (ftrop (V )) = n.}

\NormalTok{Remarque 12.2.9. — Soit X un espace k{-}analytique et soit x ∈ X. Soit f : X → Gnm}
\NormalTok{un moment. Si V est un voisinage analytique compact de x dans X, il résulte de la}
\NormalTok{proposition ci{-}dessus que le germe de ftrop (V ) en x ne dépend pas de V . Nous le noterons}
\NormalTok{ftrop (X, x).}
\NormalTok{   La preuve du théorème 3.4 de [50] en donne une description en termes de réduction}
\NormalTok{de Temkin, que nous allons rappeler ici, en supposant que le germe est (X, x) séparé.}
\NormalTok{   Le théorème 10.5.7 assure (en reprenant le vocabulaire et les notations de 10.5.6)}
\NormalTok{l’existence d’un sous{-}foncteur D de Dn\{1\} tel que pour tout H ∈ D et tout n{-}uplet}
\NormalTok{(h1 , . . . , hn ) d’éléments de H, les assertions suivantes soient équivalentes :}
\NormalTok{   (i) h ∈ D(H) ;}
\NormalTok{   (ii) il existe H ′ ∈ DH et une valuation w appartenant à (X,   \^{}   x) à valeurs dans H ′ telle}
\NormalTok{que w(f]     i (x)) = hi pour tout i.}
\NormalTok{Comme le groupe des paramètres est ici \{1\}, le foncteur D est simplement définiQpar}
\NormalTok{conjonctions et disjonctions d’inégalités de la forme φ ⩽ 1 où φ est un monôme xri i}
\NormalTok{à exposants rationnels. L’ensemble D(R×      + ) est alors un cône rationnel, et ftrop (X, x) =}
\NormalTok{                ×}
\NormalTok{(ftrop (x) · D(R+ ), ftrop (x)) en notation multiplicative.}
\NormalTok{                                                      \^{}}
\NormalTok{  En particulier, si x est intérieur, c’est{-}à{-}dire si (X, x) = P       , ce germe de polytope}
\NormalTok{                                                                 H}
\NormalTok{                                                                 \^{} (x)/e}
\NormalTok{                                                                       k}

\NormalTok{ne dépend que de la classe d’isomorphie de l’extension graduée F de e       k dans H\^{}  (x)}
\NormalTok{engendrée par les f]}
\NormalTok{                   i (x), munie de cette famille. Dans ce cas, il est équidimensionnel, de}
\NormalTok{dimension égale au degré de transcendance (gradué) de F sur e    k.}
\NormalTok{Définition 12.2.10. — Soit X un espace analytique et soit f : X → T un moment.}
\NormalTok{Soit x un point de X. On appelle dimension tropicale de f en x la borne inférieure de di{-}}
\NormalTok{mensions des polyèdres ftrop (V ), où V est un domaine analytique compact qui contient x.}
\NormalTok{On la note dimtrop,x (f ).}
\NormalTok{  Soit d le degré de transcendance résiduel de f en x. D’après le théorème 12.2.4 et la}
\NormalTok{proposition 12.2.8, on a}
\NormalTok{(12.2.10.1)                 dimtrop,x (f ) ⩽ d ⩽ dk (x) ⩽ dimx (X),}
\NormalTok{et dimtrop,x (f ) = d si x ∈}
\NormalTok{                           / ∂X. Si X est compact, il résulte de la définition que}
\NormalTok{(12.2.10.2)                   dim(ftrop (X)) = sup dimtrop,x (f ).}
\NormalTok{                                                x∈X}

\NormalTok{Proposition 12.2.11. — Soit X un espace analytique, soit x un point de X ∂(X).}
\end{Highlighting}
\end{Shaded}

\subsection{CLD.txt, newline lines
18040--18119}\label{cld.txt-newline-lines-1804018119}

\begin{Shaded}
\begin{Highlighting}[]
\NormalTok{      § 12.8. POLYÈDRES CARACTÉRISTIQUES ET DIMENSION}
\NormalTok{                            TROPICALE}

\NormalTok{12.8.1. — Soit X un espace k{-}analytique de niveau n, soit m un entier et soit f =}
\NormalTok{X → Gm   m un moment sur X. On note (f1 , . . . , fm ) les composantes de f . Pour toute}
\NormalTok{partie I de \{1, . . . , m\}, on note qI la projection canonique de Gm      I}
\NormalTok{                                                                   m sur Gm .}
\NormalTok{  Rappelons (proposition 12.6.8) que le polyèdre caractéristique de f est une partie}
\NormalTok{paralinéaire de X et que l’application ftrop |Σf est une immersion par morceaux d’espaces}
\NormalTok{paralinéaires.}
\NormalTok{  Le but de ce numéro est de relier la dimension tropicale de ftrop aux propriétés de Σf .}
\NormalTok{Théorème 12.8.2. — Soit x un point de X.}
\NormalTok{                                                                    (n)}
\NormalTok{   a) On a dimtrop,x (f ) ⩽ n, avec égalité si et seulement si x ∈ Σf .}
\NormalTok{   b) Supposons x ∈ Int(X). Les trois assertions suivantes sont équivalentes : (i) x ∈ Σf ;}
\NormalTok{                                                        (n)}
\NormalTok{(ii) f est de dimension tropicale n en x ; (iii) x ∈ Σf .}
\NormalTok{   c) Si X est compact ou plus généralement si l’application ftrop est topologiquement}
\NormalTok{                                    (n)}
\NormalTok{propre, on a ftrop (X)(n) = ftrop (Σf ).}
\NormalTok{Démonstration. — a) Soit d le degré de transcendance résiduel de f en x. On a}
\NormalTok{dimtrop,x (f ) ⩽ d ⩽ dimx X ((12.2.10.1)). Comme dimx X ⩽ n, il vient dimtrop,x (f ) ⩽}
\NormalTok{d ⩽ n.}
\NormalTok{                                                              (n)}
\NormalTok{   Supposons que dimtrop,x (f ) = n, et montrons que x ∈ Σf . Puisque dimtrop,x (f ) ⩽}
\NormalTok{d ⩽ n, on a d = n. Il existe donc un sous{-}ensemble I de cardinal d de \{1, . . . , m\}}
\NormalTok{tels que les f] i (x) soient algébriquement indépendants sur k. ll s’ensuit que x ∈ Σf . Il}
\NormalTok{                                                             e}
\NormalTok{reste à vérifier que dimx (Σf ) = n. Soit V un voisinage analytique compact arbitraire}
\NormalTok{de x dans X et démontrons que dim(Σf ∩ V ) = n. Comme dimtrop,x (f ) = n, on a}
\NormalTok{dim(ftrop (V )) = n. Soit C une cellule compacte de dimension n de ftrop (V ) et soit}
\NormalTok{             § 12.8. POLYÈDRES CARACTÉRISTIQUES ET DIMENSION TROPICALE                     357}

\NormalTok{q : Gm        n}
\NormalTok{      m → Gm un morphisme de tores tel que la restriction de qtrop à ⟨C⟩ soit injective.}
\NormalTok{                                                                            (n)}
\NormalTok{Alors, dim(q ◦ f )trop (V ) = n. D’après le lemme 12.2.13, (q ◦ f )trop (V )sq est contenu dans}
\NormalTok{(q ◦f )(V ), donc dans (q ◦f )(Σf ∩V ). Comme dim(Σf ∩V ) ⩾ n, il vient dim(Σf ∩V ) = n,}
\NormalTok{ce qu’il fallait démontrer.}
\NormalTok{                                          (n)}
\NormalTok{   Inversement, supposons que x ∈ Σf et montrons que dimtrop,x (f ) = n. Soit V un}
\NormalTok{                                         (n)                                (n)}
\NormalTok{voisinage compact de x. Alors, V ∩ Σf est un voisinage de x dans Σf , donc contient}
\NormalTok{une cellule de dimension n puisque dimx Σf = n par hypothèse. Comme la restriction}
\NormalTok{               (n)}
\NormalTok{de ftrop à Σf est une immersion par morceaux, cela prouve que dim(ftrop (V )) ⩾ n,}
\NormalTok{et donc que dim(ftrop (V )) = n puisque dim V ⩽ n. Comme V est arbitraire, il vient}
\NormalTok{dimtrop,x (f ) = n.}
\NormalTok{   b) Faisons l’hypothèse que x ∈ Int(X). Supposons x ∈ Σf . Il existe alors un sous{-}}
\NormalTok{ensemble I de \{1, . . . , m\} tels que les f]   i (x) pour i ∈ I soient algébriquement indépen{-}}
\NormalTok{dantes sur k ; par conséquent, d ⩾ n, si bien que d = n puisqu’on sait que d ⩽ n. Et}
\NormalTok{              e}
\NormalTok{comme x ∈ Int(X), la dimension tropicale de f en x est égale à d (12.2.10.1), et donc}
\NormalTok{à n.}
\NormalTok{   Supposons réciproquement que x est de dimension tropicale n en x. Le point x appar{-}}
\NormalTok{                  (n)}
\NormalTok{tient alors à Σf d’après (a).}
\NormalTok{   L’implication (iii)⇒(i) est évidente.}
\NormalTok{   c) Supposons ftrop topologiquement propre. En raisonnant G{-}localement sur Rm on}
\NormalTok{                                           (n)}
\NormalTok{voit que ftrop (X), ftrop (Σf ) et ftrop (Σf ) sont des parties paralinéaires fermées de Rm .}
\NormalTok{                                                                                           (n)}
\NormalTok{Comme ftrop induit une immersion par morceaux de Σf sur ftrop (Σf ), l’image ftrop (Σf )}
\NormalTok{est purement de dimension n et est donc contenue dans ftrop (X)(n) .}
\NormalTok{                                                                         (n)}
\NormalTok{  Soit maintenant ξ un point de ftrop (X) qui n’appartient pas à ftrop (Σf ). Soit x un}
\NormalTok{                                                     (n)}
\NormalTok{point de X tel que ftrop (x) = ξ ; on a donc x ̸∈ Σf . D’après a), ftrop est de dimension}
\NormalTok{tropicale \textless{} n en x. Comme ftrop est topologiquement propre, il existe un voisinage}
\NormalTok{                             −1}
\NormalTok{analytique compact V de ftrop   (ξ) et un voisinage paralinéaire U de ξ dans ftrop (X) tels}
\NormalTok{      −1}
\NormalTok{que ftrop (U ) ⊂ V et que dim(ftrop (V )) ⩽ n−1. Puisque U ⊂ ftrop (V ), il est de dimension}
\NormalTok{au plus n − 1 et ξ n’appartient à ftrop (X)(n) .}

\NormalTok{Corollaire 12.8.3. — Soit X un espace k{-}analytique de niveau n et soit f : X → T un}
\NormalTok{                  (n)                                             (n)}
\NormalTok{moment. Alors Σf coïncide avec le polyèdre de variation maximale Σftrop de l’application}
\NormalTok{paralinéaire ftrop (définition 5.2.1).}
\NormalTok{Démonstration. — C’est l’assertion a) du théorème 12.8.2.}
\end{Highlighting}
\end{Shaded}

\subsection{CLD.txt, newline lines
19153--19240}\label{cld.txt-newline-lines-1915319240}

\begin{Shaded}
\begin{Highlighting}[]
\NormalTok{Lemme 13.2.8 (Longueur générique et extension des scalaires)}
\NormalTok{    Soit X un espace k{-}analytique et soit K une extension complète de k. Soit F un}
\NormalTok{faisceau cohérent sur X. Soit Z une composante irréductible de X, soit ℓ la longueur}
\NormalTok{générique de X le long de Z, soit ℓK la longueur générique de FK le long de ZK , et soit}
\NormalTok{µ la multiplicité générique de Xred,K le long de ZK . On a alors ℓK = µℓ.}
\NormalTok{                                                   gen     gen}
\NormalTok{Démonstration. — Choisissons un point y sur ZK         ⊂ XK    dont on note x l’image sur}
\NormalTok{X ; par construction, x ∈ Z. L’anneau local OXK,G ,y est artinien car y ∈ XK    gen}
\NormalTok{                                                                                    , et il est}
\NormalTok{plat sur OXG ,x (13.1.5) ; ce dernier est donc également artinien, ce qui montre que x}
\NormalTok{appartient à X gen , et donc à Z gen car x ∈ Z.}
\NormalTok{   Pour alléger les notations posons A = OXG ,x , B = OXK,G ,y et M = FXG ,x . On note m}
\NormalTok{l’idéal maximal de A et κ = A/m son corps résiduel. Enfin on pose B ′ = B/mB = κ⊗A B.}
\NormalTok{Remarquons que κ = OXred,G ,x et B ′ = OXred,K,G ,y (proposition 13.1.9 d ) et c)).}
\NormalTok{   Par définition de ℓ il existe une filtration}

\NormalTok{                             \{0\} = M0 ⊂ M1 ⊂ . . . ⊂ Mℓ = M}

\NormalTok{du A{-}module M dont les quotients successifs sont A{-}isomorphes à κ. Comme B est plat}
\NormalTok{sur A, cette filtration en induit une}

\NormalTok{                \{0\} = B ⊗A M0 ⊂ B ⊗A M1 ⊂ . . . ⊂ B ⊗A Mℓ = B ⊗A M}

\NormalTok{dans la catégorie des B{-}modules dont les quotients successifs sont tous B{-}isomorphes à}
\NormalTok{B ′ . Il vient}
\NormalTok{                    ℓK = ℓB (B ⊗A M ) = ℓ · ℓB (B ′ ) = ℓ · ℓB ′ (B ′ ) = ℓµ.}
\NormalTok{378              CHAPITRE 13. COMPLÉMENTS DE GÉOMÉTRIE ANALYTIQUE}

\NormalTok{§ 13.3. DEGRÉ ET DEGRÉ CALIBRÉ D’UN MORPHISME QUASI{-}FINI}

\NormalTok{13.3.1. — Soit f : Y → X un morphisme d’espaces k{-}analytiques. Soit y ∈ Y . Si}
\NormalTok{dimy (f −1 (f (y)) \textgreater{} 0, on pose degy (f ) = +∞. Sinon, le point y est isolé dans sa fibre}
\NormalTok{f −1 (f (y)), l’anneau local de f −1 (f (y)) en y est une H (f (y)){-}algèbre de dimension finie}
\NormalTok{et on note degy (f ) sa dimension.}
\NormalTok{   Lorsque f est fini et plat en y, cet entier degy (f ) coïncide avec l’entier degy (f ) défini}
\NormalTok{en 11.6.3.}
\NormalTok{   Nous dirons que degy (f ) est le degré de f en y.}
\NormalTok{   Soit V un domaine analytique de Y contenant y et soit U un domaine analytique}
\NormalTok{de X tel que f (V ) ⊂ U ; soit g : V → U le morphisme déduit de f . Pour tout y ∈ V , on}
\NormalTok{a degy (f ) = degy (g), car pour tout y ∈ V d’image x ∈ X, les anneaux locaux en y de}
\NormalTok{f −1 (f (y)) et de g −1 (g(y)) coïncident.}
\NormalTok{   Soit Z un sous{-}espace analytique fermé de X contenant x et soit g : Y ×X Z → Z le}
\NormalTok{morphisme déduit de f par changement de base. Comme les espaces H (x){-}analytiques}
\NormalTok{f −1 (x) et g −1 (x) coïncident, on a degy (f ) = degy (g).}
\NormalTok{13.3.2. — Soit f : Y → X un morphisme d’espaces k{-}analytiques. Pour tout x ∈ X,}
\NormalTok{on définit de le degré de f au{-}dessus de x comme la somme}
\NormalTok{                                            X}
\NormalTok{                               degx (f ) =     degy (f ).}
\NormalTok{                                              y∈f −1 (x)}


\NormalTok{Il est fini si et seulement si f −1 (x) est compact et de dimension nulle ; alors f −1 (x) est}
\NormalTok{un espace analytique fini sur H (x), et degx (f ) = dimH (x) (O(f −1 (x))).}
\NormalTok{   Lorsque f est globalement fini et plat, cet entier degx (f ) coïncide avec l’entier degx (f )}
\NormalTok{défini en 11.6.1.}

\NormalTok{Lemme 13.3.3. — Soit f : Y → X un morphisme d’espaces analytiques. Soit K une}
\NormalTok{extension valuée complète de k ; notons p : XK → X et q : YK → Y les morphismes de}
\NormalTok{changement de base, et fK : YK → XK le morphisme déduit de f par changement de}
\NormalTok{base. Soit xK ∈ p−1 (x).}
\NormalTok{  a) Pour tout y ∈ Y d’image x = f (y) ∈ X, tout point xK ∈ p−1 (x), on a}
\NormalTok{                                        X}
\NormalTok{                         degy (f ) =              degz (fK ).}
\NormalTok{                                           −1}
\NormalTok{                                        z∈fK  (xK )∩q −1 (y)}


\NormalTok{  b) Pour tout x ∈ X et tout point xK ∈ p−1 (x), on a degx (f ) = degxK (fK ).}

\NormalTok{Démonstration. — Tout se passe dans les fibres f −1 (x) et fibre fK−1 (xK ) = f −1 (x) ×}
\NormalTok{M (H (xK )). Cela permet de supposer que X = M (k) = \{x\} et XK = M (K) = \{xK \}.}
\NormalTok{  Prouvons l’égalité a). Si dimy (Y ) \textgreater{} 0, il en est de même de dimz (YK ) pour tout}
\NormalTok{antécédent z de y, si bien que les deux membres de l’égalité sont infinis. Supposons}
\NormalTok{maintenant dimy (Y ). Alors OY,y est une k{-}algèbre artinienne et OY,y ⊗k K = z OYK ,z ,}
\NormalTok{                                                                              Q}
\NormalTok{où z parcourt l’ensemble des antécédents de y dans YK . L’égalité s’ensuit.}
\NormalTok{  L’égalité b) s’en déduit en sommant sur tous les points de Y .}
\NormalTok{               § 13.3. DEGRÉ ET DEGRÉ CALIBRÉ D’UN MORPHISME QUASI{-}FINI                    379}

\NormalTok{13.3.4. — Soit A un anneau local artinien et soit M un A{-}module de longueur finie.}
\NormalTok{Nous noterons δA (M ) le nombre rationnel ℓA (M )/ℓA (A). Indiquons quelques propriétés}
\NormalTok{élémentaires de cet invariant.}
\NormalTok{   Si A est réduit, c’est un corps, on a ℓA (A) = 1 et δA (M ) est donc simplement ℓA (M ),}
\NormalTok{c’est{-}à{-}dire la dimension de M comme A{-}espace vectoriel.}
\NormalTok{   Supposons M plat sur A, c’est{-}à{-}dire libre comme A{-}module ; soit r son rang. Comme}
\NormalTok{M ≃ Ar , ℓA (M ) = rℓA (A), si bien que δA (M ) = r.}
\NormalTok{   Si M = Ared , alors ℓA (M ) = 1 et δA (M ) = 1/ℓA (A).}
\NormalTok{   Soit B une A{-}algèbre finie. Soit κ le corps résiduel de A et λ celui de B. Choisis{-}}
\end{Highlighting}
\end{Shaded}

\subsection{CLD.txt, newline lines
18355--18407}\label{cld.txt-newline-lines-1835518407}

\begin{Shaded}
\begin{Highlighting}[]
\NormalTok{Théorème 12.9.11. — Soit X un espace k{-}analytique de niveau n. Soit L une exten{-}}
\NormalTok{sion complète de k et notons p : XL → X la projection canonique. Soit Σ une partie sque{-}}
\NormalTok{lettique de X ; notons ΣL la réunion des polyèdres caractéristiques ΣfL , où f : V → T}
\NormalTok{parcourt l’ensemble des moments définis sur un domaine analytique V de X tels que}
\NormalTok{Σf ⊂ Σ.}
\NormalTok{   L’ensemble ΣL est une partie squelettique de XL et le morphisme p induit une immer{-}}
\NormalTok{sion par morceaux surjective de ΣL sur Σ qui préserve la dimension locale.}
\NormalTok{   De plus, pour tout moment f : V → T tel que Σf ⊂ Σ, on a ΣL ×Σ Σf = ΣfL .}

\NormalTok{Démonstration. — La première assertion est G{-}locale sur X ; comme elle est vérifiée}
\NormalTok{pour l’application ΣfL → Σf , il suffit de démontrer la seconde assertion.}
\NormalTok{   Soit f : V → T un moment défini sur un domaine analytique V de X. Démontrons}
\NormalTok{l’égalité PL ×P Σf = ΣfL . Cela revient à démontrer que si g : W → T ′ est un moment et}
\NormalTok{x ∈ Σf ∩Σg , alors les fibres de ΣfL et ΣgL au{-}dessus de x coïncident. Quitte à remplacer X}
\NormalTok{par V ∩ W , on suppose que les moments f et g sont définis sur X.}
\NormalTok{   Soit x ∈ Σf ∩ Σg . Soit C la complétion d’une clôture algébrique de k et soit M une}
\NormalTok{extension complète composée de L et C. Considérons le diagramme}

\NormalTok{                                               XM   →}
\NormalTok{                                                           σ}
\NormalTok{                                         q′}
\NormalTok{                                         →    ←      ←}
\NormalTok{                                                       →}
\NormalTok{                                                           ←}

\NormalTok{                                    XL                  XC}
\NormalTok{                                         ←             ←}
\NormalTok{                                         p →}
\NormalTok{                                                   →   q}

\NormalTok{                                               X}
\NormalTok{dans lequel σ est la section de Poineau de l’application XM → XC dont l’existence}
\NormalTok{résulte de ce que C est algébriquement clos. Comme l’extension C/k contient une sous{-}}
\NormalTok{extension algébrique dense, on a ΣfC = q −1 (Σf ). Comme C est algébriquement clos, on}
\NormalTok{a ΣfM = σ(ΣfC ). Enfin, on a ΣfL = q ′ (ΣfM ). Par suite,}
\NormalTok{                                  ΣfL = q ′ (σ(q −1 (Σf ))),}
\NormalTok{donc la fibre de ΣfL au{-}dessus de x est donnée par la formule}
\NormalTok{                                   ΣfL ,x = q ′ (σ(q −1 (x)))}
\NormalTok{dans laquelle Σf n’intervient plus ; cette fibre coïncide donc avec la fibre de ΣgL au{-}dessus}
\NormalTok{de x. Cela conclut la démonstration du théorème.}
\NormalTok{  Le cas des polyèdres caractéristiques et la définition de ΣL entraînent de plus la}
\NormalTok{caractérisation suivante :}

\NormalTok{Proposition 12.9.12. — Conservons les notations du théorème 12.9.11.}
\NormalTok{   a) Si l’extension L/k possède une sous{-}extension algébrique dense, alors ΣL = p−1 (Σ) ;}
\NormalTok{   b) Si l’extension L/k est universellement multiplicative, ou si le corps le corps k est}
\NormalTok{algébriquement clos, on a ΣL = σL (Σ).}

\NormalTok{Démonstration. — Compte tenu de la seconde assertion du théorème, l’assertion a)}
\NormalTok{découle de la proposition 12.7.2, et l’assertion b) du corollaire 12.7.5.}
\NormalTok{                  § 12.9. PARTIES SQUELETTIQUES D’UN ESPACE ANALYTIQUE                         363}

\NormalTok{Lemme 12.9.13. — Soit X un espace k{-}analytique de niveau n, soit Σ une partie}
\end{Highlighting}
\end{Shaded}

\subsection{CLD.txt, newline lines
20320--20360}\label{cld.txt-newline-lines-2032020360}

\begin{Shaded}
\begin{Highlighting}[]
\NormalTok{                           § 14.2. MOMENTS ET DEGRÉS}

\NormalTok{Lemme 14.2.1. — Soit X un espace analytique de niveau n, soit f : X → Gnm un}
\NormalTok{moment. Alors, f est quasi{-}fini et plat en tout point de Σf .}
\NormalTok{Démonstration. — Soit x ∈ Σf . On a f (x) ∈ S(Gnm ), donc dk (f (x)) = n. D’après}
\NormalTok{le lemme 11.5.1, le morphisme f est quasi{-}fini en x. Comme Gnm est réduit, le théo{-}}
\NormalTok{rème 10.3.7 de [52] entraîne que f est plat en x.}

\NormalTok{Proposition 14.2.2. — Soit X un espace analytique de niveau n, soit f : X → T un}
\NormalTok{moment sur X et soit C une f {-}cellule commode de dimension n de X. Soit q : T → Gnm}
\NormalTok{un morphisme de tores tel que qtrop |ftrop (C) est injectif. Alors le morphisme q ◦ f est de}
\NormalTok{dimension relative nulle en tout point de C et il existe un unique entier d \textgreater{} 0 tel que ce}
\NormalTok{morphisme soit fini et plat de degré d en tout point de C̊.}
\NormalTok{Démonstration. — Soit q : T → Gnm un morphisme de tores tel que qtrop soit injectif sur}
\NormalTok{ftrop (C). La dimension tropicale de q ◦ f est égale à n en tout point de C, donc C ⊂}
\NormalTok{(q ◦f )−1 (S(Gnm )) d’après la prop. 12.2.8, (iv). D’après le lemme 11.5.1, le morphisme q ◦f}
\NormalTok{est de dimension relative nulle en tout point de C.}
\NormalTok{   Comme C̊ ne rencontre pas ∂(X), il résulte du lemme 11.6.5 que le morphisme q ◦f est}
\NormalTok{fini et plat en tout point de C̊ ; on a noté degx (q ◦ f ) son degré et il s’agit de démontrer}
\NormalTok{que l’entier degx (q ◦ f ) est indépendant du point x ∈ C̊.}
\NormalTok{   Soit D la cellule sq(qtrop (ftrop (C))) de Gnm . Si un point x ∈ Int(X) vérifie q ◦ f (x) ∈ D̊,}
\NormalTok{                                                                                 (n)}
\NormalTok{alors q◦f est de dimension tropicale n en x, donc f aussi, si bien que x ∈ Σf . Autrement}
\NormalTok{dit, on a les inclusions}
\NormalTok{                                                                (n)}
\NormalTok{                             C̊ ⊂ (q ◦ f )−1 (D̊) ∩ Int(X) ⊂ Σf .}
\NormalTok{                                                                        (n)}
\NormalTok{   Soit maintenant x ∈ C̊. Par hypothèse, C̊ est ouvert dans Σf ; il est donc ouvert}
\NormalTok{dans (q ◦ f )−1 (D̊) ∩ Int(X). Comme q ◦ f est fini et plat en x, il existe un voisinage}
\NormalTok{ouvert U de x dans Int(X) et un voisinage ouvert V de q ◦ f (x) dans Gnm tels que}
\NormalTok{q ◦ f induise un morphisme fini et plat de degré degx (q ◦ f ) de U sur V , et tel que}
\NormalTok{(q ◦ f )−1 (D̊) ∩ U ⊂ C̊. Un point y ∈ C̊ ∩ U est alors le seul antécédent dans U de}
\NormalTok{q ◦ f (y) par q ◦ f ; on a donc degy (q ◦ f ) = degx (q ◦ f ) pour tout y ∈ C̊ ∩ U . La}
\NormalTok{fonction x 7→ degx (q ◦ f ) est localement constante sur l’ensemble C̊ ; comme celui{-}ci est}
\NormalTok{connexe, elle est constante.}
\NormalTok{400           CHAPITRE 14. CALIBRAGE CANONIQUE D’UN ESPACE ANALYTIQUE}

\NormalTok{Proposition 14.2.3. — Soit X un espace k{-}analytique paracompact de niveau n et}
\NormalTok{soit f : X → Gnm un moment. Soit Σ un sous{-}espace paralinéaire paracompact de Σf .}
\NormalTok{Il existe une décomposition cellulaire C de Σ en f {-}cellules telle que pour toute cellule}
\NormalTok{C ∈ C , l’application x 7→ degx (f ) soit finie et constante sur C̊.}
\end{Highlighting}
\end{Shaded}

\subsection{CLD.txt, newline lines
20700--20870}\label{cld.txt-newline-lines-2070020870}

\begin{Shaded}
\begin{Highlighting}[]
\NormalTok{de l’anneau local de (g ′ )−1 (g ′ (y)) en g ′ (y). Considérons X × Gdm comme Gdm,X ; alors}
\NormalTok{406           CHAPITRE 14. CALIBRAGE CANONIQUE D’UN ESPACE ANALYTIQUE}

\NormalTok{g ′ (y) s’identifie au point g|Yx (y) de Gdm,H (x) . De même, la fibre (g ′ )−1 (g ′ (y)) s’identifie à}
\NormalTok{la fibre (g|Yx )−1 (g|Yx (y)). Il en résulte l’égalité degy (g ′ ) = degy (g|Yx ).}
\NormalTok{                                   ′                                            ′}
\NormalTok{     Démontrons enfin que degg (y) (f ′ ) = degx (f ). Par définition, degg (y) (f ′ ) est la dimen{-}}
\NormalTok{sion sur f ′ (g ′ (y)) = h(y) de l’anneau local de (f ′ )−1 (h(y)) en g ′ (y). Le morphisme f ′}
\NormalTok{est déduit de f par changement de base. Ainsi, (f ′ )−1 (h(y)) = f −1 (f (x)) ⊗H (f (x))}
\NormalTok{H (h(y)). L’anneau local de f −1 (f (x)) en x est une algèbre locale artinienne de di{-}}
\NormalTok{mension degx (f ) sur H (x). Son produit tensoriel par H (h(y)) est le produit des an{-}}
\NormalTok{neaux locaux de (f ′ )−1 (h(y)) en tous les antécédents de x. Pour conclure la démons{-}}
\NormalTok{tration de l’égalité voulue, il suffit de démontrer que le point x n’a qu’un seul an{-}}
\NormalTok{técédent dans (f ′ )−1 (h(y)). Cela résulte du fait que H (f (x)) est séparablement clos}
\NormalTok{dans H (h(y)) : le point h(y) appartient à S(Gnm,H (f (x)) ), le corollaire 12.1.8 s’applique}
\NormalTok{donc.}



\NormalTok{                        § 14.3. CALIBRAGES CANONIQUES}

\NormalTok{  Soit X un espace analytique de niveau n.}
\NormalTok{Proposition 14.3.1. — Soit f : X → T un moment et soit C une f {-}cellule de dimen{-}}
\NormalTok{                                                                                 ′}
\NormalTok{sion n. Soit q et q ′ des morphismes de tores de T dans Gnm tels que qtrop et qtrop  soient in{-}}
\NormalTok{jectifs sur ftrop (C). Soit σq l’unique section affine de qtrop dont l’image contient ftrop (C) ;}
\NormalTok{on définit σq′ de façon analogue.}
\NormalTok{   Soit µ un vecteur{-}volume sur Rn . Pour presque tout x ∈ C, on a l’égalité}
\NormalTok{(14.3.1.1)                  degx (q ◦ f )σq,∗ (µ) = degx (q ′ ◦ f )σq′ ,∗ (µ)}
\NormalTok{de vecteurs volumes sur l’espace affine ⟨ftrop (C)⟩.}
\NormalTok{Démonstration. — Notons Cf = ftrop (C).}
\NormalTok{  Supposons d’abord que q ′ soit de la forme φ ◦ q, où φ : Gnm → Gnm est un morphisme}
\NormalTok{surjectif de tores. D’après la proposition 13.3.10, on a}
\NormalTok{                  degx (q ′ ◦ f ) = degx (φ ◦ q ◦ f ) = degq◦f (x) (φ) degx (q ◦ f ).}
\NormalTok{Le morphisme φ est fini et plat et pour tout t ∈ S(Gnm ), on a degt (φ) = deg(φ), en vertu}
\NormalTok{du corollaire 12.1.7. Par suite, on a degx (q ′ ◦ f ) = deg(φ) degx (q ◦ f ) pour tout x ∈ C.}
\NormalTok{   Par ailleurs, φtrop est une application affine bijective de Rn dans lui{-}même et l’on a}
\NormalTok{σq′ = σq ◦φ−1                                                                    →}
\NormalTok{                                                                                 −}
\NormalTok{              trop . La valeur absolue du déterminant de l’application linéaire φ trop est égale}
\NormalTok{au degré de φ. Par suite, φ−1                                                 −1}
\NormalTok{                                trop applique le vecteur volume µ sur deg(φ) µ, de sorte que}
\NormalTok{σq′ ,∗ (µ) = deg(φ)−1 σq,∗ (µ). La relation voulue est donc vérifiée dans ce cas particulier.}
\NormalTok{                               ′}
\NormalTok{   Supposons ensuite que qtrop    − qtrop soit constante sur Cf . D’après le théorème 14.2.7,}
\NormalTok{          x                 x ′}
\NormalTok{on a deg (q ◦ f ) = deg (q ◦ f ) pour tout x ∈ C ∩ Int(X) ; d’autre part, C ∩ ∂(X)}
\NormalTok{est localement négligeable. De plus, σq et σq′ diffèrent d’une constante, si bien que}
\NormalTok{σq,∗ (µ) = σq′ ,∗ (µ). Cela prouve la relation voulue dans ce cas.}
\NormalTok{   Montrons maintenant comment la conjonction de ces deux cas permet d’obtenir le}
\NormalTok{cas général. Notons ψ l’unique application linéaire de Rn dans lui{-}même telle que ψ ◦}
\NormalTok{          ′}
\NormalTok{qtrop et qtrop diffèrent d’une translation sur Cf . Compte tenu de la convention faite}
\NormalTok{au paragraphe 12.1.13, l’espace affine ⟨Cf ⟩ est Q{-}rationnel, si bien que ψ appartient}
\NormalTok{                                 § 14.3. CALIBRAGES CANONIQUES                                  407}

\NormalTok{à GLn (Q). Il existe donc deux morphismes de tores φ et φ′ de Gnm dans lui{-}même tels}
\NormalTok{que φ′trop ◦ ψ = φtrop . Ainsi, (φ′ ◦ q ′ ◦ f )trop et (φ ◦ q ◦ f )trop diffèrent d’une constante}
\NormalTok{                       −1}
\NormalTok{sur C. Posons V = ftrop   (Cf ) ; c’est un domaine analytique fermé de X qui contient C.}
\NormalTok{                ′  ′}
\NormalTok{Sur V , on a (φ ◦ q ◦ f )trop = (φ ◦ q ◦ f )trop .}
\NormalTok{  D’après le premier cas particulier, on a donc presque partout les relations}

\NormalTok{                     degx (φ ◦ q ◦ f |V )σφ◦q,∗ (µ) = degx (q ◦ f |V )σq,∗ (µ)}

\NormalTok{et}
\NormalTok{                   degx (φ′ ◦ q ′ ◦ f |V )σφ′ ◦q′ ,∗ (µ) = degx (q ′ ◦ f |V )σq′ ,∗ (µ).}
\NormalTok{D’autre part, il résulte du second cas particulier traité que}

\NormalTok{                degx (φ ◦ q ◦ f |V )σφ◦q,∗ (µ) = degx (φ′ ◦ q ′ ◦ f |V )σφ′ ◦q′ ,∗ (µ).}

\NormalTok{Comme V est un domaine de X contenant x, on a les égalités de degrés}

\NormalTok{             degx (q ◦ f |V ) = degx (q ◦ f ) et         degx (q ′ ◦ f |V ) = degx (q ′ ◦ f )}

\NormalTok{pour tout x ∈ V . Cela conclut la démonstration de la proposition.}

\NormalTok{14.3.2. — Soit Σ une partie squelettique de X, paracompacte et contenue dans le}
\NormalTok{polyèdre caractéristique Σf d’un moment f : V → T défini sur un domaine analytique V}
\NormalTok{de X. Nous allons définir un calibrage µΣ de Σ.}
\NormalTok{   Soit C une décomposition cellulaire de Σ en f {-}cellules commodes (lemme 14.1.5).}
\NormalTok{Soit C une cellule de dimension n appartenant à C . Soit q : T → Gnm un morphisme de}
\NormalTok{tores tel que qtrop soit injectif sur ftrop (C). D’après la proposition 14.2.2, le morphisme}
\NormalTok{q ◦f est fini et plat sur C̊ et son degré x 7→ degx (q ◦f ) prend une valeur finie, strictement}
\NormalTok{positive et constante dC sur C̊. Notons e = |e1 ∧ · · · ∧ en | le vecteur{-}volume canonique}
\NormalTok{de Rn , considéré comme un calibrage constant, et posons µC = dC (q ◦ f )∗trop (e) ; c’est}
\NormalTok{un calibrage standard de la cellule C. Il ne dépend pas du choix de q d’après la propo{-}}
\NormalTok{sition 14.3.1. La famille (µC ) définit un calibrage standard de Σ ; on le note µΣ,f,C .}
\NormalTok{   Vérifions que ce calibrage est indépendant du choix de (V, f, C ).}
\NormalTok{   Soit V ′ un domaine analytique de X, soit f ′ : V ′ → T ′ un moment tel que Σ ⊂ Σf ′ et}
\NormalTok{soit C ′ une décomposition cellulaire de Σ en f ′ {-}cellules commodes. Posons W = V ∩ V ′}
\NormalTok{et soit g le moment (f, f ′ ) sur W et considérons une décomposition cellulaire D de Σ en}
\NormalTok{g{-}cellules commodes qui est adaptée aux décompositions cellulaires C et C ′ .}
\NormalTok{   Soit D une cellule de dimension n appartenant à D. Soit q : T → Gnm un morphisme}
\NormalTok{de tores tel que qtrop est injectif sur ftrop (D). Notons aussi p : T × T ′ → T la première}
\NormalTok{projection.}
\NormalTok{   Soit C une cellule de C contenant D. Comme D̊ ne rencontre pas ∂(W ), il est contenu}
\NormalTok{dans Int(W/V ), si bien que q ◦ f |W est fini et plat de degré dC en tout point de D̊. Il}
\NormalTok{existe aussi un entier dD \textgreater{} 0 tel que (q ◦ p) ◦ g soit fini et plat de degré dD en tout point}
\NormalTok{de D̊ (proposition 14.2.2). Comme on a (q ◦ p) ◦ g = q ◦ f |W , on a dC = dD . Par ailleurs,}
\NormalTok{(q ◦ g)∗trop (e) = (q ◦ f )∗trop (e). Par suite, les triplets (V, f, C ) et (W, g, D) induisent le}
\NormalTok{même calibrage sur D.}
\NormalTok{   Cela entraîne que le calibrage µΣ,f,C de Σ déduit de (V, f, C ) coïncide avec celui déduit}
\NormalTok{de (W, g, D) ; par symétrie, il coïncide avec celui déduit de (V ′ , f ′ , C ′ ). Notons{-}le µΣ .}
\NormalTok{408          CHAPITRE 14. CALIBRAGE CANONIQUE D’UN ESPACE ANALYTIQUE}

\NormalTok{14.3.3. — Soit Σ′ un sous{-}espace paralinéaire de Σ. Soit C ′ une décomposition cellu{-}}
\NormalTok{laire de Σ en f {-}cellules commodes qui raffine C ∩ Σ′ . La description de µΣ′ ,f,C ′ entraîne}
\NormalTok{alors l’égalité µΣ |Σ′ = µΣ′ .}
\NormalTok{Lemme 14.3.4. — Soit X un espace analytique de niveau n et soit Σ une partie sque{-}}
\NormalTok{lettique de X. Il existe un unique calibrage µΣ de Σ vérifiant la propriété suivante : pour}
\NormalTok{tout domaine analytique V de X, tout moment f sur V et tout sous{-}espace paralinéaire Σ′}
\NormalTok{de Σf ∩ Σ qui est paracompact, on a µΣ |Σ′ = µΣ′ . Ce calibrage µΣ est standard.}
\NormalTok{Démonstration. — Par définition, l’espace squelettique Σ est G{-}recouvert par des sous{-}}
\NormalTok{espaces paralinéaires Σ′ comme dans l’énoncé. Les calibrages formant un G{-}faisceau, il}
\NormalTok{existe au plus un tel calibrage. Pour démontrer son existence, il suffit alors de vérifier}
\NormalTok{la condition de coïncidence. Or, si Σ′ est Σ′′ sont des sous{-}espaces paralinéaires de Σ}
\NormalTok{satisfaisant les conditions de l’énoncé, on a µΣ′ |Σ′ ∩Σ′′ = µΣ′ ∩Σ′′ = µΣ′′ |Σ′ ∩Σ′′ d’après le}
\NormalTok{paragraphe 14.3.3.}

\NormalTok{Définition 14.3.5. — Soit X un espace analytique de niveau n, topologiquement sé{-}}
\NormalTok{paré, et soit Σ une partie squelettique de X. On appelle calibrage canonique de Σ le}
\NormalTok{calibrage dont le lemme 14.3.4 affirme l’existence et l’unicité.}
\NormalTok{   Par construction, ce calibrage est standard. Si Σ′ ⊂ Σ, on a µΣ |Σ′ = µΣ′ . Il résulte}
\NormalTok{aussi des constructions que si Σ est contenu dans un domaine analytique V de X, le}
\NormalTok{calibrage canonique µΣ de Σ, vu comme partie squelettique de V , coïncide avec son}
\NormalTok{calibrage canonique comme partie squelettique de X.}
\NormalTok{   Ce calibrage canonique munit l’espace vertébré associé à l’espace analytique X d’une}
\NormalTok{structure d’espace G{-}tropical calibré de niveau n. Nous supposerons toujours implicite{-}}
\NormalTok{ment que X est muni de cette structure.}
\NormalTok{Proposition 14.3.6. — Soit X un espace analytique de niveau n. Soit Σ une partie}
\NormalTok{squelettique de X, soit µ son calibrage canonique. Soit L une extension complète de k,}
\NormalTok{soit p : XL → X le morphisme de changement de base, soit ΣL la partie squelettique dé{-}}
\NormalTok{duite de Σ par changement de base (théorème 12.9.11) et soit µL son calibrage canonique.}
\NormalTok{On a p∗ (µL ) = µ.}
\NormalTok{Démonstration. — Il suffit de traiter le cas où X est (para)compact et où il existe un}
\NormalTok{moment f : X → T tel que Σ = Σf . Alors, ΣL = ΣfL (théorème 12.9.11). Le morphisme}
\NormalTok{de changement de base p étant compact, il existe une décomposition cellulaire C de Σ en}
\NormalTok{f {-}cellules commodes et une décomposition cellulaire D de ΣL en fL {-}cellules commodes}
\NormalTok{telles que le couple (D, C ) soit adapté au morphisme p. Soit C une cellule de dimension n}
\NormalTok{de C . et soit q : T → Gnm un morphisme de tores tel que qtrop soit injectif sur ftrop (C).}
\NormalTok{Soit x ∈ C̊ et soit t = q ◦ f (x) ∈ S(Gnm ). Soit tL ∈ (Gnm )L l’unique antécédent de t}
\NormalTok{dans S(Gnm )L . D’après le lemme 13.3.3, on a}
\NormalTok{                                             X}
\NormalTok{                       degx (q ◦ f ) =                 degy (q ◦ fL ).}
\NormalTok{                                      y∈p−1 (x)∩(q◦fL )−1 (tL )}

\NormalTok{Les points y qui apparaissent au second membre de cette égalité sont des antécédents}
\NormalTok{de x dans ΣfL . Inversement, soit y ∈ ΣfL un point tel que p(y) = x ; prouvons que}
\NormalTok{q ◦ fL (y) = tL . Soit D la cellule de D telle que y ∈ D̊ ; on a p(D) = C, donc (q ◦ fL )trop =}
\NormalTok{(q ◦ f )trop ◦ p est injectif sur D, si bien que q ◦ fL est de dimension tropicale n en y. Cela}
\NormalTok{                  § 14.4. CALIBRAGES CANONIQUES ET MORPHISMES FINIS                         409}

\NormalTok{entraîne que q ◦ fL (y) ∈ S((Gnm )L ) ; comme c’est un antécédent de t, on a q ◦ fL (y) = tL .}
\NormalTok{On a donc                                    X}
\NormalTok{                          degx (q ◦ f ) =            degy (q ◦ fL ).}
\NormalTok{                                            y∈p−1 (x)∩ΣfL}

\NormalTok{Par conséquent, µC = D p∗ µD , où D parcourt l’ensemble des n{-}cellules D de D au{-}}
\NormalTok{                         P}
\NormalTok{dessus de C. Ainsi, µΣ = p∗ (µΣL ).}




\NormalTok{       § 14.4. CALIBRAGES CANONIQUES ET MORPHISMES FINIS}

\NormalTok{14.4.1. — Soit p : Y → X un morphisme entre espaces k{-}analytiques de niveau n.}
\NormalTok{Soit Σ une partie squelettique de X. Supposons que les fibres de p au{-}dessus de tout}
\NormalTok{point de Σ sont des ensembles finis.}
\end{Highlighting}
\end{Shaded}

\subsection{CLD.txt, newline lines
20924--20976}\label{cld.txt-newline-lines-2092420976}

\begin{Shaded}
\begin{Highlighting}[]
\NormalTok{Remarque 14.4.3. — Soit F une extension finie de k et soit X un espace F {-}analytique}
\NormalTok{de niveau n. On peut aussi voir X comme un espace k{-}analytique de niveau n ; pour}
\NormalTok{éviter les confusions nous le noterons ici X0 ; on a dim(X0 ) = dim(X).}
\NormalTok{   Comme l’extension F de k est algébrique, le squelette canonique S(Gnm,F ) est l’image}
\NormalTok{réciproque de S(Gnm ) par la projection de Gnm,F sur Gnm , et S(Gnm ) est l’image de S(Gnm,F )}
\NormalTok{([50], (0.15)). On en déduit qu’une partie de X est squelettique si et seulement si elle}
\NormalTok{l’est en tant que partie de X0 .}
\NormalTok{   Soit donc Σ une partie squelettique de X ; notons µ son calibrage canonique et µ0 celui}
\NormalTok{de Σ vu comme partie squelettique de X0 . La projection canonique de Gnm,F sur Gnm est}
\NormalTok{finie et plate de degré [F : k]. Il découle donc de la définition du calibrage canonique}
\NormalTok{que µ0 = [F : k]µ.}



\NormalTok{                § 14.5. CALIBRAGES ET TROPICALISATIONS}

\NormalTok{Définition 14.5.1. — Soit n un entier, soit X un espace k{-}analytique de niveau n et}
\NormalTok{soit f : X → T un moment agréable sur X. On appelle polyèdre tropical du moment f}
\NormalTok{                    (n)                                                            (n)}
\NormalTok{le polyèdre ftrop (Σf ) muni du calibrage image directe du calibrage canonique de Σf .}
\NormalTok{On le note Πf .}
\NormalTok{  Ce n’est autre que le polyèdre tropical de la carte tropicale agréable ftrop .}
\NormalTok{  Lorsque X est compact, rappelons (théorème 12.8.2) que l’on a l’égalité de polyèdres}
\NormalTok{Πf = ftrop (X)(n) .}
\NormalTok{                                       (n)}
\NormalTok{  Comme le calibrage canonique de Σf est standard, il en est de même du calibrage}
\NormalTok{canonique de Πf .}
\NormalTok{  Comme la formation du calibrage canonique de Σf commute à l’extension des scalaires}
\NormalTok{(proposition 14.3.6), il en est de même de la formation du calibrage canonique de Πf .}
\NormalTok{14.5.2. — Voyons maintenant comment décrire explicitement le calibrage de Πf . Soit}
\NormalTok{                                       (n)}
\NormalTok{C une décomposition cellulaire de Σf vérifiant les propriétés de la proposition 14.1.12}
\NormalTok{et soit D = \{ftrop (C) ; C ∈ C \} la décomposition cellulaire de Πf qui lui correspond.}
\NormalTok{Soit D une n{-}cellule appartenant à D et soit q : T → Gnm un morphisme de tores tel que}
\NormalTok{qtrop |D soit injectif. Soit CD l’ensemble des n{-}cellules C ∈ C telles que ftrop (C) = D ;}
\NormalTok{pour C ∈ CD , soit dC l’entier naturel que q ◦ f soit fini et plat de degré dC en tout}
\NormalTok{                           § 14.5. CALIBRAGES ET TROPICALISATIONS                              411}
\NormalTok{                                          P}
\NormalTok{point de C̊ (prop. 14.2.2). Notons dD = C∈CD dC . Notons e = |e1 ∧ · · · ∧ en | le vecteur{-}}
\NormalTok{volume canonique de Rn , considéré comme un calibrage constant. Soit σD : Rn → Ttrop}
\NormalTok{                                                                                        (n)}
\NormalTok{l’unique section affine de qtrop d’image ⟨D⟩. Par définition du calibrage canonique de Σf}
\NormalTok{et de l’image directe d’un calibrage, on a µΠf ,D = dD σD,∗ (e).}

\NormalTok{Proposition 14.5.3. — Soit X un espace k{-}analytique de dimension n, soit f : X → T}
\NormalTok{un moment agréable et soit x ∈ Πf . Soit V un voisinage de x dans Πf qui est une n{-}cellule}
\NormalTok{ouverte de Ttrop . Soit q : T → Gnm un morphisme de tores tel que qtrop |V est injectif.}
\NormalTok{Supposons que q ◦ f |ftrop}
\NormalTok{                       −1}
\NormalTok{                           (V ) est fini et plat au{-}dessus de qtrop (x)sq , de degré dx . Notons}
\NormalTok{        n}
\NormalTok{σq : R → Ttrop l’unique section affine de qtrop d’image ⟨V ⟩ et notons e = |e1 ∧ · · · ∧ en |}
\NormalTok{le vecteur{-}volume canonique de Rn , considéré comme un calibrage constant. Alors, le}
\end{Highlighting}
\end{Shaded}

\subsection{CLD.txt, newline lines
21545--21664}\label{cld.txt-newline-lines-2154521664}

\begin{Shaded}
\begin{Highlighting}[]
\NormalTok{Proposition 15.1.1. — Soit X un espace k{-}analytique topologiquement séparé, consi{-}}
\NormalTok{déré comme espace G{-}tropical affiné.}
\NormalTok{   a) Soit ω une G{-}forme de type (p, q) sur X. Il existe un G{-}recouvrement (Vi ) de X et,}
\NormalTok{pour tout i, un moment fi : Vi → Ti sur Vi , une forme αi de type (p, q) sur une partie}
\NormalTok{                                                                 ∗}
\NormalTok{paralinéaire de Ti,trop contenant fi,trop (Vi ) tels que ω|Vi = fi,trop αi .}
\NormalTok{   b) Soit ω une forme de type (p, q) sur X. Il existe un recouvrement ouvert (Vi ) de X}
\NormalTok{et, pour tout i, un moment fi : Vi → Ti sur Vi , une forme αi de type (p, q) sur une partie}
\NormalTok{                                                                 ∗}
\NormalTok{paralinéaire de Ti,trop contenant fi,trop (Vi ) tels que ω|Vi = fi,trop αi .}
\NormalTok{422 CHAPITRE 15. FORMES DIFFÉRENTIELLES RÉELLES EN GÉOMÉTRIE ULTRAMÉTRIQUE}

\NormalTok{Démonstration. — On ne traite que le cas des G{-}formes, l’autre est analogue. Par hy{-}}
\NormalTok{pothèse, il existe un G{-}recouvrement (Wi ) de X et, pour tout i, une application parali{-}}
\NormalTok{néaire gi sur Wi à valeurs dans un espace affine Rni , une forme βi de type (p, q) sur une}
\NormalTok{partie paralinéaire de Ei contenant gi (Wi ), tels que ω|Wi = gi∗ βi pour tout i. Par défini{-}}
\NormalTok{tion d’une application paralinéaire, il existe pour tout i un G{-}recouvrement (Vi,j ) de Wi}
\NormalTok{et, pour tout j, un moment fi,j sur Vi,j , à valeurs dans le tore Gnmi un nombre rationnel}
\NormalTok{ai,j et un élément ci,j ∈ Rni tels que gi,j = ai,j fi,j,trop + ci,j . Notons pi,j : Rntrop}
\NormalTok{                                                                                        i}
\NormalTok{                                                                                           → Rni}
\NormalTok{l’application affine t 7→ at + ci,j ; on a donc gi,j = pi,j ◦ fi,j,trop , de sorte que}
\NormalTok{                                               ∗}
\NormalTok{                                     ω|Vi,j = fi,j,trop (p∗i,j βi ).}

\NormalTok{L’assertion en résulte.}


\NormalTok{Corollaire 15.1.2. — Une fonction u sur X est lisse au voisinage d’un point x s’il}
\NormalTok{existe un voisinage U de x dans X, des fonctions inversibles f1 , . . . , fn sur U et une}
\NormalTok{fonction φ sur Rn au voisinage de (f1,trop (x), . . . , fn,trop (x)) telle que u = φ ◦ ftrop sur U .}

\NormalTok{  Dans la suite du texte, nous travaillerons essentiellement avec ce type de cartes (G){-}}
\NormalTok{tropicales.}

\NormalTok{Proposition 15.1.3. — Soit G une classe de coefficients. Soit X un espace analytique}
\NormalTok{topologiquement séparé et soit u une fonction numérique sur X.}
\NormalTok{  a) Pour que u soit G{-}localement de classe G , il faut et il suffit que u|Xred le soit.}
\NormalTok{  b) Si X est presque bon, pour que u soit localement de classe G , il faut et il suffit que}
\NormalTok{u|Xred le soit.}

\NormalTok{  Une variante de cet énoncé vaut pour les formes et les G{-}formes.}
\NormalTok{Démonstration. — Si u est localement (resp. G{-}localement) de classe G , il en est de}
\NormalTok{même de sa restriction à Xred . Démontrons la réciproque.}
\NormalTok{   a) Supposons que u|Xred soit G{-}localement de classe G . Soit V un domaine affinoïde}
\NormalTok{de Xred sur lequel il existe un moment f : V → Gnm et une fonction φ de classe G sur Rn}
\NormalTok{telle que u|V = φ ◦ ftrop . On déduit du théorème de Gerritzen–Grauert (voir par exemple}
\NormalTok{[128], théorème 3.1) que V est recouvert par une réunion de domaines affinoïdes de la}
\NormalTok{forme Wred , où W est un domaine affinoïde de X. Sur Wred , chacune des fi se relève}
\NormalTok{en une fonction gi sur W ; ces fonctions sont inversibles et ont même norme que fi ;}
\NormalTok{elles définissent un moment g : W → Gnm tel que uW = φ ◦ gtrop . Cela prouve que u est}
\NormalTok{G{-}localement de classe G .}
\NormalTok{   b) Soit x un point de X et soit V un voisinage analytique de x dans Xred sur lequel}
\NormalTok{il existe un moment f : V → Gnm et une fonction φ de classe G sur Rn telle que u|V =}
\NormalTok{φ ◦ ftrop . Comme X est presque bon en x, chaque fi (x) ∈ H (x) est approchée par une}
\NormalTok{fonction gi ∈ OX,x . En particulier fi et gi ont même norme au voisinage de g. Ainsi, les}
\NormalTok{gi définissent un moment g défini au voisinage de X sur lequel on a u = φ ◦ gtrop ; cela}
\NormalTok{prouve que u est de classe G au voisinage de x.}
\NormalTok{                              § 15.2. SÉPARATION HOLOMORPHE                                423}

\NormalTok{                      § 15.2. SÉPARATION HOLOMORPHE}

\NormalTok{Proposition 15.2.1. — Soit X un espace k{-}analytique topologiquement séparé, consi{-}}
\NormalTok{déré comme espace G{-}tropical affiné. Soit ω une G{-}forme de type (p, q) sur X.}
\NormalTok{  a) Le support de ω est contenu dans l’ensemble des points x ∈ X tels que dk (x) ⩾}
\NormalTok{max(p, q).}
\NormalTok{  b) Soit V un domaine analytique de X et soit Y un fermé de Zariski de V . Si}
\NormalTok{dim(Y ) \textless{} max(p, q), le support de ω ne rencontre pas Y .}

\NormalTok{Démonstration. — Démontrons l’assertion a). Soit x ∈ X un point tel que dk (x) \textless{}}
\NormalTok{max(p, q). Comme dtrop (x) ⩽ dk (x), on a donc dtrop (x) \textless{} max(p, q) ; l’assertion résulte}
\NormalTok{donc de la proposition 3.3.6.}
\NormalTok{  L’assertion b) en résulte car tout point y de Y vérifie dk (y) ⩽ dim(Y ).}

\NormalTok{15.2.2. Séparation holomorphe. — On dira qu’un espace analytique X est holo{-}}
\NormalTok{morphiquement séparé si pour tout couple (x, y) de points de X avec x ̸= y il existe une}
\NormalTok{fonction analytique f sur X telle que |f (x)| =}
\NormalTok{                                              ̸ |f (y)|. Un tel espace est topologiquement}
\NormalTok{séparé.}
\NormalTok{  Si X est un espace holomorphiquement séparé, il en va trivialement de même de tout}
\NormalTok{domaine analytique de X.}
\NormalTok{  Tout espace affinoïde est par définition holomorphiquement séparé ; en conséquence,}
\NormalTok{tout domaine analytique d’un espace affinoïde est holomorphiquement séparé.}
\NormalTok{  Un espace analytique sera dit localement holomorphiquement séparé s’il est topologi{-}}
\NormalTok{quement séparé et si chacun de ses points possède un voisinage analytique holomorphi{-}}
\NormalTok{quement séparé.}
\NormalTok{  Un espace analytique topologiquement séparé qui est bon, ou plus généralement dont}
\NormalTok{tout point possède un voisinage isomorphe à un domaine analytique d’un espace affi{-}}
\NormalTok{noïde, est localement holomorphiquement séparé. C’est en particulier le cas d’un espace}
\NormalTok{analytique topologiquement séparé sans bord.}

\NormalTok{Lemme 15.2.3. — Soit X un espace k{-}analytique holomorphiquement séparé. L’an{-}}
\NormalTok{neau des fonctions lisses sur X sépare les points.}

\NormalTok{Démonstration. — Soit x, y des points distincts de X. Par définition de la séparation}
\NormalTok{holomorphe, il existe une fonction analytique f sur X telle que |f (x)| =     ̸ |f (y)|.}
\NormalTok{  Supposons ces deux valeurs non nulles ; soit φ une fonction C ∞ à support compact}
\NormalTok{sur R valant 0 au voisinage de log|f (x)| et 1 au voisinage de log|f (y)|. La fonction}
\NormalTok{φ ◦ log|f |, définie sur l’ouvert d’inversibilité de f , est lisse ; prolongeons{-}la par 0 sur X}
\NormalTok{tout entier. Comme φ est à support compact, la fonction obtenue est nulle au voisinage}
\NormalTok{de l’ensemble des zéros de f ; elle est donc lisse et prend des valeurs différentes en x et y.}
\NormalTok{  Si f (x) = 0, on a alors f (y) ̸= 0 et on raisonne de même à partir d’une fonction φ à}
\NormalTok{support compact valant 1 au voisinage de log|f (y)|.}
\NormalTok{  Le cas où f (y) = 0 est analogue par symétrie.}

\NormalTok{Corollaire 15.2.4. — Soit X un espace k{-}analytique localement holomorphiquement}
\NormalTok{séparé. L’espace G{-}tropical affiné X est localement différentiellement séparé. En parti{-}}
\NormalTok{culier, le faisceau des fonctions lisses sur X sépare localement les points.}
\NormalTok{424 CHAPITRE 15. FORMES DIFFÉRENTIELLES RÉELLES EN GÉOMÉTRIE ULTRAMÉTRIQUE}

\NormalTok{Corollaire 15.2.5. — Soit X un espace analytique localement holomorphiquement sé{-}}
\NormalTok{paré. Soit K une partie compacte de X et soit U un voisinage ouvert de K. Il existe}
\NormalTok{une fonction lisse sur X qui vaut 1 au voisinage de K et dont le support est compact et}
\NormalTok{contenu dans U .}
\NormalTok{Démonstration. — Comme X est localement différentiellement séparé, c’est une consé{-}}
\NormalTok{quence du corollaire 6.2.4.}

\NormalTok{Théorème 15.2.6 (Stone–Weierstraß). — Soit X un espace analytique localement}
\NormalTok{holomorphiquement séparé et soit U un ouvert de X. Soit f une fonction continue à}
\NormalTok{valeurs réelles sur X dont le support est compact et contenu dans U . Soit a et b des}
\end{Highlighting}
\end{Shaded}

\subsection{CLD.txt, newline lines
21712--21776}\label{cld.txt-newline-lines-2171221776}

\begin{Shaded}
\begin{Highlighting}[]
\NormalTok{coefficients tropicalement boréliens définit une mesure de Borel ωBor sur X, et lorsque}
\NormalTok{cette G{-}forme est localement intégrable, une mesure de Radon ωRad .}

\NormalTok{Proposition 15.3.1. — Soit L une extension complète de k, soit p : XL → X le mor{-}}
\NormalTok{phisme de changement de base. Pour toute forme ω de type (n, n) sur XG à coefficients}
\NormalTok{tropicalement boréliens, on a p∗ (ωL,Bor ) = ωBor . Lorsque ω est localement intégrable, on}
\NormalTok{a p∗ (ωL,Rad ) = ωRad .}

\NormalTok{Démonstration. — En revenant aux définitions, on déduit de la compatibilité du cali{-}}
\NormalTok{brage canonique à l’extension des scalaires (proposition 14.3.6) l’égalité p∗ ω}
\NormalTok{                                                                              fL = ωL de}
\NormalTok{formes{-}volume n{-}dimensionnelles. La proposition en résulte.}

\NormalTok{Proposition 15.3.2. — Soit L une extension complète de k, soit p : XL → X le mor{-}}
\NormalTok{phisme de changement de base. Pour toute forme ω de type (n−1, n) sur XG à coefficients}
\NormalTok{tropicalement boréliens, on a p∗ (ωL,Bor ) = ωBor . Lorsque ω est localement intégrable, on}
\NormalTok{a p∗ (ωL,Rad ) = ωRad .}

\NormalTok{Démonstration. — En revenant aux définitions, on déduit de la compatibilité du cali{-}}
\NormalTok{brage canonique à l’extension des scalaires (proposition 14.3.6) l’égalité p∗ ω}
\NormalTok{                                                                              fL = ωL de}
\NormalTok{formes{-}volume (n − 1){-}dimensionnelles. La proposition en résulte.}




\NormalTok{               § 15.4. INTÉGRALES ET MORPHISMES FINIS}

\NormalTok{15.4.1. — Soit n ∈ N et soit X et Y des espaces k{-}analytiques de niveau n et soit}
\NormalTok{f : Y → X un morphisme qui est fini au{-}dessus du complémentaire d’un fermé de Zariski}
\NormalTok{de X de dimension \textless{} n. On a défini au §14.4.1 une fonction x 7→ dcx (f ) de X dans N ;}
\NormalTok{nous la noterons dcf .}
\NormalTok{   Soit p ∈ \{n − 1, n\} et soit ω une G{-}forme de type (p, n) sur X. Posons η = f ∗ ω.}

\NormalTok{Lemme 15.4.2. — Soit Σ un partie squelettique de X qui supporte fortement ω ; notons}
\NormalTok{ω̃ la forme{-}volume p{-}dimensionnelle sur Σ déduite de ω et du calibrage canonique de Σ.}
\NormalTok{   Alors p−1 (Σ) est une partie squelettique de Y qui supporte fortement f ∗ ω ; notons fg}
\NormalTok{                                                                                         ∗ω}

\NormalTok{la forme{-}volume p{-}dimensionnelle sur p−1 (Σ) déduite de f ∗ ω et du calibrage canonique}
\NormalTok{de p−1 (Σ). On a l’égalité :}
\NormalTok{                                          ∗ ω = dc ·ω̃.}
\NormalTok{                                      f∗ fg       f}
\NormalTok{426 CHAPITRE 15. FORMES DIFFÉRENTIELLES RÉELLES EN GÉOMÉTRIE ULTRAMÉTRIQUE}

\NormalTok{Démonstration. — Posons T = f −1 (Σ). Notons µΣ et µT les calibrages canoniques de Σ}
\NormalTok{et T. On a l’égalité f∗ µT = dcf ·µΣ (proposition 14.4.2). L’assertion découle alors de la}
\NormalTok{proposition 4.5.8.}
\NormalTok{  Compte tenu de la définition des mesures de Borel et, éventuellement de Radon,}
\NormalTok{associées aux formes volumes p{-}dimensionnelles, on en déduit la proposition suivante.}

\NormalTok{Proposition 15.4.3. — Supposons que ω soit localement intégrable. Alors, f ∗ ω est}
\NormalTok{localement intégrable et l’on a l’égalité}
\NormalTok{                                    f∗ (f ∗ ω)Rad = dcf ·ωRad}
\NormalTok{de mesures de Radon sur X.}

\NormalTok{   Remarquons que les supports des mesures ωRad et f∗ (f ∗ ω)Rad évitent tout fermé de Za{-}}
\NormalTok{riski d’intérieur vide de X ; en particulier, ils sont contenus dans le domaine de définition}
\NormalTok{du degré calibré.}



\NormalTok{§ 15.5. EXISTENCE DE TROPICALISATIONS (PRESQUE) GLOBALES}

\NormalTok{   La définition d’une forme (resp. d’une G{-}forme) est locale (resp. G{-}locale), mais lorsque}
\end{Highlighting}
\end{Shaded}

\subsection{CLD.txt, newline lines
22100--22151}\label{cld.txt-newline-lines-2210022151}

\begin{Shaded}
\begin{Highlighting}[]
\NormalTok{§ 15.7. L’ESPACE TROPICAL ASSOCIÉ À UN ESPACE ANALYTIQUE}

\NormalTok{  Soit X un espace analytique purement de dimension n, topologiquement séparé. On le}
\NormalTok{considère comme espace G{-}tropical calibré de niveau n. Il est riche (théorème 12.10.6).}

\NormalTok{Théorème 15.7.1. — Soit X un espace analytique purement de dimension n, topolo{-}}
\NormalTok{giquement séparé. Muni de son bord ∂(X) au sens de la théorie des espaces analytiques,}
\NormalTok{l’espace G{-}tropical affiné et calibré de niveau n associé à X est un espace tropical de}
\NormalTok{niveau n.}

\NormalTok{Démonstration. — Vérifions les conditions de la définition 6.6.2.}
\NormalTok{  L’intérieur Int(X) de X est localement holomorphiquement séparé (15.2.2) ; il est donc}
\NormalTok{localement différentiellement séparé (corollaire 15.2.4). Il est décent (proposition 12.4.7)}
\NormalTok{car il est sans bord, donc bon, donc presque bon. Il est très riche (théorème 12.11.1) et}
\NormalTok{harmonieux (théorème 15.6.1).}
\NormalTok{  Et il résulte du corollaire 12.9.14 que Σ ∩ ∂X est n{-}négligeable pour toute partie}
\NormalTok{paralinéaire Σ de X, de même que Σ∩∂(V /X) pour tout domaine V de X car ∂(V /X) ⊂}
\NormalTok{∂V .}
\NormalTok{  On munit ainsi tout espace analytique X, topologiquement séparé, et purement de}
\NormalTok{dimension n, d’une structure d’espace tropical de niveau n.}
\NormalTok{              § 15.7. L’ESPACE TROPICAL ASSOCIÉ À UN ESPACE ANALYTIQUE              433}

\NormalTok{15.7.2. — Soit X un espace analytique topologiquement séparé, purement de dimen{-}}
\NormalTok{sion n et soit V un domaine de X. Le domaine V est muni de deux façons, a priori}
\NormalTok{distinctes, d’une structure d’espace tropical de niveau n, soit comme espace analytique,}
\NormalTok{soit comme domaine de l’espace tropical X. La formule ∂(V ) = (∂(X) ∩ V ) ∪ ∂(V /X)}
\NormalTok{montre que ces deux structures coïncident.}
\NormalTok{Remarque 15.7.3. — Soit X un espace analytique de niveau n ; pour que l’on puisse}
\NormalTok{le munir d’une structure d’espace tropical de niveau n, il faut et il suffit que X soit}
\NormalTok{purement de dimension n. En effet, dans le cas contraire, il possède un domaine ana{-}}
\NormalTok{lytique compact non vide V tel que dim(V ) \textless{} n. L’espace X n’est donc pas riche : la}
\NormalTok{tropicalisation constante de V n’est couverte par aucun squelette de V puisque la seule}
\NormalTok{partie squelettique de V est vide.}
\NormalTok{CHAPITRE 16}


\NormalTok{COURANTS SUR LES ESPACES ANALYTIQUES}




\NormalTok{   Soit X un espace k{-}analytique topologiquement séparé et purement de dimension n.}
\NormalTok{Pour alléger la terminologie, nous dirons plus simplement que X est un n{-}espace k{-}}
\NormalTok{analytique. Nous l’avons muni d’une structure d’espace G{-}tropical de niveau n, calibré}
\NormalTok{et affiné, et démontré (théorème 15.7.1) que muni de son bord au sens de la théorie des}
\NormalTok{espaces k{-}analytiques, c’est un espace tropical de niveau n.}
\NormalTok{   Il en résulte des notions de formes et G{-}formes différentielles de type (p, q) sur X, à}
\NormalTok{coefficients dans une classe de coefficients, une théorie de l’intégration des (n, n){-}formes}
\NormalTok{et des (n − 1, n){-}formes ainsi qu’une théorie de courants.}
\NormalTok{   Dans ce chapitre, nous explorons plus avant les propriétés spécifiques des courants}
\NormalTok{sur les espaces analytiques qui semblent difficilement formulables dans le contexte des}
\NormalTok{espaces tropicaux.}
\end{Highlighting}
\end{Shaded}

\subsection{CLD.txt, newline lines
29080--29148}\label{cld.txt-newline-lines-2908029148}

\begin{Shaded}
\begin{Highlighting}[]
\NormalTok{pour tout entier n ⩾ 1, posons Kn = W n Wn−1 ; les ensembles Kn sont compacts.}
\NormalTok{L’ensemble ouvert Ωn = Vn+1 Vn−2 est un voisinage de Kn ; par construction, Ωn ne}
\NormalTok{rencontre pas Ωm si |m − n| ⩾ 3. Pour tout x ∈ Kn , il existe un élément Vx de B contenu}
\NormalTok{dans un des ensembles Ui et contenu dans cet ouvert Ωn et tel que x ∈ V̊x . Comme Kn est}
\NormalTok{compact, il existe une partie finie Jn de Kn telle que les V̊x , pour x ∈ Jn , recouvrent Kn .}
\NormalTok{La famille des Vx , pour n ∈ N et x ∈ Jn , convient.}




\NormalTok{      § A.2. SÉPARATION, DENSITÉ ET PARTITIONS DE L’UNITÉ}

\NormalTok{A.2.1. — Nous démontrons ici un certain nombre de résultats généraux et relative{-}}
\NormalTok{ment standard concernant la séparation, la densité de certains sous{-}espaces de fonctions}
\NormalTok{continues et les partitions de l’unité par de telles fonctions.}
\NormalTok{   Nous considérons un espace topologique localement compact X (en particulier séparé)}
\NormalTok{et un sous{-}faisceau A en R{-}algèbres du faisceau des fonctions continues à valeurs réelles}
\NormalTok{sur X. On suppose que pour tout ouvert U de X, tout élément f ∈ A (U ) et toute}
\NormalTok{fonction lisse φ sur R définie au voisinage de f (U ), la fonction φ ◦ f appartient à A (U ).}
\NormalTok{   Notons que cette hypothèse est vérifiée dans les deux cas suivants :}
\NormalTok{   – L’espace X est un espace topologique et A est le faisceau des fonctions continues}
\NormalTok{sur X ;}
\NormalTok{   – L’espace X est un espace paralinéaire affiné ou un espace analytique, plus généra{-}}
\NormalTok{lement un espace G{-}tropical affiné, et A est le faisceau des fonctions lisses sur X.}
\NormalTok{   Pour alléger la rédaction, nous appellerons fonctions lisses les sections du faisceau A .}
\NormalTok{A.2.2. — On dit que A (X) sépare les points si pour tout couple (x, y) de points de X}
\NormalTok{tels que x ̸= y, il existe une section globale f de AX telle que f (x) ̸= f (y).}
\NormalTok{   On dit que A sépare localement les points si pour tout point x de X, il existe un}
\NormalTok{voisinage ouvert U de X tel que A (U ) sépare les points de U .}
\NormalTok{                   § A.2. SÉPARATION, DENSITÉ ET PARTITIONS DE L’UNITÉ                      565}

\NormalTok{Lemme A.2.3. — Supposons que A (X) sépare les points. Soit K un sous{-}ensemble}
\NormalTok{compact de X et soit x ∈ X K. Il existe une fonction lisse sur X qui vaut 1 au voisinage}
\NormalTok{de x et nulle au voisinage de K.}
\NormalTok{Démonstration. — On commence par traiter le cas où K est réduit à un point y. Par}
\NormalTok{hypothèse, il existe une fonction lisse f sur X telle que f (x) ̸= f (y). On choisit une}
\NormalTok{fonction lisse u sur R qui est constante égale à 0 au voisinage de f (x) et constante égale}
\NormalTok{à 1 au voisinage de y ; l’existence de telles fonctions est classique. La fonction u ◦ f sur X}
\NormalTok{est lisse et convient.}
\NormalTok{   D’après ce qui précède, pour tout y ∈ K, il existe une fonction lisse fy sur X valant 1}
\NormalTok{au voisinage de x et nulle dans un voisinage ouvert Vy de    Qy. Soit S une partie finie de K}
\NormalTok{telle que les Vy , pour y ∈ S, recouvrent K ; la fonction y∈S fy convient.}

\NormalTok{Corollaire A.2.4. — On suppose que A sépare localement les points. Soit K une}
\NormalTok{partie compacte de X et soit U un voisinage ouvert de K. Il existe une fonction lisse}
\NormalTok{sur X qui vaut 1 au voisinage de K et dont le support est compact et contenu dans U .}
\NormalTok{Démonstration. — On commence par traiter le cas où le compact K est réduit à un}
\NormalTok{point x. Soit V un voisinage compact de x tel que A (V̊ ) sépare les points de V̊ . L’en{-}}
\NormalTok{semble F = V (V̊ ∩U ) est une partie fermée de X contenue dans V . Soit f une fonction}
\NormalTok{lisse sur V qui vaut 1 au voisinage de x et qui est nulle au voisinage de F . Soit g la}
\NormalTok{fonction sur X qui coïncide avec f sur V et qui est nulle en dehors de V . Elle est lisse}
\NormalTok{sur l’intérieur de V , sur l’ouvert X V ainsi que sur un voisinage ouvert de F dans X ;}
\NormalTok{ces trois ouverts recouvrent X. La fonction g est donc lisse. Son support est contenu}
\NormalTok{dans U ; comme il est contenu dans V , il est aussi compact.}
\NormalTok{   Traitons maintenant le cas général. Pour tout x ∈ K, soit fx une fonction lisse sur X}
\NormalTok{qui vaut 1 au voisinage de x et dont le support est compact et contenu dans U . Quitte à}
\NormalTok{remplacer fx par son carré, on suppose que cette fonction est positive ou nulle. Comme K}
\NormalTok{est compact, la somme g d’une sous{-}famille finie de ces fonctions est strictement positive}
\NormalTok{en tout point de K. Elle est lisse et son support est compact et contenu dans U . Soit φ}
\NormalTok{une fonction lisse de R vers [0; 1], nulle en 0 et égale à 1 sur un voisinage de g(K) ; la}
\NormalTok{fonction f = φ ◦ g convient.}

\NormalTok{Théorème A.2.5 (Stone–Weierstraß). — On suppose que A sépare localement les}
\NormalTok{points. Soit U un ouvert de X et soit f une fonction continue à valeurs réelles sur X}
\NormalTok{dont le support est compact et contenu dans U . Soit a et b des nombres réels tels que}
\NormalTok{a ⩽ f (x) ⩽ b pour tout x ∈ X.}
\NormalTok{  Pour tout ε \textgreater{} 0, il existe une fonction lisse g sur X dont le support est compact,}
\NormalTok{contenu dans U , et telle que a ⩽ g(x) ⩽ b et |f (x) − g(x)| \textless{} ε pour tout x ∈ X.}
\NormalTok{Démonstration. — Soit X ′ le compactifié d’Alexandroff de X ; comme X est localement}
\end{Highlighting}
\end{Shaded}

\subsection{GR.txt, newline lines
302--368}\label{gr.txt-newline-lines-302368}

\begin{Shaded}
\begin{Highlighting}[]
\NormalTok{2.11. Relative boundary. An important notion in this paper is that of the relative bound{-}}
\NormalTok{ary of a morphism ϕ : X → Y of k{-}analytic spaces. This is a closed subset of X which we}
\NormalTok{denote by ∂ (X /Y ). We will not give the details of this definition, but let us try to explain}
\NormalTok{the literature on the subject. In [Ber90, §2.5], Berkovich defines the relative boundary}
\NormalTok{∂ (X /Y ) for k{-}affinoid spaces X and Y . Then in [Ber90, §3.1], he generalizes ∂ (X /Y ) for}
\NormalTok{good k{-}analytic spaces using that all points of X and Y have k{-}affinoid neighborhoods.}
\NormalTok{Berkovich calls ϕ closed if ∂ (X /Y ) = ;. Berkovich defines ∂ (X /Y ) for arbitrary k{-}analytic}
\NormalTok{spaces in [Ber93, 1.5.3(ii)], as follows. A morphism ϕ is closed if for every morphism}
\NormalTok{Y ′ → Y from a good k{-}analytic space Y ′ , the fiber product X ′ := X ×Y Y ′ is good and the}
\NormalTok{canonical morphism X ′ → Y ′ is closed in the previous sense. Finally in [Ber93, Defini{-}}
\NormalTok{tion 1.5.4], a point x ∈ X is defined to be in Int(X /Y ) if x has an open neighborhood U}
\NormalTok{such that ϕ induces a closed morphism U → Y in the above sense. We refer to [Tem15,}
\NormalTok{Definition 4.2.4.1] for a more explicit equivalent definition of ∂ (X /Y ). Note that in this}
\NormalTok{paper, we will use the notion closed map in the usual topological sense meaning that it}
\NormalTok{maps closed sets to closed sets. This is different from Berkovich’s terminology.}
\NormalTok{   The complement Int(X /Y ) := X \textbackslash{} ∂ (X /Y ) is called the relative interior. We will usually}
\NormalTok{apply this in the case Y = M (k), in which case we call ∂ X := ∂ (X /M (k)) the boundary}
\NormalTok{of X and Int(X ) := X \textbackslash{} ∂ X the interior. If ∂ X = ;, we say that X boundaryless. A bound{-}}
\NormalTok{aryless space is good by definition. Note that Berkovich calls boundaryless k{-}analytic}
\NormalTok{spaces closed.}
\NormalTok{   We list here some properties of the boundary and interior which we will use throughout}
\NormalTok{the paper:}
\NormalTok{   (1) A point x ∈ X is contained in Int(X /Y ) if and only if there is an open neighbor{-}}
\NormalTok{       hood U of x such that Int(U/Y ) = U: this is part of the above definition.}
\NormalTok{   (2) If ϕ : X → Y is finite, then Int(X /Y ) = X [Ber93, Example 1.5.3(iii)].}
\NormalTok{   (3) If X is an analytic domain in Y , then Int(X /Y ) is the topological interior of X in}
\NormalTok{       Y [Ber93, Proposition 1.5.5(i)].}
\NormalTok{             ϕ     ψ}
\NormalTok{   (4) If X → Y → S are morphisms of k{-}analytic spaces, then}

\NormalTok{(2.11.1)                   Int(X /Y ) ∩ ϕ −1 Int(Y /S) ⊂ Int(X /S)}

\NormalTok{        [Ber93, Proposition 1.5.5(ii)], with equality if ψ is locally separated [Tem04,}
\NormalTok{        Corollary 5.7]. Taking S = M (k), this becomes}

\NormalTok{(2.11.2)                      Int(X /Y ) ∩ ϕ −1 Int(Y ) ⊂ Int(X ).}

\NormalTok{   (5) If k′ /k is an analytic extension field, then}

\NormalTok{                                π−1}
\NormalTok{                                 k′ /k}
\NormalTok{                                       Int(X /Y ) ⊂ Int(X k′ /Yk′ )}
\NormalTok{8                                      W. GUBLER AND J. RABINOFF}

\NormalTok{         [Ber93, Proposition 1.5.5(iii)]. Taking Y = M (k), this becomes}
\NormalTok{                                       π−1}
\NormalTok{                                        k′ /k}
\NormalTok{                                              Int(X ) ⊂ Int(X k′ ).}
\NormalTok{For convenience, we restate these properties in terms of the relative boundary:}
\NormalTok{   (1′ ) A point x ∈ X is contained in ∂ (X /Y ) if and only if ∂ (U/Y ) ̸= ; for every open}
\NormalTok{         neighborhood U of x.}
\NormalTok{   (2′ ) If ϕ is finite, then ∂ (X /Y ) = ;.}
\NormalTok{   (3′ ) If X is an analytic domain in Y , then ∂ (X /Y ) is the topological boundary of X}
\NormalTok{         in Y .}
\NormalTok{               ϕ     ψ}
\NormalTok{     ′}
\NormalTok{   (4 ) If X → Y → S are morphisms of k{-}analytic spaces, then}
\NormalTok{(2.11.3)                        ∂ (X /S) ⊂ ∂ (X /Y ) ∪ ϕ −1 (∂ (Y /S)),}
\NormalTok{         with equality if ψ is locally separated. Taking S = M (k), this becomes}
\NormalTok{(2.11.4)                             ∂ X ⊂ ∂ (X /Y ) ∪ ϕ −1 (∂ Y ).}
\NormalTok{    (5′ ) If k′ /k is an analytic extension field, then}
\NormalTok{                                     ∂ (X k′ /Yk′ ) ⊂ π−1}
\NormalTok{                                                       k′ /k}
\NormalTok{                                                             ∂ (X /Y ).}
\NormalTok{         Taking Y = M (k), this becomes}
\NormalTok{                                           ∂ X k′ ⊂ π−1}
\NormalTok{                                                     k′ /k}
\end{Highlighting}
\end{Shaded}

\subsection{Additional complete body: CLD.txt, newline lines
15061--15081}\label{additional-complete-body-cld.txt-newline-lines-1506115081}

\begin{Shaded}
\begin{Highlighting}[]
\NormalTok{les objets sout les couples (X, L) où L est une extension complète de k et X un espace L{-}}
\NormalTok{analytique. Un mrophisme de (X, L) vers (Y, M ) est un couple formé d’un k{-}morphisme}
\NormalTok{de M vers L et d’un morphisme d’espaces L{-}analytiques de X vers YL .}
\NormalTok{   Cette catégorie n’admet pas tous les produits fibrés (par exemple M (Cp ) ×M (Qp )}
\NormalTok{M (Cp ) n’existe pas en tant qu’espace analytique au{-}dessus de Qp ) mais si Y → X est}
\NormalTok{un morphisme d’espaces k{-}analytiques et X ′ → X un morphisme d’espaces analytiques}
\NormalTok{au{-}dessus de k alors Y ×k X ′ existe : si L est le corps de définition fourni avec X ′ , le}
\NormalTok{produit fibté Y ×k X ′ est simplement le produit fibré YL ×XL X ′ , calculé dans la catégorie}
\NormalTok{des espaces L{-}analytiques.}
\NormalTok{11.1.10. — Soit X un espace k{-}analytique, et soit x ∈ X.}
\NormalTok{   On note OX,x la fibre en x de la restriction de OX à la catégorie des ouverts de X,}
\NormalTok{c’est{-}à{-}dire la colimite des OX (U ) où U parcourt l’ensemble des voisinages ouverts de}
\NormalTok{X. C’est un anneau local dont nous noterons κ(x) le corps résiduel.}
\NormalTok{   Lorsque X est bon, cet anneau est noethérien ([9], th. 2.1.4), hensélien ([9], th. 2.1.5)}
\NormalTok{et excellent ([49], th. 2.13).}
\NormalTok{   Au point x est également associé une extension ultramétrique complète H (x) de k,}
\NormalTok{qu’on appelle le corps résiduel complété de x ; pour tout domaine analytique V de X}
\NormalTok{302                  CHAPITRE 11. RAPPELS DE GÉOMÉTRIE ANALYTIQUE}

\NormalTok{contenant x, on dispose d’un morphisme d’évaluation f 7→ f (x) de OX (V ) dans H (x).}
\NormalTok{Si X = M (A), le corps H (x) et l’application f 7→ f (x) de A vers H (x) coïncident avec}
\end{Highlighting}
\end{Shaded}

\subsection{Additional complete body: CLD.txt, newline lines
15356--15361}\label{additional-complete-body-cld.txt-newline-lines-1535615361}

\begin{Shaded}
\begin{Highlighting}[]
\NormalTok{kr {-}analytique de M (kr ) est nulle. Par ailleurs on a kr ⊗kb r = (kr )r ≃ kr \{T, S\}/(ST − 1),}
\NormalTok{qui est strictement kr {-}affinoïde de dimension de Krull n. La dimension k{-}analytique de}
\NormalTok{M (kr ) est donc égale à n.}


\NormalTok{11.2.5. Les composantes irréductibles. — Lorsque nous parlerons d’un espace k{-}}
\end{Highlighting}
\end{Shaded}

\subsection{Additional complete body: CLD.txt, newline lines
6978--6985}\label{additional-complete-body-cld.txt-newline-lines-69786985}

\begin{Shaded}
\begin{Highlighting}[]
\NormalTok{Prouvons inversement que cette relation dans des cartes tropicales permet de reconstituer}
\NormalTok{la mesure de Radon associée à une G{-}forme de type (n, n) et en donne ainsi une définition}
\NormalTok{tropicale.}
\NormalTok{   Soit ω une G{-}forme de type (n, n) sur X ; supposons{-}la à support compact et à coef{-}}
\NormalTok{ficients tropicalement boréliens localement bornés.}
\NormalTok{   Soit (Vi )i∈I une famille finie de domaines compacts de X tropicalisant ω et dont dont}
\NormalTok{la réunion}
\NormalTok{      T       V recouvre le support de ω. Pour toute partie non vide J de I, posons}
\end{Highlighting}
\end{Shaded}

\subsection{Additional complete body: CLD.txt, newline lines
20871--20923}\label{additional-complete-body-cld.txt-newline-lines-2087120923}

\begin{Shaded}
\begin{Highlighting}[]
\NormalTok{  Soit x un point de Σ et soit y ∈ f −1 (x). Comme dk (x) = n, on a dim(f −1 (x)) = 0, de}
\NormalTok{sorte que f est quasi{-}finie en y. Par suite, f −1 (x) est un H (x){-}espace analytique fini.}
\NormalTok{On a ainsi défini au §13.3.6 le nombre rationnel dcx (f ).}

\NormalTok{Proposition 14.4.2. — Soit p : Y → X un morphisme d’espaces k{-}analytiques de}
\NormalTok{niveau n. Soit Σ une partie squelettique de X, soit µΣ son calibrage canonique. Soit T la}
\NormalTok{partie squelettique p−1 (Σ) de Y (proposition 12.9.6) et soit µT son calibrage canonique.}
\NormalTok{  Supposons que le morphisme pΣ : T → Σ est topologiquement compact.}
\NormalTok{  a) La restriction de la fonction x 7→ dcx (p) à toute partie paralinéaire paracompacte}
\NormalTok{de Σ est localement constante en dehors d’une partie paralinéaire n{-}négligeable.}
\NormalTok{  b) On a p∗ µT = dc(p)µΣ .}

\NormalTok{   Sous les hypothèses de b), les fibres de f en tout point de Σ sont finies, de sorte que}
\NormalTok{la fonction x 7→ dcx (p) sur Σ est définie. De plus, l’assertion a) donne sens au produit}
\NormalTok{dc(p)µΣ et montre que c’est un calibrage standard de Σ.}
\NormalTok{Démonstration. — On sait (proposition 12.9.6) que p induit une immersion par mor{-}}
\NormalTok{ceaux compacte de T sur Σ. Ainsi, l’image directe p∗ µT est définie, et est un calibrage}
\NormalTok{standard de Σ.}
\NormalTok{   L’espace X est G{-}recouvert par des domaines analytiques compacts V sur lesquels}
\NormalTok{il existe un moment f tel que Σ ∩ V ⊂ Σf . Puisque la formation des calibrages cano{-}}
\NormalTok{niques et des degrés calibrés commute à la restriction à un domaine analytique, on peut}
\NormalTok{remplacer X par V et Y par p−1 (V ), ce qui nous ramène à prouver la proposition sous}
\NormalTok{l’hypothèse supplémentaire que X est compact et qu’il existe un moment f : X → Gnm tel}
\NormalTok{que Σ soit une partie squelettique de Σf . En remplaçant Y par un voisinage analytique}
\NormalTok{compact de T, on se ramène au cas où Y est compact.}
\NormalTok{                                                                                    (n)}
\NormalTok{   Les moments f et f ◦ p sont agréables et l’on a p−1 (Σf ) = Σf ◦p et p−1 (Σf ) =}
\NormalTok{ (n)                                                                        (n)}
\NormalTok{Σf ◦p (prop. 12.9.6). Choisissons une décomposition cellulaire C de Σf            en f {-}cellules}
\NormalTok{                                                     (n)}
\NormalTok{commodes, et une décomposition cellulaire D de Σf ◦p en f ◦ p{-}cellules commodes, telles}
\NormalTok{                                             (n)     (n)}
\NormalTok{que le couple (D, C ) soit adapté à pΣ(n) : Σf ◦p → Σf , à Σ et à T.}
\NormalTok{                                        f}
\NormalTok{  Soit C une n{-}cellule de C et soit D1 , . . . , Dr les n{-}cellules de D dont l’image par p}
\NormalTok{                                                   (n)        (n)}
\NormalTok{est égale à C ; l’immersion par morceaux de Σf ◦p vers Σf induit au{-}dessus de C̊ un}
\NormalTok{410          CHAPITRE 14. CALIBRAGE CANONIQUE D’UN ESPACE ANALYTIQUE}

\NormalTok{revêtement dont les feuillets sont les D̊i . Soit x ∈ C̊ ; pour tout i, on note yi l’unique}
\NormalTok{antécédent de x sur D̊i , et l’on pose ξ = f (x).}
\NormalTok{  Par définition d’une cellule commode, le point x appartient à Int(X). De même, chacun}
\NormalTok{des yi appartient Int(Y ). Le morphisme f ◦ p est fini et plat en chacun des yi , de degré}
\NormalTok{degyi (f ◦ p). Le morphisme f est fini et plat en x, de degré degx (f ).}
\NormalTok{  Pour tout i, on a degyi (f ◦ p) = dcyi (f ◦ p), car f ◦ p est fini et plat en yi . Par}
\NormalTok{multiplicativité des degrés calibrés, on a degyi (f ◦p) = dcyi (p) dcx (f ), donc degyi (f ◦p) =}
\NormalTok{dcyi (p) degx (f ) puisque f est fini et plat en x. On a donc}
\NormalTok{                   Xr                 X r}
\NormalTok{                          yi}
\NormalTok{                       deg (f ◦ p) =      dcyi (p) · degx (f ) = dcx (p) degx (f ).}
\NormalTok{                 i=1                 i=1}
\NormalTok{Compte tenu de la définition des calibrages canoniques, cela entraîne l’égalité voulue.}
\end{Highlighting}
\end{Shaded}

\subsection{Additional complete body: Ducros-S1793744212500077.txt,
newline lines
2077--2154}\label{additional-complete-body-ducros-s1793744212500077.txt-newline-lines-20772154}

\begin{Shaded}
\begin{Highlighting}[]
\NormalTok{Théorème 3.2. Soit X un espace k{-}analytique compact et Γ{-}strict de dimension d,}
\NormalTok{et soit f = (f1 , . . . , fn ) une famille de fonctions analytiques inversibles sur X ;}
\NormalTok{notons |f | l’application (|f1 |, . . . , |fn |) de X vers (R∗+ )n . Posons c = (Q, (|k ∗ | · Γ)Q ).}
\NormalTok{1) L’image |f |(X) est un c{-}polytope de (R∗+ )n , qui est de dimension inférieure ou}
\NormalTok{   égale à d.}
\NormalTok{2) L’image |f |(∂X) est contenue dans un c{-}polytope de (R∗+ )n qui est de dimension}
\NormalTok{   inférieure ou égale à d − 1 ; lorsque X est affinoı̈de |f |(∂X) est elle{-}même un c{-}}
\NormalTok{   polytope de (R∗+ )n de dimension inférieure ou égale à d − 1.}
\NormalTok{    Démonstration. La seconde assertion se déduira de la première.}
\NormalTok{(3.2.1) Preuve de 1) : réduction à une situation algébrique. On peut supposer}
\NormalTok{que X est affinoı̈de et irréductible, et l’on procède alors par récurrence sur d. Si d = 0}
\NormalTok{il n’y a rien à démontrer. Supposons que d \textgreater{} 0 et que 1) est vraie en dimension \textless{} d.}
\NormalTok{     Le théorème à établir est purement topologique. On peut donc supposer k parfait}
\NormalTok{(en étendant les scalaires au complété de la clôture radicelle de k), et X réduit (en}
\NormalTok{le remplaçant par Xred ). L’espace X est alors génériquement quasi{-}lisse. Soit Y son}
\NormalTok{lieu de non{-}quasi{-}lissité. C’est un fermé de Zariski de X dont toutes les composantes}
\NormalTok{irréductibles sont de dimension \textless{} d. Par hypothèse de récurrence, |f |(Y ) est un}
\NormalTok{c{-}polytope de (R∗+ )n de dimension ⩽ d − 1. Son image réciproque |f |−1 (|f |(Y )) est}
\NormalTok{un domaine analytique Γ{-}strict de X contenant Y . Comme Y → X est sans bord, Y}
\NormalTok{est contenu dans l’intérieur topologique de |f |−1 (|f |(Y )). Ce dernier apparaı̂t ainsi}
\NormalTok{comme un voisinage de Y dont l’image par |f | est un c{-}polytope de dimension}
\NormalTok{⩽ d − 1.}
\NormalTok{     Soit x ∈ X − Y . Comme X est quasi{-}lisse en x, le lemme 0.21 assure l’existence}
\NormalTok{d’un voisinage affinoı̈de V de x dans X, et d’une k{-}variété affine X de dimension d,}
\NormalTok{telle que V s’identifie à un domaine affinoı̈de Γ{-}strict de X an (on peut imposer à X}
\NormalTok{d’être lisse, mais nous ne nous en servirons pas).}
\NormalTok{     Par compacité de X, on se ramène finalement au cas où X lui{-}même est est un}
\NormalTok{domaine affinoı̈de Γ{-}strict de l’analytification X an d’une k{-}variété algébrique X}
\NormalTok{affine de dimension d.}
\NormalTok{     Écrivons X ≃ Spec A où A est une k{-}algèbre de type fini, et soit (a1 , . . . , ar )}
\NormalTok{un système de générateurs de A. Il existe R ∈ Γ · |k ∗ | tel que X soit contenu}
\NormalTok{dans le domaine affinoı̈de W de X an défini par les conditions |ai | ⩽ R. En vertu}
\NormalTok{du théorème de Gerritzen{-}Grauert (0.20.7), le domaine X est réunion finie de}
\NormalTok{domaines Γ{-}rationnels de W ; il suffit dès lors de traiter le cas où X est lui{-}même}
\NormalTok{un domaine Γ{-}rationnel de W .}
\NormalTok{     Il est alors défini par une conjonction d’inégalités de la forme |gi | ⩽ λi |h| où les λi}
\NormalTok{appartiennent à Γ et où (g1 , . . . , gr , h) engendre l’idéal des fonctions analytiques}
\NormalTok{sur W . Cette dernière condition implique que h ne s’annule pas sur X (tout zéro}
\NormalTok{de h sur X est un zéro commun à h et aux gi ), et l’on peut donc, par compacité,}
\NormalTok{                                            1250007{-}40}
\NormalTok{                            Espaces de Berkovich, polytopes, squelettes et théorie des modèles}


\NormalTok{rajouter une condition de la forme |h| ⩾ r à la définition de X, où r est un élément}
\NormalTok{convenable de (Γ · |k ∗ |)Q . Il est dès lors immédiat que V ne change pas si l’on}
\NormalTok{perturbe un peu les gi et h ; on peut par conséquent supposer que toutes ces}
\NormalTok{fonctions appartiennent à A.}
\NormalTok{    L’anneau des fonctions de la forme a/h , où a ∈ A, et  ⩾ 0, est dense dans}
\NormalTok{l’anneau des fonctions de X. Cela permet de supposer, quitte à remplacer X}
\NormalTok{par D(h), à remplacer gi par gi /h et à approcher convenablement chacune des fi ,}
\NormalTok{que fi ∈ A pour tout i et que h = 1.}
\NormalTok{(3.2.2) Nature c{-}polytopale de |f |(X) dans le cas algébrique. Soit (F, |.|) une}
\NormalTok{extension valuée de k (on ne suppose pas que le groupe ordonné |F ∗ | est de}
\NormalTok{rang 1). On peut donner un sens naturel à la notation X(F ), coı̈ncidant avec sa}
\NormalTok{signification usuelle lorsque F est un corps ultramétrique complet : c’est le sous{-}}
\NormalTok{ensemble de X (F ) formé des points P tels que |gi (P )| ⩽ λi pour tout i.}
\NormalTok{    Le théorème 1.2 assure l’existence d’un sous{-}foncteur fermé et définissable D}
\NormalTok{de GnΓ·|k∗ | , de dimension majorée par d (les notations et définitions sont celles de}
\NormalTok{0.32 et sq.) tel que}
\NormalTok{                             |f |(X(F )) = D(|F ∗ |) ⊂ |F ∗ |n}
\NormalTok{pour tout F ∈ Ck,Γ·|k∗ | .}
\NormalTok{     Si F est une extension ultramétrique complète et algébriquement close de k}
\NormalTok{munie d’un |k ∗ |{-}isomorphisme |F ∗ | ≃ R∗+ , tout point P de X(F ) définit un point x}
\NormalTok{de l’espace k{-}affinoı̈de X, et |f |(x) = |f |(P ) ∈ (R∗+ )n .}
\NormalTok{     Réciproquement si x ∈ X il existe une extension complète et algébriquement}
\NormalTok{close F de H (x) et un |H (x)∗ |{-}isomorphisme |F ∗ | ≃ R∗+ , et x est alors induit par}
\NormalTok{un F {-}point canonique de X.}
\NormalTok{     En conséquence, le compact |f |(X) est le sous{-}ensemble D(R∗+ ) de (R∗+ )n , ce}
\NormalTok{qui achève de prouver 1).}
                                                                                  
\NormalTok{(3.2.3) Démonstration de l’assertion 2. Si X s’écrit comme une réunion finie Xi}
\NormalTok{où les Xi sont affinoı̈des et Γ{-}stricts, son bord est contenu dans la réunion des bords}
\NormalTok{des Xi . On peut donc supposer que X est affinoı̈de, auquel cas le résultat découle}
\NormalTok{de l’assertion 1) déjà établie et du lemme 3.1.}
\NormalTok{(3.3) Remarque. Faute d’avoir réussi à exhiber une description suffisamment}
\NormalTok{maniable de ∂X en général, l’auteur ignore s’il peut arriver que |f |(∂X) ne soit}
\NormalTok{pas un polytope, lorsque X n’est pas affinoı̈de. Il serait évidemment intéressant de}
\NormalTok{disposer d’un tel exemple — ou, au contraire, d’une preuve que |f |(∂X) est toujours}
\end{Highlighting}
\end{Shaded}

\subsection{Application of the underlying Ducros polyhedrality
theorem}\label{application-of-the-underlying-ducros-polyhedrality-theorem}

Ducros Theorem3.2(1) is used only for the compact strict affinoid moment
domains in the finite separable curve argument. With Γ=\{1\}, K
nontrivially valued, these are compact Γ-strict analytic spaces of
dimension1; finite analytic unit tuples have finite polyhedral image of
dimension at most1. The proof's preliminary reductions concern only the
topological moment image and do not replace the actual algebra in
canonical measures. Its model-theoretic Theorem1.2 is explicitly an
underlying foundational dependency of the cited polyhedrality theorem,
not a new result proved here. The cited primary theorem applies
precisely under these hypotheses. No claim of universal moment-image
dimension equality is used.

\subsection{CLD-curve, newline lines
910--978}\label{cld-curve-newline-lines-910978}

\begin{Shaded}
\begin{Highlighting}[]
\NormalTok{qui est croissante et d’image I sur \{1, . . . , p\} et est croissante sur \{p + 1, . . . , n\}. Si K}
\NormalTok{et L sont des parties à n − p éléments de \{1, . . . , n\}, on a}
\NormalTok{                                      (}
\NormalTok{                                        (−1)p(n−p) εI εJ λ ∧ Jλ si K = ∁I et L = ∁J,}
\NormalTok{           ω ∧ d′ xK ∧ J d′ xL =}
\NormalTok{                                        0                            sinon.}
\NormalTok{En outre, φeK = εK d′ x∁K , si bien que}
\NormalTok{                                (}
\NormalTok{                                 (−1)p(n−p) λ ∧ Jλ si I = K et J = L,}
\NormalTok{                bω (eK , eL ) =}
\NormalTok{                                 0                 sinon.}
\NormalTok{                                § 1.2. SUPERFORMES POSITIVES                                  25}

\NormalTok{Enfin,                                                      (}
\NormalTok{                                                             1 si I = K et J = L,}
\NormalTok{               qω (eK , eL ) = ⟨d′ xI , eK ⟩⟨d′ xJ , eL ⟩ =}
\NormalTok{                                                             0 sinon,}
\NormalTok{d’où la formule annoncée.}
\NormalTok{   En outre, les formules précédentes entraînent que les signes de ces diverses formes sont}
\NormalTok{reliées par les équivalences suivantes :}
\NormalTok{La forme ω est positive au sens de la définition 1.2.3 ;}
\NormalTok{⇔ L’application bilinéaire (α, β) 7→ ω ∧ α ∧ Jβ est de signe σn−p ;}
\NormalTok{⇔ L’application bilinéaire bω est de signe σn−p ;}
\NormalTok{⇔ La forme bilinéaire qω est de signe σn−p (−1)p(n−p) σn = σp .}
\NormalTok{   Enfin, si f : W → V est une application linéaire, on a qf ∗ ω = qω ◦ (f, f ).}
\NormalTok{   Nous dirons qu’une (p, p){-}forme symétrique est définie positive si la forme bilinéaire}
\NormalTok{symétrique σp qω est définie positive.}
\NormalTok{Lemme 1.2.7. — Soit ω ∈ Ap,p (V ) une forme symétrique de type (p, p), soit}
\NormalTok{(x1 , . . . , xn ) une base de V et soit}
\NormalTok{                                                X}
\NormalTok{                               ω = (−1)p(p−1)/2   ωI,J d′ xI ∧ d′′ xJ}
\NormalTok{                                               I,J}

\NormalTok{la décomposition de ω dans la base correspondante de Ap,p (V ). Alors, les quatre propriétés}
\NormalTok{suivantes sont équivalentes :}
\NormalTok{   a) La forme ω est positive ( resp. définie positive) ;}
\NormalTok{   b) La matrice symétrique (ωI,J ) (indexée par les parties à p éléments de \{1, . . . , n\})}
\NormalTok{est positive ( resp. définie positive) ;}
\NormalTok{   c) Il existe une base (αI ) de Λp V ∗ et une famille (cI ) de nombres réels positifs ou nuls}
\NormalTok{( resp. strictement positifs), indexées par les parties à p éléments de \{1, . . . , n\}, telles}
\NormalTok{que                                             X}
\NormalTok{                              ω = (−1)p(p−1)/2     cI αI ∧ J(αI ).}
\NormalTok{                                                     I}

\NormalTok{Démonstration. — L’équivalence des assertions (1) et (2) résulte de la discussion qui}
\NormalTok{précède le lemme et de ce qu’avec les notations introduites, la matrice (ωIJ ) est celle de}
\NormalTok{la forme bilinéaire (−1)p(p−1)/2 qω dans la base duale de la base (d′ xI )I de Λp V ∗ . Celle des}
\NormalTok{propriétés (2) et propriétés (3) découle de la diagonalisation des formes quadratiques.}

\NormalTok{Corollaire 1.2.8. — L’ensemble des formes positives dans Ap,p (V ) est le cône convexe}
\NormalTok{fermé engendré par les formes du type (−1)p(p−1)/2 α ∧ Jα, où α parcourt l’ensemble des}
\NormalTok{formes de type (p, 0).}
\NormalTok{Démonstration. — Le lemme précédent montre que les formes positives sont contenues}
\NormalTok{dans ce cône. Inversement, si α ∈ Ap,0 (V ) est une forme de type (p, 0) et β est une}
\NormalTok{(n − p, 0){-}forme, la formule}
\NormalTok{     (−1)p(p−1)/2 α ∧ Jα ∧ (−1)(n−p)(n−p−1)/2 β ∧ Jβ = (−1)n(n−1)/2 (α ∧ β) ∧ J(α ∧ β)}
\NormalTok{montre que (−1)p(p−1)/2 α ∧ Jα est une forme positive, d’où l’autre inclusion.}

\NormalTok{Corollaire 1.2.9. — La condition de positivité est auto{-}duale : une forme α, symé{-}}
\NormalTok{trique de type (p, p), est positive si et seulement si α ∧ β est positive pour toute forme}
\NormalTok{positive β de type (n − p, n − p).}
\NormalTok{26 CHAPITRE 1. SUPERFORMES SUR UN ESPACE VECTORIEL RÉEL (D’APRÈS LAGERBERG)}

\NormalTok{Démonstration. — Cela résulte immédiatement du corollaire précédent.}

\NormalTok{Corollaire 1.2.10. — Les trois notions de positivité coïncident pour les formes de}
\NormalTok{bidegrés (0, 0), (1, 1), (n − 1, n − 1) et (n, n).}
\NormalTok{   Elles ne coïncident pas en général, pas plus que leurs analogues en géométrie complexe.}
\NormalTok{Démonstration. — Le cas des formes de type (0, 0) et (n, n) a déjà été expliqué. Par}
\end{Highlighting}
\end{Shaded}

\subsection{CLD-curve, newline lines
15082--15286}\label{cld-curve-newline-lines-1508215286}

\begin{Shaded}
\begin{Highlighting}[]
\NormalTok{bon en x. Un espace bon est presque bon. Le point x est traditionnellement dit rigide si}
\NormalTok{H (x) est une extension finie de k (la terminologie provient du fait que ces points sont}
\NormalTok{ceux considérés par Tate dans sa théorie des espaces analytiques rigides).}
\NormalTok{   Soit f : Y → X un morphisme entre espaces k{-}analytiques et soit x ∈ X. La fibre}
\NormalTok{f (x) possède une structure naturelle d’espace H (x){-}analytique, et sera toujours im{-}}
\NormalTok{  −1}

\NormalTok{plicitement considérée comme munie de cette structure. Nous la noterons souvent Yx .}
\NormalTok{11.1.11. — Berkovich a construit ([9], 2.6) un foncteur d’analytification X 7→ X an}
\NormalTok{qui va de la catégorie des A{-}schémas localement de type fini vers celle des bons espaces}
\NormalTok{k{-}analytiques. L’espace X an est muni d’un morphisme d’espaces localement annelés}
\NormalTok{X an → X qui est fidèlement plat ([9], Prop. 2.6.2). L’application π induit une bijection}
\NormalTok{entre l’ensemble X0an des points rigides de X an et l’ensemble X0 des points fermés}
\NormalTok{de X , qui préserve les corps résiduels ; on en tire en particulier une bijection entre}
\NormalTok{\{x ∈ X an , H (x) = k et X (k), qui est un homéomorphisme (la source étant munie de}
\NormalTok{la topologie induite par celle de X an , et le but de la topologie déduite de celle de k). Si}
\NormalTok{|k × | ̸= \{1\}, l’ensemble X0an est dense dans X an .}
\NormalTok{   Si x est un point de X an d’image ξ sur X , le morphisme local OX ,ξ → OX an ,x est}
\NormalTok{régulier ([49], théorème 3.3).}
\NormalTok{   Donnons quelques détails sur la construction de X an lorsque X est affine. Écrivons}
\NormalTok{X = Spec A où A est une k{-}algèbre de type fini. Topologiquement, X an est l’en{-}}
\NormalTok{semble des semi{-}normes multiplicatives sur A qui induisent la valeur absolue structurale}
\NormalTok{sur k, muni de la topologie induite par la topologie produit sur RA ; le morphisme de}
\NormalTok{X an vers X envoie une semi{-}norme sur son noyau. Fixons un système fini de géné{-}}
\NormalTok{rateurs a1 , . . . , am de A, et soit I le noyau de la surjection Ti 7→ ai de k[T1 , . . . , Tm ]}
\NormalTok{sur A. Pour tout R \textgreater{} 0 le sous{-}ensemble VR de X an constitué des points x tels que}
\NormalTok{|ai (x)| ⩽ R pour tout i est un domaine affinoïde de X an , d’algèbre affinoïde associée}
\NormalTok{k\{T1 /R, . . . , Tn /R\}/I. L’espace X an est la réunion croissante des VR , qui en constituent}
\NormalTok{un G{-}recouvrement.}
\NormalTok{   On dispose d’un théorème GAGA dans ce contexte : si X est propre et si π désigne}
\NormalTok{le morphisme structural de X an vers X alors F 7→ π ∗ F établit une équivalence entre}
\NormalTok{la catégorie des faisceaux cohérents sur X et celle des faisceaux cohérents sur X an , qui}
\NormalTok{préserve les groupes de cohomologie ([109], annexe A).}
\NormalTok{   Le foncteur d’analytification commute à la formation des produits fibrés.}
\NormalTok{   Pour tout espace analytique X et tout entier n, l’application f 7→ (f ∗ T1 , . . . , f ∗ Tn )}
\NormalTok{met en bijection Hom(X, An,an   k  ) et OX (X)n .}

\NormalTok{                                                                                        aI T I 7→}
\NormalTok{                                                                                      P}
\NormalTok{Exemple 11.1.12. — Soit r = (r1 , . . . , rn ) un polyrayon. L’application}
\NormalTok{max|aI | · rI de k[T1 , . . . , Tn ] vers R+ est une norme multiplicative prolongeant la valeur}
\NormalTok{absolue de k. C’est donc un point ηr de An,an     k  , situé au{-}dessus du point générique de Ank .}
\NormalTok{Le corps H (ηr ) est le complété de k(T1 , . . . , Tn ) pour la valeur absolue induite par ηr .}
\NormalTok{  Lorsque r est k{-}libre le corps H (ηr ) s’identifie au corps k{-}affinoïde kr décrit en 11.1.3}
\NormalTok{(ce dernier contient en effet k(T1 , . . . , Tn ) comme sous{-}corps dense), le singleton \{ηr \} est}
\NormalTok{     § 11.1. GÉOMÉTRIE ANALYTIQUE : CONVENTIONS, NOTATIONS, PREMIERS RAPPELS                     303}

\NormalTok{alors le domaine affinoïde de An,an}
\NormalTok{                               k    défini par les égalités |Ti | = ri pour tout i, et l’on a}
\NormalTok{OAk (\{ηr \}) = kr .}
\NormalTok{   n,an}




\NormalTok{Exemple 11.1.13      P (La droite affine). —i Pour tout a ∈ k et tout réel r, notons}
\NormalTok{ηa,r l’application      ai (T − a) 7→ max|ai |r de k[T ] vers R+ . C’est une semi{-}norme}
\NormalTok{multiplicative, et donc un point de Ak1,an ; si r = 0 c’est l’application P 7→ |P (a)| de}
\NormalTok{noyau (T − a), et ηa,0 est donc le k{-}point standard d’équation T = a ; si r \textgreater{} 0 alors ηa,r}
\NormalTok{est une norme, ce qui revient à dire que ηa,r est situé au{-}dessus du point générique de}
\NormalTok{A1k . On a ηa,r = ηb,s si et seulement si r = s ⩾ |b − a|. Si |k × | est dense dans R×  + on a}
\NormalTok{|P (ηa,r )| = supb∈k,|b−a|⩽r |P (b)|.}
\NormalTok{   On a H (ηa,0 ) = k. Supposons r \textgreater{} 0 et posons τ = (T −  \^{} a)(ηa,r ).Il résulte de l’exemple}
\NormalTok{10.8.4 que τ est transcendant sur e k et que H\^{}}
\NormalTok{                                              (ηa,r ) = e}
\NormalTok{                                                        k(τ ).}
\NormalTok{    Donnons quelques précisions supplémentaires dans deux cas extrêmes. Si r = |λ| avec}
\NormalTok{                                                                           1}
\NormalTok{λ ∈ k × alors H (η}
\NormalTok{                 \^{}  a,r ) = k(τ /λ). Comme |τ /λ| = 1 on voit que H (ηa,r ) est égal à}
\NormalTok{                            e e                e                    \^{}}
\NormalTok{k 1 (τ /λ),}
\NormalTok{e       e et est donc une extension rationnelle de degré de transcendance 1 de ek 1 ; de}
\NormalTok{                                                                 1}
\NormalTok{                                        / |k × |Q alors H\^{}}
\NormalTok{même, on a |H (ηa,r )× | = |k × |. Si r ∈                          k 1 et |H (ηa,r )× | = |k × |·rZ .}
\NormalTok{                                                         (ηa,r ) = e}
\NormalTok{  Supposons maintenant que k est algébriquement clos (auquel cas |k × | est divisible).}
\NormalTok{On démontre alors (voir [8], 1.4.4) que tout point de A1,an    k  est de l’un des quatre types}
\NormalTok{suivants.}
\NormalTok{   – Les points de type I sont les points rigides, c’est{-}à{-}dire ceux pour lesquels H (x) = k,}
\NormalTok{ou encore les points de la forme ηa,0 pour a ∈ k.}
\NormalTok{   – Les points de type II sont les points de la forme ηa,r où a ∈ k et où r ∈ |k × |. Si x}
\NormalTok{                                1}
\NormalTok{est de type II le corps H \^{}  (x) est une extension rationnelle de degré de transcendance 1}
\NormalTok{de ke1 et |H (x)× | = |k × |. En particulier dk (x) = 1.}
\NormalTok{   – Les points de type III sont les points de la forme ηa,r où a ∈ k et où r ∈         / |k × |. Si}
\NormalTok{                                     1}
\NormalTok{x est de type III le corps H   \^{}  (x) est égale à ek 1 et |H (x)× | = |k × | · rZ . En particulier}
\NormalTok{dk (x) = 1.}
\NormalTok{   – Les points de type IV n’existent que si le corps k n’est pas maximalement complet,}
\NormalTok{c’est{-}à{-}dire s’il existe des familles totalement ordonnées de boules fermées de k d’inter{-}}
\NormalTok{section vide (c’est par exemple le cas si k = Cp , ou si k est le complété du corps des}
\NormalTok{séries de Puiseux sur C). Si (B(ai , ri ))i est une telle famille de boules, le point x de type}
\NormalTok{IV qu’on lui associe est la semi{-}norme inf i ηai ,ri , qui est en fait une norme. Le point x}
\NormalTok{est donc situé au{-}dessus du point générique de A1,an     k . Le corps complet H (x) est une}
\NormalTok{extension immédiate de k, c’est{-}à{-}dire que H (x) = e}
\NormalTok{                                                  \^{}         k, ou encore que H (x) a même}
\NormalTok{groupe des valeurs et même corps résiduel (au sens classique) que k ; et cette extension}
\NormalTok{est stricte. On a dk (x) = 0.}

\NormalTok{11.1.14. — Rappelons ([36], chap. 1, §9, no 10, définition 6, p. 69) qu’un espace to{-}}
\NormalTok{pologique X est dit paracompact s’il est séparé et si tout recouvrement ouvert possède}
\NormalTok{un raffinement qui est localement fini. Un espace compact est paracompact. Une partie}
\NormalTok{fermée d’un espace paracompact est paracompacte.}
\NormalTok{   Un espace localement compact est dit dénombrable à l’infini s’il est réunion dénom{-}}
\NormalTok{brable d’ensembles compacts ; un espace localement compact dénombrable à l’infini est}
\NormalTok{304                CHAPITRE 11. RAPPELS DE GÉOMÉTRIE ANALYTIQUE}

\NormalTok{paracompact. Plus généralement, un espace topologique localement compact est pa{-}}
\NormalTok{racompact si et seulement si c’est l’espace somme d’une famille d’espaces localement}
\NormalTok{compacts dénombrables à l’infini (loc. cit., théorème 5, p. 70).}
\NormalTok{Lemme 11.1.15 (cf. [9], remarque 1.2.4 (iii)). — Un espace analytique possède}
\NormalTok{une base de voisinages ouverts dénombrables à l’infini (et donc paracompacts).}
\NormalTok{   Cet énoncé figure dans [8, 2.2.8 et 3.2.1] pour les espaces affinoïdes ; nous détaillons}
\NormalTok{ci{-}dessous les arguments de la remarque 1.2.4, (iii), de [9] qui l’en déduit.}
\NormalTok{Démonstration. — Soit x un point de X. Nous allons montrer que tout voisinage de x}
\NormalTok{dans X contient un voisinage ouvert de x qui est dénombrable à l’infini. On peut pour}
\NormalTok{ce faire restreindre X et donc le supposer topologiquement séparé. Le point x possède}
\NormalTok{alors une base de voisinages compacts dans X qui sont des réunions finies de domaines}
\NormalTok{affinoïdes contenant x. On peut donc supposer que X est compact et est une union finie}
\NormalTok{V1 ∪V2 . . .∪Vn de domaines affinoïdes contenant x ; on se donne un voisinage ouvert U de}
\NormalTok{x dans X et on veut montrer qu’il contient un voisinage ouvert de x qui est dénombrable}
\NormalTok{à l’infini. Raisonnons par récurrence sur n. Si n = 1 alors X est affinoïde et le résultat}
\NormalTok{est connu dans ce cas comme on l’a rappelé. Supposons n \textgreater{} 1 et le résultat vrai pour}
\NormalTok{n − 1. Posons W = V1 ∪ · · · ∪ Vn−1 . Par récurrence, x possède un voisinage ouvert et}
\NormalTok{dénombrable à l’infini Ω dans W ∩ U , dont l’intersection ω avec Vn est un voisinage}
\NormalTok{ouvert et dénombrable à l’infini de x dans W ∩ Vn ∩ U . Soit (Ki )i⩾0 une suite de parties}
\NormalTok{compactes de W ∩ Vn telles que ω = i⩾0 Ki et Ki ⊂ K˚}
\NormalTok{                                       S}
\NormalTok{                                                           i+1 pour tout i. Soit i ⩾ 0. Soit}
\NormalTok{y ∈ Ki ; il existe un voisinage Ωy de y dans Vn tel que Ωy ∩ W ∩ Vn ⊂ K̊i+1 ; par le cas}
\NormalTok{affinoïde on peut supposer, quitte à le restreindre, que Ωy est un ouvert dénombrable}
\NormalTok{à l’infini de Vn ∩ U . L’espace Ki étant compact, il est recouvert par un nombre fini de}
\NormalTok{ces ouverts Ωy de Vn ; leur réunion est un voisinage Ωi de Ki dans Vn ∩ U , ouvert et}
\NormalTok{dénombrable à l’infini, tel que Ωi ∩ W ∩ Vn ⊂ K̊i+1 . La réunion Ω′ des Ωi est un voisinage}
\NormalTok{ouvert de x dans Vn ∩ U , ouvert et dénombrable à l’infini, et Ω′ ∩ W ∩ Vn = ω. La réunion}
\NormalTok{Ω ∪ Ω′ est dénombrable à l’infini et contient x. On peut écrire Ω = Υ ∩ W et Ω′ = Υ′ ∩ Vn}
\NormalTok{où Υ et Υ′ sont deux ouverts de X contenus dans U . Quitte à remplacer Υ par son}
\NormalTok{intersection avec le complémentaire dans X du compact Vn ω, on peut supposer que}
\NormalTok{Υ∩Vn = ω ; et quitte à remplacer Υ′ par son intersection avec le complémentaire dans X}
\NormalTok{du compact W ω, on peut supposer que Υ′ ∩ W = ω. Mais on a alors Ω ∪ Ω′ = Υ ∪ Υ′ ,}
\NormalTok{si bien que Ω ∪ Ω′ est un voisinage ouvert de x dans U .}




\NormalTok{      § 11.2. TOPOLOGIE DE ZARISKI, THÉORIE DE LA DIMENSION}

\NormalTok{11.2.1. Topologie de Zariski. — Soit X un espace k{-}analytique. Pour tout faisceau}
\NormalTok{cohérent d’idéaux I sur X on note V (I ) son lieu des zéros, c’est{-}à{-}dire l’ensemble}
\NormalTok{des point x de X tel que f (x) = 0 pour toute section f de I définie sur un domaine}
\NormalTok{analytique de X contenant x. Les parties de X de la forme V (I ) sont les fermés d’une}
\NormalTok{topologie appelée topologie de Zariski de X.}
\NormalTok{   Supposons que X est affinoïde, et écrivons X = M (A). Se donner un faisceau cohérent}
\NormalTok{d’idéaux I sur X revient alors à se donner un idéal I de A (on a I = Ie et I = I (X)),}
\NormalTok{et V (I ) est alors l’image réciproque de V (I) par la surjection continue de M (A) vers}
\NormalTok{                § 11.2. TOPOLOGIE DE ZARISKI, THÉORIE DE LA DIMENSION               305}

\NormalTok{Spec A (il est en particulier compact). On dispose donc dans ce cas d’une bijection}
\NormalTok{croissante entre l’ensemble des fermés de Zariski de M (A) et l’ensemble des fermés de}
\NormalTok{Zariski de Spec A. En particulier la topologie de X est noethérienne.}
\NormalTok{  On ne suppose plus X affinoïde. Les fermés de Zariski de X sont fermés dans X (cela se}
\NormalTok{vérifie G{-}localement, ce qui perrmet de se ramener au cas affinoïde). Si X est compact,}
\NormalTok{sa topologie de Zariski est noethérienne (on le déduit du cas affinoïde en écrivant X}
\NormalTok{comme une union finie de domaines affinoïde).}
\NormalTok{11.2.2. Sous{-}espaces analytiques fermés. — Soit X un espace k{-}analytique et soit}
\NormalTok{I un faisceau cohérent d’idéaux sur X. Supposons d’abord que X est affinoïde, et}
\NormalTok{écrivons X = M (A). Posons I = I (X). Le fermé V (I ) coïncide avec l’ensemble des}
\NormalTok{x dans M (A) tel que f (x) = 0 pour tout x ∈ I, qui est homéomorphe à M (A/I) via}
\NormalTok{le morphisme induit par la surjection canonique de A vers A/I. Ainsi, V (I ) peut être}
\NormalTok{muni d’une structure d’espace k{-}analytique, dépendant du choix de l’idéal I , et d’un}
\NormalTok{morphisme de V (I ) dans X dont l’application continue sous{-}jacente est l’inclusion.}
\NormalTok{   On ne suppose plus X affinoïde. En procédant G{-}localement à la construction pré{-}}
\NormalTok{cédente puis en recollant on obtient une structure d’espace k{-}analytique au{-}dessus de}
\NormalTok{X sur V (I ), l’application continue sous{-}jacente au morphisme de V (I ) dans X étant}
\NormalTok{l’inclusion i : V (I ) ,→ X, et i∗ OV (I ) = OX /I .}
\NormalTok{   Un sous{-}espace analytique fermé de X est un fermé de Zariski Y de X muni d’une}
\NormalTok{structure d’espace k{-}analytique au{-}dessus de X induite selon le procédé décrit ci{-}dessus}
\NormalTok{par un faisceau cohérent d’idéaux I tel que Y = V (I ) ; notons que I se reconstitue}
\NormalTok{alors à partir de cette structure comme le noyau de OX → i∗ OY , où i : Y → X est le}
\NormalTok{morphisme structural, qui est topologiquement l’inclusion.}
\NormalTok{   Tout fermé de Zariski Y possède au moins une structure de sous{-}espace analytique}
\NormalTok{fermé : celle donnée par le faisceau}

\NormalTok{                 J := V 7→ \{f ∈ O(V ), f (x) = 0 pour tout ∈ Y ∩ V}

\NormalTok{qui est bien cohérent ([49], Prop. 4.2.1). C’est la plus petite structure de sous{-}espace}
\NormalTok{analytique fermé sur Y , au sens où J est le plus grand faisceau cohérent d’idéaux}
\NormalTok{sur X définissant le fermé Y ; c’est aussi la seule structure de sous{-}espace analytique}
\NormalTok{fermée sur Y dont le faisceau structural n’a pas de sections nilpotentes non nulles ;}
\NormalTok{on l’appelle la structure réduite sur Y . Si Y = X la structure réduite est associée au}
\NormalTok{faisceau d’idéaux dont les sections sont les fonctions localement nilpotentes ; le sous{-}}
\NormalTok{espace analytique fermé correspondant est noté Xred . Le morphisme Xred ,→ X induit}
\NormalTok{un homéomorphisme entre Xred et X pour leur topologies usuelles, mais aussi pour leurs}
\NormalTok{topologies de Zariski.}
\NormalTok{   Si Y est un sous{-}espace analytique fermé de X, il résulte du théorème de Gerritzen{-}}
\NormalTok{Grauert dans sa version due à Temkin ([128], Th. 3.1) que les domaines analytiques de}
\NormalTok{Y sont les sous{-}ensembles de Y qui admettent un G{-}recouvrement par des parties de la}
\NormalTok{forme V ∩ Y où V est un domaine affinoïde de X.}
\NormalTok{   Si V est un domaine analytique de X la topologie de Zariski de V est plus fine que la}
\NormalTok{topologie de V induite par la topologie de Zariski de X, et en général strictement plus}
\NormalTok{fine.}
\NormalTok{   Si E est une partie de X et si (Xi ) est un G{-}recouvrement de X par des domaines}
\NormalTok{analytiques, alors E est un fermé (resp. un ouvert) de Zariski de X si et seulement si}
\end{Highlighting}
\end{Shaded}

\subsection{CLD-curve, newline lines
15362--15566}\label{cld-curve-newline-lines-1536215566}

\begin{Shaded}
\begin{Highlighting}[]
\NormalTok{                                                                                    k      est}
\NormalTok{un ensemble infini de points rigides tendant vers l’infini. C’est donc un fermé de Zariski}
\NormalTok{de A1,an possédant une infinité de composantes connexes.}
\NormalTok{   On peut néanmoins développer une théorie des composantes irréductibles en géométrie}
\NormalTok{analytique sur k. Brian Conrad l’a fait dans [43] (dans le cas strictement k{-}analytique)}
\NormalTok{en procédant comme en géométrie complexe, c’est{-}à{-}dire en construisant d’abord le nor{-}}
\NormalTok{malisé X e d’un espace k{-}analytique, puis en définissant les composantes irréductibles de}
\NormalTok{X comme les images des composantes connexes de X.        e Dans ce texte, nous nous ap{-}}
\NormalTok{puierons sur une approche différente présentée dans [49], que nous allons brièvement}
\NormalTok{rappeler.}
\NormalTok{   On y démontre tout d’abord que si X est un espace k{-}analytique irréductible, il est}
\NormalTok{purement de dimension d pour un certain entier d, et tout fermé de Zariski strict de}
\NormalTok{X est alors de dimension}
\NormalTok{                  S          \textless{} d ([49], cor. 4.14) puis que si X s’écrit comme une union}
\NormalTok{localement finie i Xi de fermés de Zariski irréductibles alors X est égal à l’un des Xi .}
\NormalTok{   Si X = M (A) est un espace affinoïde, sa topologie de Zariski est noethérienne et il}
\NormalTok{possède donc un ensemble fini de composantes irréductibles pour cette topologie, que l’on}
\NormalTok{appelle simplement les composantes irréductibles de X. Ce sont les images réciproques}
\NormalTok{des composantes irréductibles de Spec A.}
\NormalTok{   Si X est un espace k{-}analytique quelconque, on appelle composante irréductible de X}
\NormalTok{toute partie de X qui est l’adhérence pour la topologie de Zariski de X d’une composante}
\NormalTok{irréductible d’un domaine affinoïde de X. On démontre ([49], Th. 4.20) que les compo{-}}
\NormalTok{santes irréductibles de X constituent un ensemble localement fini de fermés de Zariski}
\NormalTok{irréductibles de X qui recouvrent X, sont deux à deux non comparables pour l’inclusion}
\NormalTok{et sont maximaux ; plus précisément, tout fermé irréductible de X est contenu dans une}
\NormalTok{composante irréductible}
\NormalTok{                S          de X. Si E est un ensemble de composantes irréductibles de}
\NormalTok{X, la réunion Z∈E Z est un fermé de Zariski de X dont les éléments de E sont les}
\NormalTok{composantes irréductibles.}
\NormalTok{   Si x est un point de X l’entier dimx X est le maximum des dimensions des composantes}
\NormalTok{irréductibles de X contenant x ([49], Lemme 4.21).}
\NormalTok{308                 CHAPITRE 11. RAPPELS DE GÉOMÉTRIE ANALYTIQUE}

\NormalTok{   Soit V un domaine analytique de X. Pour toute composante irréductible Y de X,}
\NormalTok{l’intersection Y ∩ V est une union (éventuellement vide, éventuellement infinie) de com{-}}
\NormalTok{posantes irréductibles de V de même dimension que Y ; pour toute composante irréduc{-}}
\NormalTok{tible T de V , l’adhérence de Zariski Z de T dans X est une composante irréductible de}
\NormalTok{X de dimension dim T , et T est une composante irréductible de Z ∩ V ([49], prop. 4.22).}
\NormalTok{   Si U est un ouvert de Zariski de X les assertions suivantes sont équivalentes ([49],}
\NormalTok{Lemme 4.24)}
\NormalTok{  (i) U est Zariski{-}dense ;}
\NormalTok{  (ii) U est dense pour la topologie usuelle de X ;}
\NormalTok{  (iii) U rencontre toutes les composantes irréductibles de X.}
\NormalTok{Et si elles sont satisfaites, U ∩ V est dense dans V pour tout domaine analytique V de}
\NormalTok{X ([52], Cor. 1.5.13).}
\NormalTok{   Soit L une extension complète de k. Soit (Xi ) la famille des composantes irréductibles}
\NormalTok{de X. Pour tout i l’image réciproque Xi,L de Xi sur XL a une famille finie (Yij )j de}
\NormalTok{composantes irréductibles. Chacune des Yij a même dimension que Xi et se surjecte sur}
\NormalTok{Xi ; de plus (Yij )i,j est la famille des composantes irréductibles de XL ([49], Lemme 4.25}
\NormalTok{et théorème 7.1.16 (i), (ii) et (iii)).}
\NormalTok{   Mentionnons qu’on dispose d’un principe GAGA pour les composantes irréductibles :}
\NormalTok{si X est un schéma localement de type fini sur une algèbre affinoïde A, et si les Xi sont}
\NormalTok{les composantes irréductibles de X , les Xian sont les composantes irréductibles de X an}
\NormalTok{([52], Prop. 2.7.16). Si X est irréductible et si Y est un fermé irréductible de X on a}
\NormalTok{alors}

\NormalTok{                          dim X an = dim Y an + codim(Y , X )}

\NormalTok{([52], Lemma 2.7.10 ; la preuve est un peu plus simple lorsque X = Spec A, voir [48],}
\NormalTok{Prop. 1.11).}


\NormalTok{Lemme 11.2.6. — Soit X un espace k{-}analytique et soit x ∈ X. Les assertions sui{-}}
\NormalTok{vantes sont équivalentes :}
\NormalTok{                                    XZar}
\NormalTok{  (i) l’adhérence de Zariski \{x\}      est une composante irréductible de X ;}
\NormalTok{                           XZar}
\NormalTok{  (ii) la dimension de \{x\}      est égale à dimx X.}
\NormalTok{                             XZar}
\NormalTok{Si elles sont satisfaites, \{x\}      est l’unique composante irréductible de X contenant x.}


\NormalTok{Démonstration. — Choisissons une composante irréductible Z de X contenant x et de}
\NormalTok{                                                 XZar}
\NormalTok{dimension dimx X. Le fermé de Zariski \{x\}             est alors contenu dans la composante}
\NormalTok{irréductible Z, et les assertions (i) et (ii) sont donc toutes deux équivalentes à l’égalité}
\NormalTok{    XZar}
\NormalTok{\{x\}      = Z. Si elles sont satisfaites et si T est une composante irréductible de X}
\NormalTok{                                                                         XZar}
\NormalTok{distincte de Z alors T ∩ Z est un fermé de Zariski strict de Z = \{x\}          , si bien que T}
\NormalTok{ne peut contenir x.}
\NormalTok{                                                                             309}
\NormalTok{ § 11.3. PROPRIÉTÉS LOCALES DES ESPACES ANALYTIQUES ET DES FAISCEAUX COHÉRENTS}

\NormalTok{  § 11.3. PROPRIÉTÉS LOCALES DES ESPACES ANALYTIQUES ET}
\NormalTok{                    DES FAISCEAUX COHÉRENTS}

\NormalTok{11.3.1. — Soit P une des propriétés suivantes des anneaux locaux noethériens : vérifier}
\NormalTok{la condition Rm , être régulier, être Gorenstein, être intersection complète. Soit Q une}
\NormalTok{des propriétés suivantes d’un module de type fini sur un anneau local noethérien : être}
\NormalTok{libre de rang r, vérifier la condition Sn , être de Cohen{-}Macaulay}
\NormalTok{   Soit X un espace k{-}analytique et soit x un point de X. Soit F un faisceau cohérent}
\NormalTok{sur X. Pour tout bon domaine analytique V de X contenant x, nous noterons FV,x la}
\NormalTok{fibre en x de F |V vu comme faisceau sur le site topologique de V , c’est{-}à{-}dire encore}
\NormalTok{la colimite des F (U ) où U parcourt l’ensemble des voisinage ouverts de x dans V . S’il}
\NormalTok{existe un bon domaine analytique V de X contenant x tel que l’anneau local noethérien}
\NormalTok{OV,x (resp. tel que le OV,x {-}module de type fini FV,x ) satisfasse P (resp. Q) ceci reste vrai}
\NormalTok{pour tout bon domaine analytique de X qui contient x (théorème 3.4 de [49]). On dit}
\NormalTok{alors que X (resp. F ) satisfait la propriété P (resp. Q) en x.}
\NormalTok{   Le lieu des point en lesquels la propriété X (resp. F ) satisfait P (resp. Q) est un}
\NormalTok{ouvert de Zariski de X ([49], théorème 4.4).}
\NormalTok{   On dit que X (resp. F ) satisfait la propriété P (resp. Q) s’il la satisfait en chacun de}
\NormalTok{ses points. Si U est un domaine analytique de X contenant x alors X (resp. F ) satisfait}
\NormalTok{P (resp. Q) en x si et seulement si c’est le cas de U (resp. F |U ).}
\NormalTok{   Mentionnons une entorse à ces principes généraux : si F est libre de rang r en tout}
\NormalTok{point de X nous ne dirons évidemment pas que F est libre de rang r, mais qu’il est}
\NormalTok{localement libre de rang r – et d’après un résultat de Berkovich ([9], 1.3.4) déjà rappelé et}
\NormalTok{discuté en 11.1.8, cela équivaut à dire que F est localement libre en tant que OX {-}module}
\NormalTok{sur le site annelé XG .}
\NormalTok{   Lorsque F = OX nous dirons souvent « X satisfait Q en x » plutôt que « OX satisfait}
\NormalTok{Q en x ».}
\NormalTok{   Nous étendrons les conventions ci{-}dessus à la conjonction de différentes propriétés d’al{-}}
\NormalTok{gèbre commutative ; nous parlerons ainsi d’espace réduit (c’est{-}à{-}dire satisfaisant (R0 )}
\NormalTok{et (S1 )) ou normal (c’est{-}à{-}dire satisfaisant (R1 ) et (S2 )) en un point x, ou globalement.}
\NormalTok{Tout espace X possède un normalisé X,    e muni d’un morphisme fini surjectif X    e → X qui}
\NormalTok{induit une bijection entre π0 (X)}
\NormalTok{                                e et l’ensemble des composantes irréductibles de X ([49],}
\NormalTok{Déf. 5.10, Thm. 5.13, Thm. 5.17) ; la normalisation commute à l’analytification ([52],}
\NormalTok{Lemma 2.7.15).}
\NormalTok{   Soit L une extension complète de k et soit π le morphisme structural de XL vers X.}
\NormalTok{Le faisceau cohérent F satisfait Q en x si et seulement si π ∗ F satisfait Q en y ; si si}
\NormalTok{P est la propriété d’être Gorenstein ou d’intersection complète, X satisfait P en x si et}
\NormalTok{seulement si XL la satisfait en y ; si P est la propriété d’être régulier ou de vérifier (Rn )}
\NormalTok{alors X satisfait P en x dès que XL satisfait P en y, et la réciproque est vraie si L est}
\NormalTok{analytiquement séparable sur k ([52], Prop. 2.6.7), mais elle ne l’est pas en général.}
\NormalTok{   Si P est la propriété d’être régulier ou de vérifier (Rm ), ou une propriété impliquant}
\NormalTok{d’une façon ou d’une autre la régularité (comme le fait d’être réduit, ou normal), on}
\NormalTok{est dès lors conduit à introduire la notion de validité géométrique de P : on dit que X}
\NormalTok{satisfait la propriété P géométriquement en x si pour toute extension complète L de k}
\NormalTok{et tout antécédent y de x sur XL , l’espace XL satisfait P en y. Pour que ce soit le cas, il}
\NormalTok{suffit qu’il existe une extension parfaite L de k et un antécédent y de x sur XL tel que}
\NormalTok{XL satisfasse P en y ([52], 2.6.9).}
\NormalTok{310                 CHAPITRE 11. RAPPELS DE GÉOMÉTRIE ANALYTIQUE}

\NormalTok{   Soit X un schéma localement de type fini sur une algèbre affinoïde A et soit G un}
\NormalTok{faisceau cohérent sur X . Soit x un point de X an d’image ξ sur X . Le schéma X (resp.}
\NormalTok{le faisceau G ) satisfait P (resp. Q) en ξ si et seulement si X an (resp. G an ) la satisfait en}
\NormalTok{x ([52], Lemma 2.4.6). En étendant les scalaires au complété d’une clôture algébrique}
\NormalTok{de k, on obtient un principe GAGA analogue pour les propriétés géométriques.}



\NormalTok{                   § 11.4. PROPRIÉTÉS DES MORPHISMES}

\NormalTok{11.4.1. Bord et intérieur. — Soit f : Y → X un morphisme d’espaces k{-}analytiques.}
\NormalTok{On note Int(Y /X) l’intérieur relatif de f et ∂(Y /X) son complémentaire, le bord relatif}
\NormalTok{de f ([9], Definition 1.5.4). Par construction, Int(Y /X) est un ouvert de Y . Si ∂(Y /X) =}
\NormalTok{0, nous dirons que f est sans bord (la terminologie de [9] est closed ). Lorsque f est}
\NormalTok{l’inclusion d’un domaine analytique, Int(Y /X) coïncide avec l’intérieur Y̊ de Y dans X}
\NormalTok{au sens de la topologie générale et ∂(Y /X) coïncide donc avec le bord topologique de Y}
\NormalTok{dans X ([9], Prop. 1.5.4).}
\NormalTok{   Si X est un schéma localement de type fini sur une algèbre affinoïde A, le morphisme}
\NormalTok{structural de X an vers M (A) est sans bord.}
\NormalTok{   Si g : Z → Y et f : Y → X sont deux morphismes d’espaces k{-}analytiques alors ([9],}
\NormalTok{Prop. 1.5.5)}
\NormalTok{(11.4.1.1)                  Int(Z/Y ) ∩ g −1 (Y /X) ⊆ Int(Z/X),}
\NormalTok{et l’on a}
\NormalTok{(11.4.1.2)                  Int(Z/Y ) ∩ g −1 (Y /X) = Int(Z/X),}
\NormalTok{si f est localement séparé ([127], Corollary 5.7), ce qui est notamment le cas dès que Y}
\NormalTok{est séparé.}
\NormalTok{   Lorsque X = M (k), on note Int(Y ) = Int(Y /k) et ∂(Y ) = ∂(Y /k), qu’on appelle}
\NormalTok{l’intérieur et le bord de Y respectivement. Nous dirons que Y est sans bord si ∂(Y ) = ∅.}
\NormalTok{Si Y est sans bord, ou plus généralement si f : Y → X est sans bord et si X est bon,}
\NormalTok{l’espace Y est bon. Si f : Y → X est un morphisme d’espaces k{-}analytiques et que X}
\NormalTok{est séparé, le cas 11.4.1.2 entraîne les égalités}
\NormalTok{(11.4.1.3)    Int(Y /X) ∩ f −1 (Int(X)) = Int(Y ) et ∂(Y ) = ∂(Y /X) ∪ f −1 (∂(X)).}
\NormalTok{   Si Y → X est un morphisme d’espaces k{-}analytiques et si si X ′ est un espace analy{-}}
\NormalTok{tique sur une extension complète de k muni d’un k{-}morphisme vers X, on a l’inclusion}
\NormalTok{Int(Y /X) ×X X ′ ⊂ Int(Y ×X X ′ /X ′ ) (cela découle de la proposition 3.13 de [8]). La}
\NormalTok{classes des morphismes sans bord est en particulier stable par changement de base ana{-}}
\NormalTok{lytique sur k. Elle est stable par composition d’après 11.4.1.2 et est G{-}locale au but}
\NormalTok{([127], Cor. 5.6).}

\NormalTok{Exemple 11.4.2. — 1) Soit r un nombre réel s \textgreater{} 0. Soit D le domaine affinoïde de}
\NormalTok{Ak1,an défini par l’inégalité |T | ⩽ r ; c’est le disque fermé de rayon r au sens de Berkovich,}
\NormalTok{isomorphe à M (k\{T /r\}). Son bord topologique dans Ak1,an est le singleton \{ηr \} (défini}
\NormalTok{à l’exemple 11.1.12). On a donc aussi \{ηr \} = ∂(D/A1,an       k ) = ∂(D), la seconde égalité}
\NormalTok{                            1,an}
\NormalTok{provenant du fait que Ak est sans bord.}
\NormalTok{                                § 11.4. PROPRIÉTÉS DES MORPHISMES                                      311}

\NormalTok{   Si E désigne le domaine analytique fermé (mais non compact) de Ak1,an défini par}
\NormalTok{l’inégalité |T | ⩾ r, le bord topologique de E dans Ak1,an est encore \{ηr \}, si bien que}
\NormalTok{\{ηr \} = ∂(E/A1,ank ) = ∂(E).}
\NormalTok{   2) Soit X un espace k{-}analytique, soit f ∈ O(X) et soit W le domaine analytique de X}
\NormalTok{défini par l’inégalité |f | ⩽ r. La fonction f induit un morphisme de X vers Ak1,an , encore}
\NormalTok{noté f ; avec cette convention, on a W = f −1 (D). Comme A1,an est sans bord, ∂(W ) =}
\NormalTok{∂(W/A1,an ). La formule (11.4.1.3) appliquée au diagramme W → D → A1,an , couplée}
\NormalTok{au fait que ∂(D/A1,an k ) = \{ηr \}, assure que ∂(W ) = f}
\NormalTok{                                                          −1}
\NormalTok{                                                             (ηr ) ∪ (∂(W/D)). Et la même}
\NormalTok{formule appliquée à W → X → M (k) entraîne que ∂(W ) = ∂(W/X) ∪ (∂(X) ∩ W ).}
\NormalTok{Cela démontre que ∂(W ) contient (∂(X) ∩ W ) ∪ f −1 (ηr ). Par ailleurs, la propriété pour}
\NormalTok{un morphisme d’être sans bord étant stable par changement de base, le morphisme}
\NormalTok{W f −1 (ηr ) → X est sans bord ; autrement dit, ∂(W/X) ⊂ f −1 (ηr ). On en déduit}
\NormalTok{(11.4.2.1)                        ∂(W ) = (∂(X) ∩ W ) ∪ f −1 (ηr ).}
\NormalTok{Si W ′ est le domaine analytique de X défini par l’inégalité |f | ⩾ r, on démontre de}
\NormalTok{même que}
\NormalTok{(11.4.2.2)                       ∂(W ′ ) = (∂(X) ∩ W ′ ) ∪ f −1 (ηr ).}

\NormalTok{11.4.3. — Soit f : Y → X un morphisme d’espaces k{-}analytiques. On dit que f est}
\NormalTok{compact si l’image réciproque de toute partie compacte est compacte. Comme tout point}
\NormalTok{possède un voisinage compact, cela revient à dire que f est topologiquement propre et}
\NormalTok{topologiquement séparé.}
\NormalTok{  Suivant [9] (1.5.3), un morphisme est dit propre s’il est compact et sans bord.}
\end{Highlighting}
\end{Shaded}

\subsection{CLD-curve, newline lines
15605--15677}\label{cld-curve-newline-lines-1560515677}

\begin{Shaded}
\begin{Highlighting}[]
\NormalTok{f en y, ou encore dimension relative de Y sur X en y, et l’on note dimy f , la dimension}
\NormalTok{de l’espace H (x){-}analytique f −1 (x) en y. Nous dirons que f est quasi{-}fini en y (resp.}
\NormalTok{que f est quasi{-}fini) si dimy f = 0 (resp. si f est partout de dimension relative nulle).}
\NormalTok{   L’application y 7→ dimy f est semi{-}continue supérieurement pour la topologie de Za{-}}
\NormalTok{riski ([48], Théorème 4.9).}
\NormalTok{   Le morphisme f est fini en un point y de Y si et seulement si f est quasi{-}fini en y et}
\NormalTok{y ∈ Int(Y /X) ([9], Cor. 3.1.10).}
\NormalTok{11.4.6. Morphismes plats. — Soit f : Y → X un morphisme d’espaces k{-}}
\NormalTok{analytiques. Soit y un point de Y et soit x son image sur X.}
\NormalTok{   Supposons tout d’abord que X et Y sont bons. On dit que f est plat en y, ou que}
\NormalTok{Y est plat sur X en y, si pour tout morphisme X ′ → X dont la source X ′ est un bon}
\NormalTok{espace analytique sur k, et tout point y ′ de Y ×X X ′ situé au{-}dessus de y et dont on}
\NormalTok{note x′ l’image sur X ′ , le OX ′ ,x′ {-}module OY ×X X ′ ,y′ est plat.}
\NormalTok{   Ne supposons plus que X et Y sont bons. On dit alors que f est plat en y s’il existe}
\NormalTok{un bon domaine analytique U de X contenant x et un bon domaine analytique V de}
\NormalTok{f −1 (U ) contenant y tels que f |V : V → U soit plat en y ; et c’est alors le cas pour tout}
\NormalTok{tel couple (U, V ) ([52], 4.1.7).}
\NormalTok{   Soit maintenant U un domaine analytique quelconque (non nécessairement bon) de}
\NormalTok{X contenant x, et soit V un domaine analytique quelconque de f −1 (U ) contenant y. Il}
\NormalTok{résulte aussitôt de la définition que f est plat en y si et seulement si f |V : V → U est}
\NormalTok{plat en y.}
\NormalTok{   On dit que f est plat s’il est plat en tout point de Y . La platitude est préservée par}
\NormalTok{composition et changement de base (la définition est conçue pour garantir cette dernière}
\NormalTok{propriété).}
\NormalTok{   Le morphisme structural d’un espace k{-}analytique vers M (k) est plat ([52], Lemma}
\NormalTok{4.1.13).}
\NormalTok{   L’ensemble des points de Y en lesquels f est plat est un ouvert de Zariski de Y ([52],}
\NormalTok{Thm. 10.3.2).}
\NormalTok{   Soit φ : Y → X un morphisme entre schémas de type fini sur une algèbre affinoïde et}
\NormalTok{soit y un point de Y an d’image η sur Y . Le morphisme φan est plat en y si et seulement}
\NormalTok{si φ est plat en η ([52], Lemma 4.2.1 et Prop. 4.2.4).}


\NormalTok{Remarque 11.4.7. — Supposons que X et Y sont bons. Si f est plat en y, l’anneau}
\NormalTok{local OY,y est plat sur OX,x . La réciproque est fausse en général, le problème étant que}
\NormalTok{la platitude de OY,y est plat sur OX,x n’est pas en général préservée par changement}
\NormalTok{                            § 11.4. PROPRIÉTÉS DES MORPHISMES                              313}

\NormalTok{de base bon (un contre{-}exemple dû à Temkin est étudié en détail au paragraphe 4.4 de}
\NormalTok{[52]).}
\NormalTok{   La platitude de OY,y sur OX,x entraîne toutefois la platitude de f en y lorsque y est}
\NormalTok{un point rigide (cf. [52], 4.2.3) ou plus généralement lorsque y appartient à Int(Y /X)}
\NormalTok{([52], Thm. 8.3.4).}

\NormalTok{11.4.8. Morphismes quasi{-}lisses. — Soit f : Y → X un morphisme entre espaces k{-}}
\NormalTok{analytiques. On définit comme en géométrie algébrique le OY {-}module des différentielles}
\NormalTok{de Kähler relatives ΩY /X ([8], §3.3) ; c’est un faisceau cohérent sur Y . On dispose en}
\NormalTok{particulier sur tout espace k{-}analytique X du faisceau cohérent ΩX/k . Si x est un point}
\NormalTok{de X, la tige ΩX/k ⊗ H (x) est de dimension ⩾ dimx X ([49], Lemma 6.2 (i)) et l’on dit}
\NormalTok{que X est quasi{-}lisse en x si cette dimension vaut exactement dimx X. L’espace X est}
\NormalTok{quasi{-}lisse en x si et seulement s’il est géométriquement régulier en x ([49], Prop. 6.3) et}
\NormalTok{si x est un k{-}point, il suffit pour que ce soit le cas que X soit régulier en x ([49], Lemma}
\NormalTok{6.2 (iii) ; l’hypothèse que |k × | ̸= \{1\} figurant dans son énoncé est inutile, et n’est pas}
\NormalTok{utilisée dans la preuve).}
\NormalTok{   Cette notion admet une variante relative : dans [52] (déf. 5.2.4) on définit la notion}
\NormalTok{de morphisme quasi{-}lisse en un point au moyen d’un critère jacobien (et un morphisme}
\NormalTok{est dit quasi{-}lisse s’il l’est en tout point de sa source). On démontre ([52], Thm. 5.3.4)}
\NormalTok{que f : Y → X est quasi{-}lisse en y si et seulement si f est plat en y et la fibre f −1 (f (y))}
\NormalTok{est quasi{-}lisse en y. Si d est un entier le lieu Ud des points de Y en lesquels f est quasi{-}}
\NormalTok{lisse et de dimension relative d est un ouvert de Zariski de Y ([52], Thm. 10.7.2 (2)) et}
\NormalTok{ΩUd /X est G{-}localement libre de rang d ([52], Cor. 5.3.2 ; la preuve dans le cas absolu où}
\NormalTok{X = M (k) est beaucoup plus simple, cf. [49], Prop. 6.6).}
\NormalTok{   On dit que f est quasi{-}étale en y s’il est quasi{-}lisse et quasi{-}fini en y, et que f est}
\NormalTok{quasi{-}étale s’il est quasi{-}étale en tout point de Y .}

\NormalTok{Remarque 11.4.9. — Berkovich a introduit au paragraphe 3 de [10] une notion de}
\NormalTok{morphisme quasi{-}étale ; même si sa définition diffère en apparence de celle que nous}
\NormalTok{donnons ici, les deux sont en fait équivalentes ([52], Lemma 5.4.11).}
\NormalTok{   Berkovich a par ailleurs défini aux paragraphes 3.3 et 3.5 de [9] des notions de mor{-}}
\NormalTok{phisme étale et lisse. Si Y et X sont bons, le morphisme f est lisse en y si et seulement}
\NormalTok{s’il est quasi{-}lisse et intérieur en y ([52], Cor. 5.4.8). Il est vraisemblable que l’hypothèse}
\NormalTok{de bonté est nécessaire pour que cette équivalence soit valable ; on peut toutefois s’en}
\NormalTok{passer en dimension relative nulle : même si l’on ne suppose plus X et Y bons, f est}
\end{Highlighting}
\end{Shaded}

\subsection{CLD-curve, newline lines
15700--15755}\label{cld-curve-newline-lines-1570015755}

\begin{Shaded}
\begin{Highlighting}[]
\NormalTok{  Suivant [9] (1.5.3), un morphisme est dit propre s’il est compact et sans bord.}
\NormalTok{11.4.12. — Rappelons que les propriétés pour f d’être fini (resp. fini et plat, resp. com{-}}
\NormalTok{pact, resp. propre) sont locales sur le but (voir [9], proposition 3.1.8 et corollary 3.2.3 ;}
\NormalTok{[127], corollary 5.6).}



\NormalTok{                § 11.5. UTILISATION DE L’INVARIANT dk (x)}

\NormalTok{   Nous nous proposons ici d’établir un certain nombre de résultats techniques en géo{-}}
\NormalTok{métrie analytique qui nous seront utiles par la suite et mettent en jeu de façon essentielle}
\NormalTok{l’invariant dk (x) introduit en 11.2.3.1.}
\NormalTok{Lemme 11.5.1. — Soit f : Y → X un morphisme d’espaces analytiques. Pour tout}
\NormalTok{x ∈ X on a}
\NormalTok{                                dim Y ⩾ dim Yx + dk (x).}
\NormalTok{En particulier si Yx est non vide et si dim Y ⩽ dk (x) on a dim Y = dk (x) et dim Yx = 0.}
\NormalTok{Démonstration. — On a pour tout y ∈ Yx la minoration}
\NormalTok{                            dim Y ⩾ dk (y) = dH (x) (y) + dk (x).}
\NormalTok{Il vient}
\NormalTok{                    dim Y ⩾ sup dH (x) (y) + dk (x) = dim Yx + dk (x).}
\NormalTok{                              y∈Yx}




\NormalTok{Lemme 11.5.2. — Soit Y un espace k{-}analytique purement de dimension n, soit d un}
\NormalTok{entier et soit f : Y → Adk un morphisme. Soit x un point de Adk . Pour tout point y de Yx ,}
\NormalTok{on a dimy Yx ⩾ n − d. Si de plus dk (x) = d, on a égalité.}
\NormalTok{Démonstration. — Pour démontrer que dimy Yx ⩾ n − d on peut remplacer Y par n’im{-}}
\NormalTok{porte quel domaine analytique de Y contenant y. Et si L est une extension complète de k}
\NormalTok{et y ′ un antécédent de y sur YL dont non note x′ l’image sur XL , il suffit de démontrer que}
\NormalTok{dimy′ Yx′′ ⩾ n − d. On peut donc supposer que Y est affinoïde et que x ∈ Adk (k) ; en com{-}}
\NormalTok{posant f avec une translation, on peut même supposer que y = 0. Notons f1 , . . . , fd les}
\NormalTok{composantes du morphisme f , et posons B = OY (Y ). La fibre Yx est l’espace affinoïde}
\NormalTok{M (B/(f1 , . . . , fd )). D’après le Hauptidealsatz de Krull, les composantes irréductibles}
\NormalTok{de Spec(B/(f1 , . . . , fd )) sont de codimension au plus d dans Spec(B). Par conséquent,}
\NormalTok{les composantes irréductibles de Yx sont de dimension au moins n − d, si bien que}
\NormalTok{dimy Yx ⩾ n − d.}
\NormalTok{   Supposons que dk (x) = d. On a alors}
\NormalTok{                         dimy Yx ⩽ dim Yx ⩽ dim Y − d = n − d,}
\NormalTok{                            § 11.5. UTILISATION DE L’INVARIANT dk (x)                              315}

\NormalTok{où la seconde inégalité provient du lemme 11.5.1. Il vient dimy Yx = n − d.}

\NormalTok{11.5.3. — Soit X un espace k{-}analytique. Nous noterons X \# l’ensemble des points x}
\NormalTok{de X tels que dk (x) = dimx X. Si x est un point de X et si V est un domaine analytique}
\NormalTok{de X contenant x on a dimx V = dimx X, si bien que V \# = X \# ∩ V .}
\NormalTok{  Soit x ∈ X \# . Soit V un domaine analytique de X et soit Y un fermé de Zariski}
\NormalTok{de V contenant x. Comme dk (x) = dimx X on a dimx Y ⩾ dimx X et donc dimx Y =}
\NormalTok{                                                    VZar}
\NormalTok{dimx X = dimx V . Ainsi l’adhérence de Zariski \{x\}       est de dimension dimx V , ce qui}
\NormalTok{                 VZar}
\NormalTok{entraîne que \{x\}      est une composante irréductible de V de dimension dimx X, et c’est}
\NormalTok{alors nécessairement la seule composante irréductible de V contenant x.}
\NormalTok{Lemme 11.5.4. — Soit x un point de X ♯ .}
\NormalTok{  a) L’extension graduée H}
\end{Highlighting}
\end{Shaded}

\subsection{CLD-curve, newline lines
19240--19275}\label{cld-curve-newline-lines-1924019275}

\begin{Shaded}
\begin{Highlighting}[]
\NormalTok{                        dcx (f ) = δOXG ,x ((OXG ,x )red ) = 1/ℓ(OXG ,x ).}
\NormalTok{Autrement dit, le degré calibré de f en x est l’inverse de la multiplicité générique de X}
\NormalTok{le long de l’unique composante irréductible de X qui contient x.}
\NormalTok{Proposition 13.3.7. — Soit f : Y → X un morphisme fini d’espaces k{-}analytiques et}
\NormalTok{soit Z une composante irréductible de X. Soit λ la longueur générique de f∗ OY le long}
\NormalTok{de Z, et soit µ la multiplicité générique de X le long de Z. Pour tout x ∈ Z gen , on a}
\NormalTok{dcx (f ) = λ/µ.}
\NormalTok{  En particulier, si f est l’immersion canonique de Xred dans X, on a dcx (f ) = 1/µ.}
\NormalTok{Démonstration. — Soit x ∈ Z gen ⊂ X gen et soit y1 , . . . , yr les antécédents de x sur Y .}
\NormalTok{Soit V un bon domaine analytique de X contenant x ; posons W = f −1 (V ). On a alors}
\NormalTok{                                  X}
\NormalTok{                       dcx (f ) =   dcyi (f )}
\NormalTok{                                        i}
\NormalTok{                                             1   X}
\NormalTok{                                   =               ℓOXG ,x (OYG ,yi )}
\NormalTok{                                     ℓ(OXG ,x ) i}
\NormalTok{                                     1X}
\NormalTok{                                   =      ℓOXG ,x (OYG ,yi )}
\NormalTok{                                     µ i}
\NormalTok{                                     1         Y}
\NormalTok{                                   = ℓOXG ,x ( OYG ,yi )}
\NormalTok{                                     µ          i}
\NormalTok{                                       λ}
\NormalTok{                                   =     ,}
\NormalTok{                                       µ}
\NormalTok{où la dernière égalité résulte du corollaire 13.1.10 et du fait que f est fini, d’où la}
\NormalTok{proposition.}

\NormalTok{Lemme 13.3.8. — Soit f : Y → X un morphisme entre espaces k{-}analytiques. Soit y}
\NormalTok{un point de Y et soit x son image sur X. Supposons que x ∈ X gen et que f est quasi{-}fini}
\NormalTok{et plat en y. On a alors l’égalité dcy (f ) = degy (f ).}
\NormalTok{Démonstration. — Pour établir le lemme, on peut remplacer X par n’importe quel do{-}}
\NormalTok{maine analytique V de X contenant x et Y par n’importe quel domaine analytique W}
\NormalTok{de f −1 (V ) contenant y. Cela permet, au vu du corollaire 13.1.8, de supposer que X et}
\NormalTok{Y ont affinoïdes et que f est fini. Comme f est plat en y, l’anneau local OY,y est un}
\NormalTok{OX,x {-}module libre de rang degy (f ). En vertu de la proposition 13.1.9, c), OYG ,y est un}
\end{Highlighting}
\end{Shaded}

\subsection{CLD-curve, newline lines
19660--19720}\label{cld-curve-newline-lines-1966019720}

\begin{Shaded}
\begin{Highlighting}[]
\NormalTok{n’est pas diviseur de zéro dans OX (Vi ). Puisque les Vi recouvrent X, il vient a = 0, ce}
\NormalTok{qui est absurde, son annulateur étant égal à p.}
\NormalTok{   Si h est un élément régulier de A qui appartient à I, la fonction méromorphe hf}
\NormalTok{est une section globale de OX (X), donc appartient à A. Par conséquent, f = (hf )/h}
\NormalTok{appartient à Frac(A).}

\NormalTok{Corollaire 13.6.6. — Soit X un bon espace k{-}analytique et soit f une fonction mé{-}}
\NormalTok{romorphe sur X. Pour tout x ∈ X il existe un voisinage ouvert U de x, une fonc{-}}
\NormalTok{tion g ∈ OX (U ) et une fonction h ∈ SX (U ) telles que f |U = g/h.}
\NormalTok{13.6.7. — Soit X un espace analytique et soit f ∈ MX (X) une fonction méromorphe}
\NormalTok{sur X. On dit que f est une fonction méromorphe régulière (cf. [72], 20.1.8) si pour}
\NormalTok{tout domaine affinoïde V de X, la restriction de f à V est un élément régulier de}
\NormalTok{l’anneau MX (V ). Alors, f est un élément inversible de l’anneau MX (X). De plus, il}
\NormalTok{existe un ouvert de Zariski U de X, dont le complémentaire est rare, tel que f définisse}
\NormalTok{une fonction holomorphe inversible sur U .}
\NormalTok{   Observons que le faisceau MX× des fonctions méromorphes régulières est le faisceau}
\NormalTok{en groupes abéliens associé au faisceau des fonctions holomorphes régulières sur X.}



\NormalTok{ § 13.7. DIVISEUR D’UNE FONCTION MÉROMORPHE RÉGULIÈRE}

\NormalTok{13.7.1. — Soit X un espace k{-}analytique. Soit f ∈ SX (X) une fonction analytique}
\NormalTok{régulière sur X. Le sous{-}espace analytique fermé de X d’équation f = 0 est purement}
\NormalTok{de codimension 1, il est vide si et seulement si f est inversible. Notons div(f ) le cycle}
\NormalTok{sur X associé ; ce cycle est effectif et nul si et seulement si f est inversible.}
\NormalTok{  Pour tout sous{-}espace analytique fermé intègre Y de X, notons mY (f ) le coefficient}
\NormalTok{de Y dans div(f ).}
\NormalTok{Lemme 13.7.2. — Soit f une fonction analytique régulière sur X. Soit Y un sous{-}}
\NormalTok{espace analytique fermé intègre de X. Soit V un bon domaine analytique de X et soit}
\NormalTok{y ∈ V ∩ Y tel que l’anneau local OV ∩Y,y soit artinien. On a mY (f ) = long(OV,y /(f )).}
\NormalTok{Démonstration. — Si f (y) ̸= 0, alors Y n’est pas une composante irréductible de \{f =}
\NormalTok{0\}. Dans ce cas, mY (f ) = 0 ; de plus, f est inversible dans OV,y donc long(OV,y /(f )) = 0.}
\NormalTok{                   § 13.8. EXTENSIONS D’ABHYANKAR ET LEUR DÉFAUT                      389}

\NormalTok{   Supposons maintenant f (y) = 0. Comme \{y\} est Zariski dense dans Y (voir les com{-}}
\NormalTok{mentaires après la définition 13.1.3), on a Y ⊂ \{f = 0\}. Puisque Y est purement de}
\NormalTok{codimension 1, c’est une composante irréductible de \{f = 0\}. Par définition, mY (f ) est}
\NormalTok{la longueur générique du faisceau cohérent OX /(f ) le long de Y . D’après la proposi{-}}
\NormalTok{tion 13.2.4, c’est la longueur de l’anneau OV,y /(f ).}

\NormalTok{Lemme 13.7.3. — Si f et g sont deux fonctions holomorphes régulières sur X,}
\NormalTok{alors div(f g) = div(f ) + div(g).}
\NormalTok{Démonstration. — Soit Y un sous{-}espace analytique fermé intègre de X, de codimen{-}}
\NormalTok{sion 1 ; démontrons que mY (f g) = mY (f ) + my (g). Soit V un domaine affinoïde de X}
\NormalTok{et soit y ∈ V ∩ Y tel que OV ∩Y,y soit artinien (il en existe). Notons A = OV,y ; la suite}
\NormalTok{exacte}
\NormalTok{                         0 → gA /(f g) → A (f g) → A /(g) → 0}
\NormalTok{peut se récrire}
\NormalTok{                           0 → A /(f ) → A /f g → A /(g) → 0}
\NormalTok{puisque g n’est pas diviseur de zéro dans A . on a donc l’égalité long(OV,y /(f g)) =}
\NormalTok{long(OV,y /(f )) + long(OV,y /(g)), d’où l’assertion d’après le lemme 13.7.2.}

\NormalTok{Proposition 13.7.4. — a) La formule f 7→ div(f ) définit un morphisme de G{-}}
\NormalTok{faisceaux en monoïdes abéliens du G{-}faisceau des fonctions analytiques régulières sur X}
\NormalTok{dans le G{-}faisceau des cycles effectifs sur X.}
\NormalTok{   b) Ce morphisme de faisceaux se prolonge de manière unique en un morphisme de G{-}}
\NormalTok{faisceaux en groupes abéliens du G{-}faisceau des fonctions méromorphes régulières sur X}
\NormalTok{dans le G{-}faisceau des cycles sur X.}
\NormalTok{Démonstration. — (1) Soit V et W des domaines analytiques de X tels que W ⊂ V .}
\NormalTok{Pour toute fonction analytique régulière f sur V , f |W est une fonction analytique ré{-}}
\end{Highlighting}
\end{Shaded}

\subsection{CLD-curve, newline lines
21280--21445}\label{cld-curve-newline-lines-2128021445}

\begin{Shaded}
\begin{Highlighting}[]

\NormalTok{                § 14.7. VARIATION DES TROPICALISATIONS}

\NormalTok{Lemme 14.7.1. — Soit X un espace analytique compact, soit f : X → Y un mor{-}}
\NormalTok{phisme d’espaces analytiques, soit g : X → Gnm un moment. Pour tout y ∈ Y , on pose}
\NormalTok{Xy = f −1 (y). Soit y un point rigide de Y . Il existe un voisinage U de y dans Y et un}
\NormalTok{polyèdre ∆ ⊂ gtrop (Xy ) de dimension ⩽ dim(Xy ) − 1 tels que gtrop (Xu ) ⊂ gtrop (Xy ) et}
\NormalTok{gtrop (Xu ∩ ∂(X)) ⊂ ∆ pour tout u ∈ U .}
\NormalTok{                                               −1}
\NormalTok{Démonstration. — Le domaine analytique gtrop      (gtrop (Xy )) de X contient le fermé de}
\NormalTok{Zariski Xy , donc en est un voisinage (lemme 11.4.10). Comme X est compact, il existe}
\NormalTok{                                           −1}
\NormalTok{un voisinage U de y tel que f −1 (U ) ⊂ gtrop (gtrop (Xy )). Ainsi, si u ∈ U , gtrop (Xu ) est}
\NormalTok{contenu dans gtrop (Xy ).}
\NormalTok{   Le bord ∂(X) de X est contenu dans une réunion finie de compacts qui sont chacun de}
\NormalTok{la forme φ−1 (ηr ), où φ est un morphisme d’un domaine affinoïde V de X vers A1 et r est}
\NormalTok{un nombre réel \textgreater{} 0 ([50], lemme 3.1). Soit φ : V → A1 un tel morphisme et soit r \textgreater{} 0.}
\NormalTok{418          CHAPITRE 14. CALIBRAGE CANONIQUE D’UN ESPACE ANALYTIQUE}

\NormalTok{                                            −1}
\NormalTok{Posons Q = gtrop (Xy ∩ φ−1 (ηr )) et W = gtrop  (Q) ∩ φ−1 (ηr ). L’espace W est un domaine}
\NormalTok{analytique compact de V ∩ φ−1 (ηr ) et en contient le fermé de Zariski Xy ∩ φ−1 (ηr ). Par}
\NormalTok{suite, W est un voisinage de Xy ∩ φ−1 (ηr ) dans X ∩ φ−1 (ηr ). Puisque X est compact, et}
\NormalTok{quitte à diminuer U , on a gtrop (Xu ∩ φ−1 (ηr )) ⊂ Q pour u ∈ U .}
\NormalTok{  Par ailleurs, démontrons que dim(Q) ⩽ dim(Xy ) − 1. Soit x ∈ Xy ∩ φ−1 (ηr ) ;}
\NormalTok{on a dk (x) ⩽ dim(Xy ) ; d’autre part, dk (x) = dH (ηr ) (x) + 1 ; par conséquent,}
\NormalTok{dH (ηr ) (x) ⩽ dim(Xy ) − 1. Ainsi, la dimension H (ηr ){-}analytique de Xy ∩ φ−1 (ηr )}
\NormalTok{est majorée par dim(Xy ) − 1 et, a fortiori, dim(Q) ⩽ dim(Xy ) − 1.}
\NormalTok{  Le lemme en résulte.}

\NormalTok{Théorème 14.7.2. — Soit d un entier, soit X un espace analytique compact purement}
\NormalTok{de dimension n qui satisfait à la condition (Sd ), soit Y un espace analytique régulier}
\NormalTok{purement de dimension d et soit f : X → Y un morphisme purement de dimension}
\NormalTok{relative n − d ; pour tout t ∈ Y , on pose Xt = f −1 (t).}
\NormalTok{   Soit g : X → GN   m un moment. Soit y un point rigide de Y et soit P le polyèdre calibré}
\NormalTok{           (n−d)}
\NormalTok{gtrop (Xy )      .}
\NormalTok{   a) Il existe un voisinage ouvert U de y dans Y tel que, pour tout point u ∈ U , le}
\NormalTok{polyèdre calibré gtrop (Xu )(n−d) soit égal à P .}
\NormalTok{   b) On suppose de plus que Y = Ad et que y est l’origine. Posons X ′ = f −1 (Gdm ) et}
\NormalTok{soit h : X ′ → GN        d}
\NormalTok{                   m × Gm le moment défini par (g, f ) ; c’est un moment agréable. Il existe}
\NormalTok{un nombre réel r \textgreater{} 0 tel que le polyèdre calibré htrop (X ′ ) ∩ RN × ]−∞, log(r)[d soit égal}
\NormalTok{au produit des polyèdre calibrés P et de la cellule ]−∞, log(r)[d munie de son calibrage}
\NormalTok{standard.}
\NormalTok{   Dans une première version de ce texte, l’énoncé de ce théorème supposait d = 1.}
\NormalTok{Comme l’a noté A. Mihatsch dans [103], la démonstration s’étend à tout d.}
\NormalTok{Démonstration. — D’après le lemme 14.7.1, il existe un voisinage ouvert U de y dans Y}
\NormalTok{et un polyèdre ∆ de dimension ⩽ n − d − 1 tels que gtrop (Xu ) ⊂ gtrop (Xy ) et gtrop (Xu ∩}
\NormalTok{∂(X)) ⊂ ∆ pour tout u ∈ U . Quitte à le remplacer par la composante connexe de y}
\NormalTok{dans U , on suppose que U est connexe.}
\NormalTok{   Choisissons une décomposition cellulaire C de gtrop (Xy ) telle que ∆ soit contenu dans}
\NormalTok{la réunion des (n − d − 1){-}cellules de C .}
\NormalTok{Lemme 14.7.3. — Soit C une (n−d){-}cellule de C . Soit p : GN       m → Gm}
\NormalTok{                                                                           n−d}
\NormalTok{                                                                                un morphisme}
\NormalTok{affine de tores tel que la restriction à C de ptrop soit injective ; soit σC : Rn−d → RN}
\NormalTok{                                                                             −1}
\NormalTok{l’unique section affine de p dont l’image contient C. Posons W = gtrop          (C) ; c’est un}
\NormalTok{domaine de Weierstrass de X, purement de dimension n.}
\NormalTok{                                                                                     n−d}
\NormalTok{   a) Il existe un entier naturel δC ̸= 0 tel que le morphisme (p◦g, f )|W : W → Gm      ×Gdm}
\NormalTok{soit fini et plat de degré δC au{-}dessus de tout point de sqU (ptrop (C̊)).}
\NormalTok{   b) Pour tout u ∈ U , on a C ⊂ gtrop (Xu ).}
\NormalTok{  (Rappelons, voir §12.1.12, que l’ensemble sqU (ptrop (C̊)) est constitué des points de X}
\NormalTok{de la forme squ (c), pour u ∈ U et c ∈ ptrop (D̊).)}
\NormalTok{                                                                       n−d}
\NormalTok{Démonstration. — a) Soit u ∈ U et soit Su le squelette de (Gm              )u , identifié à une}
\NormalTok{partie de Gn−d}
\NormalTok{             m    ×  Y . Soit s un  point de  l’image  canonique  sq u trop (C̊)) de ptrop (C̊)}
\NormalTok{                                                                      (p}
\NormalTok{dans Su . Prouvons que (p ◦ g, f )|Y est fini et plat au{-}dessus de s.}
\NormalTok{  Soit x un point de W tel que (p◦g, f )(x) = s. Le lemme 11.5.1, appliqué au morphisme}
\NormalTok{de Wu dans Gn−dm,u induit par p◦g, affirme que la dimension de ce morphisme en x est nulle.}
\NormalTok{                          § 14.7. VARIATION DES TROPICALISATIONS                             419}

\NormalTok{                                                                      −1}
\NormalTok{Cela entraîne que dimx (p ◦ g, f ) = 0. Puisque ∂(W ) ⊂ ∂(X) ∪ gtrop     (∂C), on a x ̸∈ ∂(W ).}
\NormalTok{Puisque fW est propre, le lemme 11.6.6, a), entraîne que f |W est fini au{-}dessus d’un}
\NormalTok{voisinage de s.}
\NormalTok{                                                                            n−d}
\NormalTok{   Comme dH (t) (s) = n − d, on a dk (s) ⩾ n − d. Les espaces W et Gm           × Y sont tous}
\NormalTok{                                                             n−d}
\NormalTok{deux purement de dimension n, l’espace k{-}analytique Gm × Y est régulier. Comme X}
\NormalTok{vérifie la propriété (Sd ), le lemme 11.5.8 appliqué au morphisme (p ◦ g, f )|W entraîne que}
\NormalTok{ce morphisme est plat au{-}dessus de s.}
\NormalTok{   Comme U est connexe, l’ensemble sqU (ptrop (C̊)) = \{squ (ptrop (C̊))\}u∈U est connexe,}
\NormalTok{d’après 12.1.12. Par suite, le degré de (p ◦ g, f )|W au{-}dessus de s ne dépend pas du}
\NormalTok{point s choisi. Notons{-}le δC .}
\NormalTok{   Par hypothèse, C ⊂ gtrop (Xy ), donc ptrop (C) ⊂ p ◦ gtrop (Wy ) ; d’après le théo{-}}
\NormalTok{                                             (n−d)}
\NormalTok{rème 12.8.2, on a ptrop (C) ⊂ p ◦ gtrop (Σ(p◦g)y ) donc tout point de sqy (ptrop (C̊)) possède}
\NormalTok{un antécédent par p ◦ g. Cela entraîne que δC ̸= 0 et termine la démonstration de}
\NormalTok{l’assertion a).}
\NormalTok{   b) Soit s un point de sqU (ptrop (C̊)) et soit u son image sur U . Par changement de base,}
\NormalTok{le morphisme de Wu dans (Gn−d     m )u déduit de p ◦ g est fini et plat de degré δC au{-}dessus}
\NormalTok{de s.}
\NormalTok{   Comme δC ̸= 0, il s’ensuit que ptrop (C̊) ⊂ p ◦ gtrop (Wu ) pour tout u ∈ U . Par densité}
\NormalTok{de C̊ dans C et compacité de Wu , on a ptrop (C) ⊂ p ◦ gtrop (Wu ) pour tout u ∈ U . Par}
\NormalTok{définition de W , ceci entraîne que C ⊂ gtrop (Wu ) pour tout u ∈ U .}
\NormalTok{  Reprenons la démonstration du théorème 14.7.2.}
\NormalTok{  a) Démontrons d’abord la constance des polyèdres calibrés gtrop (Xu )(n−d) pour u ∈ U .}
\NormalTok{  Soit u ∈ U .}
\NormalTok{  D’après le lemme précédent, toute (n − d){-}cellule de gtrop (Xy ) est contenue}
\NormalTok{dans gtrop (Xu ). Par suite, gtrop (Xy )(n−d) ⊂ gtrop (Xu ). Par définition de U , on a}
\NormalTok{également gtrop (Xu ) ⊂ gtrop (Xy ). Il en résulte que gtrop (X, u)(n−d) = gtrop (Xy )(n−d) .}
\NormalTok{  De plus, le calibrage canonique de gtrop (Xu ) est adapté à la décomposition cellulaire C ,}
\NormalTok{car aucune (n − d){-}cellule ouverte de C ne rencontre gtrop (Xu ∩ ∂(X)).}
\NormalTok{  Soit C une (n − d) cellule de C , soit p : GN  m → Gm}
\NormalTok{                                                         n−d}
\NormalTok{                                                              un morphisme affine de tores tel}
\NormalTok{                                                              n−d}
\NormalTok{que la restriction à C de ptrop soit injective et soit σC : R     → RN l’unique section affine}
\NormalTok{de p dont l’image contient C. Par définition des calibrages canoniques, le vecteur{-}volume}
\NormalTok{µC,u de gtrop (Xu ) sur C est donné par}
\NormalTok{                               µC,u = δC σC,∗ (|e1 ∧ · · · ∧ en−d |).}
\NormalTok{Ainsi, le polyèdre calibré gtrop (Xu )(n−d) est égal à gtrop (Xy )(n−d) = P , ainsi qu’il fallait}
\NormalTok{démontrer.}
\NormalTok{  b) Nous supposons maintenant que Y = Ad et que y = 0 est l’origine.}
\NormalTok{  Le morphisme f : X → Y est topologiquement propre, donc le morphisme f |X ′ : X ′ →}
\NormalTok{Gm est topologiquement propre. Par suite, le moment h′ = (g|X ′ , f |X ′ ) est topologique{-}}
\NormalTok{  d}

\NormalTok{ment propre, et en particulier agréable.}
\NormalTok{  On choisit un nombre réel r \textgreater{} 0 assez petit tel que l’on puisse prendre U = D(0, r)}
\NormalTok{dans a).}
\NormalTok{                                                                               −1}
\NormalTok{  Posons Γ = ]−∞; log(r)[d et notons XΓ = X ′ ∩ f −1 (D(0, r)) = X ′ ∩ ftrop      (Γ).}
\NormalTok{                           ′}
\NormalTok{  Prouvons d’abord que htrop (XΓ ) = P × Γ.}
\NormalTok{  Soit γ ∈ Γ et soit u ∈ D(0, r) tel que trop(u) = γ. Soit C une cellule de dimension (n−}
\NormalTok{d) de C . La cellule C × \{γ\} est contenue dans gtrop (Xu ), d’après a). Par suite, la cellule}
\NormalTok{420           CHAPITRE 14. CALIBRAGE CANONIQUE D’UN ESPACE ANALYTIQUE}

\NormalTok{C × Γ est contenue dans h′trop (XΓ ). Comme elle est de dimension n, elle est contenue}
\NormalTok{dans h′trop (XΓ )(n) . Inversement, la définition de U au début de la démonstration garantit}
\NormalTok{que pour tout u ∈ D(0, r)∩Gdm , gtrop (Xu′ ) ⊂ gtrop (Xy ). Par suite, h′trop (XΓ ) ⊂ gtrop (Xy )×}
\NormalTok{Γ. On a donc h′trop (XΓ )(n) ⊂ gtrop (Xy )(n−d) × Γ. On a donc l’égalité voulue.}
\NormalTok{  Démontrons maintenant l’égalité des calibrages. Soit C une (n−d){-}cellule de C . notons}
\NormalTok{                   d        n}
\NormalTok{aussi q : GN m ×Gm → Gm le morphisme de tores (p, id). L’application affine τC = (σC , id)}
\NormalTok{est l’unique section affine de qtrop dont l’image contient C × Γ. Nous allons démontrer}
\NormalTok{que le calibrage de la cellule C × Γ de h′trop (XΓ )(n) est au produit du calibrage µC de C}
\NormalTok{dans P par le calibrage standard de Γ.}
\NormalTok{  Posons XC,Γ = (h′trop )−1 (C × Γ) ; nous allons démontrer que le morphisme q ◦ h′ |XC,Γ}
\NormalTok{est fini et plat de degré δC au{-}dessus de tout point de sq(qtrop (C̊ × Γ)).}
\NormalTok{  Soit c ∈ ptrop (C̊) et γ ∈ Γ Par la construction de la décomposition cellulaire C et le}
\NormalTok{choix de r, le point sq(c, γ) n’appartient pas à q ◦ h′ (∂(XC,Γ )). D’après le lemme 11.6.6,}
\NormalTok{c), le morphisme q ◦ h′ |XC,Γ est fini et plat au{-}dessus de sq(c, γ) ; notons δ son degré.}
\NormalTok{  On a q ◦ h′ = (p, id) ◦ (g, f )|X ′ = (p ◦ g, f )|X ′ . Par changement de base, le mor{-}}
\NormalTok{phisme de XC,Γ,sq(γ) dans Gn−d  m,H (sq(γ)) déduit de p ◦ g est fini et plat de degré δ au{-}dessus}
\NormalTok{de sqH sq(γ) (sq(γ)). D’après la première partie de l’énoncé et le choix de r, ce degré est}
\NormalTok{égal à δC .}
\NormalTok{  Par conséquent (cf. §14.5.2), le calibrage de la cellule C × Γ est constant, de vecteur{-}}
\NormalTok{volume}
\NormalTok{                              µC×Γ = δC τC,∗ (|e1 ∧ · · · ∧ en−1 ∧ en |).}
\NormalTok{De plus,}
\NormalTok{         τC,∗ (|e1 ∧ · · · ∧ en−1 ∧ en )| = σC,∗ (|e1 ∧ · · · ∧ en−d |) · |en−d+1 ∧ · · · ∧ en |,}
\NormalTok{qui est le produit du vecteur{-}volume associé à C dans P et du vecteur{-}volume associé}
\NormalTok{au calibrage standard de Γ. Ceci achève la démonstration du théorème.}

\NormalTok{Remarque 14.7.4. — Le théorème 14.7.2 ne peut pas être vrai sans certaines hypo{-}}
\NormalTok{thèses sur le morphisme f : X → Y .}
\NormalTok{   Considérons par exemple le cas où f est la normalisation d’une courbe lisse en dehors}
\NormalTok{d’un unique point double ordinaire y et g est le moment trivial (N = 0). Dans ce cas,}
\NormalTok{la tropicalisation de Xy est un point avec multiplicité 2. En revanche, pour tout autre}
\end{Highlighting}
\end{Shaded}

\subsection{CLD-curve, newline lines
23340--23453}\label{cld-curve-newline-lines-2334023453}

\begin{Shaded}
\begin{Highlighting}[]

\NormalTok{  Soit X un n{-}espace k{-}analytique, considéré comme espace tropical de niveau n.}
\NormalTok{16.6.1. — Soit Y un sous{-}espace analytique fermé de X, purement de dimension p. On}
\NormalTok{le considère comme un espace tropical de niveau p et on note j : Y → X l’immersion}
\NormalTok{fermée canonique. Le courant j∗ δY sur X est de bidimension (p, p) ; on le note plus}
\NormalTok{simplement δY .}
\NormalTok{                                        P}
\NormalTok{16.6.2. — Soit p un entier et soit z =     z(Y )Y un cycle purement p{-}dimensionnel}
\NormalTok{sur X. Soit ω ∈ Ac (X) ; posons}
\NormalTok{                  p,p}

\NormalTok{                                    Z      X        Z}
\NormalTok{                          ⟨δz , ω⟩ = ω =       z(Y ) ω.}
\NormalTok{                                         z       Y            Y}


\NormalTok{C’est une somme finie : en effet, comme ω est à support compact et la famille (z(Y )) est}
\NormalTok{                                      R de Zariski intègres Y de X tels que z(Y ) ̸= 0 et}
\NormalTok{localement finie, l’ensemble des fermés}
\NormalTok{ω|Y ̸= 0 est fini. L’application ω 7→ z ω est un courant : si l’on se restreint aux formes}
\NormalTok{à support dans un domaine analytique compact donné V , il s’agit de la somme finie des}
\NormalTok{courants z(Y )δY , où Y parcourt l’ensemble fini des fermés de Zariski intègres Y de X}
\NormalTok{qui rencontrent V et tels que z(Y ) ̸= 0.}
\NormalTok{   Lorsqu’on se donne un sous{-}espace analytique fermé de X, purement de dimension p,}
\NormalTok{on peut considérer d’une part l’intégration sur ce sous{-}espace et d’autre part l’intégration}
\NormalTok{sur le cycle associé. La proposition suivante affirme que ces deux opérations coïncident.}

\NormalTok{Proposition 16.6.3. — Soit Y un sous{-}espace analytique fermé de X, purement de}
\NormalTok{dimension p. On a δY = δ[Y ] .}

\NormalTok{Démonstration. — Notons j l’immersion   R de     Y dans X. Soit ω une (p, p){-}forme à support}
\NormalTok{                                              ∗}
\NormalTok{compact sur X. On a donc ⟨δY , ω⟩ = Y j ω, ce qui permet de supposer que Y = X.}
\NormalTok{  Soit (Xi ) la famille des composantes irréductibles de X et, pour tout i, soit mi la}
\NormalTok{multiplicité de X le long de Xi . Soit X ′ l’ouvert de Zariski de X constitué des points qui}
\NormalTok{sont situés sur une unique composante irréductible de X ; il contient le support de ω.}
\NormalTok{Pour tout i, posons Xi′ = X ′ ∩ Xi ; c’est un ouvert de X et l’on a}
\NormalTok{                             Z            Z              Z}
\NormalTok{                                                       1}
\NormalTok{                                     ω=            ω=        ω,}
\NormalTok{                              Xi,red          ′}
\NormalTok{                                            Xi,red    mi Xi′}
\NormalTok{                                                              ′}
\NormalTok{en vertu de la proposition 15.4.3 appliquée à l’immersion de Xi,red dans Xi′ , dont le degré}
\NormalTok{calibré est 1/mi (proposition 13.3.7). Alors, on a}
\NormalTok{              Z      Z        XZ          X Z               X Z}
\NormalTok{   ⟨δX , ω⟩ =    ω=      ω=          ω=        mi      ω=       mi        ω = ⟨δ[X] , ω⟩,}
\NormalTok{              X       X′       i   Xi′       i}
\NormalTok{                                                      ′}
\NormalTok{                                                     Xi,red       i   Xi,red}


\NormalTok{comme il fallait démontrer.}

\NormalTok{                     z(Y )Y , on définit aussi δz∂ =   z(Y )δY∂ . Si Y est un sous{-}espace}
\NormalTok{                   P                                 P}
\NormalTok{16.6.4. — Si z =}
\NormalTok{analytique fermé de X, on démontre comme pour la proposition 16.6.3 que δY∂ = δ[Y    ∂}
\NormalTok{                                                                                       ].}
\NormalTok{460              CHAPITRE 16. COURANTS SUR LES ESPACES ANALYTIQUES}

\NormalTok{               § 16.7. LA FORMULE DE POINCARÉ–LELONG}

\NormalTok{16.7.1. — Soit X un n{-}espace k{-}analytique sans bord ; on le considère comme un espace}
\NormalTok{tropical de niveau n.}
\NormalTok{   Soit f une fonction méromorphe régulière f sur X. Le courant d’intégration δdiv(f ) sur}
\NormalTok{le cycle div(f ) est un courant symétrique de bidimension (n − 1, n − 1) sur X, donc de}
\NormalTok{bidegré (1, 1) ; il est d′ {-}fermé et d′′ {-}fermé (proposition 7.7.9).}
\NormalTok{   La formule de Poincaré–Lelong (théorème 16.7.3) va entraîner qu’il est d′ d′′ {-}exact.}
\NormalTok{16.7.2. — Soit f une fonction méromorphe régulière sur X. Soit U le sous{-}espace}
\NormalTok{de X formé des points au voisinage desquels f est définie et non nulle ; c’est un ouvert}
\NormalTok{de Zariski de X sur lequel la fonction log |f | est lisse, et il est Zariski dense. Pour tout}
\NormalTok{point x du fermé complémentaire, on a dk (x) \textless{} dim(X). D’après le paragraphe 7.1.12,}
\NormalTok{le courant [log|f |] sur U se prolonge un unique courant de bidegré (0, 0) sur X que nous}
\NormalTok{noterons encore [log|f |].}
\NormalTok{Théorème 16.7.3 (Formule de Poincaré–Lelong). — Soit X un n{-}espace k{-}}
\NormalTok{analytique sans bord et soit f une fonction méromorphe régulière sur X. On a l’égalité}
\NormalTok{                                   δdiv(f ) = d′ d′′ [log|f |].}
\NormalTok{   L’exemple 16.1.4 montre que l’hypothèse de régularité sur f est nécessaire, même si}
\NormalTok{le complémentaire de son lieu d’inversibilité est de dimension \textless{} n.}
\NormalTok{Lemme 16.7.4. — L’égalité ci{-}dessus est vraie sur l’ouvert d’inversibilité de f .}
\NormalTok{Démonstration. — On peut supposer f inversible. On a alors div f = 0 et le terme}
\NormalTok{de gauche de l’égalité est donc nul. On a par ailleurs d′ d′′ x = 0 dans A 1,1 (R) ; il}
\NormalTok{vient d′ d′′ log|f | = 0.}

\NormalTok{Lemme 16.7.5. — Il suffit de démontrer le théorème 16.7.3 si l’on suppose de plus que}
\NormalTok{X satisfait à la propriété (S1 ) et que f est une fonction holomorphe régulière qui définit}
\NormalTok{un morphisme purement de dimension relative n − 1 de X vers A1 .}
\NormalTok{Démonstration. — Le théorème est local sur X ; il suffit donc de prouver que tout point x}
\NormalTok{de X a un voisinage sur lequel cette égalité de courant est vérifiée. On peut ainsi com{-}}
\NormalTok{mencer par supposer que f est quotient de deux fonctions holomorphes régulières. Par}
\NormalTok{additivité des deux termes de l’égalité à établir, on peut donc supposer que f est ho{-}}
\NormalTok{lomorphe. Pour alléger les notations, nous noterons encore f le morphisme X → A1}
\NormalTok{qu’elle induit.}
\NormalTok{   Si f (x) ̸= 0, il existe un voisinage U de x sur lequel f est inversible, et l’égalité à}
\NormalTok{établir vaut sur U en vertu du lemme 16.7.4.}
\NormalTok{   Supposons que f (x) = 0. Comme la fonction f est régulière, le morphisme f qu’elle}
\NormalTok{induit est de dimension relative n − 1 en x. Par suite, ce morphisme est de dimension ⩽}
\NormalTok{n − 1 au voisinage de x ([48], th. 4.9). En vertu du Hauptidealsatz (lemme 11.5.2), le}
\NormalTok{morphisme f est donc purement de dimension relative n − 1 au voisinage de x.}
\NormalTok{   Démontrons que l’espace X a la propriété (S1 ) en x. Soit V un domaine affinoïde de X}
\NormalTok{contenant x. Le germe de f |V en x est un élément de l’idéal maximal de OV,x et est un}
\NormalTok{élément régulier de cet anneau. En conséquence, l’anneau local OV,x a la propriété (S1 ),}
\NormalTok{si bien que X a la propriété (S1 ) en x. Il existe dès lors un voisinage ouvert U de x}
\NormalTok{                          § 16.7. LA FORMULE DE POINCARÉ–LELONG                               461}

\NormalTok{dans X tel que U satisfasse la propriété (S1 ) et tel que f |U soit purement de dimension}
\NormalTok{relative n − 1.}
\NormalTok{   Le lemme précédent permet de supposer à partir de maintenant que X a la pro{-}}
\NormalTok{priété (S1 ) et que f est une fonction holomorphe régulière induisant un morphisme}
\NormalTok{purement de dimension relative n − 1 de X dans A1 .}
\NormalTok{   Les deux courants δdiv f et d′ d′′ [log|f |] sont symétriques ; pour s’assurer de leur égalité,}
\NormalTok{il suffit de les appliquer à une forme test symétrique.}
\NormalTok{   Soit donc ω une forme lisse de type (n − 1, n − 1) sur X, symétrique et à support}
\end{Highlighting}
\end{Shaded}

\subsection{CLD-curve, newline lines
23454--23586}\label{cld-curve-newline-lines-2345423586}

\begin{Shaded}
\begin{Highlighting}[]
\NormalTok{compact. Comme le support de ω est compact, il existe une famille finie (Vi )1⩽i⩽m de}
\NormalTok{domaines analytiques compacts de X tropicalisant ω et telle tout point du support de ω}
\NormalTok{soit contenu dans l’intérieur de l’un des Vi . Notons V la réunion des Vi . C’est un domaine}
\NormalTok{analytique compact de X. Par construction, le support de ω est contenu dans l’intérieur}
\NormalTok{topologique V̊ de V dans X, qui coïncide avec Int(V /X).                                 T}
\NormalTok{   Pour toute partie non vide E de \{1, . . . , m\} et tout s \textgreater{} 0, on pose VE = i∈E Vi ,}
\NormalTok{V0 = V ∩ \{f = 0\}, V⩾s = V ∩ \{|f | ⩾ s\}, VE,0 = VE ∩ V0 et VE,⩾s = VE ∩ V⩾s . L’espace}
\NormalTok{analytique VE est intersection d’une famille non vide de domaines analytiques compacts}
\NormalTok{de X qui tropicalisent ω ; il possède en particulier la propriété (S1 ). De plus ω est tropicale}
\NormalTok{sur VE ; soit gE : VE → T un moment et soit αE une (n − 1, n − 1){-}forme sur gE,trop (VE )}
\NormalTok{telle que ω|VE = g ∗ αE ; posons aussi hE = (gE , f ) : VE → T × A1 . Choisissons un nombre}
\NormalTok{réel r \textgreater{} 0 pour lequel les conclusions du théorème 14.7.2 sont satisfaites par chacun des}
\NormalTok{couples (VE , gE ).}
\NormalTok{   Soit (u1 , u2 ) une partition de l’unité lisse de R, subordonnée au recouvrement ou{-}}
\NormalTok{vert (]−∞, log r[, ]log r − 1, +∞[). Soit t ∈ ]0, r[ tel que le support de u1 soit disjoint}
\NormalTok{de [log t, log r[. On note v1 (resp. v2 ) la fonction lisse sur X égale à u1 (log|f |) (resp.}
\NormalTok{u2 (log|f |)) sur l’ouvert d’inversibilité de f , et à 1 (resp. 0) au voisinage de \{f = 0\}.}
\NormalTok{Posons ω1 = v1 ω et ω2 = v2 ω. Il suffit de démontrer que}
\NormalTok{                                     Z                                   Z}
\NormalTok{                   ′ ′′                                ′ ′′}
\NormalTok{                ⟨d d log|f |, ω1 ⟩ =          ω1 et ⟨d d log|f |, ω2 ⟩ =        ω2 .}
\NormalTok{                                     \{f =0\}                              \{f =0\}}

\NormalTok{   Par construction, le support de ω2 est contenu dans l’ouvert \{f ̸= 0\} sur lequel}
\NormalTok{d′ d′′ log|f | s’annule ; les deux termes de la seconde égalité sont donc nuls, et il suffit}
\NormalTok{désormais d’établir la première.}
\NormalTok{   Comme d′′ ω1 est une forme de type (n−1, n), son support ne rencontre aucun fermé de}
\NormalTok{Zariski de dimension \textless{} n (proposition 3.3.6) ; en particulier, il ne rencontre pas f −1 (0).}
\NormalTok{Comme ce support est compact et la fonction |f | continue, on a inf supp(d′′ ω1 ) (|f |) \textgreater{} 0 ;}
\NormalTok{soit s ∈ ]0, t[ tel que |f | \textgreater{} s en tout point du support de d′′ ω1 . On suppose aussi que}
\NormalTok{u1 ≡ 1 dans un voisinage de ]−∞, log(s)]. En particulier, le support de d′′ ω est contenu}
\NormalTok{dans l’intérieur Int(V⩾s /X) = Int(V⩾s ) du domaine analytique fermé V⩾s de V .}
\NormalTok{   Par définition du courant log|f | et compte{-}tenu du fait que V⩾s contient le support de}
\NormalTok{ ′ ′′}
\NormalTok{d d ω1 on a}
\NormalTok{                                                   Z}
\NormalTok{                                ′ ′′}
\NormalTok{                              ⟨d d log|f |, ω1 ⟩ =   log|f | d′ d′′ ω1 .}
\NormalTok{                                                  V⩾s}

\NormalTok{Puisque ω1 est symétrique, la formule de Green (proposition 5.11.1) assure que}
\NormalTok{  Z                    Z                       Z ∂                Z ∂}
\NormalTok{               ′ ′′          ′ ′′                           ′′}
\NormalTok{      log|f | d d ω1 =    (d d log|f |) ∧ ω1 +     log|f | d ω1 −     (d′′ log|f |) ∧ ω1 .}
\NormalTok{    V⩾s                    V⩾s                          V⩾s               V⩾s}
\NormalTok{462               CHAPITRE 16. COURANTS SUR LES ESPACES ANALYTIQUES}

\NormalTok{Comme f est inversible sur V⩾s , on a d′ d′′ log|f | = 0 sur V⩾s ; de plus, le support}
\NormalTok{de d′′ ω1 ne rencontre pas ∂(V⩾s ) ; les deux premiers termes du membre de droite de}
\NormalTok{l’égalité précédente sont donc nuls, d’où l’égalité}
\NormalTok{                        Z                     Z ∂}
\NormalTok{                                    ′ ′′}
\NormalTok{                           log|f | d d ω1 = −       (d′′ log|f |) ∧ ω1 .}
\NormalTok{                           X                       V⩾s}

\NormalTok{  Comme}
\NormalTok{     R    le support}
\NormalTok{             R       de ω1 |X0 est contenu dans V0 et comme ω1 = ω au voisinage de V0 ,}
\NormalTok{on a V0 ω = X0 ω1 . Pour terminer la preuve de la formule de Poincaré–Lelong, il suffit}
\NormalTok{donc de démontrer que l’on a}
\NormalTok{                             Z ∂                       Z}
\NormalTok{                                   ′′}
\NormalTok{                                  d (log|f |) ∧ ω1 = −   ω.}
\NormalTok{                                   V⩾s                        V0}

\NormalTok{  La famille (Vi,0 ) est un G{-}recouvrement fini de V0 ; on a donc l’égalité}
\NormalTok{                             Z       X                 Z}
\NormalTok{                                             Card(E)+1}
\NormalTok{                                 ω=    (−1)                 ω.}
\NormalTok{                                  V0      E̸=∅                VE,0}


\NormalTok{Par ailleurs, la famille (Vi,⩾s ) est un G{-}recouvrement fini de V⩾s ; la proposition 5.9.5}
\NormalTok{entraîne alors l’égalité :}
\NormalTok{              Z ∂                      X                Z ∂}
\NormalTok{                    ′′}
\NormalTok{                   d (log|f |) ∧ ω1 =     (−1)Card(E)+1}
\NormalTok{                                                             d′′ (log|f |) ∧ ω1 .}
\NormalTok{                V⩾s                       E̸=∅                VE,⩾s}


\NormalTok{Comme il n’y a qu’un nombre fini de telles parties E, il suffit de démontrer que pour}
\NormalTok{toute partie non vide E de \{1, . . . , m\}, on a}
\NormalTok{                          Z ∂                           Z}
\NormalTok{                                      ′′}
\NormalTok{                                   d (log|f |) ∧ ω1 = −   ω.}
\NormalTok{                                  VE,⩾s                       VE,0}

\NormalTok{Fixons donc une telle E. Posons VE,[s,t] = VE ∩ \{s ⩽ |f | ⩽ t\} ; comme le support}
\NormalTok{de ω1 |VE,⩾s est contenu dans l’intérieur de VE,[s,t] relativement à VE,⩾s , on a}
\NormalTok{                        Z ∂                    Z ∂}
\NormalTok{                               ′′}
\NormalTok{                             d log|f | ∧ ω1 =            d′′ log|f | ∧ ω1 .}
\NormalTok{                          VE,⩾s                    VE,[s,t]}

\NormalTok{Nous allons calculer cette dernière intégrale par voie tropicale. Posons σ = log(s) et}
\NormalTok{τ = log(t). La restriction de h à VE,[s,t] est un moment ; on a ω1 |VE,[s,t] = h∗E (u1 αE ) et}
\NormalTok{d′′ log|f | ∧ ω1 |VE,[s,t] = h∗E (d′′ t ∧ u1 αE ).}
\NormalTok{   Notons (P, µ) le polyèdre calibré gtrop (VE,0 )(n−1) et e le vecteur (0, 1) de Ttrop ×Gm,trop ;}
\NormalTok{notons (Q, ν) le polyèdre calibré htrop (VE,[s,t] )(n) . Par le choix de r fait au début de la}
\NormalTok{preuve, on a (Q, ν) = (P × [σ, τ ], µ × λ) où λ désigne le calibrage standard de [σ, τ ]. En}
\NormalTok{                R∂                            R∂}
\NormalTok{conséquence VE,[s,t] d′′ log|f | ∧ ω1 = Q ⟨d′′ t ∧ u1 α, ν⟩.}
\NormalTok{   Choisissons une décomposition cellulaire C de P adaptée à son calibrage. L’ensemble}
\NormalTok{des partie de la forme C × \{σ\}, C × \{τ \} et C × [σ, τ ], où C ∈ C , est une décomposition}
\NormalTok{cellulaire de P × [στ ] adaptée à son calibrage. est la réunion de P × \{σ\}, Si F est}
\NormalTok{une (n − 1){-}cellule de P la discordance ∂νF ×\{σ\} s’identifie au produit de (−µF ) par le}
\NormalTok{vecteur{-}volume |e| ; comme d′′ t ∧ α = α ∧ d′′ t et u1 (σ) = 1, la contribution de F × \{σ\}}
\NormalTok{                         § 16.8. SUPPORT DES COURANTS POSITIFS FERMÉS                            463}
\NormalTok{    R∂}
\NormalTok{à   Q}
\NormalTok{         ⟨d′′ t ∧ u1 α, µ⟩ est égale à}
\NormalTok{                                Z                                Z}
\NormalTok{                                                ′′}
\NormalTok{                              −        ⟨u1 α ∧ d t, µF · |e|⟩ = − ⟨α, µF ⟩,}
\NormalTok{                               F ×\{σ\}                            F}
\NormalTok{                                                                                   R}
\NormalTok{si bien que la contribution}
\NormalTok{                 R           totale de P × \{σ\} à l’intégrale étudiée est égale à − P ⟨α, µ⟩}
\NormalTok{soit encore à − VE,0 ω.}
\NormalTok{   La contribution de P × \{τ \} est quant à elle nulle car la fonction u1 est nulle au}
\NormalTok{voisinage de τ . Enfin, la contribution de ∂(P ) × [σ, τ ] elle est également}
\NormalTok{                                                                           R nulle : elle est}
\NormalTok{en effet égale par construction à une somme d’intégrales de la forme F ×[σ,τ ] u1 η où F}
\NormalTok{est une (n − 2){-}cellule orientée de C et η une (n − 1){-}forme différentielle sur F × [σ, τ ]}
\NormalTok{qui provient en réalité d’une (n − 1){-}forme différentielle sur F (l’élément différentiel d′′ t}
\NormalTok{ayant été contracté) et est dès lors nulle.}
\NormalTok{   Il vient                   Z ∂                       Z}
\NormalTok{                                    ′′}
\NormalTok{                                   d (log|f |) ∧ ω1 = −       ω,}
\NormalTok{                                   V⩾s                       VE,0}
\NormalTok{ce qui conclut la démonstration du théorème 16.7.3.}
\end{Highlighting}
\end{Shaded}

\subsection{CLD-curve, newline lines
7910--8188}\label{cld-curve-newline-lines-79108188}

\begin{Shaded}
\begin{Highlighting}[]

\NormalTok{  Soit X un espace tropical de niveau n.}
\NormalTok{Définition 7.2.1. — Soit p un entier. Un courant S de bidimension (p, p) est dit}
\NormalTok{faiblement positif, respectivement positif, respectivement fortement positif, s’il est sy{-}}
\NormalTok{métrique et si pour toute forme α ∈ Acp,p (X) qui est fortement positive, respectivement}
\NormalTok{positive, respectivement faiblement positive, on a ⟨S, α⟩ ⩾ 0.}
\NormalTok{7.2.2. — Les trois notions de positivité pour les formes de type (0, 0) et (n, n) coïn{-}}
\NormalTok{cident ; il en est donc de même pour les courants. Comme les différentes notions de}
\NormalTok{positivité coïncident pour les formes symétriques de type (1, 1) et (n − 1, n − 1), il en}
\NormalTok{est de même pour les courants symétriques.}
\NormalTok{Remarque 7.2.3. — Soit T un courant tel que T et −T sont positifs ; démontrons}
\NormalTok{que T = 0. Soit ω une forme lisse à support compact sur X ; prouvons que ⟨T, ω⟩ = 0.}
\NormalTok{Si ω est antisymétrique, cela découle de la symétrie de T ; supposons donc que ω est}
\NormalTok{symétrique. Il existe aussi des formes positives α1 et α2 , lisses à support compact telles}
\NormalTok{que ω = α1 − α2 (proposition 6.5.2). Par linéarité, on peut donc supposer que ω est}
\NormalTok{positive. Mais alors, ⟨T, ω⟩ est à la fois positif (T est positif) et négatif (−T est positif),}
\NormalTok{donc nul.}
\NormalTok{Lemme 7.2.4. — Supposons que l’espace tropical X soit localement différentiellement}
\NormalTok{séparé. Les courants de bidimension (p, p) positifs ( resp. faiblement positifs, resp. for{-}}
\NormalTok{tement positifs) forment un sous{-}faisceau du faisceau Dp,p (X) des courants de bidimen{-}}
\NormalTok{sion (p, p).}
\NormalTok{Démonstration. — Soit S un courant de bidimension (p, p) sur X.}
\NormalTok{   Supposons que S soit positif. Soit U un ouvert de X. Pour toute forme ω sur U , de}
\NormalTok{type (p, p), positive et à support compact, le prolongement par 0 de ω est une forme}
\NormalTok{positive à support compact sur X. Il en résulte que S|U est positif.}
\NormalTok{   Soit maintenant (Ui ) un recouvrement ouvert de X tel que S|Ui soit positif pour tout i ;}
\NormalTok{démontrons que S est positif. Soit ω une forme de type (p, p) sur X, positive et à support}
\NormalTok{compact. Reprenons la preuve que les faisceaux forment un courant (proposition 7.1.7).}
\NormalTok{Il existe une famille finie (φi ) de fonctions lisses, positives et à support compact sur X}
\NormalTok{                                  § 7.2. COURANTS POSITIFS                                 161}
\NormalTok{                                                                              P}
\NormalTok{telle que pour tout i, le support de φi soit contenu dans P Ui , et telle que   φi soit égale}
\NormalTok{à 1 au voisinage du support de ω. On a de plus ⟨S, ω⟩ = ⟨S, φi ω⟩ ⩾ 0. Pour tout i, on}
\NormalTok{a a ⟨S, φi ω⟩ ⩾ 0, car φi ω est une forme positive dont le support est compact et contenu}
\NormalTok{dans Ui . Cela démontre que S est positif.}
\NormalTok{   Appliqué à chaque ouvert de X, cet argument démontre que les courants positifs de bi{-}}
\NormalTok{dimension (p, p) forment un sous{-}faisceau du faisceau des courants de bidimension (p, p).}
\NormalTok{Les cas des courants faiblement positifs et fortement positifs se démontrent de même.}

\NormalTok{Lemme 7.2.5. — Supposons que X soit localement différentiellement séparé. Soit T un}
\NormalTok{courant sur X. On suppose que pour tout ouvert U de X, toute forme tropicale ω sur U}
\NormalTok{qui est positive ( resp. faiblement positive, resp. fortement positive) et toute fonction lisse}
\NormalTok{positive ρ dont le support est compact et contenu dans U , on a ⟨T, ρω⟩ ⩾ 0. Alors T est}
\NormalTok{positif ( resp. fortement positif, resp. faiblement positif ).}
\NormalTok{Démonstration. — On traite le cas des formes positives, les deux autres sont analogues.}
\NormalTok{Soit ω une forme lisse positive à support compact sur X ; démontrons que ⟨T, ω⟩ ⩾ 0.}
\NormalTok{Pour tout x ∈ X, soit Ux un voisinage ouvert de X sur lequel ω est tropicale. Soit}
\NormalTok{                                                                       que, pour tout x ∈ X,}
\NormalTok{(λx )x∈X une famille à support fini de fonctions lisses positives telleP}
\NormalTok{le support de λx soit compact et contenu dans Ux , et telle que   P      λx soit égale à 1 au}
\NormalTok{voisinage du support de ω (corollaire 6.2.8). On a alors ω = x∈X λx ω. Par hypothèse,}
\NormalTok{⟨T, λx ω⟩ ⩾ 0 pour tout x. Donc ⟨T, ω⟩ ⩾ 0.}

\NormalTok{7.2.6. — L’image directe d’un courant positif (resp. faiblement positif, fortement po{-}}
\NormalTok{sitif) par un morphisme topologiquement propre d’espaces tropicaux est de même type.}
\NormalTok{Cela résulte du fait que la même propriété vaut pour l’image inverse des formes (§3.5.2).}
\NormalTok{Lemme 7.2.7. — Soit α une forme lisse symétrique sur X et soit T un courant symé{-}}
\NormalTok{trique sur X.}
\NormalTok{   a) Supposons que la forme α soit fortement positive. Alors, si T est positif ( resp.}
\NormalTok{faiblement positif, resp. fortement positif ), il en est de même de T ∧ α.}
\NormalTok{   b) Supposons que la forme α soit positive. Si T est positif, le courant T ∧ α est positif.}
\NormalTok{   c) Supposons que la forme α soit faiblement positive. Si T est fortement positif, le}
\NormalTok{courant T ∧ α est faiblement positif.}
\NormalTok{Démonstration. — Ces assertions découlent des propriétés de positivité d’un produit de}
\NormalTok{deux formes (prop. 3.5.5).}

\NormalTok{Proposition 7.2.8. — Soit X un espace tropical de niveau n, localement différentiel{-}}
\NormalTok{lement séparé.}
\NormalTok{  a) Toute mesure de Radon positive µ sur X induit, par restriction à Ac (X), un courant}
\NormalTok{positif [µ] d’ordre 0.}
\NormalTok{  b) L’application µ 7→ [µ] ainsi définie est une bijection continue de l’espace des me{-}}
\NormalTok{sures de Radon positives sur X sur l’espace des courants positifs de bidimension (0, 0)}
\NormalTok{sur X.}
\NormalTok{Démonstration. — a) Soit µ une mesure de Radon positive sur X ; démontrons que}
\NormalTok{la forme linéaire [µ] est un courant ; ce sera alors un courant positif. Soit K une partie}
\NormalTok{compacte de X et soit D un jeu de cartes couvrant sur X. Il suffit de démontrer que}
\NormalTok{la restriction de [µ] à l’espace AD,K}
\NormalTok{                                  0,0}
\NormalTok{                                      (X) est continue pour sa topologie C 0 . Soit ∥µ∥K}
\NormalTok{162                  CHAPITRE 7. COURANTS SUR UN ESPACE TROPICAL}

\NormalTok{la norme de µ sur K : pour toute fonction continue à support compact sur K, on a}
\NormalTok{|µ(f )| ⩽ ∥µ∥K ∥f ∥K . Cette inégalité vaut a fortiori pour toute fonction de AD,K (X) et}
\NormalTok{entraîne la continuité voulue.}
\NormalTok{   b) Cette implication est injective car l’espace des fonctions lisses à support compact}
\NormalTok{est dense dans l’espace des fonctions continues à support compact (corollaire 6.2.6).}
\NormalTok{   Il reste à prouver qu’elle est surjective. Soit S un courant positif de bidimension (0, 0)}
\NormalTok{sur X. Soit K une partie compacte de X. Soit λ une fonction lisse et positive, à support}
\NormalTok{compact sur X, qui vaut 1 au voisinage de K. Pour toute fonction lisse f sur X dont}
\NormalTok{le support est contenu dans K, on a −∥f ∥K λ ⩽ f ⩽ ∥f ∥K λ. Comme S est positif, il en}
\NormalTok{résulte l’inégalité}
\NormalTok{                                     |S(f )| ⩽ ∥f ∥K S(λ).}
\NormalTok{Par suite, S est continu lorsque l’espace des fonctions lisses à support compact est muni}
\NormalTok{de la topologie de la convergence compacte. Par densité, S s’étend en une forme linéaire}
\NormalTok{positive sur Cc (X), c’est{-}à{-}dire en une mesure de Radon positive.}
\NormalTok{  La démonstration précédente n’a utilisé que la positivité de la forme linéaire S et}
\NormalTok{pas sa continuité. La proposition suivante montre comment ce fait se généralise en tout}
\NormalTok{degré.}
\NormalTok{Proposition 7.2.9. — Soit X un espace tropical de niveau n, localement différentiel{-}}
\NormalTok{lement séparé, soit p un entier et soit S une application linéaire de Acp,p (X) dans R}
\NormalTok{telle que S(Jω) = S(ω) pour tout ω ∈ Acp,p (X). Si S(ω) ⩾ 0 pour toute forme fortement}
\NormalTok{positive ω, alors S est un courant faiblement positif, d’ordre 0.}
\NormalTok{Démonstration. — Il s’agit de prouver la continuité de S lorsque Acp,p (X) est muni de la}
\NormalTok{topologie C 0 . D’après le lemme 7.1.3, il suffit de démontrer l’énoncé suivant : soit f une}
\NormalTok{carte tropicale définie sur une partie ouverte U de X, soit ρ une fonction lisse, positive}
\NormalTok{et à support compact sur U , la forme linéaire ω 7→ S(ρω) sur Afp,p (U ) est continue pour}
\NormalTok{la topologie C 0 .}
\NormalTok{   Soit W un domaine compact de U qui contient le support de ρ. Notons E l’espace}
\NormalTok{affine codomaine de f et munissons{-}le d’une norme euclidienne ; notons η = d′ d′′ ∥x∥2}
\NormalTok{la (1, 1){-}forme fortement positive sur E. Soit P l’espace paralinéaire compact f (W ) ;}
\NormalTok{considérons une norme C 0 sur l’espace A p,p (P ). On déduit de la proposition 1.2.13 qu’il}
\NormalTok{existe un nombre réel c \textgreater{} 0 tel que la forme α + c∥α∥η p est fortement positive, pour}
\NormalTok{tout α ∈ A p,p (P ).}
\NormalTok{   Pour toute ω ∈ Afp,p (U ) ; il existe une unique forme α ∈ A p,p (P ) telle que ω|W = f ∗ α,}
\NormalTok{et l’application ω 7→ ∥ω∥P,0 est une semi{-}norme C 0 sur Af (U ). On la note ∥·∥W,0 .}
\NormalTok{   Soit ω ∈ Afp,p (U ) ; il existe une unique forme α ∈ A p,p (P ) telle que ω|W = f ∗ α, si}
\NormalTok{bien que ω + c∥α∥f ∗ η p est fortement positive. Il en est de même de ρ(ω + c∥α∥P,0 f ∗ η p ).}
\NormalTok{Par l’hypothèse sur S, on a donc}
\NormalTok{                              S(ρω) + c∥ω∥W,0 S(ρf ∗ η p ) ⩾ 0.}
\NormalTok{En appliquant ce raisonnement à −ω, on en déduit l’inégalité}
\NormalTok{                                |S(ρω)| ⩽ cS(ρf ∗ η p )∥ω∥W,0 .}
\NormalTok{Cela démontre la continuité requise.}
\NormalTok{  Ainsi, S est un courant. La condition S(Jω) = S(ω) implique alors que c’est un courant}
\NormalTok{symétrique, et la condition de positivité S(ω) ⩾ 0 pour ω positive entraîne que c’est un}
\NormalTok{courant faiblement positif.}
\NormalTok{                                 § 7.2. COURANTS POSITIFS                                 163}

\NormalTok{Corollaire 7.2.10.  — Soit u une fonction localement intégrable sur X. La forme li{-}}
\NormalTok{néaire ω 7→ X uω sur Acn,n (X) définit un courant d’ordre 0 et de bidimension (n, n)}
\NormalTok{            R}
\NormalTok{sur X, fortement positif si u est positif.}
\NormalTok{   On note [u] ce courant.}
\NormalTok{Démonstration. — Les fonctions u+ = sup(u, 0) et u− = sup(−u, 0) sont localement}
\NormalTok{intégrables. Les formes linéaires qu’elles définissent sont fortement positives, donc sont}
\NormalTok{des courants d’ordre 0. Comme on a u = u+ −u− , la forme linéaire donnée est la différence}
\NormalTok{de ces deux courants, donc est un courant.}

\NormalTok{7.2.11. — Soit X un espace tropical de niveau n, localement différentiellement séparé.}
\NormalTok{Soit p un entier naturel et soit T un courant positif (resp. faiblement positif, resp.}
\NormalTok{fortement positif) de bidimension (p, p).}
\NormalTok{   Soit ω une forme positive (resp. fortement positive, resp. faiblement positive) de bide{-}}
\NormalTok{gré (p, p). Le courant T ∧ω est positif, de bidimension (0, 0) ; d’après la proposition 7.2.8,}
\NormalTok{il est donc associé à une unique mesure µTω sur X. Cette mesure applique une fonction}
\NormalTok{lisse à suppport compact u sur ⟨uT, ω⟩. Lorsque ω est à support compact, cette mesure}
\NormalTok{est à support compact. Lorsque ω est positive, cette mesure est positive.}
\NormalTok{Exemple 7.2.12. — Soit X un espace tropical de niveau n, localement différentiel{-}}
\NormalTok{lement séparé. Soit p un entier naturel et soit T un courant positif (resp. faiblement}
\NormalTok{positif, resp. fortement positif) de bidimension (p, p). Soit u une fonction borélienne et}
\NormalTok{localement bornée sur X. L’application ω 7→ µTω (u) est une forme linéaire sur Acp,p (X)}
\NormalTok{que nous noterons uT .}
\NormalTok{   Lorsque u est positive, cette forme linéaire uT est positive, donc est un courant positif}
\NormalTok{(resp. faiblement positif, resp. fortement positif) sur X.}
\NormalTok{   Lorsque u est une fonction borélienne et localement bornée arbitraire, on l’écrit u =}
\NormalTok{u − u− , où u+ = sup(u, 0) et u− = sup(−u, 0), comme différence de deux fonctions}
\NormalTok{  +}

\NormalTok{positives boréliennes localement bornées sur X. Comme uT = u+ T − u− T , c’est un}
\NormalTok{courant sur X, d’ordre 0.}
\NormalTok{   Cette notation est compatible avec celle déjà introduite lorsque u est une fonction}
\NormalTok{lisse puisque l’on a alors ⟨uT, ω⟩ = ⟨T, uω⟩ = µTω (u).}
\NormalTok{Exemple 7.2.13. — Soit (µn ) une suite de mesures de Radon positives et soit µ une}
\NormalTok{mesure de Radon positive sur X. Si [µn ] → [µ] au sens des courants, alors µn → µ au}
\NormalTok{sens des mesures. Cela résulte de la densité des fonctions lisses à support compact dans}
\NormalTok{les fonctions continues à support compact (corollaire 6.2.6).}
\NormalTok{Exemple 7.2.14. — Il n’est pas vrai que tout courant d’ordre 0 se prolonge en une}
\NormalTok{mesure. Considérons l’étoile XB associée à un ensemble infini B (cf. §6.7.3, dont nous}
\NormalTok{reprenons les notations). Soit (cb )b∈B une famille de nombres réels. Pour toute fonction}
\NormalTok{lisse f sur X, posons}
\NormalTok{                                      X                      }
\NormalTok{                            ⟨T, f ⟩ =     cb fb (0) − fb (−∞) .}
\NormalTok{                                        b∈B}

\NormalTok{Cela définit un courant sur X, de bidimension (0, 0), mais il ne provient d’une mesure}
\NormalTok{que lorsque la famille (cb ) est sommable. Cette mesure est alors égale à}
\NormalTok{                                   X         X}
\NormalTok{                                     cb δη −   cb δp(b,−∞) .}
\NormalTok{                                  b∈B         b∈B}
\NormalTok{164                  CHAPITRE 7. COURANTS SUR UN ESPACE TROPICAL}

\NormalTok{                      § 7.3. COURANTS D’INTÉGRATION}

\NormalTok{7.3.1. — Soit X un espace tropical de niveau    R n. Pour toute (n, n){-}forme lisse α à sup{-}}
\NormalTok{port compact sur X, posons alors ⟨δX , α⟩ = X α, où l’intégrale d’une (n, n){-}forme lisse}
\NormalTok{à support compact sur X a été définie en (5.8.5). On définit ainsi une forme linéaire sur}
\NormalTok{l’espace Acn,n (X). Elle est positive ; d’après la proposition 7.2.9, c’est donc un courant}
\NormalTok{sur X, qu’on appelle le courant d’intégration sur X et on le note δX . Il est de bidimen{-}}
\NormalTok{sion (n, n) et d’ordre 0. Une forme antisymétrique de type (n, n) sur X étant nulle, donc}
\NormalTok{son intégrale aussi, le courant δX est symétrique.}
\NormalTok{7.3.2. — On définit de manière analogue}
\NormalTok{                                   R∂       le courant d’intégration « sur le bord » de X,}
\NormalTok{      ∂                       ∂}
\NormalTok{noté δX , par la formule ⟨α, δX ⟩ = X α, pour toute (n − 1, n){-}forme à support compact α}
\NormalTok{sur X.}
\NormalTok{  Cette   forme linéaire est continue pour la topologie C 0 . Compte tenu de la définition}
\NormalTok{de X α, cela découle en effet de la continuité pour la topologie C 0 de l’intégrale de bord}
\NormalTok{   R}
\NormalTok{d’une forme sur un polytope calibré dans un espace affine.}
\NormalTok{  Il est de bidimension (n − 1, n) et d’ordre 0 ; son support est contenu dans le bord}
\NormalTok{de X.}
\NormalTok{7.3.3. — Soit Y un espace tropical de niveau m et soit φ : Y → X un morphisme}
\NormalTok{topologiquement propre d’espaces tropicaux. Un cas particulier important est celui d’une}
\NormalTok{partie squelettique de X munie de son n{-}calibrage canonique, plus généralement d’une}
\NormalTok{partie squelettique de dimension ⩽ d quelconque munie d’un d{-}calibrage.}
\NormalTok{  Les constructions précédentes fournissent des courants φ∗ δY et φ∗ δY∂ sur X de bidi{-}}
\NormalTok{mensions (m, m) et (m − 1, m) respectivement.}
\NormalTok{  Le courant φ∗ δY est fortement positif.}
\NormalTok{7.3.4. — Soit ω une G{-}forme de type (p, q) à coefficients tropicalement boréliens loca{-}}
\NormalTok{lement bornés sur X. ROn définit un courant [ω] de bidimension (p, q) et d’ordre 0 sur X}
\NormalTok{par la formule α 7→ X ω ∧ α définie par l’intégrale des G{-}formes. Lorsque ω est une}
\NormalTok{forme lisse, on a [ω] = δX ∧ ω.}
\NormalTok{   Si ω est positive (resp. faiblement positive, resp. fortement positive), il en va de même}
\NormalTok{du courant [ω]. Cela découle de la définition d’un courant positif, des propriétés de}
\NormalTok{stabilité de la positivité par produit (prop. 3.5.5), et de ce que l’intégrale d’une (n, n){-}}
\NormalTok{forme positive est positive.                             R∂}
\NormalTok{   On définit aussi le courant [ω]∂ par la formule α 7→ X ω∧α défini à l’aide de l’intégrale}
\NormalTok{de bord des G{-}formes. Il est d’ordre 0 et de bidimension (n − p − 1, n − q). Lorsque ω}
\NormalTok{est une forme lisse, on a [ω]∂ = δX ∂}
\NormalTok{                                      ∧ ω.}
\NormalTok{   Il découle de la définition de la topologie des espaces de courants que les applications}
\NormalTok{ω 7→ [ω] et ω 7→ [ω]∂ sont continues.}
\NormalTok{7.3.5. — Nous allons maintenant généraliser cette construction aux « G{-}formes » qui}
\NormalTok{sont combinaisons linéaires de G{-}formes usuelles à coefficients dans un anneau de fonc{-}}
\NormalTok{tions non nécessairement tropicales.}
\NormalTok{   Soit V un domaine de X et soit U un voisinage ouvert de V dans X dans lequel V}
\NormalTok{est fermé. Soit u ∈ L1loc (V ), ω ∈ AXp,q}
\NormalTok{                                       G}
\NormalTok{                                          (V ) et α ∈ Acn−p,n−q (U ) ; la G{-}forme ω ∧ α|V est à}
\NormalTok{support compact}
\NormalTok{            R     dans V , de sorte que u est intégrable pour la mesure associée. L’appli{-}}
\NormalTok{cation α 7→ V u(ω ∧ α)Rad est un courant d’ordre 0 sur U . La section correspondante de}
\NormalTok{                              § 7.3. COURANTS D’INTÉGRATION                              165}

\NormalTok{π ∗ Dp,q (V ) ne dépend pas du choix de U . On définit ainsi un morphisme de G{-}faisceaux}
\NormalTok{                                L1loc,X ⊗ AXp,q}
\NormalTok{                                             G}
\NormalTok{                                                → π ∗ (Dp,q )X .}
\NormalTok{Par adjonction, on en déduit un morphisme de faisceaux}
\NormalTok{                            π∗ (L1loc,X ⊗ AXp,q}
\NormalTok{                                             G}
\NormalTok{                                                ) → π∗ π ∗ (Dp,q )X .}
\NormalTok{Comme l’unité id → π∗ π ∗ est un isomorphisme de foncteurs, on en déduit un morphisme}
\NormalTok{de faisceaux}
\NormalTok{                               π∗ (L1loc,X ⊗ AXp,q}
\NormalTok{                                                G}
\NormalTok{                                                   ) → (Dp,q )X .}
\NormalTok{   L’image d’une section ω sera notée [ω]. La notation est cohérente avec le cas des G{-}}
\NormalTok{formes lisses ; via l’isomorphisme de la proposition 3.4.11, elle est aussi cohérente avec}
\NormalTok{les cas des G{-}formes à coefficients tropicalement boréliens localement bornés. Ce courant}
\NormalTok{est d’ordre 0 et de bidimension (n − p, n − q).}
\NormalTok{   Par composition avec l’homomorphisme canonique}
\NormalTok{                               1}
\NormalTok{                         π∗ (Lloc,X ⊗ AXp,q}
\NormalTok{                                         G}
\NormalTok{                                            ) → π∗ (L1loc,X ⊗ AXp,q}
\NormalTok{                                                                 G}
\NormalTok{                                                                    ),}
\NormalTok{on en déduit un homomorphisme}
\NormalTok{                                     1}
\NormalTok{                               π∗ (Lloc,X ⊗ AXp,q}
\NormalTok{                                               G}
\NormalTok{                                                  ) → (Dp,q )X}
\NormalTok{qu’on note encore ω 7→ [ω].}
\NormalTok{  Pour toute forme α de type (n − p, n − q), lisse à support compact, on a}
\NormalTok{                            Z             Z                   Z}
\NormalTok{                      −1           −1           −1}
\NormalTok{   ⟨J[ω], α⟩ = ⟨[ω], J α⟩ =   ω∧J α=           J (Jω ∧ α) =     Jω ∧ α = ⟨[Jω], α⟩.}
\NormalTok{                               X                  X                      X}
\NormalTok{Ainsi, J[ω] = [Jω].}
\NormalTok{  On laisse au lecteur le soin d’écrire la variante de cette construction pour les intégrales}
\NormalTok{de bord.}
\NormalTok{Proposition 7.3.6. — Soit ω une G{-}forme (à coefficients tropicalement boréliens et}
\NormalTok{localement bornés) sur X.}
\NormalTok{   a) Pour que le courant [ω] soit nul (respectivement positif, respectivement faiblement}
\NormalTok{positif, respectivement fortement positif ), il faut et il suffit que pour toute partie sque{-}}
\NormalTok{lettique Σ de X, la forme ω|Σ soit nulle (respectivement positive, respectivement. . .) en}
\NormalTok{dehors d’une partie n{-}négligeable.}
\NormalTok{   b) Si ω est à coefficients tropicalement continus, le courant [ω] est nul (respective{-}}
\NormalTok{ment positif, respectivement faiblement positif, respectivement fortement positif ) si et}
\NormalTok{seulement si ω|Σ est nulle (respectivement positive, respectivement. . .) pour toute partie}
\end{Highlighting}
\end{Shaded}

\subsection{CLD-curve, newline lines
9290--9445}\label{cld-curve-newline-lines-92909445}

\begin{Shaded}
\begin{Highlighting}[]
\NormalTok{contienne Wx .}
\NormalTok{   D’après le lemme A.1.5, il existe un recouvrement localement fini (Wj )j∈J de X par}
\NormalTok{des domaines compacts de X, qui raffine le recouvrement ouvert (Ui )i∈I . et tel que les}
\NormalTok{intérieurs Z̊j recouvrent X. Pour tout j ∈ J, on choisit un indice i tel que Wj ⊂ Ui et}
\NormalTok{on note gj la carte G{-}tropicale fi |Wj .}
\NormalTok{   On est ainsi ramené au cas où l’hypothèse (i) est satisfaite, ce qui conclut la démons{-}}
\NormalTok{tration.}



\NormalTok{            § 7.7. CALCUL DIFFÉRENTIEL SUR LES COURANTS}

\NormalTok{  Soit X un espace tropical de niveau n.}
\NormalTok{7.7.1. — L’application linéaire d′ de Acp,q (X) dans Acp+1,q (X) est continue (prop. 3.7.13).}
\NormalTok{On note b′ sa transposée de Dp+1,q (X) dans Dp,q (X). Si T est un courant d’ordre ⩽ r,}
\NormalTok{alors b′ T est un courant d’ordre ⩽ r + 1.}
\NormalTok{   On définit de même b′′ : Dp,q+1 (X) → Dp,q (X) comme la transposée de l’application d′′ .}
\NormalTok{   Il résulte de la définition de la topologie des espaces de courants que ces applications}
\NormalTok{linéaires b′ et b′′ sont continues.}
\NormalTok{   Par dualité, la relation d′′ ω = J d′ J−1 ω entraîne l’égalité}
\NormalTok{(7.7.1.1)                                   b′′ = Jb′ J−1 .}
\NormalTok{                          § 7.7. CALCUL DIFFÉRENTIEL SUR LES COURANTS                  187}

\NormalTok{  On a aussi les relations (b′ )2 = 0, (b′′ )2 = 0, et b′ b′′ = −b′′ b′ .}
\NormalTok{7.7.2. — Soit ω une G{-}forme de type (p, q) sur X et soit [ω] le courant de bidegré (p, q)}
\NormalTok{qui lui est associé. Pour toute forme α ∈ Acn−p−1,n−q (X), on a}

\NormalTok{                              d′ (ω ∧ α) = d′ ω ∧ α + (−1)p+q ω ∧ d′ α.}

\NormalTok{La formule de Stokes (proposition 5.11.1) entraîne alors l’égalité}
\NormalTok{            Z ∂         Z               Z                        Z}
\NormalTok{                             ′               ′}
\NormalTok{                ω∧α=       d (ω ∧ α) =      d ω ∧ α + (−1)   p+q}
\NormalTok{                                                                   ω ∧ d′ α,}
\NormalTok{                 X               X                   X                      X}

\NormalTok{autrement dit}
\NormalTok{                                     [d′ ω] = [ω]∂ + (−1)p+q+1 b′ [ω].}
\NormalTok{Si l’on définit l’opérateur d′ de D p,q (X) dans D p+1,q par}

\NormalTok{(7.7.2.1)                                   d′ = (−1)p+q+1 b′ ,}

\NormalTok{il vient alors}

\NormalTok{(7.7.2.2)                                 d′ [ω] = [d′ ω] − [ω]∂ .}

\NormalTok{Lorsque X est sans bord et ω est une forme sur X, le courant [ω]∂ est nul, et l’on a}

\NormalTok{(7.7.2.3)                                     d′ [ω] = [d′ ω].}
\NormalTok{7.7.3. — De même, on pose}

\NormalTok{(7.7.3.1)                                  d′′ = (−1)p+q+1 b′′ .}

\NormalTok{On en déduit que les relations du lemme 1.3.6 sont encore valides sur les espaces de}
\NormalTok{courants :}

\NormalTok{(7.7.3.2)                              d′′ = J d′ J−1 = −J−1 d′ J}
\NormalTok{(7.7.3.3)                                   (d′ )2 = (d′′ )2 = 0}
\NormalTok{(7.7.3.4)                                    d′′ d′ = − d′ d′′}
\NormalTok{(7.7.3.5)                                    d′ d′′ J = J d′ d′′ .}

\NormalTok{  Soit ω une G{-}forme de type (p, q). On a}

\NormalTok{                     d′′ [ω] = J d′ J−1 [ω] = J d′ [J−1 ω] = J[d′ J−1 ω] − J[J−1 ω]∂}
\NormalTok{                           = [J d′ J−1 ω] − J[J−1 ω]∂ ,}

\NormalTok{de sorte que}

\NormalTok{(7.7.3.6)                              d′′ [ω] = [d′′ ω] − J[J−1 ω]∂ .}

\NormalTok{En outre, si X est sans bord et ω est une forme, alors Jω est également une forme, donc}
\NormalTok{on a}

\NormalTok{(7.7.3.7)                                    d′′ [ω] = [d′′ ω].}
\NormalTok{188                  CHAPITRE 7. COURANTS SUR UN ESPACE TROPICAL}

\NormalTok{7.7.4. — Supposons que l’espace X est sans bord. Comme le courant d’intégration δX}
\NormalTok{est égal à [1], il est d′ et d′′ {-}fermé.}

\NormalTok{Lemme 7.7.5. — Pour tout courant S ∈ Dp,q (X) et toute forme α ∈ Acp,q (X), on a}

\NormalTok{                                   ⟨d′ d′′ S, α⟩ = ⟨S, d′ d′′ α⟩.}

\NormalTok{Démonstration. — En effet, d′ d′′ S = ((−1)p+q+2 b′ )((−1)p+q+1 b′′ )S = − b′ b′′ S est}
\NormalTok{l’image de S par le transposé de − d′′ d′ = d′ d′′ .}


\NormalTok{Exemple 7.7.6. — Soit X un espace espace tropical de niveau n. Soit α une G{-}forme}
\NormalTok{de type (p, q) à coefficients tropicalement de classe C 2 . On a donc}

\NormalTok{                d′ d′′ [α] = d′ [d′′ α] − J[J−1 α]∂ = [d′ d′′ α] − [d′′ α]∂ − d′ J[J−1 α]∂ .}
                                                   
\NormalTok{(7.7.6.1)}

\NormalTok{  Supposons de plus que α soit symétrique. Alors, p = q et J−1 α = α, de sorte que}

\NormalTok{(7.7.6.2)                  d′ d′′ [α] = [d′ d′′ α] − [d′′ α]∂ − d′ J[α]∂ .}

\NormalTok{  Supposons de plus que X soit sans bord et p = 0, c’est{-}à{-}dire que α soit une fonction,}
\NormalTok{G{-}localement tropicalisée par une fonction de classe C 2 , et démontrons que l’on a}

\NormalTok{(7.7.6.3)                        d′ d′′ [α] = [d′ d′′ α] − [d′′ α]∂ .}

\NormalTok{Compte tenu de la relation précédente, il suffit de vérifier que le courant [α]∂ est nul.}
\NormalTok{Or, pour toute forme β de type (n − 1, n), lisse et à support compact, le nombre réel}
\NormalTok{                  R∂}
\NormalTok{⟨[α]∂ , β⟩, égal à X αβ, est l’intégrale de α contre la mesure de bord associée à β. Comme}
\NormalTok{β est une forme et comme X est sans bord, cette mesure de bord est nulle, d’où l’égalité}
\NormalTok{[α]∂ = 0.}
\NormalTok{   On retrouve en particulier la formule classique définissant le laplacien combinatoire}
\NormalTok{d’un graphe ([120], [141]).}

\NormalTok{Lemme 7.7.7. — Soit T un courant de bidimension (p, q) sur X et soit α une forme}
\NormalTok{de type (p′ , q ′ ). On a l’égalité :}

\NormalTok{                         d′ (T ∧ α) = (d′ T ) ∧ α + (−1)p+q T ∧ d′ α.}

\NormalTok{Démonstration. — En effet, pour toute forme β de type (p − p′ , q − q ′ ) sur X, on a}
\NormalTok{                                       ′     ′}
\NormalTok{            ⟨d′ (T ∧ α), β⟩ = (−1)p−p +q−q +1 ⟨T ∧ α, d′ β⟩}
\NormalTok{                                       ′     ′}
\NormalTok{                           = (−1)p−p +q−q +1 ⟨T, α ∧ d′ β⟩}
\NormalTok{                                       ′     ′               ′   ′}
\NormalTok{                           = (−1)p−p +q−q +1 ⟨T, (−1)p +q (d′ (α ∧ β) − d′ α ∧ β)⟩}
\NormalTok{                           = (−1)p+q+1 ⟨T, d′ (α ∧ β)⟩ + (−1)p+q ⟨T, d′ α ∧ β⟩}
\NormalTok{                           = ⟨d′ T ∧ α, β⟩ + (−1)p+q ⟨T ∧ d′ α, β⟩,}

\NormalTok{d’où l’égalité voulue.}
\NormalTok{                      § 7.7. CALCUL DIFFÉRENTIEL SUR LES COURANTS                                 189}

\NormalTok{7.7.8. — Soit U un ouvert de X. Pour tout courant T sur X, on a d′ (T |U ) = (d′ T |U )}
\NormalTok{et d′′ (T |U ) = (d′′ T )|U .}
\NormalTok{   Soit Y un espace tropical de niveau m et soit φ : Y → X un morphisme compact}
\NormalTok{d’espaces tropicaux. Pour tout courant T sur Y , on a d′ (φ∗ T ) = φ∗ (d′ T ) et d′′ (φ∗ T ) =}
\NormalTok{φ∗ (d′′ T ).}
\NormalTok{Proposition 7.7.9. — Soit Y un espace tropical de niveau m et soit φ : Y → X un}
\NormalTok{morphisme compact d’espaces tropicaux. On a d′ (φ∗ δY ) = φ∗ (δY∂ ). En particulier, si X}
\NormalTok{est sans bord, le courant δX est d′ {-}fermé.}
\NormalTok{Démonstration. — Soit d = dim(Y ) et soit α une (d−1, d){-}forme lisse à support compact}
\NormalTok{sur X. On a}
\NormalTok{                                                       Z}
\NormalTok{                   ⟨d (φ∗ δY ), α⟩ = −⟨φ∗ δY , d α⟩ = − φ∗ d′ α}
\NormalTok{                     ′                          ′}
\NormalTok{                                                         Y}
\NormalTok{                                      Z               Z ∂}
\NormalTok{                                   = − d′ φ∗ α = −        φ∗ α}
\NormalTok{                                                Y                     Y}
\NormalTok{                                          = −⟨φ∗ δY∂ , α⟩,}
\NormalTok{d’où la formule annoncée.}

\NormalTok{Exemple 7.7.10. — Soit T un courant squelettique de dimension p sur X. Soit f une}
\NormalTok{carte G{-}tropicale sur un domaine V de X et soit PfT la partie squelettique p{-}calibrée}
\NormalTok{pure qui décrit T relativement à f . Soit q un entier tel que 0 ⩽ q ⩽ p − 1. Démontrons}
\NormalTok{que pour toute G{-}forme ω de type (q, p) sur V qui est f {-}tropicale, et toute G{-}forme}
\NormalTok{G{-}lisse ρ de type (p − q − 1, 0) sur V telle que ρ ∧ ω soit à support compact, on a}
\NormalTok{                                                 Z ∂}
\end{Highlighting}
\end{Shaded}

\subsection{CLD-curve, newline lines
6860--7020}\label{cld-curve-newline-lines-68607020}

\begin{Shaded}
\begin{Highlighting}[]
\NormalTok{sur le site topologique de X. Le lemme est ainsi démontré.}


\NormalTok{Proposition 5.9.3. — Notons π : XG → X le morphisme canonique de sites. Il existe}
\NormalTok{un unique morphisme de faisceaux π∗ FXn−1,n}
\NormalTok{                                          G}
\NormalTok{                                              → π∗ VXn−1}
\NormalTok{                                                       G}
\NormalTok{                                                          sur l’espace topologique X, noté}
\NormalTok{ω 7→ ω}
\NormalTok{     e , qui prolonge l’application définie en (5.9.1) lorsque ω possède un support fort.}
\NormalTok{Ce morphisme est FX0,0 {-}linéaire.}

\NormalTok{Démonstration. — Compte tenu du lemme 5.9.2, cela résulte de ce que toute G{-}forme}
\NormalTok{sur un ouvert paracompact possède un support fort, et que tout espace G{-}tropical est}
\NormalTok{recouvert par ses ouverts paracompacts.}

\NormalTok{5.9.4. — Explicitons cette construction lorsque ω est tropicale. Soit f : X → E une}
\NormalTok{carte G{-}tropicale et soit α une (n − 1, n){-}forme sur un sous{-}espace paralinéaire de E}
\NormalTok{                                             (n)}
\NormalTok{contenant f (X) telle que ω = f ∗ α. Alors, Σf est un support fort de ω (lemme 5.3.4)}
\NormalTok{                                                                             (n)}
\NormalTok{et on a l’égalité de formes{-}volume (n − 1){-}dimensionnelles sur Σf                  :}

\NormalTok{(5.9.4.1)                                 ω̃ = ⟨f ∗ α, µΣ(n) ⟩.}
\NormalTok{                                                            f}


\NormalTok{Si, de plus, l’application f |Σ(n) est topologiquement propre, la proposition 4.5.8 entraîne}
\NormalTok{                               f}
\NormalTok{alors que}

\NormalTok{(5.9.4.2)                       ftrop,∗ (ω̃) = ⟨α, f∗ µΣ(n) ⟩ = ⟨α, µΠf ⟩.}
\NormalTok{                                                        f}


\NormalTok{Rappelons alors (cf. §4.6.13) que la forme{-}volume (n − 1){-}dimensionnelle ⟨α, µΠf ⟩ se}
\NormalTok{calcule par accouplement avec la discordance du calibrage de Πf .}

\NormalTok{Proposition 5.9.5 (Compensation des bords). — Soit (Xi )i∈I un G{-}recouvrement}
\NormalTok{localement fini de X et soit ω ∈ FXn−1,n}
\NormalTok{                                     GT}
\NormalTok{                                         (X) une G{-}forme de type (n − 1, n) sur X.}
\NormalTok{Pour toute partie J de I, notons XJ = j∈J Xj et uJ l’immersion de XJ dans X. On a}
\NormalTok{l’égalité}
\NormalTok{                                X}
\NormalTok{                           ω̃ =     (−1)Card(J)−1 (uJ )∗ ω|}
\NormalTok{                                                         g XJ .}
\NormalTok{                                      ∅̸=J⊂I}


\NormalTok{Démonstration. — En raisonnant localement, on suppose que X est paracompact et}
\NormalTok{donc que ω possède un support fort. La formule découle alors de la proposition 4.6.6.}
\NormalTok{                         § 5.10. INTÉGRALES ET TROPICALISATIONS                           139}

\NormalTok{5.9.6. — Soit ω une G{-}forme de type (n − 1, n) sur X à coefficients tropicalement}
\NormalTok{boréliens. Nous venons de lui associer une forme{-}volume (n − 1){-}dimensionnelle ω̃. Nous}
\NormalTok{noterons alors ωBor la mesure de Borel sur X associée par la construction du §5.7.1.}
\NormalTok{   Lorsque la forme{-}volume ω̃ est localement intégrable, la mesure ωBor est localement}
\NormalTok{finie ; on note alors ω|Rad la mesure de Radon associée à ω̃. On a |ωRad | = ωBor .}
\NormalTok{                                                     R∂         R∂}
\NormalTok{   Si f est une fonction borélienne sur X, on notera X f |ω| et X f ω les intégrales de f}
\NormalTok{contreR les mesures ωBor et ωRad lorsqu’elles existent. Cela définit en particulier l’inté{-}}
\NormalTok{         ∂}
\NormalTok{grale X ω pour toute G{-}forme de type (n − 1, n) à support compact qui est localement}
\NormalTok{intégrable (par exemple lisse).}



\NormalTok{               § 5.10. INTÉGRALES ET TROPICALISATIONS}

\NormalTok{  Soit X un espace G{-}tropical calibré de niveau n.}
\NormalTok{Proposition 5.10.1. — Soit f : X → E une carte G{-}tropicale agréable, soit Πf le}
\NormalTok{polyèdre tropical de la carte f et soit µΠf son calibrage.}
\NormalTok{  Soit p ∈ \{n−1, n\} et soit α une (p, n){-}forme à coefficients boréliens sur un sous{-}espace}
\NormalTok{paralinéaire P de E contenant f (X).}
\NormalTok{  a) Si p = n, on a l’égalité}
\NormalTok{                                   ftrop,∗ (f ∗ α)Bor ⩾ ⟨α, µΠf ⟩Bor}
\NormalTok{                             (n)}
\NormalTok{de mesures de Borel sur Πf .}
\NormalTok{  b) Si p = n − 1, on a l’inégalité}
\NormalTok{                                   ftrop,∗ (f ∗ α)Bor ⩾ ⟨α, µΠf ⟩Bor}
\NormalTok{                             (n)}
\NormalTok{de mesures de Borel sur Πf .}
\NormalTok{  c) La mesure (f ∗ α)Bor est localement finie si et seulement si la forme α|Πf est locale{-}}
\NormalTok{ment intégrable sur Πf ; dans ce cas, on a l’égalité}
\NormalTok{                                   ftrop,∗ (f ∗ α)Rad = ⟨α, µΠf ⟩Rad}
\NormalTok{de mesures de Radon sur Πf .}
\NormalTok{   d) En particulier, si α est intégrable sur Πf , la forme f ∗ α est intégrable sur X, on a}
\NormalTok{l’égalité}
\NormalTok{                                   Z         Z}
\NormalTok{                                        ∗}
\NormalTok{(5.10.1.1)                            f α=       ⟨α, µΠf ⟩}
\NormalTok{                                        X           Πf}

\NormalTok{lorsque p = n et l’égalité}
\NormalTok{                                      Z ∂          Z ∂}
\NormalTok{                                            ∗}
\NormalTok{(5.10.1.2)                                  f α=         ⟨α, µΠf ⟩}
\NormalTok{                                        X           Πf}

\NormalTok{lorsque p = n − 1.}
\NormalTok{Démonstration. — Posons ω = f ∗ α. D’après la relation (5.8.4.2) lorsque p = n et la}
\NormalTok{relation (5.9.4.2) lorsque p = n−1, on a l’égalité (ftrop )∗ (ω̃) = ⟨α, µΠf ⟩ de formes{-}volume}
\NormalTok{p{-}dimensionnelles sur Πf . Le lemme 4.4.6 entraîne alors que f∗ (ω̃Bor ) ⩾ ⟨α, µΠf ⟩Bor .}
\NormalTok{Loreque p = n, on vérifie, en revenant aux définitions et en utilisant la positivité du}
\NormalTok{140                   CHAPITRE 5. ESPACES G{-}TROPICAUX CALIBRÉS}

\NormalTok{calibrage canonique, que l’on a en fait l’égalité, d’où (1). En appliquant le lemme 4.4.6}
\NormalTok{aux mesures de Radon correspondantes, on en déduit (2), puis (3) par intégration.}

\NormalTok{5.10.2. — La formule (5.10.1.1) définit explicitement l’intégrale d’une forme tropicale.}
\NormalTok{Prouvons inversement que cette relation dans des cartes tropicales permet de reconstituer}
\NormalTok{la mesure de Radon associée à une G{-}forme de type (n, n) et en donne ainsi une définition}
\NormalTok{tropicale.}
\NormalTok{   Soit ω une G{-}forme de type (n, n) sur X ; supposons{-}la à support compact et à coef{-}}
\NormalTok{ficients tropicalement boréliens localement bornés.}
\NormalTok{   Soit (Vi )i∈I une famille finie de domaines compacts de X tropicalisant ω et dont dont}
\NormalTok{la réunion}
\NormalTok{      T       V recouvre le support de ω. Pour toute partie non vide J de I, posons}
\NormalTok{VJ = j∈J Vj . On a alors}
\NormalTok{                          Z        Z       X                  Z}
\NormalTok{                                                    Card(J)−1}
\NormalTok{                              ω=       ω=       (−1)            ω.}
\NormalTok{                          X       V       ∅̸=J⊂I              VJ}

\NormalTok{Pour toute partie non vide J, la forme ω est tropicaleR sur VJ , et VJ est compact, si bien}
\NormalTok{que toute carte G{-}tropicale sur VJ . Ainsi, l’intégrale VJ ω peut se calculer à l’aide de la}
\NormalTok{formule (5.10.1.1).}



\NormalTok{                 § 5.11. FORMULES DE STOKES ET GREEN}

\NormalTok{  Soit X un espace G{-}tropical calibré de niveau n.}
\NormalTok{Proposition 5.11.1. — a) Soit ω une G{-}forme à coefficients C 1 et à support com{-}}
\NormalTok{pact sur X, de type (n − 1, n). On a la formule de Stokes :}
\NormalTok{                                   Z         Z ∂}
\NormalTok{                                        ′}
\NormalTok{                                       d ω=      ω.}
\NormalTok{                                      X                X}
\NormalTok{   b) Soit α et β des G{-}formes à coefficients C sur X, symétriques, de type (p, p) et (q, q)}
\NormalTok{                                                   2}

\NormalTok{respectivement, où p + q = n − 1. Si l’intersection de leurs supports est compact, on a la}
\NormalTok{formule de Green :}
\NormalTok{                 Z                               Z ∂}
\NormalTok{                          ′ ′′     ′ ′′}
\NormalTok{                    (α ∧ d d β − d d α ∧ β) =        (α ∧ d′′ β − d′′ α ∧ β).}
\NormalTok{                  X                                    X}

\NormalTok{Démonstration. — Comme le support de ω est compact, il existe une partie squelettique}
\NormalTok{fermée Σ de X qui la supporte fortement. La formule de Stokes découle alors de la}
\NormalTok{formule correspondante en géométrie paralinéaire (proposition 4.6.15), appliquée à Σ(n) ,}
\NormalTok{à la forme ω|Σ(n) et au calibrage canonique de Σ(n) .}
\NormalTok{   La formule de Green s’en déduit comme dans le lemme 1.5.10 en appliquant la formule}
\NormalTok{de Stokes à la forme à support compact ω = α ∧ d′′ β − d′′ α ∧ β. Lorsque α et β sont}
\NormalTok{toutes deux à support compact, elle découle aussi de la formule de Green en géométrie}
\NormalTok{paralinéaire, appliqué à une partie polyédrale qui supporte fortement ces deux formes.}
\NormalTok{CHAPITRE 6}
\end{Highlighting}
\end{Shaded}

\subsection{GR, newline lines 70--138}\label{gr-newline-lines-70138}

\begin{Shaded}
\begin{Highlighting}[]
\NormalTok{Jonsson [BJ22] give a far reaching global pluripotential theory on a projective variety}
\NormalTok{over a trivially valued field. Some parts work for any non{-}Archimedean field, at least in}
\NormalTok{residue characteristic zero. These psh functions might be singular, and the definition uses}
\NormalTok{decreasing limits of semipositive model functions as in Zhang’s approach. Chambert–Loir}
\NormalTok{and Ducros [CD25] have posted a second version of their work in which they give a new}
\NormalTok{definition of psh functions which is functorial and is based on a transfinite regularization}
\NormalTok{process.}
\NormalTok{   In all of the above approaches, a maximum principle for plurisubharmonic functions}
\NormalTok{is missing. In this paper, we propose another class of psh functions called classically}
\NormalTok{plurisubharmonic functions in close analogy to the complex case. The definition is local,}
\NormalTok{the functions are upper semicontinuous, and we will show that properties (1)–(4) are}
\NormalTok{satisfied. It will be clear that classically psh functions form a maximal class of reasonable}
\NormalTok{plurisubharmonic functions that is compatible with Thuillier’s definition for curves. In}
\NormalTok{particular, the psh functions from the above are classically psh (see 1.2 below).}
\NormalTok{   To get regularization by smooth psh functions in the sense of Chambert{-}Loir and Ducros}
\NormalTok{or to define a Monge–Ampère measure for psh functions, additional hypotheses as in}
\NormalTok{[CD25] might be needed, but that is not the subject of this paper.}

\NormalTok{1.1. Subharmonic functions on arbitrary k{-}analytic curves. In Section 3, we define}
\NormalTok{subharmonic functions on any k{-}analytic curve X to be the elements of the following}
\NormalTok{subsheaf of the sheaf of upper semicontinuous functions on X with range [−∞, ∞).}
\NormalTok{It is the smallest subsheaf that is functorial with respect to pullback, stable under base}
\NormalTok{change, and agrees with Thuillier’s subharmonic functions on smooth k{-}analytic curves.}
\NormalTok{We will show in Proposition 3.5 that subharmonic functions satisfy the properties (1)–(3)}
\NormalTok{as in the complex case, and the local maximum principle on the interior Int(X ) = X \textbackslash{} ∂ X .}

\NormalTok{1.2. Plurisubharmonic functions on k{-}analytic spaces. We define classically plurisub{-}}
\NormalTok{harmonic functions on k{-}analytic spaces as the elements of the smallest subsheaf of the}
\NormalTok{sheaf of upper semicontinuous functions with range [−∞, ∞) that is functorial with}
\NormalTok{             CLASSICALLY PSH AND PLURIHARMONIC FUNCTIONS ON BERKOVICH SPACES              3}

\NormalTok{respect to pullback, stable under base change, and agrees with the above defined sheaf}
\NormalTok{of subharmonic functions on k{-}analytic curves. In Section 4, we will show that classically}
\NormalTok{psh functions satisfy properties (1)–(3).}
\NormalTok{   We use the terminology classically psh to distinguish from the psh functions introduced}
\NormalTok{by Chambert{-}Loir and Ducros in [CD25]. Note that psh functions are classically psh.}
\NormalTok{Indeed, psh functions are upper semi{-}continuous and the definition is functorial, so we}
\NormalTok{are reduced to the case of a smooth analytic curve C. We have to see that a psh function}
\NormalTok{is subharmonic. By base change, we may assume that k is algebraically closed and non{-}}
\NormalTok{trivially valued, see Proposition 3.5(7). Since subharmonic functions are stable under}
\NormalTok{decreasing limits, and the psh functions are defined by a transfinite regularization process}
\NormalTok{from smooth psh functions, it is enough to show that a smooth psh function on C is}
\NormalTok{subharmonic. This follows from a result of Wanner [Wan19, Theorem 4.7].}
\NormalTok{1.3. Plurisubharmonic R{-}PL functions. For an R{-}piecewise linear function h: X → R}
\NormalTok{on a good strictly k{-}analytic space over a non{-}trivially valued field k, we have studied the}
\NormalTok{notion of semipositivity in the paper [GR25]. This notion is a local analytic version of}
\NormalTok{the global semipositive metrics introduced by Zhang [Zha95] which are induced by nef}
\NormalTok{models. We show in Theorem 4.3 that an R{-}PL function h is semipositive if and only if h}
\NormalTok{is classically psh.}
\NormalTok{1.4. Connectivity by curves. In Theorem 5.4, we show that two points x, y of a con{-}}
\NormalTok{nected k{-}analytic space X can be connected by analytic curves (assumed to be compact}
\NormalTok{and defined over a non{-}Archimedean field extension of k). If ∂ X = ;, then we even show}
\NormalTok{that x, y can be connected by overconvergent curves, which means roughly speaking that}
\NormalTok{we can connect them with the interiors of compact curves. We refer to Section 5 for the}
\NormalTok{precise definitions. The proof of Theorem 5.4 is inspired by Berkovich’s proof [Ber90,}
\NormalTok{§3.2] showing that a connected k{-}analytic space is pathwise connected. However, we}
\NormalTok{have to choose different paths, as the ones in ibid are not analytic curves. This result is}
\NormalTok{related to [dJ95, Proposition 6.1.1] and [Ber07, Theorem 4.1.1] (see the discussion after}
\NormalTok{the statement of Theorem 5.4 for more details).}
\NormalTok{Theorem 1.5 (Local Maximum Principle). Let X be a k{-}analytic space, and let u be a}
\NormalTok{classically psh function on X . If u has a local maximum on x ∈ Int(X ), then u is locally}
\NormalTok{constant at x.}
\NormalTok{  This will be shown in Theorem 6.1. For the proof, one may assume that ∂ X = ;;}
\NormalTok{then we reduce the claim to the local maximum principle for k{-}analytic curves using}
\NormalTok{connectivity by overconvergent curves.}
\NormalTok{1.6. Global maximum principle for affinoids. Berkovich has shown [Ber90, Propo{-}}
\NormalTok{sition 2.4.4] that a k{-}affinoid space X = M (A ) has a Shilov boundary, i.e. there is a}
\NormalTok{smallest subset Γ of X such that}
\NormalTok{                                 max | f (x)| = max | f (x)|}
\end{Highlighting}
\end{Shaded}

\subsection{GR, newline lines 250--405}\label{gr-newline-lines-250405}

\begin{Shaded}
\begin{Highlighting}[]
\NormalTok{finite extension of k.}
\NormalTok{   If ϕ : Y → X is a morphism, then the fiber ϕ −1 (x) over x ∈ X is an H (x){-}analytic}
\NormalTok{space, and the underlying topological space has the induced topology [Ber93, §1.4].}

\NormalTok{2.8. Fiber products. The category of k{-}analytic spaces admits fiber products by [Ber93,}
\NormalTok{§1.4]. For analytic spaces X , Y over a k{-}analytic space S, we denote the fiber product by}
\NormalTok{X ×S Y . Given x ∈ X and y ∈ Y over the same point s ∈ S, there always exists z ∈ X ×S Y}
\NormalTok{mapping to x and to y with respect to the canonical projections [Duc18, 1.2.13]. For}
\NormalTok{convenience, we repeat the argument. The completed tensor product H (x)⊗         b H (s) H ( y)}
\NormalTok{contains H (x) ⊗H (s) H ( y) by a result of Gruson [Gru66, §3.2 Théorème 1(4)], so}
\NormalTok{H (x)⊗b H (s) H ( y) ̸= 0. The Berkovich spectrum of any nonzero Banach ring is nonempty}
                                                                 
\NormalTok{[Ber90, Theorem 1.2.1], and any z ′ ∈ M H (x)⊗     b H (s) H ( y) induces a point z ∈ X ×S Y}
\NormalTok{mapping to x and y. Note that H (z) is a simultaneous analytic extension field of H (x)}
\NormalTok{and of H ( y).}
\NormalTok{  Since Gruson’s result holds for Banach algebras over a non{-}Archimedean field, the}
\NormalTok{same argument shows that for two analytic field extensions of k, there is a joint analytic}
\NormalTok{extension of both.}

\NormalTok{2.9. Extension of scalars. Let k′ /k be an an analytic extension field. Then there is a}
\NormalTok{base extension functor (·)k′ which associates to a k{-}analytic space X a k′ {-}analytic space}
\NormalTok{X k′ , given on a k{-}affinoid space X = M (A ) by the k′ {-}affinoid space X k′ = M (A ⊗b k k′ )}
\NormalTok{[Ber90, §1.4]. We denote by πk′ /k : X k′ → X the structure map, which is closed and}
\NormalTok{surjective by [GRW17, Lemma 2.19], [Duc24, (3.1.1.2)].}

\NormalTok{2.10. The Shilov section. For r = (r1 , . . . , rn ) ∈ R\textgreater{}0     n}
\NormalTok{                                                                   , we consider the k{-}affinoid do{-}}
\NormalTok{main in the closed polydisc M k⟨r1 T1 , . . . , rn Tn ⟩ given by the equations |T j (x)| = r j ,}
\NormalTok{                                          −1            −1}

\NormalTok{j = 1, . . . , n. The corresponding      k{-}affinoid algebra is denoted by k r ; it consists of the}
\NormalTok{Laurent series f = λ∈Zn aλ T λ such that aλ r λ → 0 for |λ| → ∞. The k{-}affinoid al{-}}
\NormalTok{                           P}

\NormalTok{gebra M (k r ) has a unique Shilov        boundary point equal to the weighted Gauss point}
\NormalTok{η r ∈ M k⟨r1−1 T1 , . . . , rn−1 Tn ⟩ given by the weighted Gauss norm ∥ ∥ r from 2.3. If}
\NormalTok{                                                                                       p}
\NormalTok{r1 , . . . , rn induce linearly independent elements in the Q{-}vector space R\textgreater{}0 / |k× |, then k r}
\NormalTok{is a non{-}Archimedean field; it follows that k r = H (η r ). For every k{-}affinoid algebra A ,}
\NormalTok{there is always such an r such that A ⊗        b k r is a strictly k r {-}affinoid algebra. See [Ber90,}
\NormalTok{pp. 21–22].}
\NormalTok{             CLASSICALLY PSH AND PLURIHARMONIC FUNCTIONS ON BERKOVICH SPACES                 7}

\NormalTok{   For any k{-}analytic space X and r ∈ R\textgreater{}0n}
\NormalTok{                                             , the structure map πkr /k : X kr → X has a canon{-}}
\NormalTok{ical section σ : X → X r mapping x ∈ X to the unique Shilov boundary point of the fiber}
\NormalTok{π−1}
\NormalTok{  k r /k}
\NormalTok{         (x) = M (H (x)⊗b k r ) = M (H (x) r ) (which is the weighted Gauss point η r ). We}
\NormalTok{call σ the Shilov section. The Shilov section is continuous: see [Ber90, Lemma 3.2.2(i)],}
\NormalTok{[Duc18, 1.2.16].}

\NormalTok{2.11. Relative boundary. An important notion in this paper is that of the relative bound{-}}
\NormalTok{ary of a morphism ϕ : X → Y of k{-}analytic spaces. This is a closed subset of X which we}
\NormalTok{denote by ∂ (X /Y ). We will not give the details of this definition, but let us try to explain}
\NormalTok{the literature on the subject. In [Ber90, §2.5], Berkovich defines the relative boundary}
\NormalTok{∂ (X /Y ) for k{-}affinoid spaces X and Y . Then in [Ber90, §3.1], he generalizes ∂ (X /Y ) for}
\NormalTok{good k{-}analytic spaces using that all points of X and Y have k{-}affinoid neighborhoods.}
\NormalTok{Berkovich calls ϕ closed if ∂ (X /Y ) = ;. Berkovich defines ∂ (X /Y ) for arbitrary k{-}analytic}
\NormalTok{spaces in [Ber93, 1.5.3(ii)], as follows. A morphism ϕ is closed if for every morphism}
\NormalTok{Y ′ → Y from a good k{-}analytic space Y ′ , the fiber product X ′ := X ×Y Y ′ is good and the}
\NormalTok{canonical morphism X ′ → Y ′ is closed in the previous sense. Finally in [Ber93, Defini{-}}
\NormalTok{tion 1.5.4], a point x ∈ X is defined to be in Int(X /Y ) if x has an open neighborhood U}
\NormalTok{such that ϕ induces a closed morphism U → Y in the above sense. We refer to [Tem15,}
\NormalTok{Definition 4.2.4.1] for a more explicit equivalent definition of ∂ (X /Y ). Note that in this}
\NormalTok{paper, we will use the notion closed map in the usual topological sense meaning that it}
\NormalTok{maps closed sets to closed sets. This is different from Berkovich’s terminology.}
\NormalTok{   The complement Int(X /Y ) := X \textbackslash{} ∂ (X /Y ) is called the relative interior. We will usually}
\NormalTok{apply this in the case Y = M (k), in which case we call ∂ X := ∂ (X /M (k)) the boundary}
\NormalTok{of X and Int(X ) := X \textbackslash{} ∂ X the interior. If ∂ X = ;, we say that X boundaryless. A bound{-}}
\NormalTok{aryless space is good by definition. Note that Berkovich calls boundaryless k{-}analytic}
\NormalTok{spaces closed.}
\NormalTok{   We list here some properties of the boundary and interior which we will use throughout}
\NormalTok{the paper:}
\NormalTok{   (1) A point x ∈ X is contained in Int(X /Y ) if and only if there is an open neighbor{-}}
\NormalTok{       hood U of x such that Int(U/Y ) = U: this is part of the above definition.}
\NormalTok{   (2) If ϕ : X → Y is finite, then Int(X /Y ) = X [Ber93, Example 1.5.3(iii)].}
\NormalTok{   (3) If X is an analytic domain in Y , then Int(X /Y ) is the topological interior of X in}
\NormalTok{       Y [Ber93, Proposition 1.5.5(i)].}
\NormalTok{             ϕ     ψ}
\NormalTok{   (4) If X → Y → S are morphisms of k{-}analytic spaces, then}

\NormalTok{(2.11.1)                   Int(X /Y ) ∩ ϕ −1 Int(Y /S) ⊂ Int(X /S)}

\NormalTok{        [Ber93, Proposition 1.5.5(ii)], with equality if ψ is locally separated [Tem04,}
\NormalTok{        Corollary 5.7]. Taking S = M (k), this becomes}

\NormalTok{(2.11.2)                      Int(X /Y ) ∩ ϕ −1 Int(Y ) ⊂ Int(X ).}

\NormalTok{   (5) If k′ /k is an analytic extension field, then}

\NormalTok{                                π−1}
\NormalTok{                                 k′ /k}
\NormalTok{                                       Int(X /Y ) ⊂ Int(X k′ /Yk′ )}
\NormalTok{8                                      W. GUBLER AND J. RABINOFF}

\NormalTok{         [Ber93, Proposition 1.5.5(iii)]. Taking Y = M (k), this becomes}
\NormalTok{                                       π−1}
\NormalTok{                                        k′ /k}
\NormalTok{                                              Int(X ) ⊂ Int(X k′ ).}
\NormalTok{For convenience, we restate these properties in terms of the relative boundary:}
\NormalTok{   (1′ ) A point x ∈ X is contained in ∂ (X /Y ) if and only if ∂ (U/Y ) ̸= ; for every open}
\NormalTok{         neighborhood U of x.}
\NormalTok{   (2′ ) If ϕ is finite, then ∂ (X /Y ) = ;.}
\NormalTok{   (3′ ) If X is an analytic domain in Y , then ∂ (X /Y ) is the topological boundary of X}
\NormalTok{         in Y .}
\NormalTok{               ϕ     ψ}
\NormalTok{     ′}
\NormalTok{   (4 ) If X → Y → S are morphisms of k{-}analytic spaces, then}
\NormalTok{(2.11.3)                        ∂ (X /S) ⊂ ∂ (X /Y ) ∪ ϕ −1 (∂ (Y /S)),}
\NormalTok{         with equality if ψ is locally separated. Taking S = M (k), this becomes}
\NormalTok{(2.11.4)                             ∂ X ⊂ ∂ (X /Y ) ∪ ϕ −1 (∂ Y ).}
\NormalTok{    (5′ ) If k′ /k is an analytic extension field, then}
\NormalTok{                                     ∂ (X k′ /Yk′ ) ⊂ π−1}
\NormalTok{                                                       k′ /k}
\NormalTok{                                                             ∂ (X /Y ).}
\NormalTok{         Taking Y = M (k), this becomes}
\NormalTok{                                           ∂ X k′ ⊂ π−1}
\NormalTok{                                                     k′ /k}
\NormalTok{                                                           ∂ X.}
\NormalTok{We will need the following lemma, which is a consequence of deep results in [CT21] in}
\NormalTok{general, although it has an elementary proof for good spaces.}
\NormalTok{Lemma 2.12. If X is a k{-}analytic space and k′ /k is an analytic field extension, then}
\NormalTok{                                          πk′ /k (∂ X k′ ) = ∂ X .}
\NormalTok{Proof. We have πk′ /k (∂ X k′ ) ⊂ ∂ X by 2.11(5′ ). Let x ∈ ∂ X , and let U ⊂ X be a compact}
\NormalTok{neighborhood of x. We claim that x ∈ ∂ U. If not, then there is an open neighborhood}
\NormalTok{V of x in U such that ∂ V = ; by 2.11(1′ ). If V ′ ⊂ V is a smaller open neighborhood of}
\NormalTok{x in X , then ∂ V ′ = ; as well by 2.11(3′ ) and by (2.11.4) as applied to V ′ → V → M (k).}
\NormalTok{Thus we may assume that V is open in X , which implies by 2.11(1) that x ∈ Int(X ), a}
\NormalTok{contradiction. Hence x ∈ ∂ U. If x ′ ∈ π−1      k′ /k}
\NormalTok{                                                      (x), then x ′ ∈ ∂ Uk′ if and only if x ′ ∈ ∂ X k′}
\NormalTok{for the same reason. We need to show π−1        k′ /k}
\NormalTok{                                                      (x) ∩ ∂ X k′ ̸= ;, so we may replace X by U to}
\NormalTok{assume that X is compact. In this case, the closed subset ∂ X k′ is compact; suppose that}
\NormalTok{it is disjoint from π−1}
\NormalTok{                      k′ /k}
\NormalTok{                            (x). Then πk′ /k (∂ X k′ ) is a closed subset of X not containing x, so}
\NormalTok{there exists an open neighborhood V of x such that Vk′ = π−1            k′ /k}
\NormalTok{                                                                              (V ) ⊂ Int(X k′ ). This means}
\NormalTok{that ∂ Vk′ = ; by 2.11(3,4). But x ∈ ∂ V , and boundarylessness can be checked after}
\NormalTok{extension of scalars by [CT21, Theorem 11.5]. This contradiction proves the claim. □}
\NormalTok{2.13. A k{-}analytic space X is called rig{-}smooth at x ∈ X if it is flat (in the sense of}
\NormalTok{[Duc18]) and if the sheaf of Kähler differentials ΩX (considered on the G{-}topology) is}
\NormalTok{free at x of rank equal to the dimension of X at y. We call X smooth if it is rig{-}smooth}
\NormalTok{and boundaryless. We refer to [Duc18, Chapter 5] for a discussion (note that rig{-}smooth}
\NormalTok{is called quasi{-}smooth there) and comparison to other definitions.}
\NormalTok{             CLASSICALLY PSH AND PLURIHARMONIC FUNCTIONS ON BERKOVICH SPACES                9}

\NormalTok{2.14. Models and formal geometry. A quasi{-}compact admissible formal scheme over}
\NormalTok{k◦ is a quasi{-}compact formal scheme X over k◦ with an open covering by affine formal}
\NormalTok{schemes Spf(A) for k◦ {-}algebras A which are topologically of finite type and flat over k◦ .}
\NormalTok{The generic fiber Xη of X is the strictly k{-}analytic space obtained by gluing the Berkovich}
\NormalTok{spectra M (A ⊗k◦ k), and the special fiber Xs is the scheme over ek obtained by gluing the}
\NormalTok{affine schemes Spec(A ⊗k◦ ek).}

\NormalTok{Definition 2.15. A formal k◦ {-}model X for a quasi{-}compact strictly k{-}analytic space X is a}
\end{Highlighting}
\end{Shaded}

\subsection{GR, newline lines 454--1190}\label{gr-newline-lines-4541190}

\begin{Shaded}
\begin{Highlighting}[]
\NormalTok{smooth strictly analytic curves to arbitrary k{-}analytic curves.}
\NormalTok{3.1 (k{-}Analytic Curves). By a k{-}analytic curve we mean a separated k{-}analytic space of}
\NormalTok{pure dimension 1. We gather the following facts about k{-}analytic curves:}
\NormalTok{   (1) A k{-}analytic curve is good, i.e., every point has an affinoid neighborhood. [Duc24,}
\NormalTok{        Proposition 3.3.7]}
\NormalTok{   (2) A k{-}analytic curve is paracompact. [Duc24, Théorème 4.5.10]}
\NormalTok{   (3) If X is an affinoid k{-}analytic curve, then ∂ X coincides with the Shilov boundary}
\NormalTok{        of X . [Thu05, 2.1.2]}
\NormalTok{An analytic domain in a k{-}analytic curve (for example, an open subset) is again a k{-}}
\NormalTok{analytic curve.}
\NormalTok{3.2 (Thuillier’s Harmonic Functions). Thuillier [Thu05, §2.3] defines a sheaf of harmonic}
\NormalTok{functions on a smooth, strictly k{-}analytic curve for a non{-}trivially valued k, roughly as}
\NormalTok{follows. If X is affinoid, then a continuous function h: X → R is harmonic if, after passing}
\NormalTok{to a finite, separable field extension, it coincides with the composition of a harmonic}
\NormalTok{function on a skeleton of X (in the sense of graph theory) with the retraction to the}
\NormalTok{skeleton. If X is smooth, then a harmonic function on an open subset of X is defined to be}
\NormalTok{a continuous function whose restriction to every strictly affinoid subdomain is harmonic.}
\NormalTok{3.3 (Thuillier’s Subharmonic Functions). For a non{-}trivally valued k, Thuillier [Thu05,}
\NormalTok{Définition 3.1.5] defines a subharmonic function on a smooth, strictly k{-}analytic curve X}
\NormalTok{to be an upper semicontinuous function u: X → R ∪ \{−∞\} that is not identically −∞}
\NormalTok{on any connected component of X and that satisfies the following property: for every}
\NormalTok{strictly k{-}affinoid domain U ⊂ X and every harmonic function h: U → R, if u ≤ h on ∂ U}
\NormalTok{then u ≤ h on U.}
\NormalTok{  Now we extend this definition to the case of non{-}smooth curves. The idea is that a}
\NormalTok{function should be subharmonic if it pulls back to a subharmonic function along any}
\NormalTok{morphism from a smooth curve. However, an inseparable curve over an imperfect field}
\NormalTok{does not admit any non{-}constant morphisms from a smooth curve, so one has to allow}
\NormalTok{an extension of scalars as well. We also consider the constant function −∞ to be sub{-}}
\NormalTok{harmonic, following Demailly [Dem85] in the complex case.}
\NormalTok{Definition 3.4 (Subharmonic functions). Let X be a k{-}analytic curve and let u: X →}
\NormalTok{R ∪ \{−∞\} be an upper semicontinuous function. We say that u is subharmonic provided}
\NormalTok{that, for every non{-}trivially valued analytic extension field k′ /k, every smooth, connected,}
\NormalTok{strictly k′ {-}analytic curve Y , and every morphism Y → X k′ , the composition}
\NormalTok{                                        πk′ /k     u}
\NormalTok{                              Y −→ X k′ −→ X −→ R ∪ \{−∞\}}
\NormalTok{             CLASSICALLY PSH AND PLURIHARMONIC FUNCTIONS ON BERKOVICH SPACES              11}

\NormalTok{is either subharmonic on Y in the sense of [Thu05, Définition 3.1.5], or is identically −∞.}
\NormalTok{3.4.1. If X is a smooth, connected, strictly k{-}analytic curve over a non{-}trivially valued k,}
\NormalTok{then an upper semicontinuous function u: X → R ∪ \{−∞\} is subharmonic in our sense}
\NormalTok{if and only if it is either subharmonic in Thuillier’s sense or is identically −∞. Indeed, if}
\NormalTok{u is subharmonic in our sense then one can take k = k′ and Y = X to see that it is either}
\NormalTok{subharmonic in Thuillier’s sense or is identically −∞. Conversely, if u is identically −∞}
\NormalTok{then it is clearly subharmonic in our sense, and if u is subharmonic in Thuillier’s sense}
\NormalTok{then it is subharmonic in our sense by [Thu05, Proposition 3.1.14 and Corollaire 3.4.5].}
\NormalTok{3.4.2. As a special case of (3.4.1), if X is a smooth, connected, strictly k{-}analytic curve}
\NormalTok{over a non{-}trivially valued k, then a function u: X → R (taking finite values) is subhar{-}}
\NormalTok{monic in our sense if and only if it is subharmonic in Thuillier’s sense.}
\NormalTok{  Here we collect the basic properties of subharmonic functions, which we will prove}
\NormalTok{over the course of this section.}
\NormalTok{Proposition 3.5. Let X be a k{-}analytic curve and let u: X → R ∪ \{−∞\} be an upper}
\NormalTok{semicontinuous function.}
\NormalTok{   (1) (Sheafiness) The subharmonic functions form a sheaf on X .}
\NormalTok{   (2) (Analytic functions) If f ∈ Γ (X , OX ) then log | f | : X → R ∪ \{−∞\} is subharmonic.}
\NormalTok{   (3) (Boundaries) If u is subharmonic on Int(X ), then u is subharmonic on X .}
\NormalTok{   (4) (Cone Property) If u1 , u2 : X → R ∪ \{−∞\} are subharmonic and λ1 , λ2 ∈ R≥0 then}
\NormalTok{       λ1 u1 + λ2 u2 and max\{u1 , u2 \} are subharmonic.}
\NormalTok{   (5) (Limits) The infimum of a decreasing net of subharmonic functions is subharmonic.}
\NormalTok{   (6) (Maximum Principle) If u is subharmonic and u attains a local maximum at x ∈}
\NormalTok{       Int(X ), then u is constant in a neighborhood of x.}
\NormalTok{   (7) (Extension of scalars) Let k′ /k be an analytic extension field. Then u is subharmonic}
\NormalTok{       if and only if u ◦ πk′ /k : X k′ → R ∪ \{∞\} is subharmonic.}
\NormalTok{   (8) (Functoriality) Let Y be a k{-}analytic curve and let f : Y → X be a morphism. If u}
\NormalTok{       is subharmonic then u ◦ f is subharmonic, and the converse holds if f is finite and}
\NormalTok{       surjective.}
\NormalTok{   (9) (Finite Values) If u is subharmonic and u(x) = −∞ for a non{-}rig{-}point x ∈ X ,}
\NormalTok{       then u ≡ −∞ on the irreducible component of Int(X ) ∪ \{x\} containing x.}
\NormalTok{  Let us make some remarks regarding the content of Proposition 3.5.}
\NormalTok{3.5.1. If U ⊂ X is an open subset then U is again a k{-}analytic curve, so it makes sense}
\NormalTok{to speak of subharmonic functions on U. It is immediate from the definition that a sub{-}}
\NormalTok{harmonic function on X restricts to a subharmonic function on U. Hence subharmonic}
\NormalTok{functions form a presheaf on X . Proposition 3.5(1) asserts that this presheaf is a sheaf.}
\NormalTok{3.5.2. The maximum principle (Proposition 3.5(6)) is [Thu05, Proposition 3.4.10] in the}
\NormalTok{smooth, strictly analytic case.}
\NormalTok{3.5.3. In the situation of Proposition 3.5(7), it is immediate from Definition 3.4 that}
\NormalTok{u ◦ πk′ /k is subharmonic when u is subharmonic. This is the difficult direction in the}
\NormalTok{smooth, strictly analytic case: see [Thu05, Corollaire 3.4.5]. This result is hidden in the}
\NormalTok{12                                 W. GUBLER AND J. RABINOFF}

\NormalTok{fact that Definition 3.4 essentially recovers Thuillier’s definition of subharmonicity in that}
\NormalTok{case 3.4.1.}
\NormalTok{3.5.4. In the situation of Proposition 3.5(8), it is immediate from Definition 3.4 that u◦ f}
\NormalTok{is subharmonic when u is subharmonic. The remark in 3.5.3 applies to a lesser extent in}
\NormalTok{this case: see [Thu05, Proposition 3.1.14].}
\NormalTok{3.5.5. For any x ∈ X , the subset Int(X )∪\{x\} is open in X . This is obvious if x ∈ Int(X ), and}
\NormalTok{if x ∈ ∂ X then it follows from the fact that ∂ X is discrete in X . In particular, Int(X ) ∪ \{x\}}
\NormalTok{is again a k{-}analytic curve, so it makes sense to speak of its irreducible components}
\NormalTok{in Proposition 3.5(9). Irreducible components can only meet at rig{-}points (their inter{-}}
\NormalTok{sections are Zariski{-}closed), so a non{-}rig{-}point x is contained in a unique irreducible}
\NormalTok{component of Int(X ) ∪ \{x\}; this is implicit in Proposition 3.5(9).}
\NormalTok{3.5.6. Proposition 3.5(9) implies that if u is subharmonic and u(x) = −∞ for a non{-}rig{-}}
\NormalTok{point x ∈ X , then u ≡ −∞ in a neighborhood of x.}
\NormalTok{3.6. In order to extend Thuillier’s results, we will pass from a k{-}analytic curve to a}
\NormalTok{smooth, strictly k{-}analytic curve in four steps:}
\NormalTok{      (1) Extend scalars to a field as in 2.10 to make the curve strictly analytic.}
\NormalTok{                                                                                          −∞}
\NormalTok{      (2) If char(k) = p, then extend scalars to the completion of the perfect closure k p .}
\NormalTok{      (3) Replace the analytic curve by the underlying reduced curve.}
\NormalTok{      (4) Pass to the interior of the normalization, which is smooth (see 2.13).}
\NormalTok{We begin with some lemmas about this process.}
\NormalTok{                                                                                 −∞}
\NormalTok{Lemma 3.7. Suppose that char(k) = p \textgreater{} 0. Let k′ be the completion of k p . Then k′ is}
\NormalTok{perfect, and if X is a k{-}analytic space then πk′ /k : X k′ → X is a homeomorphism such that}
\NormalTok{π(∂ X k′ ) = ∂ X .}
\NormalTok{Proof. We leave it as an exercise for the reader to verify that k′ is perfect. The fact that}
\NormalTok{πk′ /k is a homeomorphism is [Duc09, Remarque 0.5], and the assertion about boundaries}
\NormalTok{is Lemma 2.12.                                                                            □}
\NormalTok{     We make the following observation to mirror Lemma 3.7. Its proof is immediate.}
\NormalTok{Lemma 3.8. Let X be a strictly k{-}analytic space and let ι : X red ,→ X be the inclusion of the}
\NormalTok{underlying reduced space. Then ι is a homeomorphism of underlying topological spaces, and}
\NormalTok{ι(∂ X red ) = ∂ X .}
\NormalTok{  We refer the reader to [Ber90, before Proposition 3.1.8] for the construction of the}
\NormalTok{normalization.}
\NormalTok{Lemma 3.9. Let X be a reduced, strictly k{-}analytic curve and let ν: Xe → X be its normal{-}}
\NormalTok{ization. Then the following hold:}
\NormalTok{    (1) Xe is rig{-}smooth if k is perfect.}
\NormalTok{    (2) ν is an isomorphism outside of a discrete set of rig{-}points of X .}
\NormalTok{                                     ∼}
\NormalTok{    (3) ν induces a bijection ∂ Xe −→ ∂ X .}
\NormalTok{    (4) ν is a topological quotient.}
\NormalTok{              CLASSICALLY PSH AND PLURIHARMONIC FUNCTIONS ON BERKOVICH SPACES                       13}

\NormalTok{                                                                             e = ν−1 (U) →}
\NormalTok{    (5) Let x ∈ X , let U be an open neighborhood of x, and consider the map U}
\NormalTok{        U. If U is a sufficiently small connected open neighborhood of x, then:}
\NormalTok{         (a) Ue → U is injective on connected components of U.  e}
\NormalTok{                                                                 ∼}
\NormalTok{         (b) Ue → U restricts to an isomorphism U  e \textbackslash{} ν−1 (x) −→ U \textbackslash{} \{x\}.}
\NormalTok{         (c) Every connected component of U   e intersects ν−1 (x).}
\NormalTok{    (6) A function h: X → R is R{-}PL if and only if h ◦ ν is R{-}PL.}



\NormalTok{Proof. The normalization Xe is locally the Berkovich spectrum of a Dedekind domain, so it}
\NormalTok{is rig{-}smooth when k is perfect because it is regular; this proves (1). Normalization is an}
\NormalTok{isomorphism away from codimension 1 as in algebraic geometry, which yields (2). Since}
\NormalTok{ν is finite, we have ∂ Xe = ν−1 (∂ X ) by the properties (2′ ) and (4′ ) of 2.11, so (3) follows}
\NormalTok{from (2) because the boundaries do not contain any rig{-}points. Assertion (4) was noted}
\NormalTok{in [Duc24, (3.1.3)].}
\NormalTok{    Now we prove (5). Since normalization is a local construction, we may shrink X to}
\NormalTok{assume that it is normal away from x. Suppose that ν−1 (x) = \{x 1 , . . . , x n \}, and choose}
\NormalTok{disjoint, connected open neighborhoods U1′ , . . . , Un′ of x 1 , . . . , x n , respectively. Let U ′ =}
\NormalTok{Sn                                                                         ∼}
\NormalTok{   i=1}
\NormalTok{       ν(Ui′ ).FSince ν induces an isomorphism Xe \textbackslash{} \{x 1 , . . . , x n \} −→ X \textbackslash{} \{x\}, it is clear that}
\NormalTok{                      n}
\NormalTok{ν−1 (U ′ ) = i=1 Ui′ . Since ν is a topological quotient, we have that U ′ is open, and ν is}
\NormalTok{injective on each Ui′ by construction.}
\NormalTok{    Suppose now that U ⊂ U ′ is a connected open neighborhood of x. By [Ber90, Theorem}
\NormalTok{3.2.1], U is path{-}connected. Clearly U satisfies (a) and (b). Let Ui = ν−1 (U) ∩ Ui′ , so that}
\NormalTok{                      e = ν−1 (U) = n Ui . In order to prove (c), it is enough to show that Ui}
\NormalTok{                                      F}
\NormalTok{x i ∈ Ui and U                          i=1}
\NormalTok{is path{-}connected. Let y ∈ Ui \textbackslash{} \{x i \}, and let γ: [0, 1] → U be a path from π( y) to x.}
\NormalTok{We may assume that γ−1 (x) = \{1\}, so that γ restricts to a path [0, 1) → U \textbackslash{} \{x\}. Let}
\NormalTok{γ}
\NormalTok{e : [0, 1) → U       e \textbackslash{} ν−1 (x) be the composition of γ|[0,1) with the inverse of the isomorphism}
\NormalTok{                      ∼}
\NormalTok{e \textbackslash{} ν−1 (x) −→ U \textbackslash{} \{x\}. Then γ}
\NormalTok{U                                        e(0) = y, so the image of γ e is contained in Ui . Extend γ  e}
\NormalTok{to a function [0, 1] → Ui by setting γ        e(1) = x i . One checks continuity of γ e at 1 using the}
\NormalTok{fact that an open neighborhood of x is the same as a collection of open neighborhoods}
\NormalTok{of x 1 , . . . , x n , as in the previous paragraph. This shows that Ui is path{-}connected.}
\NormalTok{    In the situation of (6), if h is R{-}PL then h◦ν is R{-}PL by [GJR21, Lemma 5.4(1)]. Suppose}
\NormalTok{then that h ◦ ν is R{-}PL. Piecewise linearity of h is local on X , and ν is an isomorphism}
\NormalTok{away from a finite set of (non{-}normal) rig{-}points, so it is enough to check that h is R{-}PL}
\NormalTok{in a neighborhood of a single rig{-}point x. Let x ′ ∈ Xe be a preimage ofPx. There is a}
\NormalTok{                                                                                          n}
\NormalTok{covering of Xe by strictly affinoid domains on which h ◦ ν has the form i=1 λi log | f i |,}
\NormalTok{where λi ∈ R and f i ∈ O × . Since x ′ is      Pan rig{-}point, it is contained in the interior   of one}
\NormalTok{of these affinoid domains, so h ◦ ν = i=1 λi log | f i | in a neighborhood of x . But if f is}
\NormalTok{                                                                                          ′}

\NormalTok{invertible in a neighborhood of a rig{-}point x ′ , then | f | is constant on a neighborhood of}
\NormalTok{x ′ (this elementary fact can be seen as a very special case of [Duc12, Theorem 3.4]). It}
\NormalTok{follows that h◦ν is constant in a neighborhood of x ′ . Suppose that ν−1 (x) = \{x 1 , . . . , x n \}.}
\NormalTok{For eachSn i choose an open neighborhood Ui of x i on which h ◦ ν is constant. As above,}
\NormalTok{U = i=1 ν(Ui ) is an open neighborhood of x, and clearly h is constant, hence R{-}PL,}
\NormalTok{on U.                                                                                                 □}
\NormalTok{14                                  W. GUBLER AND J. RABINOFF}

\NormalTok{Proof of Proposition 3.5(1). This follows from the corresponding fact for subharmonic}
\NormalTok{functions on smooth, strictly analytic curves [Thu05, Corollaire 3.1.13], except that one}
\NormalTok{has to be careful about where a function takes the value −∞. Suppose then that \{Ui \}}
\NormalTok{is an open cover of X and that u: X → R ∪ \{−∞\} is subharmonic on each Ui . Note that}
\NormalTok{u is upper semicontinuous. Let k′ /k be an analytic extension field with k′ non{-}trivially}
\NormalTok{valued, let Y be a smooth, connected, strictly k′ {-}analytic curve, and let f : Y → X k′ be a}
\NormalTok{morphism. Assume that u ◦ πk′ /k ◦ f ̸≡ −∞, so we must show that u ◦ πk′ /k ◦ f is sub{-}}
\NormalTok{harmonic in Thuillier’s sense. Let Ui′ = π−1 k′ /k}
\NormalTok{                                                   (Ui ) = (Ui )k′ and let Yi = f −1 (Ui′ ). Since u}
\NormalTok{is subharmonic on Ui , the composition u ◦ πk′ /k ◦ f |Yi is either subharmonic in Thuillier’s}
\NormalTok{sense or is identically −∞ on each connected component of Yi .}
\NormalTok{   Let W1 be the set of all points y ∈ Y that admit a neighborhood on which u ◦ πk′ /k ◦ f}
\NormalTok{is subharmonic in Thuillier’s sense, and let W2 be the set of all points y ∈ Y that admit a}
\NormalTok{neighborhood on which u ◦ πk′ /k ◦ f ≡ −∞. Clearly W1 and W2 are open, and they cover}
\NormalTok{Y since Y is covered by the connected components of the Yi . They are disjoint because if}
\NormalTok{u ◦ πk′ /k ◦ f is subharmonic in Thuillier’s sense on a connected neighborhood of y ∈ Y ,}
\NormalTok{then it cannot be identically −∞ on any neighborhood of y [Thu05, Lemme 3.1.9].}
\NormalTok{Since Y is connected and u ◦ πk′ /k ◦ f ̸≡ −∞, we have Y = W1 . The result now follows}
\NormalTok{from [Thu05, Corollaire 3.1.13].                                                                  □}


\NormalTok{Proof of Proposition 3.5(2). Let f be an analytic function on X . The function | f | is con{-}}
\NormalTok{tinuous and hence u := log | f | is upper semicontinuous as a continuous function with}
\NormalTok{values in R ∪ \{−∞\}. According to our definition of subharmonic functions, we have to}
\NormalTok{show that after base change of u with respect to an analytic field extension k′ /k with k′}
\NormalTok{non{-}trivially valued and then pull{-}back to a smooth strictly k′ {-}analytic curve, the resulting}
\NormalTok{function is subharmonic. This boils down to showing that u is subharmonic in the special}
\NormalTok{case of a smooth strictly k{-}analytic curve X . For this, we may assume X connected. If}
\NormalTok{f = 0, then u is identically −∞ and hence subharmonic in the sense of Definition 3.4. If}
\NormalTok{f is not identically zero, then the Poincaré–Lelong formula [Thu05, Proposition 3.3.15]}
\NormalTok{and the characterization in terms of positive currents [Thu05, Proposition 3.4.4] show}
\NormalTok{that u is subharmonic in the sense of Thuillier.                                            □}


\NormalTok{  We will use the following lemma in the proofs of Proposition 3.5(3,9). If X is a k{-}}
\NormalTok{analytic space and k′ /k is an analytic extension field, then π−1}
\NormalTok{                                                               k′ /k}
\NormalTok{                                                                     (x) = M (H (x)⊗}
\NormalTok{                                                                                   b k k′ )}
\NormalTok{has a finite Shilov boundary for any Abhyankar point x ∈ X by [Duc24, 3.2.14]; we}
\NormalTok{denote this boundary by x [k′ ] .}

\NormalTok{Lemma 3.10. Let X be a k{-}analytic curve and let k′ /k be an analytic field extension. Then}
\NormalTok{for all x ∈ ∂ X we have (∂ X k′ ) ∩ π−1}
\NormalTok{                                     k′ /k}
\NormalTok{                                           (x) = x [k′ ] .}

\NormalTok{Proof. Replacing X by an affinoid neighborhood of x, we may assume that X = M (A )}
\NormalTok{is affinoid. By a result of Ducros [Duc12, Lemme 3.1], there exists a morphism g : X →}
\NormalTok{B1 (0, s) for some s ∈ R\textgreater{}0 such that x ∈ g −1 (ηs ), where ηs is the maximal point of B1 (0, s).}
\NormalTok{                CLASSICALLY PSH AND PLURIHARMONIC FUNCTIONS ON BERKOVICH SPACES                                 15}

\NormalTok{This is discussed in detail in Section 6. Consider the commutative square}

\NormalTok{                                                          g′}
\NormalTok{                                                           → B1k′ (0, s)}
\NormalTok{                                                      ←}
\NormalTok{                                               X k′}




\NormalTok{                                                ←}




\NormalTok{                                                                  ←}
\NormalTok{                                           πk′ /k                   πk′ /k}




\NormalTok{                                                                 →}
\NormalTok{                                               →}
\NormalTok{                                                          g → B (0, s).}
\NormalTok{                                                      ←        1}
\NormalTok{                                                X}

\NormalTok{Since ∂ B1k′ (0, s) contains (only) its maximal point ηs,k′ , we have (g ′ )−1 (ηs,k′ ) ⊂ ∂ X k′}
\NormalTok{by (2.11.4) (which is an equality because B1k′ (0, s) is separated). The morphism g is an}
\NormalTok{Abhyankar presentation of x in the sense of [Duc24, 3.2.12.2], so by Proposition 3.2.13}
\NormalTok{of ibid, we have x [k′ ] = (g ′ )−1 (ηs,k′ ) ∩ π−1}
\NormalTok{                                                k′ /k}
\NormalTok{                                                      (x) ⊂ (∂ X k′ ) ∩ π−1}
\NormalTok{                                                                         k′ /k}
\NormalTok{                                                                               (x).}
\NormalTok{   For the reverse inclusion, let x ∈ (∂ X k′ ) ∩ πk′ /k (x), and suppose that x ′ ∈}
\NormalTok{                                        ′                   −1}
\NormalTok{                                                                                         / x [k′ ] . Since}
\NormalTok{∂ X k′ is the Shilov boundary of X k′ = M (A ⊗         b k k ), by [Ber90, Corollary 2.4.5] there exist}
\NormalTok{                                                            ′}

\NormalTok{f ∈ A⊗   b k k and ϵ \textgreater{} 0 such that | f | attains its maximum value at x ′ and such that}
\NormalTok{              ′}

                          
\NormalTok{                            y ∈ X k′ | f ( y)| \textgreater{} | f (x ′ )| − ϵ ⊂ X k′ \textbackslash{} x [k′ ] .}

\NormalTok{In particular, the restriction of f to π−1}
\NormalTok{                                        k′ /k}
\NormalTok{                                              (x) does not attain its maximum value on x [k′ ] ,}
\NormalTok{which cannot happen. Thus x ∈ x [k′ ] , so (∂ X k′ ) ∩ π−1}
\NormalTok{                                 ′}
\NormalTok{                                                         k′ /k}
\NormalTok{                                                               (x) ⊂ x [k′ ] .             □}

\NormalTok{Proof of Proposition 3.5(3). Suppose that u is subharmonic on Int(X ). Let k′ /k be an}
\NormalTok{analytic extension field with k′ non{-}trivially valued, let Y be a smooth, connected, strictly}
\NormalTok{k′ {-}analytic curve, and let f : Y → X k′ be a morphism. We must prove that u ◦ πk′ /k ◦}
\NormalTok{f is subharmonic in Thuillier’s sense or is identically −∞. Since ∂ Y = ;, it follows}
\NormalTok{from 2.11(4) that f (Y ) ⊂ Int(X k′ ). If πk′ /k ◦ f (Y ) ⊂ Int(X ) then u◦πk′ /k ◦ f is subharmonic}
\NormalTok{in Thuillier’s sense or is identically −∞ because u is subharmonic on Int(X ). Suppose}
\NormalTok{then that πk′ /k ◦ f ( y) = x ∈ ∂ X for some y ∈ Y , so that f ( y) ∈ Int(X k′ ) ∩ π−1                 k′ /k}
\NormalTok{                                                                                                             (x).}
\NormalTok{    We claim that πk′ /k (x)∩Int(X k′ ) is both open and closed in Int(X k′ ). By [Duc24, Propo{-}}
\NormalTok{                         −1}

\NormalTok{sition 5.3.4], if U is a connected component of π−1               k′ /k}
\NormalTok{                                                                        (X \textbackslash{} \{x\}), then the closure of U in}
\NormalTok{X k′ is contained in πk′ /k (X \textbackslash{} \{x\}) ∪ x [k′ ] . Let U be the connected component of X k′ \textbackslash{} x [k′ ]}
\NormalTok{                            −1                                ′}

\NormalTok{containing U. The above shows that U is closed in U ′ . Since U is open in X (an analytic}
\NormalTok{space is locally connected [Ber93, Remark 1.2.4(iii)] and hence its connected compo{-}}
\NormalTok{nents are open [Sta21, Lemma 04ME]), it is open in U ′ , so U = U ′ since U ′ is connected.}
\NormalTok{Hence π−1  k′ /k}
\NormalTok{                 (X \textbackslash{} \{x\}) is a union of connected components of X k′ \textbackslash{} x [k′ ] , so the same is}
\NormalTok{true of πk′ /k (x) \textbackslash{} x [k′ ] . Since X k′ \textbackslash{} x [k′ ] is itself a k{-}analytic space as x [k′ ] is a finite set, its}
\NormalTok{           −1}

\NormalTok{connected components are open, so we conclude that π−1                       k′ /k}
\NormalTok{                                                                                   (x) \textbackslash{} x [k′ ] is open and closed}
\NormalTok{in X k′ \textbackslash{} x [k′ ] . Since x ∈ ∂ X , we have (πk′ /k (x)) ∩ ∂ X k′ = x [k′ ] by Lemma 3.10, which}
\NormalTok{                                                         −1}

\NormalTok{shows Int(X k′ ) ⊂ X k′ \textbackslash{} x [k′ ] and π−1   k′ /k}
\NormalTok{                                                  (x) \textbackslash{} x [k′ ] = π−1 k′ /k}
\NormalTok{                                                                            (x) ∩ Int(X k′ ). We conclude that}
\NormalTok{πk′ /k (x) \textbackslash{} x [k′ ] is both open and closed in Int(X k′ ), as claimed.}
\NormalTok{  −1}

\NormalTok{    Since Y is connected, f (Y ) ⊂ Int(X k′ ), and f ( y) ∈ Int(X k′ ) ∩ π−1            k′ /k}
\NormalTok{                                                                                              (x), we deduce from}
\NormalTok{the above claim that f (Y ) ⊂ Int(X k′ )∩πk′ /k (x). Hence πk′ /k ◦ f is constant, so u◦πk′ /k ◦ f}
\NormalTok{                                                      −1}

\NormalTok{is constant. It follows that u is subharmonic in Thuillier’s sense or is identically −∞. □}
\NormalTok{16                                        W. GUBLER AND J. RABINOFF}

\NormalTok{Proof of Proposition 3.5(4,5). These follow quickly from the corresponding facts for sub{-}}
\NormalTok{harmonic functions on smooth, strictly analytic curves [Thu05, Proposition 3.1.8].     □}

\NormalTok{Proof of Proposition 3.5(6). Suppose first that X is strictly k{-}analytic for k non{-}trivially}
\NormalTok{                                                                                −∞}
\NormalTok{valued. If char(k) = p \textgreater{} 0 then let k′ be the completion of k p ; otherwise let k′ = k.}
\NormalTok{Let X ′ = X k′ and let π = πk′ /k : X ′ → X , so π is a homeomorphism preserving boundaries}
\NormalTok{by Lemma 3.7. Let ι : X red    ′}
\NormalTok{                                  ,→ X ′ be the inclusion of the underlying reduced space and let}
\NormalTok{ν: Xe ′ → X red′}
\NormalTok{                  be the normalization. Note that ν and ι are finite, so ∂ Xe ′ = (ι ◦ ν)−1 (∂ X ′ ),}
\NormalTok{and hence ∂ Xe ′ = (π◦ι ◦ν)−1 (∂ X ). Let Y = Int(Xe ′ ) and let f = ι ◦ν|Y : Y → X ′ . Then Y is}
\NormalTok{smooth and strictly analytic, so u ◦ π ◦ f is by definition either subharmonic in Thuillier’s}
\NormalTok{sense or identically −∞ on every connected component of Y .}
\NormalTok{   Suppose that u attains a local maximum at x ∈ Int(X ). If u(x) = −∞ then u must}
\NormalTok{be constant in a neighborhood of x by semicontinuity of u. Assume then that u(x) \textgreater{}}
\NormalTok{−∞. Both π and ι are homeomorphisms, so it suffices to show that u ◦ π ◦ ι : X red                 ′}
\NormalTok{                                                                                                      →}
\NormalTok{R ∪ \{−∞\} is constant in a neighborhood of the unique point x ∈ Int(X red ) mapping to}
\NormalTok{                                                                                 ′       ′}

\NormalTok{x. Choose a connected open neighborhood U ⊂ Int(X red                ′}
\NormalTok{                                                                       ) of x ′ as in Lemma 3.9(5), on}
\NormalTok{which u ◦ π ◦ ι attains a maximum at x ′ . By Lemma 3.9(3), we have ν−1 (U) ⊂ Y =}
\NormalTok{Int(Xe ′ ), so u ◦ π ◦ f is again either subharmonic in Thuillier’s sense or identically −∞}
\NormalTok{on each connected component C of ν−1 (U); since u(x) \textgreater{} −∞ and C ∩ ν−1 (x ′ ) ̸= ;, the}
\NormalTok{former is true. The function u◦π◦ f attains its maximum on every connected component}
\NormalTok{of ν−1 (U), so by the maximum principle for subharmonic functions on smooth, strictly}
\NormalTok{analytic curves [Thu05, Proposition 3.1.11], the function u ◦ π ◦ f is constant on every}
\NormalTok{connected component of ν−1 (U). It follows that u ◦ π ◦ ι is constant on U.}
\NormalTok{   Now we reduce to the strictly analytic case handled above. We may assume that X}
\NormalTok{is affinoid. As explained in 2.10, there is an analytic extension field k′ = k r for some}
\NormalTok{r ∈ R\textgreater{}0n}
\NormalTok{             such that X k′ is strictly k′ {-}analytic. Let σ : X → X k′ be the Shilov section of πk′ /k .}
\NormalTok{Then x = σ(x) ∈ Int(X k′ ) because π−1}
\NormalTok{           ′}
\NormalTok{                                                k′ /k}
\NormalTok{                                                      (Int(X )) ⊂ Int(X k′ ) by 2.11(5). It is immediate}
\NormalTok{from the definition that u ◦ πk′ /k is subharmonic, and it attains a local maximum at x ′ ,}
\NormalTok{so u ◦ πk′ /k is constant in an open neighborhood U ′ of x ′ . Then u is constant on the open}
\NormalTok{set U = σ−1 (U ′ ).                                                                                   □}

\NormalTok{Proof of Proposition 3.5(7). It is immediate from the definition that subharmonicity of}
\NormalTok{u implies subharmonicity of u ◦ πk′ /k . Suppose then that u ◦ πk′ /k is subharmonic. Let}
\NormalTok{k′′ /k be another analytic extension field with k′′ non{-}trivially valued, and let f : Y →}
\NormalTok{X k′′ be a morphism from a smooth, connected, strictly k′′ {-}analytic curve. Assume that}
\NormalTok{u ◦ πk′′ /k ◦ f ̸≡ −∞, so we need to show that u ◦ πk′′ /k ◦ f is subharmonic in Thuillier’s}
\NormalTok{sense. Let K be a non{-}Archimedean field that is simultaneously an analytic extension of}
\NormalTok{k′ and k′′ 2.2, and consider the following commutative diagram:}

\NormalTok{                               ←   fK            πK/k′}
\NormalTok{                                                 ←}
\NormalTok{                          YK            → XK          → X k′}
\NormalTok{                          ←}



\NormalTok{                                           ←}



\NormalTok{                                                            ←}




\NormalTok{(3.10.1)                                     πK/k′′             πk′ /k}
\NormalTok{                          →}



\NormalTok{                                          →}



\NormalTok{                                                            →}




\NormalTok{                               ←                 ←                ←}
\NormalTok{                           Y       f}
\NormalTok{                                       → X k′′         →}
\NormalTok{                                                  πk′′ /k   X            u → R ∪ \{−∞\}.}
\NormalTok{             CLASSICALLY PSH AND PLURIHARMONIC FUNCTIONS ON BERKOVICH SPACES                 17}

\NormalTok{Since u ◦ πk′ /k is subharmonic, the composition u ◦ πK/k ◦ f K is, on each connected compo{-}}
\NormalTok{nent C of YK , either subharmonic in Thuillier’s sense or identically −∞; since C surjects}
\NormalTok{onto Y [Duc18, 1.5.6] and u ◦ πk′′ /k ◦ f ̸≡ −∞, the latter does not happen. Replacing}
\NormalTok{X by Y and k′ /k by K/k′′ , we are reduced to showing that if X is smooth and strictly}
\NormalTok{k{-}analytic for k non{-}trivially valued and u ◦ πk′ /k is subharmonic in Thuillier’s sense, then}
\NormalTok{so is u.}
\NormalTok{   We can do this directly using Thuillier’s definition. (The converse is the difficult di{-}}
\NormalTok{rection: see [Thu05, Corollaire 3.4.5].) Let U ⊂ X be a strictly k{-}affinoid domain and}
\NormalTok{let h: U → R be a harmonic function such that u ≤ h on ∂ U. Then h ◦ πk′ /k : Uk′ =}
\NormalTok{π−1}
\NormalTok{  k′ /k}
\NormalTok{        (U) → R is harmonic by [Thu05, Proposition 2.3.18]. We have πk′ /k (∂ Uk′ ) = ∂ U by}
\NormalTok{Lemma 2.12, so u ◦ πk′ /k ≤ h ◦ πk′ /k on ∂ Uk′ . Since u ◦ πk′ /k is subharmonic in Thuillier’s}
\NormalTok{sense, this implies u ◦ πk′ /k ≤ h ◦ πk′ /k on Uk′ , so u ≤ h on U since πk′ /k : Uk′ ↠ U is}
\NormalTok{surjective. Thus u is subharmonic in Thuillier’s sense, as claimed.                          □}

\NormalTok{Proof of Proposition 3.5(8). It is immediate that subharmonicity of u implies subharmonic{-}}
\NormalTok{ity of u ◦ f . Suppose then that f is finite and surjective and that u ◦ f is subharmonic. Let}
\NormalTok{k′ /k be an analytic extension field with k′ non{-}trivially valued and let g : Z → X k′ be a}
\NormalTok{morphism from a smooth, connected, strictly k′ {-}analytic curve. Suppose that u◦πk′ /k ◦g ̸≡}
\NormalTok{−∞, so we need to show that u ◦ πk′ /k ◦ g is subharmonic in Thuillier’s sense. Let}
\NormalTok{W = Z ×X k′ Yk′ , and consider the following commutative diagram:}

\NormalTok{                                   ←   g′           πk′ /k}
\NormalTok{                                                    ←}
\NormalTok{                          W             → Yk′           → Y}
\NormalTok{                           ←}




\NormalTok{                                                             ←}
\NormalTok{                                            ←}




\NormalTok{(3.10.2)                f ′′           □       f′                f}
\NormalTok{                          →}



\NormalTok{                                            →}



\NormalTok{                                                             →}




\NormalTok{                                   ←                ←                ←}
\NormalTok{                               Z       g→   X k′        →}
\NormalTok{                                                    πk′ /k   X           u → R ∪ \{−∞\}.}


\NormalTok{The composition (u ◦ f ) ◦ πk′ /k ◦ g ′ is subharmonic because it is pulled back from u ◦ f by}
\NormalTok{an extension of scalars and the morphism g ′ : W → Yk′ . Since f is finite and surjective,}
\NormalTok{so are f ′ and f ′′ by 2.2, so we may replace X by Z and Y by W to assume X is a smooth,}
\NormalTok{connected, strictly k{-}analytic curve for k non{-}trivially valued and that u ̸≡ −∞. We can}
\NormalTok{check subharmonicity of u (in Thuillier’s sense) after extending scalars as in the proof of}
\NormalTok{Proposition 3.5(7), so we may assume in addition that k is perfect.}
\NormalTok{   Let ι : Yred ,→ Y be the inclusion of the underlying reduced space, and let ν: Ye → Yred be}
\NormalTok{the normalization. Since k is perfect, the normal curve Ye is rig{-}smooth, and since Ye → X}
\NormalTok{is finite and ∂ X = ; (as X is smooth), the curve Ye is smooth as well: see Lemmas 3.8}
\NormalTok{and 3.9. Both ι and ν are finite and surjective, so f ◦ ι ◦ ν: Ye → X is again finite and}
\NormalTok{surjective. By Definition 3.4 (with k = k′ ), subharmonicity of u ◦ f means that, on each}
\NormalTok{connected component C of Ye , the function (u◦ f )◦ι◦ν is either subharmonic in Thuillier’s}
\NormalTok{sense or is identically −∞. Since C → X is surjective (it is finite, hence it is a closed}
\NormalTok{map, and it is an open map by [Ber90, Lemma 3.2.4], so its image is all of X since X is}
\NormalTok{connected), the latter does not happen, so u ◦ f ◦ ι ◦ ν is subharmonic in Thuillier’s sense.}
\NormalTok{Thus we replace Y by Ye and f by f ◦ ι ◦ ν to assume X and Y are both smooth, and that}
\NormalTok{u ◦ f is subharmonic in Thuillier’s sense.}
\NormalTok{18                                W. GUBLER AND J. RABINOFF}

\NormalTok{   At this point, one proves that u is subharmonic in Thuillier’s sense directly from Thuil{-}}
\NormalTok{lier’s definition, as in the proof of Proposition 3.5(7), this time using [Thu05, Proposi{-}}
\NormalTok{tion 2.3.19]. The converse is the difficult direction: see [Thu05, Proposition 3.1.14]. □}
\NormalTok{  We will need the following lemma in the proof of Proposition 3.5(9). (Note that}
\NormalTok{Lemma 3.11 itself follows from Proposition 3.5(9) as we have seen in 3.5.5.)}
\NormalTok{Lemma 3.11. Let X be a strictly k{-}analytic curve over a non{-}trivially valued field k and let}
\NormalTok{u: X → R ∪ \{−∞\} be a subharmonic function. If x ∈ ∂ X and u(x) = −∞, then u ≡ −∞}
\NormalTok{on a neighborhood of x.}
\NormalTok{                                                          −∞}
\NormalTok{Proof. If char(k) = p \textgreater{} 0 let k′ be the completion of k p , and otherwise let k′ = k. Let}
\NormalTok{X ′ = X k′ and π = πk′ /k : X ′ → X , let ι : X red}
\NormalTok{                                                ′}
\NormalTok{                                                     ,→ X ′ be the inclusion of the underlying}
\NormalTok{reduced space, and let ν: Xe ′ → X red}
\NormalTok{                                     ′}
\NormalTok{                                        be the normalization. Then π and ι are homeomor{-}}
\NormalTok{phisms preserving boundaries by Lemmas 3.7 and 3.8, and ν is a boundary{-}preserving}
\NormalTok{isomorphism in a neighborhood of the unique point x ′ ∈ Xe ′ lying over x by Lemma 3.9.}
\NormalTok{The composition u ◦ π ◦ ι ◦ ν is subharmonic, so we may replace X by Xe ′ to assume that}
\NormalTok{X is rig{-}smooth. After shrinking X , we may assume that X is strictly affinoid. By [Ber90,}
\NormalTok{Corollary 1.3.6], if K is the completion of an algebraic closure of k then X K → X is an}
\NormalTok{open map, and by Lemma 2.12 there is a point of ∂ X K lying over x, so we may extend}
\NormalTok{scalars to assume that k is algebraically closed.}
\NormalTok{   By the semistable reduction theorem [Thu05, Théorème 2.3.8], there exists a semistable}
\NormalTok{(formal) model X of X . Let red: X → Xs denote the reduction map, which is anti{-}}
\NormalTok{continuous. Let ξ ∈ Xs be any point. As explained in [BPR13, Theorem 4.6] (see also}
\NormalTok{[BL85, Propositions 2.2 and 2.3]), if ξ is a generic point then red−1 (ξ) is a single type{-}2}
\NormalTok{point of X ; if ξ is a smooth closed point then red−1 (ξ) is isomorphic to an open unit ball;}
\NormalTok{and if ξ is a node then red−1 (ξ) is isomorphic to an open annulus. In particular, if ξ is a}
\NormalTok{closed point then red−1 (ξ) is connected and open.}
\NormalTok{   By [CD12, Lemme 6.5.1], the closure of ζ = red(x) is not proper, so ζ is a generic}
\NormalTok{point of Xs . Let Y ⊂ Xs denote its closure. Then red−1 (Y ) is an open neighborhood of}
\NormalTok{x in X . Let ξ ∈ Y be a closed point and let U = red−1 (ξ). Then x is a limit point of U:}
\NormalTok{one can use [BPR13, Lemma 3.2], or one can see this using retraction to the skeleton}
\NormalTok{S(X) from [Thu05, Théorème 2.2.10]; in this case, x is a vertex of S(X), and either U}
\NormalTok{deformation retracts onto x if ξ is smooth, or U contains an edge of S(X) adjacent to}
\NormalTok{x when ξ is a node. It follows that red−1 (Y ) is contained in the union of \{x\} with all}
\NormalTok{connected components U of Int(X ) having x as a limit point. We will show that u ≡ −∞}
\NormalTok{on every such U.}
\NormalTok{   We adapt the proof of [Thu05, Lemme 3.1.9]. Suppose that u ̸≡ −∞ on some such}
\NormalTok{U. Then u|U is subharmonic in Thuillier’s sense as U is smooth and connected and u is}
\NormalTok{subharmonic in our sense. By [Thu05, Lemme 2.3.9], for every N ∈ Z≥0 , there exists a}
\NormalTok{unique harmonic function hN : X → R such that hN (x) = −N and}
\NormalTok{                                           u( y) if u( y) \textgreater{} −∞}
                                         
\NormalTok{                               hN ( y) =}
\NormalTok{                                           0        otherwise}
\NormalTok{for every y ∈ ∂ X \textbackslash{} \{x\}. By construction, we have u − hN ≤ 0 on ∂ X . By [Thu05,}
\NormalTok{Corollaire 3.1.12], if the function u − hN attains a local maximum at a point y ∈ U, then}
\NormalTok{              CLASSICALLY PSH AND PLURIHARMONIC FUNCTIONS ON BERKOVICH SPACES                       19}

\NormalTok{u − hN takes a constant value C on U. By semicontinuity, we have}
                                                             
\NormalTok{                   −∞ = u(x) − hN (x) ≥ lim sup u(z) − hN (z) = C,}
\NormalTok{                                                   z→x}
\NormalTok{                                                   z∈U}

\NormalTok{which is a contradiction. Thus u−hN does not attain a local maximum on U. It does attain}
\NormalTok{a maximum on the closure U, which is compact. We have U ⊂ U ∪ ∂ X and u − hN ≤ 0 on}
\NormalTok{∂ X , so u−hN ≤ 0 on U. Clearly −hN (x) → ∞ as N → ∞, so by [Thu05, Lemme 2.3.10],}
\NormalTok{we have −hN → ∞ on U. Since u ≤ hN on U for all N , it follows that u ≡ −∞ on U. □}
\NormalTok{Proof of Proposition 3.5(9). We replace X by Int(X )∪\{x\} to assume that ∂ X = ; or ∂ X =}
\NormalTok{\{x\}. Replacing X by the irreducible component of X containing x (see 3.5.5), we assume}
\NormalTok{in addition that X is irreducible. Let k′ = k r /k be an analytic extension field for some}
\NormalTok{r ∈ R\textgreater{}0 n}
\NormalTok{           as in 2.10 such that X k′ is strictly k′ {-}analytic. Note that X k′ is again irreducible:}
\NormalTok{any irreducible component of X k′ is defined over a finite, separable extension of k inside}
\NormalTok{k′ [Duc18, 1.5.6], but k is algebraically closed in k′ [Duc24, 3.2.11]. The fiber π−1         k′ /k}
\NormalTok{                                                                                                    (x)}
\NormalTok{has a unique Shilov boundary point x ′ = σ(x), so by Lemma 3.10, if ∂ X = \{x\} then}
\NormalTok{∂ X k′ = \{x ′ \}; of course, if ∂ X = ; then ∂ X k′ = ; as well. Replacing X by X k′ and x by}
\NormalTok{x ′ , we may assume that X is strictly k{-}analytic and irreducible over a non{-}trivially valued}
\NormalTok{field k, and that ∂ X = \{x\} or ∂ X = ;.}
\NormalTok{                                                             −∞}
\NormalTok{    If char(k) = p \textgreater{} 0 let k′ be the completion of k p , and otherwise let k′ = k. Let X ′ =}
\NormalTok{X k′ and let π = πk′ /k : X ′ → X , let ι : X red}
\NormalTok{                                              ′}
\NormalTok{                                                  ,→ X ′ be the inclusion of the underlying reduced}
\NormalTok{space, and let ν: Xe ′ → X red′}
\NormalTok{                                 be the normalization. Since π and ν are homeomorphisms,}
\NormalTok{the space X red is again irreducible, so Xe ′ is connected by [Duc09, Proposition 5.16]. Since}
\NormalTok{              ′}

\NormalTok{x is not a rig{-}point of X , there exists a unique point x ′ ∈ Xe ′ lying over x by Lemmas 3.7,}
\NormalTok{3.8 and 3.9. We must show that u◦π◦ι ◦ν ≡ −∞ on Xe ′ . Since u◦π◦ι ◦ν is subharmonic}
\NormalTok{by Proposition 3.5(7,8), we may replace X by Xe ′ to assume that X is rig{-}smooth and}
\NormalTok{connected (and still that ∂ X = ; or ∂ X = \{x\} using the quoted lemmas again).}
\NormalTok{    If ∂ X = ; then X is smooth and connected, so u is either identically −∞ or it is sub{-}}
\NormalTok{harmonic in Thuillier’s sense. The latter case is impossible when u(x) = −∞ by [Thu05,}
\NormalTok{Proposition 3.4.10].}
\NormalTok{    Suppose then that X is rig{-}smooth and connected, and that ∂ X = \{x\} and u(x) = −∞.}
\NormalTok{By Lemma 3.11, there is an open neighborhood U of x on which u ≡ −∞. Since X}
\NormalTok{is connected, every connected component of Int(X ) contains x in its closure. Thus U}
\NormalTok{intersects every connected component C of Int(X ). This implies that u takes the value}
\NormalTok{−∞ on a non{-}rig point of C, so the previous paragraph shows that u ≡ −∞ on C. It}
\NormalTok{follows that u ≡ −∞ on X .                                                                           □}
\NormalTok{3.12. Now we turn to harmonic functions. Let X be a smooth, strictly k{-}analytic curve}
\NormalTok{with k non{-}trivially valued and let h: X → R be a function. We claim that h is harmonic}
\NormalTok{in the sense of (3.2) if and only if h and −h are subharmonic (in our sense or Thuillier’s—}
\NormalTok{see (3.4.2)). This is implicit in Thuillier’s thesis, but it requires justification. Suppose then}
\NormalTok{that h: X → R is harmonic. Then ±h are subharmonic by [Thu05, Proposition 2.3.13 and}
\NormalTok{Corollaire 3.1.12]. Conversely, suppose that ±h are subharmonic. This implies that h is}
\NormalTok{continuous. Let U ⊂ X be a strictly affinoid domain. By [Thu05, Lemme 2.3.9], there}
\NormalTok{exists a unique harmonic function h′ : U → R that coincides with h on ∂ U. By definition}
\NormalTok{20                                 W. GUBLER AND J. RABINOFF}

\NormalTok{of subharmonicity, we have h ≤ h′ on U. By the same argument as applied to −h, we see}
\NormalTok{that h = h′ on U, so that h is harmonic on U. As U was arbitrary, h is harmonic.}
\NormalTok{  By 3.12, the following definition recovers Thuillier’s definition of harmonicity on a}
\NormalTok{smooth, strictly k{-}analytic curve, and even on a strictly affinoid rig{-}smooth curve us{-}}
\NormalTok{ing [Thu05, Corollaire 2.3.14].}
\NormalTok{Definition 3.13 (Harmonic functions). Let X be a k{-}analytic curve for a possibly trivially}
\NormalTok{valued non{-}Archimedean field k. A function h: X → R is harmonic if h and −h are both}
\NormalTok{subharmonic.}
\NormalTok{   The advantage of this definition is that the following basic properties of harmonic func{-}}
\NormalTok{tions follow immediately from the corresponding properties of subharmonic functions.}
\NormalTok{Proposition 3.14. Let X be a k{-}analytic curve and let h: X → R be a continuous function.}
\NormalTok{     (1) (Sheafiness) The harmonic functions form a sheaf on X .}
\NormalTok{     (2) If f ∈ Γ (X , OX× ) then log | f | is harmonic.}
\NormalTok{     (3) The harmonic functions on X form a real vector space.}
\NormalTok{     (4) (Limits) If h is locally a uniform limit of harmonic functions, then it is harmonic.}
\NormalTok{     (5) (Maximum Principle) If h is harmonic and h attains a local extremum at x ∈ Int(X ),}
\NormalTok{         then h is constant in a neighborhood of x.}
\NormalTok{     (6) (Extension of Scalars) Let k′ /k be an analytic extension field. Then h is harmonic}
\NormalTok{         if and only if h ◦ πk′ /k : X k′ → R is harmonic.}
\NormalTok{     (7) (Functoriality) Let Y be a k{-}analytic curve and let f : Y → X be a morphism. If h is}
\NormalTok{         harmonic then h ◦ f is harmonic, and the converse holds if f is finite and surjective.}
\NormalTok{     (8) (Boundaries) If h is harmonic on Int(X ), then it is harmonic on X .}
\NormalTok{   We will later see in Theorem 7.6 that (4) holds even for pointwise limits. The fol{-}}
\NormalTok{lowing proposition asserts that Definition 3.13 is equivalent to the notion of harmonicity}
\NormalTok{in [GJR21, Definition 7.2]. We assume that k is non{-}trivially valued as in ibid.}
\NormalTok{Proposition 3.15. Let X be a strictly k{-}analytic curve and let h: X → R be a continuous}
\NormalTok{function. The function h is harmonic if and only if it is R{-}PL and the line bundle Lh (x) ∈}
\NormalTok{Pic(red(X , x))R defined in (2.17) is numerically trivial for every point x ∈ X .}
\NormalTok{Proof. This is [GJR21, Proposition 15.7] in the rig{-}smooth case; we reduce to this case}
\NormalTok{                                                                         −∞}
\NormalTok{in the usual way. If char(k) = p \textgreater{} 0 let k′ be the completion of k p , and otherwise let}
\NormalTok{k′ = k. Let X ′ = X ⊗b k k′ and let π: X ′ → X be the structure morphism. Let ι : X red}
\NormalTok{                                                                                      ′}
\NormalTok{                                                                                        ,→ X ′}
\NormalTok{be the underlying reduced space and let ν: Xe ′ → X red  ′}
\NormalTok{                                                            be the normalization, so that Xe ′ is}
\NormalTok{rig{-}smooth. Let f : Xe ′ → X denote the composition π ◦ ι ◦ ν.}
\NormalTok{   Suppose that h is harmonic. Then h ◦ f is harmonic by Proposition 3.14(6,7), so h ◦ f}
\NormalTok{is R{-}PL and Lh (x ′ ) is numerically trivial for all x ′ ∈ Xe ′ by the rig{-}smooth case. Then}
\NormalTok{h is R{-}PL by Lemma 3.9(6) and [GJR21, Lemma 5.3(3), Corollary 5.11], and Lh (x) is}
\NormalTok{numerically trivial for all x ∈ X by [GJR21, Proposition 7.6(1′ , 2′ )]. Conversely, suppose}
\NormalTok{that h is R{-}PL and that Lh (x) is numerically trivial for all x ∈ X . The same holds for}
\NormalTok{h ◦ f by [GJR21, Proposition 7.6(1′ , 2′ )], so h ◦ f is harmonic by the rig{-}smooth case, and}
\NormalTok{hence h is harmonic by Proposition 3.14(6,7).                                                  □}
\NormalTok{             CLASSICALLY PSH AND PLURIHARMONIC FUNCTIONS ON BERKOVICH SPACES               21}

\NormalTok{Remark 3.16. We still assume that X is a strictly k{-}analytic curve over a non{-}trivially}
\NormalTok{valued field k and that h: X → R is continuous. Proposition 3.15 and Proposition 3.14(6)}
\NormalTok{imply in particular that if h◦πk′ /k is harmonic for some analytic extension field k′ /k, then}
\NormalTok{h◦πk′′ /k is R{-}PL for any analytic extension k′′ /k, which is by no means obvious. Likewise,}
\NormalTok{Proposition 3.14(7) implies that if f : Y → X is finite and surjective and h◦ f is harmonic,}
\NormalTok{then h is R{-}PL, and Proposition 3.14(8) implies that if h|Int(X ) is harmonic then h is R{-}PL.}

\NormalTok{                     4. CLASSICALLY PLURISUBHARMONIC FUNCTIONS}
\NormalTok{  In this section we define the notion of a classically plurisubharmonic function. The}
\NormalTok{definition follows Definition 3.4.}
\NormalTok{Definition 4.1 (Classically Plurisubharmonic Functions). Let X be a k{-}analytic space and}
\NormalTok{let u: X → R ∪ \{−∞\} be an upper semicontinuous function. We say that u is classically}
\NormalTok{plurisubharmonic or classically psh provided that, for every analytic extension field k′ /k,}
\NormalTok{every k′ {-}analytic curve Y , and every morphism f : Y → X k′ , the function u◦πk′ /k ◦ f : Y →}
\NormalTok{R ∪ \{−∞\} is subharmonic on Y in the sense of Definition 3.4.}
\NormalTok{4.1.1. Any class of psh functions should respect extension of scalars and should be stable}
\NormalTok{under pullback, so the conditions in Definition 4.1 are necessary. In complex analysis,}
\NormalTok{this condition (without the extension of scalars) is the definition of plurisubharmonicity,}
\NormalTok{which is why we call the notion “classically psh”.}
\NormalTok{4.1.2. If X is a k{-}analytic curve, then a function u: X → R∪\{−∞\} is classically psh if and}
\NormalTok{only if it is subharmonic. Indeed, if u is classically psh then taking k′ = k and Y = X in}
\NormalTok{Definition 4.1 shows that u is subharmonic, and if u is subharmonic then it is classically}
\NormalTok{psh by Proposition 3.5(7,8).}
\NormalTok{4.1.3. In Definition 4.1, one can assume that Y is a connected, smooth, strictly k′ {-}analytic}
\NormalTok{curve: this follows from Definition 3.4.}
\NormalTok{  Classically psh functions satisfy the basic properties of subharmonic functions on curves.}
\NormalTok{Proposition 4.2. Let X be a k{-}analytic space.}
\NormalTok{   (1) (Sheafiness) The classically psh functions form a sheaf on X .}
\NormalTok{   (2) If f ∈ Γ (X , OX ) then log | f |: X → R ∪ \{−∞\} is classically psh.}
\NormalTok{   (3) (Cone Property) If u1 , u2 : X → R∪\{−∞\} are classically psh and λ1 , λ2 ∈ R≥0 then}
\NormalTok{       λ1 u1 + λ2 u2 and max\{u1 , u2 \} are classically psh.}
\NormalTok{   (4) (Limits) The infimum of a decreasing net of classically psh functions is classically}
\NormalTok{       psh. For arbitrary pointwise limits, we refer to Theorem 6.5.}
\NormalTok{   (5) (Extension of Scalars) Let k′ /k be an analytic extension field. A function u: X →}
\NormalTok{       R ∪ \{−∞\} is classically psh if and only if u ◦ πk′ /k : X k′ → R ∪ \{−∞\} is classically}
\NormalTok{       psh.}
\NormalTok{   (6) (Functoriality) Let Y be a k{-}analytic space and let f : Y → X be a morphism. If}
\NormalTok{       u: X → R ∪ \{−∞\} is classically psh then u ◦ f is classically psh, and the converse}
\NormalTok{       holds if f is finite and surjective.}
\NormalTok{Proof. Properties (1,3,4) follow immediately from Proposition 3.5(1,4,5). For Proposi{-}}
\NormalTok{tion 4.2(2), we consider an analytic function f on X . As | f | is continuous, we deduce}
\NormalTok{22                                  W. GUBLER AND J. RABINOFF}

\NormalTok{that log | f | is a continuous function with values in R ∪ \{−∞\} and hence it is upper semi{-}}
\NormalTok{continuous. Then Proposition 4.2(2) follows from Proposition 3.5(2).}
\NormalTok{    To prove Proposition 4.2(5), it is clear from the definitions that if u is classically psh,}
\NormalTok{then u ◦ πk′ /k is classically psh for every analytic extension field k′ /k. To prove the con{-}}
\NormalTok{verse, we assume that k′ /k is an analytic extension field and that u′ := u ◦ πk′ /k is a}
\NormalTok{classically psh function on X k′ . Since the morphism πk′ /k is closed and surjective 2.9,}
\NormalTok{it is clear that u is upper semicontinuous. Let k′′ /k be an analytic extension field and}
\NormalTok{let f : Y → X k′′ be a morphism over k′′ from a k′′ {-}analytic curve Y . We have to show}
\NormalTok{that u ◦ πk′′ /k ◦ f is subharmonic. As in the proof of Proposition 3.5(7), we consider}
\NormalTok{a simultaneous analytic extension field K of k′ and k′′ . It follows from the above that}
\NormalTok{u ◦ πK/k = u′ ◦ πK/k′ is classically psh and hence u ◦ πK/k ◦ f K is subharmonic. We con{-}}
\NormalTok{clude from the commutative diagram (3.10.1) that u ◦ πK/k ◦ f K = u ◦ πk′′ /k ◦ f ◦ πK/k′′ is}
\NormalTok{subharmonic and hence u ◦ πk′′ /k ◦ f is subharmonic by Proposition 3.5(7).}
\NormalTok{    Likewise, the proof of Proposition 4.2(6) follows the same lines as the proof of Propo{-}}
\NormalTok{sition 3.5(8). Again, it follows from the definitions if u is a classically psh function on X ,}
\NormalTok{then u ◦ f is a classically psh function on Y . Conversely, we assume that u ◦ f is a clas{-}}
\NormalTok{sically psh function for a finite surjective morphism f : Y → X . Since finite morphisms}
\NormalTok{are proper and hence closed, we deduce from upper semicontinuity of u ◦ f and from}
\NormalTok{surjectivity of f that u is upper semicontinuous. It remains to show that u ◦ πk′ /k ◦ g is}
\NormalTok{subharmonic for any analytic extension field k′ /k and any morphism g : Z → X k′ from a}
\NormalTok{connected k′ {-}analytic curve Z. Setting W := Z ×X k′ Yk′ , we get a morphism g ′ : W → Yk′}
\NormalTok{from f by base change with respect to f ′ : Yk′ → X k′ as in (3.10.2). Since u ◦ f is classi{-}}
\NormalTok{cally psh, it follows from Proposition 4.2(5) that u ◦ πk′ /k ◦ f ′ = u ◦ f ◦ πk′ /k is classically}
\NormalTok{psh. The morphism f is finite and surjective, hence the same holds for f ′ . We con{-}}
\NormalTok{clude that the base change f ′′ : W → Z of the morphism f ′ is finite and surjective. The}
\NormalTok{k′ {-}analytic space W has irreducible components of dimension 1 and maybe isolated irre{-}}
\NormalTok{ducible components of dimension 0. Removing the latter, we obtain a k′ {-}analytic curve}
\NormalTok{W ′ as an open and closed subset of W . Since Z is a connected k′ {-}analytic curve and f ′′}
\NormalTok{is a finite surjective map from W onto Z, we get an induced finite surjective morphism}
\NormalTok{W ′ → Z. Indeed, as W ′ is also Zariski{-}closed in W , we deduce from finiteness of f ′′ that}
\NormalTok{f ′′ (W ′ ) is a Zariski{-}closed subset in Z of dimension 1 and hence f ′′ (W ′ ) = Z as Z is a}
\NormalTok{connected k′ {-}analytic curve. Since the inclusion W ′ → W is obviously finite, the com{-}}
\NormalTok{position W ′ → Z is also finite. In the following, we replace W by W ′ , which preserves}
\NormalTok{the commutativity of the diagram (3.10.2) and has the advantage that W is a k′ {-}analytic}
\NormalTok{curve. It follows that u ◦ πk′ /k ◦ f ′ ◦ g ′ = u ◦ πk′ /k ◦ g ◦ f ′′ is subharmonic on W . The}
\NormalTok{morphism of k′ {-}analytic curves f ′′ : W → Z is finite and surjective, so we deduce from}
\NormalTok{Proposition 3.5(8) that u ◦ πk′ /k ◦ g is subharmonic.                                           □}

\NormalTok{  In the paper [GR25], we have studied semipositive R{-}PL functions on a good strictly}
\NormalTok{k{-}analytic space over a non{-}trivially valued non{-}archimedean field k. This turns out to}
\NormalTok{be the same as an R{-}PL function that is classically psh, as we now show.}

\NormalTok{Theorem 4.3. Suppose that the valuation on k is nontrivial. Let X be a good, strictly k{-}}
\NormalTok{analytic space and let h: X → R be an R{-}PL function. Then h is semipositive if and only if h}
\NormalTok{is classically psh.}
\NormalTok{              CLASSICALLY PSH AND PLURIHARMONIC FUNCTIONS ON BERKOVICH SPACES                  23}

\NormalTok{Proof. For x ∈ X , an R{-}PL function h: X → R induces a residue line bundle Lh (x) ∈}
\NormalTok{Pic(red(X , x))R on the reduction red(X , x) of the germ (X , x): see 2.17. Recall from}
\NormalTok{[GJR21, Definition 7.2] that h is semipositive at x ∈ X if and only if the associated residue}
\end{Highlighting}
\end{Shaded}

\subsection{GR, newline lines
1832--1859}\label{gr-newline-lines-18321859}

\begin{Shaded}
\begin{Highlighting}[]
\NormalTok{  In this section, we consider a k{-}analytic space X over any non{-}Archimedean field k.}
\NormalTok{Following the discussion of harmonic functions on curves in 3.12, we generalize this}
\NormalTok{notion as follows to higher dimensions.}
\NormalTok{Definition 7.1. A function h: X → R is called pluriharmonic if h and −h are classically}
\NormalTok{psh functions (see Definition 4.1).}
\NormalTok{Proposition 7.2. Pluriharmonic functions are continuous and have the properties:}
\NormalTok{     (1) (Sheafiness) The pluriharmonic functions form a sheaf on X .}
\NormalTok{     (2) If f ∈ Γ (X , OX× ) then log | f | is pluriharmonic.}
\NormalTok{     (3) (Vector Space) The pluriharmonic functions on X form a real vector space which}
\NormalTok{         we denote by PH(X ).}
\NormalTok{     (4) (Limits) If a function h: X → R is locally a uniform limit of pluriharmonic functions,}
\NormalTok{         then h is pluriharmonic.}
\NormalTok{             CLASSICALLY PSH AND PLURIHARMONIC FUNCTIONS ON BERKOVICH SPACES              35}

\NormalTok{   (5) (Maximum Principle) If h ∈ PH(X ) and h attains a local extremum at x ∈ Int(X ),}
\NormalTok{       then h is constant in a neighborhood of x.}
\NormalTok{   (6) (Extension of Scalars) Let k′ /k be an analytic extension field. Then h pluriharmonic}
\NormalTok{       if and only if h ◦ πk′ /k : X k′ → R is pluriharmonic.}
\NormalTok{   (7) (Functoriality) Let Y be a k{-}analytic space and let f : Y → X be a morphism. If}
\NormalTok{       h ∈ PH(X ) then h ◦ f ∈ PH(Y ), and the converse holds if f is finite and surjective.}
\NormalTok{Proof. It is clear that a pluriharmonic function h is continuous as h and −h are upper}
\NormalTok{semicontinuous. Properties (1), (2), (3), (6) and (7) follow immediately from the corre{-}}
\NormalTok{sponding properties of classically psh functions given in Proposition 4.2. In (4), we note}
\NormalTok{that a uniform limit of pluriharmonic functions can be easily written as a decreasing}
\NormalTok{(resp. increasing) limit of pluriharmonic functions and hence (4) follows from Proposi{-}}
\NormalTok{tion 4.2(4). The maximum principle (5) follows from the maximum principle for classi{-}}
\NormalTok{cally psh functions shown in Theorem 6.1.                                               □}
\NormalTok{   From now on, we denote by PH(X ) the space of pluriharmonic functions on X . The}
\end{Highlighting}
\end{Shaded}

\subsection{Wanner, newline lines
168--202}\label{wanner-newline-lines-168202}

\begin{Shaded}
\begin{Highlighting}[]
\NormalTok{Terminology. In this paper, let K be an algebraically closed field endowed with a com{-}}
\NormalTok{plete, non{-}archimedean, non{-}trivial absolute value | |. We write val(K × ) := − log |K × |}
\NormalTok{and R := \{x ∈ K | |x| ≤ 1\} with maximal ideal m ⊂ R and residue field K   e := R/m. By}
\NormalTok{suitably normalizing the absolute value | |, we may assume that Z, and so Q, is contained}
\NormalTok{in val(K × ). A variety over K is an irreducible separated reduced scheme of finite type}
\NormalTok{over K and a curve is a 1{-}dimensional variety over K.}

\NormalTok{Acknowledgements. The author would like to thank Walter Gubler for very carefully}
\NormalTok{reading drafts of this work and for the helpful discussions. The author is also grateful}
\NormalTok{to Philipp Jell for providing a simplification of the proof of Proposition 4.4 and further}
\NormalTok{useful comments. Finally, the author would like to thank the referee for their very precise}
\NormalTok{report and helpful suggestions.}


\NormalTok{     2. Preliminaries on Thuillier’s notion of subharmonic functions}
\NormalTok{   Let X be a smooth algebraic curve over K and X an be its analytification. In this}
\NormalTok{section, we define subharmonic functions in the sense of Thuillier on X an and give some}
\NormalTok{properties of these functions. Subharmonic functions are defined with the help of har{-}}
\NormalTok{monic functions as in the complex case. We therefore recall the definition of harmonic}
\NormalTok{functions on strictly affinoid domains from [Thu05, §2.3]. For this definition, one needs}
\NormalTok{the skeleton S(Y) corresponding to a strictly semistable formal model Y of a strictly}
\NormalTok{affinoid domain Y , which has the structure of a finite graph.}
\NormalTok{   Furthermore, we study lisse (subharmonic) functions, which are crucial for the proofs}
\NormalTok{in Section 4. These functions are Thuillier’s smooth functions, but note that they are}
\NormalTok{totally different from those smooth functions introduced by Chambert{-}Loir and Ducros,}
\NormalTok{which we define in Section 3. Hence, we use the term lisse instead of smooth here.}

\NormalTok{Definition 2.1. A semistable R{-}curve is an integral admissible formal R{-}curve Y whose}
\NormalTok{special fiber has only ordinary double points as singularities. A semistable R{-}curve is}
\NormalTok{called strictly semistable if in addition the irreducible components of the special fiber}
\NormalTok{are smooth. Let Y be a strictly analytic domain in X an , i.e. a subset that has a locally}
\NormalTok{finite covering by strictly affinoid domains, then we say that Y is a (strictly) semistable}
\NormalTok{formal model of Y if Y is a (strictly) semistable proper formal R{-}curve with generic fiber}
\NormalTok{Y η isomorphic to Y .}
\end{Highlighting}
\end{Shaded}

\subsection{Wanner, newline lines
447--519}\label{wanner-newline-lines-447519}

\begin{Shaded}
\begin{Highlighting}[]
\NormalTok{Definition 3.11. Let X be an algebraic variety over K and W be an open subset of}
\NormalTok{X an . An element of Ap,q   X (W ) is given by a family (Vi , ϕUi , αi )i∈I such that}
\NormalTok{                                                                             S}
\NormalTok{       i) for all i ∈ I the pair (Vi , ϕUi ) is a tropical chart and W = i∈I Vi ;}
\NormalTok{      ii) for all i ∈ I we have αi ∈ Ap,qTrop(Ui ) (tropUi (Vi ));}
\NormalTok{     iii) for all i, j ∈ I the restrictions αi |Vi ∩Vj = αj |Vi ∩Vj agree;}
\NormalTok{   If α′ is another differential form of bidegree (p, q) on W given by a triple (Vj′ , ϕUj′ , α′j )j∈J ,}
\NormalTok{then we consider α and α′ as the same differential form if and only if}
\NormalTok{                                         αi |Vi ∩Vj′ = α′j |Vi ∩Vj′}
\NormalTok{for every i ∈ I and j ∈ J.}
\NormalTok{   Then W 7→ Ap,q                              an                          p,q}
\NormalTok{                   X (W ) defines a sheaf on X , which we denote by AX and we write}
\NormalTok{Ap,q}
\NormalTok{  X,c (W ) for the sections with compact support in W . If the space of definition is clear,}
\NormalTok{                                                                                  p,q}
\NormalTok{we often just use the notations Ap,q and Ap,q c . By Theorem 3.7, we have AX = 0 if}
\NormalTok{max(p, q) \textgreater{} n = dim(X).}
\NormalTok{   The differentials d′ and d′′ and the wedge product carry Rover. Moreover, for every}
\NormalTok{open subset W of X an there is a non{-}trivial integration map : An,n   c (W ) → R which is}
\NormalTok{compatible with pullback.}
\NormalTok{   Note that smooth differential forms of bidegree (0, 0) are well{-}defined continuous func{-}}
\NormalTok{tions.}
\NormalTok{Definition 3.12. Let W be an open subset of X an . A function f : W → R is called}
\NormalTok{smooth if f ∈ A0,0 (W ). We use the notations C ∞ (W ) := A0,0 (W ) and Cc∞ (W ) :=}
\NormalTok{A0,0}
\NormalTok{ c (W ).}

\NormalTok{   Next to (0, 0){-}forms, we also work a lot with (n, n){-}forms. It is important to know the}
\NormalTok{following fact about these forms.}
\NormalTok{Proposition 3.13. Let W be an open subset of X an and   R α ∈ A (W}
\NormalTok{                                                                  n,n}
\NormalTok{                                                                      R ). Then there}
\NormalTok{is a unique signed Radon measure µα on W such that W f dµα = W f α for every}
\NormalTok{f ∈ Cc∞ (W ). If α has compact support on W , so has µα and |µα |(W ) \textless{} ∞.}
\NormalTok{Proof. See [Gub16, Proposition 6.8].                                                                }
\NormalTok{Definition 3.14. Let U be an open subset of Rr . A superform α ∈ Ap,p(U ) is called}
\NormalTok{strongly positive if there exist finitely many superforms αj,s of type (0, 1) and non{-}}
\NormalTok{negative smooth functions fs on U such that}
\NormalTok{                          X}
\NormalTok{                      α=      fs α1,s ∧ J(α1,s ) ∧ . . . ∧ αp,s ∧ J(αp,s ).}
\NormalTok{                                s}
\NormalTok{10                                       V. WANNER}


\NormalTok{   Let C be a polyhedral complex and Ω an open subset of |C |. A superform α ∈ Ap,p (Ω)}
\NormalTok{is called strongly positive if there is a polyhedral decomposition of C such that the}
\NormalTok{restriction of α to relint(σ) ∩ Ω is strongly positive for every polyhedron σ ∈ C .}
\NormalTok{   For an open subset W of X an , a smooth form ω ∈ Ap,p (W ) is called strongly positive if}
\NormalTok{for every point x in W there is a tropical chart (V, ϕU ) with x ∈ V such that ω = α◦tropU}
\NormalTok{on V for a strongly positive form α ∈ Ap,p Trop(U ) (tropU (V )).}
\NormalTok{   Note that for forms of type (0, 0), (1, 1), (n − 1, n − 1) and (n, n) the notion of}
\NormalTok{strongly positivity defined here coincides with the other positivity notions from [CD12,}
\NormalTok{§5.1]. Thus, we just say that a smooth form ω in Ap,p (W ) is positive if it is of one of}
\NormalTok{these types and it is strongly positive.}
\NormalTok{   A smooth function f is (strongly) positive as a form if and only if f ≥ 0.}
\NormalTok{Definition 3.15. Let W be an open subset of X an . We call a function f : W → [−∞, ∞]}
\NormalTok{locally integrable if f is integrable}
\NormalTok{                             R          withR respect to every measure µα associated to a form}
\NormalTok{α ∈ An,n}
\NormalTok{       c  (W  ). We  write     W  f ∧ α := W f dµα .}
\NormalTok{   Then for every locally integrable (e.g. continuous) function f : W → [−∞, ∞] one can}
\NormalTok{define a current in the sense of [CD12, §4.2] by}
\NormalTok{                                                             Z}
\NormalTok{                        ′ ′′        n−1,n−1}
\NormalTok{                       d d [f ] : Ac        (W ) → R, α 7→      f ∧ d′ d′′ α.}
\NormalTok{                                                          W}
\NormalTok{Definition 3.16. Let W be an open subset of            an}
\NormalTok{                                                      X . A locally integrable  function}
\NormalTok{f : W → [−∞, ∞) is called psh if f is upper semi{-}continuous and d′ d′′ [f ] is a positive}
\NormalTok{current, i.e. d′ d′′ [f ](ω) ≥ 0 for all positive forms ω ∈ An−1,n−1}
\NormalTok{                                                             c       (W ).}
\NormalTok{  Note that we do not require that a psh function has to be continuous, contrary to}
\NormalTok{[CD12, Définition 5.5.1].}
\end{Highlighting}
\end{Shaded}

\subsection{Wanner, newline lines
680--742}\label{wanner-newline-lines-680742}

\begin{Shaded}
\begin{Highlighting}[]
\NormalTok{                              XZ                              Z}
\NormalTok{                        =                 ηj g dMA(f ) =              g dMA(f ).}
\NormalTok{                              j∈J    Vj                           W}

\NormalTok{For assertion ii), we assume that f is locally psh{-}approximable, i.e. we can cover W}
\NormalTok{by open subsets (Vi )i∈I such that f is the uniform limit f = limk→∞ fik of smooth}
\NormalTok{psh functions fik on Vi . As above let (ηj )j∈J be a partition of unity subordinated to}
\NormalTok{(Vi )i∈I . For every non{-}negative function g ∈ Cc∞ (W ), the smooth function ηj g is also}
\NormalTok{non{-}negative as ηj ≥ 0 and has compact support on Vj . From the calculations above,}
\NormalTok{we get}
\NormalTok{                                         X     Z}
\NormalTok{                        d′ d′′ [f ](g) =   lim   ηj g ∧ d′ d′′ fjk ≥ 0.}
\NormalTok{                                                       k→∞ Vj}
\NormalTok{                                                 j∈J}

\NormalTok{This proves that f is psh on W .                                                                                    }
\NormalTok{Lemma 3.24. Let X be a curve, f : W → R be a globally psh{-}approximable function,}
\NormalTok{i.e. f is the uniform limit of smooth psh functions fi on W , and g ∈ Cc0 (W ), then}
\NormalTok{                           Z                     Z}
\NormalTok{                                g dMA(f ) = lim     g ∧ d′ d′′ fi .}
\NormalTok{                                    W                     i→∞ W}

\NormalTok{Proof. By [CD12, Proposition 3.3.5], there are smooth functions gk ∈ Cc∞ (W ) converging}
\NormalTok{uniformly to g. Then}
\NormalTok{                     Z                        Z}
\NormalTok{                          g dMA(f ) = lim        gk dMA(f )}
\NormalTok{                       W                 k→∞ W}
\NormalTok{                                                  Z}
\NormalTok{                                      = lim lim       gk ∧ d′ d′′ fi}
\NormalTok{                                         k→∞ i→∞ W}
\NormalTok{                                                  Z}
\NormalTok{                                      = lim lim       gk ∧ d′ d′′ fi}
\NormalTok{                                         i→∞ k→∞ W}
\NormalTok{                                              Z}
\NormalTok{                                      = lim      g ∧ d′ d′′ fi .}
\NormalTok{                                                        i→∞ W}


\NormalTok{                                                                           R           ′ ′′}
\NormalTok{                                                                                                              R}
\NormalTok{Note that we may change the order of the limits since                          W gk ∧ d d fi converges to         W g∧}
\NormalTok{d′ d′′ fi uniformly in i ∈ N.                                                                                       }

\NormalTok{   Model functions and their Monge{-}Ampère measures play an important role in the}
\NormalTok{proofs of the main theorems in Section 4. We therefore recall the definition of metrics,}
\NormalTok{formal metrics and model functions.}
\NormalTok{14                                      V. WANNER}


\NormalTok{Definition 3.25. Let X be a proper and normal variety over K and L be a line bundle}
\NormalTok{on X. A continuous metric k k on Lan associates to every section s ∈ Γ(U, L) on a Zariski}
\NormalTok{open subset U of X a continuous function ksk : U an → [0, ∞) such that kf · sk = |f | · ksk}
\NormalTok{holds for every f ∈ OX (U ) and ks(x)k = 0 if and only if s vanishes in x.}
\NormalTok{   We call a continuous metric k k on Lan smooth (resp. psh) if for every open subset U}
\NormalTok{of X and every invertible section s of L on U the function − log ksk is smooth (resp. psh)}
\NormalTok{on U an .}
\NormalTok{   A continuous metric k k on Lan is psh{-}approximable if there is a sequence of smooth}
\NormalTok{psh metrics k kk on Lan such that supx∈X an | log(ksx (x)k/ksx (x)kk )| converges to zero}
\NormalTok{for any local section sx of L that does not vanish in x. Clearly, this is independent of}
\NormalTok{the choice of sx .}
\NormalTok{Definition 3.26. Let X be a proper variety over K and let L be a line bundle on X.}
\NormalTok{A semistable formal model of (X, L) is a pair (X , L) consisting of a semistable formal}
\NormalTok{model X of X an and a line bundle L on X such that L |X an ≃ Lan . Note that we always}
\end{Highlighting}
\end{Shaded}

\subsection{Wanner, newline lines
991--1606}\label{wanner-newline-lines-9911606}

\begin{Shaded}
\begin{Highlighting}[]
\NormalTok{Lemma 4.6. Let W be an open subset of X an and f : W → R be a continuous function.}
\NormalTok{If f is not subharmonic, then there is a strictly semistable formal model (X , L) of (X, O X )}
\NormalTok{such that the model function g := − log k1kL satisfies the following properties:}
\NormalTok{      i) supp(g) is a connected strictly affinoid domain Y ⊂ W ;}
\NormalTok{     ii) g is strictly positive on Y \textbackslash{}∂Y ;}
\NormalTok{    iii)                                Z}
\NormalTok{                                               f ddc g \textless{} 0.}
\NormalTok{                                           W}
\NormalTok{Proof. By Corollary 2.13, thereR is a connected strictly affinoid domain Y and a type II}
\NormalTok{or III point x in Y \textbackslash{}∂Y with W f ddc gxY \textless{} 0. Since gxY is lisse, we can find a strictly}
\NormalTok{semistable formal model Y of Y and a piecewise affine function GYx on the corresponding}
\NormalTok{skeleton S(Y) such that gxY = GYx ◦τY on Y . Note that S(Y) is a closed connected subset}
\NormalTok{of X an with the structure of a finite metric graph. By blowing up Y, we may assume}
\NormalTok{that there is no edge in S(Y) whose endpoints are all contained in the boundary ∂Y . We}
\NormalTok{will explain in steps, why there is a strictly semistable formal model (X , L) such that}
\NormalTok{   COMPARISON OF TWO NOTIONS OF SUBHARMONICITY ON NON{-}ARCHIMEDEAN CURVES 19}


\NormalTok{the corresponding}
\NormalTok{               R model  function g = − log k1kL on X an satisfies g \textgreater{} 0 on Y \textbackslash{}∂Y , g = 0}
\NormalTok{                     c}
\NormalTok{elsewhere, and W f dd g \textless{} 0.}

\NormalTok{  0. Step: Strategy of the proof.}

\NormalTok{   We construct a val(K × ){-}rational function (see definition below) G on S(Y) with the}
\NormalTok{required properties}
\NormalTok{                                                    Z}
\NormalTok{(4.3)               G|S(Y)\textbackslash{}∂Y \textgreater{} 0, G|∂Y = 0 and           f ddc G \textless{} 0.}
\NormalTok{                                                           S(Y)}

\NormalTok{Then the function g := G ◦ τY is piecewise Q{-}linear in the sense of [GK15] on Y . Setting}
\NormalTok{g ≡ 0 on X an \textbackslash{}Y , we get a well{-}defined continuous non{-}negative function g on X an and}
\NormalTok{we have a G{-}covering of X an on which g is piecewise Q{-}linear. Thus, there is a semistable}
\NormalTok{                                         1}
\NormalTok{formal model (X , L) such that g = − m     log k1kL on X an by [GK14, Proposition 8.11\&}
\NormalTok{8.13]. Note that we always may assume X to be strictly semistable by the semistable}
\NormalTok{reduction Theorem [BL93, Ch. 7] since K is algebraically closed. Then the formal model}
\NormalTok{(X , L⊗m ) of (X, O X ) gives the claim.}
\NormalTok{   Before we start with the construction, note that a val(K × ){-}rational function is a}
\NormalTok{piecewise affine function on S(Y) (we refine the vertex set such that G is affine on every}
\NormalTok{edge) such that the following properties are satisfied}
\NormalTok{   (a) ddc G is only supported on points of type II.}
\NormalTok{   (b) G has values in val(K × ) at every vertex of S(Y).}
\NormalTok{   (c) G has rational slopes.}


\NormalTok{    1. Step: Replace GYx by a piecewise affine function G′ such that (4.3) is still}
\NormalTok{             satisfied and (a) additionally holds.}

\NormalTok{   If x is of type II, the support of ddc GYx consists only of type II points. If x is of type III,}
\NormalTok{we use that the points of type II are dense in X an . Let e = [y1 , y2 ] be the edge of S(Y)}
\NormalTok{having x in its interior. Let xn be a sequence of type II points in e converging to x with}
\NormalTok{respect to the skeletal metric. Consider the piecewise affine functions Gn on S(Y) that}
\NormalTok{are given by the affine function on [y1 , xn ] (resp. on [xn , y2 ]) connecting the points GYx (y1 )}
\NormalTok{and GYx (xn ) (resp. GYx (xn ) and GYx (y2 )) and Gn ≡ GYx on S(Y)\textbackslash{}e. It is easy to see that}
\NormalTok{the slopes converge to the ones of GYx , i.e. ddc Gn → ddc GYx for n → ∞. Furthermore, f}
\NormalTok{is continuous, so we can find n Rbig enough such that |f (xn )ddc Gn (xn ) − f (x)ddc GYx (x)|}
\NormalTok{is so small that we still have S(Y) f ddc Gn \textless{} 0 (cf. [Thu05, Proposition 3.3.4]). Set}
\NormalTok{G′ := Gn for such an n.}

\NormalTok{    2. Step: Replace G′ by a piecewise affine function G′′ such that (4.3) and (a)}
\NormalTok{             are still satisfied and (b) additionally holds.}
\NormalTok{   Due to normalizing the absolute value | |, we assume that Q is contained, and so dense,}
\NormalTok{in val(K × ). Let z be a vertex of S(Y). If G′ (z) is not in Q and (an )n is a sequence of}
\NormalTok{rational points converging to G′ (z), then the slopes of the piecewise affine functions G′n}
\NormalTok{on S(Y) resulting by replacing G′ (z) by an converge to the slopes of G′ . Thus, we can}
\NormalTok{choose an n ∈ N such that G′′ := G′n still satisfies (4.3) and takes only values in Q at}
\NormalTok{every vertex of S(Y). We choose these values such that G′′ is still strictly positive on}
\NormalTok{S(Y)\textbackslash{}∂Y .}
\NormalTok{20                                             V. WANNER}

\NormalTok{      3. Step: Replace G′′ by a piecewise affine function G such that (4.3), (a) and (b)}
\NormalTok{               are still satisfied and (c) additionally holds.}

\NormalTok{    Now, consider an edge e = [y1 , y2 ] of S(Y). Denote by d(y1 , y2 ) the distance between}
\NormalTok{these two points with respect to the skeletal metric. If d(y1 , y2 ) ∈ Q, we are done. If}
\NormalTok{not, we can find points y1′ and y2′ of type II in e with distance d(y1′ , y2′ ) ∈ Q arbitrary}
\NormalTok{close to d(y1 , y2 ). These points are chosen so that we can decompose the edge e into}
\NormalTok{[y1 , y1′ ] ∪ [y1′ , y2′ ] ∪ [y2′ , y2 ]. Let G be the piecewise affine function on e defined by the}
\NormalTok{data G ≡ G′′ (y1 ) on [y1 , y1′ ], G ≡ G′′ (y2 ) on [y2′ , y2 ] and G|[y1′ ,y2′ ] is affine (see Figure 1).}
\NormalTok{Since G(yi′ ) = G′′ (yi ) ∈ Q, the constructed function G on e has rational slopes. Choose}
\NormalTok{yi′ with d(y1′ , y2′ ) close enough to d(y1 , y2 ) such that (4.3) is still satisfied. Note that this}
\NormalTok{is possible since f is continuous (cf. [Thu05, Proposition 3.3.4]).}
\NormalTok{    We do this for all edges except the ones containing the boundary points ∂Y , where}
\NormalTok{we just move the other vertex. Then the function G on S(Y) has slopes in Q and is}
\NormalTok{consequently the required function.}


\NormalTok{                                                    G′′ (y1 )}
\NormalTok{                                                        |}
\NormalTok{              G′′ (y1 )}
\NormalTok{                  |}
\NormalTok{                                                                                         G′′ (y2 )}
\NormalTok{                                                                                             |}
\NormalTok{                          ỹ0                 ỹ1               y1′                y2′}
\NormalTok{                 •         •                   •       •        •                  •        •}
\NormalTok{                 y0                                    y1                                   y2}
\NormalTok{                                  ∈Q                                   ∈Q}

\NormalTok{                                    ẽ                                   e}

\NormalTok{                                                FIGURE 1}

                                                                                                        
\NormalTok{Theorem 4.7. Let W be an open subset of X an and f : W → R be a continuous function.}
\NormalTok{If f is psh, then f is subharmonic.}
\NormalTok{Proof. The proof is by contradiction.R We assume that there is a strictly semistable formal}
\NormalTok{model (X , L) of (X, O X ) such that W f ddc g \textless{} 0 for g := − log k1kL as in Lemma 4.6.}

\NormalTok{     0. Step: Strategy of the proof.}

\NormalTok{  We know that g = G ◦ τX on X an for the piecewise affine function G = g|S(X ) on}
\NormalTok{S(X ), the Monge{-}Ampère measure exists and ddc g = ddc G = MA(g) by Theorem 3.28.}
\NormalTok{  Assume that we have the following situation: We can write g = g + − g− on an open}
\NormalTok{subset V of W that contains the connected strictly affinoid domain Y = supp(g) such}
\NormalTok{that}
\NormalTok{     (a) g ± is the uniform limit of smooth psh functions gk± on V , i.e. g is globally ap{-}}
\NormalTok{         proximable on V ,}
\NormalTok{     (b) gk := gk+ − gk− ≥ 0 on V , and}
\NormalTok{  COMPARISON OF TWO NOTIONS OF SUBHARMONICITY ON NON{-}ARCHIMEDEAN CURVES 21}


\NormalTok{    (c) gk has compact support on V .}
\NormalTok{Moreover, we want to have a connected open subset V ′ of V with Y ⊂ V ′ , V ′ ⊂ V and}
\NormalTok{gk ∈ Cc∞ (V ′ ), and a continuous map η on V such that η ≡ 1 on V ′ and η has compact}
\NormalTok{support in V .}
\NormalTok{   In this situation, we have ηf ∈ Cc0 (V ), f = ηf on V ′ and by the definition of the}
\NormalTok{Monge{-}Ampère measure and Theorem 3.28}
\NormalTok{                            ddc g = MA(g) = MA(g + ) − MA(g− )}
\NormalTok{is a measure with finite support contained in Y ⊂ V ′ ⊂ V . Due to our assumption, this}
\NormalTok{implies}
\NormalTok{          Z                 Z                 Z                 Z}
\NormalTok{                                                           +}
\NormalTok{       0\textgreater{}     f dMA(g) =       ηf dMA(g) =      ηf dMA(g ) −       ηf dMA(g − ).}
\NormalTok{             W                   V                        V   V}
\NormalTok{Applying Lemma 3.24 to the right hand side and using gk ∈ Cc (V ′ ), we get}
\NormalTok{                                                           ∞}
\NormalTok{                   Z                            Z                           Z}
\NormalTok{             0\textgreater{}         f dMA(g) = lim              ηf ∧ d′ d′′ gk+ − lim       ηf ∧ d′ d′′ gk−}
\NormalTok{                    W                   k→∞ V                        k→∞    V}
\NormalTok{                                                Z}
\NormalTok{                                                           ′ ′′}
\NormalTok{                                     = lim          ηf ∧ d d gk}
\NormalTok{                                        k→∞ V}
\NormalTok{                                                Z}
\NormalTok{                                     = lim       f ∧ d′ d′′ gk}
\NormalTok{                                      k→∞ V ′}
\NormalTok{                                     = lim d′ d′′ [f |V ′ ](gk ).}
\NormalTok{                                      k→∞}

\NormalTok{We know that d′ d′′ [f ] is positive on W , and so it is on V ′ . Thus,}
\NormalTok{                                        lim d′ d′′ [f |V ′ ](gk ) ≥ 0,}
\NormalTok{                                       k→∞}
\NormalTok{and so we have a contradiction. Hence, f has to be subharmonic.}
\NormalTok{  We explain in several steps how to construct V , V ′ , η and the functions gk+ , gk− such}
\NormalTok{that one has the described situation.}

\NormalTok{  1. Step:       Show that the function g is globally approximable on W .}
\NormalTok{  The curve X is projective, so we may assume X to be projective as well and can}
\NormalTok{therefore find very ample line bundles L1 , L2 such that L = L1 ⊗ L−1}
\NormalTok{                                                                   2 . Thus, we can}
\NormalTok{write}
\NormalTok{                       g = log k1kL = − log ks1 kL1 + log ks2 kL2}
\NormalTok{on W for global sections s1 , s2 that coincide on the generic fiber. Since we may work on}
\NormalTok{every open subset of W containing the compact subset Y by Proposition 3.19, we may}
\NormalTok{assume that s1 and s2 do not vanish on W . Due to the very ampleness, we can find}
\NormalTok{closed embeddings ϕi : X → PnRi such that Li ≃ ϕ∗i OPni (1) and si = ϕ∗i xji for i = 1, 2.}
\NormalTok{                                                            R}
\NormalTok{Here, let x0 , . . . , xni be the coordinates of PnRi and ji ∈ \{0, . . . , ni \}. Without loss of}
\NormalTok{generality, we assume xj1 = xn1 and xj2 = xn2 . These closed embeddings induce closed}
\NormalTok{embeddings ϕi : X → PnKi into the toric varieties PnKi . Then}
\NormalTok{  g = − log ks1 kL1 + log ks2 kL2}
\NormalTok{    = ((ϕan ∗}
\NormalTok{         1 ) (− log |xn1 | +      max          log |xj |) − (ϕan ∗}
\NormalTok{                                                              2 ) (− log |xn2 | +      max          log |xj |))}
\NormalTok{                               j∈\{0,...,n1 \}                                        j∈\{0,...,n2 \}}

\NormalTok{on W . We approximate the functions φi := − log |xni | + maxj∈\{0,...,ni \} log |xj | for i = 1, 2}
\NormalTok{                                    ni ,an}
\NormalTok{by smooth convex functions on \{η ∈ PK      | |xni |η 6= 0\} as in [CD12, Proposition 6.3.2].}
\NormalTok{22                                                   V. WANNER}


\NormalTok{   For every k ∈ N\textgreater{}0 and n ∈ N, there is a smooth, convex function Mn, 1 on Rn+1 that}
\NormalTok{                                                                       k}
\NormalTok{is non{-}decreasing in every variable and has the following properties}
\NormalTok{       i) max(t0 , . . . , tn ) ≤ Mn, 1 (t0 , . . . , tn ) ≤ max(t0 , . . . , tn ) + k1 ;}
\NormalTok{                                         k}
\NormalTok{       ii) If tl + k2 ≤ maxj6=l tj , then Mn, 1 (t0 , . . . , tn ) = Mn−1, 1 (t0 , . . . , tbl , . . . , tn );}
\NormalTok{                                              k                            k}
\NormalTok{      iii) For all t ∈ R, we have Mn, 1 (t0 + t, . . . , tn + t) = Mn, 1 (t0 , . . . , tn ) + t.}
\NormalTok{                                                k                                   k}

\NormalTok{     Define for k ∈ N\textgreater{}0 and i ∈ \{1, 2\} the following function}
\NormalTok{                            φi,k := − log |xni | + Mni , 1 (log |x0 |, . . . , log |xni |)}
\NormalTok{                                                              k}

\NormalTok{    ni ,an                            ni ,an}
\NormalTok{on PK      For every point ζ in \{η ∈ PK}
\NormalTok{           .                                 | |xni |η 6= 0\} there is an open neighborhood}
\NormalTok{of ζ such that φi,k is smooth and psh on this neighborhood. Both are local properties}
\NormalTok{and si does not vanish on W , so (ϕan ∗}
\NormalTok{                                   i ) φi,k is a smooth psh function on W (see Remark}
\NormalTok{3.21) converging uniformly to − log ksi kLi on W . We use in the following the notations}
\NormalTok{g+ := − log ks1 kL1 and g− := − log ks2 kL2 . Furthermore, we set gk+ := (ϕan  ∗}
\NormalTok{                                                                            1 ) φ1,k and}
\NormalTok{gk− := (ϕan ∗}
\NormalTok{         2 ) φ2,k . Note that we have by construction}

\NormalTok{                                              g ± ≤ gk± ≤ gk± + 1/k.}

\NormalTok{     These functions do not necessarily satisfy (b) and (c), so we need to modify gk± .}

\NormalTok{      2. Step: Construct a suitable V and study the behavior of gk± outside of Y :}


\NormalTok{   At the end of Step 2, one can find an illustration of the construction in Figure 2. The}
\NormalTok{boundary of the strictly affinoid domain Y is a finite set of points of type II. By blowing}
\NormalTok{up our model X , we may assume that the points ∂Y are vertices in S(X ). Note that}
\NormalTok{we always may assume that an admissible blowing up is strictly semistable again by the}
\NormalTok{semistable reduction Theorem [BL85, Ch. 7]. Consider a point y ∈ ∂Y . Since y is of type}
\NormalTok{II, [BPR13, Corollary 4.27 \& Theorem 4.11] tells us that there are a strictly semistable}
\NormalTok{formal model X y of X and a star{-}shaped open neighborhood Ωy of y in S(X y ) such that}
\NormalTok{τX−1y (Ωy ) is an open neighborhood of y in W . Here, a star{-}shaped open neighborhood}
\NormalTok{Ωy of y in S(X y ) is a simply{-}connected open neighborhood of y in S(X y ) such that the}
\NormalTok{intersection of Ωy with any edge e in S(X y ) emanating from y is a half{-}open interval}
\NormalTok{Iy,e = [y, xe ) with endpoints y and xe of type II. By blowing up X and modifying Ωy , we}
\NormalTok{may assume that we can find this star{-}shaped open neighborhood Ωy in S(X ). We explain}
\NormalTok{how to do this. The model X has to be blown up such that every vertex of S(X y ) is a}
\NormalTok{vertex in the new skeleton S(X ). Then we can modify Ωy in the following way. Consider}
\NormalTok{an edge e of the new skeleton S(X ). Then the interior of e is either contained in an}
\NormalTok{edge ẽ of S(X y ) or lies in a connected component of X an \textbackslash{}S(X y ) isomorphic to an open}
\NormalTok{ball. In the first case, we shrink Iy,ẽ so that Iy,e is a half{-}open interval in e. Note that}
\NormalTok{τX−1 (Iy,e ) = τX−1y (Iy,e ), and so it is still contained in W . In the second case, we just add}
\NormalTok{a new half{-}open interval Iy,e to Ωy . Then τX−1 (Iy,e ) ⊂ τX−1y (y) ⊂ W . We do this blowing}
\NormalTok{ups and modifications for all boundary points. Moreover, we always choose Ωy such that}
\NormalTok{Ωy ⊂ W . Before we can construct V , we have to blow up X one more time. We find}
\NormalTok{this admissible formal blowing up X ′ of X such that Y = Y η for a formal open subset Y}
\NormalTok{of X ′ , which is possible by [BL93, Lemma 4.4]. Then we have Y = τX−1′ (S(X ′ ) ∩ Y ). As}
\NormalTok{described above, we modify Ωy such that it is a star{-}shaped open neighborhood of y in}
\NormalTok{S(X ′ ). Altogether, we have a strictly semistable formal model X ′ of X such that every}
\NormalTok{   COMPARISON OF TWO NOTIONS OF SUBHARMONICITY ON NON{-}ARCHIMEDEAN CURVES 23}


\NormalTok{boundary point y of Y has an open neighborhood τX−1′ (Ωy ) in W for an open star{-}shaped}
\NormalTok{neighborhood Ωy of y in S(X ′ ) and Y = τX−1′ (S(X ′ ) ∩ Y ).}
\NormalTok{   For every point z ∈ (S(X ′ ) ∩ Y )\textbackslash{}∂Y , we choose some star{-}shaped open neighborhood}
\NormalTok{Ωz of z in S(X ′ ) with Ωz ⊂ Y \textbackslash{}∂Y . Then τX−1′ (Ωz ) is automatically contained in W due}
\NormalTok{to Y = τX−1′ (S(X ′ ) ∩ Y ) and Y ⊂ W .}
\NormalTok{   We have constructed for every point in S(X ′ )∩Y an open neighborhood of it in S(X ′ ),}
\NormalTok{Ωy for y ∈ ∂Y and Ωz for z ∈    / ∂Y , and so these open subsets clearly cover our compact}
\NormalTok{subset S(X ′ ) ∩ Y . Thus, there is a finite subset Y0 of S(X ′ ) ∩ Y such that}
\NormalTok{                                                        [}
\NormalTok{                             Y = τX−1′ (S(X ′ ) ∩ Y ) ⊂   τX−1′ (Ωz ).}
\NormalTok{                                                               z∈Y0}

\NormalTok{By construction, the set of boundary points ∂Y is contained in Y0 . Furthermore, we}
\NormalTok{choose Y0 minimal, i.e. removing one open subset τX−1′ (Ωz ) from the covering would no}
\NormalTok{longer cover Y . Set}
\NormalTok{                                       [}
\NormalTok{                                 V :=      τX−1′ (Ωz ),}
\NormalTok{                                                 z∈Y0}
\NormalTok{then V is an open subset of W containing Y = supp(g).}
\NormalTok{   Let y be a point in ∂Y , e an edge emanating from y in S(X ′ ) not contained in Y and}
\NormalTok{Iy,e the corresponding half{-}open interval in the star{-}shaped open neighborhood Ωy . Note}
\NormalTok{that g|e = 0. We may shrink the half{-}open interval Iy,e in e, and so Ωy and V , such that}
\NormalTok{        g + = (ϕan ∗}
\NormalTok{                1 ) (− log |xn1 | +       max          log |xj |) = (ϕan ∗}
\NormalTok{                                                                      1 ) (− log |xn1 | + log |xl1 |)}
\NormalTok{                                       j∈\{0,...,n1 \}}

\NormalTok{        g − = (ϕan ∗}
\NormalTok{                2 ) (− log |xn2 | +       max          log |xj | = (ϕan ∗}
\NormalTok{                                                                     2 ) (− log |xn2 | + log |xl2 |))}
\NormalTok{                                       j∈\{0,...,n2 \}}

\NormalTok{on Iy,e for some li ∈ \{0, . . . , ni \}. Define the map Ni,max : Iy,e → N as follows}
\NormalTok{            Ni,max (η) := |\{j ∈ \{0, . . . , ni \}|(ϕan ∗                an ∗}
\NormalTok{                                                   i ) (log |xj |) = (ϕi ) (log |xli |)\}|.}

\NormalTok{Shrink Iy,e again such that (ϕan   ∗}
\NormalTok{                                i ) (log |xj |) is affine on Iy,e for every j ∈ \{0, . . . , ni \}. Con{-}}
\NormalTok{sequently, the function Ni,max is constant on Iy,e \textbackslash{}\{y\} as (ϕan       ∗                         an ∗}
\NormalTok{                                                                  i ) (log |xli |) = max(ϕi ) (log |xj |)}
\NormalTok{on Iy,e . Hence, we write Ni,max for this constant value.}
\NormalTok{   Define}
\NormalTok{                                                       2}
\NormalTok{          Ωi,k :=\{η ∈ Iy,e | (ϕan  ∗}
\NormalTok{                               i ) (log |xj |)(η) +      \textless{} (ϕan   ∗}
\NormalTok{                                                              i ) (log |xli |)(η)}
\NormalTok{                                                       k}
\NormalTok{                   if (ϕan  ∗                      an ∗}
\NormalTok{                        i ) (log |xj |)|Iy,e 6= (ϕi ) (log |xli |)|Iy,e , j ∈ \{0, . . . , ni \}\}}

\NormalTok{Let N ′ ∈ N such that Ωk := (Ω1,k ∩ Ω2,k )\textbackslash{}\{y\} is a non{-}empty connected open subset}
\NormalTok{of Iy,e for every k ≥ N ′ . We have Ωk ⊂ Ωk+1 , and we work in the following with}
\NormalTok{Ωy,e := ΩN ′ . Then we get for every k ≥ N ′ using properties i){-}iii) that}
\NormalTok{(4.4)    gk+ = (ϕan ∗                            an ∗                       an ∗}
\NormalTok{                 1 ) (− log |xn1 |) + Mn1 , 1 ((ϕ1 ) (log |x0 |), . . . , (ϕ1 ) (log |xn1 |))}
\NormalTok{                                                 k}

\NormalTok{            = (ϕan ∗                                  an ∗                        an ∗}
\NormalTok{                1 ) (− log |xn1 |) + MN1,max −1, 1 ((ϕ1 ) (log |xl1 |), . . . , (ϕ1 ) (log |xl1 |))}
\NormalTok{                                                           k}

\NormalTok{            = (ϕan ∗                   an ∗}
\NormalTok{                1 ) (− log |xn1 |) + (ϕ1 ) (log |xl1 |) + MN1,max −1, 1 (0, . . . , 0)}
\NormalTok{                                                                                k}
\NormalTok{                 +}
\NormalTok{            = g + MN1,max −1, 1 (0, . . . , 0)}
\NormalTok{                                   k}

\NormalTok{on Ωy,e . Set Ck+ := MN1,max −1, 1 (0, . . . , 0) ∈ [0, k1 ]. Analogously,}
\NormalTok{                                 k}

\NormalTok{(4.5)                           gk− = g− + MN2,max −1, 1 (0, . . . , 0)}
\NormalTok{                                                                k}
\NormalTok{24                                                   V. WANNER}


\NormalTok{on Ωy,e . Set Ck− := MN2,max −1, 1 (0, . . . , 0) ∈ [0, k1 ]. Due to g = g+ − g − = 0 on Ωy,e , we}
\NormalTok{                                 k}
\NormalTok{have}
\NormalTok{                      gk+ − gk− = MN1,max −1, 1 (0, . . . , 0) − MN2,max −1, 1 (0, . . . , 0)}
\NormalTok{                                                 k                             k}

\NormalTok{on Ωy,e .}



\NormalTok{                                                                                gk+}
\NormalTok{                                                                                gk−}
\NormalTok{                                                                                       g+ = g−}




\NormalTok{            •                      •[}
\NormalTok{                                   y                                                  xe}
\NormalTok{                S(X ′ ) ∩ Y                                Ωy,e}

\NormalTok{     W                                                     Iy,e}


\NormalTok{                                                      FIGURE 2}


\NormalTok{ 3. Step: Modify the constructed functions gk± on the open subset V ⊂ W}
\NormalTok{          such that (b) and (c) are satisfied.}

\NormalTok{   We start with (b). To ensure that the difference gk+ − gk− = (ϕan  ∗         an ∗}
\NormalTok{                                                                   1 ) φ1,k − (ϕ2 ) φ2,k is}
\NormalTok{non{-}negative on V , we work with gk+ := (ϕan ∗       1                   an ∗}
\NormalTok{                                          1 ) φ1,k + k on V instead of (ϕ1 ) φ1,k , which is}
\NormalTok{                                                                         +}
\NormalTok{still a smooth psh function on V converging uniformly to the function g = − log ks1 kL1 ,}
\NormalTok{and}
\NormalTok{                       1                               1}
\NormalTok{          (ϕan ∗}
\NormalTok{            1 ) φ1,k +   − (ϕan  ∗         an ∗}
\NormalTok{                              2 ) φ2,k ≥ (ϕ1 ) φ1,k +    − g+ + g − − (ϕan  ∗}
\NormalTok{                                                                         2 ) φ2,k}
\NormalTok{                       k                               k}
\NormalTok{                                             1 1}
\NormalTok{                                       ≥0+ − =0}
\NormalTok{                                             k k}
\NormalTok{on V . Note that we have used g = g + − g− ≥ 0 on V and property i) of Mni ,k .}
\NormalTok{   Next, we deal with (c), i.e. we modify gk+ and gk− such that gk := gk+ − gk− has compact}
\NormalTok{support on V . As in Step 2, let y be a boundary point of Y , Ωy a star{-}shaped open}
\NormalTok{neighborhood of y and Iy,e = [y, xe ) be a half{-}open interval of Ωy contained in an edge e}
\NormalTok{of S(X ′ ) emanating from y and pointing outwards of Y . Recall that we have constructed}
\NormalTok{in the second step the open subset Ωy,e of Iy,e . Here, we start with the construction}
\NormalTok{of an affine function on the closed annulus A := τX−1′ ([y, xe ]) whose graph intersects the}
\NormalTok{graphs of all our functions gk± on Ωy,e for k big enough. The annulus A is isomorphic to}
\NormalTok{a closed annulus A′ := S(a, b) = trop−1 ([val(b), val(a)]) in G1,an                     ×}
\NormalTok{                                                                  m for some a, b ∈ K with}
\NormalTok{                                                 ∼     ′}
\NormalTok{|a| \textless{} |b| and trop := − log |T |. Let Φy,e : A − → A be an isomorphism. For simplicity, we}
\NormalTok{may assume that val(b) = 0. Then we can identify Iy,e = [y, xe ) with the real half{-}open}
\NormalTok{                                                                              ′ ∈Ω}
\NormalTok{interval [0, val(a)) via tropy,e := − log |T | ◦ Φy,e . Choose points ζy,e , ζy,e  y,e ⊂ e, and}
\NormalTok{   COMPARISON OF TWO NOTIONS OF SUBHARMONICITY ON NON{-}ARCHIMEDEAN CURVES 25}


\NormalTok{m ∈ N, c ∈ R such that the function}

\NormalTok{                              ψy,e : Iy,e → R; ζ 7→ m · tropy,e (ζ) + c}

\NormalTok{satisfies}

\NormalTok{            ψy,e (ζy,e ) = ((ϕan ∗                   an ∗                          +}
\NormalTok{                              1 ) (− log |xn1 |) + (ϕ1 ) (− log |xl1 |))(ζy,e ) = g (ζy,e ),}

\NormalTok{           ′ ) = g + (ζ ′ ). Recall that N ′ was fixed in Step 2 to define Ω . Since g ±}
\NormalTok{and ψy,e (ζy,e    N ′ y,e                                                   y,e       k}
\NormalTok{converges uniformly to g+ = g − on Ωy,e , there is an N ′′ ≥ N ′ such that}

\NormalTok{                             sup (gk± − g ± ) = sup |gk± − g± | ≤ CN}
\NormalTok{                                                                   +}
\NormalTok{                                                                     ′}
\NormalTok{                            x∈Ωy,e                   x∈Ωy,e}

\NormalTok{                                        +       ± = C + + 1/N ′ , and hence}
\NormalTok{for every k ≥ N ′′ . By (4.4), we have gN ′ − g      N′}

\NormalTok{                               gk± − gN}
\NormalTok{                                      +      ±}
\NormalTok{                                        ′ = gk − g}
\NormalTok{                                                   ±         +}
\NormalTok{                                                     + g± − gN ′ \textless{} 0}


\NormalTok{on Ωy,e for every k ≥ N ′′ . Thus, for every k ≥ N ′′ there is a point ζk+ in Ωy,e such that}
\NormalTok{ψy,e (ζk+ ) = gk+ (ζk+ ). Due to g + = g− ≤ gk− ≤ gk+ by (4.5) and construction, there is also}
\NormalTok{for every k ≥ N ′′ a point ζk− in Ωy,e such that ψy,e (ζk− ) = gk− (ζk− ). Recall that g− = g+}
\NormalTok{and gk+ and gk− with k ≥ N ′ are affine on Ωy,e .}
\NormalTok{   We choose ε such that}

\NormalTok{                           Ωy,e,\textless{} := \{ζ ∈ Ωy,e | ψy,e (ζ) + 2ε \textless{} g+ (ζ)\}}
\NormalTok{                           Ωy,e,\textgreater{} := \{ζ ∈ Ωy,e | ψy,e (ζ) \textgreater{} g+ (ζ) + 2ε\}}

\NormalTok{are non{-}empty open subsets of Ωy,e , and we set Γy,e := Ωy,e \textbackslash{}(Ωy,e,\textless{} ∪ Ωy,e,\textgreater{} ).}
\NormalTok{   In the following, we smoothen the piecewise affine functions max(gk± , ψy,e ) in a proper}
\NormalTok{way. One can construct a smooth symmetric convex 1{-}Lipschitz continuous function}
\NormalTok{θε : R → (0, ∞) such that θε (a) = |a| if |a| ≥ ε. We set}
\NormalTok{                                                    a + b + θε (a − b)}
\NormalTok{(4.6)                                mε (a, b) :=                      .}
\NormalTok{                                                            2}
\NormalTok{  Then the smooth function mε : R2 → R satisfies the following properties:}
\NormalTok{      i) mε is convex;}
\NormalTok{     ii) max(a, b) ≤ mε (a, b) ≤ max(a, b) + 2ε ;}
\NormalTok{    iii) mε (a, b) = max(a, b) whenever |a − b| ≥ ε;}
\NormalTok{    iv) mε is increasing in every variable.}
\NormalTok{We define the functions}

\NormalTok{                                          gek+ := mε (gk+ , ψy,e ),}
\NormalTok{                                           gek− := mε (gk− , ψy,e )}

\NormalTok{on Ωy,e . Then e gk+ (resp. gek− ) coincides with gk+ (resp. with gk− ) on Ωy,e,\textless{} for every}
\NormalTok{k ≥ N since g = g− ≤ gk− ≤ gk+ on Iy,e by (4.5) and by construction. The functions}
\NormalTok{       ′′       +}

\NormalTok{gk+ converge uniformly to g − = g+ , so we can choose Ny,e ≥ N ′′ such that for all k ≥ Ny,e ,}
\NormalTok{we have g + + ε ≥ gk+ ≥ gk− on Ωy,e . Then gek+ and gek− coincide with ψy,e on Ωy,e,\textgreater{} for all}
\NormalTok{                 gk+ − gek− = 0 on Ωy,e,\textgreater{} for every k ≥ Ny,e . We do this for every y ∈ ∂Y}
\NormalTok{k ≥ Ny,e . Thus, e}
\NormalTok{and for every edge e in S(X ′ ) emanating from y and pointing outwards of Y .}
\NormalTok{26                                              V. WANNER}


\NormalTok{                                                                                 ψy,e}




\NormalTok{                                                                           +}
\NormalTok{                                                                          gN  ′}

\NormalTok{                                                                          gk+}
\NormalTok{                                                                          gk− g+ = g−}
\NormalTok{                        g+ + ε}
\NormalTok{               max(gk± , ψy,e )}


\NormalTok{                                                     ζy,e}
\NormalTok{          •                       •[            [     •      ]}
\NormalTok{                                                      Γy,e}
\NormalTok{                                  y                                          xe}
\NormalTok{              S(X ′ ) ∩ Y              Ωy,e,\textless{}                    Ωy,e,\textgreater{}}
\NormalTok{                                                     Ωy,e}
\NormalTok{      W                                               Iy,e}

\NormalTok{                                             FIGURE 3}
\NormalTok{                                                   S}
\NormalTok{   Recall that we work on the open subset V = z∈Y0 τX−1′ (Ωz ) of W containing Y , where}
\NormalTok{Ωz is a star{-}shaped open neighborhood of z in S(X ′ ) and Iz,e = Ωz ∩ e for every edge}
\NormalTok{e in S(X ′ ) emanating from z. We have ∂Y ⊂ Y0 and the sets Ωy \textbackslash{}Y with y ∈ ∂Y are}
\NormalTok{disjoint of all other star{-}shaped open neighborhoods Ωz . We write Sy for the set of edges}
\NormalTok{e in S(X ′ ) emanating from y and pointing outwards of Y . Then for every y ∈ ∂Y and}
\NormalTok{e ∈ Sy , we have constructed the open subset ΩS     y,e of Iy,e = [y, xe ) ⊂ Ωy .}
\NormalTok{   Let VY be the connected component of V \textbackslash{}( y∈∂Y ;e∈Sy Γy,e ) containing Y and Ve be}
\NormalTok{the connected component of V \textbackslash{}Γy,e containing xe . Note that by the construction of V}
\NormalTok{all connected components Ve are pairwise disjoint. We can extend gek± to a continuous}
\NormalTok{function on V by}
\NormalTok{                                \{}
\NormalTok{                                |      ±                       −1}
\NormalTok{                                \{mε (gk , ψy,e ◦ τX ′ ) on τX ′ (Ωy,e ),}
\NormalTok{                         gek± := gk±                      on VY ,}
\NormalTok{                                |}
\NormalTok{                                \{}
\NormalTok{                                 ψy,e ◦ τX ′              on Ve .}

\NormalTok{  Then the functions gek+ − gek− are non{-}negative and continuous on V with compact}
\NormalTok{support for every k ≥ maxy,e Ny,e . Recall that by property iv), mε is increasing in every}
\NormalTok{variable. From now on we only consider k ∈ N with k ≥ maxy∈∂Y,e∈Sy Ny,e .}

\NormalTok{                                                gk+ (resp. gek− ) converge uniformly to a function}
\NormalTok{      4. Step: Show that the modified functions e}
\NormalTok{               e+         −                g − ge− on V .}
\NormalTok{               g (resp. ge ) such that g = e+}


\NormalTok{     Define the following functions on V}
\NormalTok{                                \{}
\NormalTok{                                |       ±}
\NormalTok{                                \{mε (g , ψy,e ◦ τX ′ )       on τX−1′ (Ωy,e ),}
\NormalTok{                         ge± := g±                           on VY ,}
\NormalTok{                                |}
\NormalTok{                                \{}
\NormalTok{                                  ψy,e ◦ τX ′                on Ve .}
\NormalTok{   COMPARISON OF TWO NOTIONS OF SUBHARMONICITY ON NON{-}ARCHIMEDEAN CURVES 27}


\NormalTok{  By construction, these functions are well{-}defined and continuous. Since g = g + − g−}
\NormalTok{on V with g = g + − g− = 0 on V \textbackslash{}Y , we clearly have g = ge+ − e    g− on V .}
\NormalTok{                           +         −}
\NormalTok{  Next, we show that e   gk (resp. gek ) converge uniformly to ge (resp. to ge− ) on V . We}
\NormalTok{                                                                 +}

\NormalTok{                                                                                           gk− )}
\NormalTok{                                                                                gk+ (resp. e}
\NormalTok{know that gk+ (resp. gk− ) converge uniformly to g + (resp. to g− ) on V , thus e}
\NormalTok{                           +             −                  ±                      ±}
\NormalTok{converge uniformly to g (resp. to g ) on VY . Since gek = ψy,e ◦ τX ′ = e        g on Ve , it}
\NormalTok{                                          −1                        −1}
\NormalTok{remains to consider the open subset τX ′ (Ωy,e ). For every x ∈ τX ′ (Ωy,e ), we have}
\NormalTok{ gk± (x) − ge± (x)| = mε (gk± (x), ψy,e (τX ′ (x)) − mε (g± (x), ψy,e (τX ′ (x)))}
\NormalTok{|e}
\NormalTok{                       gk± (x) − g ± (x) + θε (gk± (x) − ψy,e (τX ′ (x))) − θε (g± (x) − ψy,e (τX ′ (x)))}
\NormalTok{                   =}
\NormalTok{                                                               2}
\NormalTok{                       gk± (x) − g ± (x)   g± (x) − g ± (x)}
\NormalTok{                   ≤                     + k}
\NormalTok{                               2                  2}
\NormalTok{                   ≤ |gk± (x) − g ± (x)|}
\NormalTok{where we used that θε is 1{-}Lipschitz continuous to get the inequality. Due to the uniform}
\NormalTok{convergence of gk± to g± on V , which contains τX−1′ (Ωy,e ), we get}
\NormalTok{                               lim      sup            gk± (x) − e}
\NormalTok{                                                      |e         g ± (x)| = 0.}
\NormalTok{                              k→∞ x∈τ −1 (Ω}
\NormalTok{                                       X′     y,e )}


\NormalTok{  5. Step: Show that the modified functions gek+ and e}
\NormalTok{                                                     gk− are smooth and psh on V .}


\NormalTok{  Note that both properties are local. We already know that gk± ∈ C ∞ (VY ), so it}
\NormalTok{remains to find for every point x in τX−1′ (Ωy,e ) ∪ Ve an open neighborhood Vx in V such}
\NormalTok{that gek± is smooth and psh on Vx .}
\NormalTok{  We start with a point x ∈ Ve . Choose an open neighborhood Vx of x in the open}
\NormalTok{subset Ve ⊂ V , then gek± is given by ψy,e ◦ τX ′ on Vx . For every ζ ∈ Vx ⊂ A, we have}
\NormalTok{            gek± (ζ) = (ψy,e ◦ τX ′ )(ζ) = (ψy,e ◦ τA )(ζ) = Φ∗y,e (−m · log |T | + c)(ζ)}
\NormalTok{on Vx . We first show}
\NormalTok{                   −m · log |T | + c ∈ ker(d′ d′′ : C ∞ (G1,an   1,1 1,an}
\NormalTok{                                                          m ) → A (Gm )).}

\NormalTok{Consider the tropical chart (V, ϕU ) = (G1,an                  1,an}
\NormalTok{                                               m , id) of Gm . Then tropU = log |T |, and}
\NormalTok{                                                    1,an}
\NormalTok{so −m · log |T | can be written as the triple (Gm , id, λ), where λ : R → R is the affine}
\NormalTok{                                                                               1,an}
\NormalTok{function t 7→ mt + c. Thus, −m · log |T | + c is a smooth function on Gm            (cf. Definition}
\NormalTok{                                                                                    1,an}
\NormalTok{3.11). The (1, 1) form d d (−m · log |T | + c) is given by the triple (Gm , id, d′ d′′ λ).}
\NormalTok{                              ′ ′′}

\NormalTok{Since λ is affine, the form d′ d′′ λ is zero, and so is d′ d′′ (−m · log |T | + c). Consequently,}
\NormalTok{−m · log |T | + c ∈ ker(d′ d′′ : C ∞ (G1,an         1,1    1,an}
\NormalTok{                                         m ) → A (Gm )). This implies that e              gk± |Vx =}
\NormalTok{Φ∗y,e (−m · log |T | + c) is a smooth psh function on Vx (see Remark 3.21).}
\NormalTok{   Now, consider x ∈ τX−1′ (Ωy,e ). We have just seen that ψy,e ◦ τX ′ is a smooth psh}
\NormalTok{function on τX−1′ ((y, xe )). Using Proposition 3.20, there is a chart (Vx , ϕUx ) with x ∈}
\NormalTok{Vx ⊂ τX−1′ ((y, xe )) such that ψy,e ◦ τX ′ = φ ◦ tropUx on Vx for a smooth function φ on}
\NormalTok{Rr that is convex restricted to every polyhedron contained in tropUx (Vx ). On the other}
\NormalTok{hand, we know that the function gk± is smooth and psh on τX−1′ ((y, xe )) as well. Hence,}
\NormalTok{there is also a chart (Vx′ , ϕUx′ ) with x ∈ Vx′ ⊂ τX−1′ ((y, xe )) such that gk± = φ′ ◦ tropUx′}
\NormalTok{                                            ′}
\NormalTok{on Vx′ for a smooth function φ′ on Rr that is convex restricted to every polyhedron}
\NormalTok{contained in tropUx′ (Vx′ ). Working on the intersection (Vx ∩ Vx′ , ϕUx × ϕUx′ ), which is a}
\NormalTok{28                                         V. WANNER}


\NormalTok{subchart of both [Gub16, Proposition 4.16], we get}
\NormalTok{                           ψy,e ◦ τX ′ = (φ ◦ Trop(π)) ◦ tropUx ∩Ux′}
\NormalTok{                                  gk± = (φ′ ◦ Trop(π ′ )) ◦ tropUx ∩Ux′}
\NormalTok{for the corresponding transition functions π, π ′ satisfying ϕUx = π ◦ (ϕUx × ϕUx′ ) and}
\NormalTok{ϕUx′ = π ′ ◦ (ϕUx × ϕUx′ ). Since Trop(π) and Trop(π ′ ) are integral affine functions on}
\NormalTok{     ′                                                                                     ′}
\NormalTok{Rr+r , the composition φ◦Trop(π) (resp. φ′ ◦Trop(π ′ )) is still a smooth function on Rr+r}
\NormalTok{with a convex restriction to every polyhedron. Thus, mε (φ ◦ Trop(π), φ′ ◦ Trop(π ′ )) is a}
\NormalTok{                           ′}
\NormalTok{smooth function on Rr+r , and the properties i) and iv) of mε imply that the restriction}
\NormalTok{to every polyhedron is convex since the restriction of φ ◦ Trop(π) and φ′ ◦ Trop(π ′ ) are.}
\NormalTok{We have}
\NormalTok{          gk± = mε (gk± , ψy,e ◦ τX ′ ) = mε (φ ◦ Trop(π), φ′ ◦ Trop(π ′ )) ◦ tropUx ∩Ux′}
\NormalTok{          e}
\NormalTok{on Vx ∩ Vx′ for every x ∈ τX−1′ (Ωy,e ), so gek± is a smooth psh function on τX−1′ (Ωy,e ) by}
\NormalTok{Proposition 3.20, which proves Step 5.}
\NormalTok{  Altogether, egk± are smooth psh function on V . We set}
\NormalTok{                                         gk := gek+ − gek− ,}
\NormalTok{                                         e}
\NormalTok{then the functions e}
\NormalTok{                   gk satisfy by construction all the required properties in Step 0.}

\NormalTok{     6. Step: Construction of V ′ and η.}

\NormalTok{   By the construction of V and e  gk , it is no problem to construct the required set V ′ ,}
\NormalTok{                       ′}
\NormalTok{i.e. an open subset V of V containing Y such that V ′ ⊂ V and e        gk ∈ Cc∞ (V ′ ). Fur{-}}
\NormalTok{thermore, let V ′′ be an open neighborhood of V ′ in V with V ′′ ⊂ V . The topological}
\NormalTok{space X an is a compact Hausdorff space. Urysohn’s Lemma states the existence of a}
\NormalTok{continuous function η : X an → [0, 1] with η ≡ 1 on V ′ and η ≡ 0 on X an \textbackslash{}V ′′ . Thus, η}
\NormalTok{has compact support in V , i.e. it is the required function in Step 0.}
\end{Highlighting}
\end{Shaded}

\subsection{van-der-Put, newline lines
1--230}\label{van-der-put-newline-lines-1230}

\begin{Shaded}
\begin{Highlighting}[]
\NormalTok{A NNALES DE L’ INSTITUT F OURIER}




\NormalTok{ M ARIUS VAN D ER P UT}
\NormalTok{ The class group of a one{-}dimensional affinoid space}
\NormalTok{ Annales de l’institut Fourier, tome 30, no 4 (1980), p. 155{-}164}
\NormalTok{ \textless{}http://www.numdam.org/item?id=AIF\_1980\_\_30\_4\_155\_0\textgreater{}}




\NormalTok{    © Annales de l’institut Fourier, 1980, tous droits réservés.}
\NormalTok{    L’accès aux archives de la revue « Annales de l’institut Fourier »}
\NormalTok{    (http://annalif.ujf{-}grenoble.fr/) implique l’accord avec les conditions gé{-}}
\NormalTok{    nérales d’utilisation (http://www.numdam.org/conditions). Toute utilisa{-}}
\NormalTok{    tion commerciale ou impression systématique est constitutive d’une in{-}}
\NormalTok{    fraction pénale. Toute copie ou impression de ce fichier doit conte{-}}
\NormalTok{    nir la présente mention de copyright.}




\NormalTok{                  Article numérisé dans le cadre du programme}
\NormalTok{                Numérisation de documents anciens mathématiques}
\NormalTok{                              http://www.numdam.org/}
\NormalTok{Ann. Inst. Fourier, Grenoble}
\NormalTok{30, 4 (1980), 155{-}164.}




\NormalTok{                THE CLASS GROUP}
\NormalTok{      OF A ONE{-}DIMENSIONAL AFFINOID SPACE}

\NormalTok{                           by Marius van der PUT}


\NormalTok{                                Introduction.}

\NormalTok{    The field k is supposed to be complete with respect to a non{-}}
\NormalTok{archimedean valuation. Moreover we will assume that k is algebraically}
\NormalTok{closed. An afflnoid space Y over k is the set of maximal ideals of an}
\NormalTok{affmoid algebra. The standard affinoid algebra is fc\textless{}T\^{},...,T\^{}\textgreater{}= the}
\NormalTok{set of all power series SflaT\^{} • • • T\^{} converging on the closed polydisk}
\NormalTok{                          \{(\^{}....rJe\^{}alllr.Kl\}.}

\NormalTok{An affinoid algebra is a residue class ring of some fc\textless{}Ti,.. .,T\^{}\textgreater{}. An}
\NormalTok{algebraic variety over k can be studied locally by its analytic structure over}
\NormalTok{k, that is by means of afflnoid spaces.}
\NormalTok{    We show that a one{-}dimensional, normal, connected afflnoid space Y is}
\NormalTok{an afflnoid subset ofa non{-}singular, complete curve C over k (Thm 1.1). If}
\NormalTok{Y has a trivial classgroup then Y is in fact an afflnoid subset of P1}
\NormalTok{(Thm 2.1). A curve is locally a unique factorization domain (U.F.D.}
\NormalTok{for short) if and only the curve is a Mumford curve (i.e. can be}
\NormalTok{parametrized by a Schottky group). In general the class group of Y}
\NormalTok{can be expressed in terms of the Jacobi{-}variety of C (prop. 3.1).}
\NormalTok{    Some examples show the connection between the class group of Y and}
\NormalTok{the class group of the (stable) reduction of Y. For fc{-}analytic spaces we refer}
\NormalTok{to}
\NormalTok{    [2], [3]. I thank A. Escassut for bringing the problem on unique}
\NormalTok{factorization on afflnoid spaces to my attention. Related questions are treated}
\NormalTok{in [1].}

\NormalTok{                 1. Affinoid subspaces of an algebraic curve.}

\NormalTok{     A curve C (non{-}singular and complete) over k has a natural structure as}
\NormalTok{(rigid) analytic space over k. This structure is given by a collection}
\NormalTok{156                       MARIUS VAN DER PUT}


\NormalTok{of subspaces Y of C, called affmoid, and a sheaf 0 = (9\^{} with}
\NormalTok{respect to the Grothendieck topology of finite coverings by afflnoids.}
\NormalTok{For any Y, \^{}P(Y) is an affinoid algebra (1{-}dim. and normal) over k}
\NormalTok{with Sp((P(Y)) = Y. We want to show :}

\NormalTok{   1.1. — THEOREM. — Every 1{-}dimensional, normal, connected affinoid space}
\NormalTok{Y = sp(A) is an affinoid subspace of a non{-}singular complete curve.}

\NormalTok{     Proof. — Y is called connected and normal if the algebra A has no}
\NormalTok{idempotents \^{} 0, 1 and A is integrally closed. We use the notations}
\NormalTok{A° = \{/eA| 11/H \textless{}!\}, A00 = \{/eA| ||/||\textless{}1\} and A = A7A00, where}
\NormalTok{11/11 = max \{\textbackslash{}f(y)\textbackslash{} \textbackslash{}y e Y\} is the spectral norm on Y. The algebra A is of}
\NormalTok{finite type over k = the residue field of k and the algebraic variety}
\NormalTok{Yc = Max (A) is called the canonical reduction of Y. There is a natural}
\NormalTok{surjective map R : Y {-}\textgreater{} Y\^{}, also called the canonical reduction. A pure}
\NormalTok{covering of an analytic space X, is an allowed covering \^{} = (U\^{}) by}
\NormalTok{affinoid spaces, such that for every i + j with U, n U, \^{} 0, the set}
\NormalTok{U; n \textbackslash{}Jj is the inverse image of a Zariski open set V\^{} in (U\^{} under the}
\NormalTok{map U, {-}\textgreater{} (U\^{}. The reduction X\^{} of X with respect to \^{} is}
\NormalTok{obtained by glueing the afflne algebraic varieties (U\^{} over the open}
\NormalTok{sets V\^{}.. The result is an algebraic variety over k. If X is separated}
\NormalTok{then the U\^{} n Uj are also affinoid, the V\^{} are afflne and equal to}
\NormalTok{(U\^{}oU), and X\^{} is separated. If X is non{-}singular, 1{-}dimensional,}
\NormalTok{connected and if X\^{} is complete then X is a non{-}singular complete}
\NormalTok{curve over k (see [2] ch. IV 2.2).}
\NormalTok{    Our proof consists of glueing affinoid spaces Y i , . . . , Y , to Y}
\NormalTok{such that the reduction of X = Y u Y\^{} u . . . u Y, with respect}
\NormalTok{to the pure covering \{Y,Yi,...,Yj is complete. Then clearly Y is}
\NormalTok{an affinoid domain of the algebraic curve X. The 1{-}dimensional}
\NormalTok{space Yc lies in a complete 1{-}dimensional Z such that F = Z {-} Y\^{}}
\NormalTok{is a finite set of non{-}singular points. Suppose that we can find for}
\NormalTok{every p e F an affinoid space Yp with canonical reduction}
\NormalTok{Rp : Yp {-}\textgreater{} (Yp)\^{} \textless{}= Z where (Y\^{} is a neighbourhood of p and}
\NormalTok{such that}
\NormalTok{             Y, \^{} R;1\^{} n Y,) \^{} R{-}\^{}Y;), n Y,) c Y.}

\NormalTok{   Then we can glue Yp to Y. The space X = YU u Yp has reduction Z}
\NormalTok{which is complete. So the glueing has to be done locally on Y and Y\^{}. The}
\NormalTok{component C of Z on which p lies can be projected into P\^{}k) such that}
\NormalTok{         THE CLASS GROUP OF A ONE{-}DIMENSIONAL AFFINOID SPACE              157}


\NormalTok{ (the image of) p is still non{-}singular. A good projection onto P1 maps p}
\NormalTok{ onto o and o is an unramified point for the projection. Replacing Y and}
\NormalTok{ Y\^{} by neighbourhoods of p we may therefore suppose :}

\NormalTok{                     W) = d)(\^{}) = feEr.Mrr\^{}JAP),}
\NormalTok{ where}

\NormalTok{     1) e(t) = (t{-}a,) . . . (t{-}a,) with \^{}, .. .,\^{} different points of \^{} ;they}
\NormalTok{ are the residues of \^{}4, . . . , a, e fe°.}
\NormalTok{     2) P is a monic irreducible polynomial of degree n with coefficients in}
\NormalTok{ \textasciitilde{}k[f\}.}
\NormalTok{    . dP}
\NormalTok{     \textquotesingle{} d\textasciitilde{}\textasciitilde{} ls mvertlble as Cement of H[\_t,(e(t))\textasciitilde{}1,5]/(P).}
\NormalTok{    4) the point « p » corresponds to t = 0.}
\NormalTok{ Then \^{}(Y)° has the form fe°\textless{}T,U,S\textgreater{}/(TE(T)U{-}l,Q) where}

\NormalTok{              E(T) = ( T {-} a i ) ; . . (T{-}a,)   and     Q = P.}

\NormalTok{Since Q is general with respect to the variable S, we can apply Weierstrass{-}}
\NormalTok{division and assume that Q is a monic polynomial of degree n in S with}
\NormalTok{coefficients in fe°\textless{}T,U\textgreater{}/(TE(T)U{-}l). Suppose that we can find a monic}
\NormalTok{polynomial Q* of degree n in S and coefficients in fe°\textless{}T,V\textgreater{}/(E(T)V{-}l)}
\NormalTok{such that}
\NormalTok{                   fe°\textless{}T,U,S\textgreater{}/(TE(T)U{-}l,Q*) \^{} \^{}(Y)°.}

\NormalTok{Then Y\^{} = Sp(fe\textless{}T,V,S\textgreater{}/(E(T)V{-}l,Q*)) has the required properties. So we}
\NormalTok{have to get rid of the negative powers of T in the coefficients of}

\NormalTok{                    C\^{}S\^{}\^{}S"{-}\^{} ... +\^{}.}


\NormalTok{    1.2. {-} LEMMA. {-} If Q* = S" + a;{-}\^{}{-}1 \^{} . . . + a? \^{}5}
\NormalTok{coefficients in A = fe°\textless{}T,U\textgreater{}/(TE(T)U{-}l) \^{} Q* = Q = p, then}
\NormalTok{   a) Q* i5 irreducible}
\NormalTok{   b) Q* /ia5 a z\^{}ro in \^{}(Y)°}
\NormalTok{   c) fe\textless{}T,U,S\textgreater{}/(TE(T)U{-}l,Q*) \^{} (P(Y).}

\NormalTok{    Proo/ {-} a) Let Q* be reducible over the quotient field of A.}
\NormalTok{Since A is normal, Q* is a product of monic polynomials with}
\NormalTok{coefficients in A. This contradicts the irreducibility of Q* = P.}
\NormalTok{158                       MARIUS VAN DER PUT}


\NormalTok{                             f      \^{}Q*1}
\NormalTok{   b) First we show that \textbackslash{}Q*,——\^{} generates the unit ideal in A[S]. Let}
\NormalTok{                          I   rfS J}
\NormalTok{                                          dQ*}
\NormalTok{m be a maximal ideal containing Q* and ——. If m n k° \^{} 0 then m}
\NormalTok{                                           dS}
\NormalTok{induces a maximal ideal of H\^{}t,(te(t))\textasciitilde{}\^{}][S] = A[S] containing P and}
\NormalTok{dP}
\NormalTok{—. This contradicts our assumptions on P. So m corresponds to a}
\NormalTok{aS}
\NormalTok{maximal ideal m\^{}, of fe\textless{}T,U\textgreater{}/(TE(T)U{-}l)[S], containing Q* and}
\NormalTok{dQ*}
\NormalTok{\textasciitilde{}dS\textasciitilde{}\textquotesingle{}}
\NormalTok{   If mi nk\textless{}T,U\textgreater{}/(TE(T)U{-}l) + 0 then m\^{}, is the kernel of a}
\NormalTok{homomorphism in k given by T i—\textgreater{} \^{} e k, S i—\textgreater{} \^{} 6 k with}

\NormalTok{                  |?4|\^{} 1,          P4E(5iO|= 1,     |\^{}1 \textless{} 1}

\NormalTok{                                 / dP\textbackslash{}}
\NormalTok{since Q*\^{}) = 0. Ffom              P.— = \^{}[\^{}(\^{}(O\textquotesingle{}\^{}S] it follows that}
\NormalTok{                                 \textbackslash{} »S/}
\NormalTok{                                 dO*}
\NormalTok{                   Z,(S)Q alt +Z,(S){-}\^{}=l + S \^{}S1}
\NormalTok{                                  "\^{}       i\textgreater{}o}

\NormalTok{for certain Z\^{}, Z2 e A[S] and \^{} 6 A with ||flf|| \textless{} 1. The substitution}
\NormalTok{T i—\^{} \^{} i ; S i—\textgreater{} \^{}2 makes 0 = 1 + \^{} \^{}Ai)\^{}2» which is impossible. So}
\NormalTok{                                   i\textgreater{}o}
\NormalTok{m and m\^{} correspond to an ideal of L[S] with L the quotient field of A.}
\NormalTok{                                              dQ*}
\NormalTok{Since Q* is irreducible, this means that —— = 0. This is obviously in}
\NormalTok{                                         dS}
\NormalTok{contradiction with (P,—) = K[\^{}?(0)\textasciitilde{}1].}
\NormalTok{                   \textbackslash{} \^{}S/}
\NormalTok{      We conclude the existence of Z\^{}, Z\^{} e A[S] with}

\NormalTok{                                                   dQ*}
\NormalTok{                        1 =Z,(S)Q*+Z2(S)——.}
\NormalTok{                                         aS}

\NormalTok{    By Newton\textquotesingle{}s method we will show that Q* has a zero in \^{}(Y)°. Let}
\NormalTok{T| e \^{}(Y)° satisfy IIQ\^{}rOII \textless{} 1 (e.g. TI is the residue of S mod Q in}
\NormalTok{                                                dO*}
\NormalTok{W).         Then       1 {-} Z,(Ti)Q*(ri) = Z\^{}) {-}\^{} (TI)       and  since}
\NormalTok{                                                 aS}
\NormalTok{        THE CLASS GROUP OF A ONE{-}DIMENSIONAL AFFINOID SPACE                    159}


\NormalTok{                                               dO*}
\NormalTok{l|Zi(ri)Q*(il)|| \textless{} 1   it   follows    that    ——(r|)      is   invertible.   Put}
\NormalTok{                                               aS}
\NormalTok{               /d0*   \textbackslash{}\textasciitilde{}1}
\NormalTok{Th = TI {-} Q*(\^{})({-}\^{}{-}(Ti)t . Then ||Q*(Tii)|| \textless{} ||Q*(Ti)||2. The usual}

\NormalTok{procedure and the completeness of (P(Y)0 show the existence of a root of Q*}
\NormalTok{in (P(Y)°.}
\NormalTok{    c) The quotient field of A[S]/Q* is contained in that of A[S]/Q,}
\NormalTok{because of (b). Both fields are extensions of degree n of the quotient field}
\NormalTok{of A. So they are equal. The rings fe\textless{}T,U,S\textgreater{}/(TE(T)U{-}l,Q*) and \^{}P(Y)}
\NormalTok{are both the integral closure of k\textless{}T,U\textgreater{}/(TE(T)U{-}l) in that field. So they}
\NormalTok{are equal.}

\NormalTok{  End of the proof of 1.1. — We choose Q* with coefficients in}
\NormalTok{k°\textless{}T,V\textgreater{}/(VE(T){-}l) and Q* = P.}

\NormalTok{   1.3. — COROLLARY. — Let Y be as in (1.1); then Y is affinoid in a curve X}
\NormalTok{(complete non{-}singular\} such that X — Y\^{}. is a finite set of non{-}singular points.}

\NormalTok{                            2. Unique factorization.}

\NormalTok{    We want to show the following :}

\NormalTok{   2.1. — THEOREM. — Let Y = Sp A be a 1{-}dimensional connected affinoid}
\end{Highlighting}
\end{Shaded}

\subsection{GJR, newline lines
1610--1799}\label{gjr-newline-lines-16101799}

\begin{Shaded}
\begin{Highlighting}[]
\NormalTok{   We claim that mϕ (X , P) = 1 for P ∈ \{A, B, C, D\}. This can be checked directly from}
\NormalTok{the definition of mϕ (X , P). We give the proof in case P = A which reduces to the generic}
\NormalTok{point of the line x = 0 in the special fiber of Xmin . By [Ber90, Proposition 2.4.4(ii)],}
\NormalTok{we have H  f(A) = K e ( y). Recall that trop(A) = (1, 0). If M = Z2 be the character lattice}
\NormalTok{of T = G2m , then Ke ( M(A))}
\NormalTok{                       e      is the subfield of H f(A) generated by the reductions of x/p}
\NormalTok{and of y/1 = y. This shows H     f(A) = K e (M}
\NormalTok{                                             e (A)), and hence 4.2(8) yields that mϕ (X , A) =}
\NormalTok{[Hf(A) : K}
\NormalTok{         e (Me (A))] = 1. The proof for the points P = B, C, D is similar.}
\NormalTok{   We conclude that mrel (P) = mϕ (X , P) = 1 for every point P of Sϕ (X ). It follows that}
\NormalTok{mϕ (X , ω) just counts the preimages of ω in Sϕ (X ) under ϕtrop ; these tropical multiplicities}
\NormalTok{are indicated in Figure 1.}

\NormalTok{                              5. PIECEWISE LINEAR FUNCTIONS}
\NormalTok{   In this section we discuss piecewise linear functions on analytic spaces. Most of these}
\NormalTok{results are standard in one form or another: for instance, see [CD12, Gub98, GM19].}
\NormalTok{However, the terminology in the literature is somewhat nonstandard, so it is important}
\NormalTok{to fix our ideas. We will find it useful to work slightly more generally with piecewise}
\NormalTok{linear functions with slopes contained in an additive subgroup Λ ⊂ R, and in some cases}
\NormalTok{new arguments are needed. In practice we will usually take Λ = Z, Q, or R, in which case}
\NormalTok{either Λ = Z or Λ is divisible.}
\NormalTok{   In this section, K is a non{-}Archimedean field with a non{-}trivial valuation v and value}
\NormalTok{group Γ := v(K × ) ⊂ R. We consider a good strictly K{-}analytic space X and an additive}
\NormalTok{subgroup Λ ⊂ R. Recall from 2.3 that a strictly analytic space is good provided that that}
\NormalTok{every point has a strictly affinoid neighbourhood. Most strictly analytic spaces occurring}
\NormalTok{in practice are good: for example, any strictly affinoid space, a space with no boundary,}
\NormalTok{or an analytic subdomain of the analytification of an algebraic variety. We restrict our}
\NormalTok{attention to good spaces as then the coherent modules on the Berkovich space X are}
\NormalTok{the same as the coherent modules on the space X G endowed with the G{-}topology [Ber93,}
\NormalTok{Proposition 1.3.4], which makes it reasonable to work with line bundles on X . We always}
\NormalTok{               FORMS ON BERKOVICH SPACES BASED ON HARMONIC TROPICALIZATIONS                       31}

\NormalTok{use the G{-}topology on X generated by the strictly affinoid domains. Let C(X ) be the space}
\NormalTok{of continuous real functions on X .}
\NormalTok{Definition 5.1. A function h: X → R is called piecewise Λ{-}linear or Λ{-}PL if there exists a}
\NormalTok{                                                                                                    ×}
\NormalTok{G{-}covering \{XP i \} of X such that for each i, there are finitely many λi j ∈ Λ and f i j ∈ O (X i )}
\NormalTok{with h|X i = j λi j log | f i j |. A Z{-}PL function is called piecewise linear or PL.}
\NormalTok{  We denote the group of PL functions on X by PL(X ), and the group of Λ{-}PL functions}
\NormalTok{by PLΛ (X ). These are subgroups of C(X ).}
\NormalTok{5.2. Some remarks:}
\NormalTok{   (1) A function h: X → R is PL if and only if there exists a G{-}covering \{X i \} of X and}
\NormalTok{       f i ∈ O (X i )× such that h|X i = log | f i |.}
\NormalTok{   (2) The set of Λ{-}PL functions on X forms a group since any two G{-}coverings \{X i \} and}
\NormalTok{       \{X ′j \} admit the common refinement \{X i ∩ X ′j \}.}
\NormalTok{   (3) The restriction of a Λ{-}PL function to a strictly analytic domain is Λ{-}PL.}
\NormalTok{   (4) The presheaf U 7→ PLΛ (U) is a sheaf in the G{-}topology, and therefore in the ana{-}}
\NormalTok{       lytic topology as well.}
\NormalTok{   (5) A function h: X → R is Λ{-}PL if and only if there exists a G{-}covering \{X i \} of}
\NormalTok{       X suchPthat for each i, there are finitely many λi j ∈ Λ and hi j ∈ PL(X i ) with}
\NormalTok{       h|X i = j λi j hi j .}
\NormalTok{   (6) A function h: X → R is Q{-}PL if and only if, locally in the analytic topology, there}
\NormalTok{       exists a positive integer N such that N h is PL.}
\NormalTok{   (7) If Λ = Z or Λ = Q then the minimum and maximum of two Λ{-}PL functions is}
\NormalTok{       again Λ{-}PL by [GM19, Proposition 2.12(d)].}
\NormalTok{Example 5.3. The assertion of 5.2(7) is not true for Λ strictly containing Q. For example,}
\NormalTok{if the value group of K is Q and X = Gan m}
\NormalTok{                                            = Spec(K[t ±1])an is the split torus of rank 1, then}
\NormalTok{the tropicalization of any strictly affinoid domain in X is a union of compact intervals with}
\NormalTok{rational vertices. We conclude that for Λ = R, the function h := π log |t|+1 = −π trop +1}
\NormalTok{is Λ{-}PL, but max(h, 0) is not Λ{-}PL as max(h, 0) = ϕ ◦ trop for ϕ(u) = max(−πu + 1, 0),}
\NormalTok{and as the piecewise linear function ϕ has a bend at the irrational number 1/π.}
\NormalTok{  Here we list some basic functoriality properties of Λ{-}PL functions.}
\NormalTok{Lemma 5.4. Let h: X → R be a function.}
\NormalTok{   (1) Let f : X ′ → X be a morphism of good strictly K{-}analytic spaces. If h ∈ PLΛ (X ) then}
\NormalTok{       h ◦ f ∈ PLΛ (X ′ ).}
\NormalTok{   (2) Let K ′ /K be a non{-}Archimedean field extension, let X ′ = X ⊗}
\NormalTok{                                                                    b K K ′ , and let π: X ′ → X}
\NormalTok{       be the structure morphism. If h ∈ PLΛ (X ) then h ◦ π ∈ PLΛ (X ′ ).}
\NormalTok{   (3) Let ι : X red → X be the inclusion of the reduction of X , and let h′ = h ◦ ι. Then}
\NormalTok{       h ∈ PLΛ (X ) if and only if h′ ∈ PLΛ (X red ).}
\NormalTok{Proof. The first two assertions are obvious. For the third, by (1), the only thing to show is}
\NormalTok{that if h′ is Λ{-}PL then h is Λ{-}PL. Let U = M (A ) ⊂ X be a strictly    Paffinoid domain, and let}
\NormalTok{                                                           ′                                    ×}
\NormalTok{Ured = M (Ared ) ⊂ X red be its reduction. Suppose that P h |Ured = λi log |g i | for g i ∈ Ared  .}
\NormalTok{If f i ∈ A lifts g i then f i ∈ A × is a unit, and h|U = λi log | f i |. Since X red has a G{-}cover}
\NormalTok{by such affinoid domains, the same is true for X , so h is Λ{-}PL.                                 }
\NormalTok{32                              W. GUBLER, P. JELL, AND J. RABINOFF}

\NormalTok{   At this point we assume that the reader is familiar with formal geometry as recalled}
\NormalTok{in 2.4. We give an interpretation of Λ{-}PL functions in terms of formal geometry which}
\NormalTok{we will use to prove some more subtle properties of piecewise Λ{-}linearity. Recall that}
\NormalTok{a formal K ◦ {-}model of a paracompact strictly K{-}analytic space X is an admissible formal}
\NormalTok{scheme X over K ◦ equipped with an identification Xη = X .}

\NormalTok{     A frame of a line bundle on an open set is a nowhere vanishing section.}

\NormalTok{5.5 (Functions from Formal Models of OX ). We assume now that X is paracompact. This}
\NormalTok{allows us to use Raynaud’s theory of formal K ◦ {-}models, as in 2.4. We will later apply}
\NormalTok{this only locally on a good, strictly analytic space, where we can always find a compact}
\NormalTok{strictly analytic neighbourhood (even a strictly affinoid neighbourhood as X is good).}
\NormalTok{Hence the reader loses nothing by assuming here and in the following that X is compact.}
\NormalTok{(Paracompactness has the advantage that one can apply these considerations also to the}
\NormalTok{analytification of algebraic varieties over K.)}
\NormalTok{   Let X be a formal K ◦ {-}model of X . The isomorphism classes of line bundles L on X}
\NormalTok{equipped with an identification Lη = OX form an abelian group denoted by M (X). Then}
\NormalTok{L ∈ M (X) determines a canonical metric k kL on OX , defined as follows: if s is a frame of}
\NormalTok{L over a formal open subset U ⊂ X, then we set ksk = 1 on Uη . This metric is well{-}defined}
\NormalTok{because any two frames differ by a unit with constant absolute value 1.}
\NormalTok{   Equivalently, we get a canonical function hL = − log k1kL : X → R with − log k f kL =}
\NormalTok{− log | f |+hL for an analytic function f on any open subset U of X . If U = Uη for a formal}
\NormalTok{open subset U of X and if f ∈ Γ (U, L) ⊂ Γ (Uη , OX ) is a frame of L, then by definition}
\NormalTok{0 = − log | f | + hL , so hL is PL because X has a G{-}cover by such U = Uη .}

\NormalTok{5.6 (PL Functions are Model Functions). If X is paracompact, then [GM19, Proposi{-}}
\NormalTok{tion 2.10] yields that h: X → R is PL if and only if there exists a formal K ◦ {-}model (X, L)}
\NormalTok{of (X , OX ) such that h = − log k1kL .}

\NormalTok{5.7 (Λ{-}PL Functions from Formal Models). For paracompact X , we extend 5.6 to the}
\NormalTok{case of Λ{-}PL functions as follows. Choose a formal K ◦ {-}model X of X , and set M (X)Λ :=}
\NormalTok{M (X)⊗Z Λ. We have a unique homomorphism h: M (X)Λ −→ C(X ), written        Pm L 7→ hL , such}
\NormalTok{that hλL =P   −λ log k1kL for all L ∈ M (X) and λ ∈ Λ. Note that if L = i=1 λi Li ∈ M (X)Λ}
\NormalTok{                m}
\NormalTok{then hL = i=1 λi hLi ∈ PLΛ (X ) by 5.2(5).}
\NormalTok{   If f : X′ → X is a morphism of formal K ◦ {-}models of X , then h f ∗ L = hL for all L ∈ M (X)Λ .}
\NormalTok{More generally, if f : X′ → X is any morphism of separated admissible formal schemes}
\NormalTok{then h f ∗ L = hL ◦ fη for all L ∈ M (X)Λ . (See 2.4 for separatedness of formal models.)}
\NormalTok{   For L ∈ M (X)Λ , we get a unique continuous metric k kL on OX defined by − log k f kL =}
\NormalTok{− log | f | + hL for an analytic function f on any open subset U of X .}

\NormalTok{   We want to extend the result of 5.6 locally in the analytic topology to Λ{-}coefficients as}
\NormalTok{in 5.7. To do this we need to show that a Λ{-}PL function is locally in the analytic topology a}
\NormalTok{linear combination of PL functions with Λ{-}coefficients. (For our purposes, we could have}
\NormalTok{taken the statement of Proposition 5.8 as the definition of Λ{-}PL; we chose Definition 5.1}
\NormalTok{because it is more compatible with the terminology in the literature.)}
\NormalTok{                 FORMS ON BERKOVICH SPACES BASED ON HARMONIC TROPICALIZATIONS                             33}

\NormalTok{Proposition 5.8. Suppose that Λ = Z or that Λ is divisible. A function h: X → R is Λ{-}}
\NormalTok{       Pm only if for every x ∈ X , there exists an open neighbourhood U of x such that}
\NormalTok{PL if and}
\NormalTok{h|U = i=1 λi hi for finitely many λ1 , . . . , λm ∈ Λ and h1 , . . . , hm ∈ PL(U).}
\NormalTok{  Before proving the proposition, we state the desired corollary.}
\NormalTok{Corollary 5.9. Suppose that Λ = Z or that Λ is divisible. A real function h: X → R is Λ{-}PL}
\NormalTok{if and only if for every point x ∈ X , there exists a compact strictly analytic neighbourhood}
\NormalTok{U of x, a formal K ◦ {-}model U of U, and L ∈ M (U)Λ , such that h|U = hL . In this case, such a}
\NormalTok{neighbourhood U can be chosen arbitrarily small.}
\NormalTok{Proof.}
\NormalTok{     PmSuppose that h is Λ{-}PL. By Proposition 5.8, after shrinking X , we may assume that}
\NormalTok{h = i=1 λi hi for λ1 , . . . , λm ∈ Λ and PL functions h1 , . . . , hm : X → R. Since every point of}
\NormalTok{X has a compact strictly analytic neighbourhood, we may also assume that X is compact.}
\NormalTok{By 5.6, for each i there is a model (Xi , Li ) of (X , OX ) such that hi = hLi . Finding a model X}
\NormalTok{                                                                              Pm}
\NormalTok{dominating}
\NormalTok{       Pm    all X i , we may   assume  all X i = X. Then  h  =  h  L for L =     λ L . Conversely,}
\NormalTok{                                                                               i=1 i i P}
\NormalTok{                                                                                         m}
\NormalTok{if L = i=1 λi Li on U for λ1 , . . . , λm ∈ Λ and L1 , . . . , Lm ∈ M (U) then hL = i=1 λi hLi is}
\NormalTok{Λ{-}PL by 5.2(5). We can choose U sufficiently small because U can be taken arbitrarily}
\NormalTok{small in Proposition 5.8.                                                                         }
\NormalTok{  We will use the following lemma in the proof of Proposition 5.8; we will use it again}
\NormalTok{several times.}
\NormalTok{Lemma 5.10. Suppose that Λ is a divisible additive subgroup of R. Let P be an abelian}
\NormalTok{group, and let x 1 , . . . , x r ∈ Λ ⊗Z P be a finite collection of elements. Then P            there is a Q{-}}
\NormalTok{                                                                                                   n}
\NormalTok{linearly independent subset \{λ1 , . . . , λn \} ⊂ Λ such that each x i has the form j=1 λi ⊗ x i j}
\NormalTok{for x i j ∈ P.}
\NormalTok{                       P mi}
\NormalTok{Proof. Write x i = j=1        λi j ⊗ yi j . Let λ1 , . . . , λn be a basis of the Q{-}subspace of Λ generated}
\NormalTok{by the λi j . Since Λ is divisible, we may assume that every λi j is a Z{-}linear combination of}
\NormalTok{                                                 P mi}
\NormalTok{λ1 , . . . , λn . Expanding λi j in the sum j=1            λi j ⊗ yi j and rearranging proves the lemma. }
\NormalTok{Proof of Proposition 5.8. The “if” direction is trivial because a covering by open sets is a}
\NormalTok{G{-}covering. Suppose then that h is Λ{-}PL. If Λ = Z then h is itself PL, so we may assume}
\NormalTok{Λ is divisible. Fix x ∈ X , and replace X by a compact strictly analytic neighbourhood of}
\NormalTok{x to assume X is compact.PChoose a G{-}covering \{X i \}i∈I of X by compact strictly analytic}
\NormalTok{                              mi}
\NormalTok{domains such that h|X i = j=1    λi j log | f i j | for finitely many λi j ∈ Λ and f i j ∈ O (X i )× . We}
\NormalTok{may assume that every X i contains x and that I is a finite set. By Lemma 5.10, there}
\NormalTok{are Q{-}linearly}
\NormalTok{       Pm       independent elements λ1 , . . . , λm ∈ Λ such that for every i ∈ I, we have}
\NormalTok{h|X i = k=1 λk hik with hik = log |g ik | for g ik ∈ O (X i )× .}
\NormalTok{  We claim that there is a neighbourhood U of x such that for all i, j ∈ I and all y ∈}
\NormalTok{U ∩ X i ∩ X j we have}
\NormalTok{(5.10.1)                         hik ( y) − hik (x) = h jk ( y) − h jk (x).}
\NormalTok{To prove this, we pick i, j ∈ I, set X i j := X i ∩ X j and consider the smooth tropicalization}
\NormalTok{map H i j : X i j → R2m given by y 7→ (hik ( y), h jk ( y))k=1,...,m . By 3.2 and 3.3, the tropical}
\NormalTok{variety of the germ (X i j , x) is the germ at H i j (x) of a finite union of (Z, Γ ){-}polytopes ∆.}
\NormalTok{34                                     W. GUBLER, P. JELL, AND J. RABINOFF}

\NormalTok{Choosing them sufficiently small, we may       S assume every ∆ is contained in H i j (X i j ). It is}
\NormalTok{enough to show (5.10.1) for y ∈ Ui j := ∆ H i−1         j}
\NormalTok{                                                          (∆) as Ui j is a neighbourhood of x in X i j ,}
\NormalTok{hence has the form U ∩ X i j for a neighbourhood U of x in X .}
\NormalTok{  We fix i, j ∈ I and pick a polytope ∆ as above. Since ∆ is a Z{-}polytope, the set (H i j (x)+}
\NormalTok{Q2m ) ∩ ∆ is dense in ∆. By construction, there are coordinate functions ℓik , ℓ jk of R2m}
\NormalTok{such that hik = ℓik ◦ H i j and h jk = ℓ jk ◦ H i j . Hence it is enough to show (5.10.1) for all}
\NormalTok{y ∈ X i j with H i j ( y) ∈ ∆ and with ω = H i j ( y) − H i j (x) ∈ Q2m . Using}
\NormalTok{                        X                X                X                X}
\NormalTok{                0=         λk hik ( y) −   λk h jk ( y) −     λk hik (x) +     λk h jk (x)}
\end{Highlighting}
\end{Shaded}

\subsection{GJR, newline lines
1800--1985}\label{gjr-newline-lines-18001985}

\begin{Shaded}
\begin{Highlighting}[]
\NormalTok{                        k                   k                    k                      k}
\NormalTok{                        X                                                      }
\NormalTok{                   =         λk hik ( y) − hik (x) − h jk ( y) + h jk (x)}
\NormalTok{                         k}

\NormalTok{and hik ( y) − hik (x) − h jk ( y) + h jk (x) = ℓik (ω) − ℓ jk (ω) ∈ Q, the Q{-}linear independence}
\NormalTok{of λ1 , . . . , λm yields (5.10.1) for all such y. This finishes the proof of the claim.}
\NormalTok{    Let us fix i ∈ I. We deduce from (5.10.1) that there are constants ci jk such that hik =}
\NormalTok{h jk + ci jk on U ∩ X i j for all j ∈ I and k = 1, . . . , m. Evaluating at a rig{-}point of U ∩ X i j , we}
\NormalTok{                    p}
\NormalTok{see that ci jk ∈ Γ . Replacing h by an integer multiple, we may assume ci jk = −v(γi jk ) ∈ Γ}
\NormalTok{for some γi jk ∈ K × . We conclude that γi jk g jk is an invertible function on X j . We claim}
\NormalTok{that there is a unique piecewise linear function hk : U → R that agrees with log |γi jk g jk | =}
\NormalTok{ci jk +h jk on U ∩X j for every j ∈ I. To see that the functions ci jk +h jk indeed glue, we know}
\NormalTok{from (5.10.1) that they agree on intersections up to constants and the claim follows as}
\NormalTok{ci jk + h jk (x) = hik (x) does not depend on j ∈ I. We note that}
\NormalTok{                     X                        X                  X                      }
\NormalTok{                         λk h jk (x) = h(x) =    λk hik (x) =       λk h jk (x) + ci jk}
\NormalTok{                    k                                 k                    k}
\NormalTok{             P}
\NormalTok{and hence        k λ k c i jk = 0. For y ∈ U ∩ X j , we deduce that}
\NormalTok{                             X                    X                    X                X}
\NormalTok{                 h( y) =         λk h jk ( y) =       λk h jk ( y) +       λk ci jk =       λk hk ( y)}
\NormalTok{                             k                    k                    k}
\NormalTok{                   P}
\NormalTok{and hence h =           λk hk on U. This proves the proposition.                                         }
\NormalTok{  With these foundations in place, we can use formal models to show that piecewise}
\NormalTok{Λ{-}linearity can sometimes be checked after pull{-}back, at least when Λ is divisible.}
\NormalTok{Proposition 5.11. Let f : X ′ → X be a finite, flat, surjective morphism of good strictly K{-}}
\NormalTok{analytic spaces. Suppose that Λ is divisible. If h: X → R is a function and h′ = h ◦ f , then}
\NormalTok{h ∈ PLΛ (X ) if and only if h′ ∈ PLΛ (X ′ ).}
\NormalTok{Proof. By Lemma 5.4(1), we only need to show that if h′ is Λ{-}PL then so is h. By [Ber93,}
\NormalTok{Proposition 3.2.7], the map f is open. Let x ∈ X , choosePx ′ ∈ X ′ mapping to x, and}
\NormalTok{                                                                  m}
\NormalTok{let U ′ be an open neighbourhood of x ′ such that h′ |U ′ = i=1 λi h′i for λ1 , . . . , λm ∈ Λ}
\NormalTok{and h′1 , . . . , h′m ∈ PL(U ′ ). Such a neighbourhood exists by Proposition 5.8. By [Ber93,}
\NormalTok{Lemma 3.1.2, Definition 3.2.5], after shrinking U ′ , the morphism f |U ′ : U ′ →Pf (U ′ ) is}
\NormalTok{                                                                                        m}
\NormalTok{finite and flat. Thus we may replace X ′ by U ′ and X by f (U ′ ) to assume h′ = i=1 λi h′i}
\NormalTok{on X ′ . Shrinking further, we may assume X and X ′ are compact. By 5.6, each PL function}
\NormalTok{h′i is induced by a line bundle on a formal model of X ′ . Passing to a dominating formal}
\NormalTok{                FORMS ON BERKOVICH SPACES BASED ON HARMONIC TROPICALIZATIONS                                         35}

\NormalTok{model, we may assume that there is a formal K ◦ {-}model X′ of X ′ and L′i ∈ M (X′ ) such}
\NormalTok{that h′i = hL′i for each i.}
\NormalTok{   Using Raynaud’s theorem [Bos14, Theorem 8.4.3, Lemma 8.4.4], after passing to an}
\NormalTok{admissible formal blowup, we can extend f to a morphism f : X′ → X for some for{-}}
\NormalTok{mal K ◦ {-}model X of X . By [BL93b, Theorem 5.2, Corollary 5.3(b)], after passing to}
\NormalTok{admissible formal blowups of X and X′ , we may assume f : X′ → X is flat and quasi{-}}
\NormalTok{finite. By [Tem00, Corollary 4.4] the morphism f : X′s → Xs is proper, so it is finite; by}
\NormalTok{Nakayama’s lemma, the morphism f : X′ → X is finite. Working locally, we may thus}
\NormalTok{assume X = Spf(A) and X′ = Spf(A′ ) are formal affines, with X connected. Then A′ is a}
\NormalTok{finite, flat A{-}module, which is locally of finite presentation by [Bos14, Theorem 7.3.4],}
\NormalTok{so it is a locally free A{-}module of some constant rank n by [Sta21, Lemma 02KB]. Thus}
\NormalTok{we can define the norm N (L′ ) of a line bundle L′ on X′ as in [EGAII, 6.5]. This is a line}
\NormalTok{bundle on X and a model of OX⊗n = OX , so we get a homomorphism N : M (X′ ) → M (X).}
\NormalTok{Tensoring with Λ, we have an induced homomorphism N : M (X′ )Λ → M (X)Λ .}
\NormalTok{   Now we can transfer the considerations in [BE21, 8.5] to the analytic setting. For}
\NormalTok{L′ ∈ M (X′ ) and x ∈ X , we claim that}
\NormalTok{                                                Y         my}
\NormalTok{(5.11.1)                        k1(x)kN (L′ ) =    k1( y)kL′ ,}
\NormalTok{                                                                 y7→ x}

\NormalTok{where m y = dimH (x) (O f −1 (x), y ) is the multiplicity of y in f −1 (x). With this in place, letting}
\NormalTok{L = N (L′ ), it follows that}
\NormalTok{                                                  X}
\NormalTok{                                         hL (x) =     m y hL′ ( y).}
\NormalTok{                                                             y7→ x}

\NormalTok{Extending by linearity,}
\NormalTok{                     Pmthe above      equation holds for L′ ∈ M (X′ )Λ and L = N (L′ ) ∈}
\NormalTok{                   ′}
\NormalTok{M (X)Λ . Letting L = i=1 λi Li and L = N (L′ ), we have h′ = hL′ , and for x ∈ X we have}
\NormalTok{                               ′}

\NormalTok{                     X                 X                         X}
\NormalTok{            hL (x) =    m y hL′ ( y) =     m y h ◦ f ( y) = h(x)   m y = n h(x).}
\NormalTok{                          y7→ x                      y7→ x                                 y7→ x}

\NormalTok{It follows that h = hn−1 L is Λ{-}PL.}
\NormalTok{   It remains to prove (5.11.1). Let x ∈ X , and let y1 , . . . , yr be the points of f −1 (x). Let}
\NormalTok{R j := O f −1 (x), y j , let F := H (x), and let m j := dim F (R j ) = m y j . We have the decomposi{-}}
\NormalTok{                                 Q}
\NormalTok{tion R := O ( f −1 (x)) =          j R j into local Artinian rings. Working locally on X = Spf(A),}
\NormalTok{by [Sta21, Lemma 0BUT] we may assume that L′ is trivial over X′ = Spf(A′ ). Choose a}
\NormalTok{frame s′ for L′ and let γ = 1/s′ ∈ (A′ ⊗K ◦ K)× , so k1kL′ = |γ| on X ′ . The norm s := N (s′ )}
\NormalTok{is a frame of N (L), and 1/s = NA′ /A(γ) ∈ (A ⊗K ◦ K)× , so}
\NormalTok{                                  k1(x)kN (L′ ) = |NA′ /A(γ)(x)| = |NR/F (γ)|.}
\NormalTok{Let K j be the residue field of R j , let γ j be the image of γ in R j , and let γ j be the residue}
\NormalTok{class in K j . A Jordan–Hölder argument (see [Ser97, 2.1c, p.9]) shows that NR j /F (γ j ) =}
\NormalTok{NK j /F (γ j )[R j :F ]/[K j :F ] . Since the canonical valuation on H ( y j ) extends the canonical valu{-}}
\NormalTok{ation on the complete non{-}archimedean field F = H (x), we have}
\NormalTok{                 |γ( y j )|[K j :F ] = |γ j |[K j :F ] = |NK j /F (γ j ))| = |NR j /F (γ j )|[K j :F ]/[R j :F ] .}
\NormalTok{36                               W. GUBLER, P. JELL, AND J. RABINOFF}

\NormalTok{Putting all these together, we deduce}
\NormalTok{                                     Y                    Y                 Y           mj}
\NormalTok{        k1(x)kN (L′ ) = |NR/F (γ)| =   |NR j /F (γ j )| =   |γ( y j )|m j =   k1( y j )kL′ .}
\NormalTok{                                        j                   j                j}

\NormalTok{This proves the claim in (5.11.1) and hence the proposition.                                      }
\NormalTok{     Proposition 5.11 gives the following variant of Lemma 5.4(2).}
\NormalTok{Corollary 5.12. Let X be a good strictly K{-}analytic space, let K ′ be the completion of an}
\NormalTok{algebraic field extension of K, let X ′ = X ⊗}
\NormalTok{                                            b K K ′ , and let π: X ′ → X be the structure morphism.}
\NormalTok{Suppose that Λ is divisible. If h: X → R is a function and h′ = h ◦ π, then h ∈ PLΛ (X ) if and}
\NormalTok{only if h′ ∈ PLΛ (X ′ ).}
\NormalTok{Proof. By Lemma 5.4(2), we only need to show that if h′ is Λ{-}PL then so is h. When K ′}
\NormalTok{is a finite extension of K then this is a special case of Proposition 5.11. We reduce to this}
\NormalTok{case as follows.}
\NormalTok{    Since piecewise Λ{-}linearity is a local property in the analytic topology, we may assume}
\NormalTok{that X is compact. Then X ′ is also compact, so by Proposition 5.8 there are finitely many}
\NormalTok{                                       ′}
\NormalTok{compact strictly analytic domains UP   i}
\NormalTok{                                         such that the interiors Vi′ of the Ui′ cover X ′ and such}
\NormalTok{that hi := h |Ui′ has the form hi = j λi j h′i j for finitely many h′i j ∈ PL(Ui′ ) and λi j ∈ Λ.}
\NormalTok{        ′      ′                   ′}

\NormalTok{By [GM19, Propositions 2.7 and 2.10], the functions h′i j extend to PL functions on X ′ .}
\NormalTok{                                                                           P}
\NormalTok{We denote the extensions also by h′i j leading to an extension j λi j h′i j of h′i which we}
\NormalTok{denote also by h′i . By [GM19, Proposition 2.18(b)], there is a finite subextension F /K of}
\NormalTok{K ′ /K such that for all i, j, there are PL functions hi j on X F := X ⊗KPF whose pull{-}backs}
\NormalTok{to X ′ agree with hi j . Then the pull{-}back of the Λ{-}PL function hi := j λi j hi j agrees with}
\NormalTok{h′i on X ′ . As seen at the beginning, it is enough to show that the pull{-}back of h to X F is}
\NormalTok{Λ{-}PL, so we may replace K by F and X by X F . Piecewise Λ{-}linearity is a local property, so}
\NormalTok{we show it in a neighbourhood of any y ∈ X . It follows from [Ber90, Proposition 1.3.5,}
\NormalTok{Corollary 1.3.6] that the structure map p : X ′ → X is open and surjective, and hence the}
\NormalTok{sets p(Vi′ ) form an open covering of X . For y ∈ X , we choose i such that y ∈ p(Vi′ ). For}
\NormalTok{all x ′ ∈ Vi′ and x = p(x ′ ), we have h(x) = h′ (x ′ ) = h′i (x ′ ) = hi (x), proving that h is Λ{-}PL}
\NormalTok{on the neighbourhood p(Vi′ ) of y.                                                                 }
\NormalTok{Corollary 5.13. Suppose that Λ is divisible. Let L/K be a (potentially infinite) normal}
\NormalTok{extension with automorphism group G := Aut(L/K). Let K ′ be the completion of L and let}
\NormalTok{X ′ := X ⊗}
\NormalTok{         b K K ′ , with structure morphism π: X ′ → X . Then h 7→ h◦π induces an isomorphism}
\NormalTok{          ∼}
\NormalTok{PLΛ (X ) −→ PLΛ (X ′ )G , where PLΛ (X ′ )G is the space of G{-}invariant Λ{-}PL functions on X ′ .}
\NormalTok{Proof. By [Ber90, Proposition 1.3.5, Corollary 1.3.6], we have X = X ′ /G, so the claim}
\NormalTok{follows from Corollary 5.12.                                                          }

\NormalTok{                 6. REDUCTION OF GERMS AND RESIDUES OF PL FUNCTIONS}
\NormalTok{   We will characterize semipositive and harmonic functions in 7 in the language of}
\NormalTok{Temkin’s theory of reductions of germs [Tem00], using residues of line bundles with}
\NormalTok{piecewise linear metrics from [CD12, Section 6]. In this section, we recall the theory}
\NormalTok{following quite closely [CD12, 6.1], and prove then the preliminary results we need.}
\NormalTok{                  FORMS ON BERKOVICH SPACES BASED ON HARMONIC TROPICALIZATIONS                                   37}

\NormalTok{  By a variety over a field k we mean an integral, separated k{-}scheme V of finite type.}
\NormalTok{The field of rational functions on V is denoted R(V ).}

\NormalTok{6.1 (Riemann–Zariski spaces). Let k be a field, let F /k be a field extension, and let PF /k}
\NormalTok{be the set of valuations on F that are trivial on k, up to equivalence. For ν ∈ PF /k we let}
\NormalTok{Oν ⊂ F denote the corresponding valuation ring. For a subset A ⊂ F we set PF /k \{A\} = \{ν ∈}
\NormalTok{PF /k | ν( f ) ≥ 0 for all f ∈ A\}. Subsets of the form PF /k \{ f1 , . . . , f n \} are said to be affine.}
\NormalTok{Note that PF /k = PF /k \{;\} is itself affine. We endow PF /k with the topology generated by}
\NormalTok{affine subsets. We equip PF /k with the sheaf of rings O : U 7→ \{ f | ν( f ) ≥ 0 for all ν ∈ U\}.}
\NormalTok{For ν ∈ PF /k the localization of O at ν is canonically identified with Oν , so PF /k is a}
\NormalTok{locally ringed space. It is shown in [Tem00, Corollary 1.3] (see [Tem04, Section 2] for}
\NormalTok{an erratum) that any affine open subset of PF /k is quasi{-}compact.}

\NormalTok{6.2 (Models). We keep the notation in 6.1. Let V be a variety equipped with a k{-}algebra}
\NormalTok{homomorphism from the function field R(V ) to F . We let PF /k \{V \} denote the set of all val{-}}
\NormalTok{uations ν ∈ PF /k such that Spec(F ) → V extends to a morphism Spec(Oν ) → V . Such an}
\NormalTok{extension is unique if it exists, so there is a map of sets πV : PF /k \{V \} → V sending ν to the}
\NormalTok{image of the closed point of Spec(Oν ) → V . It is easy to see that PF /k \{V \} is a finite union}
\NormalTok{of affine open subsets, so it is quasi{-}compact. The morphism πV : PF /k \{V \} → V is a surjec{-}}
\NormalTok{tive morphism of ringed spaces by [CD12, 6.1.1]. Conversely, if U is any quasi{-}compact}
\NormalTok{open subset of PF /k then U = PF /k \{V \} for some such V by [Tem00, Proposition 1.4]; we}
\NormalTok{call V a model of U.3}
\NormalTok{   Let V and V ′ be varieties equipped with k{-}algebra homomorphisms R(V ) → F and}
\NormalTok{R(V ′ ) → F . A dominant morphism V ′ → V , compatible with the morphisms R(V ) → F}
\NormalTok{and R(V ′ ) → F , is unique if it exists, in which case PF /k \{V ′ \} ⊂ PF /k \{V \}. In this case}
\NormalTok{we say that V ′ dominates V . We have equality PF /k \{V ′ \} = PF /k \{V \} if and only if V ′ →}
\NormalTok{V is proper by the valuative criterion of properness for integral schemes (see [EGAII,}
\NormalTok{Corollaire 7.3.10(ii)] and [Tem11, Lemma 3.2.1]).}

\NormalTok{6.3 (Functoriality). Let F ′ /k′ be an extension of F /k: that is, F ′ is an extension field}
\NormalTok{of F and k′ is a subfield of F ′ containing k. Restriction of valuations gives a canonical}
\NormalTok{map ϕ : PF ′ /k′ → PF /k . The inverse image of PF /k \{ f1 , . . . , f n \} is PF ′ /k′ \{ f1 , . . . , f n \}, so ϕ is}
\NormalTok{continuous. One checks that ϕ induces a morphism of locally ringed spaces. If k′ = k}
\NormalTok{then ϕ is surjective, and if F ′ = F then ϕ is injective.}
\NormalTok{   Let V (resp. V ′ ) be a k{-}variety (resp. k′ {-}variety) equipped with a k{-}algebra homomor{-}}
\NormalTok{phism R(V ) → F (resp. a k′ {-}algebra homomorphism R(V ′ ) → F ′ ). As in 6.2, a domi{-}}
\NormalTok{nant morphism V ′ → V of k{-}schemes, compatible with the morphisms R(V ) → F and}
\NormalTok{R(V ′ ) → F ′ , is unique if it exists, in which case ϕ(P F ′ /k′ \{V ′ \}) ⊂ P F /k \{V \}. In this case we}
\NormalTok{also say that V ′ dominates V .}
\end{Highlighting}
\end{Shaded}

\subsection{GJR, newline lines
1986--2307}\label{gjr-newline-lines-19862307}

\begin{Shaded}
\begin{Highlighting}[]

\NormalTok{Lemma 6.4. Let F ′ /k′ be an extension of F /k, and let ϕ : PF ′ /k′ → PF /k be the canonical}
\NormalTok{morphism. Let V be a k{-}variety equipped with a k{-}algebra homomorphism R(V ) → F , and}
\NormalTok{let V ′ be the closure of the image of the canonical k′ {-}morphism Spec(F ′ ) → V ⊗k k′ . Then}
\NormalTok{PF ′ /k′ \{V ′ \} = ϕ −1 (PF /k \{V \}).}
\NormalTok{   3}
\NormalTok{    In [CD12] these are called premodels.}
\NormalTok{38                             W. GUBLER, P. JELL, AND J. RABINOFF}

\NormalTok{Proof. This is similar to [Tem11, Lemma 3.2.1]. Let ν′ ∈ PF ′ /k′ and let ν = ϕ(ν′ ) = ν′ | F .}
\NormalTok{If ν′ ∈ PF ′ /k′ \{V ′ \} then we have a k′ {-}morphism Spec(Oν′ ) → V ′ extending Spec(F ′ ) →}
\NormalTok{V ′ . Let x ∈ V be the image of the closed point of Spec(Oν′ ) under the composition}
\NormalTok{Spec(Oν′ ) → V ′ → V . Then Spec(Oν′ ) → V factors through Spec(OV,x ), so we have a}
\NormalTok{k{-}algebra homomorphism OV,x → Oν′ . It follows that the image of OV,x is contained in the}
\NormalTok{intersection Oν = Oν′ ∩ F , so Spec(F ) → V extends to a morphism Spec(Oν ) → V , and}
\NormalTok{hence ν ∈ PF /k \{V \}. Conversely, suppose that ν ∈ PF /k \{V \}. Then Spec(Oν′ ) maps to V via}
\NormalTok{Spec(Oν′ ) → Spec(Oν ) → V , so Spec(F ′ ) → V ⊗k k′ extends to Spec(Oν′ ) → V ⊗k k′ . The}
\NormalTok{image of this morphism is contained in V ′ , so ν′ ∈ PF ′ /k′ \{V ′ \}.                      }
\NormalTok{  As a special case of the next lemma, we see that the poset of models of a quasi{-}compact}
\NormalTok{open subset U ⊂ PF /k is filtered.}
\NormalTok{Lemma 6.5. Let F ′ /k′ be an extension of F /k, and let ϕ : PF ′ /k′ → PF /k be the canonical}
\NormalTok{morphism. Let U ′ ⊂ PF ′ /k′ and U ⊂ PF /k be quasi{-}compact open subsets such that ϕ(U ′ ) ⊂ U.}
\NormalTok{Let V ′ and V be models of U ′ and U, respectively. Then there is a model W ′ of U ′ that}
\NormalTok{dominates both V and V ′ .}
\NormalTok{Proof. Let W be the closure of the image of the canonical k′ {-}morphism Spec(F ′ ) → V ⊗k k′ ,}
\NormalTok{as in Lemma 6.4. Then this lemma shows that W → V is dominant and ϕ −1 (U) =}
\NormalTok{PF ′ /k′ \{W \}, so we can replace V by W to assume k = k′ and F = F ′ . The given morphisms}
\NormalTok{Spec(F ) → V and Spec(F ) → V ′ determine a morphism Spec(F ) → V ×k V ′ ; let W ′ be the}
\NormalTok{closure of the image of this morphism. This is a variety. The generic point of W ′ maps to}
\NormalTok{the generic points of V and V ′ , so W ′ → V and W ′ → V ′ are dominant. We must show that}
\NormalTok{PF /k \{W ′ \} = U. For ν ∈ U we have unique morphisms Spec(O v ) → V and Spec(O v ) → V ′ ,}
\NormalTok{which induces a morphism p : Spec(O v ) → V ×k V ′ extending Spec(F ) → V ×k V ′ . The}
\NormalTok{image of p is contained in W ′ , so ν ∈ PF /k \{W ′ \}.                                     }
\NormalTok{6.6 (Invertible sheaves on Riemann–Zariski spaces). With the notation in 6.1, let U ⊂}
\NormalTok{PF /k be a quasi{-}compact open subset. Then U is a ringed space, so it makes sense to}
\NormalTok{discuss invertible sheaves on U. We let Pic(U) denote the group of isomorphism classes}
\NormalTok{of invertible sheaves, and for an additive subgroup Λ ⊂ R we set Pic(U)Λ = Pic(U) ⊗Z Λ.}
\NormalTok{   Let V be a model of U as in 6.2, with canonical morphism πV : U → V . If L V is an}
\NormalTok{invertible sheaf on V , then L := π∗V L V is an invertible sheaf on U. It is an important fact}
\NormalTok{that every invertible sheaf on U arises in this way. This is explained in [CD12, 6.1.3]. We}
\NormalTok{call (V, L V ) a model of (U, L).}
\NormalTok{   More generally, for an additive subgroup Λ of R we have a group homomorphism}
\NormalTok{Pn )Λ → Pic(U)Λ . Any element of Pic(U)Λ can be expressed as a finite sum L =}
\NormalTok{Pic(V}
\NormalTok{      λ L for λi ∈ Λ and L i ∈ Pic(U). Choosing a model for each L i and passing to a}
\NormalTok{   i=1 i i}
\NormalTok{dominating model, we see that any L ∈ Pic(U)Λ has the form π∗V L V for a model V of U}
\NormalTok{and L V ∈ Pic(V )Λ . We call such (V, L V ) a model of (U, L) as above.}
\NormalTok{6.7 (Nef line bundles). Let X be a separated k{-}scheme of finite type and let L be a line}
\NormalTok{bundle on X . For an integral, proper closed curve C ⊂ X we have a well{-}defined inter{-}}
\NormalTok{section number c1 (L).C := deg(L|C ) ∈ Z. We say that L is numerically effective or nef if}
\NormalTok{c1 (L).C ≥ 0 for all such curves C, and that L is numerically trivial if c1 (L).C = 0 for all C}
\NormalTok{(equivalently, if L and L −1 are both nef).}
\NormalTok{               FORMS ON BERKOVICH SPACES BASED ON HARMONIC TROPICALIZATIONS                     39}

\NormalTok{   More generally, let Λ ⊂ R be an additive subgroup and let C ⊂ X be an integral, proper}
\NormalTok{closed curve. We have an induced homomorphism L 7→ c1 (L).C : Pic(X )Λ → Λ, and we}
\NormalTok{say that L is nef (resp. numerically trivial) if c1 (L).C ≥ 0 (resp. c1 (L).C = 0) for all such C.}
\NormalTok{   The set of nef line bundles in Pic(X )Λ forms a monoid (resp. a cone if Λ = R), and the}
\NormalTok{set of numerically trivial line bundles form a subgroup (resp. an R{-}subspace).}
\NormalTok{  Here we list some classical properties of nefness. We sketch some proofs because in}
\NormalTok{the literature, the scheme X is generally assumed to be proper.}
\NormalTok{Lemma 6.8. Let X be a separated k{-}scheme of finite type and let L ∈ Pic(X )Λ .}
\NormalTok{    (1) The line bundle L is nef if and only if its image in Pic(X )R is nef.}
\NormalTok{    (2) Let k′ /k be a field extension and let X ′ = X ⊗k k′ , with structure morphism π: X ′ →}
\NormalTok{        X . Then L is nef if and only if π∗ L is nef.}
\NormalTok{    (3) Let f : X ′ → X be a morphism of separated, finite{-}type k{-}schemes. If L is nef then}
\NormalTok{        f ∗ L is nef, and the converse holds if f is proper and surjective.}
\NormalTok{    (4) Let ι : X red → X be the reduction map. Then L is nef if and only if ι ∗ L is nef.}
\NormalTok{The same statements hold with “nef” replaced by “numerically trivial.”}
\NormalTok{Proof. The last assertion holds because L is numerically trivial if and only if L and L −1}
\NormalTok{are nef. Assertion (1) is obvious because nefness is a numerical property, and (4) is clear}
\NormalTok{because X and X red have the same integral subcurves.}
\NormalTok{   (2) Since intersection numbers are compatible with base extension, it is clear that}
\NormalTok{if π∗ L is nef, then also L is nef. To prove the converse, we assume that L is nef. First}
\NormalTok{suppose that k′ /k is algebraic. Since every integral proper subcurve C ′ of X ′ is defined}
\NormalTok{over a finite subextension, we may assume that k′ /k is a finite extension. Then we may}
\NormalTok{view X ′ as a scheme of finite type over k and π as a finite morphism over k. Then π(C ′ )}
\NormalTok{is a proper subcurve of X and we get c1 (π∗ L).C ′ = c1 (L).π∗ (C ′ ) ≥ 0 by the projection}
\NormalTok{formula. This proves the algebraic case of (2), so we may assume k is algebraically}
\NormalTok{closed to prove the claim in general. Let C ′ again be an integral proper subcurve of X ′ .}
\NormalTok{Since C ′ is defined over a finitely generated subextension, we may assume that k′ is a}
\NormalTok{finitely generated field extension of k. Then k′ is the function field of a variety V over k.}
\NormalTok{Let πV : X ×k V → V and πX : X ×k V → X be the structure morphisms. The morphism}
\NormalTok{π: X ′ → X is the generic fiber of πV . Since C ′ is projective and flat over k′ , the spreading}
\NormalTok{out property shows that there is a non{-}empty open subset U of V such that the closure C ′}
\NormalTok{of C ′ in π−1}
\NormalTok{            V}
\NormalTok{              (U) is flat and projective over U. By [Kol96, Proposition VI.2.9], the degree}
\NormalTok{of the fiber Cu′ with respect to π∗X (L) is independent of the choice of u ∈ U. Let η be the}
\NormalTok{generic point of U and let u be a closed point of U. Since Cu′ ∼      = C for a proper closed}
\NormalTok{subcurve C of X , the claim follows from}
\NormalTok{            c1 (π∗ L).C ′ = degπ∗V (L) (Cη′ ) = degπ∗V (L) (Cu′ ) = deg L (C) = c1 (L).C ≥ 0.}

\NormalTok{   (3) If L is nef then f ∗ L is nef because if C ′ ⊂ X ′ is an integral proper k{-}curve then}
\NormalTok{f (C ′ ) ⊂ X is integral and proper of dimension ≤ 1, so the claim follows from the projec{-}}
\NormalTok{tion formula. Conversely, if f ∗ L is nef and if f is a proper surjective map, then L is nef}
\NormalTok{because every integral proper k{-}curve C ⊂ X is the image of an integral proper k{-}curve}
\NormalTok{C ′ ⊂ X ′ (take the closure of a closed point in the fiber over the generic point of C). }
\NormalTok{40                              W. GUBLER, P. JELL, AND J. RABINOFF}

\NormalTok{Definition 6.9. Let U be a quasi{-}compact open subset of PF /k , let Λ ⊂ R be an additive}
\NormalTok{subgroup, and let L ∈ Pic(U)Λ . We say that L is nef if there exists a model (V, L V ) of}
\NormalTok{(U, L) with L V nef. We say that L is numerically trivial if both L and L −1 are nef.}
\NormalTok{6.10. The set of nef elements of Pic(U)Λ forms a monoid (pass to a dominating model}
\NormalTok{and use Lemma 6.8(3)), resp. a cone if Λ = R; the set of numerically trivial elements}
\NormalTok{forms a group (resp. an R{-}subspace). An element of Pic(U)Λ is nef if and only if its image}
\NormalTok{in Pic(U)R is nef.}
\NormalTok{     The following lemma shows that nefness can be checked on any model.}
\NormalTok{Lemma 6.11. Let U be a quasi{-}compact open subset of PF /k and let L ∈ Pic(U)Λ be nef. If}
\NormalTok{(V, L V ) is a model of L, then L V is nef.}
\NormalTok{Proof. By definition, there is a model (V ′ , L V ′ ) of L such that L V ′ is nef. Let W be a model}
\NormalTok{dominating both V and V ′ as in Lemma 6.5, and let LW and LW         ′}
\NormalTok{                                                                        be the pullbacks of L V and}
\NormalTok{L V ′ to W , respectively. It is shown in [CD12, 6.1.3] that Pic(U) is the inductive limit of}
\NormalTok{the groups Pic(V ) over all models V of U. Since tensor products commute with inductive}
\NormalTok{limits, we have Pic(U)Λ = lim Pic(V )Λ . Hence we can replace W by a larger model such}
\NormalTok{               ′}
\NormalTok{                               −→}
\NormalTok{that LW = LW     in Pic(W ). Then LW is nef, so L V ′ is nef because W → V ′ is proper, using}
\NormalTok{Lemma 6.8(3) twice.                                                                               }
\NormalTok{     The next lemma is an analogue of Lemma 6.8.}
\NormalTok{Lemma 6.12. Let F ′ /k′ be an extension of F /k, and let ϕ : PF ′ /k′ → PF /k be the canonical}
\NormalTok{morphism. Let U ′ ⊂ PF ′ /k′ and U ⊂ PF /k be quasi{-}compact open subsets such that ϕ(U ′ ) ⊂ U.}
\NormalTok{If L ∈ Pic(U)Λ is nef, then ϕ ∗ L ∈ Pic(U ′ )Λ is nef. The converse is true if U ′ = ϕ −1(U), and}
\NormalTok{if one of the following conditions holds:}
\NormalTok{    (1) k′ /k is algebraic, or}
\NormalTok{    (2) tr.deg(F /k) \textless{} ∞ and tr.deg(k′ F /k′ ) = tr.deg(F /k).}
\NormalTok{The same statements hold with “nef” replaced by “numerically trivial.”}
\NormalTok{  Here k′ F denotes the smallest subfield of F ′ containing both k′ and F . It is always true}
\NormalTok{that tr.deg(k′ F /k′ ) ≤ tr.deg(F /k).}
\NormalTok{Proof. Suppose that L is nef. Let (V, L V ) be a model of (U, L), and let W be the closure}
\NormalTok{of the image of the canonical k′ {-}morphism Spec(F ′ ) → V ⊗k k′ , as in Lemma 6.4. The}
\NormalTok{pullback of L to W is nef by Lemma 6.8(2,3), so we may replace V by W to assume}
\NormalTok{k = k′ . By Lemma 6.5, there is a model V ′ of U ′ dominating V ; let f : V ′ → V be the}
\NormalTok{dominating morphism. Then f ∗ L V is nef by Lemma 6.8(3), so ϕ ∗ L is nef.}
\NormalTok{   Now suppose that ϕ ∗ L is nef and that U ′ = ϕ −1 (U). First we treat the case when k′ is}
\NormalTok{algebraic over k. If k′ = k then any model (V, L V ) of (U, L) is also model of (U ′ , ϕ ∗ L), so}
\NormalTok{L V is nef by Lemma 6.11. In general, the map PF ′ /k′ → PF /k factors as PF ′ /k′ → PF ′ /k →}
\NormalTok{PF /k , so by the previous sentence we may assume F = F ′ . Note that PF /k = PF /k′ and}
\NormalTok{U = U ′ . Let (V, L V ) be a model of L and let V ′ be a model of U ′ . We may assume that V ′}
\NormalTok{dominates V by Lemma 6.5; let f : V ′ → V be the dominating morphism. Then (V ′ , f ∗ L V )}
\NormalTok{is a model of (U ′ , ϕ ∗ L), so f ∗ L V is nef. The model V ′ is defined over a finite extension of}
\NormalTok{k: that is, there exists a finite extension k′′ /k contained in k′ and a variety V ′′ over k′′}
\NormalTok{                FORMS ON BERKOVICH SPACES BASED ON HARMONIC TROPICALIZATIONS                           41}

\NormalTok{such that V ′ = V ′′ ⊗k′′ k′ . After enlarging k′′ , we may assume that V ′′ also dominates V ;}
\NormalTok{let f ′ : V ′′ → V be the dominating morphism. By Lemma 6.8(2), the line bundle f ′∗ L is}
\NormalTok{nef. Replacing k′ with k′′ and V ′ with V ′′ , we may assume that [k′ : k] \textless{} ∞. In this case,}
\NormalTok{the variety V ′ is also a model of U. Any proper k{-}curve in V ′ is also a proper k′ {-}curve, so}
\NormalTok{f ∗ L V is nef when we consider V ′ as a k{-}variety. It follows from Lemma 6.11 that L is nef.}
\NormalTok{   Now assume that d = tr.deg(F /k) \textless{} ∞, and that tr.deg(k′ F /k′ ) = d. First we reduce}
\NormalTok{to the case when k is algebraically closed. Let F ′′ be a field containing F ′ and an algebraic}
\NormalTok{closure k of k. Pulling back via PF ′′ /k′ → PF ′ /k′ , we may replace F ′ with F ′′ to assume}
\NormalTok{k ⊂ F ′ . We have tr.deg(kF /k) = tr.deg(F /k) = d, and}

\NormalTok{                  tr.deg(k′ kF /k′ k) = tr.deg(k′ kF /k′ ) = tr.deg(F k′ /k′ ) = d.}

\NormalTok{Applying case (1) to the extension kF /k of F /k and using the above displayed equation,}
\NormalTok{we may replace k by k, F by kF , and k′ by k′ k to assume k = k.}
\NormalTok{   Choose a model (V, L V ) of (U, L) with dim(V ) = d. Since k is algebraically closed, the}
\NormalTok{scheme V ′ := V ⊗k k′ is a variety. As F /R(V ) is algebraic, we have tr.deg(k′ R(V )/k′ ) =}
\NormalTok{tr.deg(k′ F /k′ ) = d. Let x be the image of Spec(F ′ ) → V ′ , so x maps to the generic point}
\NormalTok{of V . The residue field κ(x) contains both k′ and R(V ), so tr.deg(κ(x)/k′ ) = d. It follows}
\NormalTok{that x is the generic point of V ′ . By Lemma 6.4, the variety V ′ is a model of U ′ , so ϕ ∗ L V}
\NormalTok{is nef by Lemma 6.11, hence L V is nef by Lemma 6.8(3).                                        }
\NormalTok{6.13 (Temkin’s reductions of germs). Let X be a good strictly K{-}analytic space and let}
\NormalTok{x ∈ X . Temkin’s reduction red(X , x) of the germ (X , x) is an affine open subset of}
\NormalTok{PHf(x)/ek , where H    f(x) is the reduction of the completed residue field H (x) at x as in 2.3.}
\NormalTok{The reduction red(X , x) can be constructed as follows.4 Let U be a compact strictly}
\NormalTok{analytic neighbourhood of x in X and let U be a formal K ◦ {-}model of U. The closure}
\NormalTok{VU,x of the reduction redU (x) is variety over ek by [BL93a, Proposition 4.7], and there}
\NormalTok{is a canonical homomorphism R(VU,x ) → H             f(x) by [Ber90, 2.4]. We define red(X , x)}
\NormalTok{to be PHf(x)/ek \{VU,x \}. Then red(X , x) is an affine (open) subset of PHf(x)/ek by [Tem00,}
\NormalTok{Lemma 2.2]. Let πU,x : red(X , x) → VU,x be the canonical morphism of 6.2}
\NormalTok{   Reduction of germs (X , x) 7→ red(X , x) is functorial in (X , x). It also respects extension}
\NormalTok{of the ground field, in the following sense. Let K ′ /K be a non{-}Archimedean extension}
\NormalTok{field, let X ′ = X ⊗   b K K ′ with structure morphism π: X ′ → X , and let x ′ ∈ X ′ and x = π(x ′ ).}
\NormalTok{Let ϕ : PHf(x ′ )/Ke ′ → PHf(x)/Ke be the natural morphism. Then ϕ −1 (red(X , x)) = red(X ′ , x ′ ).}
\NormalTok{See [Tem00, Lemma 4.2] (and [Tem04, Section 5] for a correction).}
\NormalTok{   We will need the fact that models of the form VU,x as in 6.13 are cofinal in the system}
\NormalTok{of all models of red(X , x).}
\NormalTok{Lemma 6.14. Let X be a good strictly K{-}analytic space, let x ∈ X , and let V be a model of}
\NormalTok{red(X , x). Then there exists a compact strictly analytic neighbourhood U of x and a formal}
\NormalTok{K ◦ {-}model U of U such that VU,x dominates V .}
\NormalTok{  4}
\NormalTok{    Reductions of germs are much more subtle when X is not good or strictly analytic, and a major moti{-}}
\NormalTok{vation of Temkin’s work is to understand goodness and strictness by studying such reductions. In our case,}
\NormalTok{the construction is relatively simple.}
\NormalTok{42                                 W. GUBLER, P. JELL, AND J. RABINOFF}
\NormalTok{                                                          S}
\NormalTok{Proof. Choose a finite affine open covering V = Vi . Then Vi = Spec(K                e [ fei j : j ∈ Ji ])}
\NormalTok{for finitely many fei j ∈ H  f(x). We pick lifts f i j ∈ OX ,x . Let U be a strictly affinoid neigh{-}}
\NormalTok{bourhood of x with f i j ∈ O (U), and let Ui := \{x ∈ U | | f i j (x)| ≤ 1, j ∈   SJi \}. By [Tem00,}
\NormalTok{Proposition 2.3(ii,iii)], the reductionsS   red(Ui , x) cover red(X , x), so ( Ui , x) = (X , x).}
\NormalTok{We may shrink U to assume that U = Ui is a (strictly affinoid) neighbourhood of x.}
\NormalTok{   We claim that we can find a formal K ◦ {-}model U of U that is covered by formal opens}
\NormalTok{Ui with (Ui )η = Ui and f i j ∈ O (Ui ). Write Ui = M (Ai ) for a strictly affinoid algebra Ai .}
\NormalTok{Choose a surjection from a Tate algebra T onto Ai such that each f i j is the image of a}
\NormalTok{generator. The image Ai of T ◦ in Ai is a flat K ◦ {-}algebra topologically of finite presentation}
\NormalTok{containing f i j ; hence Ui := Spf(Ai ) is an admissible formal model of Ui with f i j ∈ O (Ui ).}
\NormalTok{By [Bos14, Lemma 8.4.5], there exists an admissible formal K ◦ {-}model U of U and a}
\NormalTok{covering by formal open subschemes U′i ⊂ U such that (U′i )η = Ui . By Lemma 8.4.4(d)}
\NormalTok{of loc. cit., there is a formal K ◦ {-}model U′′i of Ui that is an admissible formal blowup of}
\NormalTok{both Ui and U′i . By Proposition 8.2.13 of loc. cit., there is an admissible formal blowup}
\NormalTok{ϕ : U′ → U such that ϕ −1 (U′i ) dominates U′′i . Replacing U by U′ and Ui by ϕ −1 (U′i ) proves}
\NormalTok{the claim.}
\NormalTok{   Consider the closure V ′ := VU,x of redU (x) in Us . By construction the function field}
\NormalTok{R(V ′ ) contains all fei j , so R(V ) ⊂ R(V ′ ) and hence there is a dominant rational map}
\NormalTok{V ′ → V . Moreover, for each i we have O (Vi ) ⊂ O ((Ui )s ∩ V ′ ) because f i j ∈ O (Ui ), so}
\NormalTok{(Ui )s ∩ V ′ → Vi is a morphism, and hence V ′ → V is a morphism.                                       }
\NormalTok{     We will also use the following theorem of Temkin.}
\NormalTok{Theorem 6.15 ([Tem00, Theorem 4.1]). Let f : X ′ → X be a morphism of good strictly}
\NormalTok{K{-}analytic spaces, let x ′ ∈ X ′ , and let x = f (x ′ ). Let ϕ : PHf(x ′ )/Ke → PHf(x)/Ke be the natural}
\NormalTok{morphism. Then x ∈ Int(X ′ /X ) if and only if ϕ −1 (red(X , x)) = red(X ′ , x ′ ).}
\NormalTok{6.16 (Residues of PL functions). Let h: X → R be a PL function on a good strictly K{-}}
\NormalTok{analytic space X . For x ∈ X , Chambert–Loir and Ducros define a canonical invertible}
\NormalTok{sheaf Lh (x) on the ringed space red(X , x), which is called the residue of h at x. It can be}
\NormalTok{described as follows, using [CD12, Exemple 6.6.8]. Let U be a compact strictly analytic}
\NormalTok{neighbourhood of x. By 5.6, there exists a formal K ◦ {-}model (U, L) of (U, OX |U ) such that}
\NormalTok{h = − log k1kL on U. This L restricts to an invertible sheaf LU,x on VU,x , and we get an}
\NormalTok{invertible sheaf Lh (x) = π∗U,x LU,x on red(X , x).}
\NormalTok{  Let PL(X , x) denote the group of germs of PL functions on X at x; this is intrinsic to}
\NormalTok{the germ (X , x). Applying the above to germs of PL functions at x, we obtain a homo{-}}
\NormalTok{morphism L : PL(X , x) → Pic(red(X , x)), denoted h 7→ Lh (x).}
\NormalTok{Lemma 6.17. Let (X , x) be a good strictly K{-}analytic germ. If h ∈ PL(X , x) has the form}
\NormalTok{log | f | for an invertible function f , then Lh (x) = 0 in Pic(red(X , x)).}
\NormalTok{Proof. Choose a compact strictly analytic space X and a point x ∈ X representing (X , x)}
\NormalTok{such that there exists an invertible function f on X with h = log | f |. Since X is compact,}
\NormalTok{there exists a formal K ◦ {-}model X of X . Let L = OX · f , regarded as a subsheaf of the sheaf}
\NormalTok{U 7→ Γ (Uη , OX ) on X. Then k1kL = | f |−1 , and hence h = − log k1kL . Since L is trivial, so}
\NormalTok{is LX,x , so we conclude that Lh (x) = 0.                                                    }
\NormalTok{               FORMS ON BERKOVICH SPACES BASED ON HARMONIC TROPICALIZATIONS                      43}

\NormalTok{   Now we want to extend 6.16 to Λ{-}PL functions. Let PLΛ (X , x) denote the group of}
\NormalTok{germs of Λ{-}PL functions on (X , x). WePcall h ∈ PLΛ (X , x) Λ{-}linear if there are f1 , . . . , f r ∈}
\NormalTok{                                             r}
\NormalTok{OX×,x and λ1 , . . . , λ r ∈ Λ such that h = j=1 λ j log | f j | in a neighbourhood of x.}

\NormalTok{Lemma 6.18. Suppose that Λ = Z or that Λ is divisible. Then there is a unique homomor{-}}
\NormalTok{phism}
\NormalTok{                          h 7→ Lh (x): PLΛ (X , x) −→ Pic(red(X , x))Λ}
\NormalTok{                        Pm                                                           Pm}
\NormalTok{such that, if h = i=1 λi hi with λi ∈ Λ and hi ∈ PL(X , x), then Lh (x) = i=1 λi Lhi (x).}
\NormalTok{The kernel of this homomorphism consists precisely of the Λ{-}linear functions in PLΛ (X , x).}
\NormalTok{                                              Pm                                   Pm}
\NormalTok{Proof. We have to show that Lh (x) := i=1 λi Lhi (x) is well{-}defined. If i=1 λi hi = 0 for}
\NormalTok{                                                     Pm}
\NormalTok{λi ∈ Λ and hi ∈ PL(X , x), then we need i=1 λi Lhi (x) = 0. By Lemma 5.10, we may}
\NormalTok{assume \{λ1 , . . . , λm \} is linearly independent over Q. Choose a compact strictly analytic}
\NormalTok{space X and a point x ∈ X representing (X , x) such that each P         hi extends to a PL function}
\NormalTok{                                                                           m}
\NormalTok{hi : X → R. Shrinking X , we may assume X is connected and i=1 λi hi = 0 on X .}
\NormalTok{   We claim that each hi is constant on X . Using [Ber93, Lemma 1.6.2], we see that X}
\NormalTok{is G{-}connected and hence the claim is G{-}local on X . It follows from [BGR84, Proposi{-}}
\NormalTok{tion 9.1.4/8] that X is G{-}locally connected, so we may assume that hi = − log | f i | for}
\NormalTok{an invertible analytic function f i on a connected strictly analytic domain          PU  ⊂ X . Let}
\NormalTok{           m                                      an                                   m}
\NormalTok{T = Gm and let ϕ = ( f1 , . . . , f m ): U → T , so ϕtrop = (h1 , . . . , hm ). Since i=1 λi hi = 0,}
\NormalTok{the tropical variety ϕtrop (U) is contained in the hyperplane H = \{(x 1, . . . , x m ) ∈ Rm |}
\NormalTok{λ1 x 1 + · · · + λm x m = 0\}. If ϕtrop is not constant, then ϕtrop (U) has positive dimension}
\NormalTok{because U is connected. It follows from 3.2 that ϕtrop (U) contains a (Z, Γ ){-}polytope of}
\NormalTok{dimension one, which is a line segment with rational slopes. The tangent space of this}
\NormalTok{segment contains a non{-}zero vector in H of the form (a1 , . . . , an ) ∈ Qn , which contradicts}
\NormalTok{the linear independence of \{λ1 , . . . , λn \} over Q. This proves the claim; moreover,p        since}
\NormalTok{ϕtrop (U) is a (Z, Γ ){-}point in Rn , this shows that hi is a constant function γi for γi ∈ Γ .}
\NormalTok{   Choose a positive integer N such that N γi ∈ Γ for all i, so that N γi = − log |βi | for}
\NormalTok{βi ∈ K × . Then N hi = −N log |βi |, so L N hi (x) = 0 by Lemma 6.17. Thus N Lh (x) =}
\NormalTok{Pm}
\NormalTok{         λ L (x) = 0, so Lh (x) = 0 because Λ is divisible.}
\NormalTok{   i=1 i N hi}
\NormalTok{   To prove the last claim, it is obvious from the definition of the residue Lh (x) that Λ{-}}
\NormalTok{linear functions h ∈ PLΛ (X , x) are in the kernel of the homomorphism. Conversely, let}
\NormalTok{h ∈ PLΛ (X , x) with Lh (x) = 0. We note first that for Λ = Z, the claim follows from [CD12,}
\NormalTok{6.6.2]. Indeed, Chambert{-}Loir and Ducros show that h induces a metric on OX which is}
\NormalTok{isometric to the trivial metric on a neighbourhood U of x and hence there is f ∈ O (U)×}
\NormalTok{such that h = log | f | on U. It is clear that the case Λ = Z readily yields the case Λ = Q}
\NormalTok{by clearing denominators.}
\NormalTok{   For Λ divisible, we reduce Pr to Λ = Q as follows. There are h1 , . . . , h r ∈ PLQ (X , x) and}
\NormalTok{λ1 , . . . , λ r ∈ Λ with h = i=1 λi hi . By Lemma 5.10 again, we may assume that λ1 , . . . , λ r}
\NormalTok{are Q{-}linearly independent. It is enough to prove that Lhi (x) = 0 for all i ∈ \{1, . . . , r\} as}
\NormalTok{then the special case Λ = Q above shows that hi is Q{-}linear, and hence that h is Λ{-}linear.}
\NormalTok{We argue by contradiction. Then there is a Q{-}linear form ℓ: Pic(red(X , x))Q → Q with}
\NormalTok{ℓ(Lhi (x)) 6= 0 for at least one i ∈ \{1, . . . , r\}. This extends by linearity to a homomorphism}
\NormalTok{44                                 W. GUBLER, P. JELL, AND J. RABINOFF}

\NormalTok{ℓ: Pic(red(X , x))Λ → R, and}
\NormalTok{                                                     X}
\NormalTok{                                                     r}
\NormalTok{                                   0 = ℓ(Lh (x)) =          λi ℓ(Lhi (x))}
\NormalTok{                                                      i=1}

\NormalTok{is a non{-}trivial Q{-}linear combination of λ1 , . . . , λ r which is a contradiction.                     }
\NormalTok{6.19 (Functoriality of residues). Let f : (X ′ , x ′ ) → (X , x) be a morphism of good strictly}
\NormalTok{K{-}analytic germs. For h ∈ PLΛ (X , x), we get h′ = h ◦ f ∈ PLΛ (X ′ , x ′ ). We have an induced}
\NormalTok{morphism fe : red(X ′ , x ′ ) → red(X , x), and Lh′ (x ′ ) = fe∗ Lh (x) by [CD12, 6.6.7]. A similar}
\NormalTok{statement holds for extension of scalars.}

\NormalTok{                          7. SEMIPOSITIVITY AND HARMONIC FUNCTIONS}
\NormalTok{   In this section, we study positivity conditions on piecewise linear functions. We will}
\NormalTok{introduce the crucial notion of a harmonic function, on which our theory of differential}
\NormalTok{forms rests. We consider a good strictly K{-}analytic space X and an additive subgroup}
\NormalTok{Λ ⊂ R. We will always assume that Λ = Z or that Λ is divisible.}
\NormalTok{   We can characterize semipositivity either in terms of formal geometry or residues of}
\NormalTok{Λ{-}PL functions, as elaborated in 6. A line bundle L on an admissible formal scheme X}
\NormalTok{is called nef (resp. numerically trivial) if its restriction to Xs has the same property. In}
\NormalTok{other words, the line bundle L is nef if the degree of L|C is nonnegative for every closed}
\NormalTok{integral curve C ⊂ Xs which is proper over K   e .5}
\NormalTok{   Let X be a formal K ◦ {-}model of a paracompact good strictly K{-}analytic space X . Recall}
\NormalTok{from 5.5 that M (X) denotes the group of models of OX , and that we have defined a}
\NormalTok{homomorphism h 7→ hL : M (X)Λ → PL(X )Λ . We will also use the residue Lh (x) of h at}
\NormalTok{x ∈ X from 6.16.}
\NormalTok{Proposition 7.1. Let h: X → R be a Λ{-}PL function and x ∈ X . The following are equivalent:}
\NormalTok{         (1) There exists a paracompact neighbourhood U of x in X , a formal K ◦ {-}model U of U,}
\NormalTok{             and a nef (resp. numerically trivial) element L ∈ M (U)Λ , such that h|U = hL .}
\NormalTok{         (2) The residue Lh (x) ∈ Pic(red(X , x))Λ is nef (resp. numerically trivial).}
\NormalTok{Proof. The first assertion implies the second by the description of Lh (x) in 6.16. We will}
\NormalTok{prove the nef case, as the numerically trivial case is identical. The following argument}
\NormalTok{is similar to [GK19, Proposition 6.5]. Suppose that Lh (x) is nef. By Corollary 5.9, there}
\NormalTok{exists a compact strictly analytic neighbourhood U of x, a formal K ◦ {-}model U of U, and an}
\NormalTok{element L ∈ M (U)Λ , such that h = hL . As in 6.16 we let VU,x be the closure of redU (x) and}
\NormalTok{we let LU,x be the restriction of L to VU,x . Then (VU,x , LU,x ) is a model of (red(X , x), Lh (x)),}
\end{Highlighting}
\end{Shaded}

\subsection{GJR, newline lines
2308--2435}\label{gjr-newline-lines-23082435}

\begin{Shaded}
\begin{Highlighting}[]
\NormalTok{so LU,x is nef by Lemma 6.11. Let W be a compact, strictly analytic neighbourhood of x}
\NormalTok{contained in the open set red−1 U}
\NormalTok{                                    (VU,x ). By [BL93a, Lemma 4.4], there exists an admissible}
\NormalTok{formal blowup f : U′ → U such that W = Wη for a formal open subset W ⊂ U′ . Note that}
\NormalTok{h|W = h f ∗ L|W . We claim that f ∗ L|W is nef.}
\NormalTok{     5}
\NormalTok{    If X is a proper K ◦ {-}model of a proper variety X over K and L is a line bundle on X with L = L |X}
\NormalTok{such that the restriction of L to the special fiber of X is nef, then L is a nef line bundle on X by [GM19,}
\NormalTok{4.8].}
\NormalTok{              FORMS ON BERKOVICH SPACES BASED ON HARMONIC TROPICALIZATIONS                 45}

\NormalTok{  By functoriality of reductions, the set{-}theoretic image of W under f is contained in}
\NormalTok{VU,x . In particular, if C is any integral, proper subcurve of Ws , then f (C) ⊂ VU,x . Since}
\NormalTok{C is reduced, the schematic image of C is contained in VU,x , so we have a commutative}
\NormalTok{square of morphisms of ek{-}schemes:}
\NormalTok{                                            C    /W}
\NormalTok{                                                       s}
\NormalTok{                                        f              f}
                                                  
\NormalTok{                                        VU,x     /U}
\NormalTok{                                                       s}

\NormalTok{Since LU,x is nef, we have f ∗ LU,x .C ≥ 0 by the projection formula.                       }
\NormalTok{Definition 7.2. A Λ{-}PL function h: X → R is called semipositive at a point x ∈ X if it}
\NormalTok{satisfies the equivalent conditions of Proposition 7.1, and it is harmonic at x if both h}
\NormalTok{and −h are semipositive at x. We say that h is semipositive if it is semipositive at all}
\NormalTok{points of X .}
\NormalTok{   A harmonic function on X is an R{-}PL function h: X → R which is harmonic at all points}
\NormalTok{of X . A Λ{-}PL harmonic function is called Λ{-}harmonic.}
\NormalTok{7.3. Some remarks:}
\NormalTok{   (1) A Λ{-}PL function h is harmonic at x ∈ X if and only if Lh (x) is numerically trivial.}
\NormalTok{       Equivalently, there exists a paracompact neighbourhood U of x, a formal K ◦ {-}}
\NormalTok{       model U of U, and a numerically trivial element L ∈ M (U)Λ , such that h = hL .}
\NormalTok{   (2) Semipositivity and harmonicity are insensitive to the group Λ, in the following}
\NormalTok{       sense: if h ∈ PLΛ (X ), then h is semipositive at a point x if and only if its image in}
\NormalTok{       PLR (X ) is semipositive at x.}
\NormalTok{   (3) Semipositivity and harmonicity of h at x is intrinsic to the germ of h in PLΛ (X , x).}
\NormalTok{       In particular, the semipositive (resp. harmonic) functions on X form a sheaf in}
\NormalTok{       the analytic topology. (They do not form a sheaf in the G{-}topology).}
\NormalTok{   (4) It follows from Proposition 7.1(1) that the locus of all points x ∈ X at which h is}
\NormalTok{       semipositive (resp. harmonic) is open.}
\NormalTok{   (5) The set of semipositive functions at a point forms a monoid (resp. a cone if}
\NormalTok{       Λ = R), and the set of harmonic functions at a point forms a group (resp. an}
\NormalTok{       R{-}subspace).}
\NormalTok{   (6) If h = log | f | for an invertible analytic function f in a neighbourhood of x, then h}
\NormalTok{       is harmonic at x by Lemma 6.17 (the converse does not hold: see Remark 15.10).}
\NormalTok{   (7) Harmonic functions generalize smooth functions in the same way that numerical}
\NormalTok{       equivalence generalizes rational equivalence in algebraic geometry.}
\NormalTok{Lemma 7.4. Suppose that X is paracompact. Let X be a formal K ◦ {-}model of X and let L ∈}
\NormalTok{M (X)Λ . Then hL is semipositive (resp. harmonic) if and only if L is nef (resp. numerically}
\NormalTok{trivial).}
\NormalTok{Proof. The harmonic case follows from the semipositive case. If L is nef then hL is semi{-}}
\NormalTok{positive by definition. Suppose then that h := hL is semipositive. Let C ⊂ Xs be an}
\NormalTok{integral, proper closed curve. By surjectivity of redX , there exists a point x ∈ X reducing}
\NormalTok{to the generic point of C. Then (C, L|C ) is a model of (red(X , x), Lh (x)), so L|C is nef by}
\NormalTok{Lemma 6.11. In particular, c1 (L).C ≥ 0, so L is nef.                                       }
\NormalTok{46                               W. GUBLER, P. JELL, AND J. RABINOFF}

\NormalTok{  Combined with 5.6, we see that a PL function h: X → R on a paracompact space X is}
\NormalTok{semipositive (resp. harmonic) if and only if there exists a formal K ◦ {-}model X of X and a}
\NormalTok{nef (resp. numerically trivial) line bundle L ∈ M (X) such that h = hL = − log k1kL .}
\NormalTok{  We have the following alternative characterization of harmonic functions. We do not}
\NormalTok{know if the same characterization holds for semipositive functions.}
\NormalTok{Proposition 7.5. Let h: X → R be a Λ{-}PL function and let x ∈ X . Then h is harmonic Pm at x}
\NormalTok{if and only if there exists a strictly analytic neighbourhood U of x such that h|U = i=1 λi hi}
\NormalTok{for Z{-}harmonic functions h1 , . . . , hm on U and λ1 , . . . , λm ∈ Λ.}
\NormalTok{Proof. The proposition is vacuous when Λ = Z, so assume Λ is divisible. The “if” direc{-}}
\NormalTok{tion is trivial. Suppose then thatPmh is harmonic at x. By Proposition 5.8, after shrinking}
\NormalTok{X , we may assume that h = i=1 λi hi for λ1 , . . . , λm ∈ Λ and h1 , . . . , hm ∈ PL(X ). By}
\NormalTok{Lemma 5.10, we may assume that   Pm λ1 , . . . , λm are Q{-}linearly independent. By Lemma 6.18,}
\NormalTok{the residue Lh (x) is equal to i=1 λi Lhi (x). Let V be a model for red(X , x) in the sense}
\NormalTok{of 6.2 such that each Lhi (x) is the pullback of a line bundle L V,i on V . Let C ⊂ V be an inte{-}}
\NormalTok{                                                                            Pm}
\NormalTok{gral, proper closed curve. Since Lh (x) is numerically trivial, we have i=1 λi c1 (L V,i (x)).C =}
\NormalTok{0. But c1 (L V,i (x)).C ∈ Z for each i, so by linear independence, we have c1 (L V,i (x)).C = 0}
\NormalTok{for each i. Thus L V,i is numerically trivial, so hi is harmonic.                              }
\NormalTok{     Now we gather some functoriality properties of semipositive and harmonic functions.}
\NormalTok{Proposition 7.6. Let X be a good strictly K{-}analytic space and let h ∈ PLΛ (X ).}
\NormalTok{   Let f : X ′ → X be a morphism of good strictly K{-}analytic spaces and let h′ = h ◦ f .}
\NormalTok{     (1) For x ′ ∈ X ′ , let x = f (x ′ ). If h is semipositive at x, then h′ is semipositive at x ′ .}
\NormalTok{           The converse holds if x ′ ∈ Int(X ′ /X ).}
\NormalTok{    (1′ ) If h is semipositive, then h′ is semipositive. The converse holds if f is proper and}
\NormalTok{           surjective.}
\NormalTok{Let K /K be a non{-}Archimedean field extension, let X ′ = X ⊗}
\NormalTok{       ′                                                             b K K ′ , let π: X ′ → X be the}
\NormalTok{structure morphism, let x ′ ∈ X ′ , let x = π(x ′ ), and let h′ = h ◦ π.}
\NormalTok{     (2) If h is semipositive at x, then h′ is semipositive at x ′ . The converse holds if we have}
\NormalTok{           tr.deg(Ke ′H  f(x)/K}
\NormalTok{                              e ′ ) = tr.deg(Hf(x)/K e ), where K}
\NormalTok{                                                                e ′H}
\NormalTok{                                                                   f(x) is the smallest subfield of}
\NormalTok{           Hf(x ) containing H}
\NormalTok{                 ′                f(x) and K e.′}

\NormalTok{    (2 ) The function h is semipositive if and only if h′ is semipositive.}
\NormalTok{       ′}

\NormalTok{Let ι : X red → X be the inclusion of the reduction of X and let h′ = h ◦ ι.}
\NormalTok{     (3) For x ′ ∈ X red , the function h is semipositive at x = ι(x ′ ) if and only if h′ is semipos{-}}
\NormalTok{           itive at x ′ .}
\NormalTok{    (3 ) The function h is semipositive if and only if h′ is semipositive.}
\NormalTok{       ′}

\NormalTok{The statements (1)–(3) and (1′ )–(3′ ) hold with “semipositive” replaced by “harmonic.”}
\NormalTok{Proof. In each case, harmonicity follows easily from semipositivity. In all cases, the func{-}}
\NormalTok{tion h′ is Λ{-}PL by Lemma 5.4.}
\NormalTok{   We prove first (1). If h is semipositive at x then h′ is semipositive at x ′ by Lemma 6.12}
\NormalTok{and functoriality of reduction. The converse holds if x ∈ Int(X ′ /X ) by Theorem 6.15}
\NormalTok{and Lemma 6.12(1), using functoriality of residues 6.19. Assertion (1′ ) follows because}
\NormalTok{a proper morphism has no boundary.}
\NormalTok{                 FORMS ON BERKOVICH SPACES BASED ON HARMONIC TROPICALIZATIONS                             47}

\NormalTok{  To prove (2), let ϕ : PHf(x ′ )/Ke ′ → PHf(x)/Ke be the natural morphism. We have seen in 6.13}
\NormalTok{that ϕ −1 (red(X , x)) = red(X ′ , x ′ ). The assertions now follow from Lemma 6.12 and}
\NormalTok{functoriality of residues 6.19. To prove (2′ ), by Corollary 5.9 we may shrink X such that}
\NormalTok{X is compact with formal model X and such that there is L ∈ M (X)Λ with h = hL . Then}
\NormalTok{the claim follows from Lemma 6.8(2) and Lemma 7.4.}
\NormalTok{  The claims (3) and (3′ ) follow from (7.6,1′ ) since X red → X is proper and surjective. }
\NormalTok{ The following proposition allows us to pass to Galois extensions when studying har{-}}
\NormalTok{monic and semipositive functions.}
\NormalTok{Proposition 7.7. Suppose that Λ is divisible. Let L/K be a (potentially infinite) normal}
\NormalTok{extension with automorphism group G := Aut(L/K). Let K ′ be the completion of L and let}
\NormalTok{X ′ := X ⊗}
\NormalTok{         b K K ′ , with structure morphism π: X ′ → X . Then h 7→ h ◦ π induces an isomor{-}}
\NormalTok{phism from the monoid of semipositive Λ{-}PL functions on X (resp. the group of Λ{-}harmonic}
\NormalTok{functions on X ) to the monoid of G{-}invariant semipositive Λ{-}PL functions on X ′ (resp. the}
\NormalTok{group of G{-}invariant Λ{-}harmonic functions on X ′ ).}
\NormalTok{Proof. This is an immediate consequence of Corollary 5.13 and Proposition 7.6(2′ ).                       }
\end{Highlighting}
\end{Shaded}

\subsection{GJR, newline lines
3890--4047}\label{gjr-newline-lines-38904047}

\begin{Shaded}
\begin{Highlighting}[]
\NormalTok{                                 ′ ∗           ∗ ∗             ∗}
\NormalTok{Proposition 8.11 andR that (h ) (α) = π (h (α)) = π (ω), it is clear that both integrals in}
\NormalTok{(1) coincide with h(X ) α. Proposition 11.1 then yields (1) in general.}
\NormalTok{                          d}
\NormalTok{   To prove (2), we choose a finite covering (X i′ )i∈I of supp(ω′ ) by compact strictly ana{-}}
\NormalTok{lytic domains X i′ such that ω′ |X i′ = h∗i (α′i ) for some α′i ∈ Aps}
\NormalTok{                                                                   d,d}
\NormalTok{                                                                       (h(X i′ )) and a piecewise linear}
\NormalTok{tropicalization hi : X i′ → Rni . Then the compact strictly analytic domains σ−1 (X i′ ) cover}
\NormalTok{the support of σ∗ (ω′ ), and the restriction of σ∗ (ω′ ) to σ−1 (X i′ ) is given by (σ ◦ hi )∗ (α′i ).}
\NormalTok{By Proposition 11.1(3), we have}
\NormalTok{                              Z                        Z          Z}
\NormalTok{                                               σ∗ (ω) =                α′i =          ω.}
\NormalTok{                                 σ−1 (X i′ )              hi (X i′ )           X i′}

\NormalTok{The integrals in (2) can be computed by the inclusion{-}exclusion formula (11.1.1) as in}
\NormalTok{the proof of Proposition 11.1, so (2) follows.                                      }
\NormalTok{                             15. HARMONIC FUNCTIONS ON CURVES}
\NormalTok{                                                                                  p}
\NormalTok{   Let K be a non{-}Archimedean field with nontrivial value group Γ ⊂ R, and let Γ :=}
\NormalTok{\{t ∈ R | ∃k ∈ N\textgreater{}0 such that kt ∈ Γ \} be the divisible hull. We will show that our notion}
\NormalTok{              FORMS ON BERKOVICH SPACES BASED ON HARMONIC TROPICALIZATIONS                77}

\NormalTok{of harmonic functions introduced in Section 7 agrees in the case of a rig{-}smooth curve}
\NormalTok{with the harmonic functions in Thuillier’s thesis [Thu05]. This will be important in the}
\NormalTok{companion paper [GJR21].}
\NormalTok{   First, we recall some notions from the theory of graphs and their relationship to the}
\NormalTok{structure of non{-}Archimedean analytic curves. For more details, we refer to [GJR21].}

\NormalTok{15.1 (Metric Graphs with Boundary). We consider a weighted metric graph (Σ, ∂ Σ) with}
\NormalTok{boundary. This is a finite multigraph with no loop edges endowed with positive edge}
\NormalTok{lengths assumed to be in Γ . The boundary ∂ Σ is any subset of the vertices of Σ. Every}
\NormalTok{edge is endowed with the additional datum of a weight w(e) ∈ Z\textgreater{}0 . If we ignore the}
\NormalTok{boundary, then we get a compact piecewise (Z, Γ ){-}linear space of dimension at most 1 as}
\NormalTok{introduced in 2.5. (Note that isolated vertices are allowed.)}
\NormalTok{   A function h: Σ → R is called harmonic if the restriction}
\NormalTok{                                                        P         of h to any edge is affine}
\NormalTok{R{-}linear, and if for any vertex x of Σ \textbackslash{} ∂ Σ, we have e− =x w(e)de (h) = 0, where e ranges}
\NormalTok{over all edges with tail vertex x and de (h) is the outgoing slope of h at x in the direction}
\NormalTok{of e.}
\NormalTok{   For a subgroup Λ of R containing Γ , we let Σ(Λ) denote the subset of Σ consisting of}
\NormalTok{all points x such that f (x) ∈ Λ for all (Z, Γ ){-}linear functionspf on an edge e containing}
\NormalTok{x. Let U ⊂ Σ be a subgraph with vertices contained in Σ( Γ ). A function h: U → R}
\NormalTok{is called piecewise (Z, Γ ){-}linear if there}
\NormalTok{                                        p is a subdivision of the subgraph U such that the}
\NormalTok{vertices of the subdivision are in Σ( Γ ) and such that the restriction of h to any edge e}
\NormalTok{of the subdivision extends to a (Z, Γ ){-}linear function on the edge of Σ containing e.}

\NormalTok{15.2 (Skeletons of Curves). Metric graphs with boundary arise in non{-}Archimedean ge{-}}
\NormalTok{ometry as follows. Let X be a compact rig{-}smooth curve over K, i.e., a compact rig{-}smooth}
\NormalTok{strictly analytic space of pure dimension 1. We assume that X has a strictly semistable}
\NormalTok{model X over the valuation ring K ◦ . This is a formal K ◦ {-}model of X (see 2.4) whose spe{-}}
\NormalTok{cial fiber Xs is geometrically reduced with smooth irreducible components and at most}
\NormalTok{double points as singularities. This is equivalent to [Thu05, Definition 2.2.8], where}
\NormalTok{such models are called simplement semi{-}stable. To such a model is associated a weighted}
\NormalTok{metric graph with boundary Σ := S(X), called the skeleton of X. There is a canonical}
\NormalTok{inclusion Σ ,→ X . The vertices of Σ correspond to the irreducible components of Xs and}
\NormalTok{the edges correspond to the singular points of Xs , so that Σ is the incidence graph of Xs .}
\NormalTok{The boundary ∂ Σ consists of the vertices of Σ corresponding to non{-}proper irreducible}
\NormalTok{components of Xs . If x ∈ Xs is a singular point with corresponding edge e, then the}
\NormalTok{weight of e is defined to be w(e) = [Ke (x) : K}
\NormalTok{                                              e ]. There is an étale neighbourhood of x of}
\NormalTok{the form Spf(K\{X , Y \}/(X Y − ̟)) for some ̟ ∈ K ◦◦ \textbackslash{} \{0\}; the length of e is then defined}
\NormalTok{to be v(̟). There is a canonical deformation retraction map τ: X → S(X).}

\NormalTok{15.3 (Subdivisions of the Skeleton). Consider}
\NormalTok{                                      p         a subdivision Σ′ of the skeleton Σ = S(X)}
\NormalTok{such that all vertices of Σ′ are in Σ( Γ ). Vilsmeier [Vil21, Construction 2.6] has de{-}}
\NormalTok{fined a canonical formal scheme X′ with generic fiber X and reduced special fiber, and a}
\NormalTok{canonical morphism ι : X′ → X extending IdX , such that for any edge e′ of Σ′ , the analytic}
\NormalTok{subdomain (τ)−1(e′ ) of X is the generic fiber of an open subset of X′ . Note that X′ is not}
\NormalTok{necessarily an admissible formal scheme, i.e. not necessarily of topologically finite type.}
\NormalTok{78                             W. GUBLER, P. JELL, AND J. RABINOFF}

\NormalTok{By [Gub13b, Propositions 6.7 and 6.9], the formal scheme X′ is admissible if the value}
\NormalTok{group Γ is discrete or if the vertices of Σ′ are contained in Σ(Γ ).}
\NormalTok{15.4 (Piecewise Affine Functions on the Skeleton).}
\NormalTok{                                             p         Let Σ = S(X) as in 15.2, let U be a}
\NormalTok{subgraph of Σ with vertices contained in Σ( Γ ), and let h: U → R be a piecewise (Z, Γ ){-}}
\NormalTok{linear function in the sense of 15.1. Passing to a subdivision of U, we may assume that}
\NormalTok{h is (Z, Γ ){-}linear on each edge of U. The vertices of U induce a subdivision Σ′ of Σ; we}
\NormalTok{denote the associated formal K ◦ {-}model of X by X′ as in 15.3, with canonical morphism}
\NormalTok{ι : X′ → X. When Γ dense in R, we assume that all vertices of U are in Σ(Γ ); then}
\NormalTok{X′ is an admissible formal scheme. The subset Y := τ−1 (U) is the generic fiber of a}
\NormalTok{unique open formal subscheme Y′ in X′ . By the arguments used in the proof of [Vil21,}
\NormalTok{Proposition 2.11], there is a canonical vertical Cartier divisor D on Y′ with associated}
\NormalTok{formal metric on OY determined by − log k1kO (D) = h ◦ τ. See 5.5.}
\NormalTok{Proposition 15.5. Let X be a compact rig{-}smooth curve over K and let h: X → R.}
\NormalTok{     (1) If X is a strictly semistable model of X over K ◦ , then h = hL for a line bundle}
\NormalTok{         L ∈ M (X) if and only if there is a function f : S(X) → R which is (Z, Γ ){-}linear on}
\NormalTok{         each edge such that h = f ◦ τ.}
\NormalTok{     (2) If x ∈ ∂ X , then h is harmonic at x if and only if h is R{-}PL at x.}
\NormalTok{Proof. The “if” part in (1) follows from 15.4. Conversely, assume that h = hL for L ∈}
\NormalTok{M (X). Using the identification L|X = OX , the meromorphic section 1 of L induces a}
\NormalTok{vertical Cartier divisor D on X. There is a covering of X by open subsets U such that S(U)}
\NormalTok{is either an edge or a vertex of S(X) and such that D|U = div(γ) for some γ ∈ O (U)×,}
\NormalTok{where U = Uη . It follows from [Ber99, Theorem 5.3 and Step 13 in its proof] that}
\NormalTok{h|U = − log |γ| factorizes through a (Z, Γ ){-}linear function on S(U). This proves (1).}
\NormalTok{   Since x ∈ ∂ X , the curve ∂ X is affine and hence the nef condition in Proposition 7.1 is}
\NormalTok{vacuous. This proves (2).                                                                 }
\NormalTok{15.6 (Thuillier’s Harmonic Functions). We consider now an arbitrary rig{-}smooth curve}
\NormalTok{X over K. We call a function h: X → R harmonic in the sense of Thuillier if every x ∈ X}
\NormalTok{has a strictly affinoid neighbourhood V , a finite Galois extension K ′ /K, and a strictly}
\NormalTok{semistable model V′ of V ′ := V ⊗K K ′ , such that the pullback h′ of h|V factorizes through}
\NormalTok{the skeleton S(V′ ) and such that h′ |S(V′ ) is a harmonic function on the skeleton S(V′ )}
\NormalTok{considered as a weighted metric graph with boundary as in 15.2.}
\NormalTok{   For x ∈ ∂ X , Proposition 15.5 and Corollary 5.12 show that the above harmonicity}
\NormalTok{condition at x is equivalent to piecewise R{-}linearity at x, using that the boundary of}
\NormalTok{S(X) is ∂ X . Note that Thuillier [Thu05, Section 2.3] used this definition only for strictly}
\NormalTok{affinoid or boundaryfree X , so the above definition is slightly more general than his.}
\NormalTok{Proposition 15.7. Let X be a rig{-}smooth curve. Then a function h: X → R is harmonic in}
\NormalTok{the sense of Definition 7.2 if and only if h is harmonic in the sense of Thuillier.}
\NormalTok{Proof. The claim is local on X , so we may assume X strictly affinoid. We first show}
\NormalTok{that the harmonicity notions agree for a piecewise linear function h ∈ PL(X ). By 5.6,}
\NormalTok{we have that h = − log k1kL for a model L of OX on a formal K ◦ {-}model X of X . Using}
\NormalTok{base change, Proposition 7.6 and [Thu05, Remarque 2.3.17], we are free to pass to a}
\NormalTok{              FORMS ON BERKOVICH SPACES BASED ON HARMONIC TROPICALIZATIONS                79}

\NormalTok{finite Galois extension of K. Arguing locally, the semistable reduction theorem ([Thu05,}
\NormalTok{Théorème 2.3.8] under these hypotheses; see also [BL85, Section 7] and [Duc24, Theo{-}}
\NormalTok{rem 5.1.14(iv), 6.1.3]) shows that we may assume that X is a strictly semistable model.}
\NormalTok{By Proposition 15.5, there is a function f on the skeleton S(X) which is (Z, Γ ){-}linear on}
\NormalTok{each edge such that h = f ◦ τ. For x ∈ S(X) \textbackslash{} ∂ X with corresponding proper irreducible}
\NormalTok{component C := redX (x) of Xs , it follows as in [KRZB16, Theorem 2.6] that the sum of}
\NormalTok{outgoing slopes of f at x is equal to deg(L|C ). We conclude that h is harmonic in the}
\NormalTok{sense of Thuillier if and only if L is numerically equivalent to 0. The latter is equivalent}
\NormalTok{to harmonicity of h in the sense of Definition 7.2, as we have shown in Lemma 7.4.}
\NormalTok{   Now we prove the proposition for any function h: X → R. Assume that h is a harmonic}
\NormalTok{function on X as in Definition 7.2. By Proposition 7.5, the function h is locally an R{-}linear}
\NormalTok{combination of Z{-}harmonic functions, so we are reduced to the case handled above since}
\NormalTok{the set of Thuillier{-}harmonic functions forms a real vector space.}
\NormalTok{   Conversely, let h be a harmonic function in the sense of Thuillier. We have to check}
\NormalTok{for any x ∈ X that h is harmonic at x. As above, we are free to pass to a finite Galois}
\NormalTok{extension of X , so we may extend scalars and shrink X to assume that X has a strictly}
\NormalTok{semistable model X and that h factorizes through a harmonic function f : S(X) → R. It}
\NormalTok{is enough to show that f is an R{-}linear combination of piecewise (Z, Γ ){-}linear harmonic}
\NormalTok{functions f i on S(X). Indeed, by Proposition 15.5, we have that f i ◦ τ is piecewise linear,}
\NormalTok{so f i ◦ τ is harmonic by the the first paragraph of this proof, and hence h is harmonic.}
\NormalTok{   Suppose first that x ∈ ∂ X = ∂ S(X). After shrinking X, we may assume that S(X)}
\NormalTok{is a star around the vertex x, so that f is affine R{-}linear on each edge. Translating by}
\NormalTok{a constant function, we may assume f (x) = 0. Then f is an R{-}linear combination of}
\NormalTok{piecewise Z{-}linear functions f e on S(X): namely, the functions f e that have slope 1 on a}
\NormalTok{given edge e and slope 0 on the others. Since x ∈ ∂ S(X), Proposition 15.5 yields that}
\NormalTok{each f e is harmonic and hence h is harmonic.}
\NormalTok{   If x ∈ X \textbackslash{} ∂ X , then as above, we may assume that S(X) is either an edge containing}
\NormalTok{τ(x) in its interior, or that S(X) is a star around the vertex τ(x). In the first case, f is}
\NormalTok{an R{-}multiple of a (Z, Γ ){-}linear function on the edge, so h is harmonic. In the second}
\NormalTok{case, we may assume that f (τ(x)) = 0. Then f is harmonic if and only if f is linear on}
\NormalTok{each edge and the weighted sum of the outgoing slopes at τ(x) is zero. This is a linear}
\NormalTok{equation for the slopes which is defined over Q, so there is a basis of solutions defined}
\NormalTok{over Q and even over Z. The elements f i of such a basis induce piecewise (Z, Γ ){-}linear}
\NormalTok{functions f i which are harmonic on S(X). Since f is an R{-}linear combination of the f i ,}
\NormalTok{we have that h is harmonic.                                                                }

\NormalTok{   Note that Γ is either dense inpR or is a discrete subgroup. In the following, we set}
\NormalTok{Λ := Γ if Γ is dense and Λ := Γ (the divisible hull in R) if Γ is discrete. When Γ is}
\NormalTok{discrete, we may rescale the valuation to assume Γ = Z.}
\end{Highlighting}
\end{Shaded}

\subsection{Berkovich, newline lines
1350--1444}\label{berkovich-newline-lines-13501444}

\begin{Shaded}
\begin{Highlighting}[]
\NormalTok{through the strictly affinoid domains in X, are compatible on intersections, and to they}
\NormalTok{glue together to form a sheaf of rings (9\^{} on the G{-}topological space Xo. The locally}
\NormalTok{G{-}ringed space (Xo, (P\^{}) satisfies the conditions of Definition 9.3.1/4 from [BGR],}
\NormalTok{and so we get a rigid \^{}{-}analytic space. We remark that the rigid yfe{-}analytic space cons{-}}
\NormalTok{tructed is quasiseparated. (A rigid A{-}analytic space is called quasiseparated if the inter{-}}
\NormalTok{section of two open affinoid domains is a finite union of open affinoid domains.)}

\NormalTok{       1.6.1. Theorem. —The correspondence X i{-}\textgreater{} XQ is a fully faithful junctor from the}
\NormalTok{category of Hausdorff strictly k{-}analytic spaces to the category of quasiseparated rigid k{-}andlytic}
\NormalTok{36}
\NormalTok{                                    VLADIMIR G. BERKOVICH}

\NormalTok{spaces. Furthermore, this functor induces an equivalence between, the category of paracompact strictly}
\NormalTok{k\textquotesingle{}analytic spaces and the category of quasiseparated rigid k{-}analytic spaces that have an admissible}
\NormalTok{affinoid covering of finite type.}

\NormalTok{      A collection of subsets of a set is said to be of finite type if each subset of the col{-}}
\NormalTok{lection meets only a finite number of other subsets of the collection.}

\NormalTok{      Proof. — Let X be a Hausdorff strictly A{-}analytic space. First of all we establish}
\NormalTok{the following fact.}

\NormalTok{      1.6.2. Lemma. — (i) Any open affinoid domain in the rigid k{-}analytic space XQ is of the}
\NormalTok{formVo, where V is a strictly affinoid domain in X.}
\NormalTok{      (ii) Let \{VJ\^{} be a system of strictly affinoid domains in X. Then \{ V , o \} \^{} i is}
\NormalTok{ an admissible covering of XQ if and only if each point ofX has a neighborhood of the form}
\NormalTok{ V\^{} U ... uV\^{} (i.e., \{ V J \^{} i is a quasinet on X).}

\NormalTok{       Proof. — (i) An open affinoid domain in XQ is an open immersion of rigid \^{}{-}analytic}
\NormalTok{spaces/: Uo {-}\textgreater{} XQ, where U is a strictly \^{}{-}affinoid space. In particular, for any strictly}
\NormalTok{affinoid domain V C X , /{-}\^{}(Vo) is a finite union of affinoid domains in UQ, and}
\NormalTok{\{/"\^{}(Vo)\}, where V runs through strictly affinoid domains in X, is an affinoid covering.}
\NormalTok{It follows that we can find strictly affinoid domains Ui, ..., U\^{} C U and V\textbackslash{}, ..., V C X}
\NormalTok{such that U = Ui U ... u U\^{} and f\textbackslash{}\^{} \^{} comes from a morphism of strictly affinoid}
\NormalTok{spaces 9,: U, {-}\textgreater{}V, that identifies U, with an affinoid domain in V,. Moreover, all 9,}
\NormalTok{are compatible on intersections. Therefore, we get a morphism of strictly \^{}{-}analytic}
\NormalTok{spaces \textless{}p : U {-}\textgreater{} X. Since 9, as a map of topological spaces, is compact and induces an}
\NormalTok{injection on the everywhere dense subset Uo C U, it follows that 9 induces a homeo{-}}
\NormalTok{morphism of U with its image in X. Finally, 9 identifies U, with a strictly affinoid}
\NormalTok{domain in V,, and therefore 9 identifies U with a strictly analytic domain in X. It is}
\NormalTok{clear that this is a strictly affinoid domain.}
\NormalTok{       (ii) Suppose first that \{V, o h e i ls an admissible covering of XQ. This means}
\NormalTok{that, for any strictly affinoid domain V C X, \{ V, o n V\^{} \}, \^{}; is an admissible covering}
\NormalTok{of Vo. It follows that V is contained in a finite union V\^{} U .. . U V\^{}, and therefore}
\NormalTok{each point of X has a neighborhood of the required form. Conversely, assume that}
\NormalTok{the latter property is true. Then any strictly affinoid domain is contained in a finite}
\NormalTok{union V\^{} U . .. u V\^{}, and therefore \{ V, \^{} n VQ \}\^{} j is an admissible covering of Vo. •}
\NormalTok{       Let 9 : Y {-}\textgreater{} X be a morphism of strictly A{-}analytic spaces. First of all we claim}
\NormalTok{that the induced map 90 : Yg {-}\textgreater{} XQ is continuous with respect to the G{-}topologies on XQ}
\NormalTok{and Yo. Let \^{}C XQ be an admissible open subset, and let VC Y be a strictly affinoid}
\NormalTok{domain. By [BGR], 9.1.4/2, and Lemma 1.6.2 (i), the set \^{} has an admissible cove{-}}
\NormalTok{ring \{ U \^{} p J \^{} j , where U\^{} are strictly affinoid domains in X. By Corollary 1.2.14,}
\NormalTok{for each i e I one has 9{-}1(U,) n V = U,gj. V\^{}., where V,, are strictly affinoid domains}
\NormalTok{in V and J, is finite. We get a covering \{ V,, o \} of the set 9o\textquotesingle{}lW \^{} VQ. To verify that}
\NormalTok{            fiTALE COHOMOLOGY FOR NON{-}ARCHIMEDEAN ANALYTIC SPACES                             37}

\NormalTok{the latter set is admissible open in VQ, it suffices to show that for any morphism of}
\NormalTok{strictly \^{}{-}affinoid spaces + : W —Vwith\^{}(Wo) C 9o{-}l(\^{}) n Vg, the covering \{ \^{} \textquotesingle{}(V,, o)\}}
\NormalTok{has a finite affinoid covering that refines it. But this follows from the fact that the latter}
\NormalTok{condition is satisfied by the covering \{(9\^{})0" \^{}U\^{}o h e r Thus, the set 9o\textasciitilde{}l(\textless{}\%\^{}) is admis{-}}
\NormalTok{sible open in Yo. In the same way one shows that the preimage of an admissible covering}
\NormalTok{of an admissible open set is an admissible open covering. Hence the map of G{-}topological}
\NormalTok{spaces 90 : Y\^{} {-}\textgreater{} X\^{} is continuous. That 9 induces a morphism of locally G{-}rmged}
\NormalTok{spaces easily follows from this.}
\NormalTok{       Let now/: Y() {-}\textgreater{} X\^{} be a morphism between the above rigid \^{}{-}analytic spaces.}
\NormalTok{We have to show that it comes from a unique morphism of strictly \^{}{-}analytic spaces}
\NormalTok{9 : Y {-}\textgreater{} X. First of all, from Proposition 1.3.2 it follows that it suffices to verify the}
\NormalTok{required fact only in the case when Y is strictly \^{}{-}affinoid. For this we remark that the}
\NormalTok{system \{ Uo\}, where U runs through strictly affinoid domains in X, is an admissible}
\NormalTok{covering of X\^{}. Therefore \{/"\^{}(Uo)\} is an admissible covering of Yg by admissible open}
\NormalTok{subsets. By [BGR], 9.1.4/2, the latter covering has a finite affinoid covering that refines}
\NormalTok{it. In this way we get strictly affinoid domains V,, . . . , V \^{} C Y and Ui, ..., U\^{}C X}
\NormalTok{such that YQ = V\^{}o U . . . u V\^{} (and therefore Y = V\^{} u . . . u VJ and/(V\^{}) C U,\^{}.}
\NormalTok{The induced morphisms of strictly affinoid spaces V, {-}\textgreater{} U, are obviously compatible}
\NormalTok{on intersections, and therefore we get a morphism of strictly yfe{-}analytic spaces 9 : Y {-}\textgreater{} X.}
\NormalTok{It is easy to see that 90 ==/and that 9 is a unique morphism satisfying this property.}
\NormalTok{       If X is a paracompact strictly A{-}analytic space, then it has a strictly A{-}affinoid}
\NormalTok{atlas with a locally finite net, and therefore the rigid \^{}{-}analytic space X\^{} has an admis{-}}
\NormalTok{sible affinoid covering of finite type. It is also evident that it is quasiseparated. Conversely,}
\NormalTok{let °K be a quasiseparated rigid \^{}{-}analytic space that has an admissible affinoid cove{-}}
\NormalTok{\^{}g \{ \^{} \} i e i of finite tyP0{-} First of \^{} let \^{} = L\^{}o, where the U, are strictly}
\NormalTok{A{-}affinoid spaces. Since 3E is quasiseparated, for any pair i,j e I the intersection \^{} n \%.}
\NormalTok{is a finite union of open affinoid domains in SK\textbackslash{} Thus, there are strictly special domains}
\NormalTok{U,, C U, and U\^{} C U, that correspond to \^{} n \^{}, under the identifications \^{} == U, p}
\NormalTok{and \^{}. = U, o. Let \^{} denote the induced isomorphism U \^{} \^{} U \^{} . It is clear that}
\NormalTok{v}
\NormalTok{  ii = U? \^{}(U, n UJ = U,, n U,\^{} and \^{} == v\^{} o v,, on U\^{}. n U\^{}. By Proposi{-}}
\NormalTok{tion 1.3.3, we can glue all U, along U,, and get a paracompact strictly \^{}{-}analytic}
\NormalTok{space X. It is easy to see that XQ is isomorphic to S\textbackslash{} •}
\NormalTok{       Let X be a Hausdorff strictly A{-}analytic space. From Lemma 1.6.2 it follows}
\NormalTok{that there is an isomorphism of topoi X\^{} \^{} X\^{}. In particular, there is a morphism of}
\NormalTok{topoi (7\^{},71;*) : X\^{} {-}\^{}X\^{} such that the functor TC* is fully faithful (and TC, is not).}
\NormalTok{Furthermore, from Proposition 1.3.6 it follows that if F is an abelian sheaf on X, then}
\NormalTok{H\^{}X, F) \^{}\textgreater{} H\^{}Xo, TT* F), q \^{} 0, and if X is good and F is a coherent 0\^{} module, then}
\NormalTok{H\^{}(X,F)\^{}IP(Xo,Fo), q \^{} O , where F o = = \^{} F ® \^{} \^{} \^{} o {-} Finally, there are equi{-}}
\end{Highlighting}
\end{Shaded}

\subsection{Berkovich, newline lines
2200--2347}\label{berkovich-newline-lines-22002347}

\begin{Shaded}
\begin{Highlighting}[]
\NormalTok{3.1. Quasifinite morphisms}

\NormalTok{       3.1.1. Definition. — A morphism of\^{}{-}analytic spaces 9 : Y {-}\textgreater{} X is said to be finite}
\NormalTok{at a pointy e Y if there exist open neighborhoods Y\^{} ofy and \^{}ofy(j/) such that 9 induces}
\NormalTok{a finite morphism i\^{} {-}\textgreater{} W; 9 is said to be quasifinite if it is finite at any point y e Y.}
\NormalTok{       It follows from the definition that quasifinite morphisms are locally separated}
\NormalTok{and closed.}

\NormalTok{          3.1.2. Lemma. — If a morphism 9 : Y {-}\textgreater{} X is finite at a point y e Y, then the neigh{-}}
\NormalTok{borhoods V and W from the Definition 3 . 1 . 1 can he found arbitrary small.}

\NormalTok{     Proof. — We may assume that the morphism y is finite. Let x == y(jy) and let}
\NormalTok{(p{-}\^{}) == \{j/i ===J\^{}25 • • {-}\^{}A \}• Since the map [ Y [ {-}\^{} [ X [ is compact, we can find a}
\NormalTok{            fiTALE COHOMOLOGY FOR NON{-}ARCHIMEDEAN ANALYTIC SPACES                                55}

\NormalTok{sufficiently small open neighborhood \^{} of x such that y"1\^{}) =11\^{}1\^{}, where}
\NormalTok{y, e \^{} and \^{} n \^{}. = 0 for i \^{}j. It follows that all of the morphisms \^{} {-}\textgreater{} \%}
\NormalTok{are finite, and the neighborhood V\^{} ofj\^{} is sufficiently small. •}

\NormalTok{      3.1.3. Corollary. — Quasifinite morphisms are preserved under compositions, under any}
\NormalTok{base change functor, and under any ground field extension functor. •}

\NormalTok{     We remark that if 9 : Y {-}\textgreater{} X is a quasifinite morphism and the space X is good,}
\NormalTok{then Y is also good. Quasifinite morphisms between good \^{}{-}analytic spaces can be}
\NormalTok{characterized as follows.}

\NormalTok{      3.1.4. Proposition. — Let 9 : Y {-}\textgreater{}X be a morphism of good k{-}analytic spaces, and let}
\NormalTok{y e Y and x == \textless{}p(j/). The following are equivalent:}
\NormalTok{      a) 9 is finite at y\textbackslash{}}
\NormalTok{       b) there exist sufficiently small affinoid neighborhoods V ofy and U ofx such that 9 induces}
\NormalTok{a finite morphism V {-}\textgreater{} U;}
\NormalTok{      c) the point y is isolated in the fibre \^{}\textquotesingle{}\textasciitilde{}l\{x) and y eInt(Y/X).}

\NormalTok{      Proof. — The implications a) =\textgreater{} c) and b) =\textgreater{} c) are trivial. The implication}
\NormalTok{b) =\textgreater{} a) follows from the following Lemma.}

\NormalTok{      3.1.5. Lemma. — Let 9 : Y {-}\textgreater{} X be a morphism, and let V C Y and U C X be affinoid}
\NormalTok{subdomains such that 9 induces a finite morphism \^{} : V {-}\textgreater{} U. Then Int(V/Y) = \^{}(In\^{}U/X)).}
\NormalTok{In particular, y induces a finite morphism Int(V/Y) {-}\textgreater{}Int(U/X).}
\NormalTok{      Proof. — From [Ber], 3.1.3, it follows that}
\NormalTok{                  Int(V/X) = Int(V/Y) n Int(Y/X) = Int(V/Y).}

\NormalTok{On the other hand, Int(V/X) = Int(V/U) n \^{}(In\^{}U/X)) = \^{}(In\^{}U/X)). There{-}}
\NormalTok{fore Int(V/Y) = \^{}(In\^{}U/X)). •}
\NormalTok{     To verify the implication c ) =\textgreater{} b), it suffices to assume that X==\^{}(j\^{}) and}
\NormalTok{Y = Jl\{SS\} are \^{}{-}affinoid. Furthermore, since X and Y are compact, we can decrease}
\NormalTok{them and assume that \^{}\textasciitilde{}l(x) ===\{j»\}. Since y e!nt(Y/X), there exists an admissible}
\NormalTok{epimorphism}
\NormalTok{                  Tr:\^{}\^{}{-}{-} 1 \^{},...\^{} 1 \^{}\}{-}{-}\^{}}

\NormalTok{such that | 7r(T,) \{y) \textbackslash{} \textless{} r,, 1 \^{} i \^{} n. For any affinoid neighborhood U of A?, n induces}
\NormalTok{an admissible epimorphism}
\NormalTok{                  \^{};\^{}\^{}\{\^{}lTl,...,r„{-}lT„\}\^{}\^{}\^{}.}

\NormalTok{      If U is sufficiently small, then y\textquotesingle{}\textquotesingle{}l(U) is a sufficiently small affinoid neighborhood}
\NormalTok{of the point j\textbackslash{} Therefore we can find sufficiently small U such that [ 7r(TJ (y)| \textless{} r}
\NormalTok{56                                VLADIMIR G. BERKOVICH}

\NormalTok{for all points y \textquotesingle{} ey\^{}U). This means that the induced morphism of\^{}{-}affinoid spaces}
\NormalTok{\^{}(U) {-}\^{}U is closed. By [Ber], 2.5.13, the latter morphism is finite. •}

\NormalTok{      3.1.6. Corollary. — If a morphism of good k{-}analytic spaces 9 : Y {-}\textgreater{} X is finite at a point}
\NormalTok{y e Y, then tf\^{} y is a finite 0\^{} \^{}{-}algebra. •}

\NormalTok{    Recall that a morphism of schemes 9 : W {-}\textgreater{} S£ is called quasifinite if any point}
\NormalTok{y eW is isolated in the fibre y\textquotesingle{}\^{}x) = \^{}, where x = (p(y) (see [SGA1], 1.2).}

\NormalTok{      3.1.7. Corollary. — A morphism 9 : 3£ —\^{}W between schemes of locally finite type over}
\NormalTok{3E == Spec(J\^{}), where \^{} is a k{-}affinoid algebra\^{} is quasifinite if and only if the corresponding}
\NormalTok{morphism \^{}n: 2£\^{} {-}\textgreater{} \^{}an is quasifinite.}

\NormalTok{      Proof. — For any pointy e \^{}an, there is an isomorphism of \^{}\{y) {-}analydc spaces}

\NormalTok{                  (\^{}\^{}\^{}(J))\^{}\^{}.}

\NormalTok{Therefore our assertion follows from the fact that the morphism 9\^{} is closed (Pro{-}}
\NormalTok{position 2.6.1). •}

\NormalTok{     3.1.8. Proposition. — Let 9 : Y {-}\textgreater{}• X be a morphism of k{-}analytic spaces\^{} and let y e Y}
\NormalTok{and x = \textless{}p(j0. Then the following are equivalent:}
\NormalTok{       a) 9 is finite at y\textbackslash{}}
\NormalTok{       b) there exist analytic domains X\^{}, ..., X\^{}C X such that x e X\^{} n ... n X\^{},}
\NormalTok{X\^{} u . . . u X\^{} is a neighborhood of x and the morphisms (p\textquotesingle{}\^{}X\^{}) {-}\textgreater{} X\^{} are finite at y\textbackslash{}}
\NormalTok{       c) there exist affinoid domains Vi, . . . , V\^{} C Y and U\^{}, ..., U\^{} C X such that}
\NormalTok{y e Vi n . . . n V\^{}, V\^{} u ... u V\^{} fl\^{} Ui U ... u U\^{} flr\^{} neighborhoods of y and x,}
\NormalTok{respectively\^{} \textless{}p(VJ C U,, and the induced morphisms 9,: V\^{} —\^{} U, flwrf \textless{}p\^{} : V, n V, {-}\textgreater{} U\^{} n U,}
\NormalTok{\textless{}?r\^{} \^{}TZZ\^{} \^{} j.}

\NormalTok{      We remark that if the spaces X and Y are separated, then the finiteness atj\^{} of}
\NormalTok{all the morphisms 9, from c ) implies the same property for the morphisms \textless{}p\^{}. Indeed,}
\NormalTok{in this case U, n U, and V, n V, are affinoid domains and, by Proposition 3.1.4,}
\NormalTok{it suffices to verify that y e Int(V, n V,/U, n U,). If V,, := y.{-}\^{}U, n U,), then}
\NormalTok{y e!nt(V,,/U, n U,), and therefore y e Int(V,, n V,JU, n U,). It remains to note}
\NormalTok{that V , , n V , , = V , n V , .}
\NormalTok{       proof. — The implication a) =\textgreater{} b) is trivial. Furthermore, we remark that all}
\NormalTok{ three properties remain true if we replace X and Y by sufficiently small open neigh{-}}
\NormalTok{ borhoods of the point x andjy, respectively. In particular, we may assume that X and Y}
\NormalTok{ are Hausdorff.}
\NormalTok{       b) =\textgreater{} c ) . We may assume that X, == U,\textquotesingle{} are affinoid domains. Then we can}
\NormalTok{ find affinoid neighborhoods V,\textquotesingle{} of y in \^{}\textasciitilde{}\textasciitilde{}l(Vf,) such that the induced morphisms}
\NormalTok{            fiTALE COHOMOLOGY FOR NON{-}ARCHIMEDEAN ANALYTIC SPACES                             57}

\NormalTok{9,\textquotesingle{} : V,\textquotesingle{} {-}\textgreater{} U,\textquotesingle{} are finite at j. Furthermore, shrinking X and Y, we may assume that}
\NormalTok{the morphisms 9,\textquotesingle{} are finite and \^{}\textasciitilde{}\textbackslash{}x) =\{\^{}\}. Thus we get a morphism \^{} : V {-}\textgreater{} U,}
\NormalTok{where V = V; u . . . u V, and U = U, u . . . u U, are compact analytic neigh\^{}}
\NormalTok{borhoods ofy and x, respectively. It follows that we can find on open neighborhood \%}
\NormalTok{of x in X with \^{} C U such that \^{} := \^{}{-}1\^{}) is an open neighborhood ofy in Y and,}
\NormalTok{for every 1 \^{} \^{} n, \^{}(U; n \^{}) C V;. If now U, is an affinoid neighborhood of x in}
\NormalTok{U; n\^{}, then V,:== \^{}(U,) is an affinoid neighborhood of y in V; n \^{}, and the}
\NormalTok{induced morphism 9,: V, {-}\^{} U, is finite. Since V, n V, = (p\^{}W \^{} U,), the induced}
\NormalTok{morphisms 9,, : V, n V, {-}\textgreater{} U, n U, are finite. Hence, 9 satisfies c).}

\NormalTok{      3.1.9. Lemma. — If Y is an analytic domain in X and the canonical morphism Y {-}\textgreater{} X}
\NormalTok{is quasifinite, then Y is open in X.}

\NormalTok{        Proof. — Since quasifinite morphisms are closed, Int(Y/X) = Y. By Proposi{-}}
\NormalTok{tion 1.5.5 (ii), Int(Y/X) coincides with the topological interior of Y in X. It follows}
\NormalTok{that Y is open in X. •}
\NormalTok{        c ) =\textgreater{} a). We can shrink all the affinoid domains V, and assume that 9,{-}1(A:) == \{y \}.}
\NormalTok{Then for sufficiently small affinoid neighborhoods U,\textquotesingle{} of A: in U,, 9\^{}1(\^{}/) and}
\NormalTok{Pi\^{}W n u \^{} ) = Pz\textquotesingle{}\^{}UQ \^{}r1\^{}.) are sufficiently small neighborhoods of y in V}
\NormalTok{and V, n V,, respectively. Thus, we can shrink all the affinoid domains U, and assume}
\NormalTok{that all of the morphisms 9\^{} and 9\^{}. are finite.}
\NormalTok{        Furthermore, the morphism 9, induces finite morphisms 9\^{}.: V, n V {-}\textgreater{} U, n U}
\NormalTok{and V , , : = 9 , {-} l ( \^{} n U , . ) {-} \textgreater{} \^{} n U , , and therefore the canonical embedding of}
\NormalTok{special domains V, n V, {-}\textgreater{} V,, is finite. From Lemma 3.1.9 it follows that V, n V,}
\NormalTok{is open in V,,, and therefore V,, == (V, n V,) U W,, for some special domain W\^{}CV,!}
\NormalTok{We get 9{-} l (\^{}) = = V , U W , , where W , = U , \^{} W , , . Since y e V, n . . . n V,, we}
\NormalTok{have x \^{} 9(W,), and therefore we can find, for each 1 \^{} i \^{} n, an affinoid neighborhood U,\textquotesingle{}}
\NormalTok{of A; in U, such that U; n 9(W,) = 0. The latter implies that 9\textasciitilde{}1(U;) == V; := 9\^{}1(U;)t.}
\NormalTok{Hence, the morphism V\^{} U . .. u V\^{} {-}\textgreater{} U[ u . . . u U\^{} is finite, and the required}
\NormalTok{statement follows. •}

\NormalTok{      3 . 1 . 1 0 . Corollary. — A morphism of k{-}analytic spaces 9 : Y {-}\textgreater{} X is finite at a point}
\NormalTok{y e Y if and only if the point y is isolated in the fibre 9{-}\textquotesingle{}l(9(\^{})) and y e!nt(Y/X).}

\NormalTok{       Proof. — The direct implication is clear. Suppose thatj\^{} is isolated in 9\textasciitilde{} l (9(\^{}))}
\NormalTok{andjy eInt(Y/X). We can shrink Y and assume that 9 is closed. Let U\^{}, . .., U\^{} be}
\NormalTok{affinoid domains in X such that \^{}\{y\} e Ui n . . . n U\^{} and Ui u ... u U\^{} is a neigh{-}}
\NormalTok{borhood of 9(j/). From Proposition 3.1.4 is follows that 9 {-}l (\^{}) {-}\textgreater{} U, are closed}
\NormalTok{morphisms of good \^{}{-}analytic spaces. In particular, the property b ) of Proposition 3.1.8}
\NormalTok{holds. It follows that 9 is finite at y. •}
\NormalTok{58                                VLADIMIR G. BERKOVICH}


\NormalTok{3.2. Flat quasifinite morphisms}

\NormalTok{      A morphism of good \^{}{-}analytic spaces 9 : Y —»{-} X is said to be flat at a pointy e Y}
\end{Highlighting}
\end{Shaded}

\subsection{Thuillier, newline lines
3070--3455}\label{thuillier-newline-lines-30703455}

\begin{Shaded}
\begin{Highlighting}[]
\NormalTok{                             0         si c est ontenue dans S(Y) − S(X );}
\NormalTok{                λx,t (Φ) =}
\NormalTok{                             λx,t (Φ) si c est ontenue dans S(X );}

\NormalTok{l\textquotesingle{}égalité}
\NormalTok{                                     ddc Φ = S(q)∗ (ddc S(q)∗ Φ)}
\NormalTok{en dé oule immédiatement.                                                                            2}
\NormalTok{Proposition 2.2.27.  Soient X , Y deux S {-} ourbes simplement semi{-}stables et f : Xη →}
\NormalTok{Yη un morphisme quasi{-}ni. Il existe une extension nie séparable k′ /k et un é latement}
\NormalTok{admissible q : X ′ → X ×S S ′ , S ′ = Spf(k′ ◦ ), tels que}
\NormalTok{  (i ) X ′ soit une S {-} ourbe simplement semi{-}stable et le morphisme f ′ = f ⊗ 1 : Xη ⊗k k′ →}
\NormalTok{Yη ⊗k k′ se prolonge en un Spf(S ′ ){-}morphisme f ′ : X ′ → Y ′ , où Y ′ = Y ×S Spf(S ′ )}
\NormalTok{  (ii ) l\textquotesingle{}appli ation S(f ′ ) = τY ′ ◦ f|S(X ′ ) de S(X ′ ) dans S(Y ′ ) soit un morphisme de polyèdres}
\NormalTok{entiers,}
\NormalTok{  (iii ) le diagramme}
\NormalTok{                                                          ′}
\NormalTok{                                                         τX}
\NormalTok{                                            Xη′                     / S(X ′ )}

\NormalTok{                                            f                               S(f ′ )}
                                                                       
\NormalTok{                                            Yη′          τY ′}
\NormalTok{                                                                    / S(Y ′ )}

\NormalTok{soit ommutatif.}
\NormalTok{                                                                                                  49}



\NormalTok{  Dans es onditions :}
\NormalTok{  (iv ) pour tout point x ∈ S(X ′ ) − ∂Xη′ et toute fon tion Φ ∈ A0 (S(Y ′ )),}
\NormalTok{                          Z                              Z}
\NormalTok{                                   c   ′ ∗             ′}
\NormalTok{                              dd S(f ) Φ = (degx f )            ddc Φ.}
\NormalTok{                             \{x\}                                      \{f ′ (x)\}}

\NormalTok{Démonstration.}
\NormalTok{(i) Tous les points du sous{-}ensemble Σ = f −1 (S0 (Y)) de Xη sont de type (2) ( ar f est quasi{-}}
\NormalTok{ni) et, par appli ation du lemme 2.2.22, il existe une extension nie k′ /k et un é latement}
\NormalTok{admissible q : X ′ → X ×S S ′ , ave S ′ = Spf(k′ ◦ ), tels que X ′ soit une S {-} ourbe simplement}
\NormalTok{semi{-}stable et p−1}
\NormalTok{                 k ′ (Σ) ⊂ S0 (X ). Sous es onditions, la proposition 2.1.15 montre que le}
\NormalTok{                                 ′}

\NormalTok{morphisme f = f ⊗ 1 : Xη = Xη ⊗k k′ → Yη ⊗k k′ se prolonge en un S ′ {-}morphisme f : X ′ →}
\NormalTok{              ′            ′}

\NormalTok{Y ′ = Y ×S S ′ .}
\NormalTok{  Nous notons dorénavant X , Y , f et f en lieu et pla e de X ′ , Y ′ , f ′ et f ′ respe tivement.}
\NormalTok{(ii) Considérons deux points x ∈ Sing(Xs ), y ∈ Sing(Ys ) tels que f (x) = y et soient U ⊂ Xs ,}
\NormalTok{V ⊂ Ys des voisinages ouverts anes onnexes de x, y respe tivement tels que U − \{x\} et}
\NormalTok{V − \{y\} soient lisses. Par onstru tion des polyèdres entiers S(X ) et S(Y),}
                                                               
\NormalTok{                  S(X ) ∩ X|U η = S(X|U ) et S(Y) ∩ Y|V η = S(Y|V )}

\NormalTok{sont des polytopes entiers, de groupes stru turaux Λ(X|U ) et Λ(Y|V ) respe tivement. Il nous}
\NormalTok{faut nous assurer que, pour toute fon tion ϕ ∈ Λ(Y|V ), la fon tion S(f )∗ ϕ sur S(X|U ) appar{-}}
\NormalTok{tient à Λ(X|U ) ; e i est ontenu dans la formule}
\NormalTok{                                        S(f )∗ ϕ = Λ(f )(ϕ),}
\NormalTok{où Λ(f ) est l\textquotesingle{}homomorphisme de Λ(Y|V ) dans Λ(X|U ) induit par le morphisme f : X|U → Y|V}
\NormalTok{(2.1.1). La véri ation en est fa ile :}
\NormalTok{    il n\textquotesingle{}y a rien à dire si ϕ est onstante ;}
\NormalTok{    si la fon tion ϕ appartient au sous{-}groupe Λk (Y|V ) de Λ(Y|V ), elle est de la forme − log |u|,}
\NormalTok{                              ×}
\NormalTok{     u ∈ Γ(Y|V , OY ) ⊗k◦ k , et alors}
                                                                  
\NormalTok{                          S(f )∗ ϕ = f ∗ τY∗ ϕ |S(X ) = f ∗ log |u| |S(X )}
\NormalTok{                                    = log |f ∗ u||S(X )}
\NormalTok{                                    = Λ(f )(ϕ).}

\NormalTok{(iii) L\textquotesingle{}appli ation S(f ) : S(X ) → S(Y) étant par dénition la omposée τY ◦ f|S(X ) , le dia{-}}
\NormalTok{gramme}
\NormalTok{                                                 τX}
\NormalTok{                                         Xη           / S(X )}

\NormalTok{                                         f                    S(f )}
                                                         
\NormalTok{                                          Yη     τY}
\NormalTok{                                                      / S(Y)}

\NormalTok{est ommutatif si et seulement si l\textquotesingle{}appli ation τY ◦f est onstante sur les bres de la rétra tion}
\NormalTok{τX .                                                                                       }
\NormalTok{    Soit Ω l\textquotesingle{}ouvert Ys − Sing(Ys ) ; la rétra tion τY étant lo alement onstante sur Y|Ω η et}
                                                                                               
\NormalTok{     à valeurs dans S0 (Y), τY ◦ f est onstantes sur les bres de τY au{-}dessus de S Y|f −1 (Ω)}
\NormalTok{     puisque es dernières sont onnexes.}
\NormalTok{50}



\NormalTok{      Soient y un point dans Sing(Ys ) et V un voisinage ane onnexe de y dans Ys tel que}
\NormalTok{       V − \{y\} soit lisse. Étant donnée une fon tion ϕ ∈ Λ(Y|V ), la fon tion f ∗ ϕ = Λ(f )(ϕ)}
\NormalTok{       appartient à Λ(X|f −1 (V ) ) ; elle est don onstante sur les bres de τX . Comme}
\NormalTok{                                        f ∗ ϕ = (τY ◦ f )∗ ϕ|S(Y)}
\NormalTok{       et que les éléments de Λ(Y|V ) séparent les points de S(Y|V ), ela établit que l\textquotesingle{}appli ation}
                                                                          
\NormalTok{       τY ◦ f est onstante sur les bres de τX au{-}dessus de S X|f −1 (V ) .}
\NormalTok{(iv) Traitons tout d\textquotesingle{}abord le as d\textquotesingle{}un point x ∈ S(X ) appartenant à S0 (X ).}
\NormalTok{   Quitte à ee tuer un é latement admissible sur Y , on peut supposer que le point y ap{-}}
\NormalTok{partient à S0 (Y). Notons Cx (resp. Cy ) la omposante irrédu tible de Xs (resp. Ys ) de point}
\NormalTok{générique rX (x) (resp. rY (y)).}
\NormalTok{   Le morphisme f s : Xs → Ys induit un morphisme quasi{-}ni et dominant de Cx dans Cy et}
\NormalTok{l\textquotesingle{}homomorphisme orrespondant κ(rY (y)) → κ(rX (x)) entre les orps de fon tions oïn ide}
\NormalTok{ave elui provenant de l\textquotesingle{}homomorphisme H(y) → H(x), déni par f , via les identi ations}
\NormalTok{  anoniques}
\NormalTok{                             ] = κ(rX (x)),}
\NormalTok{                            H(x)                 ] = κ(rY (y))}
\NormalTok{                                                H(y)}
\NormalTok{                                          / ∂Xη est équivalente à la propreté de la e}
\NormalTok{(proposition 2.1.2, (iii)). L\textquotesingle{}hypothèse x ∈                                         k{-} ourbe}
\NormalTok{algébrique Cx ; il en est alors de même pour la ourbe Cy et le morphisme induit par f s de}
\NormalTok{Cx dans Cy est ni, de degré}
\NormalTok{                              ] : H(y)]}
\NormalTok{                             [H(x) ] = [H(x) : H(y)] = deg f}
\NormalTok{                                                          x}

\NormalTok{(l\textquotesingle{}extension H(x)/H(y) est non ramiée puisque, les S {-} ourbes X , Y étant simplement semi{-}}
\NormalTok{stables, |H(y)× | = |k× | = |H(x)× |).}
\NormalTok{   Soit p un point singulier de Ys appartenant à Cy ; à p orrespond un sous{-}polytope entier}
\NormalTok{P de S(Y) tel que y ∈ ∂P et on note t la fon tion ane positive sur P qui s\textquotesingle{}annule en x et}
\NormalTok{engendre le groupe abélien Λ(P )/R. Notons ep′ (f ) l\textquotesingle{}indi e de ramid ation de f en un point}
\NormalTok{fermé p′ de Cx .}
\NormalTok{   Nous avons d\textquotesingle{}une part}
\NormalTok{                                          ep′ (f ) = 1}
\NormalTok{pour tout point fermé ξ ′ ∈ Cx ∩ (Xs − Sing(Xs )) : il passe en eet né essairement deux}
\NormalTok{ omposantes irrédu tibles distin tes de Xs par tout point de Cx au{-}dessus d\textquotesingle{}un point de}
\NormalTok{Cy ∩ Sing(Ys ).}
\NormalTok{  D\textquotesingle{}autre part, la dénition de la stru ture entière sur S(X ) est telle que}
\NormalTok{                                        λx,p′ (S(f )∗ t) = ep′ (f )}
\NormalTok{pour tout point p′ ∈ Cx ∩ Sing(Xs ).}
\NormalTok{  Nous en déduisons :}
\NormalTok{                Z                              X}
\NormalTok{                    ddc S(f )∗ t =                           λx,p′ (S(f )∗ t)[κ(p′ ) : e}
\NormalTok{                                                                                       k]}
\NormalTok{                    \{x\}                  p′ ∈Cx ∩Sing(Xs )}
\NormalTok{                                               X}
\NormalTok{                                    =                        ep′ (f )[κ(p′ ) : κ(p)][κ(ξ) : e}
\NormalTok{                                                                                            k]}
\NormalTok{                                         p′ ∈Cx ∩Sing(Xs )}

\NormalTok{                                    = (deg (f |Cx ))[κ(p) : e}
\NormalTok{                                                            k]}
\NormalTok{                                                Z}
\NormalTok{                                    = (degx f )      ddc t}
\NormalTok{                                                      \{y\}}
\NormalTok{                                                                                               51}



\NormalTok{et l\textquotesingle{}égalité}
\NormalTok{                            Z                               Z}
\NormalTok{                                     c      ∗}
\NormalTok{                                   dd S(f ) Φ = (degx f )                 ddc Φ}
\NormalTok{                             \{x\}                                \{f (x)\}}

\NormalTok{pour toute fon tion Φ ∈ A0 (S(Y)) s\textquotesingle{}en déduit immédiatement par linéarité.}
\NormalTok{    Le as général se déduit du pré édent en appliquant la proposition 2.2.21 et le lemme 2.2.22.}
\NormalTok{2}


\NormalTok{                                   2.3. Fon tions harmoniques}
\NormalTok{2.3.1. L\textquotesingle{}espa e H(X ).}
\NormalTok{   Étant donnée une S {-} ourbe X , le bord de l\textquotesingle{}espa e k{-}analytique Xη (2.1.2) est le sous{-}}
\NormalTok{ensemble ∂Xη de S0 (X ) formé des points au{-}dessus des points génériques des omposantes}
\NormalTok{irrédu tibles anes de Xs .}
\NormalTok{Définition 2.3.1.  Soit X une S {-} ourbe simplement semi{-}stable ; le sous{-}espa e ve toriel}

\NormalTok{                                           τX∗ H(S(X ), ∂Xη )}
\NormalTok{de C0 (|Xη |, R) est noté H(X ).}

\NormalTok{   Nous nous restreindrons, dans e qui suit, à la onsidération de S {-} ourbes simplement}
\NormalTok{simplement semi{-}stables X qui sont quasi{-} ompa tes ; ette hypothèse équivaut à la ompa ité}
\NormalTok{de l\textquotesingle{}espa e topologique |Xη |.}
\NormalTok{Proposition 2.3.2}
\NormalTok{  (i ) Lorsque la e}
\NormalTok{                  k{-} ourbe Xs est propre, H(X ) est l\textquotesingle{}espa e des fon tions réelles lo alement}
\NormalTok{ onstantes sur Xη .}
\NormalTok{  (ii ) Lorsque ∂Xη ren ontre haque omposante onnexe de |Xη |, l\textquotesingle{}appli ation de restri tion}
\NormalTok{                                         H(X ) → Hom(∂Xη , R)}
\NormalTok{est bije tive.}
\NormalTok{Démonstration.}
\NormalTok{(i) Si la e}
\NormalTok{          k{-} ourbe Xs est propre, l\textquotesingle{}ensemble ∂Xη est vide ; l\textquotesingle{}assertion dé oule don du orollaire}
\NormalTok{1.2.11.}
\NormalTok{(ii) Si l\textquotesingle{}ensemble ∂Xη ren ontre ha une des omposantes onnexes de |Xη |, il ren ontre éga{-}}
\NormalTok{lement haque omposante onnexe de S(X ) (remarque 2.2.19,3) et l\textquotesingle{}assertion dé oule de la}
\NormalTok{proposition 1.2.15.                                                                           2}

\NormalTok{    Comme annon é en introdu tion de e hapitre, nous allons établir le résultat suivant.}
\NormalTok{Proposition 2.3.3.  Soient X , X}
\NormalTok{                                       ′ deux S {-} ourbes simplement semi{-}stables et f : X ′ →}
\NormalTok{                                                                                             η}
\NormalTok{Xη un morphisme quasi{-}ni. L\textquotesingle{}appli ation R{-}linéaire f ∗ : C0 (|Xη |, R) → C0 (|Xη′ |, R) envoie le}
\NormalTok{sous{-}espa e H(X ) dans le sous{-}espa e H(X ′ ).}
\NormalTok{Démonstration. Vu la proposition 2.2.26, le résultat s\textquotesingle{}obtient en ombinant les trois lemmes}
\NormalTok{suivants.                                                                                       2}
\NormalTok{Lemme 2.3.4.  Soient X         une S {-} ourbe simplement semi{-}stable et q : Y → X un é late{-}}
\NormalTok{ment admissible tel que Y soit une S {-} ourbe simplement simplement semi{-}stable ; l\textquotesingle{}isomor{-}}
\NormalTok{                            ∼}
\NormalTok{phisme qη∗ : C0 (|Xη |, R) −→ C0 (|Yη |, R) identie les sous{-}espa es H(X ) et H(Y).}
\NormalTok{52}



\NormalTok{Démonstration. L\textquotesingle{}isomorphisme qη identie ∂Yη et ∂Xη .}
\NormalTok{  D\textquotesingle{}après la proposition 2.2.26, le polyèdre entier S(X ) s\textquotesingle{}identie à un domaine polyédral}
\NormalTok{dans le polyèdre entier S(Y) tel que, posant S(q) = τX ◦ qη |S(Y) , le diagramme}
\NormalTok{                                                    τY}
\NormalTok{                                           Yη            / S(Y)}

\NormalTok{                                       qη                           S(q)}
                                                               
\NormalTok{                                           Xη       τX}
\NormalTok{                                                         / S(X )}

\NormalTok{soit ommutatif, haque omposante onnexe de S(Y) − S(X ) ne ren ontre S(X ) qu\textquotesingle{}en un}
\NormalTok{seul point et l\textquotesingle{}appli ation S(q) soit la rétra tion envoyant ha une de es omposantes sur le}
\NormalTok{point orrespondant de S(X ).}
\NormalTok{   Les espa es k{-}analytiques Yη et Xη sont identiés via l\textquotesingle{}isomorphisme qη ; on déduit du}
\NormalTok{diagramme ommutatif pré édent que, pour toute fon tion F sur S(X ), la fon tion f = F ◦τX}
\NormalTok{sur Xη = Yη se fa torise par τX sous la forme :}
\NormalTok{                                     f = (F ◦ S(q)) ◦ τY .}
\NormalTok{L\textquotesingle{}appli ation linéaire}
\NormalTok{                              S(q)∗ : A0 (S(X )) → A0 (S(Y))}
\NormalTok{induisant un isomorphisme de H(S(X ), ∂Xη ) sur H(S(Y), ∂Yη ) (proposition 1.2.20), il en}
\NormalTok{résulte que l\textquotesingle{}isomorphisme anonique de C0 (|Xη |, R) sur C0 (|Yη |, R) envoie le sous{-}espa e}
\NormalTok{H(X ) dans H(Y) et induit un isomorphisme de l\textquotesingle{}un sur l\textquotesingle{}autre.                            2}
\NormalTok{Lemme 2.3.5.  Soient X         une S {-} ourbe simplement semi{-}stable (quasi{-} ompa te ) et K/k}
\NormalTok{une extension non{-}ar himédienne ; une fon tion h ∈ C0 (|Xη |, R) appartient au sous{-}espa e}
\NormalTok{H(X ) si et seulement si la fon tion p∗K h appartient au sous{-}espa e H(X ×S Spf(K ◦ )).}
\NormalTok{Démonstration. Le morphisme anonique pK : Xη ⊗            b k K → Xη induit un isomorphisme}
\NormalTok{G{-}lo al surje tif S(X )K → S(X ) de polyèdres entiers pondérés (proposition 2.2.21) et, les}
\NormalTok{ omposantes irrédu tibles anes de la K   e {-} ourbe Xs ⊗e K   e étant les images ré iproques des}
\NormalTok{                                                           k}
\NormalTok{                                                             −1}
\NormalTok{ omposantes irrédu tibles anes de Xs , ∂(Xη ⊗  b k K) = pK (∂Xη )) ∩ S0 (X ×S Spf(K ◦ )). Ce i}
\NormalTok{implique, pour tout fon tion Φ ∈ A (S(X )), que la fon tion p∗K Φ appartient à A0 (S(X )K ) et}
\NormalTok{                                     0}

\NormalTok{que le support de la mesure ddc Φ = (pK )∗ ddc p∗K Φ est ontenu dans ∂Xη si et seulement si}
\NormalTok{ elui de la mesure ddc p∗K Φ est ontenu dans ∂(Xη ⊗  b k K).                                  2}
\NormalTok{Lemme 2.3.6.  Soient X , Y deux S {-} ourbes simplement semi{-}stables et f : Yη → Xη un}
\NormalTok{morphisme quasi{-}ni ; l\textquotesingle{}homomorphisme f ∗ : C0 (|Yη |, R) → C0 (|Xη |, R) envoie le sous{-}espa e}
\NormalTok{H(Y) dans H(X ).}
\NormalTok{Démonstration. Appliquant la proposition 2.2.27 et les deux lemmes pré édents, on peut}
\NormalTok{supposer que f induit un morphisme de Y dans X . On a alors un diagramme ommutatif}
\NormalTok{                                                    τX}
\NormalTok{                                       Xη                / S(X )}

\NormalTok{                                       f                            S(f )}
                                                           
\NormalTok{                                       Yη           τY}
\NormalTok{                                                         / S(Y),}

\NormalTok{où S(f ) = τY ◦ f est un morphisme de polyèdres entiers.}
\NormalTok{  Étant donnée une fon tion H dans H(S(Y), ∂Yη ),}
\NormalTok{                                  f ∗ (H ◦ τY ) = S(f )∗ H ◦ τX}
\NormalTok{                                                                                                53}



\NormalTok{et S(f )∗ H ∈ A0 (S(X )). Il reste à vérier que la fon tion S(f )∗ H est harmonique sur S(X ) −}
\NormalTok{∂Xη . C\textquotesingle{}est une onséquen e dire te du point (iv) de la proposition 2.2.27 : f (x) ∈ / ∂Yη pour}
\NormalTok{tout point x ∈ S(X ) − ∂Xη et}
\NormalTok{                         Z                              Z}
\NormalTok{                               ddc S(f )∗ H = (degx f )       ddc H = 0.}
\NormalTok{                          \{x\}                             \{f (x)\}}
\NormalTok{                                                                                                 2}
\NormalTok{  Si X est un espa e stri tement k{-}anoïde et ϕ : X → ˜ Xη est un isomorphisme de X sur la}
\NormalTok{bre générique d\textquotesingle{}une S {-} ourbe simplement semi{-}stable X , le sous{-}espa e ve toriel réel ϕ∗ H(X )}
\NormalTok{de C0 (|X|, R) ne dépend pas du hoix de X et de ϕ ; il est noté H(X).}
\NormalTok{   De manière générale, étant donnés un espa e k{-}analytique X et une lture algébrique ka de}
\NormalTok{k, le groupe de Galois G de la fermeture séparable ksep de k dans ka opère par automorphismes}
\NormalTok{sur l\textquotesingle{}espa e c                b kc}
\NormalTok{             ka {-}analytique X ⊗  ka . On dispose d\textquotesingle{}un morphisme G{-}équivariant anonique}
\NormalTok{                                          p : X⊗  ca → X}
\NormalTok{                                               b kk}
\NormalTok{d\textquotesingle{}espa es analytiques et l\textquotesingle{}appli ation ontinue sous{-}ja ente est un quotient de |X ⊗        ca | par}
\NormalTok{                                                                                        b kk}
\NormalTok{l\textquotesingle{}a tion de G dans la atégorie des espa es topologiques : \textquotesingle{}est en eet un quotient dans la}
\NormalTok{  atégorie des ensembles d\textquotesingle{}après le orollaire 1.3.6 de [3[centerline] et l\textquotesingle{}appli ation ontinue p est ouverte.}
\NormalTok{Proposition 2.3.7.  Quelles que soient la S {-} ourbe simplement semi{-}stable X et l\textquotesingle{}ex{-}}
\NormalTok{tension nie galoisienne k′ /k, l\textquotesingle{}a tion du groupe G = Gal(k′ /k) sur C0 (|Xη ⊗k k′ |, R) stabilise}
\NormalTok{le sous{-}espa e H(Xη ) et l\textquotesingle{}isomorphisme anonique}
\NormalTok{                                 C0 (|Xη |, R) → C0 (|Xη ⊗k k′ |, R)G}
\NormalTok{induit un isomorphisme de H(Xη ) sur H(Xη ⊗k k′ )G .}
\NormalTok{Démonstration. On peut supposer Xη onnexe.}
\NormalTok{   Si Xη est propre, H(Xη ) (resp. H(Xη ⊗k k′ )) est l\textquotesingle{}espa e des fon tions onstantes sur Xη}
\NormalTok{(resp. lo alement onstantes sur Xη ⊗k k′ ) (proposition 2.3.2, 1) et l\textquotesingle{}assertion est évidente.}
\NormalTok{   Si Xη n\textquotesingle{}est pas propre, alors il en est de même pour haque omposante onnexe de Xη ⊗k k′}
\NormalTok{et, d\textquotesingle{}après la proposition 2.3.2,2, les restri tions}
\NormalTok{                                      H(Xη ) → Hom(∂Xη , R)}
\NormalTok{et}
\NormalTok{                                H(Xη ⊗k k′ ) → Hom(∂(Xη ⊗k k′ , R)}
\NormalTok{sont bije tives. L\textquotesingle{}assertion résulte ainsi de l\textquotesingle{}identité ∂Xη = ∂(Xη ⊗k k′ )G .                   2}

\NormalTok{2.3.2. Le théorème de rédu tion semi{-}stable.}
\NormalTok{   C\textquotesingle{}est i i qu\textquotesingle{}intervient le théorème de rédu tion semi{-}stable.}
\NormalTok{Théorème 2.3.8.  Soit X un espa e stri tement k {-}anoïde purement de dimension 1}
\NormalTok{et rig{-}lisse ; il existe une extension nie séparable k′ de k telle que le domaine k′ {-}anoïde}
\NormalTok{X ⊗k k′ soit isomorphe à la bre générique d\textquotesingle{}une Spf(k′ ◦ ){-} ourbe simplement semi{-}stable.}
\NormalTok{Démonstration. Il existe une extension nie séparable k1 de k et une ourbe algébrique C}
\NormalTok{propre et lisse sur k1 telles que X ⊗k k1 soit isomorphe à un domaine anoïde de C ([31[centerline],}
\NormalTok{theorem 1). Appliquant le théorème de rédu tion semi{-}stable sous la forme rigide{-}analytique}
\NormalTok{établie par S. Bos h et W. Lütkebohmert dans [9[centerline], il existe une extension nie séparable k′}
\NormalTok{de k1 telle que C ⊗k1 k′ soit isomorphe à la bre générique d\textquotesingle{}une Spf(k′ ◦ ){-} ourbe C dont la}
\NormalTok{bre spé iale est simplement semi{-}stable ; \textquotesingle{}est automatiquement une S {-} ourbe semi{-}stable}
\NormalTok{54}



\NormalTok{(remarque 2.2.9, 3). En hoisissant onvenablement k′ , on peut en outre supposer que tout}
\NormalTok{point x du bord de X ⊗k k′ satisfait à la ondition |κ(x)× | = |k′ × |.}
\NormalTok{    Envisageons X ′ = X ⊗k k′ omme un domaine stri tement k′ {-}anoïde de Cη et soit q :}
\NormalTok{C ′ → C un é latement admissible tel que C ′ soit une Spf(k′ ◦ ){-} ourbe simplement semi{-}stable,}
\NormalTok{ave ∂X ′ ⊂ S0 (C ′ ) (lemme 2.2.22). L\textquotesingle{}image de X ′ par l\textquotesingle{}appli ation de rédu tion Cη′ → Cs}
\NormalTok{est un ouvert U et X ′ est isomorphe à la bre générique de la Spf(k′ ◦ ){-} ourbe simplement}
\NormalTok{semi{-}stable X = C|U′ .                                                                       2}
\NormalTok{  Soient X un espa e stri tement k{-}anoïde purement de dimension 1 rig{-}lisse et k′ , k′′ deux}
\NormalTok{extensions nies galoisiennes de k telles que X ⊗k k′ et X ⊗k k′′ soient isomorphes aux bres}
\NormalTok{génériques de ourbes stri tement semi{-}stables ; étant donnée une extension nie galoisienne}
\NormalTok{K de k oiant k′ et k′′ , la proposition 2.3.7 implique l\textquotesingle{}identité des espa es H(X ⊗k k′ ) et}
\NormalTok{                   ′                                          ′′}
\NormalTok{H(X ⊗k K)Gal(K/k ) (resp. H(X ⊗k k′′ ) et H(X ⊗k K)Gal(K/k ) ), d\textquotesingle{}où dé oule l\textquotesingle{}égalité}
\NormalTok{                               ′                                              ′′}
\NormalTok{              H(X ⊗k k′ )Gal(k /k) = H(X ⊗k K)Gal(K/k) = H(X ⊗k k′′ )Gal(k /k) .}
\NormalTok{Le théorème pré édent permet ainsi de dénir, pour tout espa e stri tement k{-}anoïde X}
\NormalTok{purement de dimension 1 et rig{-}lisse, un sous{-}espa e ve toriel H(X) de C0 (|X|, R) : H(X) =}
\NormalTok{                ′}
\NormalTok{H(X ⊗k k′ )Gal(k /k) ⊂ C0 (|X|, R), k′ étant une extension nie galoisienne de k telle que X ⊗k k′}
\NormalTok{soit isomorphe à la bre générique d\textquotesingle{}une Spf(k′ ◦ ){-} ourbe simplement semi{-}stable.}
\NormalTok{   Désignant par C la atégorie dont les objets sont les espa es stri tement k{-}anoïdes, pure{-}}
\NormalTok{ment de dimension 1 et rig{-}lisses, et les morphismes les immersions anoïdes, il dé oule de la}
\NormalTok{proposition 2.3.3 que l\textquotesingle{}on vient ainsi de dénir un fon teur H sur Cop à valeurs dans VectR .}
\NormalTok{  Nous aurons besoin des deux résultats suivants.}
\NormalTok{Lemme 2.3.9.  Pour tout espa e stri tement k {-}anoïde Y rig{-}lisse et purement de di{-}}
\NormalTok{mension 1, l\textquotesingle{}appli ation de restri tion}
\NormalTok{                                   C0 (|Y |, R) → Hom(Γ(Y ), R)}
\NormalTok{induit un isomorphisme du sous{-}espa e H(Y ) sur Hom(Γ(Y ), R).}
\NormalTok{Démonstration. Soit k′ une extension galoisienne nie de k telle que l\textquotesingle{}espa e stri tement}
\NormalTok{k′ {-}anoïde Y ⊗k k′ soit isomorphe à la bre générique d\textquotesingle{}une Spf(k′ ◦ ){-} ourbe simplement semi{-}}
\NormalTok{                                                                                            ′}
\NormalTok{stable Y . La proje tion anonique p : Y ⊗k k′ → Y induit une bije tion de Γ(Y ⊗k k′ )Gal(k /k)}
\NormalTok{sur Γ(Y ) et l\textquotesingle{}appli ation de restri tion}
\NormalTok{                                   C0 (|Y |, R) → Hom(Γ(Y ), R)}
\NormalTok{ oïn ide ave l\textquotesingle{}appli ation induite par la restri tion}
\NormalTok{                            C0 (|Y ⊗k k′ |, R) → Hom(Γ(Y ⊗k k′ ), R)}
\NormalTok{                                                                                           ′}
\NormalTok{entre les sous{-}espa es des invariants sous Gal(k′ /k) ; puisque H(Y ) = H(Y ⊗k k′ )Gal(k /k) , la}
\NormalTok{ on lusion dé oule de la proposition 2.3.2.                                                    2}
\NormalTok{Lemme 2.3.10.  Soient Y un espa e stri tement k {-}anoïde onnexe, de dimension 1 et}
\NormalTok{rig{-}lisse, et (hn ) une suite roissante de fon tions appartenant à H(Y ). L\textquotesingle{}alternative suivante}
\NormalTok{est vraie :}
\NormalTok{    soit la suite (hn ) onverge, uniformément sur Y , vers une fon tion h ∈ H(Y ),}
\NormalTok{    soit la suite (hn ) onverge, lo alement uniformément sur Y − Γ(Y ), vers ∞.}
\NormalTok{Démonstration. Soient k′ une extension nie galoisienne de k telle que Y ′ = Y ⊗k k′ soit}
\NormalTok{isomorphe à la bre générique d\textquotesingle{}une Spf(k′ ◦ ){-} ourbe simplement nodale Y′ et p : Y ′ → Y la}
\NormalTok{proje tion anonique. Si l\textquotesingle{}assertion est démontrée pour la suite (p∗ hn ) dans H(Y ′ ), elle l\textquotesingle{}est}
\NormalTok{                                                                      ′}
\NormalTok{également pour la suite (hn ) : en eet, puisque H(Y ) = H(Y ′ )Gal(k /k) , il sut de s\textquotesingle{}assurer}
\NormalTok{                                                                                               55}



\NormalTok{que la fon tion limite h ∈ H(Y ′ ) du premier as de l\textquotesingle{}alternative est invariante sous Gal(k′ /k) ;}
\NormalTok{utilisant l\textquotesingle{}isomorphisme Gal(k′ /k){-}équivariant de restri tion}
\NormalTok{                                   H(Y ′ ) → Hom(Γ(Y ′ ), R),}
\NormalTok{                                                    ′}
\NormalTok{ ela dé oule de l\textquotesingle{}identité entre Γ(Y ) et Γ(Y ′ )Gal(k /k) .}
\NormalTok{   Vu la dénition de l\textquotesingle{}espa e H(Y ′ ) = H(Y ′ ) (dénition 2.3.1), le résultat pour une suite}
\NormalTok{ roissante dans H(Y ′ ) s\textquotesingle{}obtient par appli ation de la proposition 1.2.12.                 2}
\NormalTok{2.3.3. Le fais eau HX .}
\NormalTok{  Soit X une ourbe stri tement k{-}analytique lisse. Le fon teur H : Cop → VectR donne}
\NormalTok{naturellement naissan e à un préfais eau HX sur le site rigide XR de X (2.1.1) : étant donné}
\NormalTok{un objet V de XR , les domaines stri tement k{-}anoïdes de X ontenus dans V forment un}
\NormalTok{système indu tif ltrant I(V ) (les morphismes sont les in lusions) et l\textquotesingle{}on pose}
\NormalTok{                                      HX (V ) = ←}
\NormalTok{                                                lim}
\NormalTok{                                                  −I(V ) H.}
\NormalTok{Remarque 2.3.11}
\NormalTok{  1. Pour tout domaine stri tement k{-}anoïde Y dans X , l\textquotesingle{}homomorphisme anonique}
\NormalTok{     HX (Y ) → H(Y ) est un isomorphisme.}
\NormalTok{  2. Le préfais eau HX sur le site XR n\textquotesingle{}est lairement pas un fais eau. Il sut en eet de}
\NormalTok{      onsidérer, sur X = P1k = A1k ∪ \{∞\}, le re ouvrement ni par les domaines stri tement}
\NormalTok{     k{-}anoïdes Y1 = \{|T1 ≤ r\}, Y2 = \{r ≤ |T1 | ≤ r −1 \} et Y3 = \{|T1 | ≥ r −1 \}  r étant un}
\NormalTok{     élément de ]0, 1[∩|k× |  et la fon tion ontinue h sur X telle que}
\end{Highlighting}
\end{Shaded}

\subsection{Thuillier, newline lines
3550--3765}\label{thuillier-newline-lines-35503765}

\begin{Shaded}
\begin{Highlighting}[]

\NormalTok{  Étant donné un domaine k{-}anoïde Y d\textquotesingle{}une ourbe k{-}analytique lisse X , nous noterons}
\NormalTok{H(Y ) le sous{-}espa e ve toriel}
\NormalTok{                                   Γ(Y − ∂Y, HX ) ∩ C0 (Y, R)}
\NormalTok{de C0 (|Y |, R) ; vu le orollaire 2.3.14, ette notation oïn ide ave elle introduite pré édem{-}}
\NormalTok{ment lorsque Y est stri tement k{-}anoïde.}
\NormalTok{   La orrespondan e Y 7→ H(Y ) est un préfais eau sur le G{-}site XG de X qui prolonge}
\NormalTok{évidemment le préfais eau sur le site XR déni au début de 2.3.3.}
\NormalTok{Proposition 2.3.18.  Soient X une ourbe stri tement k {-}analytique lisse et K une ex{-}}
\NormalTok{tension non ar himédienne de k ; notant XK pour X ⊗       b k K et pK le morphisme anonique}
\NormalTok{XK → X , pK HX ⊂ HX ⊗}
\NormalTok{              ∗}
\NormalTok{                          b kK .}
\NormalTok{Démonstration. C\textquotesingle{}est une onséquen e immédiate de la dénition du fais eau HX et du}
\NormalTok{lemme 2.3.5.                                                                                      2}
\NormalTok{   A hevons e paragraphe par le résultat suivant :}
\NormalTok{Proposition 2.3.19.  Soient X, Y deux ourbes stri tement k {-}analytiques lisses et f :}
\NormalTok{Y → X un morphisme. Pour toute fon tion harmonique h ∈ Γ(X, HX ), la fon tion f ∗ h = h◦f}
\NormalTok{est harmonique sur Y : f ∗ h ∈ Γ(Y, HY ).}
\NormalTok{Démonstration. On peut supposer Y onnexe et f quasi{-}ni. Étant donnés des domaines}
\NormalTok{stri tement k{-}anoïdes W ⊂ Y et V ⊂ X tels que f (W ) ⊂ V , h|V ∈ H(V ) et don , par appli{-}}
\NormalTok{ ation des propositions 2.3.3 et 2.3.7, (f ∗ h)|W ∈ H(W ). La fon tion f ∗ h est don harmonique}
\NormalTok{sur Y .                                                                                      2}

\NormalTok{2.3.4. Le préfais eau log |OX× |.}
\NormalTok{  Considérant toujours une ourbe stri tement k{-}analytique lisse X , examinons mainte{-}}
\NormalTok{nant les liens entre les fon tions harmoniques sur X et les fon tions de la forme log |f |,}
\NormalTok{            ×}
\NormalTok{f ∈ Γ(X, OX   ).}
\NormalTok{                                                   ×}
\NormalTok{Proposition 2.3.20.  Pour toute se tion f ∈ Γ(X, OX ), la fon tion log |f | est harmo{-}}
\NormalTok{nique sur X .}
\NormalTok{Démonstration. Il sut de vérier que la restri tion de la fon tion log |f | à tout domaine}
\NormalTok{stri tement k{-}anoïde Y ⊂ X appartient au sous{-}espa e H(Y ) ⊂ C0 (|Y |, R) ; vu le lemme}
\NormalTok{2.3.5 et le théorème 2.3.8, on peut en outre supposer que Y est la bre générique d\textquotesingle{}une S {-}}
\NormalTok{ ourbe simplement semi{-}stable, auquel as il ne reste plus qu\textquotesingle{}à invoquer le point (iii) de la}
\NormalTok{proposition 2.2.24.                                                                              2}
\NormalTok{Théorème 2.3.21.  Soit X une ourbe stri tement k {-}analytique lisse. Le fais eau de R{-}}
\NormalTok{                                                                                                 ×}
\NormalTok{ve toriels engendré par les germes de fon tions log |f |, f étant un germe de se tion de OX        ,}
\NormalTok{est un sous{-}fais eau de HX ; le quotient est un fais eau gratte{-} iel sur |X| et il est nul si l\textquotesingle{}une}
\NormalTok{des deux onditions suivantes est vériée :}
\NormalTok{58}



\NormalTok{      le orps e}
\NormalTok{               k est algébrique sur un orps ni ;}
\NormalTok{      la ourbe X ⊗  ca est lo alement isomorphe à P1 .}
\NormalTok{                  b kk                               a        k}
\NormalTok{Démonstration. Soit FX le sous{-}fais eau de CX asso ié au préfais eau sur |X| faisant orres{-}}
\NormalTok{pondre à un ouvert U le sous{-}espa e ve toriel de C0 (U, R) engendré par les fon tions log |f |,}
\NormalTok{               ×}
\NormalTok{f ∈ Γ(U, OX      ) ; \textquotesingle{}est un sous{-}fais eau de HX en vertu de la proposition pré édente. Les germes}
\NormalTok{de FX et HX oïn ident en tout point x de X de type (1), (3) ou (4) :}
\NormalTok{     si x est de type (1) ou (4), e point admet un système fondamental de voisinages k{-}}
\NormalTok{       anoïdes ayant un unique point de Shilov et les se tions de FX et HX sont don}
\NormalTok{        onstantes au voisinage de x ;}
\NormalTok{     si x est de type (3), e point admet un système fondamental de voisinages k{-}anoïdes}
\NormalTok{       ayant exa tement deux points de Shilov ; les espa es ve toriels FX,x et HX,x sont alors}
\NormalTok{       de dimension 2, engendrés par les germes de la fon tion onstante 1 et de la fon tion}
\NormalTok{       log |f |, f étant une se tion quel onque de OX au voisinage de x telle que |f | ne soit pas}
\NormalTok{       lo alement onstante en x.}
\NormalTok{    Soient maintenant x ∈ X un point de type (2), k′ une extension nie galoisienne de k et x′ un}
\NormalTok{point de X ′ = X⊗k k′ au{-}dessus de x : l\textquotesingle{}égalité FX ′ ,x′ = HX ′ ,x′ implique l\textquotesingle{}égalité FX,x = HX,x .}
\NormalTok{En eet, si V est un voisinage stri tement k{-}anoïde de x dans X , V ′ est la omposante}
\NormalTok{ onnexe de x′ dans V ⊗k k′ , et f ′ ∈ A×                                                        ×}
\NormalTok{                                                    V ′ , la norme f de f est un élément de AV tel que}
\NormalTok{                                                                              ′}
\NormalTok{                      ′}
\NormalTok{p∗ |f ||V ′ = |f ′ |[k :k] ; si une fon tion h ∈ H(V ) est telle que p∗ h|V ′ soit une ombinaison linéaire}
\NormalTok{de fon tions log |f |, f ∈ A×                                                                           ×}
\NormalTok{                                    V ′ , \textquotesingle{}est don une ombinaison linéaire de fon tions log |f |, f ∈ AV .}
\NormalTok{Cette remarque préliminaire permet de supposer que X est isomorphe à la bre générique}
\NormalTok{d\textquotesingle{}une S {-} ourbe simplement semi{-}stable X et que le point x orrespond à une omposante}
\NormalTok{irrédu tible propre et géométriquement onnexe Cx de Xs .}
\NormalTok{Lemme 2.3.22.  Ave les notations en vigueur, l\textquotesingle{}espa e ve toriel HX,x /FX,x est anoni{-}}
\NormalTok{quement isomorphe à l\textquotesingle{}espa e ve toriel Pic0 (Cx ) ⊗Z R.}
\NormalTok{Démonstration. Soit V un voisinage stri tement k{-}anoïde de X . Il existe une extension}
\NormalTok{galoisienne nie k′ de k et un é latement admissible q : X ′ → X ⊗S Spf(k′ ◦ ) tel que V ⊗k k′}
\NormalTok{soit isomorphe à la bre générique du s héma formel induit par X ′ sur un ouvert U de Xs′ . La}
\NormalTok{ omposante irrédu tible Cx de Xs étant supposée géométriquement onnexe, il y a un unique}
\NormalTok{point x′ ∈ X ⊗k k′ au{-}dessus de x, orrespondant à la omposante irrédu tible Cx ⊗ek ke′ de}
\NormalTok{X ⊗S Spf(S ′ ) et, elle{-} i étant propre et lisse, le morphisme q induit un isomorphisme de la}
\NormalTok{ omposante irrédu tible de X ′ asso iée à x′ sur Cx ⊗ek ke′ .}
\NormalTok{                                                               ′}
\NormalTok{    À toute fon tion h ∈ H(V ) = H(V ⊗k k′ )Gal(k /k) orrespond naturellement un R{-}diviseur}
\NormalTok{div(h) sur la ke′ {-} ourbe Cx ⊗ek ke′ , vue omme omposante irrédu tible de Xs′ : en eet, l\textquotesingle{}ensemble}
\NormalTok{Tx′ S(X ′ ) est anoniquement identié à l\textquotesingle{}ensemble des points singuliers de Xs′ appartenant à}
\NormalTok{Cx ⊗ek ke′ et, notant x      et le point orrespondant à l\textquotesingle{}élément t de Tx′ S(X ′ ), on pose :}
\NormalTok{                                                          X}
\NormalTok{                                            div(h) =             λx,t (h)[e}
\NormalTok{                                                                          xt ].}
\NormalTok{                                              t∈Tx′ S(X ′ )}

\NormalTok{La ondition d\textquotesingle{}harmoni ité de h en x′ équivaut à la ondition}
\NormalTok{                                            deg div(h) = 0.}
\NormalTok{   Le groupe de Galois Gal(k′ /k) opère naturellement sur la ke′ {-} ourbe Cx ⊗ek ke′ et son a tion}
\NormalTok{est transitive sur haque bre du morphisme anonique Cx ⊗ek ke′ → Cx ; puisque h est, par}
\NormalTok{hypothèse, invariante sous Gal(k′ /k), ela implique que le R{-}diviseur div(h) sur Cx ⊗ek ke′}
\NormalTok{des end en un R{-}diviseur de degré 0 sur sur Cx , en ore noté div(h).}
\NormalTok{                                                                                                59}



\NormalTok{  Nous obtenons ainsi une appli ation linéaire div : HX,x → Div(Cx ) ⊗Z R s\textquotesingle{}insérant dans}
\NormalTok{une suite}
\NormalTok{                  0       /R      / HX,x   div   / Div(Cx ) ⊗Z R deg   /R     / 0,}

\NormalTok{laquelle est exa te. Seule la surje tivité de div sur le sous{-}espa e des R{-}diviseurs de degré 0}
\NormalTok{requiert une démonstration. Soit}
\NormalTok{                                              X}
\NormalTok{                                        D=        n(e}
\NormalTok{                                                    x)[ex]}
\NormalTok{                                             e∈|D|}
\NormalTok{                                             x}

\NormalTok{un R{-}diviseur de degré 0 sur Cx ; soit I l\textquotesingle{}Idéal (admissible) dénissant le sous{-}s héma fermé}
\NormalTok{réduit Z = |D|∩(Xs −Sing(Xs )) dans Xs et soit q : X ′ → X son é latement. Puisque X est lisse}
\NormalTok{aux points de Z , X ′ est en ore une S {-} ourbe simplement semi{-}stable et on note Y l\textquotesingle{}ouvert de}
\NormalTok{X ′ obtenu en retirant, à ha une des omposantes irrédu tibles du diviseur ex eptionnel de X ′ ,}
\NormalTok{un point fermé disjoint (du transformé stri t) de X . La S {-} ourbe Y est simplement semi{-}stable}
\NormalTok{et Yη est un voisinage de x dans X = Xη . L\textquotesingle{}ensemble Tx S(Y) s\textquotesingle{}indentie anoniquement à}
\NormalTok{une partie nie de l\textquotesingle{}ensemble des points fermés de Cx et nous avons fait e qu\textquotesingle{}il fallait pour}
\NormalTok{que}
\NormalTok{                                         |D| ⊂ Tx S(Y).}
\NormalTok{Considérons alors une fon tion ane par mor eaux H sur S(Y) telle que, pour tout polytope}
\NormalTok{entier P ⊂ S(Y) de dimension un orrespondant à un élément t ∈ Tx S(Y),}
\NormalTok{{-} λx,t (H) = 0 si x  / |D|,}
\NormalTok{                  et ∈}
\NormalTok{{-} λx,t (H) = n(e}
\NormalTok{               xt ) si x}
\NormalTok{                       et ∈ |D|.}
\NormalTok{   Une telle fon tion H est harmonique au point x en vertu de la nullité du degré de D ; elle}
\NormalTok{est don harmonique sur un voisinage de x dans S(Y) et la fon tion h = τY∗ H sur X est par}
\NormalTok{ onséquent harmonique au voisinage de x dans X . L\textquotesingle{}identité}
\NormalTok{                                           div(h) = D}
\NormalTok{est vériée par onstru tion et nous avons ainsi établi l\textquotesingle{}exa titude de la suite onsidérée}
\NormalTok{ i{-}dessus.}
\NormalTok{  Soit Pr(Cx ) le groupe des diviseurs prin ipaux sur Cx . On dispose d\textquotesingle{}une autre suite exa te}

\NormalTok{                                  /R       / FX,x    div   / Pr(Cx )   / 0,}
\NormalTok{                           0}
\NormalTok{                                            ×}
\NormalTok{ ar, pour tout germe log |f | ∈ FX,x , f ∈ OX,x , le diviseur de la fon tion harmonique N log |f |}
\NormalTok{s\textquotesingle{}identie à elui de la fon tion méromorphe induite par f N /α sur Cx , où N est un entier}
\NormalTok{supérieur à 1 tel que |f (x)|N ∈ |k× | et α est un élément de k tel que |f (x)|N = |α|. Il en}
\NormalTok{résulte un isomorphisme d\textquotesingle{}espa es ve toriels réels}
\NormalTok{                                                 ∼}
\NormalTok{                                HX,x /FX,x −→ Pic0 (Cx ) ⊗Z R,}
\NormalTok{ e qui a hève la preuve du lemme.                                                                2}
\NormalTok{   Terminons la démonstration du théorème. Le fais eau HX /FX est un fais eau gratte{-} iel}
\NormalTok{dont le support est un ensemble dis ret de points de type (2) ; sa nullité équivaut à la propriété}
\NormalTok{suivante : pour tout point x ∈ X ⊗k c ka de type (2), le groupe de Pi ard Pic0 (Cx ) de la e   ka {-}}
\NormalTok{                                                                   g , est de torsion. Il y a au}
\NormalTok{ ourbe irrédu tible, propre et lisse Cx , de orps des fon tions κ(x)}
\NormalTok{moins deux onditions impliquant ette propriété :}
\NormalTok{    le orps ek est algébrique sur un orps ni : Pic0 (Cx ) est l\textquotesingle{}ensemble des e  ka {-}points d\textquotesingle{}un}
\NormalTok{     e}
\NormalTok{     ka {-}s héma en groupes de type ni et don est de ardinal ni ;}
\NormalTok{60}



\NormalTok{      pour tout point x ∈ X ⊗k k ca de type (2), la eka {-} ourbe irrédu tible, propre et lisse de}
\NormalTok{                           g est rationnelle  de manière équivalente, X ⊗}
\NormalTok{        orps des fon tions κ(x)                                             b kb}
\NormalTok{                                                                               ka est lo alement}
\NormalTok{       isomorphe à Pba (si X est irrédu tible et propre, \textquotesingle{}est une ourbe de Mumford ).}
\NormalTok{                     1}
\NormalTok{                     k}
\NormalTok{                                                                                              2}
\NormalTok{                                                                                              61}



\NormalTok{                              3. Théorie du potentiel}



\NormalTok{                              3.1. Le problème de Diri hlet}
\NormalTok{  Dans e qui suit, X est une ourbe stri tement k{-}analytique lisse.}
\NormalTok{3.1.1. Propriétés fondamentales des fon tions harmoniques.}
\NormalTok{  Les se tions du fais eau HX satisfont au prin ipe du maximum.}
\NormalTok{Proposition 3.1.1.  Étant donné un ouvert Ω de X , une fon tion h ∈ Γ(Ω, HX ) admet}
\NormalTok{un extremum lo al en un point x de Ω si et seulement si elle est onstante au voisinage de x.}
\NormalTok{  C\textquotesingle{}est la onjon tion de la proposition 2.3.13 et du orollaire 2.3.14.}
\NormalTok{   Le problème de Diri hlet est lo alement résoluble : tout point x ∈ X admet un système}
\NormalTok{fondamental de voisinages ouverts relativement ompa ts Ω tels que l\textquotesingle{}appli ation de restri tion}
\NormalTok{                              C0 (Ω, R) ∩ Γ(Ω, HX ) → C0 (∂Ω, R)}
\NormalTok{soit un isomorphisme. Il sut en eet de prendre pour Ω l\textquotesingle{}intérieur d\textquotesingle{}un voisinage stri tement}
\NormalTok{k{-}anoïde Y de x dans X et d\textquotesingle{}appliquer le lemme 2.3.9.}
\NormalTok{  Le prin ipe de Harna k est valide}
\NormalTok{Proposition 3.1.2.  Soient Ω un ouvert onnexe de X et (Ωn ) une suite d\textquotesingle{}ouverts de}
\NormalTok{X telle que haque point de Ω admette un voisinage ontenu dans presque tous les Ωn : quelle}
\NormalTok{que soit la suite roissante de fon tions hn ∈ Γ(Ωn , HX ), l\textquotesingle{}alternative est la suivante :}
\NormalTok{   soit (hn ) onverge vers +∞ uniformément sur tout ompa t de Ω,}

\NormalTok{   soit (hn ) onverge vers une fon tion h ∈ Γ(Ω, HX ), uniformément sur tout ompa t de}
\NormalTok{     Ω.}
\NormalTok{Démonstration. Puisque tout point de x admet un système fondamental de voisinages stri {-}}
\NormalTok{tement k{-}anoïdes, il dé oule du lemme 2.3.10 que l\textquotesingle{}ensemble U (resp. U ′ = Ω−U ) des points}
\NormalTok{de Ω en lesquels la suite (hn ) est majorée (resp. n\textquotesingle{}est pas majorée) est une partie ouverte ; de}
\NormalTok{plus, toujours par appli ation de e lemme, la suite (hn ) onverge}
\NormalTok{   lo alement uniformément vers ∞ sur U ′ ;}
\NormalTok{   lo alement uniformément sur U vers une fon tion h  né essairement ontinue  telle que,}
\NormalTok{     pour tout domaine stri tement k{-}anoïde Y de Ω, la restri tion de h à Y appartienne à}
\NormalTok{     H(Y ) ; vu la dénition du fais eau HX (2.3.3), la fon tion h est harmonique sur Ω.}
\NormalTok{                                                                                               2}
\NormalTok{   Notons également le résultat de onvergen e suivant :}
\NormalTok{Proposition 3.1.3.  Soient Ω un ouvert de X et (hn ) une suite de fon tions harmo{-}}
\NormalTok{niques sur Ω ; si la suite (hn ) est lo alement bornée, on peut en extraire une sous{-}suite onver{-}}
\NormalTok{geant, lo alement uniformément sur Ω, vers une fon tion harmonique.}
\NormalTok{Démonstration. Nous démontrerons au orollaire 3.3.7 que toute famille lo alement bornée}
\NormalTok{de fon tions harmoniques sur Ω est automatiquement équi ontinue en tout point de Ω ; bien}
\NormalTok{entendu, la proposition ne sera pas utilisée d\textquotesingle{}i i là. Appliquant le théorème d\textquotesingle{}As oli, on peut}
\end{Highlighting}
\end{Shaded}

\subsection{Thuillier: visually checked primary
transcription}\label{thuillier-visually-checked-primary-transcription}

\section{Thuillier primary-source
transcription}\label{thuillier-primary-source-transcription}

Controlling PDF: https://theses.hal.science/tel-00010990v1/document
SHA-256:
\nolinkurl{9c57da951e8dca2676d762e9e8c6ace302675938d72434a5b0390f15af7ea29a}. The
root viewed PDF physical pages 67 and 68, printed pages 53 and 54. The
following passages are transcribed from those rendered pages to repair
missing characters in the PDF's text extraction. The surrounding
extracted proof bodies accompany the review separately. These are
primary-source quotations, not newly proved lemmas or grants of
assumptions.

\subsection{Printed page 53, proof of Theorem 2.3.8, first
paragraph}\label{printed-page-53-proof-of-theorem-2.3.8-first-paragraph}

« Il existe une extension finie séparable \(k_1\) de \(k\) et une courbe
algébrique \(C\) propre et lisse sur \(k_1\) telles que
\(X\otimes_k k_1\) soit isomorphe à un domaine affinoïde de \(C\)
({[}31{]}, theorem 1). Appliquant le théorème de réduction semi-stable
sous la forme rigide-analytique établie par S. Bosch et W. Lütkebohmert
dans {[}9{]}, il existe une extension finie séparable \(k'\) de \(k_1\)
telle que \(C\otimes_{k_1} k'\) soit isomorphe à la fibre générique
d'une \(\operatorname{Spf}(k'^\circ)\)-courbe \(\mathcal C\) dont la
fibre spéciale est simplement semi-stable ; c'est automatiquement une
\(S\)-courbe semi-stable ». The final parenthesis and continuation
appear on printed page 54.

\subsection{Printed page 54, Lemma 2.3.9 and its
proof}\label{printed-page-54-lemma-2.3.9-and-its-proof}

« Pour tout espace strictement \(k\)-affinoïde \(Y\) rig-lisse et
purement de dimension \(1\), l'application de restriction
\(C^0(|Y|,\mathbb R)\to\operatorname{Hom}(\Gamma(Y),\mathbb R)\) induit
un isomorphisme du sous-espace \(H(Y)\) sur
\(\operatorname{Hom}(\Gamma(Y),\mathbb R)\).

Démonstration. Soit \(k'\) une extension galoisienne finie de \(k\)
telle que l'espace strictement \(k'\)-affinoïde \(Y\otimes_k k'\) soit
isomorphe à la fibre générique d'une
\(\operatorname{Spf}(k'^\circ)\)-courbe simplement semi-stable
\(\mathcal Y\). La projection canonique \(p:Y\otimes_k k'\to Y\) induit
une bijection de \(\Gamma(Y\otimes_k k')^{\operatorname{Gal}(k'/k)}\)
sur \(\Gamma(Y)\) et l'application de restriction
\(C^0(|Y|,\mathbb R)\to\operatorname{Hom}(\Gamma(Y),\mathbb R)\)
coïncide avec l'application induite par la restriction
\(C^0(|Y\otimes_k k'|,\mathbb R)\to
\operatorname{Hom}(\Gamma(Y\otimes_k k'),\mathbb R)\) entre les
sous-espaces des invariants sous \(\operatorname{Gal}(k'/k)\) ; puisque
\(H(Y)=H(Y\otimes_k k')^{\operatorname{Gal}(k'/k)}\), la conclusion
découle de la proposition 2.3.2. »

The superscript notation on the boundary is retained as printed; it is
not silently reinterpreted as a pointwise fixed-boundary assertion.

\section{Bibliography}\label{bibliography}

\begin{itemize}
\tightlist
\item
  A. Chambert-Loir and A. Ducros, \emph{Formes différentielles réelles
  et courants sur les espaces de Berkovich}, arXiv:1204.6277v2, 25 July
  2025.
\item
  W. Gubler and J. Rabinoff, \emph{Classically psh and pluriharmonic
  functions on Berkovich spaces}, arXiv:2506.13548v2, 17 September 2025.
\item
  W. Gubler, P. Jell, and J. Rabinoff, \emph{Forms on Berkovich spaces
  based on harmonic tropicalizations}, cached arXiv:2111.05741v3. Cached
  numbering, not an uninspected journal renumbering, controls
  applications here.
\item
  V. Wanner, \emph{Comparison of two notions of subharmonicity on
  non-archimedean curves}, cached arXiv:1801.04713v3; published Math. Z.
  293 (2019), 443--474.
\item
  M. van der Put, \emph{The class group of a one-dimensional affinoid
  space}, Ann. Inst. Fourier 30(4) (1980), 155--164, Theorem 1.1.
\item
  A. Thuillier, thesis, HAL tel-00010990v1, Theorem 2.3.8 and the
  explicit algebraic-domain assertion in its proof; Lemma 2.3.9 and the
  scalar harmonic functoriality used in GR's supplied proof.
\item
  V. Berkovich, \emph{Étale cohomology for non-Archimedean analytic
  spaces}, Publ. Math. IHÉS 78 (1993), Sections 1.5--1.6 and 3.1.
\item
  S. Bosch, U. Güntzer, R. Remmert, \emph{Non-Archimedean Analysis},
  Springer (1984), 3.2.4/2 and 6.1.2/1, the finite-valued-field and
  strict Nullstellensatz foundations stated exactly in the text.
\item
  A. Ducros, \emph{Variation de la dimension relative en géométrie
  analytique p-adique}, Compositio Math. 143 (2007), 1511--1532, Theorem
  4.9, in the local fibre-dimension formulation stated by CLD 11.4.5.
  The additional cached polyhedrality body is separately included with
  the corrected transport sources; no universal dimension equality is
  used.
\end{itemize}

\end{document}
